\section{NN-C001 --- Quasilinear sparse systems compress on their original variables}
\begin{proposition}[NN-C001]
Fix integers $k,a,b\geq1$. Given $m\leq n\lceil\log_2 n\rceil^a$ polynomials $f_e\in\mathbb F_2[x_1,\ldots,x_n]$, each listed with at most $\lceil\log_2 n\rceil^b$ monomials of degree at most $k$, one can deterministically construct at most $6n$ linear combinations of the $f_e$ having precisely the same common zeros in $\mathbb F_2^n$. Each combination has polylogarithmically many monomials, of degree at most $k$. Construction and exact counting take
\[
 \operatorname{poly}_{k,a,b}(n+m)\,
 2^{O_{k,a,b}(n\log\log n/\sqrt{\log n})}
\]
bit operations and polynomial space. The construction introduces no auxiliary Boolean variables. This is preservation of the Boolean zero set, not equality of ideals in the unrestricted polynomial ring.
\end{proposition}
\paragraph{Exact inputs.}
E: The Exponential Time Hypothesis is False, supplied LaTeX. Section 6, Lemmas 6.1--6.2; Appendix D Lemma D.2 (lem:tensor-sparse-query).

\begin{proof}
Handle $m=0$ and bounded $n$ directly. Put $B=\max(1,\lceil m/n\rceil)$, $\mu=\lfloor\log_2m\rfloor+1$, $P=20000B^2\mu^2$ and $D_0=320000B^2\mu^3$. Apply E, Lemma 6.1, to construct $A^{(0)}$ with at most $3n$ rows and row weight at most $D_0$, such that every nonzero vector in its kernel has weight at least $m/P$. This lemma concerns arbitrary vectors of bits; it never uses the clause origin of those bits.

For $x\in\mathbb F_2^n$, set $v(x)=(f_1(x),\ldots,f_m(x))$. Write $S=\{x:v(x)=0\}$ and let $D(A)$ be the number of $x\notin S$ with $Av(x)=0$. The count returned by a common-zero query is $N(A)=|S|+D(A)$. At a stage, enumerate all $P$-tuples of pairs $(e,c)\in[m]\times\mathbb F_2$, interpreted as a row $h$ with repeated indices combined modulo two. Every false survivor has at least $m/P$ nonzero coordinates. A uniform tuple hits such a coordinate with probability at least $1-(1-1/P)^P>1/2$. Conditional on one hit, the independent coefficient of that occurrence makes $h\cdot v(x)$ a uniform bit, even if other occurrences repeat the same coordinate. Thus its rejection probability is at least $1/4$.

Consequently $\mathbb E_hD(A;h)\leq(3/4)D(A)$. Evaluate $N(A;h)$ for every tuple and append a minimizer. Since $|S|$ is the same for all candidates, this also minimizes $D$. After $3n$ rounds the false survivor count is at most $2^n(3/4)^{3n}=(27/32)^n<1$, and hence is zero. Every original zero survives every row, proving exact equality of the Boolean zero sets.

A tested system has at most $6n$ equations, degree at most $k$, and at most $\lceil\log_2n\rceil^b\max(P,D_0)$ monomials per equation. This is the interface of Lemma D.2 for one fixed constant depending only on $k,a,b$. There are at most $3n(2m)^P+1=\exp((\log n)^{O_{a,b}(1)})$ counter calls. Initial code construction has the same kind of overhead. Every fixed polylogarithm is $o(n\log\log n/\sqrt{\log n})$, so multiplying these costs gives the stated bound. Enumerate candidates sequentially, retaining only the current matrix, candidate, best candidate and counts; Lemma D.2 then preserves polynomial space.
\end{proof}
\paragraph{Status and complexity impact.}
Proved from designated input E; self-check recorded. Complete deduction-proof review by root: passed; imported input theorems are not claimed independently proved. Extends the linear-number-of-equations core to a quasilinear number of sparse equations while keeping the original witness parameter $n$. It supplies no polynomial-time hitting set: selection itself uses the expensive exact counter and depends on the particular map $v$. It does not preserve all zeros over extension fields and does not turn a signed circuit into a monotone circuit.

\paragraph{Geometric and representation dictionary.}
In the Boolean coordinate algebra $B_n=\mathbb F_2[x]/(x_i^2-x_i)$, evaluation gives $B_n\cong\prod_{x\in\mathbb F_2^n}\mathbb F_2$. An ideal in this product is determined by the coordinates on which it vanishes. Hence the original and compressed equations do generate the same ideal after passage to $B_n$, and define the same finite reduced, finite etale zero-dimensional scheme. The resulting isomorphism and its count are computed without listing its possibly $2^n$ components. The trace dictionary is SGA7 II XXII Corollary 3.2; the relative finite trace is SGA4 XVII 6.2.3. There is no new reductive-group action or representation multiplicity attached by this compression.



\section{NN-C002 --- Fixed characteristic-two fields and exact preimage tasks}
\begin{proposition}[NN-C002]
For every fixed $r,k,a,b$, NN-C001 extends to sparse systems over $\mathbb F_{2^r}$ in $n$ field variables, with $m\leq n(\log n)^a$ equations and at most $(\log n)^b$ monomials per equation. Counts, lexicographically first witnesses and exact uniform sampling from a nonempty zero set are available in $\operatorname{poly}(n)2^{O(n\log\log n/\sqrt{\log n})}$ time and polynomial space; the sampling time is expected. The same bounds apply to the fibers of a map with those coordinate-polynomial bounds.
\end{proposition}
\paragraph{Exact inputs.}
NN-C001. full proof in this register.

E: The Exponential Time Hypothesis is False, supplied LaTeX. Appendix D Lemma D.2; Theorem 1.1 for background self-reduction.

\begin{proof}
Fix once and for all an $\mathbb F_2$-basis of $K=\mathbb F_{2^r}$ and its multiplication table. Write each variable as $x_i=\sum_{j=1}^r x_{ij}b_j$. A monomial of degree at most $k$ expands in each coordinate as a sum of at most $r^k$ products of at most $k$ bits, with coefficients determined by this fixed multiplication table. Thus each equation gives $r$ equations in $N=rn$ bits, of degree at most $k$ and polylogarithmic sparsity; the number of equations remains at most $N(\log N)^{O(1)}$. This is a bijective coordinate encoding, so it is parsimonious. Apply NN-C001 and absorb the fixed factor $r$ in the asymptotic constants.

For witness search keep all $N$ declared bits. At stage $j$, add the equations fixing the already chosen prefix and test the count after setting the next bit to zero. If that count is positive choose zero; otherwise choose one. Provided the original count is positive, at least one branch is nonempty. At most $N$ further linear equations and $N+1$ counts are required, so the same bounds hold. For a fiber $F(x)=y$, subtract the fixed coordinate $y_e$ from each polynomial; this adds at most one monomial.

For exact sampling, let $A_0,A_1$ be the two child counts at a prefix. Draw an integer uniformly from $\{0,\ldots,A_0+A_1-1\}$ by rejection from $\lceil\log_2(A_0+A_1)\rceil$ random bits. Its expected number of trials is below two. Take branch zero precisely if the integer is below $A_0$. The probability of any full satisfying assignment telescopes to $1/|S|$. The largest counts have at most $N+1$ bits and only $O(N)$ queries are needed. When the count is zero the algorithm reports that sampling is unavailable. This proves the expected-time and space claims.
\end{proof}
\paragraph{Status and complexity impact.}
Proved from NN-C001; exact field/running-time parameters stated. Complete deduction-proof review by root: passed; imported input theorems are not claimed independently proved. The field degree $r$ is fixed. Letting $r$ grow changes both the Boolean witness count and the monomial expansion $r^k$. The proof uses characteristic two essentially: field coordinates are $\mathbb F_2$-linear. It is not a same-parameter theorem for sparse equations over arbitrary odd-characteristic fields. Search and sampling do not imply that an NP language has a polynomial-time unambiguous verifier.

\paragraph{Geometric and representation dictionary.}
The field-coordinate construction is the elementary affine coordinate description of restriction of scalars for the finite separable extension $K/\mathbb F_2$, restricted to rational points. To obtain the finite counting scheme one also imposes $x_i^{2^r}-x_i=0$, or equivalently all bit relations $x_{ij}^2-x_{ij}=0$. Its coordinate algebra is a product indexed by the actual field-valued solutions. Counting is rank/trace; conditional sampling is repeated restriction to clopen subsets. No purity beyond the trivial weight-zero Frobenius representation on these rational components is required. SGA4 XVII 6.2.3 supplies the finite trace formalism; no claim of a canonical GCT representation is made.



\section{NN-C003 --- Exact rational weighted partition functions with controlled bit cost}
\begin{proposition}[NN-C003]
Fix a finite alphabet size $q$, arity $w$ and density constant $C$. Given $n$ variables with $q$ states and at most $Cn$ explicitly tabulated rational factors, each depending on at most $w$ variables, their sum-product partition function is computable exactly in
\[
 \operatorname{poly}(L+n)\,
 2^{O_{q,w,C}(n(\log\log n)^2/\sqrt{\log n})}
\]
time and polynomial space, where $L$ includes every numerator and denominator bit. Signs and zero weights are allowed. The number of factors, rather than only the number of variables, is part of the hypothesis.
\end{proposition}
\paragraph{Exact inputs.}
E: The Exponential Time Hypothesis is False, supplied LaTeX. Proposition 3.4; Lemmas 4.7--4.8; high-incidence branching in Section 5; Appendix B Lemma B.1 and its bit-height method.

\begin{proof}
Encode each $q$-state variable by $t=\lceil\log_2q\rceil$ bits and add a unary-in-the-state validity factor which is one on the $q$ chosen codes and zero on the other codes. Replace each local table by its multilinear interpolation polynomial on its at most $wt$ bits, using
\[
 \Phi(x)=\sum_{a\in\{0,1\}^{I}}\Phi(a)
 \prod_{i:a_i=1}x_i\prod_{i:a_i=0}(1-x_i).
\]
Here $t,w$ are fixed, so table sizes, polynomial degrees and scopes are bounded constants. Invalid encodings contribute zero and valid encodings correspond bijectively to assignments in the original alphabet. Let $N=tn$. The number of factors and total number of factor--variable incidences are $O(N)$, and the total product degree in independent incidence variables is $O(N)$.

If a bit belongs to more than $b_0=\max(1,\lceil M/N\rceil)\lceil\sqrt{\log_2N}\rceil$ scopes, branch on it. There are at most $M/b_0\leq N/\sqrt{\log_2N}$ such bits. Each branch has $N'=N-O(N/\sqrt{\log N})$, still $\Theta(N)$, and maximum incidence bounded by a fixed power of $\log N'$. Empty scopes are scalars. Add factors equal to one to declare any unused bits. Lemmas 4.7--4.8 supply a layout for Proposition 3.4. Its incidence degree is $O(N')$ and its coefficient algebra can be $\mathbb Q$. The resulting sum-product cost is $\operatorname{poly}(N)2^{O(N'(\log\log N')^2/\sqrt{\log N'})}$. Branching multiplies by $2^{O(N/\sqrt{\log N})}$ and hence does not change the displayed exponent.

For bit complexity, let $D_{\rm in}$ be the product of all positive table-entry denominators. Then $\log_2D_{\rm in}\leq L$, and every table coefficient has denominator dividing $D_{\rm in}$. Let $f=O(N)$ be the number of factors after validity and neutral factors are included, and let $w'$ be their fixed maximum bit arity. E Section 5, Proposition 5.3, and Appendix B use the explicit common denominator $d_{\rm den}=(R-1)!Q!$ for replacement-row coefficients. Taylor substitution introduces binomial integers but no additional denominators. Each evaluated factor therefore has denominator dividing $D_{\rm in}d_{\rm den}^{w'}$, and every group-product term and final sum can use the single denominator
\[
 \Omega=(D_{\rm in}d_{\rm den}^{w'})^f.
\]
Its bit length is at most $f(L+w'\log_2d_{\rm den})$, polynomial in $N+L$. After this common denominator is cleared, a term's numerator is a product of $O(fw')$ input and substitution coefficients, each with polynomial bit length in $N+L$. Summing the enumerated terms adds only the logarithm of their number to that bit bound, even if signs cancel.

For readout, Appendix B's Hermite matrix $\mathsf V$ has dimension $M_H=QJ=\operatorname{poly}(N)$ and entries $\binom uj t^{u-j}$, of $O(M_H\log(Q+2))$ bits. Multiplying the Taylor data by $\Omega$ makes them integral. Cramer's rule gives the requested constant term as $\det\mathsf V^{(0)}/(\Omega\det\mathsf V)$. An $h\times h$ minor with entry bit bound $B$ has bit length at most $hB+\log_2(h!)+O(1)$; this bounds all fraction-free elimination intermediates and gives polynomial bit cost. Unlike the unweighted integer application, the final answer here may remain rational, which is allowed. Each branch and its interpolation therefore use polynomial-bit arithmetic; sequential processing preserves polynomial space. The final denominator divides the product of the individual factor denominators and the output numerator also has polynomial bit length. No unit-cost assumption on large rational weights is used.
\end{proof}
\paragraph{Status and complexity impact.}
Proved extension of E Proposition 3.4; uses its weaker exponent, not an unproved extension of D.1. Complete deduction-proof review by root: passed; imported input theorems are not claimed independently proved. Gives signed or rational weighted CSP partition functions, bounded-arity spin models and bounded-density weighted model counting. Arbitrary weights of $B$ bits do not automatically fit D.1 with precision $P_0=O(n)$: a product can need $\Theta(nB)$ bits. This proof keeps that cost in the polynomial arithmetic overhead by using the rational-algebra routine and its weaker exponent. Dense instances with superlinear numbers of factors are not covered by this statement.

\paragraph{Geometric and representation dictionary.}
For the finite split scheme $X=\operatorname{Spec}(\mathbb Q^{q^n})$, the product of local factors is an element $W\in\Gamma(X,\mathcal O_X)$, and its sum over assignments is the trace of multiplication by $W$ in that finite etale algebra. In the idempotent basis of this product algebra, multiplication by $W$ is diagonal with entries $W(x)$, so its trace is exactly the displayed sum. EGA IV4, printed p.111,18.2.1--18.2.2 is the original finite-algebra trace construction; equation18.2.1.2 states base-change compatibility. SGA4 XVII 6.2.3 provides the relative-dimension-zero trace formalism. This is a weighted endomorphism trace, not automatically a point count or a pure Frobenius trace: arbitrary rational or negative weights need not be eigenvalues of a geometric Frobenius with the desired weight. The natural basis is indexed by assignments; there is no promised short basis for the whole exponential-dimensional algebra.



\section{NN-C004 --- Fixed algebraic tensor alphabets and exact quantum amplitudes}
\begin{proposition}[NN-C004]
Fix a finite set $\Omega$ of algebraic numbers and a constant maximum tensor arity. A tensor network with $s$ Boolean wires and entries from $\Omega$ has an exact deterministic contraction algorithm in $\operatorname{poly}(s+L)2^{O_\Omega(s/\sqrt{\log s})}$ time and polynomial space. Consequently a $q$-qubit, $g$-gate circuit from any fixed bounded-arity algebraic gate set has each specified computational-basis amplitude computable exactly in $\operatorname{poly}(q+g)2^{O((q+g)/\sqrt{\log(q+g)})}$ time and polynomial space. A fixed field representation and complex embedding specify the exact output.
\end{proposition}
\paragraph{Exact inputs.}
E: The Exponential Time Hypothesis is False, supplied LaTeX. Appendix D Theorem D.1, coefficient ring and precision hypotheses; its equality-tree and bipartite subdivision encoding.

J.S. Milne, Fields and Galois Theory, v5.10 (2022). Theorem5.1, pp.61--62, primitive element theorem.

\begin{proof}
Enlarge the fixed alphabet by its complex conjugates if probabilities will be requested, and let $K$ be the resulting fixed number field. By Milne Theorem5.1 choose a primitive element $\beta$ for $K$. If $a_dT^d+\cdots+a_0$ is an integer polynomial annihilating $\beta$ with $a_d\ne0$, then $\alpha=a_d\beta$ is integral: multiply the relation by $a_d^{d-1}$ to obtain a monic integer equation for $\alpha$. It still generates $K$. Its degree $d$ is constant. There is a fixed positive integer $c$ such that every $\omega\in\Omega$ can be written
\[
 \omega=c^{-1}\sum_{j=0}^{d-1}a_{\omega,j}\alpha^j,
 \qquad a_{\omega,j}\in\mathbb Z.
\]
Remove isolated tensors, multiplying their entries separately at the end. The remaining number of vertices $v$ is at most $2s$. Introduce a formal variable $V$ and replace each tensor entry by its integer polynomial $\sum_j a_{\omega,j}V^j$. The contraction is an integer polynomial $H(V)$ of degree at most $(d-1)v$. If $A=\max(1,\max_\omega\sum_j|a_{\omega,j}|)$, then
\[
 \|H\|_1\leq 2^s A^v.
\]
Thus choose $\ell>(d-1)v$ and $P_0$ with $2^{P_0-1}>2^sA^v$; both are $O_\Omega(s)$. Reduction modulo $V^\ell$ loses no coefficient, and reduction modulo $2^{P_0}$ permits unique signed-integer recovery of every coefficient.

Subdivide every edge with an equality tensor to make the graph bipartite. Its wire count changes by a factor two, its maximum degree stays bounded and it has $O(s)$ nonisolated vertices. Its local entries are computable in polynomial time in the required ring $(\mathbb Z/2^{P_0}\mathbb Z)[V]/(V^\ell)$. All hypotheses of D.1 hold after increasing its fixed constant $C$. Recover $H$ coefficientwise, evaluate $H(\alpha)/c^v$, reduce to the fixed basis of $K$, and multiply the removed scalar tensors. This last arithmetic has polynomial output height and polynomial cost in the full description length.

For the quantum specialization, insert each gate as its input/output tensor and each fixed initial or final basis state as a unary tensor. A bounded-arity circuit has $O(q+g)$ wire segments: the sum of numbers of incidences at gates is $O(g)$, and each qubit contributes its initial segment. The tensor contraction is the path-sum formula for the specified amplitude. Apply the preceding result. A probability for a fixed outcome is obtained by multiplying the exact amplitude by its complex conjugate in the fixed field, with polynomial extra work.
\end{proof}
\paragraph{Status and complexity impact.}
Proved from E D.1 with explicit precision and algebraic-height bounds. Complete deduction-proof review by root: passed; imported input theorems are not claimed independently proved. This is a subexponential algorithm in total circuit size $q+g$, not an efficient simulation in the input-qubit parameter for arbitrary polynomial-size circuits. It gives no BQP=P consequence. The fixed finite algebraic alphabet is essential to the stated exponent: variable coefficient heights or growing number-field degree require separate accounting. The result concerns numerical contraction, not a small symbolic arithmetic circuit for all tensor entries.

\paragraph{Geometric and representation dictionary.}
The contraction is composition of evaluation/coevaluation pairings in finite-dimensional $K$-vector spaces, and can also be written as the finite-algebra trace of a diagonal weighted assignment operator as in NN-C003. The diagonal matrix calculation in NN-C003 supplies that algebraic trace dictionary directly; SGA4 XVII 6.2.3 is the verified finite-trace comparison. The temporary $V$ is a coefficient-encoding device; it is not a Frobenius coordinate, deformation of a motive, or canonical-basis parameter. Fixed algebraicity ensures exact arithmetic but does not imply purity or a representation-theoretic obstruction.



\section{NN-C005 --- Parity is a top Boolean coefficient, with a hard sparse quadratic boundary}
\begin{proposition}[NN-C005]
For polynomials $f_1,\ldots,f_m\in\mathbb F_2[x_1,\ldots,x_n]$, let $N$ count their common zeros in $\mathbb F_2^n$ and let $\operatorname{ML}$ denote reduction modulo $(x_i^2-x_i)$. Then
\[
 N\bmod2=[x_1\cdots x_n]\operatorname{ML}\!\left(\prod_{i=1}^m(1+f_i)\right).
\]
If $\sum_i\deg f_i<n$, this parity is zero. Nevertheless computing this coefficient is $\oplus\mathrm P$-complete already for $m=O(n)$ quadratic equations with at most three monomials per equation. Exact counting on the same class is $\#\mathrm P$-complete under parsimonious polynomial-time reductions.
\end{proposition}
\paragraph{Exact inputs.}
Elementary Boolean-coordinate-algebra calculation. full proof below.

Katz, On a theorem of Ax, AJM 93 (1971), 485--499. Theorem 1.0 and definition 0.3.

E: The Exponential Time Hypothesis is False, supplied LaTeX. Appendix D Lemma D.2, algorithm for this hard class.

\begin{proof}
On each bit assignment, $\prod_i(1+f_i(x))$ equals one exactly at a common zero. For a squarefree monomial $x^S$, its sum over all bit assignments equals $2^{n-|S|}$, reduced modulo two. This is one for $S=[n]$ and zero otherwise. Summing the multilinear representative proves the coefficient identity. Reduction to squarefree monomials never increases degree, so the top coefficient vanishes when the displayed degree sum is below $n$.

For hardness, take an arbitrary polynomial-time Boolean witness verifier and compile it, on a fixed instance, into an AND/NOT circuit. Keep its $u$ witness-input bits and introduce one bit $z_g$ for each of its $s$ gates. Encode an AND gate by $z_g+ab=0$, a NOT gate by $z_g+a+1=0$, and acceptance by $z_{\rm out}+1=0$. Constants can be substituted directly. There are $N'=u+s$ variables and $s+1=O(N')$ equations, each of degree at most two and with at most three monomials. For every assignment to the witness bits, all gate bits are uniquely determined, in topological order. The equations have one extension if that witness accepts and none otherwise. This is a parsimonious reduction from every $\#\mathrm P$ verifier; taking parity proves the $\oplus\mathrm P$ claim. Membership follows by guessing the bit assignment and checking the equations.

This also explains why the low-degree divisibility theorem does not solve the general hard instances. Over $\mathbb F_q$, Katz's Theorem 1.0 states $q^\mu\mid N$ with $\mu=\max(0,\lceil(n-\sum_i d_i)/\max_i d_i\rceil)$ for nonconstant defining polynomials of degrees $d_i$. If $\mu=0$ it imposes no divisibility. Even when $\mu>0$, the congruence does not distinguish zero from a positive multiple of $q^\mu$. In the circuit encoding the sum of degrees generally meets or exceeds the number of variables. Adding all Boolean relations to make a finite scheme also destroys the small degree-sum hypothesis. E provides the stated subexponential algorithm for the hard class, not a polynomial coefficient-extraction algorithm.
\end{proof}
\paragraph{Status and complexity impact.}
Identity and reduction proved in full; Katz divisibility used by exact classical locator. Complete deduction-proof review by root: passed; imported input theorems are not claimed independently proved. Identifies the precise computational gate behind a mod-two cohomological count. Low individual degree and constant monomial sparsity do not make parity easy when the number of equations grows linearly. E changes the exponential running time, while $\oplus\mathrm P$-completeness remains valid.

\paragraph{Geometric and representation dictionary.}
The Boolean quotient is the coordinate algebra of a finite split etale scheme. The coefficient functional above is its sum-of-evaluations trace over $\mathbb F_2$, expressed in the squarefree monomial basis. This basis has $2^n$ elements. Its finite scheme has weight-zero Frobenius but the trace modulo two is still $\oplus\mathrm P$-complete from a succinct equation presentation. SGA7 II XXII Corollary 3.2 gives exactly the coherent trace congruence, with properness satisfied for this finite scheme. It gives no polynomial algorithm for constructing an explicit cohomology basis.



\section{NN-C006 --- The effective Frobenius-data gate and the correct affine boundary}
\begin{proposition}[NN-C006]
Suppose a deterministic polynomial-time procedure, on every sparse quadratic system in NN-C005, outputs polynomially many explicitly represented $\mathbb F_2$ matrices of polynomial total dimension whose alternating trace equals the number of common bit zeros modulo two. Then $\mathrm P=\oplus\mathrm P$ and $\mathrm{NP}=\mathrm{RP}$. If it instead outputs the exact integer rank/count in polynomial time, then $\mathrm{FP}=\#\mathrm P$ and $\mathrm{PH}=\mathrm P$. The mere existence of finite coherent or etale cohomology is insufficient. For a nonproper affine scheme, use a proper compactification and its boundary, or compactly supported etale cohomology, rather than the ordinary affine coherent $H^0$.
\end{proposition}
\paragraph{Exact inputs.}
NN-C005. parsimonious sparse-quadratic reduction.

Valiant--Vazirani, NP is as easy as detecting unique solutions, TCS 47 (1986), 85--93. Theorem 1.1; Corollaries 3.1--3.4, especially Corollary 3.2 p.91.

Toda, PP is as Hard as the Polynomial-Time Hierarchy, SIAM J. Comput. 20 (1991), 865--877. main theorem PH subset P\textasciicircum{}PP; publisher abstract verified, internal numbering not claimed.

SGA 7 II, Expose XXII. Theorem 3.1 (3.1.1); Corollary 3.2 (3.2.1).

SGA 5, Expose III. equation 6.13.2; Theorem 6.13.3.

\begin{proof}
The trace of an explicit matrix is the sum of its diagonal entries. Therefore the assumed construction followed by polynomially many field operations decides the parity problem in polynomial time. By NN-C005 every language in $\oplus\mathrm P$ reduces to that problem in polynomial time, yielding $\oplus\mathrm P\subseteq\mathrm P$; the reverse containment is immediate. Valiant--Vazirani's one-sided randomized isolation reduction maps an unsatisfiable formula only to unsatisfiable formulas and maps a satisfiable formula to a uniquely satisfiable formula with inverse-polynomial probability. A parity oracle always rejects the former and accepts the latter. Polynomial amplification gives SAT in RP. Since RP is contained in NP and SAT is NP-complete, NP=RP.

For the integer version, the parsimonious reduction makes all $\#\mathrm P$ functions polynomial-time computable; all polynomial-time nonnegative polynomial-output-length functions can also be counted by choosing a binary integer below their value, so the standard functional equality follows. A PP comparison is then polynomial-time computable. Toda's main theorem puts PH in $\mathrm P^{\mathrm{PP}}$, which now equals P. This integer conclusion is stronger than the parity conclusion and must not be inferred from a single residue.

The classical trace expressions are mathematically valid without any efficiency claim. For the finite Boolean solution scheme, $H^0(\mathcal O)$ is $\mathbb F_2^N$, all higher coherent cohomology vanishes, and Frobenius acts as the identity. That matrix can have exponential dimension. An arbitrary affine scheme does not satisfy the proper coherent formula: for example $H^0(\mathbb A^1,\mathcal O)=\mathbb F_q[x]$ is infinite-dimensional. If $U\subseteq\overline U$ is an open immersion into a proper finite-type scheme and $D=(\overline U\setminus U)_{\rm red}$, then $|U(\mathbb F_q)|=|\overline U(\mathbb F_q)|-|D(\mathbb F_q)|$. Applying SGA7 XXII Corollary 3.2 to both proper schemes gives the correct difference of coherent traces modulo $p$. Alternatively SGA5 III 6.13.3 supplies the compact-support $\mathbb F_p$ etale trace formula. Neither operation supplies polynomial-size matrices or an algorithm to construct them.
\end{proof}
\paragraph{Status and complexity impact.}
Conditional gate proved; classical trace hypotheses audited; no unconditional complexity collapse claimed. Complete deduction-proof review by root: passed; imported input theorems are not claimed independently proved. Separates purity/existence from effective presentation. A, M and a Boolean support object do not by themselves provide the hypothesized matrix construction. The parity reduction would derandomize an NP decision only to RP through Valiant--Vazirani; claiming NP=P at this point would require an additional deterministic isolation gate.

\paragraph{Geometric and representation dictionary.}
The dictionary is literal: solution set to finite etale scheme, count to rank, parity to coherent Frobenius trace, and affine restriction to proper object minus boundary. On the representation side $H^0$ is a direct sum of trivial Frobenius representations indexed by solutions. Weight zero controls eigenvalue magnitude but not the cost of extracting the number of summands from equations. SGA7 II XXII 3.1--3.2 and SGA5 III 6.13 are the exact locators; no Grothendieck archive document is invented.



\section{NN-C007 --- Exact deterministic-isolation gates and why logspace isolation does not transfer automatically}
\begin{proposition}[NN-C007]
A polynomial-time computable polynomial-length list of witness restrictions, each forbidding all witnesses on no-instances and containing an odd number of witnesses for at least one list member on each yes-instance, would imply $\mathrm{NP}\subseteq\oplus\mathrm P$. If all these parity tests were in P it would imply NP=P. A single polynomial-time restriction that always has at most one witness and preserves nonemptiness would imply NP=UP. E supplies such a single isolating restriction only at its subexponential running time. Hoza's short seed generator does not give the required polynomial-size list for arbitrary NP computations.
\end{proposition}
\paragraph{Exact inputs.}
Valiant--Vazirani, NP is as easy as detecting unique solutions, TCS 47 (1986), 85--93. Theorem 1.1; Corollaries 3.1--3.4, especially Corollary 3.2 p.91.

Hoza, NL = UL, ECCC TR26-323, supplied PDF. Theorems 1.1 and 1.3; Definition 1.2; Section 1.2.

E: The Exponential Time Hypothesis is False, supplied LaTeX. Theorem 1.1 and Appendix D Lemma D.2, conditional-count self-reduction.

\begin{proof}
Let the list have $t=t(|x|)$ members with witness counts $a_1,\ldots,a_t$. In $\mathbb F_2$ the OR of their parities is
\[
 1+\prod_{i=1}^t(1+a_i).
\]
This expression is a $\#\mathrm P$ count modulo two: a nondeterministic machine has one distinguished accepting branch, and on the other branch independently chooses, for each $i$, either one dummy option or an accepting witness for restriction $i$. It accepts exactly the valid tuples. There are $1+\prod_i(1+a_i)$ accepting paths, and the total nondeterministic witness length is polynomial. On no-instances every $a_i=0$, so the parity is zero; on yes-instances some $a_i$ is odd, so the parity is one. This proves the $\oplus\mathrm P$ inclusion. If each list parity is deterministically computable in polynomial time, ordinary OR decides the NP language in P.

For the single-restriction assertion, guess its restricted witness and check it. The global promise of at most one witness makes this an unambiguous polynomial-time verifier, and preservation of nonemptiness gives the original language. Hence every NP language lies in UP; the reverse inclusion is definitional. A list with some uniquely satisfiable member is not enough for this argument: other members can accept and there is no supplied unambiguous way to select the first good member. The parity-list argument does not cure that uniqueness problem.

E counts under added unit equations. Repeatedly count the zero branch of the next prefix bit, choosing a nonempty branch, until one satisfying assignment remains. At most $n$ count queries and $n$ unit constraints suffice; an unsatisfiable input can be sent to a fixed inconsistent restriction. The construction thus preserves nonemptiness and yields at most one witness, but its time is E's subexponential bound, not polynomial.

Hoza's Theorem 1.3 gives a seed length $O(\log(wn))$ for width-$w$, length-$n$ layered graphs, hence a polynomial-size seed list when the graph has polynomial width and length. An NP verifier can have a configuration width $2^{\operatorname{poly}(N)}$ on inputs of length $N$. For that graph the seed bound is polynomial in $N$, allowing exponentially many seeds. The min-isolation test also uses the specific Reinhardt--Allender unambiguous reachability machinery cited in Hoza Section 1.2; it is not a generic polynomial-time test that an arbitrary NP witness restriction is unique. These two parameter/interface gaps prevent the proposed automatic NP transfer.
\end{proof}
\paragraph{Status and complexity impact.}
Conditional implications and E self-reduction proved; Hoza parameter gate audited. Complete deduction-proof review by root: passed; imported input theorems are not claimed independently proved. Makes the desired NP analogues concrete and reviewable. Randomized Valiant--Vazirani plus efficient parity yields NP=RP; deterministic polynomial-list odd isolation yields NP subset parity-P; globally unique deterministic isolation yields NP=UP. These are different conclusions with different extra hypotheses.

\paragraph{Geometric and representation dictionary.}
A restriction cuts a finite Boolean solution scheme by additional equations, equivalently selects a clopen collection of its rational components. Odd isolation asks that one selected finite scheme have odd degree; unique isolation asks for degree at most one. The finite trace of SGA4 XVII 6.2.3 records these degrees. A weight function's unique minimum is a combinatorial selection on this set, not a consequence of Frobenius purity. In the representation dictionary, direct sums encode alternative restrictions and tensor products encode simultaneous witness choices; the displayed parity OR uses their mod-two rank arithmetic.



\section{NN-C008 --- Subexponential computation-transcript proofs, without a natural-system size claim}
\begin{proposition}[NN-C008]
There exists a Cook--Reckhow propositional proof system in which an unsatisfiable fixed-width $k$-CNF with $n$ variables and $m$ clauses has a refutation of length $\operatorname{poly}_k(n+m)2^{O_k(n\log\log n/\sqrt{\log n})}$; at fixed clause density the sharper exponent $O(n/\sqrt{\log n})$ applies. Every tautological Boolean formula of size $L$ has a proof in this system of length $\operatorname{poly}(L)2^{O(L/\sqrt{\log L})}$. Analogous computation certificates certify emptiness of the sparse systems of NN-C001.
\end{proposition}
\paragraph{Exact inputs.}
E: The Exponential Time Hypothesis is False, supplied LaTeX. Theorem 1.1, including polynomial space and fixed-density bound.

Cook--Reckhow proof-system definition. Polynomial-time decidable proof relation, constructed explicitly below.

\begin{proof}
Fix E's deterministic exact counting machine, including a standard binary implementation of its arithmetic, and encode each complete configuration by a polynomial number of bits in the input length. A certificate for a CNF consists of its full computation transcript ending with count zero. Check that the first configuration is the prescribed start configuration, every successive pair obeys the fixed local transition relation, and the final output is zero. Verification is polynomial in the transcript length. Soundness follows from determinism and correctness of E: an accepted transcript is its actual computation and a zero count means unsatisfiability. Completeness follows by recording that computation for an unsatisfiable formula. If the run has $T$ steps and space $S$, the transcript length is $O(TS)$ up to encoding overhead. Since $S$ is polynomial, this has the asserted size.

To obtain a proof system for all tautologies, an input proof additionally contains a formula $\varphi$ and the canonical Tseitin encoding of a circuit for $\neg\varphi$, with the output forced to one. For a size-$L$ formula this has $O(L)$ variables and $O(L)$ clauses of width at most three. Local gate constraints have a unique extension for every assignment to the original variables, so the CNF is unsatisfiable exactly when $\varphi$ is a tautology. The verifier checks this deterministic encoding and an E transcript for it. Fixed-density Theorem 1.1 gives $2^{O(L/\sqrt{\log L})}$ up to polynomial factors. On invalid proof strings the Cook--Reckhow output-function version returns a fixed tautology, so its range is exactly TAUT. Replacing E's CNF solver by NN-C001 similarly gives polynomially checkable computation certificates of sparse-system emptiness.
\end{proof}
\paragraph{Status and complexity impact.}
Proved existence of an algorithm-dependent proof system; no claim about resolution, PC, Frege or IPS. Complete deduction-proof review by root: passed; imported input theorems are not claimed independently proved. The proofs are superpolynomial in general, so this does not imply NP=coNP. The verification procedure checks a Turing computation; it supplies no bound on width, degree, monomial count or size in a pre-existing algebraic or propositional proof system.

\paragraph{Geometric and representation dictionary.}
Unsatisfiability is emptiness of the finite Boolean solution scheme. A computation transcript certifies this from its equations; it is not automatically an effective Nullstellensatz identity or a short free resolution of its coordinate algebra. In the finite etale model, emptiness means rank zero and zero $H^0$, while SGA7 II XXII 3.2 only identifies a residue of that rank. E computes the integer rank through its own algorithm, not through a claimed small cohomological presentation.



\section{NN-C009 --- Small semantic quotient certificates do not supply short polynomial-calculus proofs}
\begin{proposition}[NN-C009]
Every unsatisfiable Boolean polynomial system has a small shared-circuit identity proving $1$ belongs to its generated ideal in the Boolean function algebra. However checking that identity modulo the Boolean ideal is coNP-hard, even when the defining equations are the sparse quadratics in NN-C005. E does not supply a polynomial bound for ordinary Nullstellensatz or polynomial-calculus certificates.
\end{proposition}
\paragraph{Exact inputs.}
NN-C005. Boolean quotient and parsimonious gate encoding.

E: The Exponential Time Hypothesis is False, supplied LaTeX. Theorem 1.1 / Lemma D.2 for the algorithmic distinction.

\begin{proof}
In characteristic two put $g_i=\prod_{j<i}(1+f_j)$. A telescoping product gives the formal polynomial identity
\[
 \sum_{i=1}^m f_i g_i=1+\prod_{i=1}^m(1+f_i).
\]
If there is no common bit zero, the product on the right vanishes at every Boolean assignment and is zero in $B_n=\mathbb F_2[x]/(x_i^2-x_i)$. Thus $\sum_i f_i g_i=1$ in $B_n$. All prefix products can share intermediate gates, so the collection of the $g_i$ and the displayed sum has size linear in the total supplied circuit size plus $m$.

This compact expression is not a polynomial-time checkable proof of the crucial vanishing assertion. Given any Boolean verifier circuit, construct the equations of NN-C005 and their product $P=\prod_i(1+f_i)$. The circuit for $P$ has polynomial size. It represents the zero element of $B_n$ if and only if the original verifier has no accepting input. Therefore deterministic polynomial-time verification of these semantic quotient identities would decide UNSAT in P, and hence P=NP. The construction itself does not solve that check.

In the ordinary polynomial ring, an actual Nullstellensatz certificate requires additional polynomials $h_j$ satisfying
\[
 1=\sum_i f_i g_i+\sum_{j=1}^n h_j(x)(x_j^2-x_j).
\]
The vanishing product is indeed in the Boolean ideal because $B_n$ is the algebra of all Boolean functions, but the above small circuit expression has not constructed small $h_j$ or a small derivation of membership. Expansion or multilinear reduction can involve exponentially many monomials. A transcript of E as in NN-C008 is a valid algorithmic certificate; changing its proof system to polynomial calculus, resolution or an ideal-proof system requires an additional simulation theorem with quantitative bounds. NN-C029 subsequently supplies subexponential ordinary ideal multipliers over $\mathbb Q$ for bounded-density local verifier systems by a different construction. NN-C033 supplies the fixed-characteristic circuit transfer, and NN-C034 constructs integral constant-free certificates valid after every base change. None of these gives polynomial-size multipliers or a subexponential ordinary monomial-PC simulation; NN-C021 and NN-C034 instead exhibit an exponential ratio between monomial-PC and circuit-IPS size on the same family.
\end{proof}
\paragraph{Status and complexity impact.}
Explicit identity and verification hardness proved; later NN-C029--NN-C034 supply ordinary circuit certificates while retaining the distinction from monomial polynomial calculus. Complete deduction-proof review by root: passed; imported input theorems are not claimed independently proved. Blocks the inference that detecting or removing supported arithmetic zeros automatically produces short algebraic refutations. It identifies the missing effective ideal-membership/derivation gate, without claiming an impossibility result for every potential simulation of E.

\paragraph{Geometric and representation dictionary.}
The quotient $B_n$ is the coordinate algebra of a finite split scheme; its idempotent decomposition makes semantic ideal membership transparent after all components are listed. The telescoping identity is an ideal-theoretic statement on that finite scheme. Producing polynomial-ring correction terms is a presentation problem, analogous to replacing an abstract zero object by an explicit exact complex. Rank or Frobenius-trace identities in SGA7 XXII do not control those presentation sizes. The translation is algebraic; no unproved categorical purity theorem is used.



\section{NN-C010 --- Burgisser transfer: the exact Dedekind family and the one-way conclusion}
\begin{proposition}[NN-C010]
The characteristic-zero Cook--Valiant transfer needs RH for Dedekind zeta functions of the bounded-degree, bounded-height number fields furnishing algebraic solutions to the constant-replacement systems in its proof. These fields are generally arbitrary number fields, not necessarily abelian or normal. Conditional on M, the needed RH is available, and Burgisser's transfer gives
\[
 \mathrm{VP}_{\mathbb C}=\mathrm{VNP}_{\mathbb C}
 \ \Longrightarrow\ \#\mathrm P/\mathrm{poly}=\mathrm{FP}/\mathrm{poly},
\]
with the stronger Boolean-part and nonuniform-collapse conclusions at the cited classical locators. Its contrapositive is a Boolean nonuniform counting lower bound implying an algebraic separation. A suffices for the restricted proof route when those witness fields can be chosen abelian with the same quantitative bounds; no supplied theorem establishes that restriction in general. No converse transfer is supplied.
\end{proposition}
\paragraph{Exact inputs.}
Burgisser, Completeness and Reduction in Algebraic Complexity Theory, Springer 2000, Chapter 4. Theorem 4.5(A), Corollary 4.6(1), pp.64--65; Theorem 4.12 and Remark 4.17, pp.71--75; Theorem 4.18 and Corollary 4.19, p.76; proof of 4.5(A3), pp.77--78.

A: supplied Alpoge Lean archive. DirichletLNonvanishingRightOfHalf; DirichletLZerosOnLine; LFunction\_ne\_zero\_of\_half\_lt\_re; LFunction\_eq\_zero\_re; Remark p.4 step20.

M: Masullo claimed general Hecke GRH. Thm12.5.1.1 pp.519--520; Thm13.3.1.1 p.522; all uses conditional on M.

Prior CX register. CX-K01--CX-K03: analogous Koiran witness-field audit, used as background only.

\begin{proof}
The classical proof to be imported is Burgisser, Chapter 4, proof of Theorem 4.5(A3), with Theorem 4.12 and the root-modulo-primes estimate of Corollary 4.19. It replaces arbitrary circuit constants by variables and imposes agreement on all Boolean inputs. A bounded algebraic solution is represented in one number field $K=\mathbb Q(\alpha)$. The prime-ideal theorem used in Section 4.6 assumes RH for $\zeta_K$; degree-one prime ideals yield roots modulo rational primes. The proof does not pass to a Galois closure. The constants of one sufficiently large finite-field specialization become polynomial-length nonuniform advice. Theorem 4.5(A) and Corollary 4.6(1) state the resulting transfer and collapse. These are exact classical proof locators, rather than an asserted new proof of that theorem.

For the new premise audit, if $K/\mathbb Q$ is finite abelian, its Dedekind zeta function is a product of the corresponding primitive Dirichlet $L$-functions. A therefore gives its RH. If $K$ is arbitrary, the trivial Hecke character of $K$ has $L$-function $\zeta_K$, so M supplies precisely the missing input, conditional on M. Neither statement requires holomorphy of a nontrivial Artin $L$-function. The existence of algebraic solutions alone does not make their fields abelian: for example $X^3-2=0$ has roots generating non-Galois cubic fields; none is contained in an abelian extension of $\mathbb Q$, since every subextension of an abelian Galois extension is Galois. This example illustrates the field-class gap, not a proof that any logically possible consequence of A cannot cover it.

Finally take the contrapositive of the displayed implication. E's supplied lower bound concerns near-linear circuits for a language in $\mathrm E^{\mathrm{NP}}$. It is not a superpolynomial lower bound for the Boolean functions computing $\#\mathrm P$ values, and hence does not instantiate the antecedent. Uniform $\mathrm P\ne\mathrm{NP}$ alone also does not state a nonuniform counting lower bound. Conversely, an algebraic separation does not reverse an implication whose antecedent is algebraic equality. Thus the combination closes the analytic premise only under M (or a separately proved abelian-witness promise), leaving the Boolean lower-bound gate intact.
\end{proof}
\paragraph{Status and complexity impact.}
Classical transfer by exact locator; field-family audit proved; general characteristic-zero use conditional on M. Complete deduction-proof review by root: passed; imported input theorems are not claimed independently proved. Use function notation $\#\mathrm P\not\subseteq\mathrm{FP}/\mathrm{poly}$, or explicitly encode output bits as Boolean functions, rather than silently writing a counting class as a decision class. Under the same classical hypothesis, Corollary 4.6(1) also gives $\mathrm{NC}^3/\mathrm{poly}=\mathrm P/\mathrm{poly}=\mathrm{NP}/\mathrm{poly}=\mathrm{PH}/\mathrm{poly}$ and a uniform second-level collapse of PH if VP=VNP. These do not assert VP=VNP, and E does not choose an alternative.

\paragraph{Geometric and representation dictionary.}
The constant-replacement equations define an affine scheme over $\mathbb Z$. A closed point of its characteristic-zero fiber supplies $\operatorname{Spec}K$; reductions with residue degree one give rational points over finite fields. Its finite permutation representation has Dedekind zeta function $\zeta_K$. The relevant trace counts prime-ideal splitting, not satisfying assignments. The elementary matrix trace of multiplication supplies the finite-algebra language, and Burgisser Section 4.6 supplies the exact arithmetic estimate. A covers abelian character decompositions; the general field need not admit a decomposition into Dirichlet characters. This is the same type of analytic field gate as Koiran, but the output here is nonuniform circuit simulation.



\section{NN-C011 --- Inversion and collision attacks for sparse low-degree maps}
\begin{proposition}[NN-C011]
Fix $r,k,a,b$. Let $F_n:\mathbb F_{2^r}^n\to\mathbb F_{2^r}^{m(n)}$ have $m(n)\leq n(\log n)^a$ explicitly given coordinate polynomials of degree at most $k$ and sparsity at most $(\log n)^b$. There is a deterministic $2^{O(n\log\log n/\sqrt{\log n})}\operatorname{poly}(n)$ algorithm that inverts every image point. The same bound finds a collision whenever one exists, and counts ordered unequal colliding pairs. Thus these maps cannot have $2^{\Omega(n)}$ inversion or collision security against all algorithms in a security parameter proportional to $n$. This does not rule out ordinary polynomial-time one-wayness.
\end{proposition}
\paragraph{Exact inputs.}
NN-C001--NN-C002. quasilinear sparse-system counting and search.

E: The Exponential Time Hypothesis is False, supplied LaTeX. Appendix C Corollary C.2 is the previously known constant-locality special case, not newly claimed here.

\begin{proof}
For inversion solve the equations $F_{n,e}(x)-y_e=0$. The parameter bounds and exact preimage search are those of NN-C002. The algorithm succeeds for every $y$ in the image, so in particular its success probability is one when $y$ is sampled as $F_n(x)$ under any distribution on $x$.

For collisions introduce two input vectors $x,x'$, giving $2n$ field variables, and impose $F_{n,e}(x)-F_{n,e}(x')=0$. The coordinate degree remains at most $k$ and sparsity at most twice the original bound. After conversion to the fixed binary basis there are $2rn$ bits. Partition the unequal pairs according to the first bit position $j$ where the binary encodings differ. In branch $j$, add linear equations $x_i+x'_i=0$ for $i<j$ and $x_j+x'_j+1=0$. These at most $rn$ equations preserve the quasilinear equation bound. Every unequal pair belongs to exactly one branch. Count each branch with NN-C001 and sum to obtain the ordered unequal-pair count; search any nonempty branch for a collision. There are only $rn$ branches and replacing $n$ by $2rn$ changes the exponent by a fixed factor. If the field-output length is smaller than the field-input length, a collision exists by the pigeonhole principle; otherwise the algorithm can correctly report none.

For every fixed $c>0$, the algorithm's exponent divided by $n$ tends to zero, so its running time is eventually below $2^{cn}$. This contradicts a claimed lower bound against every algorithm at that exponential security scale, including an average-case claim requiring negligible attack success, because the inversion attack succeeds on every image. A security parameter with an unrelated scaling must be substituted explicitly. A subexponential algorithm is still superpolynomial, so the standard definition of a one-way function, which excludes only polynomial-time adversaries, is not contradicted.
\end{proof}
\paragraph{Status and complexity impact.}
Proved attack consequence for the stated representation and parameter; no blanket claim about deployed cryptosystems. Complete deduction-proof review by root: passed; imported input theorems are not claimed independently proved. Extends E's constant-locality inversion corollary to polylogarithmic monomial sparsity and a quasilinear output length, and adds a direct collision-counting construction. Dense multivariate quadratics with $\Theta(n^2)$ monomials per coordinate are outside this theorem. A generic polynomial-size Boolean circuit encoding has $n+s$ variables and can lose subexponentiality in the intended security parameter; the input representation is part of the result.

\paragraph{Geometric and representation dictionary.}
An inversion fiber is a rational-point fiber of the polynomial map $F_n$. Collisions are rational points of its fiber product with itself, with the diagonal removed; the first-differing-bit decomposition is a disjoint clopen stratification of this finite Boolean scheme. Counts are degrees of those finite etale pieces, with trace formalism SGA4 XVII 6.2.3. This construction supplies an algorithm for the succinct fibers, not a theorem that GRH or purity makes arbitrary fibers easy.



\section{NN-C012 --- PIT transfer fails at both the identity interface and the time quantifiers}
\begin{proposition}[NN-C012]
E's Boolean sparse-system algorithm does not, through the stated reductions, satisfy the Kabanets--Impagliazzo formal-PIT hypothesis. First, vanishing on the Boolean cube differs from polynomial identity. Second, $2^{O(N\log\log N/\sqrt{\log N})}$ is a mildly subexponential bound, not the stronger $\mathrm{NSUBEXP}=\bigcap_{\epsilon>0}\mathrm{NTIME}(2^{N^\epsilon})$ bound used by the cited PIT transfer. Third, general circuit encodings charge their gates in the witness variable count. A current independent unconditional arithmetic/Boolean lower-bound disjunction remains a disjunction and does not close these gates.
\end{proposition}
\paragraph{Exact inputs.}
E: The Exponential Time Hypothesis is False, supplied LaTeX. Appendix D Lemma D.2; Corollary 1.4 and Appendix C Corollary C.1.

Kabanets--Impagliazzo, Derandomizing Polynomial Identity Tests Means Proving Circuit Lower Bounds, Computational Complexity 13 (2004), 1--46. Main theorem as stated on author research page; refined statement at CGGM26 Theorem 1.1.

Chukhin--Gajulapalli--Golovnev--Mihajlin, A Little Derandomization Goes a Long Way, ECCC TR26-313, 5 Nov 2026. Theorem 1 p.5; Theorem 1.1 and definition of NSUBEXP p.3; Theorem 8.8 and proof of Theorem 1 pp.46--47.

\begin{proof}
The nonzero polynomial $x^2-x$ is zero on both Boolean inputs. Thus any procedure whose only assertion is whether a circuit or polynomial vanishes on the Boolean cube cannot serve as formal PIT. Over a finite field one can instead evaluate in a sufficiently large extension field, but then a reduction must encode that extension and every arithmetic gate; the constant-degree sparse-polynomial interface is not an automatic encoding of an arbitrary formal circuit.

The time issue is independent. Set $g(N)=N\log\log N/\sqrt{\log N}$. For every fixed $0<\epsilon<1$,
\[
 \frac{g(N)}{N^\epsilon}
 =N^{1-\epsilon}\frac{\log\log N}{\sqrt{\log N}}\longrightarrow\infty.
\]
Therefore the upper bound supplied by E cannot be replaced by $2^{N^\epsilon}$ for every $\epsilon>0$. This is a statement about what the known bound guarantees, not a lower bound against possible future faster algorithms. Even a linear reduction into E's interface would not supply the required NSUBEXP guarantee.

For circuit encodings, an $n$-input circuit with $s$ gates has the parsimonious sparse quadratic encoding in $n+s$ variables from NN-C005. Taking $s=n^d$ gives a bound exponential in approximately $n^d/\sqrt{\log n}$, not the all-polynomial-circuit-size $2^n/\operatorname{poly}(n)$ type saving required by standard stronger SAT-to-lower-bound transfers. The already supplied E circuit corollary is valid at its own exact near-linear scale; it cannot be promoted by dropping $s$.

For current-source completeness, CGGM26 Theorem 1 states the unconditional alternative $\mathrm{NE}\not\subseteq\mathrm{SIZE}[2^n/(2n)]$ or $\mathrm{VP}_{\mathbb Z}^{0}\ne\mathrm{VNP}_{\mathbb Z}^{0}$. This theorem is independent of E and A. It does not select the algebraic alternative, and its constant-free formal-degree-bounded classes must not be silently replaced by ordinary $\mathrm{VP}_{\mathbb C}$ and $\mathrm{VNP}_{\mathbb C}$. Its precise quantitative PIT gates are useful further research targets, but E has not met them here.
\end{proof}
\paragraph{Status and complexity impact.}
Explicit counterexample and asymptotic/interface audit proved; no formal independence claim. Complete deduction-proof review by root: passed; imported input theorems are not claimed independently proved. Prevents conflating ETH-style $2^{o(n)}$ savings with the stronger SUBEXP quantifier, and numerical Boolean evaluation with symbolic arithmetic identity. No PRG, hitting set or BPP=P conclusion follows from the supplied E circuit lower bound alone. CGGM26 supplies an up-to-date background disjunction, not a new consequence counted as ours.

\paragraph{Geometric and representation dictionary.}
Formal identity asks whether a regular function is zero in an affine coordinate ring. Boolean testing asks whether its restriction to the finite reduced subscheme cut out by $x_i^2-x_i$ is zero. The restriction map has a nonzero kernel, as the example shows. A trace on the finite subscheme forgets still more data. Finite-algebra matrix traces and the SGA7 XXII trace statement therefore do not identify formal PIT with Boolean zero testing. On the representation side, one must retain the polynomial family and its arithmetic circuit model; the finite evaluation representation is only a quotient.



\section{NN-C013 --- Full energy histograms and sharper Hamming-weight enumerators}
\begin{proposition}[NN-C013]
Fix $k,w,C,d$. Let $F$ be any $k$-CNF on $n$ declared variables. Suppose an energy $E(x)$ is the sum of at most $Cn$ explicitly tabulated nonnegative integer local costs, each on at most $w$ bits, with total possible range at most $M\leq n^d$. The entire polynomial
\[
 Z_F(T)=\sum_{x:F(x)=1}T^{E(x)}
\]
is computable coefficientwise in $\operatorname{poly}(n+m+M)2^{O(n(\log\log n)^2/\sqrt{\log n})}$ time and polynomial space. Hence minimum/maximum feasible energy and the numbers of optimizers are computable within that bound. For Hamming weight $E(x)=\sum_i x_i$, a separate linear-size addition-tree encoding improves the exponent to $O(n\log\log n/\sqrt{\log n})$. In particular, the full independence polynomial of an arbitrary graph has that sharper bound in its number of vertices.
\end{proposition}
\paragraph{Exact inputs.}
E: The Exponential Time Hypothesis is False, supplied LaTeX. Proposition 3.4; Lemmas 4.7--4.8, 5.1--5.2, 6.1--6.3; Appendix B; Theorem 1.1.

NN-C001. compression preserves the full solution set, not merely its cardinality.

\begin{proof}
Apply E's disjoint sparsification lemma with a rational parameter of reciprocal polylogarithmic size, for example $\sigma=1/\lceil\sqrt{\log_2n}\rceil$. For fixed $k$, the leaves have $n(\log n)^{O_k(1)}$ clauses on the original variables, their satisfying sets partition that of $F$, and their number is at most a polynomial times $2^{\sigma n}$. Compress each leaf as in NN-C001. The final equations $P_i=0$ have count $a=O(n)$, fixed degree and polylogarithmic sparsity, and their Boolean zero set is exactly that leaf's satisfying set.

Set $\ell=n+1$ and work in the finite-dimensional commutative $\mathbb Q$-algebra
\[
 \mathcal A=\mathbb Q[V,T]/(V^\ell,T^{M+1}).
\]
For each $P_i$ use the parity-lift factor from E Lemmas 5.1--5.2,
\[
 B_i(x,V)=\sum_{j<\ell}e_j(L_i(x))\,[L_i(x)(L_i(x)-1)]^jV^j,
\]
where $L_i$ is its integer monomial lift. For a local cost table $c_h$ insert the multilinear interpolation factor $W_h(x,T)$ whose value on a bit assignment is exactly $T^{c_h(x)}$. Its input-variable degree is at most $w$, regardless of the exponent of $T$. The product of parity factors has incidence degree at most $k(a+2(\ell-1))=O(n)$ after the global $V$ truncation, and the cost factors add at most $wCn=O(n)$ incidence degree. If $M_{\rm inc}$ is the total scope size, then $M_{\rm inc}=n(\log n)^{O(1)}$. Set
\[
 b_0=\max\{1,\lceil M_{\rm inc}/n\rceil\}\lceil\sqrt{\log_2n}\rceil.
\]
Branch on all variables whose incidence exceeds $b_0$. There are at most $M_{\rm inc}/b_0\leq n/\sqrt{\log_2 n}$ such variables, and the remaining incidence is polylogarithmic. Lemmas 4.7--4.8 and Proposition 3.4 then evaluate the sum of the full product in the asserted weaker-exponent time.

After evaluating on Boolean assignments, apply the coefficient-sum map in $V$ only. This map is used after multiplication and is not treated as a ring homomorphism. Every dropped term of $V$-order at least $\ell$ is divisible by $2^\ell$ at integer arguments. Consequently the resulting integer polynomial in $T$ is coefficientwise congruent modulo $2^\ell$ to the leaf energy enumerator. Each genuine coefficient is between zero and $2^n$, so the least nonnegative residues recover it exactly. The $T$ truncation loses nothing because the sum of the local maximum energies is at most $M$. The algebra dimension is $(n+1)(M+1)$, polynomial in the stated parameters. E's rational height argument applies coordinatewise: only polynomially many bits are needed, and algebra operations have polynomial overhead. Process leaves and branches sequentially, adding their recovered coefficient arrays. Disjointness proves the final polynomial is exactly $Z_F$ and keeps every final coefficient at most $2^n$.

For the sharper Hamming specialization, form a balanced tree adding the $n$ one-bit inputs. At level $j$ there are $O(n/2^j)$ additions of $O(j)$-bit integers, using $O(j)$ Boolean gates per addition. The total gate count is $O(n\sum_{j\geq1}j/2^j)=O(n)$. Equality of the final binary sum to a chosen $h\in\{0,\ldots,n\}$ uses $O(\log n)$ further gates. Its Tseitin clauses have a unique extension for every original assignment and introduce $O(n)$ variables. Conjoin these clauses with $F$, obtaining fixed width $\max(k,3)$ and $O(n)$ variables, even when $F$ is dense. E Theorem 1.1 computes the coefficient of $T^h$ with the sharper general-density exponent. Repeat for $n+1$ values of $h$. For a graph use $F=\bigwedge_{uv\in E}(\neg x_u\vee\neg x_v)$; the coefficient counts independent sets of each cardinality. The polynomial input prefactor includes the edge list.
\end{proof}
\paragraph{Status and complexity impact.}
Proved coefficient-algebra extension; Hamming specialization separately sharpened by a parsimonious linear-size circuit. Complete deduction-proof review by root: passed; imported input theorems are not claimed independently proved. Gives complete polynomial-size distributions and exact optimization, not only a single unweighted count. Polynomially bounded total energy and $O(n)$ bounded-scope cost tables are essential. Arbitrarily many cost factors or exponentially many output coefficients fall outside the bound. The independence polynomial can be evaluated at any supplied rational activity with polynomial extra bit cost; this uses its $n+1$ coefficients and does not assert a general arbitrary-weight shortcut for dense CSPs.

\paragraph{Geometric and representation dictionary.}
On the finite Boolean solution scheme, the energy partitions the rational components into clopen subsets $X_j$. The enumerator is $\sum_j\operatorname{rank}H^0(X_j,\mathcal O)T^j$, equivalently the character of the graded vector space with one basis vector per solution in its energy degree. The finite-algebra trace is immediate in the component idempotent basis; SGA4 XVII 6.2.3 supplies the verified finite-trace formalism. The variable $T$ records this chosen grading; it does not turn $T^j$ into a Tate twist or a geometric Frobenius eigenvalue without a separate realization. Thus the result produces an effective graded rank computation but no plethysm multiplicity formula or GCT obstruction.



\section{NN-C014 --- Two unconditional limits on universal hitting-set interpretations of E}
\begin{proposition}[NN-C014]
A set $H\subseteq\mathbb F_2^n$ meeting the nonempty zero set of every system of $n$ sparse linear Boolean equations must equal the entire cube. Also, a fixed $a\times m$ binary matrix with $a<m$ cannot preserve the Boolean zero set of every $m$-tuple of sparse linear polynomials by replacing $f$ with $Af$. Therefore E's small collection of parity checks is necessarily instance dependent; its theorem is not a universal subexponential assignment hitting set or a universal rank-deficient linear sketch of violations.
\end{proposition}
\paragraph{Exact inputs.}
Elementary singleton and kernel arguments. complete proof below.

E: The Exponential Time Hypothesis is False, supplied LaTeX. Section 6, Lemmas 6.1--6.2, instance-dependent greedy completion.

\begin{proof}
For each $b=(b_1,\ldots,b_n)\in\mathbb F_2^n$, take the $n$ equations $x_i+b_i=0$. Their common zero set is exactly $\{b\}$, and each equation has degree one and at most two monomials. If $H$ meets every such nonempty set then $b\in H$ for every $b$, so $H=\mathbb F_2^n$. This is an information-theoretic obstruction to the specified point-query hitting-set model; it says nothing against algorithms that manipulate the supplied equations.

For the second assertion, let $A$ have fewer rows than columns. Choose a nonzero vector $z\in\ker A$. For $n\geq1$ set $f_i(x)=z_i x_1$. Every $f_i$ has degree at most one and at most one monomial. The original system vanishes exactly on $x_1=0$, since some $z_i=1$. Yet $Af(x)=(Az)x_1=0$ identically, so the compressed system accepts the whole cube. Thus no such fixed matrix preserves zero sets for all instances. Equivalently a universal linear sketch separating every possible nonzero violation vector must have trivial kernel and at least $m$ rows.

E's Lemma 6.1 is consistent with this obstruction: its universal initial checks exclude only vectors of small support. Lemma 6.2 then selects additional checks using the particular image $\{f(x):x\in\mathbb F_2^n\}$ and exact survivor counts. The completed matrix only needs its kernel to avoid the nonzero vectors in that image. NN-C001 retains precisely this dependence and computational cost.
\end{proof}
\paragraph{Status and complexity impact.}
Proved structural limitation. Complete deduction-proof review by root: passed; imported input theorems are not claimed independently proved. Separates E's conditional-expectation construction from a black-box SAT hitting set and from formal PIT hitting sets. A uniform efficient constructor depending on the instance is a different question; no P versus NP lower bound is inferred from the two elementary impossibility statements.

\paragraph{Geometric and representation dictionary.}
Each singleton system defines one rational reduced point, a proper finite etale scheme of degree one. A universal set of evaluation points must meet every such closed subscheme of the Boolean scheme, forcing every component to occur. In its coordinate algebra $\mathbb F_2^{2^n}$, these are the primitive coordinate idempotents. The sketch counterexample changes the ideal and doubles the selected clopen locus. Finite trace and degree are as in SGA4 XVII 6.2.3; weight-zero purity does not remove the need to retain each possible component in a universal evaluation set.



\section{NN-C015 --- The constant-free model removes the analytic constant-elimination gate}
\begin{proposition}[NN-C015]
For constant-free integer arithmetic circuits with polynomial size and polynomial formal degree, the implication
\[
 \mathrm{VP}_{\mathbb Z}^{0}=\mathrm{VNP}_{\mathbb Z}^{0}
 \quad\Longrightarrow\quad \#\mathrm P\subseteq\mathrm{FP}/\mathrm{poly}
\]
holds without A or M. Thus the arbitrary-constant characteristic-zero transfer in NN-C010 and the constant-free alternative of CGGM26 have different analytic requirements. A separation for these constant-free classes does not itself establish $\mathrm{VP}_{\mathbb C}\ne\mathrm{VNP}_{\mathbb C}$. This is a classical model refinement, with a direct proof supplied here, not a claimed new theorem from E.
\end{proposition}
\paragraph{Exact inputs.}
Definitions of VP\textasciicircum{}0\_Z and VNP\textasciicircum{}0\_Z. CGGM26 p.5: constant-free circuits of polynomial size and formal degree.

Chukhin--Gajulapalli--Golovnev--Mihajlin, A Little Derandomization Goes a Long Way, ECCC TR26-313, 5 Nov 2026. Proposition 8.1 p.40, coefficient-norm bound; Theorem 1 p.5.

Burgisser, Completeness and Reduction in Algebraic Complexity Theory, Springer 2000, Chapter 4. Theorem 4.5(A), Corollary 4.6(1), pp.64--65; Theorem 4.12 and Remark 4.17, pp.71--75; Theorem 4.18 and Corollary 4.19, p.76; proof of 4.5(A3), pp.77--78.

\begin{proof}
Let $R(x,w)$ be any polynomial-time witness verifier defining a $\#\mathrm P$ function. For each input length, build a polynomial-size AND/NOT verifier circuit and introduce a Boolean variable for each gate. Over $\mathbb Z$, the gate-validity indicator for $z=ab$ is
\[
 zab+(1-z)(1-ab),
\]
and that for $z=1-a$ is $z(1-a)+(1-z)a$. Each is a constant-free bounded-degree polynomial, equal to one exactly when the corresponding Boolean gate is correct. Multiply all these indicators and the output bit. The resulting polynomial $Q(x,w,z)$ has a constant-free circuit of polynomial size and formal degree. For each $x,w$ there is exactly one consistent assignment to the gate variables, and its product is one precisely when $R(x,w)$ accepts. Hence $\sum_{w,z\in\{0,1\}^{\operatorname{poly}(n)}}Q(x,w,z)$ belongs to $\mathrm{VNP}_{\mathbb Z}^{0}$ and equals the desired count on Boolean $x$.

Assume equality of the two constant-free classes. Its counting polynomial has a constant-free arithmetic circuit with $s=\operatorname{poly}(n)$ gates and formal degree $d=\operatorname{poly}(n)$. Here constants have formal degree one, multiplication adds formal degrees, and addition or subtraction takes their maximum. Write $\|P\|_1$ for the sum of absolute coefficient values. A direct induction in topological gate order gives $\|P_g\|_1\leq2^{t d_g}$ at gate number $t$: addition increases the larger norm by at most a factor two, while multiplication multiplies the two bounds and adds formal degrees. The leaves $0,1,-1$ and variables satisfy the base case. Thus every used gate has coefficient norm at most $2^{sd}$. At Boolean inputs its integer value has at most $sd+O(1)$ bits. Ordinary integer addition and schoolbook multiplication evaluate the whole circuit in polynomial time. The circuit description is polynomial-length advice, proving the functional containment.

The formal-degree restriction matters for this characteristic-zero height argument: repeatedly squaring the constant two produces $2^{2^s}$ from $O(s)$ gates despite ordinary polynomial degree zero. Its formal degree grows exponentially and is excluded from that argument. NN-C020 gives a different modular simulation of bounded integer outputs which does not require this restriction. Finally, constant-free polynomial-formal-degree circuits form a more restrictive model than arithmetic circuits with arbitrary complex constants. Hardness against the former does not imply hardness against the latter. The proof above uses explicit integer arithmetic and never chooses a prime or a residue-field specialization, so it has no hidden GRH assumption.
\end{proof}
\paragraph{Status and complexity impact.}
Classical model refinement proved directly. Complete deduction-proof review by root: passed; imported input theorems are not claimed independently proved. Explains the restricted transfer underlying the current unconditional win--win. It does not convert E's near-linear $\mathrm E^{\mathrm{NP}}$ lower bound into the required nonuniform counting lower bound. It also does not supply a converse from ordinary algebraic separation to Boolean separation.

\paragraph{Geometric and representation dictionary.}
The verifier product gives a finite Boolean scheme together with its acceptance idempotent; summing the idempotent is the finite-algebra trace. Over $\mathbb Z$, the supplied integer presentation and formal-degree bound make all intermediate coefficients effectively bounded. The trace dictionary follows directly by using the component idempotent basis, as in NN-C003; compare SGA4 XVII 6.2.3. By contrast NN-C010 starts with arbitrary complex constants, whose replacement by finite-field data uses number-field points and Dedekind zeta estimates. These are different effective presentations of related rank/count constructions, not different purity statements.



\section{NN-C016 --- A genuine symbolic transfer for bounded-arity partition polynomials}
\begin{proposition}[NN-C016]
Fix arity $w$ and factor density $C$. For a factor hypergraph with $n$ Boolean variables and $m\leq Cn$ factors of scope at most $w$, attach an independent formal weight $y_{f,a}$ to every local assignment $a$ of factor $f$. Its full partition polynomial
\[
 Z_H(y)=\sum_{x\in\{0,1\}^n}\prod_f y_{f,x|I_f}
\]
has a deterministically constructible arithmetic circuit over $\mathbb Q$ of size
\[
 \operatorname{poly}(n)\,2^{O_{w,C}(n(\log\log n)^2/\sqrt{\log n})}.
\]
The rational constants have polynomial bit length. In particular the symbolic perfect-matching polynomial of a graph of bounded maximum degree has subexponential arithmetic circuit size in its vertex/edge parameter. This is a restricted symbolic consequence of E's full common routine; it is stronger than what follows merely by regarding E as a numerical SAT counter.
\end{proposition}
\paragraph{Exact inputs.}
E: The Exponential Time Hypothesis is False, supplied LaTeX. Proposition 3.4, especially algebra-operation count and rational Hermite interpolation; Lemmas 4.7--4.8; Section 3.1 product interface and measurement-separation proof.

\begin{proof}
Use the coefficient algebra $\mathcal A=\mathbb Q[y_{f,a}]$. For each factor take its Boolean multilinear interpolation
\[
 \Phi_f(x_{I_f})=\sum_{a\in\{0,1\}^{I_f}}y_{f,a}
       \prod_{i:a_i=1}x_i\prod_{i:a_i=0}(1-x_i).
\]
It is linear in the formal weights and has degree at most $w$ in its Boolean inputs. The sum-product is exactly $Z_H(y)$, and its independent-incidence input degree is at most $wm=O(n)$. The number of incidences is also $O(n)$. Branch on variables of incidence above $\max(1,\lceil M_{\rm inc}/n\rceil)\lceil\sqrt{\log_2n}\rceil$, as in NN-C003, and retain the original scopes even if a substitution makes some coefficient zero. There are at most $n/\sqrt{\log_2n}$ branched bits and every resulting layout satisfies Lemmas 4.7--4.8.

Proposition 3.4 explicitly allows every commutative $\mathbb Q$-algebra, counts polynomially many operations in that algebra per enumerated vector, and permits final interpolation by the inverse of a rational Hermite matrix. Implement an algebra element as a pointer to a gate in a shared arithmetic directed acyclic graph. Each addition, subtraction or multiplication in $\mathcal A$ creates one arithmetic gate; truncated Taylor arrays have polynomial length and require polynomially many such gates per operation. Factor evaluation takes $O(2^w)$ times a polynomial number of algebra operations because $w$ is fixed. For interpolation compute the inverse matrix over $\mathbb Q$ once and use its entries as rational constants, followed by linear combinations of the symbolic Taylor data. Do not run fraction-free elimination with symbolic pivot entries: no division by an element involving $y$ is needed.

All control choices are independent of symbolic coefficients. The layout construction uses only the scopes. The measurement-separation scans select representatives by their integer measurements and fixed input bits, not by comparing values in $\mathcal A$. Branching fixes declared Boolean variables and does not ask whether a symbolic factor is zero. Therefore the complete computation can be unrolled into a circuit without a symbolic identity oracle. Proposition 3.4 and Lemma 4.8 bound the number of operations by the displayed exponential times a polynomial, including all high-incidence branches. Layout and Hermite constants have polynomial bit length by E's row-coefficient and Appendix B bounds. Circuit construction time is polynomial in its output size; storing algebra elements as gate references avoids expansion into monomials. The output polynomial has degree $m$ in the weights, even though it may have exponentially many monomials.

For a bounded-degree graph, put one bit $x_e$ on each edge. At every vertex insert the fixed-arity indicator that exactly one incident edge is selected. On each edge insert the unary factor $(1-x_e)+y_e x_e$. The product is zero off perfect matchings and equals $\prod_{e\in M}y_e$ on a perfect matching. With maximum degree fixed, the number of bits, factors and incidences is $O(v+e)$; isolated vertices immediately give the zero polynomial when they must be matched. Thus this is a specialization of the partition construction, preserving circuit size. In the bipartite case it is the permanent polynomial with the graph's forbidden entries set to zero.
\end{proof}
\paragraph{Status and complexity impact.}
Proved symbolic extension from E Proposition 3.4. Complete deduction-proof review by root: passed; imported input theorems are not claimed independently proved. This corrects an overly broad counting-only objection in this restricted setting. It does not give a polynomial-size circuit or VP=VNP. A dense $r\times r$ permanent encoded by edge bits has $\Theta(r^2)$ Boolean variables, so substituting that size into the result gives no useful subexponential bound in $r$ (and is worse than the usual $2^{O(r)}$ upper bound). Compression through parity lifts does not automatically preserve this characteristic-zero symbolic argument: coefficientwise reduction modulo a power of two is not an arithmetic-circuit operation over $\mathbb Q$. This original construction uses rational constants. NN-C033 subsequently proves a modified construction over every prime field, and NN-C034 gives integral constant-free circuits of polynomial formal degree over Z itself. These are strengthened constructions; the original rational circuit is not simply reduced modulo a denominator prime.

\paragraph{Geometric and representation dictionary.}
The partition polynomial is the trace of multiplication by $\prod_f\Phi_f$ on the finite free algebra $\mathbb Q[y][x]/(x_i^2-x_i)$ over the parameter scheme $\operatorname{Spec}\mathbb Q[y]$. In the orthogonal idempotent basis indexed by bit assignments, this multiplication matrix is diagonal with the corresponding factor products on the diagonal. Its trace is their sum. Applying a coefficient-ring homomorphism entrywise commutes with summing the diagonal, proving the claimed base-change compatibility directly. The original EGA IV4, printed p.111,18.2.1--18.2.2 records this trace, with base change in equation18.2.1.2; compare the relative-dimension-zero formalism in SGA4 XVII 6.2.3. Specializing formal weights commutes with the trace and with the constructed arithmetic circuit. This is an actual family-level symbolic computation, unlike a single numerical rank. It still does not provide orbit-closure multiplicity obstructions, border lower bounds, or a finite Weil realization of nonstandard canonical bases.



\section{NN-C017 --- Sparse equations over every fixed finite field by prime-adic total-order truncation}
\begin{proposition}[NN-C017]
Fix a finite field $K=\mathbb F_{p^r}$ and constants $k,a,b$. Given $n$ variables over $K$ and at most $n(\log n)^a$ polynomial equations of degree at most $k$, with at most $(\log n)^b$ listed monomials per equation, their exact number of common zeros is computable deterministically in
\[
 \operatorname{poly}(L+n)\,
 2^{O_{K,k,a,b}(n(\log\log n)^2/\sqrt{\log n})}
\]
bit operations and polynomial space. Here $L$ is the explicit input length. Witness search and exact uniform sampling from a nonempty zero set have the same bound, with expected time for sampling. This extends NN-C002 to odd characteristic at the weaker Proposition 3.4 exponent; it does not assert the sharper Appendix D bound there.
\end{proposition}
\paragraph{Exact inputs.}
E: The Exponential Time Hypothesis is False, supplied LaTeX. Proposition 3.4; Lemmas 4.7--4.8; Lemmas 5.1--5.2; Proposition 5.3 and Appendix B; complete proofs of Lemmas 6.1--6.2.

NN-C003. explicit common-denominator and Hermite-minor bit-cost method.

NN-C002. prefix self-reduction and exact rejection sampling.

\begin{proof}
We first count common zeros of $m\leq CN$ sparse bounded-degree polynomials in $N$ Boolean variables with coefficients in the fixed prime field $\mathbb F_p$. Choose integer representatives of their coefficients and let $L_i(x)$ be the resulting integer polynomials. On Boolean inputs they reduce to the prescribed polynomials modulo $p$. Set $Y_i=L_i^{p-1}$. Fermat's elementary identity gives $Y_i(x)\equiv0$ or $1\pmod p$, with the first case exactly when $L_i(x)\equiv0\pmod p$. For completeness, for $u\not\equiv0\pmod p$, multiplication by $u$ permutes the $p-1$ nonzero residues; comparing their products and cancelling $(p-1)!$ gives $u^{p-1}=1$ in $\mathbb F_p$.

Set $\ell=N+1$ and use exactly E Lemma 5.1's integer polynomial
\[
 p_\ell(Y)=(1-Y)^\ell\sum_{j=0}^{\ell-1}\binom{\ell+j-1}{j}Y^j
           =\sum_{j=0}^{\ell-1}e_j(Y)h(Y)^j,
 \qquad h(Y)=Y(Y-1),\quad\deg e_j\leq1.
\]
The proof of that lemma establishes the polynomial congruences $p_\ell\equiv1\pmod{Y^\ell}$ and $p_\ell\equiv0\pmod{(1-Y)^\ell}$, independently of characteristic. Therefore $p_\ell(Y_i(x))$ is the zero indicator modulo $p^\ell$, and $p\mid h(Y_i(x))$. In $\mathcal A=\mathbb Q[V]/(V^\ell)$ use
\[
 B_i(x,V)=\sum_{j<\ell}e_j(Y_i(x))h(Y_i(x))^jV^j.
\]
After multiplying all factors, sum the retained $V$ coefficients. Each discarded term has total $h$-order at least $\ell$ and is divisible by $p^\ell$ at every Boolean input. Thus the sum-product recovers the number of simultaneous zeros modulo $p^\ell$. Since $p^\ell>2^N$, its least nonnegative residue is the exact count. This readout is performed after multiplication, not as a ring homomorphism.

If the original degree is at most $d$, then $\deg Y_i\leq d(p-1)$. Even with independent input copies for every factor, every retained coefficient has degree at most
\[
 d(p-1)\bigl(m+2(\ell-1)\bigr)=O_{p,d,C}(N).
\]
Scopes remain polylogarithmic and total incidence is $N(\log N)^{O(1)}$. Add neutral unary factors for unused declared bits. If total incidence is $M$, branch on bits appearing in more than
$\max(1,\lceil M/N\rceil)\lceil\sqrt{\log_2N}\rceil$ scopes. At most $N/\sqrt{\log_2N}$ bits are branched; every branch retains $\Theta(N)$ declared bits and has polylogarithmic maximum incidence. E Lemmas 4.7--4.8 and Proposition 3.4 now apply. Factor evaluation computes the sparse $L_i$, its fixed power $p-1$, and successively the $h$ powers in the truncated algebra; it uses polynomially many rational operations and never enumerates a factor's entire scope.

The bit bound is explicit. With E's row denominator $d_{\rm den}=(R-1)!Q!$, a substituted $Y_i$ has denominator dividing $\Delta=d_{\rm den}^{d(p-1)}$. Each $B_i$ has denominator dividing $\Delta^{2\ell-1}$ and their product has denominator dividing $\Omega=\Delta^{m(2\ell-1)}$. Coefficient bits are polynomial by E Lemma 5.1 and Appendix B; the additional fixed power $p-1$ changes only constants. Clearing $\Omega$ bounds every Taylor datum by a polynomial-bit integer: a product has polynomially many polynomial-bit factors and summing the enumerated terms adds only the logarithm of their number. The integer Hermite matrix and its fraction-free minor bound are exactly as in NN-C003. Final integer reduction is modulo $p^\ell$, which has $O_p(N)$ bits. The time is the displayed weaker exponent and sequential evaluation gives polynomial space.

Next extend the equation count from $CN$ to $N(\log N)^{O(1)}$. E Lemma 6.1's graph, its shifts, its real-valued Walsh mixing estimate, and its shortening do not depend on the coefficient field. Keep its parameters
\[
 B=\max(1,\lceil m/N\rceil),\quad \mu=\lfloor\log_2m\rfloor+1,
 \quad P=20000B^2\mu^2,\quad D_0=320000B^2\mu^3.
\]
Use the same $d_G,q_G,c_G,\delta$ and replace only the binary local matrix by a uniformly random $c_G\times d_G$ matrix over $\mathbb F_p$. A fixed nonzero vector belongs to its kernel with probability $p^{-c_G}$. If $h=\lceil\delta d_G\rceil-1\geq1$, the number of nonzero vectors of weight at most $h$ is at most
\[
 \sum_{j=1}^h\binom{d_G}{j}(p-1)^j
 \leq(h+1)(p d_G)^h
 \leq(p d_G)^{2h}
 <p^{d_G/(25B)}.
\]
Indeed $d_G\leq m<2^\mu$, so $\log_p(p d_G)\leq1+\mu\leq2\mu$, while $h<d_G/(100B\mu)$. Since $c_G\geq d_G/(8B)$, the union bound is below one. If $h=0$ there is nothing to exclude. Exhaustive enumeration of local matrices and the low-weight vectors costs $\exp(\operatorname{poly}(B,\log(m+2)))$ for fixed $p$. The rest of Lemma 6.1 uses only supports and graph mixing, so it produces at most $3N$ rows, each of support at most $D_0$, whose nonzero kernel vectors have support at least $m/P$.

For the remaining greedy stage, use lists in $([m]\times\mathbb F_p)^P$. A list hits the support of a false survivor with probability at least $1-e^{-1}>1/2$. Conditional on a hit, one uniform coefficient makes the tested sum uniform in $\mathbb F_p$, even with repeated indices. Its rejection probability is consequently at least $(1-e^{-1})(1-1/p)>1/4$. A candidate minimizing the exact count leaves at most $3/4$ of the false survivors, exactly as in E Lemma 6.2. There are at most $2^N$ Boolean assignments, so $3N$ stages remove all false survivors. Every candidate has at most $6N$ equations and polylogarithmic monomial sparsity. The number of core calls is at most $3N(pm)^P=\exp((\log N)^{O(1)})$, absorbed by the claimed running time. All enumeration is sequential and uses polynomial space. This compression preserves the common zero set on the Boolean cube.

Finally encode the original $\mathbb F_{p^r}$ variables. Fix an $\mathbb F_p$ basis; each original variable has $r$ coordinates and each equation gives $r$ coordinate equations. A degree-$k$ monomial expands into at most a fixed multiple of $r^k$ coordinate monomials. Encode each prime-field coordinate by $t=\lceil\log_2p\rceil$ bits, interpreted as the integer $\sum_{j<t}2^jx_j$. Add one polynomial equation per coordinate saying that this bit code is below $p$: explicitly, take one minus the sum of the bit-pattern indicator monomials for $0,\ldots,p-1$. On bits it is zero exactly for valid codes. It has constant degree and at most $2^t$ expanded pattern terms times a fixed constant. There are $N=rtn=\Theta_K(n)$ bits, $rm+rn=N(\log N)^{O(1)}$ equations, fixed degree and polylogarithmic monomial sparsity. Valid codes are in bijection with the original field assignments. The Boolean prime-field counter therefore gives exactly the desired number. Prefix fixing adds only $O(N)$ linear equations, so NN-C002's search and exact sampling proofs apply unchanged.
\end{proof}
\paragraph{Status and complexity impact.}
Complete new deduction proved; complete self proof-check passed. Complete deduction-proof review by root: passed; imported input theorems are not claimed independently proved. This removes an actual algorithmic restriction: sparse systems over fixed odd-characteristic fields now have a deterministic mildly subexponential exact counter with the original variable parameter, even when their scopes are polylogarithmic. The modulus and field extension degree must remain fixed; their powers and coordinate expansion enter the constants and cannot be hidden if they grow. The stronger characteristic-two Appendix D exponent has not been transferred. Neither polynomial-time counting nor a uniform identity test follows.

\paragraph{Geometric and representation dictionary.}
After adjoining $x_i^{|K|}-x_i$, the original zero set is a finite reduced split scheme over $K$ whose rank is the count. The basis expansion and valid-code equations give a bijection of rational components with a finite Boolean scheme over $\mathbb F_p$; they are an effective coordinate presentation, not an identification of unrestricted affine zero schemes. The $p$-adic order in $h(Y)$ is a divisibility filtration on integer-valued indicators. The formal nilpotent $V$ records this filtration and is not a Frobenius coordinate or a Tate twist. On the resulting finite scheme, the component-idempotent basis makes counting the trace of the identity; SGA7 II XXII Corollary 3.2 supplies its mod-$p$ coherent trace expression and SGA4 XVII 6.2.3 the finite-trace formalism. The algorithm computes this succinctly presented rank without constructing the exponential basis and supplies no new representation-theoretic positivity statement.



\section{NN-C018 --- Sparse polynomial equations over every fixed finite commutative ring}
\begin{proposition}[NN-C018]
Let $R$ be any fixed finite commutative ring. Fix $k,a,b$. Exact counting, witness search, and exact uniform sampling for systems of at most $n(\log n)^a$ degree-at-most-$k$ polynomials over $R$, with at most $(\log n)^b$ monomials per equation, have the NN-C017 time bound
\[
 \operatorname{poly}(L+n)\,
 2^{O_{R,k,a,b}(n(\log\log n)^2/\sqrt{\log n})}
\]
and polynomial space. Sampling has expected running time. In a fixed $p$-primary ring, the same bound computes a weighted sum over the solutions when there are $O(n)$ bounded-arity factors whose rational table entries lie in a fixed finite alphabet. For a general $R$, this weighted assertion holds when the factors separate as products over its primary components. Arbitrary input-height rational weights and weights coupling distinct primary components are not included in this weighted assertion.
\end{proposition}
\paragraph{Exact inputs.}
NN-C017. complete prime-field Boolean-system counter and p-ary equation compression.

J.S. Milne, Group Theory, v4.01 (2025). Theorems1.54 and1.57(c), pp.25--27, finite abelian groups as products of prime-power cyclic groups.

E: The Exponential Time Hypothesis is False, supplied LaTeX. Proposition3.4; Lemmas4.7--4.8 and5.1--5.2; AppendixB.

\begin{proof}
The zero ring is trivial, so assume $R$ has at least two elements. Write its characteristic as $M=\prod_p p^{a_p}$. The elementary integer Chinese remainder construction gives integers $e_p$ which are one modulo $p^{a_p}$ and zero modulo every other prime-power divisor of $M$. Their images in $R$ are orthogonal central idempotents with sum one. Hence
\[
 R\longrightarrow\prod_p R_p,\qquad R_p=e_pR,\qquad
 x\longmapsto(e_px)_p
\]
is a ring isomorphism: its inverse adds the components, and the identities just stated check both compositions and multiplication. Every equation therefore holds exactly when all its component equations hold. Variables in distinct components are independent. Thus the solution set is the Cartesian product of the component solution sets, and the exact counts multiply. There are only a fixed number of components. It suffices to handle a fixed $p$-primary ring.

By the stated finite-abelian-group theorem, fix additive generators $b_1,\ldots,b_d$ of this ring with unique coordinates
\[
 x=\sum_{j=1}^d z_jb_j,
 \qquad0\leq z_j<p^{a_j},\qquad
 (R,+)\cong\bigoplus_{j=1}^d\mathbb Z/p^{a_j}\mathbb Z.
\]
The multiplication table $b_ib_j=\sum_t c_{ijt}b_t$, with fixed integer representatives for the $c_{ijt}$, is part of this fixed presentation. Expanding one degree-$k$ monomial in these coordinates introduces only a constant depending on $R,k$ many integer monomials, all of degree at most $k$. Constant coefficients in $R$ are handled by the same fixed multiplication table. Consequently a ring equation is equivalent to $d$ congruences
\[
 L_t(z)\equiv0\pmod{p^{a_t}},
\]
where $L_t$ is an integer polynomial of fixed degree and polylogarithmic monomial sparsity. Intermediate reductions modulo the coordinate moduli are unnecessary: the additive relations $p^{a_t}b_t=0$ ensure that the final coordinate congruence gives the correct element.

Encode each coordinate by $\lceil\log_2(p^{a_j})\rceil$ bits and impose its fixed-size validity table, as in NN-C017. This gives $N=\Theta_R(n)$ Boolean variables. Every coordinate polynomial after substitution still has fixed degree and polylogarithmic sparsity. We next replace divisibility by $p^{a_t}$ by a fixed number of genuinely sparse $\mathbb F_p$ polynomial equations, taking care not to invert a factorial divisible by $p$.

For every integer $u$ and $j\geq0$, define the integer-valued polynomial
\[
 D_j(u)=\binom{u}{p^j}.
\]
For fixed $a$ and $j<a$, this is periodic modulo $p$ with period $p^a$. Indeed the formal power series identity
$(1+T)^{u+p^a}=(1+T)^u(1+T)^{p^a}$ holds also for negative integer $u$, with integral binomial coefficients. Modulo $p$, the last factor is $1+T^{p^a}$, by repeated Frobenius. Thus coefficients below degree $p^a$ do not change. For $0\leq u<p^a$, write $u=\sum_{i<a}u_ip^i$ with $0\leq u_i<p$. Then
\[
 (1+T)^u\equiv\prod_{i<a}(1+T^{p^i})^{u_i}\pmod p.
\]
The coefficient of $T^{p^j}$ is $u_j$: the maximum total exponent available from the lower factors is $p^j-1$, and any higher factor has exponent greater than $p^j$. By periodicity, for every integer $u$ the residues $D_0(u),\ldots,D_{a-1}(u)$ are exactly the base-$p$ digits of $u\bmod p^a$. In particular,
\[
 p^a\mid u\quad\Longleftrightarrow\quad
 D_j(u)\equiv0\pmod p\quad(0\leq j<a).
\]

Apply this to each $L_t(x)$. As a rational polynomial,
$D_j(L_t)=\prod_{h=0}^{p^j-1}(L_t-h)/(p^j)!$ has fixed degree and polylogarithmic monomial sparsity: $p^j$ is a constant depending only on $R$. Expand it over $\mathbb Q$, replace every positive bit power $x_i^h$ by $x_i$, and collect like monomials. The resulting multilinear polynomial has \emph{integer} coefficients. To prove this, for any rational polynomial $Q$ taking integer values on the Boolean cube, its unique multilinear representative has coefficient
\[
 [x^S]\operatorname{ML}(Q)
   =\sum_{T\subseteq S}(-1)^{|S|-|T|}\,Q(\mathbf1_T),
\]
where coordinates outside $T$ are zero. This is the inversion of
$Q(\mathbf1_T)=\sum_{S\subseteq T}[x^S]\operatorname{ML}(Q)$ and is an integer. Here $Q=D_j(L_t)$ is integer-valued because $L_t$ has integer Boolean values. Computationally one simply expands and collects the sparse polynomial; the coefficient formula proves integrality and does not require enumerating all assignments on the polylogarithmic scope. All rational denominators before collection divide the fixed factorial, and coefficient bit lengths are polynomial in the input length. Reduce the resulting integer coefficients modulo $p$ only after this collection.

The final Boolean system has $N(\log N)^{O(1)}$ equations, fixed degree, and polylogarithmic monomial sparsity. The number of coordinate and digit equations per original equation is fixed. It includes the constant-size validity equations. NN-C017 applies and is parsimonious with respect to valid coordinate encodings. The primary-component counts multiply as above; their outputs have $O_R(n)$ bits. Self-reduction and sampling work componentwise, and independent uniform component samples give the uniform distribution on their Cartesian product. This proves the unweighted claims.

For the weighted $p$-primary assertion, first compress the constraint system by the \emph{unweighted} NN-C017 counter; the selected equations preserve its exact Boolean zero set. Let the fixed rational weight alphabet have a common positive denominator $D$, and let $A\geq1$ bound the absolute values of all the integer weights obtained by multiplying by $D$. If there are $f\leq Cn$ weight factors, the numerator of the weighted sum after this scaling has absolute value at most $2^N A^f$. Choose $\ell=O_{R,C,A}(N)$ so that
\[
 p^\ell>2^{N+1}A^f.
\]
Express each scaled table by its integer multilinear interpolation polynomial on its fixed number of coordinate bits. Fill invalid codes by zero; the validity constraints already exclude them. Multiply these weight factors by the constraint factors in the same algebra $\mathbb Q[V]/(V^\ell)$ used in NN-C017. At Boolean points every omitted order-$\ell$ term remains divisible by $p^\ell$ after multiplication by the integer weight. The total incidence degree is still $O(N)$: the compressed constraint contribution is $O(m+\ell)$ and the $O(n)$ weight factors each have bounded degree and scope. The added scopes preserve polylogarithmic total incidence. The same layout, coefficient-height, interpolation, and space proof applies. The residue has a unique representative in $(-p^\ell/2,p^\ell/2)$ within the stated magnitude bound; this is the desired signed numerator. Divide by $D^f$ to recover the rational weighted sum.

If the ring has several primary components and each weight factor is a product of factors in the individual components, the full weighted summand and the solution domain separate, so the total partition function is the product of the computed component partition functions. Without that hypothesis, the product-of-counts argument does not apply. Likewise a variable $B$-bit weight alphabet can require $\ell=\Theta(n+nB)$, so this proof cannot simply hide arbitrary weight height in a polynomial overhead while retaining the same exponent.
\end{proof}
\paragraph{Status and complexity impact.}
Complete deduction proved; complete self proof-check passed. Complete deduction-proof review by root: passed; imported input theorems are not claimed independently proved. Nilpotents and fixed composite characteristic do not by themselves obstruct the sparse-equation algorithm: the digit reduction keeps the number of variables, equation degree and monomial sparsity within the required fixed-parameter interface. This is a theorem for each fixed ring, including fixed prime-power residue rings, not an algorithm uniform in a growing modulus such as an RSA modulus. The weighted restrictions are substantive, and the full rational bounded-arity theorem NN-C003 remains a different interface.

\paragraph{Geometric and representation dictionary.}
A polynomial solution over a finite nonreduced ring is an $R$-point of the affine equation scheme, not in general a geometric point of a reduced scheme. The primary decomposition is a decomposition of $\operatorname{Spec}R$ into clopen components; its functor of points turns the equation solution set into the displayed Cartesian product. Additive coordinates and binomial digits then present this finite set by a finite reduced Boolean scheme over $\mathbb F_p$. That presentation preserves the set and its rank, not the nilpotent structure, tangent spaces or deformation theory of the original equation scheme. Counting is the trace of the identity on the product algebra indexed by these points, and weighted counting is the diagonal multiplication trace; these assertions follow directly in the idempotent basis, with SGA4 XVII6.2.3 as the finite-trace dictionary. SGA7 II XXII Corollary3.2 applies to the resulting finite scheme for its mod-$p$ trace, not automatically to the original nonreduced ring-valued functor. No canonical-basis or orbit-closure positivity is inferred.



\section{NN-C019 --- E constructs one polynomial-height isolating weight vector for each instance}
\begin{proposition}[NN-C019]
Let $\mathcal F\subseteq\{0,1\}^n$ be a witness set whose prefix-restricted counts admit one of the exact algorithms NN-C001, NN-C002 or NN-C017. If $\mathcal F$ is nonempty, a deterministic algorithm in the same time bound up to a polynomial factor constructs a vector
\[
 w\in\{1,n+1\}^n\subseteq\{1,\ldots,2n\}^n
\]
and a unique minimum-weight witness $x^*\in\mathcal F$, with minimum weight at most $n$. It also identifies that witness. For fixed-width CNF, E Theorem1.1 gives $\operatorname{poly}(n+m)2^{O_k(n\log\log n/\sqrt{\log n})}$ time and polynomial space. For a fixed-density CNF, the construction retains the sharper $2^{O_{k,C}(n/\sqrt{\log n})}$ exponent.
\end{proposition}
\paragraph{Exact inputs.}
NN-C002. complete exact prefix-count self-reduction.

NN-C017. fixed-field extension and prefix constraints.

E: The Exponential Time Hypothesis is False, supplied LaTeX. Theorem1.1, including fixed-density statement.

Mulmuley--Vazirani--Vazirani, Matching is as Easy as Matrix Inversion, STOC1987. Section3, Lemma1, printed p.347; compare Lemma2, p.348 for the unique-lowest-term valuation argument.

\begin{proof}
Use exact prefix counts, with the bit order $0<1$, to find the lexicographically first witness $x^*$. Let $S=\{i:x_i^*=1\}$. No proper subset of $S$ is the support of a witness: its bit string would agree until its first removed one, where it has zero and $x^*$ has one, contradicting lexicographic minimality. Define $w_i=1$ for $i\in S$ and $w_i=n+1$ for $i\notin S$.

The weight of $x^*$ is $|S|\leq n$. If another witness has support $T\subseteq S$, then $T\ne S$ would contradict the preceding inclusion minimality; equality gives the same bit string. Every other witness therefore has an element outside $S$ and has weight at least $n+1>|S|$. Thus $x^*$ is the unique minimum. The argument also covers $S=\varnothing$. If the initial count is zero, the algorithm reports emptiness instead of claiming an isolated witness.

There are at most $n+1$ counting calls. In the sparse-system interfaces each prefix adds at most $n$ linear equations, with fixed degree and one or two monomials, so the corresponding bounds are preserved. For fixed-width CNF, a prefix consists of at most $n$ unit clauses; padding a unit clause by repeated literals keeps it within the width bound, or one substitutes its bit directly. The number of variables never increases, and an initial bound $m\leq Cn$ becomes at most $(C+1)n$ clauses if all variables are retained. Thus E's fixed-density exponent $2^{O_{k,C}(n/\sqrt{\log n})}$ is preserved as well. All counts, the current prefix, and the final weights use polynomial space.

MVV Lemma1 states that independent uniform weights in $[2n]$ isolate any nonempty family with probability at least $1/2$. The construction here selects a weight in that same range deterministically by first solving the instance at E's running time. It is not a polynomial-time replacement for the MVV random choice, and it is not one small weight family that works uniformly for all succinctly represented witness sets. It also does not eliminate signed cancellation at an unweighted specialization: isolating a term in a signed expression produces a nonzero leading term of a new weight-generating polynomial, whereas the original specialization can still vanish.
\end{proof}
\paragraph{Status and complexity impact.}
Complete deduction proved; complete self proof-check passed. Complete deduction-proof review by root: passed; imported input theorems are not claimed independently proved. This makes the available deterministic NP-style isolation precise: instance-dependent isolation and witness recovery are mildly subexponential for the stated sparse interfaces, with polynomially bounded weights. A polynomial-size output does not imply polynomial construction time. It supplies neither $\mathrm{NP}=\mathrm{UP}$ nor $\mathrm{NP}\subseteq\oplus\mathrm P$ with polynomial uniformity. Hoza's logarithmic-seed layered-graph construction remains a different polynomial-time uniform theorem; NN-C007 states the missing NP transfer gates.

\paragraph{Geometric and representation dictionary.}
The finite witness scheme has a diagonal grading operator with one basis vector of degree $w\cdot x$ for every witness. Its character is $\sum_{x\in\mathcal F}T^{w\cdot x}$, and the least-degree graded piece has rank one. Equivalently, $w$ exposes a single vertex of the convex hull of the witness bit vectors and selects one monomial in the initial form of the positive generating polynomial. The count/character interpretation follows directly from the component-idempotent basis and the finite-trace formalism of SGA4 XVII6.2.3. These are finite-set and torus-grading statements; they do not construct a GCT canonical basis or make a signed structure-constant formula positive. A numerical signed specialization forgets the grading and can reintroduce cancellation.



\section{NN-C020 --- Unconditional Boolean simulation from a small explicit coefficient algebra}
\begin{proposition}[NN-C020]
Let $C_n$ be a division-free arithmetic circuit of polynomial size over $B_n\otimes_{\mathbb Z}\mathbb Q$, where $B_n$ is a commutative finite free $\mathbb Z$-algebra supplied by a basis $e_0=1,e_1,\ldots,e_{r(n)-1}$ of polynomial rank and polynomial-bit integer multiplication constants. Suppose every circuit constant is supplied by polynomial-bit rational coordinates in this basis, and the total circuit and algebra description has polynomial size. Assume that on every Boolean input the output is $g_n(x)e_0$, where $g_n(x)$ is an integer in $[0,2^{w(n)})$ for a polynomial $w$. Then $g\in\mathrm{FP}/\mathrm{poly}$, without GRH and without a formal-degree bound. If these descriptions and the output bound are constructible in polynomial time, then $g\in\mathrm{FP}$.

In particular, the conclusion applies when constants lie in a number field of polynomial degree supplied by an integral primitive element with polynomial-bit monic minimal polynomial and polynomial-bit rational power-basis coordinates. A collapse giving such circuits for all $\#\mathrm P$ counting families implies $\#\mathrm P\subseteq\mathrm{FP}/\mathrm{poly}$ unconditionally.
\end{proposition}
\paragraph{Exact inputs.}
Elementary quotient-algebra simulation. complete proof below.

NN-C015. constant-free counting-family construction; coefficient-height proof refined here.

Burgisser, Completeness and Reduction in Algebraic Complexity Theory, Springer 2000, Chapter 4. Theorem 4.5(A), Corollary 4.6(1), pp.64--65; Theorem 4.12 and Remark 4.17, pp.71--75; Theorem 4.18 and Corollary 4.19, p.76; proof of 4.5(A3), pp.77--78.

EGA IV, fourth part, Publications Mathematiques IHES32 (1967). Printed p.111,18.2.1--18.2.2; base-change equation18.2.1.2; primary PDF page110 (zero-based109).

\begin{proof}
Write each constant as $d_g^{-1}\sum_j a_{g,j}e_j$, with $d_g$ a positive integer. Put $D=\prod_g d_g$. This product has polynomial bit length, because the total constant description does. Every circuit value lies in the localization $B_n[D^{-1}]$. Choose the integer modulus
\[
 M=D\,2^{w(n)+1}+1.
\]
Its bit length is polynomial, it exceeds $2^{w(n)}$, and it is coprime to every $d_g$. In fact, an explicit inverse of $d_g$ modulo $M$ is
\[
 -\frac{D}{d_g}\,2^{w(n)+1}\pmod M,
\]
because multiplying this expression by $d_g$ gives $1-M$.

Therefore the canonical map $B_n\to B_n\otimes_{\mathbb Z}\mathbb Z/M\mathbb Z$ extends to a ring homomorphism
\[
 B_n[D^{-1}]\longrightarrow
 B_n\otimes_{\mathbb Z}\mathbb Z/M\mathbb Z.
\]
Evaluate the full circuit in this quotient algebra. Each stored value has $r(n)$ coordinates modulo $M$. Addition is coordinatewise and multiplication uses the supplied table; a naive multiplication needs only $O(r(n)^3)$ modular scalar operations. All integers have $O(\log M)$ bits. Consequently the whole evaluation uses polynomially many bit operations and polynomial space, independent of the magnitude of unreduced intermediate characteristic-zero values. Exponentially large formal degrees cause no difficulty: repeated squaring remains a sequence of the originally supplied polynomially many modular gates.

By the homomorphism property, the final vector is $(g_n(x)\bmod M,0,\ldots,0)$. The algebra remains free with the supplied basis after base change, so this coordinate statement is literal, not merely equality in an unspecified residue field. Since $0\leq g_n(x)<2^{w(n)}<M$, the least nonnegative representative of the first coordinate is the exact desired integer. This proves the simulation. Polynomial descriptions serve as advice in the nonuniform case; in the uniform case they are constructed before evaluation in polynomial time.

For the number-field specialization, let $f_n\in\mathbb Z[T]$ be the supplied monic minimal polynomial of the integral primitive element $\alpha_n$. Then $B_n=\mathbb Z[\alpha_n]\cong\mathbb Z[T]/(f_n)$ is free in the power basis of size $\deg f_n$. Reducing products of two basis powers by the monic relation constructs an integral multiplication table of polynomial bit length: at most $O(\deg f_n)$ successive reductions are needed, and each can increase the coefficient bit bound by at most a polynomial in the degree and the bit height of $f_n$. The given rational constant coordinates lie in $B_n[D^{-1}]$. The general argument applies even if $f_n$ becomes reducible or inseparable modulo $M$; neither a field quotient nor a root modulo $M$ is required.

Applying the simulation to the VNP counting polynomials constructed in NN-C015 proves its $\#\mathrm P$ consequence directly for this larger effective coefficient model. In particular, constant-free circuits use $B_n=\mathbb Z$ and $D=1$. Their formal-degree restriction is part of the definition of the usual constant-free Valiant classes and was needed for NN-C015's characteristic-zero height proof, but it is not needed for this modular Boolean simulation of promised bounded integer outputs. The proof also applies to explicitly supplied rational constants of polynomial total bit height.

This does not replace Burgisser's general complex-constant theorem. An arbitrary polynomial-size circuit over $\mathbb C$ need not come with a polynomial-rank finite algebra, polynomial-bit constants, or an effectively small algebraic replacement. Burgisser's constant-replacement system constructs bounded algebraic witness fields of possibly much larger degree, and its small residue-degree-one prime argument addresses that different presentation problem. The present proof uses the entire quotient algebra in polynomial rank and thus has no corresponding prime-splitting requirement.
\end{proof}
\paragraph{Status and complexity impact.}
Complete model refinement proved; complete self proof-check passed. Complete deduction-proof review by root: passed; imported input theorems are not claimed independently proved. The analytic transfer gate is now more precisely located: polynomial rank and an explicit polynomial-bit coefficient presentation suffice, with composite-modulus arithmetic and no RH input. The contrapositive transfers Boolean counting lower bounds to this restricted but broad arithmetic model unconditionally. It does not transfer ordinary $\mathrm P\ne\mathrm{NP}$ to $\mathrm{VP}_{\mathbb C}\ne\mathrm{VNP}_{\mathbb C}$, and does not give a converse to Burgisser. NN-C043 subsequently removes this theorem's global-rank restriction for explicitly presented rational radicals of any one fixed prime degree by a constant-rank local specialization. Arbitrary complex constants, absent effective presentations, general high-degree coefficient fields, and unrestricted division gates remain outside the combined simulation results.

\paragraph{Geometric and representation dictionary.}
The coefficient object is a finite locally free affine scheme $\operatorname{Spec}B_n\to\operatorname{Spec}\mathbb Z$ with an explicit polynomial-rank trivialization of its pushforward algebra. The algorithm performs base change to $\operatorname{Spec}(\mathbb Z/M\mathbb Z)$ and evaluates in the full fiber algebra, including all of its components and nilpotents. It does not seek a rational point of that fiber. The scalar-output condition makes the resulting section descend to the base ring, where the magnitude bound identifies its unique integer lift. EGA IV4, printed p.111,18.2.1--18.2.2 gives the finite-algebra trace dictionary and base-change formula; the computational proof itself is the explicit multiplication table above. A weight or purity assertion without this small presentation does not supply the simulation, and no GCT orbit-closure obstruction is constructed.



\section{NN-C021 --- E computation transcripts have no polynomial-size translation to general resolution or polynomial calculus}
\begin{proposition}[NN-C021]
For the computation-transcript proof system of NN-C008, there is no fixed polynomial upper bound converting the size of every transcript refutation into the size of a resolution refutation of the same CNF. The analogous statement holds for ordinary polynomial calculus over $\mathbb F_2$, when proof size counts explicitly written nonzero monomials. More concretely, there are fixed-density 3-CNF contradictions on $N$ variables whose E transcripts have size
\[
 2^{O(N/\sqrt{\log N})},
\]
while every refutation in the relevant natural system has size $2^{\Omega(N)}$. The resolution and polynomial-calculus lower-bound families may be different. These natural-system lower bounds are classical inputs; the stated comparison with E's new transcript system is the deduction.
\end{proposition}
\paragraph{Exact inputs.}
E: The Exponential Time Hypothesis is False, supplied LaTeX. Theorem1.1, bounded-clause-density exponent.

NN-C008. complete computation-transcript proof-system construction.

Buss--Grigoriev--Impagliazzo--Pitassi, Linear Gaps between Degrees for the Polynomial Calculus Modulo Distinct Primes, JCSS62(2001),267--289. Author-hosted journal-final manuscript: Section2 pp.5--6; Lemma3 pp.7--8; Theorem15 p.17; Corollaries21--22 and proof p.19.

Ben-Sasson--Wigderson, Short Proofs Are Narrow: Resolution Made Simple, JACM 48 (2001), 149--169. Lemma4.2 p.155; Theorem4.4 and Corollary4.5 p.156; size-width Corollary3.6 p.155; proof of Theorem4.4 Section6.1 p.161.

Impagliazzo--Pudlak--Sgall, Lower Bounds for the Polynomial Calculus and the Groebner Basis Algorithm, Computational Complexity8(1999),127--144. Author-hosted accepted manuscript: Section6, Theorem6.2 and complete proof p.13; Corollary6.3 pp.13--14.

\begin{proof}
For resolution, take 3-regular connected expanding graphs with an odd total vertex charge and the associated Tseitin CNFs. BSW Lemma4.2 gives at most four clauses per graph vertex and one variable per edge, hence $O(N)$ clauses of width three. Their Corollary4.5 gives resolution size $2^{\Omega(N)}$; equivalently, combine their linear width lower bound Theorem4.4 with Corollary3.6. E Theorem1.1 nevertheless counts each such fixed-density 3-CNF in $2^{O(N/\sqrt{\log N})}$ time and polynomial space. Its count is zero. Recording all configurations as in NN-C008 multiplies the time by only a polynomial, absorbed by the same exponent. This gives the asserted transcript upper bound on this family.

For polynomial calculus over $\mathbb F_2$, use BGIP's Boolean mod-$p$ Tseitin family with $p=3$, over a fixed-degree expanding graph. Their Corollary21 states a linear degree lower bound for this Boolean encoding over characteristic two, and Corollary22's proof explicitly states that its 3-CNF translation has $O(N)$ clauses and variables and exponential monomial-size lower bounds. There is, however, a directional omission in the printed proof of Corollary21: its appeal to Lemma2 supplies a reduction from the Boolean system to the root-of-unity system, whereas the definition and Lemma1 require the reverse reduction to transfer the cited hardness. We supply that reverse reduction here instead of using that inference as printed.

Work in the fixed field $F=\mathbb F_4$, choose a primitive cube root $\omega$, and use BGIP's root-of-unity system $\operatorname{TS}_n(3)$ in edge variables $y_e$. Define the Boolean indicator substitutions
\[
 x_{e,i}=\delta_i(y_e),\qquad
 \delta_i(T)=\frac13\sum_{j=0}^{2}\omega^{-ij}T^j
 \quad(0\leq i<3).
\]
The scalar $3$ is invertible in characteristic two. On $T=\omega^a$ this polynomial is one for $a=i$ and zero otherwise. Every substituted Boolean axiom and clause of $\operatorname{BTS}_{n,3}$ involves only a fixed number $s$ of edge variables and has fixed degree. Include the domain equations $y_e^3-1$ for these variables and the corresponding local root-of-unity vertex or reverse-edge constraints $g_1,\ldots,g_t$. These local equations force exactly the local Boolean constraint in question.

Here is a complete constant-degree derivation argument. The algebra
\[
 F[y_1,\ldots,y_s]/(y_1^3-1,\ldots,y_s^3-1)
 \cong\prod_{a\in\{1,\omega,\omega^2\}^{s}}F
\]
has distinct roots in every coordinate. Let $q$ be the substituted Boolean constraint polynomial. At a tuple where all $g_j$ vanish, $q$ vanishes. At every other tuple, choose the first nonzero $g_j$ and set the value of $h_j$ to $q/g_j$, setting the other $h$ values to zero. Lagrange interpolation gives representatives $h_j$ of individual degree at most two, hence total degree at most $2s$, and
\[
 q=\sum_j h_jg_j+\sum_{e=1}^{s}u_e(y)(y_e^3-1).
\]
The last equality is an equality of ordinary polynomials: reduce the difference successively by the monic equations $y_e^3-1$. The remainder has individual degrees below three and vanishes at every root tuple, so it is zero. Since $s$, all input degrees and all interpolation degrees are fixed, the degrees of every displayed product are bounded by a constant $C$. Expand each multiplier into its finitely many monomials; each monomial multiple of an axiom is derivable by successive single-variable multiplication, and their linear combination gives $q$ in PC degree at most $C$. This proves a $(2,C)$ reduction from $\operatorname{TS}_n(3)$ to the Boolean system in the exact sense of BGIP's definition and Lemma1. Substitution of a degree-$d$ Boolean PC refutation therefore gives a root-of-unity refutation of degree at most $\max(C,2d)$.

BGIP Theorem15 bounds that latter degree below by $\Omega(n)$. A refutation over $\mathbb F_2$ is also one over $F$, so the Boolean system requires degree $\Omega(n)$ over $\mathbb F_2$. Its numbers of variables and clauses are $O(n)$ and its clause width is fixed. BGIP Lemma3 and its proof turn this into a 3-CNF with $O(n)$ variables and clauses while preserving the degree lower bound up to a constant factor. Impagliazzo--Pudlak--Sgall Theorem6.2 says that a proof with $M$ nonzero monomials gives a proof of degree at most $\max(d_0,2\lceil\sqrt{N\log M}\rceil+1)$, where $d_0$ is the initial degree. For constant $d_0$ and minimum refutation degree $d=\Omega(N)$, this implies $M=2^{\Omega(d^2/N)}=2^{\Omega(N)}$. This is also the exact size-degree chain invoked in BGIP Corollary22's proof. Apply E Theorem1.1 and NN-C008 to this 3-CNF family exactly as before.

Now suppose a fixed polynomial $P$ bounded the length of a natural-system refutation in terms of the length $L$ of every E transcript for the same formula. On one of the corresponding hard families,
\[
 P(L)\leq2^{O(N/\sqrt{\log N})}=2^{o(N)},
\]
contradicting the classical $2^{\Omega(N)}$ lower bound. This excludes even a universal polynomial size bound, and therefore in particular excludes a polynomial-time proof translation of that kind. The argument compares proofs of the same formula under its stated encoding; it does not forbid changing the target proof system by extension axioms or new inference rules.

For the familiar parity Tseitin equations over $\mathbb F_2$ themselves, summing all vertex equations already yields the contradiction in low degree. Thus the polynomial-calculus example deliberately uses mod three constraints over characteristic two, as in BGIP; the resolution example alone cannot justify a characteristic-two polynomial-calculus lower bound. This field distinction is essential.

No lower bound on unrestricted satisfiability algorithms follows from these proof lower bounds. BGIP Corollary22's algorithmic conclusion concerns implementations that explicitly represent polynomials as sums of monomials and generate the corresponding polynomial-calculus derivation. E's interpolation and global count computation need not produce such a derivation. NN-C009's compact semantic identity likewise does not constitute a short polynomial-calculus refutation.
\end{proof}
\paragraph{Status and complexity impact.}
Complete deduction and reverse low-degree reduction proved; primary source clauses and printed directional omission checked. Complete deduction-proof review by root: passed; imported input theorems are not claimed independently proved. This gives a concrete obstruction at the proof-system interface left open by NN-C009. E removes the exponential-time hypothesis but does not remove exponential resolution or characteristic-two polynomial-calculus proof-size lower bounds. The transcript system is still not polynomially bounded on all tautologies, so this comparison does not imply $\mathrm{NP}=\mathrm{coNP}$ or separate those classes.

\paragraph{Geometric and representation dictionary.}
For every hard contradiction, the Boolean-coordinate equation scheme is empty; its finite algebra has rank and Frobenius trace zero. E computes this vanishing efficiently relative to the former exhaustive exponent. Resolution width and polynomial-calculus degree/monomial size concern the cost of deriving the unit ideal from the chosen generators in specified syntactic systems. Thus the same zero scheme can have a short computation transcript and a long ideal-derivation certificate. The finite trace dictionary is SGA7 II XXII Corollary3.2, but that identity places no bound on the chosen presentation's derivation complexity. This is an effective distinction between a geometric zero object and an explicit proof of emptiness in a fixed calculus, not a new failure of purity or a representation-theoretic multiplicity statement.



\section{NN-C022 --- Local-map inversion and image distinguishing retain the original seed parameter}
\begin{proposition}[NN-C022]
Fix $k$. Let $F:\{0,1\}^n\to\{0,1\}^m$ be specified by a list of output scopes of size at most $k$ and their Boolean truth tables. For every target $y$, the number $|F^{-1}(y)|$, image membership, an inverse when one exists, and exact uniform sampling from a nonempty fiber are computable in
\[
 \operatorname{poly}_k(n+m)\,
 2^{O_k(n\log\log n/\sqrt{\log n})}
\]
time and polynomial space; sampling has this expected-time bound. If $m\leq Cn$, replace the exponent by $O_{k,C}(n/\sqrt{\log n})$. The number of output coordinates may be polynomially larger than $n$ in the first bound.

When $m>n$, these algorithms give a deterministic distinguisher for $F(U_n)$ against $U_m$, with advantage at least $1-2^{n-m}$ and the same time bound. Thus fixed-locality maps with polynomial output length cannot have $2^{\Omega(n)}$ inversion security or pseudorandom-generator security. This conclusion concerns the original seed length and does not assume expansion, random scopes, injectivity, or regular fibers.
\end{proposition}
\paragraph{Exact inputs.}
E: The Exponential Time Hypothesis is False, supplied LaTeX. Theorem1.1, arbitrary-density and bounded-density counting bounds.

NN-C002. prefix-count self-reduction and exact conditional sampling proof.

NN-C001. quasilinear fixed-degree polylog-sparse systems, for the logarithmic-locality extension.

Goldreich, Candidate One-Way Functions Based on Expander Graphs, 2011 revision of ECCC TR00-090. Section2, equation(1), printedp.78; Section2.2, equation(2), printedp.79; fixed versus logarithmic locality explicitly distinguished in Section2.

\begin{proof}
For each output coordinate $i$, enumerate the at most $2^k$ assignments $a$ to its scope $I_i$ for which $F_i(a)\ne y_i$. Exclude each one by the single clause
\[
 C_{i,a}=\bigvee_{j\in I_i}
 \begin{cases}x_j,&a_j=0,\\ \neg x_j,&a_j=1.\end{cases}
\]
This clause is false precisely on the forbidden local assignment. The conjunction $\Phi_y$ has exactly $n$ original variables, at most $2^k m$ clauses, and width at most $k$. Constant false coordinates give an immediate empty fiber; constant true constraints may be omitted. E applies with fixed width $\max(2,k)$. The satisfying assignments of $\Phi_y$ are exactly $F^{-1}(y)$, with no auxiliary gate variables or multiplicities. This proves the displayed counting bounds.

To find an inverse, keep the same $n$ declared variables and repeatedly count after fixing the next bit to zero, choosing that branch when its count is positive. There are at most $n$ extra unit clauses and $n+1$ counting calls. The bounded-density hypothesis remains bounded density when $m=O(n)$. Exact sampling uses the two child counts as probabilities; binary rejection sampling of a uniform integer below their sum has expected fewer than two trials. The proof in NN-C002 applies verbatim, including polynomial space. A nonempty fiber is inverted with certainty, for every target in the image; the argument is not merely an average-case attack.

Set $D_F(y)=1$ exactly when the fiber count is positive. Then
\[
 \Pr[D_F(F(U_n))=1]=1,\qquad
 \Pr[D_F(U_m)=1]=\frac{|\operatorname{im}F|}{2^m}
 \leq2^{n-m}.
\]
Subtracting gives the claimed advantage. In particular one bit of stretch already gives advantage at least one half, and larger stretch improves it. For polynomially bounded $m$, the polynomial description/read-time factor is $2^{o(n)}$, so both algorithms run in $2^{o(n)}$. If the map family is uniformly constructible, the attack is uniform; otherwise its given description is part of the input, or polynomial advice for that family. This refutes an exponential-in-$n$ security claim for this model, while allowing the possibility of hardness against all polynomial-time attacks.

Goldreich's equation(1) has exactly this form when the local predicate has a fixed number of arguments, regardless of the chosen expanding graph. His alternative equation(2) publishes the predicate table as part of the output; conditioning on that published table likewise leaves only the $n$ seed bits unknown. The paper also permits locality growing like $\log n$. An arbitrary truth table of that arity produces growing-width clauses, for which E's constants are not uniform in the width; the fixed-locality result alone does not cover that regime.

There is a useful algebraic extension of that boundary. Suppose the output predicates, although they may depend on $O(\log n)$ coordinates, have multilinear algebraic normal form over $\mathbb F_2$ of a fixed degree $d$, and suppose $m\leq n(\log n)^a$. Each such form has at most $\sum_{j=0}^{d}\binom{O(\log n)}j=(\log n)^{O_d(1)}$ monomials. Subtracting $y_i$ preserves this bound, so NN-C001 applies on the same $n$ original bits. The same inverse and image distinguisher therefore have $2^{O(n\log\log n/\sqrt{\log n})}$ time. Computing the algebraic normal forms from the explicit $2^{O(\log n)}$ truth tables takes polynomial time by Boolean finite differences. No assertion is made for a logarithmic-locality predicate of growing algebraic degree.

More generally, suppose some output coordinates are fixed-locality tables while all remaining coordinates are supplied by a Boolean circuit containing $s$ gates. Give a variable only to each of these $s$ gates and impose its constant-width Tseitin constraints, together with the requested output bits. The extension is unique for every seed. All local coordinates are still encoded without auxiliary variables. The resulting width-$\max(k,3)$ formula has $N=n+s$ variables and $O_k(m+s)$ clauses, hence costs
\[
 \operatorname{poly}_k(n+m+s)\,
 2^{O_k(N\log\log N/\sqrt{\log N})},
\]
with the sharper exponent $O(N/\sqrt{\log N})$ when $m+s=O(N)$. Thus it is the size of the genuinely nonlocal circuit part that enlarges the witness parameter, not the total number of local output constraints. All search and distinguishing conclusions persist with this explicit parameter substitution.
\end{proof}
\paragraph{Status and complexity impact.}
Complete deduction proved from E and NN-C001/002; self-check recorded. Complete deduction-proof review by root: passed; imported input theorems are not claimed independently proved. This narrows the cryptographic gate left by a generic Cook--Tseitin reduction: arbitrarily many polynomially listed local observations do not enlarge the unknown seed space. In particular, expansion can obstruct specific local or sequential inversion heuristics while failing to establish exponential security against E's global counting method. The conclusion does not refute the existence of ordinary one-way functions or pseudorandom generators, does not estimate E's constants at a practical key size, and does not make a claim for arbitrary dense multivariate-quadratic maps or logarithmic-locality random predicates of growing degree. The image test uses the original nonnegative integer fiber count; parity alone cannot replace it when a fiber has positive even cardinality.

\paragraph{Geometric and representation dictionary.}
Let $X_n=\operatorname{Spec}\mathbb Z[x_1,\ldots,x_n]/(x_j^2-x_j)$, the disjoint union of $2^n$ copies of $\operatorname{Spec}\mathbb Z$. Each Boolean output table has a unique integer multilinear interpolation which is an idempotent in this algebra, so $F$ defines a morphism $X_n\to X_m$. Over a Boolean target section $y$, its fiber is the finite split scheme with coordinate algebra $\mathbb Z^{|F^{-1}(y)|}$. The local clauses give a concise presentation of this same fiber. Its finite-locally-free rank, or trace of the identity, is the inverse count; reduction to $\mathbb F_2$ gives only the parity of that rank. E computes the rank without expanding the exponential product algebra. This is the explicit multiplication-matrix trace of EGA IV4,18.2.1--18.2.2, and the finite trace formalism of SGA4 XVII6.2.3. The image distinguisher compares support in the finite target with its uniform counting measure; the image has at most the source's number of components. No zeta-function estimate, cohomological purity theorem, or orbit-closure representation is needed, and A or M does not alter this parameter calculation.



\section{NN-C023 --- Parity hardness persists when the coherent Frobenius matrix has only one entry}
\begin{proposition}[NN-C023]
There is a deterministic polynomial-time construction taking a 3-CNF $\Phi$ to a nonzero homogeneous polynomial $F$ of degree $r$ in exactly $r\geq3$ variables over $\mathbb F_2$, given by an arithmetic formula of size $O(n+m)$, such that for the proper hypersurface
\[
 X_\Phi=V(F)\subset\mathbb P^{r-1}_{\mathbb F_2}
\]
one has
\[
 H^0(X_\Phi,\mathcal O)=\mathbb F_2,\qquad
 H^{r-2}(X_\Phi,\mathcal O)=\mathbb F_2,
\]
with all other coherent cohomology zero, and absolute Frobenius acts on the latter one-dimensional space by the scalar
\[
 \#\operatorname{SAT}(\Phi)\pmod2.
\]
Here $r=O(n+m)$, and the basis of the one-dimensional space has an explicit universal description. Consequently computing this single Frobenius entry is $\oplus\mathrm P$-complete on the specified factored class. A polynomial-time algorithm for that entry would imply $\mathrm P=\oplus\mathrm P$ and $\mathrm{NP}=\mathrm{RP}$. E computes it in subexponential time on this class. The hypersurfaces are allowed to be singular and reducible; no smoothness or purity assertion is part of the theorem.
\end{proposition}
\paragraph{Exact inputs.}
NN-C005. parsimonious Boolean gate encoding and parity-counting completeness.

NN-C006. effective Frobenius-data gate being sharpened.

Katz, Algebraic Solutions of Differential Equations (p-Curvature and the Hodge Filtration), Inventiones Mathematicae18(1972),1--118. Equations2.3.7.12--2.3.7.13 and Algorithm2.3.7.14, p.35; rank-one Special Case2.3.7.17, p.36.

SGA 7 II, Expose XXII. Theorem 3.1 (3.1.1); Corollary 3.2 (3.2.1).

E: The Exponential Time Hypothesis is False, supplied LaTeX. Theorem1.1.

Valiant--Vazirani, NP is as easy as detecting unique solutions, TCS 47 (1986), 85--93. Theorem 1.1; Corollaries 3.1--3.4, especially Corollary 3.2 p.91.

\begin{proof}
First suppose $\Phi$ has no empty clause and at least one variable. Regard each clause as a Boolean relation on its distinct variable scope. Add a tautological unary relation to every unused variable; this does not change the satisfying assignments and ensures every variable occurs. Let $\mathcal I$ be the set of factor--variable incidences and put $M=|\mathcal I|$. Then $n\leq M\leq3m+n$ after these additions. For each incidence $e=(f,i)$ introduce two distinct variables $z_{e,0},z_{e,1}$. Define integer polynomials
\[
 P_f(z)=\sum_{a\in\{0,1\}^{I_f}:\,a\text{ satisfies }f}
              \prod_{i\in I_f}z_{(f,i),a_i},
 \qquad
 Q_i(z)=\prod_{f:\,i\in I_f}z_{(f,i),0}
             +\prod_{f:\,i\in I_f}z_{(f,i),1}.
\]
Every $P_f$ is homogeneous of degree $|I_f|$ with at most eight distinct monomials. Every $Q_i$ is homogeneous of its positive incidence degree and has two distinct monomials. Put
\[
 P=\prod_f P_f\prod_i Q_i,\qquad
 Z=\prod_{e\in\mathcal I}z_{e,0}z_{e,1}.
\]
Both $P$ and $Z$ have degree $2M$. In a choice of one term from every factor contributing to the coefficient $[Z]P$, the choice of the second or first term of $Q_i$ forces every adjacent $P_f$ to choose the complementary rail. Thus all occurrences encode a single Boolean value of variable $i$. The chosen local assignment must satisfy each clause. Conversely every satisfying assignment makes exactly one consistent choice of all $P_f$ and $Q_i$ terms. All coefficients are positive integers equal to one before expansion, so this proves the exact identity
\[
 [Z]P=\#\operatorname{SAT}(\Phi)
\]
over $\mathbb Z$, including its reduction modulo two.

Adjoin $z_0$, set $r=2M+1$ and $F=z_0P$, and reduce coefficients modulo two. No $P_f$ or $Q_i$ is zero: their monomials are distinct and every clause has a satisfying local assignment. A polynomial ring over $\mathbb F_2$ is an integral domain, so $F\ne0$. It is homogeneous of degree $r$ in $r$ variables. Each rail occurs in exactly two factors, and the total number of factor incidences and multiplication gates needed to describe $F$ is $O(M)$; the eight-term local sums cost only a constant per clause. Moreover
\[
 [z_0Z]F=\#\operatorname{SAT}(\Phi)\pmod2.
\]
If an empty clause was present, use instead $F=z_1^3$ in three variables, whose relevant coefficient is zero. For the finitely many zero-variable cases use $z_1^3$ or $z_1z_2z_3$ according as the count is zero or one. These are again nonzero homogeneous degree-three equations.

We give the coherent calculation, which does not require smoothness. For every nonzero homogeneous degree-$r$ form $F$ in $r$ variables, its multiplication map on projective space is injective and gives
\[
 0\longrightarrow\mathcal O_{\mathbb P^{r-1}}(-r)
 \xrightarrow{\,F\,}\mathcal O_{\mathbb P^{r-1}}
 \longrightarrow\mathcal O_X\longrightarrow0.
\]
On the standard coordinate affine cover, the top Cech cohomology of $\mathcal O(-r)$ is represented by homogeneous Laurent monomials of total degree $-r$ having every exponent negative; monomials with a nonnegative exponent lie in a lower intersection and are zero in the top quotient. There is exactly one remaining monomial,
\[
 b=\frac1{z_1\cdots z_r}.
\]
The same Cech complex, split by Laurent monomial, is acyclic in intermediate degrees: for each exponent pattern other than the all-negative top pattern or the all-nonnegative global-section pattern, its nonzero components form the augmented simplex complex on the coordinates that may be inverted. Pairing subsets by adjoining a fixed admissible coordinate gives a contracting homotopy. Consequently $H^{r-1}(\mathbb P^{r-1},\mathcal O(-r))$ is the one-dimensional span of $b$, all its other cohomology groups vanish, and $\mathcal O$ has only $H^0=\mathbb F_2$. The long exact sequence gives the displayed cohomology of $X$ and the canonical boundary identification of its top group with the span of $b$.

The Frobenius square map on $\mathcal O_X$ lifts to the square map on $\mathcal O_{\mathbb P^{r-1}}$ and the map
\[
 a\longmapsto F a^2
\]
on $\mathcal O(-r)$, because $F(Fa^2)=(Fa)^2$. The latter map has the correct degree $-r$ when applied to a homogeneous section of degree $-r$. Naturality of the boundary map therefore identifies Frobenius on the top group with
\[
 b\longmapsto\frac{F}{(z_1\cdots z_r)^2}.
\]
In the top Cech quotient a numerator monomial survives only when every exponent is at most one. Since the numerator has total degree exactly $r$, that monomial must be $z_1\cdots z_r$. Thus the map is multiplication by $[z_1\cdots z_r]F$. This is precisely the rank-one version of Katz's coefficient formula, now verified for this possibly singular hypersurface. Relabeling $z_1,\ldots,z_r$ as $z_0$ and the rails proves the Frobenius assertion. SGA7 II XXII Corollary3.2 also yields
\[
 \#X_\Phi(\mathbb F_2)\equiv1+\#\operatorname{SAT}(\Phi)\pmod2.
\]

For hardness, the construction is a polynomial-time reduction from parity-3-SAT; that problem is $\oplus\mathrm P$-complete by the parsimonious circuit encoding in NN-C005. For membership on the specified factored class, nondeterministically choose one monomial from each displayed factor, and accept exactly when their exponent vectors sum to the all-one vector. There are constantly many choices at each factor, and exponent checking takes polynomial time. The number of accepting choices is the coefficient, so its parity is in $\oplus\mathrm P$. A polynomial-time Frobenius-entry algorithm therefore decides every $\oplus\mathrm P$ language, giving $\mathrm P=\oplus\mathrm P$; Valiant--Vazirani then gives $\mathrm{NP}=\mathrm{RP}$ as in NN-C006.

For E's upper bound one may simply count the original $\Phi$. There is also a direct bound from the factored input: encode the constant number of monomial choices of each factor by a constant number of bits, add validity constraints, and impose for every rail that exactly one of its two adjacent factors selects it. These are constant-arity Boolean constraints on $O(n+m)$ choice bits. There are $O(M)$ constraints, and $M=O(n+m)$, so their truth-table CNF has fixed width and bounded density. E Theorem1.1 gives time $2^{O((n+m)/\sqrt{\log(n+m)})}$ up to polynomial factors. Thus the single entry has E's subexponential algorithm, while its polynomial-time computation remains the precise parity-counting gate.
\end{proof}
\paragraph{Status and complexity impact.}
Complete construction and singular coherent calculation proved; full self-check and finite coefficient/point-count checks recorded. Complete deduction-proof review by root: passed; imported input theorems are not claimed independently proved. This strictly sharpens the effective-presentation issue in NN-C006. Exponential cohomology rank is not the only obstruction: here the relevant basis and rank are explicitly known and the entire missing matrix is one bit, yet computing that bit is parity-counting complete. A compact equation, a one-dimensional coherent cohomology group, and the existence of a trace formula do not constitute an efficient algorithm for Frobenius. The result concerns a succinct factored input; expansion into all monomials may already cost exponential time. It gives neither a lower bound for unrestricted arithmetic circuits nor a claim that the smooth hypersurface case is hard. A and M do not supply the missing entry.

\paragraph{Geometric and representation dictionary.}
The construction is already a scheme-theoretic theorem: parity-3-SAT is encoded in the absolute-Frobenius endomorphism of an explicitly trivialized rank-one coherent cohomology space of a projective hypersurface. Its sheaf-theoretic exact sequence and Cech basis prove the dictionary directly. Katz's original pp.35--36 locate the corresponding smooth Hasse--Witt formula, and SGA7 II XXII3.1--3.2 locate the zeta congruence and trace-point-count identity for arbitrary proper schemes. These reducible hypersurfaces are not asserted to furnish pure l-adic realizations of canonical bases, and their one-dimensional coherent Frobenius space is not a GCT representation multiplicity. The theorem identifies an effective arithmetic operation inside the geometric formalism, rather than importing Boolean hardness merely by analogy.



\section{NN-C024 --- The parity gate survives smooth hypersurfaces and genuine finite Weil objects}
\begin{proposition}[NN-C024]
Let $F$ be the degree-$r$ form constructed in NN-C023, so $r$ is odd and $c=[z_1\cdots z_r]F=\#\operatorname{SAT}(\Phi)\bmod2$. Put
\[
 H=\sum_{i=1}^r z_i^r,\qquad F_u=H+uF.
\]
The hypersurface $V(F_u)$ over $\mathbb F_2(u)$ is smooth and projective, its top coherent cohomology is a known one-dimensional space, and its absolute Frobenius is semilinear with matrix coefficient $uc$. Thus deciding whether the linearized Frobenius on this line is zero is parity-counting complete, even for this guaranteed generically smooth one-parameter family.

There is also a randomized polynomial-time reduction to smooth projective hypersurfaces over explicitly constructed finite fields $\mathbb F_q$, with $\log_2q=O(r\log r)$. Its output is smooth with probability greater than three quarters; on every smooth output the $\mathbb F_q$-linear $q$-power Frobenius scalar on the top coherent line is exactly $c$. Therefore a polynomial-time procedure correct on smooth inputs for this scalar, or for the point-count parity, would imply $\oplus\mathrm P\subseteq\mathrm{BPP}$ and then $\mathrm{NP}=\mathrm{RP}$. This randomized promise reduction does not assert $\mathrm P=\oplus\mathrm P$ for the finite-field problem.
\end{proposition}
\paragraph{Exact inputs.}
NN-C023. homogeneous coefficient encoding and complete coherent Frobenius calculation.

SGA1, ExposeII, Section4. Theorem4.10(iii) and Corollary4.12, Jacobian criterion for a hypersurface; text edition labels Theorem10 and Corollary12 within Section4.

Katz, Algebraic Solutions of Differential Equations, Inventiones Mathematicae18(1972),1--118. Algorithm2.3.7.14 pp.35--36; rank-one Special Case2.3.7.17 p.36.

SGA 7 II, Expose XXII. Theorem 3.1 (3.1.1); Corollary 3.2 (3.2.1).

Deligne, La conjecture de Weil I, Publications Mathematiques IHES43(1974),273--307. Theorem1.6, printedp.276.

Valiant--Vazirani, NP is as easy as detecting unique solutions, TCS 47 (1986), 85--93. Theorem 1.1; Corollaries 3.1--3.4, especially Corollary 3.2 p.91.

E: The Exponential Time Hypothesis is False, supplied LaTeX. Theorem1.1, for the inherited evaluation bound.

\begin{proof}
Because $r$ is odd, the partial derivatives of $H$ in characteristic two are $z_i^{r-1}$. We give an explicit degree bound for bad parameters without computing a discriminant. Set
\[
 L=r(r-2)+1,\qquad
 T=\binom{L+r-1}{r-1}=\binom{r(r-1)}{r-1}.
\]
Every degree-$L$ monomial $m$ is divisible by some $z_i^{r-1}$: otherwise every exponent would be at most $r-2$ and their sum at most $L-1$. Choose the least such $i=i(m)$. For each of the $T$ degree-$L$ monomials form
\[
 \frac{m}{z_{i(m)}^{r-1}}\frac{\partial F_u}{\partial z_{i(m)}}
 =m+u\frac{m}{z_{i(m)}^{r-1}}\frac{\partial F}{\partial z_{i(m)}}.
\]
In the degree-$L$ monomial basis their coefficient matrix is $B(u)=I+uA$ for a fixed matrix $A$ over $\mathbb F_2$. Consequently
\[
 \Delta(u)=\det B(u),\qquad \Delta(0)=1,\qquad \deg\Delta\leq T.
\]
Whenever $\Delta(a)\ne0$, every degree-$L$ monomial is in the ideal generated by the partial derivatives of $F_a$. They cannot have a common nonzero zero over an algebraic closure, since a nonzero coordinate would give a nonvanishing monomial $z_i^L$ in that ideal. The hypersurface is therefore smooth by SGA1 II4.10(iii), or its hypersurface Corollary4.12. In a projective coordinate chart this is the usual nonvanishing-differential criterion; Euler's identity $\sum_i z_i\partial_iF_a=rF_a$ ensures that at a point of $V(F_a)$ some derivative in a chart coordinate is nonzero whenever the homogeneous gradient is nonzero.

This proves smoothness over $\mathbb F_2(u)$, since $\Delta$ is a nonzero polynomial. It also proves that at most $T$ field elements can be excluded by this sufficient smoothness test. The exponentially sized matrix $B$ is used only in the proof of this bound. The reduction never forms $B$, computes $\Delta$, expands $F$, or tests smoothness. The integer $T$ has $O(r\log r)$ bits and can be computed by the product formula for the binomial coefficient in polynomial time. The form $H+uF$ has a polynomial-size circuit. It is nonzero for every $u$: the constructed $F$ is divisible by its extra coordinate, whereas $H$ is not; the explicit cubic edge-case replacements in NN-C023 are also not scalar multiples of $H$.

For completeness the required finite field can be built deterministically with polynomial bit cost. Choose $b$ to be a power of two, large enough that $2^b\geq8(T+1)$ and minimal with that property. Then $b=O(r\log r)$. Build fields of degrees $1,2,4,\ldots,b$ by successive quadratic extensions. In a constructed field $K$ of size $2^d$, the absolute trace
\[
 \operatorname{Tr}(a)=\sum_{j=0}^{d-1}a^{2^j}
\]
is a nonzero $\mathbb F_2$-linear functional: it takes values in $\mathbb F_2$, and its defining nonzero polynomial has degree $2^{d-1}<2^d$, so cannot vanish on every element. Evaluate it on a supplied vector-space basis to find an $a$ of trace one. Then $Y^2+Y+a$ has no root in $K$, since $\operatorname{Tr}(y^2+y)=0$, and hence is irreducible. Adjoin a root and use the evident doubled basis. Repeated squaring, testing the basis vectors, and the tower's arithmetic all take polynomial time in $b$.

If a single defining irreducible polynomial is desired, the last adjoined element $\theta$ has degree $b$ over $\mathbb F_2$: its degree divides the power of two $b$, and every proper subfield of $\mathbb F_{2^b}$ lies in the preceding field $\mathbb F_{2^{b/2}}$, whereas $\theta$ does not. Equivalently, any smaller degree divides $b/2$ and makes $\theta$ a root of $X^{2^{b/2}}-X$, whose complete root set is that preceding field. Linear algebra on $1,\theta,\ldots,\theta^b$ gives its monic minimal polynomial and the change to the power basis in polynomial time. Thus the construction uses an explicit polynomial-length finite-field presentation, not an oracle for unspecified algebraic constants.

Retain a basis beginning with $1$. Choose the parameter $a$ uniformly from the affine half-space of field elements whose coefficient on that basis vector is one. Every such $a$ is nonzero, there are $q/2$ choices, and exactly $b-1$ fair random bits suffice. At most $T$ choices are roots of $\Delta$, so
\[
 \Pr[V(F_a)\text{ is not smooth}]\leq\frac{2T}{q}<\frac14.
\]
The complete reduction has worst-case polynomial running time; it does not use an unbounded rejection loop.

Now apply the coherent calculation of NN-C023, or Katz2.3.7.17 on a smooth fiber. The coefficient of $z_1\cdots z_r$ in $H$ is zero because $r\geq3$. Absolute $2$-power Frobenius on the top coherent line therefore sends its displayed basis vector $e$ to $ac\,e$ and is semilinear on coefficients. Over the rational function field the coefficient is $uc$, so its linearization is zero exactly when $c=0$ and is an isomorphism otherwise; one must not call this a linear endomorphism over $\mathbb F_2(u)$. Over $\mathbb F_q$, iterate this semilinear map $b$ times. Its $\mathbb F_q$-linear scalar is
\[
 (ac)^{1+2+\cdots+2^{b-1}}=(ac)^{q-1}=c,
\]
because $a\ne0$ and $c\in\{0,1\}$. The distinction between the semilinear matrix coefficient $ac$ and the linear $q$-Frobenius scalar $c$ is essential.

All other coherent groups are as in NN-C023: $H^0=\mathbb F_q$, top degree $r-2$ has rank one, and the remaining groups vanish. SGA7 II XXII3.1 and3.2 consequently give, on a smooth sampled fiber (indeed on every nonzero fiber),
\[
 \#V(F_a)(\mathbb F_q)\equiv1+c\pmod2,
 \qquad
 Z(V(F_a)/\mathbb F_q,t)\equiv\frac{1-ct}{1-t}\pmod2.
\]
Here $r-2$ is odd. In particular the full zeta series modulo two is $1$ when $c=1$ and $(1-t)^{-1}$ when $c=0$.

Suppose a polynomial-time procedure computes the top $q$-Frobenius scalar on every smooth input in this class. Clock it at its polynomial bound and treat malformed outputs or timeouts arbitrarily. Its behavior on nonsmooth inputs is unrestricted. On a sampled input, correctness nevertheless has probability greater than three quarters because the smooth event alone has that probability. Independent repetitions and majority vote therefore put parity-3-SAT, and hence all of $\oplus\mathrm P$, in BPP. The same argument applies to a smooth-input point-count parity procedure, after adding one. It is a randomized promise reduction; it gives no deterministic parity algorithm.

To justify the stated NP consequence, use Valiant--Vazirani isolation followed by the BPP parity procedure, with its error amplified below the inverse polynomial isolation success probability. This gives a BPP procedure for SAT. For a satisfiable formula, self-reduce one variable at a time with decision error at most $1/(10n)$ per query; with probability at least nine tenths all needed answers are correct and a satisfying assignment is found. Check the resulting assignment directly. On an unsatisfiable input no string passes that final check. Thus SAT is in RP, and NP=RP. This verified-witness step is required to obtain one-sided error.

The generic-family problem has a stronger deterministic implication: a deterministic polynomial-time test of whether its linearized top Frobenius is zero directly decides parity-3-SAT, so gives P=$\oplus$P. E instead computes $c$ through the original formula or the monomial-choice encoding of NN-C023. On the bounded-density family this takes $2^{O(r/\sqrt{\log r})}$ time up to polynomial factors, including the polynomial field-construction overhead.
\end{proof}
\paragraph{Status and complexity impact.}
Complete smooth-pencil, field-construction and randomized-promise proofs; complete deduction-level self-reading passed. Complete deduction-proof review by root: passed; imported input theorems are not claimed independently proved. This removes singularity from the geometric obstruction in NN-C023, with a precisely stated randomized finite-field gate. The good fibers are genuine smooth projective varieties with the pure l-adic Frobenius representations guaranteed by Deligne's theorem. Producing such objects and knowing purity still leaves parity information computationally encoded in a rank-one coherent Frobenius map. This proof itself gives a randomized specialization; NN-C030 subsequently closes its deterministic finite-field specialization gate by using the actual encoding's pairwise homogeneity. A polynomial bound on all l-adic Betti numbers or an efficient full expansion of the succinct equation is still not asserted. Neither result separates VP from VNP. A and M do not calculate this coefficient or supply the missing polynomial-time procedure.

\paragraph{Geometric and representation dictionary.}
Let $U=\operatorname{Spec}\mathbb F_2[u,\Delta(u)^{-1}]$. The pencil defines a smooth projective hypersurface family over $U$, with a canonically based coherent line $R^{r-2}\pi_*\mathcal O$; its linearized relative Frobenius is multiplication by $uc$. The open smooth family is established by an explicit homogeneous Jacobian minor, which is the local differential criterion of SGA1 II4.10--4.12. The finite-field trace and full zeta congruence are exactly SGA7 II XXII3.1--3.2. Every good closed fiber has pure l-adic cohomology by Deligne, Weil I1.6; purity constrains the complex absolute values of Frobenius eigenvalues and does not compute this coherent characteristic-two invariant. The construction supplies genuine finite Weil objects for these counting gadgets, not canonical realizations of GCT's nonstandard basis or its structure constants. It sharpens an effectiveness gate while leaving that separate functorial/canonical realization problem intact.



\section{NN-C025 --- Geometrically integral smooth quadratic affine schemes retain exact counts above parity}
\begin{proposition}[NN-C025]
A $k$-CNF $\Phi$ with $n\geq1$ variables and $m$ clauses has a deterministic polynomial-time construction of a geometrically integral smooth affine scheme $Y_\Phi/\mathbb F_2$ of dimension exactly $n$, together with a finite etale degree-two map
\[
             Y_\Phi\longrightarrow\mathbb A^n_{\mathbb F_2},
\]
such that
\[
        \#Y_\Phi(\mathbb F_2)=2\#\operatorname{SAT}(\Phi).
\]
It has a presentation by $s+1$ equations of degree at most two in $n+s+1$ coordinates, with $s=O(n+km+\log(m+2))$, and an explicitly displayed Jacobian minor identically equal to one. Thus rational-point existence is NP-complete on this promised class, exact point counting is $\#\mathrm P$-hard by one division by two, and point counting modulo four encodes parity-SAT. Point counting modulo two is identically zero on the entire class. E evaluates these counts through the original $n$ witnesses, with its exact Theorem1.1 bound; geometric auxiliary coordinates do not become independent Boolean witnesses.
\end{proposition}
\paragraph{Exact inputs.}
NN-C006 and prior CX-E07. finite Boolean solution schemes and the correct affine trace boundary; used only for contrast, not repeated.

E: The Exponential Time Hypothesis is False, supplied LaTeX. Theorem1.1, original-variable exact k-CNF counter.

SGA1, ExposeII, Section4. Theorem4.10(iii) and Corollary4.12, Jacobian smoothness criterion.

Deligne, La conjecture de Weil I, PMIHES43(1974),273--307. Section1.3--1.5, especially compact-support trace equation1.5.1, printedpp274--275; contrast projective smooth Theorem1.6 p276.

Valiant--Vazirani, NP is as easy as detecting unique solutions, TCS 47 (1986), 85--93. Theorem 1.1; Corollaries 3.1--3.4, especially Corollary 3.2 p.91.

\begin{proof}
For a clause $C$, let $v_C(x)$ be the product of its false-literal indicators: a positive literal contributes $1+x_i$, and a negative literal contributes $x_i$. Put
\[
 S_\Phi(x)=\prod_C(1+v_C(x)),\qquad P_\Phi(x)=1+S_\Phi(x).
\]
On $\mathbb F_2^n$, $P_\Phi(x)=0$ exactly for satisfying assignments. The polynomial has degree at most $km$ and has a fan-in-two arithmetic circuit of $O(n+km)$ gates over $\mathbb F_2$. Empty products and empty clauses have their usual values, so tautological and contradictory inputs are included.

Choose the odd integer $d=2km+3$ and form
\[
        f_\Phi(x)=P_\Phi(x)+x_1^d+x_1.
\]
The added two terms vanish on every bit assignment. They ensure that the total degree is exactly the odd integer $d$, with leading monomial $x_1^d$. Binary powering takes $O(\log(d+1))$ additional gates. Define
\[
 Y_\Phi=\operatorname{Spec}
 \mathbb F_2[x_1,\ldots,x_n,t]/(t^2+t+f_\Phi(x)).
\]
The coordinate algebra is free of rank two over $\mathbb F_2[x]$, with basis $1,t$, because the equation is monic in $t$. Its relative derivative in $t$ is one. The hypersurface Jacobian criterion therefore makes its projection smooth of relative dimension zero, hence etale; alternatively it is finite flat and its relative differentials vanish. Composition with the smooth affine-space projection makes $Y_\Phi$ smooth of dimension $n$.

Geometric integrality is also explicit. Over any algebraic closure $\overline k$ of $\mathbb F_2$, a factorization of $T^2+T+f_\Phi$ over $\overline k(x)$ would give a rational function $h$ with $h^2+h=f_\Phi$. Write $h=A/B$ in coprime polynomials. The relation $A^2+AB=f_\Phi B^2$ forces every irreducible factor of $B$ to divide $A$, so $B$ is constant. Thus $h$ is a polynomial. If nonconstant, $h^2+h$ has even total degree $2\deg h$; if constant, it has degree zero or is zero. Neither possibility equals the odd-degree polynomial $f_\Phi$. The quadratic is therefore irreducible over the fraction field. Its monic quotient embeds in that field extension, so is a domain. This proves geometric integrality, not merely connectedness or reducedness.

For an $\mathbb F_2$ point, $t^2+t=0$ for both possible values of $t$. Each satisfying input $x$ therefore has exactly two points above it, and every nonsatisfying input has none. This proves the exact factor-two identity. No equations $x_i^2-x_i=0$ were imposed on the scheme: its geometric dimension is $n$, and its projection has two geometric points over every geometric base point. The vanishing or splitting of the rational fiber over a bit assignment is the arithmetic distinction being used.

To obtain the quadratic complete-intersection presentation, order the gates of the circuit computing $f_\Phi$ topologically. Introduce one coordinate $y_j$ for each gate, with equation
\[
 y_j-a-b=0,\quad y_j-ab=0,\quad\hbox{or}\quad y_j-a-1=0
\]
as appropriate, where $a,b$ are inputs or earlier gates. In characteristic two subtraction equals addition. Add the final equation $t^2+t+y_{\rm out}=0$. In the columns $(y_1,\ldots,y_s,t)$ and this order of equations, the Jacobian submatrix is lower triangular with diagonal entries all one: each gate equation has derivative one in its new gate, and the final equation has derivative one in $t$. Its determinant is identically one over the whole coordinate ring. Eliminating the gates, in order, recovers precisely the displayed Artin--Schreier coordinate algebra. The gate equations are a regular sequence by successive elimination, and the final monic polynomial is a nonzerodivisor. Thus this is a complete intersection, with constant-degree sparse equations and a proved global smoothness certificate. All coordinates and equations are written in polynomial time.

The exact-count complexity claim concerns this presentation. Guessing its polynomially many bit coordinates and checking its equations is a $\#\mathrm P$ computation. From its integer count one recovers $\#\operatorname{SAT}(\Phi)$ by division by two, so hardness is under a polynomial-time numerical reduction, not a falsely claimed parsimonious reduction for an always-even output. The divided count is itself a $\#\mathrm P$ function: guess just $x$ and check $P_\Phi(x)=0$. Its parity is exactly half of the point count modulo four. Consequently a deterministic polynomial-time mod-four algorithm gives P=$\oplus$P and, by Valiant--Vazirani, NP=RP. A polynomial-time existence algorithm would directly give P=NP. All these are conditional gates, not conclusions from smoothness.

For the upper bound, the constructor retains the source $k$-CNF presentation. Each assignment to its $n$ input bits determines every gate coordinate, and the extra coordinate $t$ contributes the known multiplicity two. E Theorem1.1 therefore gives
\[
 \operatorname{poly}_{k}(n+m)\,
 2^{O_k(n\log\log n/\sqrt{\log n})}
\]
bit operations and polynomial space, with exponent $O_{k,D}(n/\sqrt{\log n})$ when $m\leq Dn$. Applying a sparse-equation counter blindly to all $n+s+1$ coordinates would instead use that larger variable parameter. The smaller bound follows from the explicit functional graph and the original clause structure; it is not claimed for arbitrary succinct arithmetic circuits. Inputs with no variables have their constant count evaluated separately.
\end{proof}
\paragraph{Status and complexity impact.}
Complete affine construction, geometric-integrality and parameter proofs; complete deduction-level self-reading passed. Complete deduction-proof review by root: passed; imported input theorems are not claimed independently proved. This strengthens the affine boundary of NN-C006: even geometric integrality, smoothness, an everywhere invertible quadratic Jacobian minor, and a finite etale cover of affine space do not remove exact counting hardness. It is different from the earlier zero-dimensional Boolean solution scheme, whose number of components already equals the count. Here the geometric object is integral of the original witness dimension, and every rational count is even. The precise information lost by reduction to characteristic two is one binary digit: modulo four retains it. The construction does not supply a smooth projective compactification in polynomial-size ordinary coordinates, nor permit the pure smooth-projective conclusion of Deligne1.6. It does not resolve P versus NP or VP versus VNP; E changes the available running time while the reductions remain valid.

\paragraph{Geometric and representation dictionary.}
For every fixed-width CNF $\Phi$ on $n\geq1$ Boolean variables there is an explicitly constructed polynomial $f_\Phi\in\mathbb F_2[x_1,\ldots,x_n]$, represented by a polynomial-size arithmetic circuit and having odd degree, such that
\[
 U_\Phi=\operatorname{Spec}\mathbb F_2[x,t]/(t^2+t-f_\Phi(x))
\]
is smooth and geometrically integral of dimension $n$, finite etale of degree two over $\mathbb A^n$, and satisfies $\#U_\Phi(\mathbb F_2)=2\#\operatorname{SAT}(\Phi)$. Its graph presentation is a polynomial-size complete intersection of quadratic equations with a displayed triangular Jacobian minor equal to one. In particular its count is always zero modulo two, whereas its count modulo four is twice the Boolean parity. The proof above constructs $f_\Phi=1+S_\Phi+x_1^d+x_1$ with $d$ odd and larger than the verifier degree; the odd pole at infinity proves that $f_\Phi$ is not $h^2+h$ over the algebraic closure. Thus this geometric integrality claim has an explicit certificate argument.

The map is an Artin--Schreier torsor under the constant group $\mathbb Z/2\mathbb Z$. Over a satisfying rational input its fiber is $\operatorname{Spec}(\mathbb F_2\times\mathbb F_2)$; over a nonsatisfying one it is $\operatorname{Spec}\mathbb F_4$. These fibers have the same geometric degree, but respectively two and zero rational points. The multiplication trace of the unit in either rank-two algebra is zero in characteristic two; its algebraic trace locator is EGA IV4,18.2.1--18.2.2, printedp111. SGA4 XVII6.2.3 is the corresponding sheaf-theoretic finite trace, a distinct formalism. The whole affine point count is the alternating compact-support etale trace, by Deligne Weil I(1.5.1); it is not the trace on infinite-dimensional ordinary affine coherent $H^0$. Its integer residue modulo four is a separate operation, not a map from an odd-adic coefficient field to $\mathbb Z/4\mathbb Z$. Exact point counting on these circuit-presented smooth objects is counting-complete by division by two, while E evaluates the constructed objects through the original $n$-variable formula. No canonical GCT multiplicity or positive polytope formula is identified.



\section{NN-C026 --- Deterministic pure projective curves encode counting in the next binary digit}
\begin{proposition}[NN-C026]
From a $k$-CNF $\Phi$ with $n\geq1$ variables and $m$ clauses one can deterministically construct, in polynomial time, an explicitly presented field $\mathbb F_q$ with $\log_2q=O(n)$ and a smooth geometrically integral projective curve $C_\Phi/\mathbb F_q$ such that
\[
               \#C_\Phi(\mathbb F_q)=1+2\#\operatorname{SAT}(\Phi).
\]
The curve is given by a polynomial-size arithmetic circuit for a monic odd-degree Artin--Schreier equation, together with its explicit weighted-projective completion. Its degree and genus are exponentially large and are encoded in binary. There is no promise of a polynomial-degree dense equation or a polynomial-genus curve.

The integer count, and the integer trace of the pure weight-one l-adic Frobenius representation, are therefore $\#\mathrm P$-hard from these succinct descriptions. Reduction of the count modulo four is parity-counting complete after subtracting one and dividing by two. Nevertheless every curve in this class has
\[
 Z(C_\Phi/\mathbb F_q,t)\equiv(1-t)^{-1}\pmod2,
\]
and its entire coherent Frobenius matrix is determined by a simple entry formula from its genus and field, independently of the satisfying-assignment count. This entry formula does not require listing an exponentially large matrix. Thus a genuine pure finite Weil object can be constructed deterministically here, while the next integer counting digit remains a separate effective gate.
\end{proposition}
\paragraph{Exact inputs.}
NN-C024. deterministic quadratic-tower finite-field construction.

NN-C025. Boolean indicator and odd-degree Artin--Schreier irreducibility arguments.

E: The Exponential Time Hypothesis is False, supplied LaTeX. Theorem1.1, evaluation through the original formula.

SGA1, ExposeII, Section4. Theorem4.10(iii) and Corollary4.12, used on the two explicit affine curve charts.

The Stacks Project, Cohomology of Schemes. Lemma30.2.2, tag01XB (affine acyclicity); Lemma30.2.6, tag01XD (affine-cover Cech comparison).

SGA 7 II, Expose XXII. Theorem 3.1 (3.1.1); Corollary 3.2 (3.2.1).

Deligne, La conjecture de Weil I, PMIHES43(1974),273--307. Trace equation1.5.1 p275; Theorem1.6 p276; Poincare duality Theorem2.3 and discussion2.4--2.5 p281.

Valiant--Vazirani, NP is as easy as detecting unique solutions, TCS 47 (1986), 85--93. Theorem 1.1; Corollaries 3.1--3.4, especially Corollary 3.2 p.91.

\begin{proof}
Let $b$ be the least power of two with $b\geq n$, and put $q=2^b$. Then $n\leq b<2n$ unless $n=1$, when $b=1$. The field tower in NN-C024 constructs $K=\mathbb F_q$, its arithmetic, and a basis $e_1,\ldots,e_b$ deterministically in polynomial time. If needed it also supplies a monic irreducible defining polynomial of degree $b$. The absolute trace $\operatorname{Tr}:K\to\mathbb F_2$ is nonzero, by the polynomial-degree proof in NN-C024. Its pairing $(a,x)\mapsto\operatorname{Tr}(ax)$ is nondegenerate: multiplication by any nonzero $a$ is a bijection, so this pairing cannot vanish for every $x$. Linear algebra over $\mathbb F_2$ therefore constructs the dual basis $\beta_i$ with
\[
 \operatorname{Tr}(\beta_i e_j)=\delta_{ij}.
\]
Define the linearized coordinate polynomials
\[
 L_i(X)=\operatorname{Tr}(\beta_i X)
       =\sum_{j=0}^{b-1}(\beta_iX)^{2^j}.
\]
They have degree at most $q/2$, polynomial-size repeated-squaring circuits, and for $x=\sum_j a_j e_j\in K$ satisfy $L_i(x)=a_i\in\mathbb F_2$.

Let $S_\Phi$ be the Boolean satisfaction polynomial in NN-C025. Use all $b$ coordinates to define
\[
 P(X)=1+S_\Phi(L_1(X),\ldots,L_n(X))
                 \prod_{i=n+1}^{b}(1+L_i(X)).
\]
For every $x\in K$, $P(x)$ is zero or one; it is zero exactly when its first $n$ coordinates satisfy $\Phi$ and all remaining coordinates are zero. Hence exactly $\#\operatorname{SAT}(\Phi)$ field elements have $P(x)=0$, with no padding multiplicity. Its formal degree is at most $(km+b)q/2$ and its circuit size is polynomial in $n+m$ for fixed $k$.

Find $\lambda\in K$ with $\operatorname{Tr}(\lambda)=1$ by testing the basis vectors. Choose
\[
 d=(km+b+1)q+1,
 \qquad
 f(X)=\lambda P(X)+(X^q+X)X^{d-q}.
\]
The integer $d$ is odd, its binary length is $O(n+\log(m+2))$, and it exceeds the degree bound for $P$. The polynomial $f$ is monic of degree exactly $d$. All powers shown can be computed by binary powering, so its circuit and coefficient presentation have polynomial bit length. On $K$ the extra term vanishes, and
\[
           \operatorname{Tr}(f(x))=P(x).
\]
The linear map $y\mapsto y^2+y$ on $K$ has kernel $\{0,1\}$ and image equal to the trace-zero hyperplane. The image is contained in that hyperplane by telescoping the trace, and both have dimension $b-1$. Therefore the affine equation $y^2+y=f(x)$ has exactly two $K$-solutions in $y$ when $P(x)=0$, and none otherwise. Its number of affine points is exactly $2\#\operatorname{SAT}(\Phi)$.

Write $d=2g+1$ and set $w=g+1$. Let $f^h(X,Z)=Z^d f(X/Z)$ be its degree-$d$ homogenization. It has a polynomial-size homogeneous arithmetic circuit: at each addition gate multiply the lower formal-degree input by the appropriate power of $Z$, and homogenize multiplication gates multiplicatively. Every needed exponent has $O(\log d)$ bits, so binary powering keeps the overhead polynomial. Define
\[
 C_\Phi=V\bigl(Y^2+YZ^w-Z f^h(X,Z)\bigr)
              \subset\mathbb P_K(1,w,1).
\]
All terms have weighted degree $2w=d+1$. Weighted projectivity does not require constructing an exponentially large ordinary embedding. For existence of such an embedding, the $w$th Veronese ring of $K[X,Y,Z]$ with weights $(1,w,1)$ is generated in degree one by $Y$ and $X^aZ^{w-a}$ for $0\leq a\leq w$. Its Proj is the same weighted projective space and is a projective closed subscheme of ordinary projective space. The latter embedding has exponentially many coordinates and is not part of the polynomial-size output.

The curve is covered by the two weight-one charts $Z\ne0$ and $X\ne0$: a hypothetical point with $X=Z=0$ would have $Y^2=0$ and hence all homogeneous coordinates zero. On $Z=1$ its equation is $y^2+y=f(x)$, whose $y$ derivative is one. On $X=1$, put $z=Z/X$ and $v=Y/X^w$; the equation is
\[
       v^2+v z^w=z^{d+1}f(1/z).
\]
At $z=0$ this has the unique point $v=0$. Its derivative in $z$ is one at that point, because $f$ is monic of degree $d$ and the right side starts with $z$. Thus the curve is smooth everywhere, over an algebraic closure as well, by the local Jacobian criterion. There is exactly one $K$-rational point at infinity. The ambient possible singular point $[0:1:0]$ is off the curve, so it supplies no exception to this chart argument.

The affine polynomial $y^2+y+f(x)$ is geometrically irreducible by the odd-degree proof in NN-C025. The smooth point at infinity lies on the closure of that affine curve: on its smooth one-dimensional chart, $z$ is a nonzero regular function and its zero set is just that point. Thus the projective curve is geometrically integral. This proves
\[
              \#C_\Phi(K)=1+2\#\operatorname{SAT}(\Phi)
\]
with no smoothness promise and no randomized parameter selection.

We compute the genus and the coherent Frobenius directly, so no unverified genus formula is needed. On the intersection of the two affine charts, the coordinate ring has the free $K[x,x^{-1}]$-basis $1,y$. Sections from the $Z=1$ chart have coefficients in $K[x]$ for both basis elements. Sections from the $X=1$ chart have constant coefficients in $K[x^{-1}]$ and $y$-coefficients in $x^{-w}K[x^{-1}]$, since $v=x^{-w}y$. The Cech quotient has no constant part, and its $y$-part has the basis
\[
       e_j=[x^{-j}y],\qquad 1\leq j\leq g.
\]
The same intersection calculation gives $H^0(C_\Phi,\mathcal O)=K$. The two affine opens and their affine intersection compute coherent cohomology by affine acyclicity and the Cech comparison (Stacks, Lemmas30.2.2 and30.2.6); hence $H^1$ has dimension $g=(d-1)/2$. In particular the genus is exponential in the original witness parameter, despite the short equation.

Absolute two-power Frobenius sends
\[
 e_j\longmapsto[x^{-2j}y^2]
       =[x^{-2j}y]+[x^{-2j}f(x)].
\]
The second summand is a Laurent polynomial in $x$ with no $y$ term, and is a Cech coboundary. The first is $e_{2j}$ if $2j\leq g$ and is a coboundary otherwise. Thus its semilinear matrix is the coefficient-independent shift
\[
 e_j\longmapsto
 \begin{cases}e_{2j},&2j\leq g,\\0,&2j>g.\end{cases}
\]
Iterating $b$ times gives the $K$-linear $q$-Frobenius shift $e_j\mapsto e_{qj}$, with out-of-range terms zero. It is nilpotent, with known rank $\lfloor g/q\rfloor$, independently of $f$ and $\Phi$. Applying SGA7 XXII3.1 to this proper curve gives the asserted complete zeta congruence $(1-t)^{-1}$, not just equality of first point-count residues. At fixed $n,m,k$, both the field and this entire based coherent Frobenius description are fixed, although satisfying-assignment counts can differ.

Finally let $\ell$ be any odd prime. The geometric l-adic cohomology is a genuine finite Weil representation. Its $H^1$ is pure of weight one by Deligne1.6. The connected curve has trace one on $H^0$ and trace $q$ on $H^2$, by the duality and degree discussion2.3--2.5. Deligne's trace formula1.5.1 gives the integer identity
\[
 \operatorname{Tr}(\operatorname{Frob}_q\mid H^1_{\rm et})
           =q+1-\#C_\Phi(K)
           =q-2\#\operatorname{SAT}(\Phi).
\]
Purity controls the eigenvalue magnitudes, but computing this integer trace, or its residue modulo four, recovers the count or its parity respectively. This is reduction of an integer trace, not reduction of an odd-adic field modulo four.

The divided count belongs to $\#\mathrm P$ directly on the whole succinct Artin--Schreier class: guess the $b$ coordinate bits of $x\in K$, evaluate the circuit for $f(x)$, and test $\operatorname{Tr}(f(x))=0$. For every monic odd-degree member, it equals $(\#C(K)-1)/2$. Our reduction proves its $\#\mathrm P$-completeness, and its parity is $\oplus\mathrm P$-complete. A deterministic polynomial-time mod-four count procedure would therefore give P=$\oplus$P and NP=RP; exact polynomial-time counting gives the stronger counting and Toda consequences recorded in NN-C006. E instead evaluates the constructed instances from the source $k$-CNF in its original $n$-variable time bound, with polynomial construction overhead. No efficient method for arbitrary high-degree circuit-presented curves is inferred.
\end{proof}
\paragraph{Status and complexity impact.}
Complete deterministic field, curve, Cech and counting proofs; complete deduction-level self-reading passed. Complete deduction-proof review by root: passed; imported input theorems are not claimed independently proved. This supplies a deterministic smooth projective finite-field counting object, strengthening the existence side of the geometric gate. It does not derandomize the low-degree rank-one-coherent construction of NN-C024: the curve degree and coherent rank are exponential, and the parity information now occurs modulo four rather than modulo two. The representation is an actual pure l-adic representation attached to a smooth curve, but obtaining its integer trace from the succinct input remains counting-hard. All characteristic-two coherent Frobenius data above are already coefficient-independent, showing by explicit construction why that residue-level realization cannot determine the next integer digit. These are valid reduction barriers, not unconditional complexity separations or a GCT canonical-basis realization.

\paragraph{Geometric and representation dictionary.}
For each fixed-width CNF $\Phi$ on $n\geq1$ variables, the proof explicitly constructs a finite field $K=\mathbb F_q$, $q=2^b$ with $b$ polynomial in the input size, and a monic odd-degree circuit-presented polynomial $f_\Phi(X)\in K[X]$. Put $d=\deg f_\Phi=2g+1$ and $w=g+1$. Its weighted projective completion
\[
 C_\Phi:\quad Y^2+YZ^w=Z^{d+1}f_\Phi(X/Z)
       \subset\mathbb P(1,w,1)
\]
is a smooth geometrically integral projective curve of genus $g$, with one rational point at infinity and
\[
 \#C_\Phi(K)=1+2\#\operatorname{SAT}(\Phi).
\]
The polynomial and the integers $d,g$ have polynomial-size circuit/binary descriptions; $d$ and $g$ themselves may be exponential. The two affine charts in the proof establish smoothness, completion and the Cech basis $e_j=[x^{-j}y]$, $1\leq j\leq g$, for $H^1(C_\Phi,\mathcal O)$. Absolute Frobenius sends $e_j$ to $e_{2j}$, with out-of-range indices zero; $q$-Frobenius sends it to $e_{qj}$. These coherent operators are nilpotent, independently of the encoded satisfying count.

For every prime $\ell\ne2$, the actual $H^1_{\rm et}(C_{\overline K},\mathbb Q_\ell)$ is pure of weight one and has Frobenius trace $q-2\#\operatorname{SAT}(\Phi)$. SGA7 XXII3.1 explains the common full zeta reduction modulo two; Deligne Weil I1.5.1,1.6 and2.3--2.5 supply the integral trace and purity interfaces. This proves a precise contrast between the characteristic-two coherent realization and the next binary digit of the integer etale trace. The divided count $(\#C(K)-1)/2$ is a counting-complete function on the succinct class, and its parity is parity-counting-complete. E evaluates the constructed instances from their original Boolean witnesses; it does not give a polynomial algorithm on arbitrary high-degree circuit-presented curves. No plethysm/Kronecker multiplicity identification or canonical-basis construction is claimed.



\section{NN-C027 --- Smooth integral cubic arithmetic schemes carry counting-hard pure Tate multiplicities}
\begin{proposition}[NN-C027]
Let $N=\#\operatorname{SAT}(\Phi)$ for a fixed-width CNF with $n\geq1$ input variables and $m$ clauses. In polynomial time one can construct a smooth affine scheme $X_\Phi/\mathbb Z$, with every geometric fiber integral, given by one sparse cubic equation in $O(n+m)$ coordinates. There are explicit integers $v,h=O(n+m)$, with $v,h\geq1$, for which
\[
 [X_\Phi]=\mathbb L^{v+h}+N\mathbb L^h,
 \qquad
 \#X_\Phi(\mathbb F_q)=q^{v+h}+Nq^h
\]
for every prime power $q$, and
\[
         \zeta_{X_\Phi}(s)=\zeta(s-v-h)\zeta(s-h)^N.
\]
The Grothendieck-class identity is proved by an explicit stratification. Separately, over any field $K$ and any prime $\ell$ invertible in $K$, the actual geometric compactly supported l-adic groups are
\[
 H_c^{2h}=\mathbb Q_\ell(-h)^N,
 \qquad H_c^{2(v+h)}=\mathbb Q_\ell(-v-h),
\]
and all others vanish, as representations of the absolute Galois group of $K$. Over finite fields these particular groups are pure Tate; this is proved, not inferred from their point counts alone.

For every nonreal nontrivial zeta zero $\rho$ in $0<\Re\rho<1$, one has
\[
 \operatorname{ord}_{s=h+\rho}\zeta_{X_\Phi}(s)
              =N\operatorname{ord}_{s=\rho}\zeta(s).
\]
The shifted zero locations are therefore controlled by A, while their multiplicity contains the original exact Boolean count. No use of M is made. E computes the multiplicity in its original $n$-variable time bound. Polynomial-time extraction of that multiplicity from these cubic presentations would compute every $\#\mathrm P$ function in polynomial time.

There is also a literal quadratic-character variant over $B=\mathbb Z[1/2]$:
\[
 X^-_\Phi:\ t^2+1=\sum_i y_i f_i(x),\qquad
 \zeta_{X^-_\Phi}(s)=\zeta^{(2)}(s-v-h)L(s-h,\chi_{-4})^N,
\]
where $\zeta^{(2)}(s)=(1-2^{-s})\zeta(s)$. Its lower group is
$H_c^{2h}=(\chi_{-4}\otimes\mathbb Q_\ell(-h))^N$ and its top group is the same Tate line. At any nonreal nontrivial zero $\rho$ of this Dirichlet L-function, its shifted zeta has order $N\operatorname{ord}_\rho L$. Both variants are smooth with geometrically integral fibers; their fixed-$\mathbb F_3$ point-count parity is $1+N\bmod2$.
\end{proposition}
\paragraph{Exact inputs.}
Prior catalogue CX-E07. constant finite etale Boolean Z-scheme with degree equal to the satisfying-assignment count; background construction only.

E: The Exponential Time Hypothesis is False, supplied LaTeX. Theorem1.1, original-variable counting bound.

A: Alpoge, Remark; designated RH and Dirichlet-GRH input theorems. riemannHypothesis and riemannZeta\_eq\_zero\_re, Main.lean lines4,6; LFunction\_eq\_zero\_re\_of\_isPrimitive, Main.lean line12, as recorded in OURS\_ALPOGE\_LEAN\_MAP\_20261120.md Section3.

SGA1, ExposeII, Section4. Theorem4.10(iii) and Corollary4.12, relative Jacobian criterion over the integer base.

SGA4, ExposeXVII. Theorem5.2.6 and Proposition5.2.8 (!-base change and geometric stalks); Theorem5.4.3 (compact-support Kunneth); localization5.1.16.3.

SGA4, ExposeXVIII. Equation2.8.1 (affine-space top trace), Theorem2.9 (relative trace, base change and normalization); Proposition1.1.6 and equation1.1.3.2 for the projective-line calculation.

NIST Digital Library of Mathematical Functions. Equation25.4.2 (zeta functional equation); equation25.2.11 (Euler product on Re(s)>1).

NN-C006. exact-count and parity complexity implications, including Toda and Valiant--Vazirani locators.

\begin{proof}
First recall the parameter content of the earlier finite Boolean construction, without counting it as a new finding. Compile the CNF into an AND/NOT verifier circuit with $O(n+km)$ gates. Let $x_1,\ldots,x_v$ be all input and gate coordinates. Impose their Boolean equations $x_j^2-x_j=0$, the quadratic gate equations $z-ab=0$ and $z+a-1=0$, and the acceptance equation $z_{\rm out}-1=0$. These are $h=O(v)$ polynomials $f_1,\ldots,f_h\in\mathbb Z[x]$ of degree at most two and constant monomial sparsity. The Boolean algebra splits as a product of copies of $\mathbb Z$ indexed by bit assignments. Each gate or acceptance equation evaluates to zero or to a unit $\pm1$ in each factor. Consequently
\[
        Y=V(f_1,\ldots,f_h)\subset\mathbb A^v_\mathbb Z
           \simeq\coprod_{N}\operatorname{Spec}\mathbb Z.
\]
There is one extension of each accepting input to the gate coordinates. Retain a Boolean polynomial, say $f_1=x_1^2-x_1$, which is nonzero over every field.

Now introduce $y_1,\ldots,y_h,t$ and define
\[
 X=X_\Phi=V\left(t^2-t-\sum_{i=1}^{h}y_i f_i(x)\right)
                  \subset\mathbb A^{v+h+1}_\mathbb Z.
\]
Its single equation has degree at most three, $O(v+h)$ monomials and constant-size integer coefficients. Denote its left side by $G$. In its coordinate ring,
\[
 1=(2t-1)^2-4\sum_i y_i f_i
   =(2t-1)\frac{\partial G}{\partial t}
       +4\sum_i y_i\frac{\partial G}{\partial y_i}.
\]
Thus the relative Jacobian ideal contains one over $\mathbb Z$ itself. The relative hypersurface criterion proves that the projection $\pi:X\to\mathbb A^v_\mathbb Z$ is smooth of relative dimension $h$, and hence $X\to\operatorname{Spec}\mathbb Z$ is smooth of relative dimension $v+h$. This argument treats characteristic two and odd characteristics uniformly.

To prove geometric integrality, work over an algebraically closed field $K$ of any characteristic and view $G$ as a polynomial in $y_1$ over
$R=K[x,t,y_2,\ldots,y_h]$. Its coefficient is $-f_1\ne0$. The other coefficient is $t^2-t-\sum_{i>1}y_i f_i$, a monic polynomial in $t$. It has no common irreducible factor with $f_1$: every such factor of $f_1$ depends only on $x$, whereas the leading coefficient in $t$ is one. Hence $G$ is primitive over the unique-factorization domain $R$. It is linear and irreducible over $\operatorname{Frac}(R)$, so Gauss's lemma makes it irreducible, and therefore prime, in $R[y_1]$. Every geometric fiber is a domain. The monic equation in $t$ also makes the original coordinate algebra free of rank two over $\mathbb Z[x,y]$.

The count and class have an elementary proof valid over every base. Let $U_i$ be the locally closed stratum where $f_1=\cdots=f_{i-1}=0$ and $f_i$ is invertible. On $U_i$, solve the equation for $y_i$. The other $h-1$ coordinates $y_j$ and $t$ are free, giving
\[
             \pi^{-1}(U_i)\simeq U_i\times\mathbb A^h.
\]
On $Y$, the equation is $t(t-1)=0$; the two factors are comaximal over every ring. Thus
\[
             \pi^{-1}(Y)\simeq
                   (Y\times\mathbb A^h)\coprod(Y\times\mathbb A^h).
\]
Here $[\ ]$ denotes the Grothendieck ring with the usual closed-open scissor relation and $\mathbb L=[\mathbb A^1_\mathbb Z]$. Summing the displayed strata gives
\[
 [X]=\mathbb L^h([\mathbb A^v]-[Y])+2\mathbb L^h[Y]
     =\mathbb L^{v+h}+\mathbb L^h[Y]
     =\mathbb L^{v+h}+N\mathbb L^h.
\]
The same stratification proves the point count, with $N$ rational points of $Y$ over every finite field. Equivalently, for a base point with some $f_i(x)\ne0$, every choice of $t$ and the other $h-1$ coordinates gives one point, so its fiber has $q^h$ points. At a point of $Y$, the fiber has $2q^h$ points. No cancellation estimate or character-sum bound is needed.

We next justify the stronger actual-cohomology assertion. This step is separate from the Grothendieck-class calculation. Fix any base field $K$ and a prime $\ell$ invertible in $K$, with rational coefficients $\mathbb Q_\ell$. Work on the geometric base for cohomology and retain the descent data from $K$. The trace, stalk and projector argument does not require a finite field; Frobenius will be used only for the subsequent finite-field interpretation. The standard affine-space calculation is
\[
 R\Gamma_c(\mathbb A^h_{\overline K},\mathbb Q_\ell)
                   =\mathbb Q_\ell(-h)[-2h].
\]
For $h=1$ this follows from Kummer theory and $\operatorname{Pic}(\mathbb P^1)=\mathbb Z$, the projective-line top trace, and the localization sequence at infinity; powers follow from compact-support Kunneth. Exact source interfaces are SGA4 XVIII1.1.3.2 and Proposition1.1.6, XVII5.1.16.3, XVII5.4.3, and XVIII2.8.1. The !-base-change theorem XVII5.2.6 and its geometric-stalk Proposition5.2.8 now show that $R\pi_!\mathbb Q_\ell$ is concentrated in degree $2h$. Write its sole sheaf as $\mathcal T$.

The involution $\iota:t\mapsto1-t$ acts on $\mathcal T$. Since two is invertible in $\mathbb Q_\ell$, the projectors $(1\pm\iota)/2$ split it into $\mathcal T^+\oplus\mathcal T^-$. Off $Y$ the fiber is affine $h$-space, and $\iota$ is an affine automorphism in the displayed coordinates; it acts trivially on its top compact cohomology. On $Y$ it exchanges the two affine-space components. The relative trace of SGA4 XVIII2.9 gives
\[
                  \mathcal T(h)\longrightarrow\mathbb Q_\ell.
\]
On the plus part it is an isomorphism on every geometric stalk: off $Y$ it is the affine-space trace, and on $Y$ it is the sum of the two component traces, restricted to the diagonal line. The latter sends $(a,a)$ to $2a$; its normalized basis is half the sum of the two component classes. This remains invertible for $\ell=2$ because the coefficient field is $\mathbb Q_2$, not $\mathbb Z_2$. A morphism of etale sheaves that is an isomorphism on every geometric stalk is an isomorphism, so this proves a global identification of the plus sheaf, not just equality of ranks.

The minus part vanishes off $Y$. Base change to $Y$ identifies it globally with the anti-diagonal in the two labeled components $t=0$ and $t=1$, namely $\mathbb Q_\ell(-h)$. If $i:Y\hookrightarrow\mathbb A^v$ is the immersion, the canonical localization map $\mathcal T^-\to i_*i^*\mathcal T^-$ is an isomorphism on every geometric stalk, and hence an isomorphism of sheaves. It follows that $\mathcal T^-=i_*\mathbb Q_\ell(-h)$. All maps are defined over $K$ and are Galois-equivariant; when $K$ is finite they commute with Frobenius. Since the original derived image has a single cohomology sheaf, it is canonically that sheaf shifted. We have proved
\[
 R\pi_!\mathbb Q_\ell
  \simeq\bigl(\mathbb Q_\ell\oplus i_*\mathbb Q_\ell\bigr)(-h)[-2h].
\]
Taking compactly supported cohomology on $\mathbb A^v$ gives exactly the two claimed groups, because $Y$ is $N$ split points. Their degrees are distinct because $v\geq1$. Involution acts as minus one on $H_c^{2h}$ and plus one on the top group. These are actual Tate representations, pure of the displayed distinct weights when the base field is finite; no inference from an arbitrary positive polynomial point count is used.

The local zeta function at a prime $p$ follows from the all-extension-field count:
\[
 Z(X_{\mathbb F_p},T)
       =\frac{1}{(1-p^{v+h}T)(1-p^hT)^N}.
\]
Taking the Euler product, initially for $\Re s>v+h+1$, gives the stated arithmetic zeta identity, and the right side gives its meromorphic continuation. A's RH theorem places the nontrivial zeros of the two Riemann-zeta factors on the shifted lines $\Re s=v+h+1/2$ and $\Re s=h+1/2$, with the second factor present only when $N>0$. For any fixed nonreal nontrivial zero $\rho$ of $\zeta$, use only $0<\Re\rho<1$ and $\Im\rho\ne0$. At $s=h+\rho$ the other factor is $\zeta(\rho-v)\ne0$: its argument has negative real part and nonzero imaginary part, and the functional equation together with the zero-free Euler product on $\Re s>1$ excludes a zero there. Hence the displayed order formula holds without A; A is needed here only to specify the shifted critical lines. No simplicity assumption is required, and no algorithm for measuring this analytic order from a cubic equation is supplied by RH.

The quadratic-character variant uses the same $f_i,v,h,N$, over $B=\mathbb Z[1/2]$, and has equation $t^2+1=\sum_i y_i f_i(x)$. Let $G^-=t^2+1-\sum_i y_i f_i$. In its coordinate ring over $B$,
\[
 1=-\frac t2\frac{\partial G^-}{\partial t}
                     -\sum_i y_i\frac{\partial G^-}{\partial y_i}.
\]
Thus the projection $\pi^-:X^-\to\mathbb A^v_B$ is smooth of relative dimension $h$, by the relative hypersurface criterion already sourced above. For geometric integrality, view $G^-$ as a polynomial linear in a coordinate $y_i$ whose coefficient $f_i$ is nonzero. That coefficient has no common irreducible factor with the remaining coefficient, which is monic in $t$. The primitive-linear-polynomial argument above applies in every residue characteristic of $B$.

Off $Y=V(f_1,\ldots,f_h)$, solving for a coordinate $y_i$ again gives affine $h$-space fibers and the same locally closed affine-bundle stratification. Over $Y$ the scheme is
\[
       Y\times_B\operatorname{Spec}B[i]\times_B\mathbb A^h_B.
\]
This is the actual quadratic etale cover with its two conjugate geometric points. For $p$ odd, the number of roots of $t^2+1$ in $\mathbb F_{p^a}$ is $1+\chi_{-4}(p)^a$: a square root of minus one is present precisely when $p^a\equiv1\pmod4$. Since $Y$ has $N$ rational points over every finite field, adding the rational fibers proves
\[
 \#X^-(\mathbb F_{p^a})=p^{a(v+h)}+N\chi_{-4}(p)^a p^{ah}.
\] The corresponding scissor identity is
\[
 [X^-]=\mathbb L^{v+h}
       +N\bigl([\operatorname{Spec}B[i]]-[\operatorname{Spec}B]\bigr)\mathbb L^h
\]
over $B$. Its difference of classes is the quadratic sign class; the entire scheme is nevertheless an honest smooth integral geometric object.

The full compact-support realization is not inferred solely from this class. Use the cohomological projector proof above for the involution $t\mapsto-t$. Over a geometric base point off $Y$, its action on the top compact cohomology of affine $h$-space is the identity. Over $Y$, it interchanges the two components indexed by $i$ and $-i$. Consequently the plus part of $R^{2h}\pi^-_!\mathbb Q_\ell$ is the constant Tate sheaf, by the global trace normalization. The minus part is supported on $Y$ and is the anti-invariant summand of the rank-two permutation sheaf of the quadratic cover, tensored with $\mathbb Q_\ell(-h)$. This is the sheaf of the character $\chi_{-4}$. Geometric-stalk base change, the global relative trace, and the canonical localization map identify these summands as sheaves, exactly as above; there is no unresolved extension between them because the idempotent projectors split it. All other direct-image sheaves vanish because the geometric fibers are unions of affine $h$-spaces.

Taking compact supports on $\mathbb A^v$ therefore gives
\[
 H_c^{2h}(X^-_{\overline{\mathbb F}_p},\mathbb Q_\ell)
       =\bigl(\chi_{-4}\otimes\mathbb Q_\ell(-h)\bigr)^N,
 \qquad
 H_c^{2(v+h)}=\mathbb Q_\ell(-v-h).
\]
Here the character notation denotes the unramified local Frobenius representation at $p$; its eigenvalue is $\chi_{-4}(p)$. Thus the lower group has Frobenius eigenvalues $\chi_{-4}(p)p^h$, of complex absolute value $p^h$, and is pure of weight $2h$. The upper group is pure of weight $2(v+h)$. The degrees differ because $v\geq1$. This explicit calculation uses only classical etale trace and base-change theorems; A is used afterward for the global analytic zero locations.

The local zeta factor is
\[
 \frac{1}{(1-p^{v+h}T)(1-\chi_{-4}(p)p^hT)^N}.
\]
Multiplication over odd primes gives the stated arithmetic zeta identity in its half-plane of absolute convergence and then meromorphically. The omitted prime-two factor of the Tate term is exactly the factor in $\zeta^{(2)}$. Its extra zeros are explicit Euler-factor zeros, not nontrivial zeros of a Dirichlet L-function.

For any nonreal nontrivial zero $\rho$ of $L(s,\chi_{-4})$, write $\sigma=\Re\rho$ and use only $0<\sigma<1$. The Tate factor at $s=h+\rho$ is $\zeta^{(2)}(\rho-v)\ne0$. Indeed $\zeta(\rho-v)\ne0$ by the functional-equation argument above, and the modulus of $2^{-(\rho-v)}$ is $2^{v-\sigma}>1$, so the removed Euler factor is also nonzero. Hence
\[
 \operatorname{ord}_{s=h+\rho}\zeta_{X^-_\Phi}(s)
              =N\operatorname{ord}_{s=\rho}L(s,\chi_{-4}).
\]
This order identity also needs no A; A identifies the shifted critical line. No assumption that the chosen zero is simple is made. Locating its line does not evaluate its multiplicity in this constructed zeta function.

At the fixed prime $p=3$, where $\chi_{-4}(3)=-1$, one obtains
\[
       N=\frac{3^{v+h}-\#X^-_\Phi(\mathbb F_3)}{3^h}.
\]
The point-count function is in $\#\mathrm P$: guess the finitely many ternary coordinates using binary codes, reject unused codes, and check the one cubic equation. It is counting-hard by the displayed numerical reduction. Its parity equals $1+N\bmod2$, so even parity counting over this fixed field is $\oplus\mathrm P$-complete on the promised class. Rational-point existence is always true, because $N\leq2^v<3^v$ gives $\#X^-(\mathbb F_3)>0$. The hard Boolean question is whether the sign-Tate summand has nonzero rank, not whether the entire smooth scheme has a rational point.

E evaluates $N$ from the original CNF with its original $n$-variable bound and polynomial bit overhead for the displayed arithmetic. A makes the nontrivial zero-location statements unconditional under the user's designated input convention; it does not compute $N$. M is unused because the only nontrivial arithmetic character is the primitive quadratic Dirichlet character modulo four.

Finally consider effectivity. A single point count at $q=2$ recovers
\[
          N=\frac{\#X(\mathbb F_2)-2^{v+h}}{2^h}.
\]
Counting those points is in $\#\mathrm P$ by guessing the $v+h+1$ bit coordinates, and is $\#\mathrm P$-hard by this numerical reduction. Since $v\geq1$, its residue modulo $2^{h+1}$ is $2^h(N\bmod2)$, whereas every lower binary digit vanishes. At any odd prime power $q$, instead, $\#X(\mathbb F_q)\equiv1+N\pmod2$. In particular parity of rational-point counting is $\oplus\mathrm P$-complete already for these smooth geometrically integral sparse cubic hypersurfaces over the fixed field $\mathbb F_3$. The multiplicity $N=\dim H_c^{2h}$ is a $\#\mathrm P$ function on the encoded class; its nonvanishing is NP-complete and its parity is $\oplus\mathrm P$-complete. Rational-point existence for $X$ itself is always true, because $\#X\geq2^{v+h}$.

There is also an exact fixed-modulus transfer. Fix an integer $m_0\geq2$ and a prime $p$ not dividing $m_0$. Since $p^h$ is a unit modulo $m_0$,
\[
 N\bmod m_0=(p^h)^{-1}\bigl(\#X(\mathbb F_p)-p^{v+h}\bigr)\bmod m_0.
\]
Conversely the same point-count identity computes the point count modulo $m_0$ from $N\bmod m_0$. Thus the two modular counting functions are polynomial-time equivalent on the encoded class. This is stated for the function returning the count residue; no convention-dependent identification with a class called $\operatorname{Mod}_{m_0}\mathrm P$ is needed. If instead $p$ divides the modulus, division by $p^h$ is not invertible at that precision. In particular modulo $p^r$ with $r\leq h$ the whole count vanishes, whereas knowing the count modulo $p^{h+r}$ and subtracting the top term recovers $N\bmod p^r$ after exact division by $p^h$. The support of the whole point count is always nonzero; it does not supply the separate hard test $N>0$ for the lower cohomological multiplicity.

E computes $N$ through the source formula in $\operatorname{poly}_k(n+m)2^{O_k(n\log\log n/\sqrt{\log n})}$ time, with the sharper exponent when $m=O(n)$. Once $N$ is known, every displayed count and compact zeta-factor description is obtained by polynomial-bit arithmetic; a count at $q$ has $O((v+h)\log q+n)$ bits. Substituting the geometric coordinate count into E would be a weaker parameter statement. The elementary smooth construction is a classical quadric/Artin--Schreier form of the Cayley counting mechanism; the new conclusion here is its explicit arithmetic, pure-sign, and original-witness effectiveness gate, not a claim to have discovered the underlying point-count identity.
\end{proof}
\paragraph{Status and complexity impact.}
Complete integral smoothness, stratification, global cohomological projector, quadratic-character and arithmetic-zeta proofs; complete deduction-level self-reading passed. Complete deduction-proof review by root: passed; imported input theorems are not claimed independently proved. The exact obstruction is now a multiplicity in a genuine pure Tate cohomology group of a smooth geometrically integral arithmetic scheme. Its degree is at most three, its defining polynomial is sparse, and its two global zeta descriptions use Riemann-zeta shifts and the quadratic Dirichlet factor covered by A. This proves that those geometric and analytic structures can coexist with a counting-complete multiplicity-extraction problem. It does not prove that polynomial-time extraction is impossible: such an algorithm would give $\mathrm{FP}=\#\mathrm P$ and $\mathrm{PH}=\mathrm P$, while parity extraction would give $\mathrm P=\oplus\mathrm P$ and $\mathrm{NP}=\mathrm{RP}$. M is unnecessary for these objects, but that fact does not remove the general Dedekind witness-field gate in Burgisser or Koiran. No GCT canonical-basis identification, positive polytope formula, or VP/VNP separation follows without an additional effective representation-theoretic construction.

\paragraph{Geometric and representation dictionary.}
Let $f_1,\ldots,f_h\in\mathbb Z[x_1,\ldots,x_v]$ be the sparse quadratic Boolean and verifier equations in the proof, so that $v,h\geq1$, $f_1=x_1^2-x_1$, and $Y=V(f_1,\ldots,f_h)\simeq\coprod_N\operatorname{Spec}\mathbb Z$. Then the sparse cubic arithmetic hypersurface
\[
 X=V\left(t^2-t-\sum_{i=1}^h y_i f_i(x)\right)
             \subset\mathbb A^{v+h+1}_\mathbb Z
\]
is smooth of relative dimension $v+h$ with geometrically integral fibers. Its class is $[X]=\mathbb L^{v+h}+N\mathbb L^h$. For every field $K$ and every invertible prime $\ell$, its actual geometric compact-support groups, with their Galois actions, are
\[
 H_c^{2h}(X_{\overline K},\mathbb Q_\ell)=\mathbb Q_\ell(-h)^N,
 \qquad
 H_c^{2(v+h)}(X_{\overline K},\mathbb Q_\ell)=\mathbb Q_\ell(-v-h),
\]
and all other groups vanish. When $K$ is finite the displayed identifications are Frobenius-equivariant; they do not require finite $K$ to hold as Galois representations. Its arithmetic zeta is
$\zeta_X(s)=\zeta(s-v-h)\zeta(s-h)^N$.
These are separate proved assertions: affine-bundle stratification proves the scissor identity, while !-base change, the global trace and the involution $t\mapsto1-t$ prove the actual sheaf decomposition. Their exact interfaces are SGA4 XVII5.2.6/5.2.8 and XVIII2.8.1/2.9. The relative Jacobian identity proves smoothness over the full integer base, using SGA1 II4.10(iii)/4.12.

Over $B=\mathbb Z[1/2]$, replace the equation by $t^2+1=\sum_i y_i f_i(x)$. The resulting $X^-$ is also smooth with geometrically integral fibers. Its lower compact-support group is $(\chi_{-4}\otimes\mathbb Q_\ell(-h))^N$, its upper group is the same Tate line, and
\[
 \zeta_{X^-}(s)=\zeta^{(2)}(s-v-h)L(s-h,\chi_{-4})^N.
\]
Here the sign is the actual anti-invariant summand of the regular permutation representation $1\oplus\chi_{-4}$ of $\operatorname{Spec}\mathbb Z[i,1/2]\to\operatorname{Spec}B$. The first construction has a separate algebraic involution; the second realizes a quadratic Galois sign. Neither asserts an injective realization of the full split-zero semiring.

The proof constructs these data from a Boolean formula with $N$ satisfying assignments, with $v,h$ and the cubic description polynomial in that formula's size. Consequently extracting the lower cohomology multiplicity $N$ from this succinct geometric input is counting-complete, its parity is parity-counting-complete, and its nonvanishing is NP-complete. The corresponding shifted nonreal nontrivial zeta-zero orders are $N$ times a fixed nonzero analytic multiplicity; this identity is unconditional, while A locates the shifted critical lines. E evaluates $N$ with its original-witness bound. Thus actual purity, a concrete sign representation, smooth integral geometry and A-controlled analytic zero locations all coexist with this precise multiplicity-extraction problem. No unconditional time lower bound or effective GCT canonical-object identification is inferred.



\section{NN-C028 --- Embedding comparison for actual pure Tate cohomology is PP-complete}
\begin{proposition}[NN-C028]
There is a polynomial-time construction from a pair of Boolean verification circuits to two smooth geometrically integral affine schemes $X_1,X_2/\mathbb Z$, each defined by one sparse cubic equation. They have equal relative dimension $V+H$, and
\[
 H_c^{2H}(X_{j,\overline{\mathbb F}_p},\mathbb Q_\ell)
          =\mathbb Q_\ell(-H)^{N_j},\qquad
 H_c^{2(V+H)}=\mathbb Q_\ell(-V-H),
\]
with all other groups zero and $N_j$ the corresponding acceptance count. For any fixed prime $p$ and any fixed $\ell\ne p$, deciding whether a Frobenius-equivariant linear injection
\[
 H_c^{2H}(X_1,\mathbb Q_\ell)
       \hookrightarrow H_c^{2H}(X_2,\mathbb Q_\ell)
\]
exists is PP-complete on this encoded family. Hardness persists when $N_1,N_2>0$, so the Frobenius eigenvalue support, the cohomological degree support and the two weights agree. The same comparison holds simultaneously at every prime, and for the literal quadratic-character variants over $\mathbb Z[1/2]$ at every odd prime.

This is existence of a linear map between actual geometric cohomology representations. It is not a claim that the map is induced by a morphism, correspondence, or containment between the schemes. Under A the global zeta factors have the same shifted critical lines; their multiplicity comparison is still the PP-complete operation.
\end{proposition}
\paragraph{Exact inputs.}
NN-C027. smooth integral sparse cubic construction and actual compact-support Tate decomposition; both arithmetic zeta identities.

NN-G015. work/gct/chapter.tex subsectionNN-G015; support-versus-count comparison is background, not re-registered.

NN-C006. Toda main-theorem locator and resulting PH collapse.

Aaronson, Quantum Computing, Postselection, and Probabilistic Polynomial-Time (2005). Author manuscript Section1, printedp1, definition of PP by strict-majority acceptance; the reduction below is proved directly.

SGA4, ExposesXVII--XVIII. XVII Theorem5.2.6 and Proposition5.2.8; XVIII equation2.8.1 and Theorem2.9; used through the complete global sheaf proof in NN-C027.

E: The Exponential Time Hypothesis is False, supplied LaTeX. Theorem1.1, original-CNF-variable count and bounded-density specialization.

\begin{proof}
Start with two verification circuits. The gate construction of NN-C027 gives systems $f_{j,1},\ldots,f_{j,h_j}$ on $v_j$ Boolean input and gate coordinates, with
\[
 Y_j=V(f_{j,1},\ldots,f_{j,h_j})\simeq
             \coprod_{N_j}\operatorname{Spec}\mathbb Z.
\]
Each gate extension is unique. The equations have degree at most two and constant monomial sparsity, and each system retains $x_1^2-x_1$ among its equations. A fixed dummy coordinate constrained to zero can be introduced first if necessary to make every coordinate list nonempty.

The geometric parameters can be matched without multiplying either count. Put $V=\max(v_1,v_2)$. For each missing coordinate append a new variable $z$ and the two equations $z^2-z=0$, $z=0$. These force its unique value zero, so $Y_j$ and $N_j$ are unchanged. Let the resulting numbers of equations be $h'_1,h'_2$, and put $H=\max(h'_1,h'_2)$. Pad the shorter list by repeated copies of its retained Boolean equation until both lists have exactly $H$ entries. Repeating an equation does not change its common-zero scheme. All of this uses only polynomially many coordinates, equations and constant-size coefficients. Apply NN-C027 to the padded lists:
\[
 X_j:\quad t^2-t=\sum_{i=1}^H y_i f_{j,i}(x).
\]
Its proof applies unchanged to redundant equations and gives the displayed actual cohomology. In particular the lower Frobenius operator is the scalar $p^H$ on a vector space of dimension $N_j$. Every linear map between these two lower groups therefore commutes with Frobenius. Such a map can be injective exactly when
\[
                         N_1\leq N_2.
\]
This proves the assertion simultaneously at all primes; no explicitly listed exponential-size cohomology basis is required in the input. In the quadratic-character version the common scalar is $\chi_{-4}(p)p^H$, so the same equivalence holds at every odd prime. The equal top Tate lines can also be included if one asks for an injection of the full graded compact-support representation.

For hardness, take a Boolean circuit $C$ on $n\geq1$ bits, with $a$ accepting assignments. On a tagged input $(b,x)\in\{0,1\}^{n+1}$ define two polynomial-size verifiers. The right verifier accepts when $b=0$ and $C(x)=1$, or when $b=1$ and $x=0^n$. Its count is $N_R=a+1$. The left verifier accepts when $b=0$ and $x_1=0$, or when $b=1$ and $x$ is one of the two distinct strings $0^n,e_1$. Its count is
\[
             N_L=2^{n-1}+2.
\]
Both counts are positive. The injection from the left lower group to the right exists precisely when
\[
  2^{n-1}+2\leq a+1
       \quad\Longleftrightarrow\quad a>2^{n-1}.
\]
Strict-majority circuit acceptance is PP-complete directly from the definition: unroll a polynomial-time probabilistic verifier on its fixed input, using its independent fair coin bits as circuit inputs, and pad the coin count if needed. The circuit accepts on exactly the machine's accepting coin strings. Thus the entire reduction to the two sparse cubics is polynomial and proves PP-hardness, including the positive-rank restriction. An input circuit with no coin bits can be given an ignored bit without changing its majority status.

For completeness, here is a direct PP-membership construction that handles equality rather than silently replacing $\leq$ by $<$. Let $P,Q$ be the two acceptance predicates, with counts $A,B$, padded to a common witness length $w$ by requiring every added bit to be zero. Add a new bit $r$ and define
$P'(r,x)=[r=0]P(x)$, $Q'(r,x)=[r=0]Q(x)$.
Their common domain has size $M=2^{w+1}$, while their counts remain $A,B$. On a selector bit $b$ and this domain, accept when
\[
          (b=0\ \hbox{and }Q'=1)
            \quad\hbox{or}\quad(b=1\ \hbox{and }P'=0).
\]
There are $M+B-A$ accepting strings among $2M$. The string $b=0,r=1,x=0^w$ is always rejecting. Add a fresh free bit $u$, doubling all paths, and override exactly the one path $u=0,b=0,r=1,x=0^w$ to accept. The number of accepting paths becomes
\[
                   2M+2(B-A)+1
\]
among $4M$. This exceeds half exactly when $B-A\geq0$. The whole test is a polynomial-size Boolean predicate. It proves PP membership on the encoded family, and hence PP-completeness. The circuit-and-equation format can be checked syntactically; the assertion is not a promised PP upper bound for arbitrary sparse cubic schemes whose provenance is absent.

The associated arithmetic zeta functions are
\[
             \zeta_{X_j}(s)
                 =\zeta(s-V-H)\zeta(s-H)^{N_j}.
\]
When $N_1,N_2>0$ the types of nontrivial factors, cohomological degrees and eigenvalue supports agree. At any fixed nonreal nontrivial zero $\rho$ of $\zeta$, NN-C027 gives
\[
 \operatorname{ord}_{s=H+\rho}\zeta_{X_j}(s)
                 =N_j\operatorname{ord}_\rho\zeta.
\]
Consequently the same embedding comparison is comparison of those two analytic orders. The order identity is unconditional; A fixes the common shifted critical lines. For the quadratic variant replace the lower factor by $L(s-H,\chi_{-4})$ and remove the upper prime-two Euler factor as in NN-C027. The same conclusion holds at its nonreal nontrivial zeros.

E decides the comparison by computing the two original counts and comparing their binary integers. For fixed-width source CNFs it keeps the original variable parameters, regardless of the padding of the geometric schemes. For a general verification circuit of size $s$ on $n$ inputs, the parsimonious gate-to-3-CNF encoding has $O(n+s)$ variables and a linear number of clauses, so Theorem1.1's bounded-density form gives a bound $\operatorname{poly}(n+s)2^{O((n+s)/\sqrt{\log(n+s)})}$ for that source circuit. No bound solely in the original coin count $n$ is inferred for arbitrary circuits with many gates.
\end{proof}
\paragraph{Status and complexity impact.}
Complete actual-cohomology comparison, parameter matching and strict-majority proofs; complete deduction-level self-reading passed. Complete deduction-proof review by root: passed; imported input theorems are not claimed independently proved. This turns NN-G015's abstract support-versus-dominance example into a comparison of actual pure geometric representations of smooth integral cubic schemes. The representation type and analytic critical lines are fixed; deciding which multiplicity is larger is PP-complete. A polynomial-time injection-existence procedure on this succinct family would give $\mathrm P=\mathrm{PP}$, and hence $\mathrm{PH}=\mathrm P$ by Toda. This is a conditional implication, not an unconditional lower bound. E supplies a subexponential evaluator of the comparison, and A does not evaluate the multiplicities. The result does not transfer PP-hardness to the permanent/determinant orbit-closure multiplicities or establish a GCT multiplicity obstruction.

\paragraph{Geometric and representation dictionary.}
The algebraic-geometric assertion is that, for the explicitly padded Boolean systems $Y_j\simeq\coprod_{N_j}\operatorname{Spec}\mathbb Z$, the two cubic schemes
\[
 X_j=V\left(t^2-t-\sum_{i=1}^H y_i f_{j,i}(x)\right)
          \subset\mathbb A^{V+H+1}_\mathbb Z
\]
are smooth with geometrically integral fibers, and their actual lower compact-support realizations are $\mathbb Q_\ell(-H)^{N_j}$. In the category of rational l-adic Frobenius representations, an injection between these realizations exists if and only if $N_1\leq N_2$. This equivalence holds for every prime with identical geometric parameters, and for the anti-invariant quadratic-character realizations at every odd prime. The scheme construction, relative trace and projectors are those proved in NN-C027, with exact SGA4 XVII5.2.6/5.2.8 and XVIII2.8.1/2.9 interfaces. Parsimonious padding proves that the geometric dictionary preserves the relevant counts and common Tate twist.

Thus even on this actual pure Tate subcategory, a succinct defining-equation presentation can make the existence of a subrepresentation PP-complete. The explicit matrices, once supplied, would be scalar and their rank comparison elementary; obtaining their multiplicities from the equations is the hard operation. Their zeta factors have A-controlled shifted critical lines, but the orders at a fixed nonreal lower-factor zero give precisely the same comparison. This is a theorem about linear injections of cohomology realizations, not about scheme morphisms, algebraic correspondences or orbit-closure containment. No fullness statement for a realization functor is assumed.



\section{NN-C029 --- E constructs subexponential ordinary ideal certificates and affine trivializations}
\begin{proposition}[NN-C029]
Fix a density constant $C$. Let $f_1,\ldots,f_h\in\mathbb Z[x_1,\ldots,x_v]$, $h\leq Cv$, be the Boolean, AND/NOT-gate and acceptance equations of NN-C027. In particular all $x_j^2-x_j$ are retained, each equation has degree at most two, scope at most three and constant coefficient size, and its Boolean values belong to $\{0,1,-1\}$. If the system is unsatisfiable, E deterministically constructs division-free arithmetic circuits over $\mathbb Q$ for integral polynomials $a_i(x)$ satisfying the ordinary polynomial identity
\[
                  \sum_{i=1}^h a_i(x)f_i(x)=1,
          \qquad \deg a_i\leq v.
\]
The total circuit size and construction time are
\[
 S(v)=\operatorname{poly}(v)\,
 2^{O_C(v(\log\log v)^2/\sqrt{\log v})},
\]
and every rational constant has polynomial bit length. Thus $\sum_i a_i(x)z_i$ is a Hilbert-like IPS certificate with a constructible $S(v)$-size proof over $\mathbb Q$. Its coefficient polynomials are integral; the circuit representation is allowed rational constants.

For fixed-width $k$-CNFs with $m\leq Cn$ clauses there is a stronger specialization on the original $n$ input variables: writing $g_i$ for the product of the clause's falsifying literals, E constructs integral $A_i$ with $\sum_i A_i g_i=1$, without Boolean axioms, in size and time $S(n)$, with constants also depending on $k$. The first-violated-clause identity itself is classical Grochow--Pitassi Lemma4.2; the deduction here is the E-based circuit bound and its ordinary-ideal extension to the quadratic verifier presentation.

In the empty case, E also constructs $S(v)$-size rational arithmetic circuit descriptions of mutually inverse integral polynomial maps
\[
 X=\left\{t^2-t=\sum_i y_i f_i(x)\right\}
       \simeq\mathbb A^{v+h}_{\mathbb Z}.
\]
Their degrees are at most $v+2$ and $v+6$. For a verifier on $n$ input bits with $s$ gates, the parameter is $v=n+s$, up to harmless constant additions. No original-$n$ bound for arbitrary long verifier circuits is asserted.

There is an explicit large-characteristic transfer. Put
$B(v)=\max(2,\lceil\log_2 v\rceil)$. The multiplier circuits can be chosen with constants in $\mathbb Z[1/B(v)!]$. Hence they give circuits of the same size over every field of characteristic $p>B(v)$, with construction bit overhead polynomial in $\log p$. The same statement holds for C016's universal partition polynomials and for the sparse-CNF specialization with $v$ replaced by $n$. This original localization argument does not cover fixed small characteristic as the input size tends to infinity. NN-C033 subsequently closes that limitation by a modified integral-swap-matrix and extension-node construction; NN-C034 then constructs integral circuits over Z itself.
\end{proposition}
\paragraph{Exact inputs.}
NN-C016. complete coefficient-independent symbolic partition construction, fixed arity and linear factor density.

E: The Exponential Time Hypothesis is False, supplied LaTeX. Proposition3.4 over every commutative Q-algebra; Lemma4.6 replacement-row denominator (R-1)!Q!; Lemmas4.7--4.8 and equation parameter-QR; AppendixB Hermite matrix in subsection rational-readout; coefficient-independent measurement separation.

Grochow--Pitassi, Circuit Complexity, Proof Complexity, and Polynomial Identity Testing: The Ideal Proof System. Author JACM manuscript, Definition1.1 and Hilbert-like terminology, printed draftp1:4; Lemma4.2 and complete proof, printed draftpp1:26--1:28, equations(4)--(6).

EGA IV, fourth part, Publications Mathematiques IHES32 (1967). Printed p.111,18.2.1--18.2.2; base-change equation18.2.1.2; primary PDF page110 (zero-based109).

NN-S016. ordinary integral Bezout identity and explicit affine-space coordinate formulas; work/complexity/c027\_degree\_addendum.tex and work/root/chapter8.tex.

NN-C027. quadratic AND/NOT verifier equations with all Boolean axioms and the associated smooth integral cubic.

\begin{proof}
Bounded values of the parameters can be handled by direct expansion. We first prove the ordinary-ideal assertion for the quadratic verifier list. For a bit vector $b$ put
\[
 e_b(x)=\prod_{j:b_j=1}x_j\prod_{j:b_j=0}(1-x_j),
 \qquad
 a_i^{(0)}(x)=\sum_{b\in\{0,1\}^v}e_b(x)f_i(b)
                  \prod_{j<i}(1-f_j(b)^2).
\]
The factors $1-f_j(b)^2$ are exactly the indicators of $f_j(b)=0$. At any bit vector $b$, unsatisfiability guarantees a first index $i$ with $f_i(b)\ne0$. Its coefficient in the displayed sum is $f_i(b)\in\{1,-1\}$, and all other first-violation coefficients vanish. Therefore
\[
                    \sum_i a_i^{(0)}(b)f_i(b)=1.
\]
Each $a_i^{(0)}$ is an integral multilinear polynomial of degree at most $v$. The Boolean equations themselves have zero values and hence zero preliminary multipliers.

The preliminary multipliers are symbolic bounded-arity partition functions. Their summation variables are the $v$ bits $b$. There are at most $h$ prefix factors with tables $1-f_j(b)^2$, one local factor with table $f_i(b)$, and $v$ unary factors with respective weights $x_l$ at $b_l=1$ and $1-x_l$ at $b_l=0$. Every scope has size at most three and there are $O(v)$ factors and incidences. The fact that the unreduced expression $1-f_j^2$ can have degree four is harmless: its Boolean table, interpolated on at most three variables, has degree at most three. The coefficient algebra is $\mathbb Q[x_1,\ldots,x_v]$, with the $x_l$ treated as formal coefficients rather than extra summation variables.

Apply NN-C016 separately to these $h$ partition functions and specialize its independent factor weights to the indicated constants or linear polynomials. These are constant-size substitutions, so they preserve the bound $S(v)$. The factor count, rather than an unproved bound on each variable's incidence, supplies the hypothesis: the total number of incidences is $O(v)$, and C016 explicitly branches on high-incidence variables. Multiplying by $h=O(v)$ changes only the polynomial prefactor. Its construction makes no identity tests on the coefficient algebra and introduces no division by a polynomial in the $x_l$. Its only inversions are of explicit rational interpolation constants, whose bit lengths are polynomial. This yields circuits for every $a_i^{(0)}$ within the claimed time and size.

An ordinary polynomial identity still requires Boolean-ideal correction; evaluating correctly on the cube alone is insufficient. Here locality makes that correction cheap without globally expanding the circuit. Let $I_i$ be the scope of $f_i$ and $t_i=|I_i|\leq3$. Since $a_i^{(0)}$ is multilinear, its local interpolation is
\[
 a_i^{(0)}(x)=\sum_{\sigma\in\{0,1\}^{I_i}}
        A_{i,\sigma}(x_{\overline I_i})e_\sigma(x_{I_i}),
 \qquad A_{i,\sigma}=a_i^{(0)}\big|_{x_{I_i}=\sigma}.
\]
At most eight circuit specializations compute these polynomials, and their degrees are at most $v-t_i$. In the constant number of local variables the polynomial
\[
       e_\sigma f_i-f_i(\sigma)e_\sigma
\]
has zero multilinear remainder. Its degree is at most $t_i+2$. Explicit integral monomial division therefore gives
\[
 e_\sigma f_i-f_i(\sigma)e_\sigma
      =\sum_{j\in I_i}r_{i,\sigma,j}(x_{I_i})(x_j^2-x_j),
      \qquad \deg r_{i,\sigma,j}\leq t_i.
\]
For completeness, this division uses only
\[
 x^\alpha-x^{\epsilon}=\sum_j
  \left(\prod_{l<j}x_l^{\epsilon_l}\right)
  (x_j^{\alpha_j}-x_j^{\epsilon_j})
  \left(\prod_{l>j}x_l^{\alpha_l}\right),
 \quad \epsilon_l=\min(\alpha_l,1),
\]
and $x_j^a-x_j=(x_j^2-x_j)\sum_{r=0}^{a-2}x_j^r$ for $a\geq2$.
There are constantly many local monomials and coefficients of bounded size. Thus all $r_{i,\sigma,j}$ have constant-size integral circuits. Empty scopes give no correction.

Write $\mathrm{ML}$ for the unique Boolean multilinear remainder and put
\[
 Q_{i,j}=\sum_{\sigma\in\{0,1\}^{I_i}}
                       A_{i,\sigma}r_{i,\sigma,j}.
\]
The local identities imply
\[
 a_i^{(0)}f_i-\mathrm{ML}(a_i^{(0)}f_i)
              =\sum_{j\in I_i}Q_{i,j}(x)(x_j^2-x_j).
\]
The asserted remainder really is multilinear: its local factor is
$\sum_\sigma f_i(\sigma)e_\sigma$, with the corresponding coefficients $A_{i,\sigma}$ involving only the other variables. Its values are those of $a_i^{(0)}f_i$. Each $Q_{i,j}$ has degree at most $(v-t_i)+t_i=v$. Summing over $i$, the sum of the remainders is the unique multilinear polynomial whose Boolean values are all one, hence is identically one over $\mathbb Z$. Consequently
\[
 1=\sum_i a_i^{(0)}f_i-
             \sum_{j=1}^v\left(\sum_i Q_{i,j}\right)(x_j^2-x_j).
\]
Subtract the indicated $Q$ sum from the multiplier of each retained Boolean equation. This produces the promised integral $a_i$ and an ordinary identity in $\mathbb Z[x]$. Constantly many specializations and local products per equation, followed by $O(v)$ additions, keep the total circuit size and construction time within $S(v)$. In particular, no exponential multilinear expansion is being silently treated as a polynomial-time operation. The preliminary coefficient norm is bounded by $3^v$ in total, and the local correction introduces only a fixed factor per equation, so the resulting integer coefficients have $O(v+\log(v+1))$ bits; integrality is a property of the output polynomials, not of every rational circuit gate.

This construction can be made unconditional as an algorithm that either reports satisfiability or produces the certificate. Each local relation $f_i(b)=0$ is a constant-scope Boolean table, convertible to at most eight forbidden-assignment clauses. The resulting fixed-width CNF has the same $v$ variables and $O(v)$ clauses. E first decides it within its bounded-density running time, which is absorbed by $S(v)$. When it is satisfiable the algorithm reports that outcome; when it is not, the displayed construction applies. The circuit $C(x,z)=\sum_i a_i(x)z_i$ satisfies $C(x,0)=0$ and $C(x,f(x))=1$, exactly Definition1.1 of IPS, and is linear in the placeholders. Its actual polynomial degree is at most $v+1$. If a syntactic variable-degree bound is desired, retain homogeneous components $0,\ldots,v$ at every gate of each multiplier circuit: addition is componentwise and multiplication uses the truncated convolution. There is no variable division, so discarded higher components cannot contribute to a later lower component. Since the output degree is at most $v$, their sum is exactly the same multiplier. This uses $O(v^2 S(v))$ gates and no identity test, with syntactic degree at most $v$ when rational constants have variable-degree zero. It is not the different constant-free formal-degree convention in which constant leaves have degree one.

For the stronger sparse-CNF specialization, delete tautological clauses and repeated literals. An empty clause has the immediate certificate $1\cdot1=1$. Otherwise let $I_i$ be the distinct-variable scope of clause $i$, let $\alpha_i$ be its unique falsifying assignment, and set $g_i=e_{\alpha_i}(x_{I_i})$. Define
\[
 A_i(x_{\overline I_i})=
    \sum_{b\in\{0,1\}^{\overline I_i}}
       e_b(x_{\overline I_i})
       \prod_{j<i}\bigl(1-g_j(b,\alpha_i)\bigr).
\]
Upon multiplying by $g_i$, this is precisely the sum of the full idempotents $e_c(x)$ over assignments whose first violated clause is $i$. For an unsatisfiable CNF these sets partition the entire cube, so
\[
              \sum_i A_i g_i=\sum_c e_c
                   =\prod_{l=1}^n(x_l+(1-x_l))=1
\]
as an ordinary integral polynomial identity. In particular $\deg A_i\leq n-|I_i|$ and every product $A_i g_i$ has degree at most $n$. This is the classical identity of Grochow--Pitassi Lemma4.2, equations(4)--(6), specialized by fixing the falsifying local assignment. Its E consequence is as follows. Each $A_i$ is a partition function on $n-|I_i|$ summation bits, with at most $m+n$ factors of arity at most $k$. For fixed $k,C$ and $n>2k$, its variable count is at least $n/2$ and the factor density remains bounded. C016 therefore constructs all of them in size and time $S(n)$. Small $n$ is handled directly. No Boolean correction, auxiliary gate variable, or expansion in the original $n$ variables is required in this specialization. For dense CNFs this particular argument does not preserve the same original-$n$ bound; compiling into verifier gates changes the relevant parameter to their total number. NN-C036 subsequently supplies an original-$n$ integral circuit certificate for arbitrary-density fixed-width CNFs by slower disjoint sparsification, with its explicitly weaker exponent.

We next prove the asserted large-characteristic transfer, including a uniform denominator bound across every high-incidence branch. In a branch with $u$ remaining summation variables, E Lemma4.8 uses
\[
 \lambda=\lceil\log_2u\rceil,\quad
 \nu=\lceil\log_2\lambda\rceil,\quad
 Q=\left\lfloor\frac{\lambda}{128\nu}\right\rfloor,
\]
and every successful layout has $R<Q$. Since $u\leq v$, both $Q$ and $R$ are at most $B(v)$. The high-incidence branching introduces no new summation variables. For fixed density and arity, only finitely many input sizes fail the layout's numerical checks; direct expansion handles them using only integral constants. The branch sizes are $v-O(v/\sqrt{\log v})$, so this fallback affects only a fixed finite range and preserves the asymptotic bound. Thus the single localization
\[
                       A_v=\mathbb Z[1/B(v)!]
\]
works for every branch, preliminary multiplier and local correction. In the sparse-CNF specialization each summation dimension is at most $n$, so $A_n$ works uniformly there.

Here is why all arithmetic-circuit constants may be chosen in that localization. E Lemma4.6 gives the exact common denominator $(R-1)!Q!$ for every replacement row. Its prime divisors are at most $B(v)$, so these row coefficients lie in $A_v$. Taylor substitution and Boolean differences introduce only binomial integers and signs. Factor interpolation on the Boolean cube uses only zero, one and minus one. All subsequent sums and products therefore remain in $A_v$ until the final Hasse-Hermite interpolation.

That interpolation also needs only this localization. Let the multiplicity of each node be $J$ and put $M=QJ$. The evaluation-of-jets map is
\[
 A_v[T]\big/\!\left(\prod_{t=1}^{Q}(T-t)^J\right)
       \longrightarrow\prod_{t=1}^{Q} A_v[z]/(z^J),
 \qquad P\longmapsto(P(t+z))_{t=1}^{Q}.
\]
For two distinct nodes, their nonzero difference is an integer of absolute value less than $Q$, and is a unit in $A_v$. The corresponding linear ideals are comaximal, and so are their $J$th powers: expanding a relation $a+b=1$ to exponent $2J-1$ expresses one as a sum of an element of each power. The Chinese remainder theorem therefore makes this map an isomorphism. Monic division gives the free basis $1,T,\ldots,T^{M-1}$ on the left; the right uses the local powers $1,z,\ldots,z^{J-1}$. In these bases its matrix is exactly E's Hasse-Hermite matrix
\[
  \mathsf V_{(t,j),u}=\binom uj t^{u-j}\quad(u\geq j),
               \qquad \mathsf V_{(t,j),u}=0\quad(u<j).
\]
Its inverse thus has entries in $A_v$. Compute those final entries over $\mathbb Q$ as before and put them in the circuit as constants; intermediate elimination pivots are not circuit gates. Their polynomial bit-height bound is unchanged. This reasoning does not divide by $j!$ and imposes no bound $J<p$.

Consequently the C016 circuit is a circuit over $A_v$ for a polynomial with integer coefficients in its independent weights. Equality with that universal partition polynomial holds in $\mathbb Q[y]$, and hence in $A_v[y]$ by the injective inclusion into $\mathbb Q[y]$. For every prime $p>B(v)$ there is a homomorphism $A_v\to\mathbb F_p$. Apply it to every constant and gate: equality survives, and integer-weight specializations, local corrections and homogeneous-component constructions survive as well. The resulting output is exactly the reduction of each integral multiplier, not merely a polynomial agreeing with it on a finite set. Reducing polynomial-bit numerators and denominators and inverting their nonzero residues uses polynomial bit cost per constant in $v+\log p$. The compiled circuit uses only prime-field constants and is valid over every extension field of that characteristic. This proves the size and construction claims. No procedure choosing a prime, and no assertion for a fixed prime below the threshold, is hidden in the argument.

Finally apply the ordinary identity to the smooth cubic of NN-C027. Retain $x$ and define
\[
 z_i=y_i+a_i(x)(2t-t^2).
\]
On $X$, the scalar product $\sum_i f_i z_i$ equals
$(t^2-t)+(2t-t^2)=t$. Conversely, for free coordinates $(x,z)$ set
\[
 t=\sum_i f_i(x)z_i,
 \qquad y_i=z_i+a_i(x)(t^2-2t).
\]
Then $\sum_i f_i y_i=t+(t^2-2t)=t^2-t$, and direct substitution shows both composites are identities. These are integral polynomial maps over $\mathbb Z$. The first has degree at most $v+2$; the second has degree at most $v+6$, because $t$ has degree at most three. Shared circuit substitution adds only polynomially many gates to the multiplier circuits. This proves the effective geometric assertion. The rational syntax in this original proof describes integral polynomial maps; it is not a license to reduce a rational constant modulo a prime dividing its denominator. The later integral compilation in NN-C034 supplies a different circuit valid after arbitrary base change.
\end{proof}
\paragraph{Status and complexity impact.}
Complete construction of ordinary ideal multipliers, local correction, sparse-CNF specialization and integral affine trivialization. Complete deduction-proof review by root: passed; imported input theorems are not claimed independently proved. This closes a specific algebraic-certificate gate left by NN-C009: for bounded-density local verifier systems, E yields actual ordinary ideal multipliers with subexponential rational arithmetic circuits. It strengthens the transcript conclusion of NN-C008 and makes NN-S016's empty-case affine trivialization effective in this circuit model. It gives neither polynomial-size certificates nor a polynomial-time verifier for arbitrary supplied IPS circuits; verification still presents polynomial identity tests. It therefore implies no $\mathrm{NP}=\mathrm{coNP}$ collapse and no VP=VNP conclusion. The exponential monomial-size polynomial-calculus lower bounds in NN-C021 remain compatible with these shared circuits. The denominator audit here gives circuits in every characteristic greater than $\max(2,\lceil\log_2 v\rceil)$. NN-C033 subsequently supplies a different construction over every prime field, including constant-free circuits of polynomial formal degree. Semantic integrality alone did not justify reducing the original rational constants at denominator primes. NN-C034 subsequently closes the integral-gate question over $\mathbb Z$ by cyclotomic interpolation and supplies constant-free circuits of polynomial formal degree.

\paragraph{Geometric and representation dictionary.}
The algebraic-geometric theorem is constructive in the following precise representation. For the split Boolean coordinate algebra
\[
 B=\mathbb Z[x_1,\ldots,x_v]/(x_j^2-x_j)
                 \simeq\prod_{b\in\{0,1\}^v}\mathbb Z,
\]
let $Y=V(f_1,\ldots,f_h)$ be the closed subscheme cut out by the local verifier equations. If $Y$ is empty, E produces rational arithmetic circuits representing integral polynomials $a_i$ of degree at most $v$ whose ordinary ideal combination is one in $\mathbb Z[x]$. The first-violation construction partitions the component idempotents of $B$; the constant-scope division above lifts that partition to the ambient affine coordinate ring. It thereby constructs an integral polynomial isomorphism
\[
 V\left(t^2-t-\sum_i y_i f_i(x)\right)
        \longrightarrow\mathbb A^{v+h}_{\mathbb Z},
 \quad (x,y,t)\longmapsto(x,y+a(2t-t^2)),
\]
with the explicitly proved inverse and degree bounds. Its circuit description has the displayed $S(v)$ bound. Moreover that description is defined over the explicit open arithmetic base $\operatorname{Spec}\mathbb Z[1/B(v)!]$: the Hasse-jet Chinese-remainder isomorphism proves the interpolation denominators are invertible there. Thus the same effective presentation of the ideal identity and affine trivialization base-changes to every field of characteristic $p>B(v)$. The integral morphisms exist at all primes; only their claimed arithmetic-circuit presentation has this restricted base.

The symbolic preliminary multipliers are finite-algebra traces of products of local weights, so C016's trace construction and base change apply, with exact classical locator EGA IV4, printedp111,18.2.1--18.2.2 and equation18.2.1.2. The passage from quotient identities to ordinary ideal multipliers is proved here by local interpolation, without invoking an unproved effectivity theorem. In the sparse-CNF case, the classical first-violation partition already gives the ordinary identity without the Boolean equations. These are integral geometric statements with localized arithmetic-circuit presentations; no efficient circuit presentation at every prime, motivic equivalence theorem, or GCT representation-theoretic lower bound is inferred.



\section{NN-C030 --- Pairwise homogeneity gives a deterministic smooth finite-field parity reduction}
\begin{proposition}[NN-C030]
The randomized finite-field specialization in NN-C024 can be replaced by a deterministic polynomial-time construction on the actual incidence-rail encoding of NN-C023. Given a Boolean formula, construct a smooth projective hypersurface over an explicitly presented finite field $\mathbb F_{2^b}$, where $b=O(r^2)$ and its factored degree-$r$ equation has polynomial description length. The top coherent cohomology has its known one-dimensional basis, and its $2^b$-power Frobenius scalar is exactly the original satisfying-assignment parity $c$.

More precisely, let $r=2M+1\geq3$ and let
\[
 F=z_0P(z_{1,0},z_{1,1},\ldots,z_{M,0},z_{M,1})
\]
be nonzero over $\mathbb F_2$, with $P$ homogeneous of degree two in every displayed pair. Put
\[
 H=z_0^r+\sum_{e=1}^M(z_{e,0}^r+z_{e,1}^r),
 \quad d=\operatorname{ord}_r(2),\quad D=dr.
\]
Every parameter $a$ for which $V(H+aF)$ is singular lies in $\mathbb F_{2^D}$. Choose a power of two $b>D$ with $b\leq2D$ and an explicitly constructed field element $\theta$ of absolute degree $b$. Then $V(H+\theta F)$ is smooth. Its top absolute Frobenius has semilinear coefficient $\theta c$, while its field-linear Frobenius scalar is $c$.

Consequently scalar extraction, and equivalently point-count parity on this encoded class, is parity-counting complete under deterministic polynomial-time reductions. A polynomial-time algorithm for either would give $\mathrm P=\oplus\mathrm P$ and $\mathrm{NP}=\mathrm{RP}$. E gives the inherited subexponential evaluator. This closes the deterministic finite-field smooth-specialization gate in NN-C024; it does not construct canonical GCT objects.
\end{proposition}
\paragraph{Exact inputs.}
NN-C023. incidence-rail factorization F=z0 P, with each rail pair of degree two and parity coefficient c.

NN-C024. complete quadratic Artin--Schreier field-tower construction and semilinear-versus-q-linear Frobenius computation.

SGA1, ExposeII, Section4. Theorem4.10(iii) and Corollary4.12, the homogeneous hypersurface Jacobian criterion used through NN-C024.

Katz, Algebraic Solutions of Differential Equations, Invent.Math.18(1972),1--118. Rank-one Special Case2.3.7.17, printedp36; exact Cech calculation supplied in NN-C023.

SGA 7 II, Expose XXII. Theorem 3.1 (3.1.1); Corollary 3.2 (3.2.1).

Deligne, La conjecture de Weil I, PMIHES43(1974),273--307. Theorem1.6, printedp276, smooth projective finite-field purity.

Valiant--Vazirani, NP is as easy as detecting unique solutions, TCS 47 (1986), 85--93. Theorem 1.1; Corollaries 3.1--3.4, especially Corollary 3.2 p.91.

E: The Exponential Time Hypothesis is False, supplied LaTeX. Theorem1.1, inherited original-formula evaluator.

\begin{proof}
We first prove the singular-parameter assertion for every nonzero $P$ with the stated pairwise homogeneity; no factorization or complexity assumption is needed for this geometric lemma. Work in a fixed algebraic closure of $\mathbb F_2$. The odd integer $r$ is one in this field, so the partial derivatives of $H$ are the corresponding $(r-1)$st powers. The zero parameter gives the smooth Fermat hypersurface. Suppose $H+aF$ has a singular projective point for some $a$.

At such a point, $z_0$ cannot be zero. If it were, all rail partial derivatives of $F=z_0P$ would vanish, and the zero-gradient equations would force every rail coordinate to be zero. This is not a projective point. Normalize $z_0=1$. The $z_0$ derivative now gives
\[
                              1+aP=0.
\]
Thus $a\ne0$, $P\ne0$, and $aP=1$ at the normalized point. If both coordinates of any rail pair were zero, pairwise degree two would force $P=0$. This too is impossible.

Euler's identity in one rail pair reads
\[
 z_{e,0}\frac{\partial P}{\partial z_{e,0}}
 +z_{e,1}\frac{\partial P}{\partial z_{e,1}}=2P=0.
\]
Multiply its two zero-gradient equations by their corresponding rail coordinates and add. The displayed identity yields
\[
                         z_{e,0}^r=z_{e,1}^r.
\]
Neither coordinate can be zero, since then both would be zero. Therefore
\[
              s_e=z_{e,0}\ne0,\qquad
              t_e=z_{e,1}/z_{e,0},\qquad t_e^r=1.
\]
Since $r$ is odd, $2$ is invertible modulo $r$ and its multiplicative order $d$ exists, with $d\leq\varphi(r)\leq r-1$. Put $q_0=2^d$. As $r$ divides $q_0-1$, every $r$th root of unity lies in the field $K=\mathbb F_{q_0}$. Thus all $t_e$ lie in $K$.

Pairwise homogeneity gives an identity
\[
 P(s_1,s_1t_1,\ldots,s_M,s_Mt_M)
                 =\left(\prod_e s_e^2\right)C(t_1,\ldots,t_M),
\]
where $C(t)=P(1,t_1,\ldots,1,t_M)$ has coefficients in $\mathbb F_2$. Each polynomial $z_{e,0}\partial P/\partial z_{e,0}$ is also homogeneous of degree two in every pair. Hence its normalized value is $(\prod_j s_j^2)C_e(t)$ for a polynomial $C_e$ over $\mathbb F_2$. Since $P\ne0$, $C(t)\ne0$. Multiplying the $z_{e,0}$ zero-gradient equation by $z_{e,0}$ and using $aP=1$ shows
\[
            s_e^r=
            \frac{z_{e,0}\,\partial P/\partial z_{e,0}}{P}
                  =\frac{C_e(t)}{C(t)}\in K^*.
\]

Every such $s_e$ lies in $K'=\mathbb F_{q_0^r}$, with no Kummer extension oracle required. Indeed $s_e^{r(q_0-1)}=1$. The integer $r(q_0-1)$ divides $q_0^r-1$: after dividing by $q_0-1$, the remaining sum $1+q_0+\cdots+q_0^{r-1}$ is divisible by $r$ because $q_0\equiv1\pmod r$. Thus $s_e^{q_0^r-1}=1$. The roots of $X^{q_0^r}-X$ form the unique field with $q_0^r$ elements in the algebraic closure, so $s_e\in K'$. Also $K\subset K'$, because $d$ divides $dr$. All coordinates of the normalized singular point therefore lie in $K'$, and
\[
                         a=1/P\in K'=\mathbb F_{2^{dr}}.
\]
This proves the claimed confinement of all bad parameters. It is stronger than a bound on the number or total degree of bad parameters: their individual field of definition is contained in one explicitly bounded finite field. The proof uses the multihomogeneous rail structure; it makes no such claim for arbitrary degree-$r$ forms.

We next construct a parameter outside this bad field with polynomial bit cost. Compute $d$ by repeated doubling modulo $r$ until the residue one returns; this takes at most $r$ steps and polynomial time in $r$. Compute $D=dr$ and choose the least power of two $b>D$. Then $b\leq2D\leq2r^2$. There is no need to construct the possibly non-power-of-two-degree field $\mathbb F_{2^D}$ itself.

Use the explicit quadratic tower proved in NN-C024: beginning with $\mathbb F_2$, at each degree $j=1,2,4,\ldots,b/2$, find a basis element of trace one in the current field and adjoin a root of $Y^2+Y+u$. Such an element exists because the absolute trace is a nonzero linear functional, and the polynomial is irreducible because $\operatorname{Tr}(y^2+y)=0$. Basis evaluation, squaring and arithmetic have polynomial cost in $b$. Let $\theta$ be the final adjoined root. It lies outside the preceding field of degree $b/2$. Its absolute degree divides $b$; every proper divisor of the power of two $b$ divides $b/2$, and every element of such degree belongs to that preceding field. Hence $\theta$ has absolute degree exactly $b$. In particular $\theta\notin\mathbb F_{2^D}$, since an element of the latter has degree at most $D<b$. The polynomial $H+\theta F$ is therefore smooth by the singular-parameter lemma.

The tower and its basis are an explicit field presentation of polynomial length. If one monic irreducible polynomial is required, compute the first linear dependence of $1,\theta,\ldots,\theta^b$ by linear algebra in the tower basis. Its degree is $b$ by the preceding argument, so it is the minimal polynomial; the same computations give the change to a power basis. This remains polynomial in $b=O(r^2)$. The parameter is nonzero because it has degree $b>1$. No random choice, smoothness test, discriminant expansion, or enumeration of the bad field is used.

We verify that the encoding always meets the geometric lemma. In NN-C023 each incidence rail pair occurs once in its clause factor and once in its variable-consistency factor. Every monomial in each of those factors uses exactly one member of the pair. Thus their product $P$ has degree exactly two in that pair, and it is independent of the extra coordinate $z_0$. The factors are nonzero over $\mathbb F_2$, so $P\ne0$. There are $2M+1=r$ total coordinates, and $F=z_0P$ is homogeneous of total degree $r$. The general hard instances therefore satisfy every hypothesis above, even when their target coefficient $c$ is zero.

For the elementary exceptional source inputs that NN-C023 handled by a non-rail cubic, first compute their parity directly. An empty clause has parity zero; a formula with no declared variables has its single assignment explicitly testable. Replace such an instance by one of two fixed one-variable formulas: the positive unit clause $x$, with parity one, or the two contradictory unit clauses $x$ and $\neg x$, with parity zero. Both have nonempty incidence sets and the same rail construction just proved. This replacement is made only for cases whose answer is already directly known. For ordinary inputs, unused variables can be made incident by the tautological factors in NN-C023, which preserve the exact count. Thus no oracle for the hard parity is used in any edge case.

The equation $H+\theta F$ has polynomial circuit description: the rail product is already succinct, and each pure $r$th power is computed by repeated squaring and multiplication. It is nonzero because $F$ is divisible by $z_0$ whereas the rail terms of $H$ are not. Its hypersurface in $\mathbb P^{r-1}_{\mathbb F_{2^b}}$ is projective, and the zero-gradient criterion proves geometric smoothness. The nontrivial coherent group, in degree $r-2$, is the known line of NN-C023. Its absolute Frobenius is semilinear, with matrix coefficient
\[
          [z_0\prod_e z_{e,0}z_{e,1}](H+\theta F)=\theta c.
\]
The pure powers in $H$ make no contribution to this squarefree monomial. Iterating semilinear Frobenius $b$ times gives the field-linear scalar
\[
                   (\theta c)^{1+2+\cdots+2^{b-1}}
                      =(\theta c)^{2^b-1}=c.
\]
This uses $\theta\ne0$ and $c\in\{0,1\}$. Thus the deterministic smooth specialization preserves the hard entry exactly, rather than only its nonvanishing with some probability.

SGA7 II XXII3.1--3.2 give the same full congruence as in NN-C024,
\[
 \#X(\mathbb F_{2^b})\equiv1+c\pmod2,
 \qquad Z(X/\mathbb F_{2^b},T)\equiv\frac{1-cT}{1-T}\pmod2.
\]
Deligne's Theorem1.6 applies to every produced hypersurface, so these are genuine smooth projective finite-field realizations with pure l-adic cohomology. That theorem is an imported classical input, not a purity conclusion inferred from the congruence.

The construction is now a deterministic polynomial-time reduction from parity-3-SAT. On the encoded family scalar membership in $\oplus\mathrm P$ follows from the factor-choice verifier for $c$ in NN-C023. Point-count parity also lies in $\oplus\mathrm P$ directly: guess a projective coordinate vector in the explicit field, require its first nonzero coordinate to equal one, and verify the equation. This normalization gives exactly one representative per rational point, and $rb=O(r^3)$ guessed bits suffice. The complement by one in the displayed formula is an elementary parity-output conversion. Hence the asserted parity-completeness holds under deterministic reductions. A polynomial-time scalar or parity algorithm gives $\mathrm P=\oplus\mathrm P$; Valiant--Vazirani then gives $\mathrm{NP}=\mathrm{RP}$, as recorded in NN-C006. This strengthens the randomized-promise consequence in NN-C024.

E computes $c$ from the original source formula, keeping its original variable parameter and its stated density dependence. The field construction and the displayed geometric description add only polynomial overhead in the formula size. Neither the $rb$ projective-point witness length nor the full l-adic Betti number is substituted into that sharper original-formula evaluator. A and M are unused in this construction.
\end{proof}
\paragraph{Status and complexity impact.}
Complete singular-parameter confinement, deterministic field construction and preserved-Frobenius proof. Complete deduction-proof review by root: passed; imported input theorems are not claimed independently proved. The open finite-field specialization gate of NN-C024 is closed for its actual hard rail family. Smoothness, properness, finite base field, a specified rank-one coherent basis and a deterministic construction now coexist with a parity-complete Frobenius entry. This original construction has field degree at most quadratic in the succinct hypersurface parameter. NN-C031 subsequently removes the growing-field limitation by incidence padding over fixed F4; the present bound is retained as the historical general-parameter result. No polynomial bound on total l-adic rank is asserted. The result strengthens the conditional algorithmic consequence to $\mathrm P=\oplus\mathrm P$ and $\mathrm{NP}=\mathrm{RP}$, rather than merely a randomized reduction. It proves no unconditional complexity lower bound, no VP/VNP separation, and no realization of a GCT canonical basis.

\paragraph{Geometric and representation dictionary.}
Let $P\in\mathbb F_2[z_{e,0},z_{e,1}]$ be nonzero of degree two in each of $M$ pairs, and put $r=2M+1$, $F=z_0P$. In the Fermat pencil $H+uF$, every singular geometric fiber occurs at a parameter in the one finite field $\mathbb F_{2^{r\operatorname{ord}_r(2)}}$. The proof is the pairwise Euler identity, normalization at $z_0=1$, and the resulting root-of-unity and power equations above. A parameter of absolute degree greater than that bounded field degree therefore gives a smooth projective fiber; a quadratic Artin--Schreier tower supplies such a parameter with degree $O(r^2)$ and a polynomial-length presentation.

For the incidence-rail polynomials, the proper smooth fiber has $H^0(\mathcal O)=\mathbb F_{2^b}$ and one other nonzero coherent group, a canonically based line in degree $r-2$. Absolute Frobenius has semilinear coefficient $\theta c$, and the linearized field Frobenius has scalar $c$, the original parity count. The complete Cech proof is NN-C023 and the classical smooth formula is Katz2.3.7.17, p36. SGA1 II4.10(iii)/4.12 is the smoothness interface, SGA7 II XXII3.1--3.2 gives the proper zeta congruence, and Deligne Weil I1.6 gives purity of the actual l-adic realizations. These provide an exact deterministic finite-field geometric dictionary for the hard coefficient. A purity or existence theorem does not evaluate it, and the construction supplies no canonical identification with GCT multiplicities.



\section{NN-C031 --- Forced-variable padding puts the smooth rank-one parity gate over fixed F4}
\begin{proposition}[NN-C031]
Fix once and for all
\[
 K=\mathbb F_4=\mathbb F_2[\omega]/(\omega^2+\omega+1).
\]
There is a deterministic polynomial-time reduction from parity-3-SAT to a family of smooth projective hypersurfaces over this fixed field. Each has a polynomial-size homogeneous arithmetic-circuit equation, and its top coherent cohomology is a canonically based one-dimensional space. The $4$-power Frobenius scalar on that line is the original satisfying-assignment parity $c$. Point-count parity is $1+c$.

The geometric statement is uniform: let $k\geq3$ be odd, $r=2^k-1$, $M=(r-1)/2$, and let $P$ be any nonzero polynomial over $\mathbb F_2$ homogeneous of degree two in every one of $M$ variable pairs. Then
\[
 X_P=V\left(z_0^r+\sum_{e=1}^{M}(z_{e,0}^r+z_{e,1}^r)
                         +\omega z_0P\right)
                       \subset\mathbb P^{r-1}_K
\]
is smooth projective. If $c$ is the squarefree coefficient of $z_0P$, its top coherent line has absolute-Frobenius coefficient $\omega c$ and $4$-Frobenius scalar $c$. The hard subclass has $P$ in the succinct rail-factor format, so its scalar and point-count parity are $\oplus\mathrm P$-complete under deterministic polynomial-time reductions. Polynomial-time extraction gives $\mathrm P=\oplus\mathrm P$ and $\mathrm{NP}=\mathrm{RP}$. No varying base field is needed.
\end{proposition}
\paragraph{Exact inputs.}
NN-C030. complete singular-parameter confinement lemma for pairwise degree-two P: every bad a lies in F\_(2\textasciicircum{}(r ord\_r2)).

NN-C023. incidence-rail coefficient encoding; tautological factors for unused declared variables; one-dimensional coherent basis.

SGA1, ExposeII, Section4. Theorem4.10(iii) and Corollary4.12; smooth hypersurface interface used in NN-C030.

Katz, Algebraic Solutions of Differential Equations, Invent.Math.18(1972),1--118. Rank-one Special Case2.3.7.17, printedp36; full direct Cech proof in NN-C023.

SGA 7 II, Expose XXII. Theorem 3.1 (3.1.1); Corollary 3.2 (3.2.1).

Deligne, La conjecture de Weil I, PMIHES43(1974),273--307. Theorem1.6, printedp276.

Valiant--Vazirani, NP is as easy as detecting unique solutions, TCS 47 (1986), 85--93. Theorem 1.1; Corollaries 3.1--3.4, especially Corollary 3.2 p.91.

E: The Exponential Time Hypothesis is False, supplied LaTeX. Theorem1.1; evaluator retains the source formula parameter before padding.

\begin{proof}
First prove the uniform smoothness statement. For $r=2^k-1$, the multiplicative order of $2$ modulo $r$ is exactly $k$. The exponent $k$ works, while for $0<i<k$ the positive integer $2^i-1$ is strictly smaller than $r$ and therefore cannot be divisible by $r$. NN-C030 confines every bad parameter in the Fermat pencil $H+u z_0P$ to
\[
                            \mathbb F_{2^{rk}}.
\]
Both $r$ and $k$ are odd. The element $\omega$ has degree two over $\mathbb F_2$, so it cannot lie in that odd-degree extension: otherwise the tower law would make two divide $rk$. Therefore the fiber at the fixed parameter $\omega$ is smooth. The same reasoning applies to its conjugate $\omega^2$. Projectivity follows from the homogeneous equation. The form is nonzero, since its rail pure-power terms are not divisible by $z_0$, whereas $z_0P$ is.

The coherent calculation from NN-C023 applies over $K$ without change. There is $H^0(\mathcal O)=K$, one other nonzero coherent group in degree $r-2$, and its canonical boundary/Cech basis has absolute-Frobenius coefficient equal to the full squarefree coefficient of the defining form. The Fermat terms do not contribute to that coefficient. Thus the coefficient is $\omega c$, with $c\in\mathbb F_2$. Iterating this semilinear map twice yields the $K$-linear $4$-power Frobenius with scalar
\[
                        (\omega c)^{1+2}=\omega^3c^3=c.
\]
There is no semilinear-versus-linear ambiguity: absolute Frobenius sends $ae$ to $a^2\omega c\,e$, while its second iterate sends $ae$ to $a c\,e$ because $a^4=a$ in $K$.

We next arrange this special degree with polynomial padding while preserving the original count. Start with the incidence-rail construction of NN-C023, with $M_0\geq1$ incidences and $r_0=2M_0+1$. Choose the least odd integer $k\geq3$ such that
\[
                  r=2^k-1\geq r_0,
                 \qquad M=(r-1)/2.
\]
This is computed from the binary length of $r_0+1$ and a parity check. The new degree is linear in the old one: if $k>3$, minimality gives $2^{k-2}-1<r_0$, hence $r<4(r_0+1)-1$; the bounded case $k=3$ also satisfies a constant-factor bound. Consequently $M-M_0=O(M_0)$.

Append exactly $M-M_0$ fresh Boolean variables, each with its own positive unit clause. Each new variable has exactly one allowed value, so every old satisfying assignment has exactly one extension. The number of satisfying assignments, and in particular its parity, is unchanged. Each such gadget adds one clause, one variable and exactly one incidence; no old occurrence is modified. Thus the new incidence count is exactly $M$. In the rail polynomial its clause factor and its two-term variable-consistency factor together have degree two in the new rail pair and squarefree coefficient one. All old factors remain unchanged. This verifies both the required multihomogeneity and the coefficient preservation directly at the polynomial level.

As in NN-C030, treat the elementary exceptional source instances before forming the rails. An empty clause has known parity zero; an instance with no declared variables has its one assignment directly testable. Replace it by the positive unit formula $x$ for parity one or the contradictory unit formula $x\wedge\neg x$ for parity zero. These are fixed instances with nonempty incidence sets. For all other inputs, tautological factors make every otherwise unused declared variable incident without restricting it. Tautological clauses and repeated literals can be simplified first. After these steps each local clause factor is nonzero and the product $P$ is nonzero over $\mathbb F_2$. No decision of the hard parity is used to choose a padding gadget.

The construction now has precisely $r$ coordinates and a degree-$r$ equation over the fixed two-bit field $K$. Every local clause factor has at most eight monomials; each variable-consistency factor has two. The number of factor incidences and their arithmetic-circuit descriptions is polynomial in the source formula size. Each pure power $z^r$ takes $O(\log r)$ multiplications by repeated squaring, so the entire sum of Fermat terms and the rail product has a polynomial-size arithmetic circuit. Its degree is itself polynomial in the source size. This assertion concerns the circuit/factor presentation; fully expanding its coefficients can still require exponentially many monomials.

The actual smooth hypersurface therefore satisfies
\[
 \#X_P(\mathbb F_4)\equiv1+c\pmod2,
 \qquad
 Z(X_P/\mathbb F_4,T)\equiv\frac{1-cT}{1-T}\pmod2,
\]
by SGA7 II XXII3.1--3.2, exactly as in NN-C024. Deligne's Theorem1.6 supplies purity of its l-adic cohomology because every output is smooth and projective. The proof uses the singular-parameter confinement to establish smoothness first; it does not infer smoothness or purity from a point-count congruence.

For a precise complete problem, use the syntactically specified incidence-rail factor format of NN-C023, with nonempty scopes, the displayed consistency factors, and an incidence count giving $r=2^k-1$ for odd $k\geq3$. These conditions are checkable in polynomial time, and fixed straight-line instructions generate the equation $H+\omega z_0P$ from that format. No identity test against an arbitrary purported equivalent circuit is part of the input validation. Every such input is smooth by the uniform theorem. Nondeterministically choose one monomial from each factor and check that the exponents sum to the full squarefree vector. The parity of the number of accepting choices is $c$, proving $\oplus\mathrm P$ membership for the Frobenius scalar on this class. The preceding padded reduction proves hardness. This is a complete problem on a specified geometric presentation family, not a general promise that every smooth hypersurface has a one-dimensional coherent top group.

Point-count parity belongs to $\oplus\mathrm P$ independently of the special coefficient format. Guess its $r$ coordinates in $K$, require the first nonzero coordinate to be one, and check the circuit equation. The chosen normalization gives exactly one vector for each projective rational point. There are $2r$ guessed bits, and all field operations have constant bit cost. The displayed parity formula converts this output to $c$ by adding one. Thus point-count parity on this smooth fixed-field family is also parity-counting complete. An algorithm correct on arbitrary smooth projective inputs in particular handles these guaranteed smooth instances, so polynomial time would give $\mathrm P=\oplus\mathrm P$. Valiant--Vazirani then gives $\mathrm{NP}=\mathrm{RP}$ with the scope already detailed in NN-C006.

E computes the parity by counting the source formula before the forced-variable padding. Its original variable count and clause-density hypotheses therefore remain the relevant parameters, with only polynomial overhead for building the padded geometric input. Running E instead on all geometric coordinates would lose that sharper statement. The field is the one fixed $\mathbb F_4$ throughout, and neither A nor M is used.
\end{proof}
\paragraph{Status and complexity impact.}
Complete fixed-field smoothness, exact padding, coherent-Frobenius and parity-completeness proofs. Complete deduction-proof review by root: passed; imported input theorems are not claimed independently proved. The remaining varying-field qualification of NN-C030 is removed: the hard rank-one coherent Frobenius entry is realized deterministically on smooth projective hypersurfaces over fixed $\mathbb F_4$. The padding is linear and parsimonious at the Boolean count level. Neither total l-adic rank nor the fully expanded defining polynomial is claimed polynomially bounded; the geometric equation is succinct. This yields a strictly smaller effective gate than arbitrary field construction or random specialization. It supplies actual pure finite Weil realizations for the counting gadgets without evaluating the hard coefficient or identifying a GCT canonical basis.

\paragraph{Geometric and representation dictionary.}
For every odd $k\geq3$, every nonzero pairwise-degree-two $P$ over $\mathbb F_2$ in $M=2^{k-1}-1$ pairs defines the smooth projective hypersurface
\[
 V\left(z_0^{2^k-1}+\sum_e(z_{e,0}^{2^k-1}+z_{e,1}^{2^k-1})
                        +\omega z_0P\right)/\mathbb F_4.
\]
The singular-parameter field from NN-C030 has odd degree $k(2^k-1)$ and therefore excludes the fixed quadratic parameter $\omega$. Its top coherent line is explicitly trivialized by the boundary/Cech construction. Absolute Frobenius has semilinear coefficient $\omega c$; $4$-power Frobenius has scalar $c$. The lower coherent groups are only the degree-zero field, and the proper zeta congruence is $(1-cT)/(1-T)$ modulo two. Exact classical interfaces are Katz2.3.7.17, SGA1 II4.10(iii)/4.12, SGA7 II XXII3.1--3.2 and Deligne Weil I1.6.

The forced-unit-clause padding turns any source formula into this geometric class while preserving its count exactly. Thus the field, smoothness and coherent basis are all supplied effectively, and the sole unknown scalar is parity-counting complete. This is an actual smooth projective finite-field construction, with pure l-adic cohomology, and a precise computational operation inside it. It does not assert a canonical representation-theoretic realization of permanent/determinant multiplicities.



\section{NN-C032 --- An explicit fixed-F2 smooth projective model has four-dimensional coherent cohomology with hard Frobenius}
\begin{proposition}[NN-C032]
From the fixed-$\mathbb F_4$ hypersurface $X=V(G)$ of NN-C031, with $G=H+\omega F$ in $r$ coordinates, one can deterministically construct a smooth geometrically integral projective scheme
\[
             W\subset\mathbb P^{r^2-1}_{\mathbb F_2},
              \qquad W\simeq\operatorname{Res}_{\mathbb F_4/\mathbb F_2}X.
\]
It has a polynomial-size homogeneous circuit-equation presentation consisting of $O(r^4)$ quadratic equations and $O(r)$ degree-$r$ equations. Put $d=r-2$, which is odd. Its coherent cohomology is
\[
 H^0(W,\mathcal O)=\mathbb F_2,\quad
 H^d(W,\mathcal O)=\mathbb F_2^2,\quad
 H^{2d}(W,\mathcal O)=\mathbb F_2,
\]
and every other group vanishes. In explicitly descended bases, absolute Frobenius is respectively
\[
                   1,\qquad
              c\begin{pmatrix}0&1\\1&0\end{pmatrix},
                   \qquad c,
\]
where $c$ is the original satisfying-assignment parity. Thus the entire coherent realization has dimension four, its basis is supplied, and one of its one-dimensional Frobenius scalars is parity-counting complete over the fixed prime field.

Moreover
\[
 \#W(\mathbb F_2)=\#X(\mathbb F_4),\qquad
 \#W(\mathbb F_2)\equiv1+c\pmod2,\qquad
 Z(W/\mathbb F_2,T)\equiv\frac{1-cT}{1-T}\pmod2.
\]
These yield deterministic parity-counting-complete problems on the encoded family. Every output is smooth projective and has pure l-adic cohomology. The presentation uses a projective system of equations; it is not asserted to be a hypersurface over $\mathbb F_2$.
\end{proposition}
\paragraph{Exact inputs.}
NN-C031. deterministic fixed-F4 smooth projective hypersurface, canonical coherent basis and absolute Frobenius coefficient omega*c.

Stacks Project, Constructions of Schemes. Lemma27.13.6, Tag01WD, Segre embedding; the exact two-by-two-minor chart computation is proved below.

Stacks Project, Criteria for Representability. Section97.11, Tag05Y8; equation97.11.0.2, Lemma97.11.1 (fppf sheaf) and Lemma97.11.2 (base change).

Stacks Project, Descent. Lemma35.23.29, Tag02VL: smoothness is fpqc local on the base.

Stacks Project, Cohomology of Schemes. Lemma30.5.2, Tag02KH, flat base change; especially its affine-base statement(2) and Cech proof.

Stacks Project, Varieties. Lemma33.29.1, Section Tag0BEC: coherent Kunneth for quasi-compact schemes with affine diagonal; cup-product normalization and complete proof.

Stacks Project, Varieties and Commutative Algebra. Lemma33.25.3, Tag056S: smooth over a field implies regular; Lemma10.106.2, Tag00NP: a regular local ring is a domain.

SGA 7 II, Expose XXII. Theorem 3.1 (3.1.1); Corollary 3.2 (3.2.1).

Deligne, La conjecture de Weil I, PMIHES43(1974),273--307. Theorem1.6, printedp276.

Valiant--Vazirani, NP is as easy as detecting unique solutions, TCS 47 (1986), 85--93. Theorem 1.1; Corollaries 3.1--3.4, especially Corollary 3.2 p.91.

E: The Exponential Time Hypothesis is False, supplied LaTeX. Theorem1.1, original-source evaluator.

\begin{proof}
Write $K=\mathbb F_4$ and let $\sigma$ be its nontrivial automorphism, so $\sigma(\omega)=\omega^2=\omega+1$. Use $r^2$ homogeneous coordinates over $\mathbb F_2$: diagonal coordinates $a_i$, and coordinates $b_{ij},c_{ij}$ for every $i<j$. Form the matrix with entries in $K[a,b,c]$
\[
 M_{ii}=a_i,\qquad
 M_{ij}=b_{ij}+\omega c_{ij},\qquad
 M_{ji}=b_{ij}+\omega^2 c_{ij}\quad(i<j).
\]
Coefficientwise conjugation, with the coordinate indeterminates held fixed, satisfies $\sigma(M)=M^{\mathsf t}$. For a polynomial $A$ over $K$ in these coordinates, write uniquely $A=A_0+\omega A_1$ with $A_0,A_1$ over $\mathbb F_2$.

Define $W$ by the following homogeneous equations over $\mathbb F_2$: both coefficient components of every two-by-two minor of $M$, and both components of
\[
                     G(M_{1j},\ldots,M_{rj})
                       \quad(1\leq j\leq r).
\]
This is a closed projective scheme. There are $2\binom r2^2=O(r^4)$ minor equations and $2r$ column equations, with degrees two and $r$ respectively; redundant or zero equations may be retained.

We identify this scheme after the faithfully flat extension $K/\mathbb F_2$, including its scheme structure. On each off-diagonal pair, the coordinate change
\[
        (b_{ij},c_{ij})\longmapsto(M_{ij},M_{ji})
\]
has determinant $\omega^2-\omega=1$, so over $K$ the entries of $M$ are independent matrix coordinates. Imposing both $\mathbb F_2$ coefficient components of a polynomial is equivalent over $K$ to imposing that polynomial and its coefficientwise conjugate. For the minors the conjugates are other minors of the transpose, already in the list. Thus their equations define the ordinary projective rank-one matrix scheme.

Here is the exact scheme-theoretic chart check. On a projective chart where $M_{ij}$ is nonzero, normalize it to one. The minors force
\[
                         M_{kl}=M_{kj}M_{il}.
\]
The remaining entries of column $j$ and row $i$ are free coordinates, giving the standard affine chart of
$\mathbb P^{r-1}\times\mathbb P^{r-1}$ under the Segre map $(z,w)\mapsto zw^{\mathsf t}$. These charts cover the projective minor scheme, so no extra embedded component is being discarded. This is the explicit minor form of the Segre embedding in Stacks27.13.6.

Under $M=zw^{\mathsf t}$, the column equations are
\[
                      w_j^rG(z)=0\quad(1\leq j\leq r).
\]
Their conjugates are the row equations
\[
                      z_i^rG^\sigma(w)=0\quad(1\leq i\leq r),
             \qquad G^\sigma=H+\omega^2F.
\]
Locally on the product of projective spaces some $z_i$ and some $w_j$ are units, so the generated sheaf of ideals is exactly $(G(z),G^\sigma(w))$. This proves an isomorphism of schemes
\[
                         W_K\simeq X\times_K X^\sigma.
\]
In these matrix coordinates the descent automorphism conjugates coefficients and transposes $M$; it therefore interchanges the two factors with their conjugation data. This is precisely the descent datum of restriction of scalars for the quadratic extension. More formally, $K\otimes_{\mathbb F_2}K\simeq K\times K$, so the restriction-of-scalars functor base-changes to $X\times X^\sigma$, with that same factor exchange. The functor is an fppf sheaf and commutes with base change by Stacks97.11.1--97.11.2. The displayed compatible isomorphism over the faithfully flat cover thus descends to an identification of its sheaf with the represented sheaf of $W$. This proves the claimed Weil-restriction identification; no general projectivity theorem for restriction of scalars is needed because the projective equations are already explicit.

Both $X$ and $X^\sigma$ are smooth by NN-C031. Their product is smooth, and smoothness descends under the finite faithfully flat extension by Stacks35.23.29. Hence $W$ is smooth over $\mathbb F_2$, of dimension $2(r-2)=2d$. Over an algebraic closure, each factor has only the ground field as its degree-zero coherent cohomology, by NN-C023's Cech calculation. The same is true of the product by coherent Kunneth. It is therefore connected, since a decomposition into two open-and-closed pieces would give a nontrivial idempotent in its global functions. A smooth scheme over a field has regular local rings by Stacks33.25.3, and regular local rings are domains by Stacks10.106.2; distinct irreducible components consequently cannot intersect. Thus this connected smooth geometric scheme is integral. This proves geometric integrality of $W$ as well.

We now compute its full coherent Frobenius, taking care with descent and semilinearity. Over $K$, let $e$ be the canonical basis of $H^d(X,\mathcal O_X)$ from NN-C023 and let $e'$ be its conjugate basis on $X^\sigma$. Coherent Kunneth (Stacks33.29.1) and flat base change (Stacks30.5.2) show that the only cohomology groups of $W_K$ are
\[
 K\quad\text{in degree }0,\qquad
 Ke\oplus Ke'\quad\text{in degree }d,\qquad
 K(e\mathbin\cup e')\quad\text{in degree }2d.
\]
Here $e,e'$ in the middle line mean their pullbacks from the two factors. The Kunneth isomorphism is the pullback-and-cup-product map, so it is compatible with Frobenius and the factor exchange, including its graded sign.

The descent automorphism is semilinear on coefficients and sends $e$ to $e'$ and $e'$ to $e$. Thus the following middle vectors are fixed:
\[
                    v_1=e+e',\qquad
                    v_2=\omega e+\omega^2e'.
\]
They are independent over $K$, because the determinant of their coefficients is $\omega^2-\omega=1$. Their fixed coefficient combinations are exactly their $\mathbb F_2$ span. The top vector $w=e\cup e'$ is also fixed: exchange gives $e'\cup e=(-1)^{d^2}e\cup e'$, and this sign is one in characteristic two. Hence $v_1,v_2,w$ are explicit descended bases over $\mathbb F_2$. This also follows by taking the invariant vectors in the flat-base-change isomorphism; it does not require computing any unknown Frobenius entry.

Let $\mathfrak F$ denote absolute Frobenius on coherent cohomology. On the two factors, NN-C031 gives
\[
                \mathfrak F(e)=\omega c\,e,
                \qquad\mathfrak F(e')=\omega^2c\,e',
 \qquad \mathfrak F(a\xi)=a^2\mathfrak F(\xi).
\]
Under flat base change from $\mathbb F_2$ to $K$, this absolute Frobenius corresponds to Frobenius on $H^*(W,\mathcal O_W)$ tensored with squaring on $K$, not with the identity on $K$. Absolute Frobenius is natural for the displayed isomorphism with the product. Restricting its semilinear action to the descended $\mathbb F_2$ vectors therefore computes the original $\mathbb F_2$-linear Frobenius on $W$ exactly.
Applying these formulas to the descended middle basis gives
\[
       \mathfrak F(v_1)=c v_2,\qquad
       \mathfrak F(v_2)=c v_1.
\]
For example, the coefficient of $e$ in the second identity is
$\omega^2\omega c=c$, and the other coefficient is
$(\omega^2)^2\omega^2c=c$. Cup-product compatibility gives
\[
                 \mathfrak F(w)=\omega\omega^2c^2w=cw.
\]
On degree zero it is the identity over $\mathbb F_2$. These computations prove exactly the three displayed matrices. As a separate consistency check, the square of the middle matrix is $cI$, agreeing with $4$-power Frobenius on both factors; the top scalar has square $c$ as required. The argument supplies a basis and a symbolic matrix description in terms of $c$, not a polynomial-time evaluation of that unknown bit.

There is also an elementary exact rational-point bijection, which does not use a trace formula. If $[z]\in X(K)$, define
\[
                            M=z\sigma(z)^{\mathsf t}.
\]
It is Hermitian and rank one. Replacing $z$ by $\lambda z$ for $\lambda\in K^*$ leaves it unchanged, since $\lambda\sigma(\lambda)=\lambda^3=1$. Every column is a scalar multiple of $z$, so the column equations vanish. Its diagonal entries belong to $\mathbb F_2$, and every off-diagonal pair has the unique displayed $(b_{ij},c_{ij})$ coordinates in $\mathbb F_2$. Thus it defines an $\mathbb F_2$ point of $W$.

Conversely, an $\mathbb F_2$ point gives a nonzero Hermitian rank-one matrix over $K$. It has a nonzero diagonal entry. If every diagonal were zero, the minor relation
$M_{ij}M_{ji}=M_{ii}M_{jj}$ would force every off-diagonal entry to vanish, since $M_{ij}M_{ji}$ is its nonzero norm whenever $M_{ij}\ne0$. This would give the zero matrix. Choose $i$ with $M_{ii}=1$, and let $z$ be column $i$. The minors give $M=z\sigma(z)^{\mathsf t}$, with $z_i=1$. The column equation implies $G(z)=0$. The column space recovers $[z]$ uniquely. These two constructions are inverse, proving
\[
                             \#W(\mathbb F_2)=\#X(K).
\]
Projective scaling on the $\mathbb F_2$ matrix coordinates introduces no additional quotient, since $\mathbb F_2^*=\{1\}$.

For the full zeta congruence, apply SGA7 II XXII3.1 to the just-computed coherent Frobenius. The degree $d$ is odd and $2d$ is even. The middle determinant is
\[
       \det\left(I-Tc\begin{pmatrix}0&1\\1&0\end{pmatrix}\right)
                      =1-c^2T^2=(1-cT)^2
\]
in $\mathbb F_2[T]$. Dividing by the degree-zero and top determinants gives
\[
     Z(W/\mathbb F_2,T)\equiv
         \frac{(1-cT)^2}{(1-T)(1-cT)}
                   =\frac{1-cT}{1-T}\pmod2.
\]
The point-count congruence is also immediate from the exact bijection and NN-C031. The full series congruence is obtained from the coherent determinants, not inferred merely from parity of the first point count. Deligne's theorem applies to the smooth projective $W$, giving the actual pure l-adic realizations independently of this characteristic-two calculation.

All constructions are effective with polynomial output length. A $K$ arithmetic gate can be replaced by constantly many $\mathbb F_2$ gates computing its two coordinate components: for example
\[
 (a+b\omega)(u+v\omega)
                  =(au+bv)+(av+bu+bv)\omega.
\]
The $O(r^4)$ minors have constant-size circuits, and substituting each of the $r$ columns into $G$ multiplies its circuit size by at most $r$, up to a constant coordinate-arithmetic factor. There are only $r^2$ projective coordinates. No expansion of $G$ into all monomials is required, and no factorization or Galois-descent oracle is hidden in these formulas.

The source-to-$W$ construction is a deterministic polynomial-time reduction. Extracting the top Frobenius scalar is exactly extracting $c$, so it is $\oplus\mathrm P$-hard and is in $\oplus\mathrm P$ on the syntactic input family through the factor-choice verifier of NN-C023. Point-count parity is directly in $\oplus\mathrm P$: guess the $r^2$ binary projective coordinates, reject the zero vector, and check the homogeneous circuit equations. Every nonzero binary vector is already a unique projective representative. Its output is $1+c$, which has the same parity-counting complexity; complementation preserves $\oplus\mathrm P$ by adding one distinguished accepting computation. Thus the point-parity decision language is also complete under deterministic many-one reductions. A polynomial-time algorithm for either geometric task gives $\mathrm P=\oplus\mathrm P$ and $\mathrm{NP}=\mathrm{RP}$ as before.

E evaluates $c$ through the original formula before its geometric padding and descent. The extra construction cost is polynomial in that source size. A and M are unused. The geometric input is a prescribed homogeneous circuit-equation family with its mechanically generated format; no polynomial-time recognition theorem for arbitrary smooth projective presentations is being asserted.
\end{proof}
\paragraph{Status and complexity impact.}
Complete explicit projective descent, scheme-structure, coherent-Frobenius, point-bijection and complexity proofs. Complete deduction-proof review by root: passed; imported input theorems are not claimed independently proved. The hard coherent-Frobenius gate now occurs over the fixed prime field, in a smooth geometrically integral projective scheme with a supplied four-dimensional coherent basis. Neither random specialization, growing coefficient fields, singular fibers, unknown coherent rank nor unknown basis is the missing operation. The single bit $c$ is still parity-counting complete. The total l-adic rank and expanded equation size are not claimed polynomial. This identifies a concrete remaining scalar-extraction gate inside an actual pure finite Weil realization; it does not supply GCT's canonical-basis functor or an unconditional circuit lower bound.

\paragraph{Geometric and representation dictionary.}
Let $X=V(H+\omega F)/\mathbb F_4$ be a hypersurface of NN-C031. The displayed Hermitian-matrix equations define a projective $\mathbb F_2$ scheme $W$ whose base change is the Segre product $X\times X^\sigma$ and whose descent datum is factor exchange. Thus $W$ represents the Weil restriction. It is smooth and geometrically integral, with coherent cohomology only in degrees $0,d,2d$, of ranks $1,2,1$. In the explicit descended bases $1$, $(e+e',\omega e+\omega^2e')$, and $e\cup e'$, its Frobenius operators are $1$, $c\left(\begin{smallmatrix}0&1\\1&0\end{smallmatrix}\right)$, and $c$. Its exact rational-point bijection with $X(\mathbb F_4)$ and its proper zeta congruence are proved above.

The scheme-theoretic interfaces are the explicit Segre chart proof with Stacks27.13.6, restriction-of-scalars sheaf/base-change Lemmas97.11.1--97.11.2, smoothness descent Lemma35.23.29, coherent flat base change Lemma30.5.2 and coherent Kunneth Lemma33.29.1. SGA7 II XXII3.1--3.2 supplies the coherent zeta congruence and Deligne Weil I1.6 supplies purity. These are exact locators, with no original Grothendieck archive cote claimed. The formula for the descended middle matrix is proved by semilinear arithmetic and cup-product compatibility, not guessed from point counts. The underlying count reduction and E evaluator are unchanged; the new construction removes the fixed-field obstruction while preserving an explicit small coherent realization.



\section{NN-C033 --- Integral swap matrices remove the fixed-characteristic symbolic-circuit gate}
\begin{proposition}[NN-C033]
Fix arity $w$ and factor density $C$. For every prime $p$ and every factor hypergraph on $n$ Boolean variables with $m\leq Cn$ factors of scope at most $w$, its universal partition polynomial
\[
 Z_H(y)=\sum_{x\in\{0,1\}^n}\prod_f y_{f,x|I_f}
                      \in\mathbb F_p[y]
\]
has a deterministically constructible, division-free arithmetic circuit over $\mathbb F_p$ of size
\[
 S(n)=\operatorname{poly}(n)\,
       2^{O_{w,C}(n(\log\log n)^2/\sqrt{\log n})}.
\]
Construction bit time is $S(n)\operatorname{poly}(\log p)$. The equality is equality of formal polynomials, with arbitrary independent weights, in every characteristic including two. The output circuit uses only prime-field constants. A small auxiliary field is used during construction and eliminated by explicit coordinate simulation.

Consequently all ordinary ideal-multiplier and affine-trivialization constructions of NN-C029 hold over every field of characteristic $p$, with the same variable parameter and subexponential bound. In particular, an unsatisfiable bounded-density local Boolean verifier has a Hilbert-like IPS certificate over $\mathbb F_2$ of that size, with ordinary polynomial identity $\sum_i a_i f_i=1$ and multiplier degrees at most $v$. For a fixed-width $O(n)$-clause CNF the original-$n$ specialization still needs no Boolean axioms. The output multipliers are exactly reductions of the explicitly specified integral first-violation and local-correction polynomials.

After polynomial overhead, the prime-field circuits can be made constant-free with syntactic formal degree at most $\max(m,1)$ for $Z_H$, at most $\max(v,1)$ for the multipliers, and at most $v+1$ for the IPS certificate. Here constant-free means that the only constant leaves are $0,1,-1$, and formal degree gives each constant or variable leaf degree one, addition the maximum, and multiplication the sum. This is a positive-characteristic statement. This prime-field argument itself does not construct integral-gate circuits over $\mathbb Z$; that remaining step is subsequently proved in NN-C034.
\end{proposition}
\paragraph{Exact inputs.}
E: The Exponential Time Hypothesis is False, supplied LaTeX. Lemma3.2 complete integer-measurement-fiber proof; Lemma3.3 ring-valued Boolean differences; Proposition3.4 Taylor routine; Lemmas4.1--4.3 swap matrix and row support; equation swap-matrix-definition; Lemmas4.5--4.8 layout, height and parameters.

NN-C016. bounded-arity universal partition polynomial, high-incidence branching and coefficient-independent circuit unrolling.

NN-C029. ordinary integral ideal-multiplier correction, sparse-CNF first-violation identity, homogeneous components and affine coordinate maps.

EGA IV, fourth part, Publications Mathematiques IHES32 (1967). Printed p.111,18.2.1--18.2.2; base-change equation18.2.1.2; primary PDF page110 (zero-based109).

Elementary finite-field construction and Chinese remainder theorem. Complete existence, deterministic bounded-size construction and Hasse-jet argument proved below; no fast irreducible-polynomial oracle is assumed.

\begin{proof}
We modify the actual common routine of E, rather than reducing its rational circuit modulo $p$. The key is an integral strengthening of its auxiliary swap matrix.

Use E's notation: commuting integral permutation involutions $\Pi_1,\ldots,\Pi_Q$, their product $V_0$, and
\[
 W_0=\frac12\sum_{s=1}^Q(I-\Pi_s),\qquad
 \mathcal I=\{A_0,\ldots,A_0+R-1\},\qquad
 M_0=(-1)^{A_0}V_0\sum_{j=0}^{R-1}
                      (-2)^j\binom{W_0-A_0I}{j}.
\]
Although this displayed construction uses rational matrices, $M_0$ is in fact integral. In the formal matrix power-series ring over $\mathbb Q$, one has
\begin{equation}\label{eq:C033-integral-generating}
 \sum_{j\geq0}(-2)^j\binom{W_0-A_0I}{j}U^j
       =(1-2U)^{-A_0}\prod_{s=1}^Q\bigl(I+U(\Pi_s-I)\bigr).
\end{equation}
To prove the identity, use the common rational eigenbasis in E Lemma4.1. If exactly $a$ involutions act by minus one, the right side acts by
$(1-2U)^{-A_0}(1-2U)^a=(1-2U)^{a-A_0}$, whose generalized binomial expansion is exactly the left side. Equality on this basis proves the matrix identity coefficient by coefficient. The scalar series $(1-2U)^{-A_0}$ has integral coefficients: for $A_0>0$ its coefficient of $U^j$ is $2^j\binom{A_0+j-1}{j}$; for $A_0=0$ the entire series is $1$. Every factor on the right is an integral matrix polynomial. Therefore each coefficient on the left is integral, and so is $M_0$.

This proof also gives a direct integral construction by truncating the product at degree $R-1$. It takes polynomially many operations in the auxiliary dimension $D$, $Q$ and $R$. Since $A_0,R\leq Q$, the coefficients have polynomial bit length, for example the norm bounds of E Lemma4.6 imply $O(Q\log(Q+2))$ bits for entries of $M_0$. The stronger conclusion is integrality, not a new asymptotic height assumption. The generating identity is proved over $\mathbb Q$ only to identify this particular integral matrix; after the integer coefficients have been formed, reduction or scalar extension is an ordinary ring homomorphism. No occurrence of $1/2$ or $1/j!$ is evaluated in characteristic two.

Keep the integer measurement matrix $C^{\mathcal B}$ of E Lemma4.1, the same incidence labels and copies, and the same integral matrices $C_{t,g}$. Two facts of E's construction remain available. First, $M_0[u,y]\ne0$ implies at most $R-1$ mismatches between $a_y(s)$ and $d_u(s)$, by Lemma4.2. Second,
\[
   C^{\mathcal B}\xi=0,\quad \xi\in\mathbb Z^{\mathcal B}
                  \quad\Longrightarrow\quad M_0\xi=\xi
                       \quad\hbox{in }\mathbb Z^{\mathcal B}.
\]
Indeed the hypothesis also holds over $\mathbb Q$, where Lemma4.1 proves the conclusion; integrality then makes it an exact integer equality. We do not claim that the reduction of $M_0$ fixes the entire kernel of the reduced matrix over $\mathbb F_p$. That stronger assertion is unnecessary and is not obtained by reducing a rational kernel calculation.

Choose an explicit extension field $L/\mathbb F_p$ containing at least $Q$ distinct nonzero elements, denoted $\alpha_1,\ldots,\alpha_Q$. If $p>Q$, take $L=\mathbb F_p$ and $\alpha_t=t$. If $p\leq Q$, let $e$ be least with $p^e>Q$. Then
\[
                 p^e\leq pQ\leq Q^2,\qquad e=O(\log(Q+1)).
\]
An explicit monic irreducible polynomial of degree $e$ can be found deterministically by enumerating the $p^e$ candidates and testing division by every monic polynomial of positive degree at most $\lfloor e/2\rfloor$. A reducible polynomial has a proper factor of degree at most half its degree, so this test is exact, and the total enumeration has cost polynomial in $Q$ and $\log p$.

For completeness, existence does not require a randomized construction theorem. In an algebraic closure the roots of $T^{p^e}-T$ are distinct and closed under the field operations, so they form a field with $p^e$ elements. Every element of smaller degree has degree a proper divisor of $e$, by the period of Frobenius, and hence belongs to the union of the root sets of $T^{p^d}-T$ for proper divisors $d$ of $e$. For $e\geq2$, the size of that union is at most
$\sum_{d=1}^{\lfloor e/2\rfloor}p^d<p^e$. Some element therefore has degree $e$, and its minimal polynomial is an irreducible candidate. The case $e=1$ is immediate. Represent $L$ by the resulting quotient with basis $1,\theta,\ldots,\theta^{e-1}$, and take any $Q$ distinct nonzero coordinate vectors. All arithmetic in this field has bit cost polynomial in $Q$ and $\log p$.

For a pair $(u,y)$ with $M_0[u,y]\ne0$, write
$E_{u,y}=\{s:a_y(s)\ne d_u(s)\}$ and replace E's scalar polynomial by
\[
       \varphi_{u,y}(T)=\prod_{s\in E_{u,y}}(1-T/\alpha_s).
\]
Form the new rows over $L$ by
\[
 \vec P_u(T)=\delta_u(I-M_0)
          +\sum_yM_0[u,y]\varphi_{u,y}(T)\delta_y.
\]
Integer matrix entries in this formula are reduced into $L$. The rows have degree at most $R-1$ and satisfy $\vec P_u(0)=\delta_u$. At $T=\alpha_t$, every term with $a_y(t)\ne d_u(t)$ vanishes. All denominators here are nonzero field elements chosen explicitly.

The separation property needed by the routine still holds on its actual enumeration domain. Suppose two Boolean assignments differ only inside an off-target group and have equal \emph{integer} measurements $C_{t,g}x_g$. Their difference, copied into each auxiliary array, is an integer vector in the kernel of $C^{\mathcal B}$ and supported in that off-target group. The term $I-M_0$ annihilates this integer vector by the exact integer equality above. Every surviving $\varphi$ term uses the target group and hence also annihilates the difference. Reducing these equalities into $L$ proves that $\langle\vec P_{f,v}(\alpha_t),x\rangle$ is constant on each such integer measurement fiber. The same reasoning applies over every commutative $L$-algebra. It uses an equality of integer vectors before scalar extension, not an assumed saturation of a reduced kernel.

We spell out why this weaker fiber statement suffices for the entire symbolic routine. E Lemma3.2 enumerates tuples of integer measurement values in $[-D,D]$, scans Boolean assignments to find one representative in each nonempty local fiber, and then scans the same fibers again. To evaluate a factor assigned to group $g$, it combines the actual group assignment with fixed representatives in all other groups. Every intermediate assembled vector is Boolean. Successively replacing an off-target group's representative by another assignment in the same integer fiber leaves that factor unchanged by the property just proved. Distributivity then gives the exact product of group sums and sum over measurement tuples. The proof never needs invariance on arbitrary vectors over $L$, and it never replaces integer measurement equality by equality modulo $p$. Thus its vector count, storage bound and coefficient-independent control remain exactly those in Lemma3.2.

Let $\mathcal A=L[y]$, with the factor weights retained as formal independent indeterminates, and use the standard multilinear factor interpolation
\[
 \Phi_f(z)=\sum_{a\in\{0,1\}^{I_f}}y_{f,a}
            \prod_{i:a_i=1}z_i\prod_{i:a_i=0}(1-z_i).
\]
Its value at a Boolean local input is the corresponding indeterminate as a formal polynomial, in every characteristic. Define the polynomial $H(T)$ using the new replacement rows as in E Section3.1. Then $H(0)=Z_H(y)$ exactly in $\mathcal A$, and $\deg_T H\leq dR$ with $d\leq wm=O(n)$.

Set $q=\lfloor dR/Q\rfloor$ and $J=q+1$. At node $\alpha_t$ write
$\Delta P(z)=P(\alpha_t+z)-P(\alpha_t)$. Its entries are divisible by $z$. Introduce the independent Boolean-substitution indeterminates $X_i$ used in E's auxiliary polynomial $\mathcal G_x(X,z)$. Every occurrence of $X_i$ carries a factor $z$, so $\deg_X[z^j]\mathcal G_x\leq j$. This is a polynomial degree statement over $\mathcal A$ and has no characteristic restriction.

E Lemma3.3 is explicitly an identity over every commutative ring. Its proof replaces positive powers of each $X_i$ by $X_i$ and performs subset Mobius inversion with coefficients $1,-1$. Apply it coefficientwise to $[z^j]\mathcal G_x$, for $j<J$. For each fixed substitution $X=\mathbf1_U$, the corrections involving $\Delta P$ are independent of the enumerated Boolean assignment $x$. Hence the integer-fiber separation just proved computes the corresponding restricted sum with $x_i=1$ for $i\in E$. The resulting identity is the exact Hasse jet
\[
                      H(\alpha_t+z)\bmod z^J
                         \quad\hbox{in }\mathcal A[z]/(z^J).
\]
All equalities are formal polynomial identities in the weight indeterminates $y$. We have not compared numerical evaluations at finitely many $y$ values.

Finally the ideals $(T-\alpha_t)^J$ are pairwise comaximal in $L[T]$. Their differences are nonzero units. The Chinese remainder theorem gives an isomorphism
\[
 L[T]\big/\!\prod_{t=1}^Q(T-\alpha_t)^J
              \longrightarrow\prod_{t=1}^Q L[z]/(z^J),
       \qquad H\longmapsto(H(\alpha_t+z))_t.
\]
Both sides have dimension $QJ$, with the monomial and jet bases. Its invertible matrix has entries
$\binom uj\alpha_t^{u-j}$, using integer binomial coefficients reduced into $L$. This is Hasse interpolation: it divides by no $j!$ and requires no inequality $J<p$. Compute its inverse over the explicit finite field and use only its final entries as scalar constants. Tensoring this matrix identity with the coefficient algebra $\mathcal A$ recovers $H$ and $H(0)$ formally. Any internal field inversions used to find the constants disappear from the resulting arithmetic circuit. Since $QJ>dR$, the recovered polynomial is the original $H$.

The number of enumerated vectors is the same as E Proposition3.4:
\[
 2Qr\left(\sum_{s=0}^{\min(q,n)}2^s\binom ns\right)
                       (2D+1)^{2\tau n}2^{2n/r}.
\]
The combinatorial layouts and their integer measurements are unchanged. For bounded arity and linear total incidence, branch on at most $n/\sqrt{\log_2n}$ variables of incidence above
$\max(1,\lceil M_{\rm inc}/n\rceil)\lceil\sqrt{\log_2n}\rceil$, retaining declared scopes. On each branch E Lemmas4.7--4.8 apply with the same $Q,R,r,D,\tau$ and the same three exponential costs. The field and matrix computations introduce polynomial overhead only. Bounded exceptional input sizes are handled by direct expansion. Summing all branches yields $S(n)$ operations in the coefficient algebra, with construction bit cost $S(n)\operatorname{poly}(\log p)$.

Unroll these operations into a shared division-free circuit over $L$. Every branch, representative and layout choice depends only on discrete scopes, Boolean bits and integer measurement data. No test of whether a symbolic polynomial is zero occurs. Replace each $L$-valued gate by its $e$ coordinates over $\mathbb F_p$. Addition is componentwise and multiplication uses the explicit structure constants for the quotient basis, at polynomial overhead in $e$. A weight input $y$ is represented by $(y,0,\ldots,0)$. The final computed element is precisely the universal partition polynomial, whose coefficients already lie in $\mathbb F_p$. Its coordinate of one is therefore exactly $Z_H(y)$ and every other coordinate is the zero formal polynomial. Taking that coordinate supplies the required prime-field circuit; no identity test is needed to know the other coordinates vanish. This also proves agreement with the reduction of the original integral universal partition polynomial, not merely agreement as a function on a finite field.

NN-C029's preliminary first-violation multipliers are constant-size substitutions into these universal partition polynomials. Its at-most-eight local specializations and integral monomial division establish the ordinary identity over $\mathbb Z$; performing those same constructions over $\mathbb F_p$ reduces each polynomial identity coefficientwise. Consequently the final circuits compute the reductions of the specified integral multipliers and satisfy $\sum_i a_if_i=1$ as a formal identity. The stronger sparse-CNF identity and affine coordinate formulas transfer for the same reason. For example the mutually inverse maps remain
\[
 z_i=y_i+a_i(x)(2t-t^2),\qquad
 t=\sum_i f_i(x)z_i,\qquad
 y_i=z_i+a_i(x)(t^2-2t),
\]
with coefficients reduced in the field. In characteristic two the terms containing $2t$ vanish, and the directly verified identities still hold. Original-$n$ versus $n+s$ verifier parameters are exactly those already proved in NN-C029.

To obtain the constant-free formal-degree claim, apply homogeneous-component truncation through the known output degree $D_0$ in the polynomial indeterminates. It has polynomial overhead and computes the same output because no division by an indeterminate is present. Every degree-zero component is a prime-field constant and can be evaluated explicitly. A multiplication with such a scalar can be implemented by binary doubling and addition of its other input, using $O(\log p)$ addition gates; this does not increase formal degree. Each nonconstant multiplication combines positive-degree components, so induction bounds its formal degree by the sum of their variable degrees. Constant outputs themselves can be built from one by additions, with formal degree one. Replacing all scalar multiplications in this manner, and using constant leaves only $0,1,-1$, gives formal degree at most $\max(D_0,1)$ and polynomial overhead in $\log p$. Take $D_0=m$ for the universal partition polynomial, $D_0=v$ for multipliers, and $D_0=v+1$ for the Hilbert-like IPS polynomial. This argument uses reduction in the prime field; it does not convert the original rational circuits into integral-gate circuits over $\mathbb Z$.
\end{proof}
\paragraph{Status and complexity impact.}
Complete integral-matrix, modified finite-characteristic routine, formal-output descent and ideal-certificate proofs. Complete deduction-proof review by root: passed; imported input theorems are not claimed independently proved. This closes the fixed-small-characteristic gap left by NN-C016 and the large-characteristic-only addendum of NN-C029. It gives formal symbolic partition circuits, ordinary ideal certificates and empty-case affine trivializations over fixed F2 at E's subexponential source-parameter bound. The proof does not merely reduce rational interpolation constants modulo two: it changes the nodes and uses integral Boolean measurement fibers. Constant-free prime-field circuits of polynomial formal degree follow after homogeneous-component processing. NN-C034 subsequently supplies integral-gate circuits over Z as well. Polynomial-size IPS certificates, efficient PIT verification and a VP/VNP collapse remain unproved. Exponential monomial-PC lower bounds are still compatible with the resulting shared circuits.

\paragraph{Geometric and representation dictionary.}
For every prime $p$, let $B=\mathbb F_p[y]$ and let
\[
 \mathscr A=B[x_1,\ldots,x_n]/(x_i^2-x_i),\qquad
       a=\prod_f\Phi_f(x).
\]
This is a finite free split algebra of rank $2^n$, with the Boolean orthogonal idempotents as a basis. Its multiplication trace $\operatorname{Tr}_{\mathscr A/B}(a)$ is the universal partition polynomial, because multiplication is diagonal in that basis. The theorem constructs a prime-field arithmetic circuit for this relative trace in size $S(n)$, naturally under arbitrary substitutions of the parameters. The auxiliary scalar extension is finite separable and free; the proof computes the trace after that extension and recovers its original scalar coordinate by an explicit basis. This descent is an identity of regular functions on the full parameter scheme, not pointwise interpolation on its rational points.

For an unsatisfiable local verifier, the same construction supplies an ordinary ideal identity in $\mathbb F_p[x]$ and polynomial maps identifying its associated smooth cubic with affine space over every extension field of characteristic $p$. The maps and inverses are explicitly represented by subexponential circuits and have the degree bounds from NN-C029. Thus the characteristic-two geometric trivialization gate is also closed for this family. The exact classical finite-locally-free trace locator is EGA IV4,18.2.1--18.2.2, printed p111, especially base-change equation18.2.1.2. The integral-matrix identity, change of interpolation nodes, integer-fiber separation and coordinate descent are proved here; no additional Grothendieck archive locator is asserted. This family-level effective trace theorem supplies neither a canonical GCT multiplicity realization nor general polynomial-time geometric recognition.



\section{NN-C034 --- Cyclotomic interpolation gives integral constant-free symbolic circuits over every ring}
\begin{proposition}[NN-C034]
Fix factor arity $w$ and density $C$. The full integral universal partition polynomial
\[
 Z_H(y)=\sum_{x\in\{0,1\}^n}\prod_f y_{f,x|I_f}
                              \in\mathbb Z[y]
\]
of every $n$-variable, $m\leq Cn$-factor instance has a deterministically constructible division-free arithmetic circuit with integer constants, of size and bit construction time
\[
 S(n)=\operatorname{poly}(n)\,
       2^{O_{w,C}(n(\log\log n)^2/\sqrt{\log n})}.
\]
All integer constants have polynomial bit length. The circuit can be made constant-free, with only $0,1,-1$ as constant leaves, at polynomial overhead and with polynomial syntactic formal degree. For the universal polynomial one may take formal degree at most $\max(m,1)$ when subtraction gates are allowed and have the same degree rule as addition.

More generally, let $m\leq Cn$ factors be explicit monomial lists $\Phi_f=\sum_j y_{f,j}x^{\alpha(f,j)}$ of full bit length $L$, with independent coefficient blocks, declared scope at most $(\lceil\log_2n\rceil)^a$ and total input-degree budget $\sum_f\max_j|\alpha(f,j)|\leq Cn$, for fixed $C,a$. Their symbolic sum-product has size and bit construction time $\operatorname{poly}(n,L)2^{O_{C,a}(n(\log\log n)^2/\sqrt{\log n})}$, with the same integral and constant-free conclusions and formal weight degree at most $\max(m,1)$. This includes fixed-degree polylogarithmically sparse factors and preserves the original $n$ Boolean variables.

As a concrete application, the symbolic permanent of an $r\times r$ matrix with at most $Cr(\lceil\log_2r\rceil)^a$ allowed entries in total, for fixed $C,a$, has such an integral constant-free circuit in the original matrix-order parameter $r$, of formal degree at most $2r$. No individual row- or column-degree restriction is needed.

This is an identity over $\mathbb Z$, so the very same circuit computes the universal partition polynomial after base change to \emph{every} commutative ring with identity. No characteristic choice, finite-field construction, modulus factorization or rational denominator remains in that circuit.

Consequently the ordinary integral Nullstellensatz multipliers of NN-C029, and its mutually inverse integral affine-space coordinate maps in the empty case, have deterministically constructible integral-gate and constant-free circuits at the same subexponential parameter bound, with polynomial formal degree. For local verifiers the parameter is $v=n+s$ including wires; for fixed-width linear-clause CNFs the stronger original-$n$ certificate still omits Boolean axioms. These are ordinary polynomial identities over $\mathbb Z$, not only separate certificates over each finite field.

On the fixed-density 3-CNF family of NN-C021, this yields an exponential ratio between ordinary monomial-PC proof size and circuit-IPS proof size over the same field $\mathbb F_2$: the former is $2^{\Omega(n)}$, while the latter is at most $S(n)=2^{o(n)}$, with $n$ the original CNF variable count. The circuit certificates retain actual degree at most $n+1$ and polynomial formal degree.
\end{proposition}
\paragraph{Exact inputs.}
NN-C033. integral swap-matrix generating identity and the complete integer-measurement-fiber Taylor routine over a coefficient ring.

NN-C016 and NN-C029. universal bounded-arity partition polynomial; ordinary integral ideal multipliers, sparse-CNF identity and affine coordinate maps.

NN-C021. fixed-density mod-three Tseitin 3-CNF family with exponential monomial-PC lower bounds over F2; full repaired source reduction and size-degree deduction.

E: The Exponential Time Hypothesis is False, supplied LaTeX. Section3.1 independent-incidence total degree; complete Lemmas3.2--3.3, Proposition3.4 and Lemmas4.1--4.8; same discrete layouts, exact parameter checks and three exponential costs.

Servedio--Wan, Computing sparse permanents faster, author manuscript18 June2005. Section2, printed p3, first display: Rysers inclusion-exclusion identity; original definition p1.

Erdos, Beweis eines Satzes von Tschebyschef, Acta Litt. Sci. Szeged5(1932),194--198. Section5, printed pp197--198: for every integer n>1 a prime lies strictly between n and2n; Q=1 handled directly.

EGA IV, fourth part, Publications Mathematiques IHES32 (1967). Printed p.111,18.2.1--18.2.2; base-change equation18.2.1.2; primary PDF page110 (zero-based109).

Integral interpolation construction below. Explicit two-prime cyclotomic quotient, integral unit inverses, Hasse cardinal polynomials and coefficient-size proof; proposed by root and developed here.

\begin{proof}
The argument uses the integral matrix and exact integer measurement fibers proved in NN-C033. The new step is to carry out all interpolation inside an explicit finite free integral algebra with sufficiently many nodes whose differences are integral units. Coordinate simulation will then return to $\mathbb Z$.

Let $Q\geq2$ be the number of interpolation nodes in one layout. Choose primes
\[
                   Q<p<2Q<q<4Q.
\]
Bertrand's theorem, in the exact primary source cited above, supplies them. Trial division of every candidate up to $4Q$ finds them deterministically in polynomial time in $Q$; no fast primality theorem is needed. For $Q=1$ take $p=2,q=3$. Put $N=pq$ and
\[
 \Phi_N(T)=\frac{(T^N-1)(T-1)}{(T^p-1)(T^q-1)},\qquad
 a=(p-1)(q-1),\qquad A=\mathbb Z[\zeta]=\mathbb Z[T]/(\Phi_N(T)).
\]
Here the notation for $A$ means this explicit quotient; no computation of a ring of integers is required. We verify the polynomial facts used about it.

Over $\mathbb C$, the denominator's roots occur in the numerator with at least the same multiplicities: the common root one has multiplicity two, and the remaining $p$th and $q$th roots occur once. Thus the quotient is a monic polynomial of degree $a$. Polynomial division by a monic integer polynomial takes place over $\mathbb Z$, and its complex remainder is zero, proving $\Phi_N\in\mathbb Z[T]$. Its roots are exactly the primitive $N$th roots, each with multiplicity one, because $N$ is the product of the two distinct primes. Its value at one is
\[
                    \Phi_N(1)=\frac{N}{pq}=1,
\]
by canceling the two factors $T-1$, or equivalently by evaluating the resulting geometric sums. Also $\Phi_N$ divides $T^N-1$ over $\mathbb Z$, so $\zeta^N=1$ in $A$. The algebra $A$ is free over $\mathbb Z$ with the explicit basis
$1,\zeta,\ldots,\zeta^{a-1}$, of rank $a=O(Q^2)$. In particular its scalar inclusion is injective. Irreducibility of $\Phi_N$ is not needed anywhere in the argument.

If $\gcd(d,N)=1$, taking $d$th powers permutes the primitive roots. Therefore $\Phi_N(T)$ divides $\Phi_N(T^d)$ for a nonnegative representative of $d$ modulo $N$. Monic division again shows the divisibility over $\mathbb Z$. Hence $\Phi_N(\zeta^d)=0$ in $A$. Define the integral polynomial
\[
                     U(T)=\frac{\Phi_N(T)-1}{T-1}.
\]
It is integral because its numerator vanishes at one. The identity
\[
                   (1-\zeta^d)U(\zeta^d)=1
                 \qquad(\gcd(d,N)=1)
\]
proves that $1-\zeta^d$ is an integral unit with a displayed inverse.

Take the $Q$ nodes $\alpha_s=\zeta^s$, $1\leq s\leq Q$. Each is a unit because $\zeta^N=1$. For $i\ne j$, the nonzero integer $i-j$ has absolute value less than $Q<p,q$, so it is coprime to $N$. Thus
$\alpha_i-\alpha_j=\zeta^j(\zeta^{i-j}-1)$ is also a unit, with an inverse obtained from the displayed formula and a power of $\zeta$. All these nodes and inverses are elements of the original integral quotient $A$; no localization has occurred. The quotient may ramify after reduction at $p$ or $q$, but this has no effect on these unit identities.

The integral swap matrix $M_0$ from NN-C033 now gives rows
\[
 \varphi_{u,y}(T)=\prod_{s\in E_{u,y}}(1-T/\alpha_s),\qquad
 P_u(T)=\delta_u(I-M_0)+\sum_yM_0[u,y]\varphi_{u,y}(T)\delta_y
                             \quad\hbox{in }A[T]^{1\times D}.
\]
Division by $\alpha_s$ here means multiplication by its explicit integral unit inverse. Their degree and zero-at-mismatched-node properties are unchanged. The exact integer-kernel identity for $M_0$ gives invariance on the same integer Boolean measurement fibers after mapping integers into $A$. This is all that the fiber proof of E Lemma3.2 uses. The ring-valued Boolean-difference identity of Lemma3.3, fixed substitutions, and Taylor truncations therefore compute all jets
\[
                  h_i(z)=H(\alpha_i+z)\bmod z^J
                              \quad\hbox{in }A[y][z]/(z^J)
\]
with the same operation bound as NN-C033. Every equality is formal in the independent weights $y$, and $H(0)=Z_H(y)$ as an element of $A[y]$.

We give the entire final Hasse interpolation explicitly, so that there is no unexamined inverse-matrix or denominator step. Put
\[
 B_i(T)=\prod_{j\ne i}
             \left((T-\alpha_j)(\alpha_i-\alpha_j)^{-1}\right)^J,
 \qquad b_i(z)=B_i(\alpha_i+z).
\]
Then $b_i(0)=1$. Its inverse modulo $z^J$ is the polynomial
$c_i(z)=\sum_{k=0}^{J-1}c_{i,k}z^k$ determined by
\[
 c_{i,0}=1,\qquad
 c_{i,k}=-\sum_{l=1}^k[z^l]b_i(z)\,c_{i,k-l}\quad(1\leq k<J).
\]
This recurrence uses only integer-algebra additions and multiplications. Define
\[
 \widetilde H(T)=\sum_{i=1}^Q B_i(T)
      \left([h_i(z)c_i(z)]_{<J}\right)_{z=T-\alpha_i}.
\]
Its degree is at most $(Q-1)J+J-1=QJ-1$. At node $\alpha_i$ its jet is $h_i$, because $b_ic_i=1\bmod z^J$; at every other node the factor $B_i$ vanishes to order at least $J$. The ideals $(T-\alpha_i)^J$ are pairwise comaximal in $A[T]$, since the differences are units. Their intersection equals their product, generated by the monic polynomial $\prod_i(T-\alpha_i)^J$ of degree $QJ$. A polynomial of smaller degree in this ideal is zero, even over a ring with zero divisors, by monic polynomial division. Therefore the jets uniquely determine a polynomial of degree less than $QJ$. As $\deg H<QJ$, one has $\widetilde H=H$ as a formal identity in $A[y][T]$.

In particular $H(0)$ is computed by the displayed formula using only fixed integral scalars and the jet entries. All interpolation constants are in $A$. This uses neither Gaussian pivots nor division by a derivative factorial, and remains valid after every coefficient-ring homomorphism. No assertion about equality of numerical functions is substituted for polynomial equality.

Here are explicit size and coefficient-height bounds for the scalar computations. The rank $a$ and $N$ are $O(Q^2)$. Since the roots of $\Phi_N$ have complex absolute value one, the sum of absolute values of its coefficients is at most $2^a$ by the elementary-symmetric bounds $\binom ak$. Monic division therefore constructs it in polynomial bit time. Division of $\Phi_N-1$ by $T-1$ gives $U$ of coefficient norm at most $a2^{a+1}$. Evaluating its finitely many needed powers of $\zeta$, and reducing to the monic basis, involves only polynomially many powers and operations of degree polynomial in $Q$.

More quantitatively, let $\|u\|$ be the sum of absolute coordinates of an element $u\in A$ in the monic basis, and let
\[
 K=\max\left(1,\max_{0\leq i,j<a}\|\zeta^{i+j}\|\right).
\]
Then $\|uv\|\leq K\|u\|\|v\|$. Repeated monic reduction of powers of degree at most $2a-2$, using coefficient bounds $2^a$, gives $\log_2 K=O(a^2+1)$. The unit inverses above have coordinate bit height polynomial in $Q$: reduce exponents first modulo $N$, using $\zeta^N=1$, and then reduce the polynomial of degree below $N$ in the monic basis. This requires at most $N$ monic reductions, each with coefficient growth bounded by $1+\|\Phi_N\|$, so its logarithmic norm is polynomial in $Q$.

The coefficients of $B_i(\alpha_i+z)$ are formed from $O(QJ)$ factors with these bounded norms. Extend $\|\cdot\|$ to polynomials by summing coefficient norms. The same multiplication inequality applies, so these coefficients have logarithmic norm polynomial in $Q,J$. Let $B\geq1$ bound the norms of the nonconstant coefficients of $b_i$, uniformly in $i$. If $C_k=\|c_{i,k}\|$, the recurrence gives
\[
 C_0=1,\qquad C_k\leq KB\sum_{l=0}^{k-1}C_l,
          \qquad C_k\leq KB(1+KB)^{k-1}\quad(k\geq1).
\]
Thus $\log(1+C_k)=O(J\log(2+KB))$, still polynomial in $Q,J$. The cardinal-polynomial formula then gives polynomial bit bounds for every final scalar coefficient as well. All scalar computations take polynomial bit time: they involve polynomially many operations on arrays of polynomial length and on integers of these polynomial bit heights. The same estimates apply to rows, shifted Taylor coefficients and a local factor evaluator, because their degrees and scope sizes are polynomially bounded. No integer of superpolynomial bit length is being hidden in a unit inverse.

Use the same high-incidence branches, labels, integer measurements and layout parameters as in NN-C033. Its proof gives $S(n)$ operations over $A[y]$. Each fixed-layout scalar table and each local coefficient evaluation has polynomial bit cost; repeated local evaluations are included in the $S(n)$ operation bound. Store each symbolic result as an arithmetic-gate reference. Replacing an $A$-valued gate by its $a$ coordinates gives an integer arithmetic circuit: addition is coordinatewise and multiplication uses the integer structure constants for the monic quotient basis. There are at most polynomially many coordinate gates per original operation. Their constants have the polynomial bit bounds just proved. Hence the total size and construction bit time remain $S(n)$.

The coordinate descent is exact. The polynomial initially specified as $Z_H(y)$ has integer coefficients, and every step proved that the computed element equals its image in $A[y]$. Since $A$ is free with basis beginning in one, that image has first coordinate $Z_H(y)$ and all other coordinates zero. Take the first output coordinate. This gives the original integral polynomial with integer-gate syntax, without computing its exponentially many coefficients and without a symbolic identity oracle. Its image under any map $\mathbb Z\to R$ is therefore the universal partition polynomial over the arbitrary target ring $R$.

We finally justify the constant-free and formal-degree assertions without assuming that semantic integrality alone eliminates denominators. The integral construction just obtained already has only integer scalar constants of polynomial bit length. It may be arranged homogeneous in the factor-weight blocks: every local factor is linear in its own block, a group product uses exactly its assigned factors, every group sum retains that block degree, and the final Taylor and interpolation sums have total weight degree $m$. In constructing a local evaluator, precompute its scalar Taylor coefficients and then form the linear combinations of its factor weights; all these scalar coefficients have the height bounds above. Empty group products contribute one, and their sums are integer counts at most $2^n$, also of polynomial bit length. Thus all degree-zero operations may be performed as explicit scalar preprocessing. After coordinate simulation, each remaining product either combines positive-degree homogeneous pieces or multiplies one piece by a known integer scalar. Its degree bookkeeping is known from this construction and requires no identity test.

For an integer scalar of bit length $L$, multiplication of a symbolic input by that scalar can be implemented with $O(L)$ doubling and addition/subtraction gates. These gates do not increase formal degree. Integer constant outputs themselves can be built from $1$ and $-1$ by such additions, with formal degree one. Replacing all scalar multiplications this way leaves only products of positive-degree pieces. Induction then bounds formal degree by the positive weight degree, and hence by $m$; the constant case uses degree one. This gives the stated constant-free bound at polynomial overhead. If one's convention permits only addition and multiplication, maintain both signs of each value; the same construction has formal degree at most $2\max(m,1)$, which is still polynomial.

Apply these integral circuits to the exact universal-partition specializations in NN-C029. Preliminary multipliers, their constantly many local specializations, and their fixed-scope integral division corrections now all use integral gates. Their ordinary identity over $\mathbb Z$ is the already-proved first-violation and local-remainder identity, with exactly the same specified multiplier polynomials. Each substituted weight is a constant or a degree-one expression, and each local correction has bounded degree and size, so the constant-free formal-degree bound stays $O(v)$, even without identifying it with the sharper actual degree bound $\deg a_i\leq v$. The IPS polynomial adds its placeholder variables with one multiplication per term. The explicit affine coordinate maps use only these multipliers and the displayed low-degree formulas of NN-C029, preserving subexponential circuit size and polynomial formal degree. This proves all assertions over $\mathbb Z$ and hence after arbitrary base change.

\paragraph{Extension to sparse polynomial factors with growing scopes.}
The same integral construction applies beyond constant arity.
Fix integers $C,a\ge1$.  Suppose there are $n\ge2$ Boolean variables
and $m\le Cn$ factors, each given by an explicit monomial list
\[
 \Phi_f(x)=\sum_{j=1}^{\ell_f}y_{f,j}x^{\alpha(f,j)},
 \qquad d_f=\max_j|\alpha(f,j)|,
 \qquad \sum_f d_f\le Cn.
\]
The coefficient symbols $y_{f,j}$ are independent; repeated
monomials in a list are permitted.  An empty list means the zero
factor and is handled immediately.  Let $L$ be the full bit length
of these lists, including their variable indices and binary
exponents.  Suppose each declared scope $I_f$ has size at most
$\lambda^a$, where $\lambda=\lceil\log_2 n\rceil$.
The degree bound is imposed on these explicitly supplied
polynomials before Boolean reduction.  Then
\[
 Z_{\Phi}(y)=\sum_{x\in\{0,1\}^n}\prod_f\Phi_f(x)
\]
has an integral division-free circuit of size and deterministic
bit construction time
\[
 \operatorname{poly}(n,L)\,
 2^{O_{C,a}(n(\log\log n)^2/\sqrt{\log n})}.
\]
Its integer constants have polynomial bit length in $n,L$.
It admits the same constant-free compilation of formal degree
at most $\max(m,1)$, with subtraction gates assigned the
addition degree rule.  The output is the formal polynomial
over $\mathbb Z[y]$, so this single circuit identity transfers
to every commutative coefficient ring.

In particular, fixed-degree factors with a polylogarithmic
number of monomials satisfy these hypotheses after changing
the fixed constants: if the degree is at most $k$ and the
list length at most $(\log n)^b$, the scope is at most
$k(\log n)^b$, and the total degree is at most $kCn$.
This preserves the original number $n$ of Boolean variables.
There are no auxiliary verifier variables, and no truth table
on the entire growing scope is constructed.

We prove the exact quantitative interface.  Append one
dummy factor equal to one, with declared scope $\{i\}$, for
each variable $i$.  These factors contribute no coefficient
symbols and have input degree zero; they only ensure that
every variable has a declared incidence.  If $M_0$ is the
resulting total scope length, then
\[
              n\le M_0\le n+Cn\lambda^a.
\]
Put
\[
 B_0=\max(1,\lceil M_0/n\rceil)\lceil\sqrt\lambda\rceil,
 \qquad H=\{i:\deg_{\rm inc}(i)>B_0\}.
\]
Double-counting incidences gives
$|H|\le M_0/B_0\le n/\sqrt\lambda$.
Sum over all $2^{|H|}$ assignments to these variables.  On a
branch, substitute the chosen bits in the explicit monomial
lists and remove them from each declared scope.  Keep the
remaining declared scopes even if some monomial becomes
zero.  This uses only the syntactic list and chosen bits,
with no test of a formal polynomial identity.

Let $n'=n-|H|$.  For sufficiently large $n$ one has
$n/2\le n'\le n$.  Each surviving variable has incidence at
most $B_0$, every scope still has size at most $\lambda^a$,
and the dummy factors ensure $n'\le M'\le B_0n'$ for the
remaining total incidence $M'$.  Factors with empty scope
can be multiplied outside the branch's counting routine;
their polynomials in the $y$ variables are explicit linear
forms.  This removal does not change $M'$.

Choose a fixed integer $\beta$ larger than $a+2$, increasing
the threshold of the exceptional finite range to absorb
$C$.  With $\lambda'=\lceil\log_2 n'\rceil$, both the maximum
scope and $B_0$ are at most $(\lambda')^\beta$ for all
sufficiently large $n$.  E Lemma~4.8, with
$\nu'=\lceil\log_2\lambda'\rceil$, therefore supplies
\[
 r=\left\lfloor\frac{\sqrt{\lambda'}}{(\nu')^2}\right\rfloor,
 \quad Q=\left\lfloor\frac{\lambda'}{128\nu'}\right\rfloor,
 \quad R=(\beta+3)\lceil\sqrt{\lambda'}\rceil,
 \quad \tau=\frac1{r\lambda'},\quad D=r^{2Q}.
\]
Its exact numerical checks and E Lemma~4.7 construct the
layout in polynomial time.  Input sizes failing those
checks form a finite range depending only on the fixed
constants; direct expansion handles that range within
the asserted asymptotic bound.

The independent-incidence degree in E Section~3.1 is bounded
by the sum of the factor degrees, because replacing the
same Boolean variable in different factors by distinct
indeterminates does not increase any factor's total degree.
Consequently, on every branch,
\[
       d\le\sum_f d_f\le Cn\le 2Cn'.
\]
This is the crucial degree hypothesis; the possibly larger
number $M'$ of incidences is not substituted for $d$.
The interpolation polynomial has degree at most $dR$.
Taking $q=\lfloor dR/Q\rfloor$ and $J=q+1$ gives $QJ>dR$.
E Proposition~3.4 has the explicit enumeration bound
\[
       2Qr\,C_q(2D+1)^{2\tau n'}2^{2n'/r},\qquad
       C_q=\sum_{s=0}^{\min(q,n')}2^s\binom{n'}s.
\]
Put $\eta=R/Q$.  The source's finite-difference estimate
$\eta^qC_q\le(1+2\eta)^{n'}$ yields
\[
 \log_2 C_q
 \le q\log_2(Q/R)+n'\log_2(1+2R/Q)
 =O_C\!\left(n'\frac RQ\log\frac{2Q}{R}\right).
\]
Together with E Lemma~4.8 this bounds all three exponential
costs by
\[
 O_{C,a}\!\left(\frac{n'(\nu')^2}{\sqrt{\lambda'}}\right).
\]
The additional branching cost has logarithm
$|H|\le n/\sqrt\lambda$.  Since $n'=\Theta(n)$, their sum
is $O_{C,a}(n(\log\log n)^2/\sqrt{\log n})$.

It remains to verify the factor and coefficient costs,
which an arity statement alone would not justify.  After
the fixed Boolean substitutions in E Lemma~3.3, evaluate
each listed monomial directly in the truncated algebra
$A[z]/(z^J)$, using its binary exponents or at most
$\sum_i\alpha(f,j)_i\le Cn$ multiplications.  Sum these
values with their formal coefficients $y_{f,j}$.
This takes polynomially many algebra operations in
$n,J,L$; it never enumerates the $2^{|I_f|}$ local
assignments.  Before multiplication by the formal
coefficients, each such substituted monomial is a
scalar Taylor polynomial.  The rows and unit inverses
have the polynomial coordinate heights already proved
for the cyclotomic construction.  A product of at most
$Cn$ such scalar inputs, followed by Taylor truncation,
still has polynomial bit height by the same multiplicative
coordinate-norm bound.  Adding at most $L$ listed terms
increases logarithmic norm by at most $O(\log(L+1))$.
All local scalar preprocessing is therefore polynomial
in $n,L$, and $Q,R,J,D$ are polynomially bounded in $n$.

Use the integral swap matrix, integer measurement fibers,
unit-separated cyclotomic nodes and explicit Hasse
interpolation proved above.  Their correctness depends on
the layout and the degree estimate, not on constant scope
size.  They give formal identities over $A[y]$ at each
stage.  The same coefficient-independent choices of
representatives and branches allow circuit unrolling,
and taking the coordinate of one in $A$ descends the
formal output to $\mathbb Z[y]$.  Each genuine factor
remains linear in its own coefficient block.  Group
products, sums, empty-scope multiplication and the final
sum over branches retain the same block degree.  The
constant-free compilation therefore has formal degree
at most $\max(m,1)$, just as in the constant-arity case.
This finishes the construction and its bit bound.

The hypotheses delimit the extension.  A general truth
table on $\operatorname{polylog} n$ variables can have
interpolation degree equal to its scope size; $O(n)$ such
factors need not have total degree $O(n)$.  In that case
the actual $q=\lfloor dR/Q\rfloor$ must be retained and
the displayed bound has not been proved.  Nor does a
coefficient-dependent compression of a numerical equation
system automatically preserve this universal formal
polynomial in independent parameters.

\paragraph{A sparse permanent application in the matrix-order parameter.}
Fix $a\ge1$ and $r\ge2$.  Let $\Omega\subseteq[r]\times[r]$ be a prescribed
matrix support with at most $(\lceil\log_2 r\rceil)^a$ entries in
each row, and give the allowed entries independent variables
$a_{ij}$.  The matrix has zero elsewhere.  Its permanent polynomial
has a deterministically constructible integral constant-free circuit
of size and bit construction time
\[
 \operatorname{poly}(r)\,
       2^{O_a(r(\log\log r)^2/\sqrt{\log r})},
\]
with formal degree $O(r)$.  No column-degree restriction is imposed.
The same claim allows a fixed multiplicative constant in the
row-support bound, by increasing the fixed scope exponent and
handling bounded sizes directly.  The parameter is the matrix
order $r$, rather than the number of allowed matrix entries.

Import Ryser's classical inclusion--exclusion identity in the
exact form reproduced in Servedio--Wan, \emph{Computing sparse
permanents faster}, author manuscript dated June~18, 2005,
Section~2, printed p.~3, first display:
\[
 \operatorname{per}(A)=(-1)^r\sum_{S\subseteq[r]}(-1)^{|S|}
                          \prod_{i=1}^r\sum_{j\in S}a_{ij}.
\]
The primary manuscript is
\url{https://www.cs.columbia.edu/~atw12/papers/perm.pdf}.
It states the identity for real or complex matrices.  Both sides
are polynomials with integer coefficients in independent matrix
entries, so agreement for all complex entries is an identity in
$\mathbb Z[a_{ij}]$ and hence after every ring base change.
This uses that a polynomial vanishing on all complex tuples is
the zero polynomial; it does not reprove the classical formula.

For the characteristic vector $x$ of $S$ one has
$(-1)^{|S|}=\prod_{j=1}^r(1-2x_j)$.  Consequently
\[
 \operatorname{per}(A)=(-1)^r
 \sum_{x\in\{0,1\}^r}
       \left(\prod_{j=1}^r(1-2x_j)\right)
       \left(\prod_{i=1}^r\sum_{j:(i,j)\in\Omega}a_{ij}x_j\right).
\]
This is a sum-product on exactly $r$ Boolean variables with $2r$
linear factors.  Its independent-incidence degree is at most
$2r$, each row factor has polylogarithmic scope, and the unary
sign factors have scope one.  The list description has polynomial
size.  A zero row gives the zero permanent immediately, and the
empty matrix has permanent one.  Apply the growing-scope theorem
of NN-C034 to the nontrivial cases, specializing the unary factor
coefficients to $1,-2$.  These are fixed integer constants, with
$-2=(-1)+(-1)$ available in the constant-free syntax.  Its integral
descent proves the full permanent identity; it does not merely
compute its values modulo a prime.  The factor-block formal degree
is at most $2r$, and the outer sign can be implemented by subtraction
from zero.  This proves the claimed circuit and construction bounds.

An allowed column may occur in every row.  The total-incidence
branching already proved in NN-C034 removes at most
$r/\sqrt{\log r}$ such high-incidence variables before applying
the layout, so unbounded column degree introduces no extra witness
parameter.  In particular, the conclusion covers supports with
$r\operatorname{polylog} r$ allowed entries at the original matrix
order.  It makes no claim for an arbitrary dense support: its
row scopes can have size $r$, outside the hypothesis.

\paragraph{Only a total support bound is necessary.}
Fix $C,a\ge1$ and $r\ge2$.  For a prescribed support
$\Omega\subseteq[r]\times[r]$ with
\[
                     |\Omega|\le Cr(\lceil\log_2r\rceil)^a,
\]
the full permanent polynomial has a deterministically
constructible integral constant-free circuit of size and
bit construction time
\[
 \operatorname{poly}(r)\,
       2^{O_{C,a}(r(\log\log r)^2/\sqrt{\log r})},
\]
and formal degree at most $2r$, with the subtraction convention
used above.  Thus no bound on any individual row or column
degree is needed.  The original parameter is still the
matrix order $r$.

Here is the full reduction to the preceding row-scope result.
Let $\lambda=\lceil\log_2r\rceil$, set
$T=\lceil\lambda^{a+2}\rceil$, and let $H$ be the set of
rows containing more than $T$ allowed entries.  If $h=|H|$,
then counting allowed entries gives
\[
                          h\le Cr/\lambda^2.
\]
For sufficiently large $r$, this is at most $r/2$.
Expand the definition of the permanent by first specifying
the columns matched to the rows of $H$.  It gives the exact
integral polynomial identity
\[
 \operatorname{per}(A)=
 \sum_{\substack{\pi:H\hookrightarrow[r]\\(i,\pi(i))\in\Omega}}
       \left(\prod_{i\in H}a_{i,\pi(i)}\right)
       \operatorname{per}\bigl(A_{[r]\setminus H,\,[r]\setminus\pi(H)}\bigr).
\]
This is simply the partition of the permutations in the
permanent definition according to their restrictions to
$H$; each full permutation occurs in exactly one summand.
There is no sign and no multiplicity factor.  If $H$ is
empty the sum has its one empty restriction, and its
weight is one.  Enumerating at most $r^h$ functions and
retaining the allowed injections constructs the sum
deterministically with polynomial overhead per term.

Every minor has order $r'=r-h\ge r/2$.  Its remaining
row supports have size at most $T$, and deleting columns
cannot increase them.  Since
$\lceil\log_2r'\rceil\ge\lambda-1$, one has
\[
                   T\le(\lceil\log_2r'\rceil)^{a+3}
\]
for sufficiently large $r$.  The preceding row-scope
permanent construction therefore applies to every minor,
with a fixed changed scope exponent.  Zero rows can be
detected from the prescribed support and give zero terms.
The finitely many exceptional matrix orders, including
orders zero and one, are computed directly.

The logarithm of the number of restrictions is at most
\[
                   h\log_2r=O_C(r/\log r),
\]
which is absorbed by
$O_{C,a}(r(\log\log r)^2/\sqrt{\log r})$.
The minor construction has that same bound because
$r'=\Theta(r)$.  Multiplication by the $h$ selected entry
variables and summation of all terms add only the stated
overhead.  Each minor circuit has formal degree at most
$2r'$, so each term has formal degree at most
$2r'+h\le2r$, and summation takes the maximum.  All scalar
constants retain their polynomial bit lengths.  This
proves the total-support theorem as an integral formal
identity, and therefore over every commutative coefficient
ring by base change.

There is a concrete proof-size separation on the family already audited in NN-C021. Use its mod-three Tseitin 3-CNF contradictions, with $n$ denoting the number of variables of the resulting CNF and with $O(n)$ clauses. Every ordinary polynomial-calculus refutation over $\mathbb F_2$ has explicitly written monomial size at least $2^{\kappa n}$ for some fixed $\kappa>0$. The sparse-CNF specialization just proved constructs a constant-free Hilbert-like IPS circuit over $\mathbb F_2$ of size $S(n)=2^{o(n)}$ for the same clause polynomials, without introducing verifier wires or Boolean axioms into that IPS identity. If $M_{\rm PC}(n)$ and $C_{\rm IPS}(n)$ are the minimum sizes in these two specified models, then for all sufficiently large $n$,
\[
       \frac{M_{\rm PC}(n)}{C_{\rm IPS}(n)}
           \geq\frac{2^{\kappa n}}{S(n)}
           \geq 2^{\kappa n/2}.
\]
This is the precise exponential ratio. The numerator counts explicit monomials throughout an ordinary PC derivation; the denominator counts gates in a circuit IPS certificate. The latter certificate has actual polynomial degree at most $n+1$ and a constant-free circuit of formal degree $O(n)$. These degree statements neither identify the two size measures nor contradict the linear PC degree lower bound. The deduction uses the classical monomial-PC lower bound and E's designated algorithmic input, and implies no polynomially bounded proof system or $\mathrm{NP}=\mathrm{coNP}$ collapse.
\end{proof}
\paragraph{Status and complexity impact.}
Complete integral cyclotomic-node, Hasse interpolation, height, symbolic descent and certificate proof. Complete deduction-proof review by root: passed; imported input theorems are not claimed independently proved. The remaining integral-gate obstacle in NN-C029 and NN-C033 is closed for the stated symbolic families, including polynomial-list factors with polylogarithmic scopes and linear total input degree. List length and coefficient bit cost are explicit; generic growing-scope truth tables need not satisfy the degree budget. The construction is over Z itself, with polynomial-bit integer constants and a constant-free presentation of polynomial formal degree, so it transfers to every coefficient ring through one universal identity. It makes actual integral ideal multipliers and empty-case geometric coordinate maps effective at E's subexponential bound. It does not give polynomial-size circuits or IPS certificates, deterministic polynomial-time PIT, NP=coNP, or $VP^0=VNP^0$. Dense permanent encodings and long verifier circuits retain their actual edge/wire parameter overhead. The finite free cyclotomic algebra is an explicit computational extension, not an assumed canonical GCT object.

\paragraph{Geometric and representation dictionary.}
Let $B=\mathbb Z[y]$ and
\[
 \mathscr A=B[x_1,\ldots,x_n]/(x_i^2-x_i),\qquad
                       a=\prod_f\Phi_f(x).
\]
The finite locally free multiplication trace $\operatorname{Tr}_{\mathscr A/B}(a)$ is the universal partition polynomial. The theorem supplies an integral arithmetic circuit for this regular function on $\operatorname{Spec}B$ in size $S(n)$, with a constant-free presentation of polynomial formal degree. The auxiliary extension $B\to A\otimes_{\mathbb Z}B$ is finite free and faithfully flat. It contains the required distinct unit-separated interpolation sections; these differences remain units after every base change, including ramified residue characteristics. Exact interpolation and coordinate projection descend the computed regular function to $B$.

Fix $C,a\ge1$.  Let $m\le Cn$ explicit polynomial lists
$\Phi_f=\sum_j y_{f,j}x^{\alpha(f,j)}$ have description
length $L$, declared scopes of size at most
$(\lceil\log_2n\rceil)^a$, and total degree budget
$\sum_f\max_j|\alpha(f,j)|\le Cn$.  Put $B=\mathbb Z[y]$ and
\[
 \mathscr A=B[x_1,\ldots,x_n]/(x_i^2-x_i),\qquad
                       u=\prod_f\Phi_f(x).
\]
The finite locally free algebra $\mathscr A$ has rank $2^n$.
Its multiplication trace
$\operatorname{Tr}_{\mathscr A/B}(u)=Z_\Phi(y)$, as is
seen in the orthogonal Boolean-idempotent basis.  This
regular function on $\operatorname{Spec}B$ has an integral
circuit of size and bit construction time
$\operatorname{poly}(n,L)2^{O_{C,a}(n(\log\log n)^2/\sqrt{\log n})}$,
with polynomial-bit integer constants and a constant-free
compilation of formal degree at most $\max(m,1)$, with $n$
the original number of Boolean coordinates.  For every map
$B\to R$ the same circuit computes the base-changed
multiplication trace over $R$; coordinate descent from
the finite free cyclotomic extension proves the identity
before any specialization.  The exact classical trace
and base-change interface is EGA IV4, 18.2.1--18.2.2,
printed p.~111, including equation 18.2.1.2.  This is an
effective theorem for these specified regular trace
functions, not a claimed polynomial-time algorithm for
general finite algebra traces or GCT multiplicities.

For a support $\Omega\subseteq[r]\times[r]$ satisfying
$|\Omega|\le Cr(\lceil\log_2r\rceil)^a$, put
$B=\mathbb Z[a_{ij}:(i,j)\in\Omega]$ and
$\mathscr A=B[x_1,\ldots,x_r]/(x_j^2-x_j)$.  The regular
function
\[
 \operatorname{Tr}_{\mathscr A/B}\left(
     (-1)^r\prod_j(1-2x_j)
                 \prod_i\sum_{j:(i,j)\in\Omega}a_{ij}x_j\right)
                              =\operatorname{per}(A)
\]
has an integral constant-free circuit of size
$\operatorname{poly}(r)2^{O_{C,a}(r(\log\log r)^2/\sqrt{\log r})}$
and formal degree at most $2r$, valid after every base
change.  Ryser's identity identifies this signed trace,
while the heavy-row decomposition constructs it using
smaller trace functions with polylogarithmic row scopes.

For the split torus with characters indexed by $\Omega$,
let $\mathcal V_\Omega$ have one basis vector of character
$\prod_i a_{i,\sigma(i)}$ per allowed perfect matching.
Its effective character is the same permanent polynomial.
The restriction partition above gives the actual direct-sum
decomposition
\[
 \mathcal V_\Omega\simeq
  \bigoplus_{\pi}
    \chi_\pi\otimes\mathcal V_{\Omega_\pi},
 \qquad \chi_\pi=\prod_{i\in H}a_{i,\pi(i)},
\]
where $\Omega_\pi$ is the remaining minor support and its
characters are included in the original torus.  Applying
the integral minor circuits to these summands gives the
stated character circuit; positivity of the decomposition
does not remove the virtual differences inside the minor
circuits.  The precise classical trace and character-ring
interfaces remain EGA IV4, 18.2.1--18.2.2, printed p.~111,
and the SGA3/SGA6 locators already stated in NN-S019.

For the unsatisfiable local-verifier scheme, the same circuits give an ordinary global ideal identity $\sum_i a_if_i=1$ over $\mathbb Z$ and the explicit mutually inverse maps identifying its smooth counting cubic with affine space. The maps have the actual degree bounds from NN-C029, integral coefficient circuits of size $S(v)$, and polynomial formal degree after removing all large constant leaves. Thus these are effective integral geometric trivializations, valid after arbitrary base change, with no denominator-prime exclusion. On the empty Boolean equation schemes of NN-C021, these circuit ideal certificates over $\mathbb F_2$ have size $2^{o(n)}$, while deriving their unit ideal by ordinary monomial polynomial calculus requires size $2^{\Omega(n)}$. This compares specified presentations of the unit-ideal proof, not different underlying zero schemes or different cohomological degrees. The exact classical multiplication-trace and base-change locator is EGA IV4,18.2.1--18.2.2, printed p111, equation18.2.1.2. All node, inverse, flat scalar-extension and polynomial-identity calculations are proved above; no unverified Grothendieck archive cote or GCT canonical-basis functor is invoked.



\section{NN-C035 --- Exact zero-set compression yields formal weighted CNF circuits in characteristic two}
\begin{proposition}[NN-C035]
Fix a clause width $k\ge2$. For any $k$-CNF $F$ on $n$ declared variables with $m$ clauses, the full satisfying-assignment polynomial
\[
 A_F(z)=\sum_{x\models F}\prod_{i=1}^n z_{i,x_i}
                       \in\mathbb F_2[z_{1,0},z_{1,1},\ldots,z_{n,0},z_{n,1}]
\]
has a deterministically constructible constant-free arithmetic circuit of size and bit construction time
\[
 (n+m)^{O_k(1)}\,
      2^{O_k(n(\log\log n)^2/\sqrt{\log n})}.
\]
Its formal degree is $O_k(n)$. The parameter is the original number of Boolean variables; no verifier wires are added. This is equality of formal polynomials, valid after base change to every commutative ring of characteristic two.

More generally, one may replace the product of the $n$ displayed unary weights by at most $Cn$ explicitly listed polynomial factors in $x$, with independent coefficient parameters, polylogarithmic scopes, and total input-degree budget at most $Cn$, for fixed $C$ and scope exponent. The same bound holds with a polynomial factor in their description length and with the constants depending on these fixed parameters.

In particular, for every graph $G$ on $n$ vertices, with no sparsity or maximum-degree restriction, its full multivariate independence polynomial
\[
                       I_G(y)=\sum_{S\text{ independent}}\prod_{i\in S}y_i
                                      \in\mathbb F_2[y_1,\ldots,y_n]
\]
has such a circuit in the original vertex parameter. The two-rail polynomial $A_F$ is formally nonzero exactly when $F$ is satisfiable, even though its value at all weights equal to one records only the parity of the satisfying count.
\end{proposition}
\paragraph{Exact inputs.}
E: The Exponential Time Hypothesis is False, supplied LaTeX. Lemma6.3 disjoint sparsification on the original declared variables; Lemmas6.1--6.2 and Proposition6.4, especially the exact no-false-survivor conclusion.

NN-C001. deterministic compression of quasilinear sparse F2 equation lists to at most6n equations, preserving their full Boolean zero set.

NN-C034. reviewed growing-scope extension: explicit polynomial lists, polylogarithmic scopes, linear total input degree, and one formal integral circuit identity over every coefficient ring.

EGA IV, fourth part, Publications Mathematiques IHES32 (1967). Printed p.111,18.2.1--18.2.2; base-change equation18.2.1.2; primary PDF page110 (zero-based109).

\begin{proof}
We first retain the exact scope of the imported compression theorem. Its rows are selected by numerical zero counts for the supplied Boolean equations. They depend on $F$, but they do not depend on any of the formal weight parameters introduced here. NN-C001 proves equality of the whole Boolean zero set, not just equality of its cardinality. That stronger assertion is what permits symbolic weights.

Handle an empty clause by outputting zero. A formula with no clauses has the explicit circuit $\prod_i(z_{i,0}+z_{i,1})$ in the two-rail case. Bounded exceptional sizes may be expanded directly. For the general case put $\lambda=\lceil\log_2 n\rceil$ and apply E Lemma6.3 with the rational parameter
\[
                         \sigma=1/\lceil\sqrt\lambda\rceil.
\]
The lemma produces a tree of at most $\operatorname{poly}_k(n)2^{O(n/\sqrt{\log n})}$ nodes on the same $n$ declared variables. Its leaf satisfying sets partition the satisfying set of $F$. In every leaf each variable has occurrence at most
\[
 B=O_k\bigl([\sigma^{-1}\log(2/\sigma)]^{k-1}\bigr)
                                  =(\log n)^{O_k(1)}.
\]
The input normalization costs $\operatorname{poly}_k(n,m)$ and the numerical parameter has $O(\log\log(n+2))$ bits. Thus every nontrivial leaf has $n(\log n)^{O_k(1)}$ clauses, while the number of Boolean variables is still $n$.

For a leaf $F_t$, take the usual clause-violation polynomial
\[
                  v_C(x)=\prod_{\ell\in C}(1-\ell(x))
                                         \quad\hbox{over }\mathbb F_2,
\]
where the affine literal polynomial has value one exactly when that literal is true. A clause has at most $k$ literals, so its expanded violation polynomial has degree at most $k$ and at most $2^k$ monomials. It vanishes exactly when the clause is satisfied. Apply NN-C001 to these equations. The resulting list
\[
                   p_{t,1},\ldots,p_{t,b_t},\qquad b_t\le6n,
\]
has degree at most $k$, polylogarithmic monomial sparsity and hence polylogarithmic scopes, and satisfies
\[
 x\models F_t\quad\Longleftrightarrow\quad
              p_{t,j}(x)=0\text{ for every }j
                     \qquad(x\in\{0,1\}^n).
\]
Its construction takes $\operatorname{poly}_k(n)2^{O_k(n\log\log n/\sqrt{\log n})}$ bit operations. This uses only the numerical equations, so all coefficient-parameter symbols remain untouched.

Now work in the polynomial coefficient algebra $\mathbb F_2[z]$. For a Boolean assignment, each $p_{t,j}(x)$ is either zero or one. Therefore
\[
                e_t(x)=\prod_{j=1}^{b_t}(1+p_{t,j}(x))
\]
is exactly the indicator of $F_t$ on the Boolean cube. Use the unary interpolation factors
\[
                   W_i(x_i)=(1-x_i)z_{i,0}+x_i z_{i,1}.
\]
The sum-product
\[
              Z_t(z)=\sum_{x\in\{0,1\}^n}e_t(x)\prod_iW_i(x_i)
\]
therefore equals $\sum_{x\models F_t}\prod_i z_{i,x_i}$ as a formal polynomial over $\mathbb F_2$. This assertion follows term by term before any evaluation of $z$.

There are at most $7n$ factors. The sum of their degrees in the Boolean inputs is at most $(6k+1)n$, and every scope is polylogarithmic. Each factor is explicitly listed, and the list lengths are polynomial in $n$. These are precisely the hypotheses of the growing-scope part of NN-C034. To use its universal coefficient form, first place independent formal coefficients on each monomial occurrence of these factors. Construct the universal integral circuit supplied there. Then apply the coefficient-ring homomorphism to $\mathbb F_2[z]$ which sends the constraint coefficients to their known zero-one values and the weight coefficients to the displayed linear parameters. Thus the descended circuit computes $Z_t$ as a formal polynomial. The intermediary integral universal identity is being specialized to characteristic two; the characteristic-two indicator equality is not asserted as an integer equality.

The block compilation in NN-C034 gives a constant-free circuit of formal degree $O_k(n)$ after these specializations: each genuine coefficient input becomes a constant or a linear expression in a weight parameter, and only $O_k(n)$ factor blocks are used. In characteristic two subtraction is addition. The resulting size and construction time per leaf are bounded by $\operatorname{poly}_k(n)2^{O_k(n(\log\log n)^2/\sqrt{\log n})}$. All representation and control choices are discrete or depend on the known clause coefficients; there is no symbolic identity test and no evaluation/interpolation in the weight parameters.

Sum these circuits over the sparsification leaves. Disjointness is essential: every satisfying assignment occurs in exactly one leaf, and every nonsatisfying assignment in none. Consequently the sum is exactly $A_F(z)$, with no multiplicity introduced by the decomposition. Adding the leaf-count exponent $O_k(n/\sqrt{\log n})$, and the smaller compression-construction exponent, preserves the displayed bound. An additive sum of the leaf circuits does not increase formal degree.

For the general weighted statement, replace $\prod_iW_i$ by the supplied factor product. The compressed constraint factors add at most $6kn$ to its total input degree and at most $6n$ to its factor count. The maximum scope remains a fixed power of $\log n$, and direct list evaluation has polynomial cost in the stated description length. NN-C034 therefore applies with changed fixed constants. The same indicator identity and disjoint leaf sum prove the formal weighted equality. In this version distinct assignments can yield the same weight monomial, whose coefficient correctly counts their contributions modulo two.

For the graph application take
\[
                         F_G=\bigwedge_{uv\in E(G)}(\neg x_u\vee\neg x_v).
\]
Its satisfying assignments are exactly the characteristic vectors of independent vertex sets. It has $n$ variables and $|E(G)|$ width-two clauses; no edge variable is introduced. Specializing $z_{i,0}=1$ and $z_{i,1}=y_i$ gives $I_G(y)$. The construction time is polynomial in $n+|E(G)|$ times the displayed original-$n$ exponential.

Finally, different assignments give different two-rail monomials, each with coefficient one. Hence $A_F$ is the zero formal polynomial if and only if the satisfying set is empty. The specialization $z_{i,b}=1$ collapses all these monomials to the same constant and records only their sum modulo two. These are different readout operations. The construction is mildly subexponential, not polynomial, so even a polynomial-time PIT algorithm applied to its output would not by itself yield a polynomial-time SAT algorithm. Nor does a circuit over $\mathbb F_2$ for a zero-one-coefficient polynomial automatically provide an integral circuit for its chosen zero-one coefficient lift: reduction loses multiples of two, and no coefficientwise lifting operation is supplied here. NN-C036 separately constructs the integral polynomial for this CNF family by slower disjoint sparsification and exact integer clause indicators, with a weaker exponent.
\end{proof}
\paragraph{Status and complexity impact.}
Complete characteristic-two symbolic-weight deduction proved. Complete deduction-proof review by root: passed; imported input theorems are not claimed independently proved. This closes the bounded-density restriction for formal weighted fixed-width CNF counting in characteristic two while retaining the original witness variables. Numerical compression is legitimately reused because it preserves the entire Boolean solution set and ignores weight parameters. The result is stronger than a parity value at a single specialization and weaker than a polynomial-size circuit theorem. This characteristic-two construction does not itself give a characteristic-zero circuit for the chosen zero-one coefficient lift; NN-C036 subsequently supplies a direct integral construction for this family at a weaker exponent. Neither result gives general PIT derandomization or a VP/VNP collapse. The dense-graph independence-polynomial example has an explicit original-vertex bound, unlike an encoding that charges every graph edge as a witness bit.

\paragraph{Geometric and representation dictionary.}
Let $Y_F$ be the finite etale subscheme of $\operatorname{Spec}(\mathbb F_2[x_1,\ldots,x_n]/(x_i^2-x_i))$ cut out by the clause-violation equations. Put $B=\mathbb F_2[z_{i,b}]$ and $\mathscr A_F=\mathcal O(Y_F)\otimes_{\mathbb F_2}B$. It is finite locally free over $B$, with one rank-one summand for every satisfying assignment. For
\[
                           w=\prod_i\bigl((1-x_i)z_{i,0}+x_i z_{i,1}\bigr),
\]
the multiplication trace $\operatorname{Tr}_{\mathscr A_F/B}(w)$ is $A_F(z)$ in the component-idempotent basis. For fixed clause width, this regular function has a deterministically constructible constant-free circuit of size $(n+m)^{O_k(1)}2^{O_k(n(\log\log n)^2/\sqrt{\log n})}$ and formal degree $O_k(n)$. This statement remains true after every base change to a characteristic-two coefficient ring. The same trace theorem holds for the explicitly listed weight factors with polylogarithmic scopes and linear total input degree stated above.

The compressed equations define the same closed subscheme inside the Boolean scheme, because its product coordinate algebra has ideals determined by their vanishing components. The sparsification leaves form a disjoint union decomposition of $Y_F$, and multiplication trace adds over that decomposition. The nonzero two-rail coefficient records an actual component; trace at all weights equal to one is only its rank modulo two. Every point of $Y_F$ is $\mathbb F_2$-rational, so Frobenius on $H^0_{\mathrm{et}}(Y_F\times\overline{\mathbb F}_2,\mathbb Q_\ell)$, for $\ell\ne2$, is the identity. This weight-zero purity does not compute the symbolic trace. The exact classical multiplication-trace and base-change locator is EGA IV4,18.2.1--18.2.2, printed p111, equation18.2.1.2. No reductive-group multiplicity or canonical GCT object is asserted.



\section{NN-C036 --- Slower sparsification gives integral symbolic weights and ideal certificates for arbitrary-density CNFs}
\begin{proposition}[NN-C036]
Fix $k\ge2$. For every $k$-CNF $F$ on $n$ declared variables with $m$ clauses, the integral satisfying-assignment polynomial
\[
       A_F(z)=\sum_{x\models F}\prod_{i=1}^n z_{i,x_i}
                      \in\mathbb Z[z_{i,b}:1\le i\le n,\ b\in\{0,1\}]
\]
has a deterministically constructible division-free integral circuit of size and bit construction time
\[
 (n+m)^{O_k(1)}
  2^{O_k\!\left(n(\log\log n)^{1+1/k}/(\log n)^{1/(2k)}\right)}
                                                    = (n+m)^{O_k(1)}2^{o(n)}.
\]
The circuit can be made constant-free with polynomial formal degree. It computes the formal polynomial over $\mathbb Z$, so the same circuit works after every commutative-ring base change. The number of original Boolean variables remains $n$.

The same assertion holds for a product of $O(n)$ explicitly listed weight factors with independent coefficient parameters, polylogarithmic scopes, and linear total input degree; include a polynomial factor in their description length and allow the constants to depend on the fixed scope and degree bounds. Thus arbitrary-density fixed-width CNFs admit symbolic integral weighted counting at a mildly subexponential original-variable bound. Arbitrary rational specializations of the weights can also be evaluated within the displayed time times a polynomial in their full bit length.

In particular the full multivariate independence polynomial of every $n$-vertex graph, with no edge-density restriction, has an integral constant-free circuit of size
\[
 \operatorname{poly}(n+|E|)\,
       2^{O\!\left(n(\log\log n)^{3/2}/(\log n)^{1/4}\right)}
\]
and polynomial formal degree. The characteristic-two construction of NN-C035 retains its sharper exponent. This integral construction does not lift an arbitrary characteristic-two circuit; it uses a different decomposition of the original Boolean problem.

For every unsatisfiable normalized width-$k$ CNF, the same bound constructs ordinary integral multipliers $h_j$ with $\sum_jv_jh_j=1$ for its clause-violation polynomials $v_j$, without Boolean axioms or extra variables. One has $\deg h_j\le n-|S_j|$ for clause scope $S_j$, and the resulting Hilbert-like IPS circuit has actual degree at most $n+1$ and polynomial formal degree. The associated hypersurface $t^2-t=\sum_jy_jv_j(x)$ admits explicit integral affine-space coordinate maps of the same circuit bound, of forward degree at most $n+2$ and inverse degree at most $n+2k+2$.
\end{proposition}
\paragraph{Exact inputs.}
E: The Exponential Time Hypothesis is False, supplied LaTeX. Lemma6.3 disjoint sparsification, exact occurrence bound B=O\_k((sigma\textasciicircum{}-1 log(2/sigma))\textasciicircum{}(k-1)); Section3.1 total input degree and Proposition3.4 with the actual q=floor(dR/Q); complete Lemmas4.7--4.8.

NN-C034. integral swap-matrix and cyclotomic Hasse routine, explicit polynomial coefficient-height bounds, formal coordinate descent and constant-free block compilation; the construction is reused below with a larger actual degree parameter.

NN-C035. faster characteristic-two formal weighted construction, retained as a distinct bound.

NN-C015. constant-free coefficient-norm induction at polynomial formal degree; used for exact rational specialization bit bounds.

Grochow--Pitassi, Circuit Complexity, Proof Complexity, and Polynomial Identity Testing: The Ideal Proof System. Author JACM manuscript Definition1.1 printed draftp1:4; Lemma4.2 full proof pp1:26--1:28, equations(4)--(6); classical first-violation identity.

NN-C029. classical first-violation coefficient factorization and the integral unit-ideal affine coordinate formulas.

EGA IV, fourth part, Publications Mathematiques IHES32 (1967). Printed p.111,18.2.1--18.2.2; base-change equation18.2.1.2; primary PDF page110 (zero-based109).

\begin{proof}
The construction uses the integral clause-satisfaction polynomial instead of compressing the constraints in characteristic two. For a clause $C$, let $\ell(x)$ be the affine integer polynomial for each literal and put
\[
                   c_C(x)=1-\prod_{\ell\in C}(1-\ell(x))\in\mathbb Z[x].
\]
For every Boolean assignment this factor is exactly one when $C$ is satisfied and zero otherwise. Its degree is at most the clause width and its expanded list has at most $2^k+1$ terms. These are literal integer equalities, valid after every coefficient-ring map. An empty clause gives output zero. A formula without clauses has the explicit two-rail circuit $\prod_i(z_{i,0}+z_{i,1})$. Bounded exceptional input sizes are expanded directly.

We choose a slower sparsification parameter than the one used for the characteristic-two theorem. Put
\[
 \lambda=\lceil\log_2 n\rceil,\qquad
 \nu=\lceil\log_2\lambda\rceil,
 \qquad s=\max\{t\in\mathbb Z_{\ge0}:t^{2k}\nu^{2k+2}\le\lambda\}.
\]
For sufficiently large $n$, one has $s\ge2$ and
\[
 \frac1s=\Theta_k\!\left(
             \frac{\nu^{1+1/k}}{\lambda^{1/(2k)}}\right)=o(1).
\]
The integer $s$ is computed exactly by integer comparisons in polynomial time in $\log n$ for fixed $k$. Sizes with $s<2$ form a finite range. Apply E Lemma6.3 with $\sigma=1/s$. Its leaf satisfying sets partition that of $F$, its tree has $\operatorname{poly}_k(n,s)2^{n/s}$ nodes, and its explicit leaf occurrence bound satisfies
\[
 B=O_k\bigl((s\log(2s))^{k-1}\bigr)
       =O_k(s^{k-1}\nu^{k-1}).
\]
We may increase $B$ to at least one. The choice of $s$ implies
\begin{equation}\label{eq:C036-balance}
     \frac{B\nu^2}{\sqrt\lambda}
       =O_k\!\left(\frac{s^{k-1}\nu^{k+1}}{\sqrt\lambda}\right)
       =O_k(1/s),
\end{equation}
because $s^k\nu^{k+1}\le\sqrt\lambda$. This estimate balances the cost per leaf against the number of leaves.

Fix a nontrivial leaf $F_t$. It has at most $nB$ nonempty clauses and total clause-incidence length at most $nB$. Its product of satisfaction factors has independent-incidence input degree at most $nB$, since each factor's degree is bounded by its number of literals. Multiply by the $n$ unary factors
\[
                       W_i(x_i)=(1-x_i)z_{i,0}+x_i z_{i,1}.
\]
The desired leaf polynomial is exactly
\[
 Z_t(z)=\sum_{x\in\{0,1\}^n}\left(\prod_{C\in F_t}c_C(x)\right)
                                      \prod_iW_i(x_i)
                          \quad\hbox{in }\mathbb Z[z].
\]
For the more general weight factors, their specified linear total-degree budget adds $O(n)$ to the same bound. In either case the number of factors is $O(n(B+1))$, the total input degree is $d=O(n(B+1))$, and the total list length is polynomial in $n$ and the supplied weight-list length. The crucial difference from the literal theorem statement of NN-C034 is that $B$ now grows. We therefore reuse its construction with the actual $d$, rather than invoking its linear-degree corollary outside its hypotheses.

The clause scopes are constant, and their maximum variable incidence is at most $B$. General weight scopes may be polylogarithmic and have uneven incidences. If needed, append the degree-zero dummy factors and branch on high-incidence variables exactly as in NN-C034, using the full total scope length $M$ and threshold
\[
                   \max(1,\lceil M/n\rceil)\lceil\sqrt\lambda\rceil.
\]
There are at most $n/\sqrt\lambda$ branched bits. Every remaining instance has $n'=n-O(n/\sqrt\lambda)=\Theta(n)$ variables, polylogarithmic scopes and maximum incidence, and input degree $d=O(n'B)$. The unbranched unary-weight case already has this layout regime. A fixed exponent $\beta$, depending only on $k$ and the fixed weight-scope exponent, bounds the maximum scope and incidence by $(\lceil\log_2n'\rceil)^\beta$ for all sufficiently large $n$. This is possible because $B$ and every added weight-incidence threshold are fixed powers of $\log n$.

Apply the exact layout parameters of E Lemma4.8 to $n'$, writing $\lambda'=\lceil\log_2n'\rceil$ and $\nu'=\lceil\log_2\lambda'\rceil$. They satisfy
\[
 \frac RQ=O_\beta(\nu'/\sqrt{\lambda'}),\qquad
 \log(2Q/R)=O(\nu'),\qquad
 \tau\log(2D+1)+\frac1r=O_\beta((\nu')^2/\sqrt{\lambda'}).
\]
The source constructs and checks the numerical hypotheses exactly; its finite exceptional range is handled by direct expansion. The interpolation polynomial has degree at most $dR$. Set
\[
                         q=\lfloor dR/Q\rfloor,\qquad J=q+1.
\]
E Proposition3.4 supplies the enumeration factor
$2QrC_q(2D+1)^{2\tau n'}2^{2n'/r}$, where
$C_q=\sum_{j\le\min(q,n')}2^j\binom{n'}j$.
Its finite-difference estimate with $\eta=R/Q$ gives
\[
 \log_2 C_q
   \le q\log_2(Q/R)+n'\log_2(1+2R/Q)
   =O\!\left(n'B\frac RQ\log\frac{2Q}{R}\right).
\]
Consequently all three exponential costs are
\[
           O_{k,\beta}\!\left(n'B\frac{(\nu')^2}{\sqrt{\lambda'}}\right)
                       =O_{k,\beta}(n/s)
\]
by \eqref{eq:C036-balance} and $n'=\Theta(n)$. The high-incidence branching contributes at most $n/\sqrt\lambda=O(n/s)$ to the exponent. In particular $q,J,Q,R,D$ remain polynomially bounded in $n$; no numerical parameter with exponentially long encoding has entered the construction.

We check the additional hypotheses of the integral implementation. Take the same finite free cyclotomic algebra, unit-separated nodes, integral swap matrix and exact integer measurement fibers as NN-C034. Their construction depends on the layout, not on having only $O(n)$ factors. A clause factor evaluates by its at most $2^k+1$ monomials. A supplied weight factor evaluates directly from its list. Every substituted monomial has degree at most $d=O(nB)$ and is a product of that many row Taylor inputs at worst. The coefficient-norm estimates in NN-C034 bound the bit height of each row and each unit inverse polynomially in $n$. Repeated multiplication of $O(nB)$ such scalar inputs, in a polynomial-length truncated Taylor array, adds only a polynomial to the logarithmic norm. The factor count $O(nB)$ and $B=\operatorname{polylog}n$ are themselves polynomial in $n$. Hence all local scalar preprocessing has polynomial bit cost. The explicit Hasse cardinal recurrence has polynomial bit cost in $Q,J$, exactly as before.

All identities in the integer-fiber separation, Boolean-difference formula, Taylor readout and cyclotomic interpolation hold formally in the weight variables. Store symbolic ring operations as shared gate references, and descend by taking the scalar coordinate in the finite free algebra. This constructs an integral circuit for $Z_t$, with size and bit construction time a polynomial factor times $2^{O_{k,\beta}(n/s)}$. There is no symbolic identity test and no reduction modulo a prime. For constant-free compilation one may first put independent coefficient symbols on every monomial of every factor, use the same factor-block construction, and then specialize each clause coefficient to its fixed integer value and each weight coefficient to its specified parameter. Its formal degree is $O(nB+n)$, since the factor count is $O(nB+n)$ and every such coefficient specialization is constant or linear. This is polynomial in $n$. The scalar-height estimates justify replacing all integer scalar multiplications by binary doubling and addition/subtraction at polynomial overhead. This is the same explicit integral compilation as NN-C034, now with a polylogarithmic factor-count multiplier.

Sum the resulting integral circuits over the disjoint sparsification leaves. Every satisfying assignment contributes its weight exactly once, as an integer polynomial, and no other assignment contributes. The output is therefore $A_F$ over $\mathbb Z$, or the stipulated general weighted sum. The tree count, leaf work and high-incidence branches all have exponent $O_{k,\beta}(n/s)$. Polynomial preprocessing and input normalization give the asserted full time bound. An additive sum of leaf circuits preserves their polynomial formal-degree bound. Formal equality over $\mathbb Z$ proves all coefficient-ring base changes simultaneously.

For completeness, rational specialization of the two-rail weights also has the asserted bit-time bound, without an unproved bound on intermediate arithmetic values. Let the circuit have $S$ gates and formal degree $D=\operatorname{poly}(n)$, and let $L$ be the total bit length of the rational inputs. Choose a common input denominator $Q_0$ by multiplying their denominators, so $\log Q_0\le L$. The constant-free coefficient-norm induction of NN-C015 gives coefficient norm at most $2^{SD}$ at every gate. After evaluation, a denominator divides $Q_0^D$ and the numerator has bit length $O(D(S+L+1))$. Exact rational arithmetic therefore evaluates all $S$ gates in $\operatorname{poly}(S,n,L)$ bit operations. Taking a fixed polynomial power of $S$ preserves the displayed $2^{O_k(n/s)}$ bound. The same argument applies to the general weight coefficients. We make no polynomial-working-space claim from this circuit evaluation argument.

For the independent-set example use the width-two CNF $\bigwedge_{uv\in E}(\neg x_u\vee\neg x_v)$ on the original vertex bits and then substitute $z_{i,0}=1$, $z_{i,1}=y_i$. Setting $k=2$ gives the stated exponent. Its coefficients are the integer zero-one indicators of independent subsets. This proves an integral circuit for their chosen lift by direct integer clause indicators; it is not a generic procedure for lifting circuits from $\mathbb F_2$ to $\mathbb Z$. The sharper characteristic-two bound of NN-C035 is therefore retained alongside this result.

\paragraph{Ordinary integral ideal certificates without a density bound.}
Let $F=C_1\wedge\cdots\wedge C_m$ be an unsatisfiable fixed-width
$k$-CNF on $n$ declared variables.  Use the standard clause encoding:
delete repeated literals, discard tautological clauses, and let
$v_j(x)$ be the product of the distinct falsifying literal factors
of $C_j$.  Empty clauses have $v_j=1$ and give an immediate
certificate.  Then one can deterministically construct integral
constant-free circuits for polynomials $h_j$ such that
\[
                 \sum_{j=1}^m v_j(x)h_j(x)=1
                              \quad\hbox{in }\mathbb Z[x_1,\ldots,x_n],
\]
without adding Boolean equations or verifier variables.  Their
total size and bit construction time are
\[
 (n+m)^{O_k(1)}
  2^{O_k(n(\log\log n)^{1+1/k}/(\log n)^{1/(2k)})}.
\]
If $S_j$ is the scope of $C_j$, the multipliers may be chosen
with $\deg h_j\le n-|S_j|$, so each product $v_jh_j$ has actual
degree at most $n$.  The circuit formal degree is polynomial
in $n$; it is not identified with the sharper actual degree.

The ideal identity is the classical first-violated-clause
certificate of Grochow--Pitassi, \emph{Circuit Complexity, Proof
Complexity, and Polynomial Identity Testing: The Ideal Proof
System}, author JACM manuscript, Lemma~4.2 and equations
(4)--(6), printed draft pp.~1:26--1:28; the primary file is
\url{https://www.cs.toronto.edu/~toni/Papers/jacm-gp.pdf}.
Its Definition~1.1 is on printed draft p.~1:4.  These are
the author manuscript's page markers, not an asserted match
with the final journal pagination.  The new consequence is
the constructive original-$n$ circuit bound for all clause
densities, obtained by the following exact factorization of
those classical multipliers.

Suppose there is no empty clause.  For each $j$, let
$b^{(j)}\in\{0,1\}^{S_j}$ be the unique assignment falsifying
$C_j$.  Restrict the prefix $C_1\wedge\cdots\wedge C_{j-1}$
by these fixed bits, keeping all the remaining $n_j=n-|S_j|$
declared variables.  Call the resulting width-at-most-$k$
CNF $R_j$.  Simplified false or empty prefixes are handled
by the same direct cases as in the main theorem.  Apply
NN-C036 to its full two-rail polynomial
\[
 H_j(z)=\sum_{a\models R_j}\prod_{i\notin S_j}z_{i,a_i},
\]
and specialize $z_{i,0}=1-x_i$, $z_{i,1}=x_i$.  The output is
\[
 h_j(x_{[n]\setminus S_j})=
 \sum_{a\models R_j}\prod_{i\notin S_j}
                       x_i^{a_i}(1-x_i)^{1-a_i}.
\]
Thus $h_j$ is exactly the coefficient of the $j$th placeholder
in the cited first-violation certificate: fixing the falsifying
tuple leaves precisely the assignments satisfying all earlier
clauses.  Multiplication by $v_j$ restores the missing factors
on $S_j$.  In the notation
$\delta_b(x)=\prod_i x_i^{b_i}(1-x_i)^{1-b_i}$, this identifies
$v_jh_j$ with the sum of $\delta_b$ over the first-violation
class of clause $j$.  Equations (4)--(6) of the cited classical
proof therefore give the displayed ordinary integral identity.
No identity merely modulo $x_i^2-x_i$ is substituted for it.

For $n>2k$, each nontrivial call uses $n-k\le n_j\le n$,
so $n_j=\Theta(n)$ and its exponent is bounded by the same
displayed expression in $n$.  There are at most $m$ calls,
and forming all restricted prefixes has polynomial cost
in $n+m$.  The fixed small range of $n$ is handled by direct
expansion.  Each coefficient substitution is a constant-free
linear expression, hence preserves the polynomial formal-degree
bound.  This proves the total circuit and construction bounds.
The displayed expression for $h_j$ proves its actual degree
bound without expanding it.  A uniform algorithm may first
use E Theorem~1.1 to decide satisfiability, reporting that
case or constructing the certificate when the answer is
negative; E's decision cost is absorbed by this weaker exponent.

With placeholders $u_1,\ldots,u_m$, the circuit
\[
                    C(x,u)=\sum_j u_jh_j(x)
\]
is a Hilbert-like IPS certificate: $C(x,0)=0$ and
$C(x,v(x))=1$.  It has actual degree at most $n+1$,
polynomial formal degree, and the same circuit-size bound.
The certificate works over every commutative ring by base
change from its integral identity.  These superpolynomial
size bounds do not make the system polynomially bounded,
give polynomial-time deterministic PIT verification, or
imply a corresponding small monomial-PC derivation.
\end{proof}
\paragraph{Status and complexity impact.}
Integral dense-CNF construction, optimized parameter bound and ordinary ideal-certificate consequence proved. Complete deduction-proof review by root: passed; imported input theorems are not claimed independently proved. This closes the specific integral formal-weight gate for arbitrary-density fixed-width CNFs and arbitrary-graph independence polynomials. The cost is a weaker exponent than the characteristic-two compressed construction: it balances disjoint sparsification with the actual larger interpolation degree and does not lift a finite-field circuit. It also constructs ordinary integral Hilbert-like IPS certificates without Boolean axioms and explicit affine coordinate maps at the original-variable bound, for every clause density. The construction is uniform and preserves the original Boolean variables, but remains superpolynomial. It supplies no polynomial-size IPS or arithmetic circuits, no general deterministic polynomial-time PIT, and no VP/VNP or P/NP collapse. Generic verifier reductions still charge their gate variables before this theorem is applied.

\paragraph{Geometric and representation dictionary.}
Let $F$ be a fixed-width CNF on $n$ declared Boolean coordinates. Its integral Boolean scheme is
\[
 Y_F=\operatorname{Spec}\!\left(
 \mathbb Z[x_1,\ldots,x_n]/(x_i^2-x_i,\ 1-c_C(x):C\in F)\right).
\]
Here $c_C=1-\prod_{\ell\in C}(1-\ell(x))$ is the integral clause-satisfaction indicator. The scheme is the disjoint union of one $\operatorname{Spec}\mathbb Z$ for each satisfying assignment. Put $B=\mathbb Z[z_{i,b}]$, $\mathscr A_F=\mathcal O(Y_F)\otimes_{\mathbb Z}B$, and
\[
                         w=\prod_i((1-x_i)z_{i,0}+x_i z_{i,1}).
\]
Then the finite locally free multiplication trace $\operatorname{Tr}_{\mathscr A_F/B}(w)$ equals the integral assignment polynomial $A_F$. For fixed width $k$, this regular trace function has a deterministically constructible integral constant-free circuit of size $(n+m)^{O_k(1)}2^{O_k(n(\log\log n)^{1+1/k}/(\log n)^{1/(2k)})}$ and polynomial formal degree. The same circuit computes the trace after every base change, including all characteristics and nonreduced coefficient rings. Explicit weight factors with polylogarithmic scopes and linear total input degree have the same bound with their description length and fixed parameters included.

Disjoint sparsification is a decomposition of $Y_F$ into clopen finite etale subschemes over $\mathbb Z$, and the trace adds over it. The construction retains the actual integer component multiplicities, rather than only their residues. Equivalently, the torus with one character $z_{i,b}$ per literal rail has an effective module with one basis vector of character $\prod_i z_{i,x_i}$ for each satisfying assignment; its character is computed by the integral circuit using virtual differences. For a graph, this yields the effective independent-subset character after setting the zero rails to one. The exact trace/base-change source is EGA IV4,18.2.1--18.2.2, printed p111, equation18.2.1.2; the split-torus character and Grothendieck-ring locators are those in NN-S019. This is an effective statement for these presented finite schemes and characters, not a theorem about general GCT multiplicities.

For an unsatisfiable normalized width-$k$ CNF with clause
polynomials $v_1,\ldots,v_m$, the closed subscheme
$V(v_1,\ldots,v_m)\subseteq\mathbb A^n_{\mathbb Z}$ is empty,
and its unit ideal has the displayed explicit integral
circuit certificate of size
$(n+m)^{O_k(1)}2^{O_k(n(\log\log n)^{1+1/k}/(\log n)^{1/(2k)})}$.
No Boolean equations occur in that ideal.  The multipliers
have degree at most $n-|S_j|$ and integral constant-free
circuits of polynomial formal degree.

For the affine hypersurface
\[
 X_F=\{t^2-t=\sum_{j=1}^m y_jv_j(x)\}
                          \subseteq\mathbb A^{n+m+1}_{\mathbb Z},
\]
the same certificate gives the explicit isomorphism
$X_F\simeq\mathbb A^{n+m}_{\mathbb Z}$ by the unit-ideal
coordinate formulas already proved in NN-C029:
\[
 z_j=y_j+h_j(x)(2t-t^2),\qquad
 t=\sum_jv_j(x)z_j,\qquad
 y_j=z_j+h_j(x)(t^2-2t).
\]
The first formula is the map from $X_F$ to coordinates
$(x,z)$, and the last two define its inverse.  The equality
$\sum_jv_jh_j=1$ makes $\sum_jv_jz_j=t$ on $X_F$ and
$\sum_jv_jy_j=t^2-t$ on the inverse image; substitution
then gives both identity composites.  These calculations
use no division and hold over every base ring.  The forward
map has degree at most $n+2$ and the inverse has degree at
most $n+2k+2$, since $\deg(\sum_jv_jz_j)\le k+1$.
Sharing the common sum and multiplier circuits gives both
maps the same mildly subexponential original-$n$ circuit
bound and polynomial formal degree.  This is an effective
integral trivialization for the stated empty clause systems,
not a general recognition theorem for affine space.



\section{NN-C037 --- Rank and divisor bounds exclude dense factors from small-state linear interpolation}
\begin{proposition}[NN-C037]
Let $K$ be any field, let $\alpha_1,\ldots,\alpha_Q\in K^\times$ be distinct, and let $I\subseteq[n]$ have size $w\ge1$. Suppose polynomial replacement rows
\[
          P_v(T)\in K[T]^{1\times n}\quad(v\in I),
          \qquad P_v(0)=e_v,\quad\deg_TP_v\le R
\]
obey E's kernel-separation interface. At node $\alpha_t$, the coordinates are partitioned into groups $S_{t,1},\ldots,S_{t,r}$, the factor is assigned to $g_t$, and integer or $K$-valued measurement matrices $C_{t,h}$ satisfy
\[
 C_{t,h}\xi=0\ \Longrightarrow\quad
               P_v(\alpha_t)\widetilde\xi=0
           \qquad(h\ne g_t,\ v\in I).
\]
Let
\[
 a_t=|I\cap S_{t,g_t}|,\qquad
 c_t=\sum_{h\ne g_t}
              \operatorname{rank}_K C_{t,h}[:,I\cap S_{t,h}].
\]
Then the following necessary bound holds:
\begin{equation}\label{eq:C037-scope-bound}
                   wR\ \ge\ \sum_{t=1}^Q
                                  \max(0,w-a_t-c_t).
\end{equation}
In particular, a full-scope factor $I=[n]$ in a layout with $|S_{t,h}|\le2n/r$ and at most $2\tau n$ measurement rows per node requires
\begin{equation}\label{eq:C037-dense-bound}
                       \frac RQ\ge1-\frac2r-2\tau.
\end{equation}
Whenever $R<Q$, the same bound gives the quantitative necessary scope condition
\[
                 \frac wn\le
                  \frac{2/r+2\tau}{1-R/Q}.
\]
In particular, if $w\ge\varepsilon n$ for fixed $\varepsilon>0$, then $R/Q\ge1-2/(\varepsilon r)-2\tau/\varepsilon$. Thus E's useful regime $r\to\infty$, $\tau\to0$, $R/Q\to0$ cannot accommodate even one factor whose scope contains a fixed positive fraction of the original variables. This conclusion is independent of that factor's polynomial degree.

For a universal dense linear factor, the same obstruction holds even if its replacement coefficients depend rationally on the formal factor coefficients, provided the interpolation nodes, group assignments and measurement matrices remain independent of those coefficients and the same all-vector kernel condition holds. It follows by differentiating the aggregate linear form, rather than assuming coefficient-independent replacement rows.

More generally, let a linear deformation of the universal scope-$w$ factor have coefficients rational in its independent weights and polynomial of degree at most $R$ in $T$. At distinct nonzero rational coefficient-dependent nodes $\alpha_t$, let $N_t$ be the number of distinct formal factor values on the Boolean cube. Then
\[
 wR\ge\sum_t\max(0,w-\lfloor\log_2N_t\rfloor-\delta_t),
 \qquad \delta_t\le1,
\]
with $\delta_t=0$ for coefficient-independent nodes. Any summary with at most $M_t$ states determining that value on a Zariski-dense set of coefficient specializations implies $N_t\le M_t$, even if the summary is coefficient-adaptive or nonlinear. In particular E's actual integer-fiber state count gives, for a full-scope factor,
\[
          R/Q\ge1-2/r-2\tau\log_2(2D+1)-1/n.
\]
Thus even generic coefficient-adaptive Boolean summaries and moving rational nodes cannot restore the useful $R/Q=o(1)$ regime while their state count is $2^{o(n)}$ and the stipulated deformation remains linear in the Boolean coordinates. Over an infinite coefficient field the same obstruction applies to finite arithmetic decision trees producing such rational families with uniform bounds: at least one leaf is Zariski dense and satisfies the generic inequality.

Consequently the integral sparse-permanent theorem NN-C034 does not extend to the dense symbolic permanent just because Ryser's row factors have degree one. This is a proved obstruction to this layout architecture, not an arithmetic-circuit lower bound for the permanent.
\end{proposition}
\paragraph{Exact inputs.}
E: The Exponential Time Hypothesis is False, supplied LaTeX. Definition3.1 (layout), items1--4 and the all-vector kernel condition; Section3.1 independent-incidence degree; Proposition3.4 exact C\_q enumeration; Lemmas4.7--4.8.

NN-C034. Integral universal polynomial construction and growing-scope extension; total-quasilinear-support permanent application, which is retained within its stated hypothesis.

Servedio--Wan, Computing sparse permanents faster, author manuscript18 June2005. Section2 printed p3 first display: Rysers inclusion-exclusion identity.

EGA IV, fourth part, Publications Mathematiques IHES32 (1967). 21.1.4 printed p256: divisor of an invertible-sheaf meromorphic section; 21.1.7 and equation21.1.7.1 p257: effectivity for a regular section; Proposition21.2.11 and full proof p262: section/divisor correspondence.

EGA IV, fourth part, Publications Mathematiques IHES32 (1967). Printed p.111,18.2.1--18.2.2; base-change equation18.2.1.2; primary PDF page110 (zero-based109).

\begin{proof}
We first prove the determinant estimate, including its exact local multiplicity. Form the $w\times w$ matrix
\[
                  A(T)=(P_v(T)|_I)_{v\in I}.
\]
Its value at zero is $I_w$, so $D(T)=\det A(T)$ is nonzero, $D(0)=1$, and $\deg D\le wR$.

Fix one node $\alpha_t$. A row vector on a group which vanishes on $\ker C_{t,h}$ is in the row span of $C_{t,h}$: its linear functional factors through the image of $C_{t,h}$, and a functional on that image extends to the measurement target by choosing a basis. Thus, on each group other than $g_t$, every row $P_v(\alpha_t)$ belongs to the measurement row span. Restricting to columns $I$ shows that the rows of $A(\alpha_t)$ lie in the direct sum of the full coordinate space on $I\cap S_{t,g_t}$ and the row spaces of the restricted $C_{t,h}$ on the other, disjoint groups. Therefore
\[
             \operatorname{rank}A(\alpha_t)
                           \le\min(w,a_t+c_t).
\]

For completeness, let $s=\operatorname{rank}A(\alpha_t)$. Constant invertible row and column operations over $K$ transform $A(\alpha_t)$ into a matrix with its last $w-s$ rows zero. Apply these same operations to $A(T)$. Every entry in those rows now vanishes at $T=\alpha_t$, hence is divisible in $K[T]$ by $T-\alpha_t$. Multilinearity of the determinant gives
\[
                 (T-\alpha_t)^{w-s}\mid D(T).
\]
The determinants of the constant row and column operations are nonzero constants, so they do not alter the order of vanishing. This argument holds in every characteristic. The distinct linear factors for the $Q$ nodes are pairwise coprime. Their product, with the displayed exponents, divides the nonzero $D(T)$. Comparing degrees proves equation~\eqref{eq:C037-scope-bound}.

In an E layout, $a_t\le2n/r$ and $c_t\le\sum_h\operatorname{rank}C_{t,h}\le2\tau n$. For $w=n$, the sum in that equation is at least $Qn(1-2/r-2\tau)$, even when the latter expression is negative; this proves equation~\eqref{eq:C037-dense-bound}. In general it gives
$wR\ge Q(w-2n/r-2\tau n)$. Rearranging when $R<Q$ gives the stated necessary upper bound on $w/n$; dividing by $Qw$ gives its consequence for $w\ge\varepsilon n$. In particular the obstruction remains after deleting $o(n)$ variables from a full-scope factor: its residual scope and its residual variable count are still equal. High-incidence branching on only $o(n)$ bits does not make a dense Ryser row sparse.

\paragraph{Why linearity in the original variables does not remove the obstruction.}
First suppose the replacement rows are independent of the coefficient symbols $u_v$ of a dense linear factor. Its deformed value is
\[
                      L(T,x;u)=\sum_{v=1}^n u_vP_v(T)x.
\]
If $L(\alpha_t,x;u)$ is unchanged under every addition of a vector in the measurement kernel on an unassigned group, as an identity in the independent symbols $u_v$, then comparing their coefficients yields the kernel condition for every $P_v$. The determinant obstruction therefore already applies to this aggregate-linear interpretation.

There is a stronger formulation which does not assume those rows are coefficient-independent. Let $u_1,\ldots,u_n$ be algebraically independent over $K$, put $F=K(u_1,\ldots,u_n)$, and suppose a deformed aggregate linear form is
\[
 L(T,x;u)=b(T,u)+\sum_{v=1}^n\ell_v(T,u)x_v,
 \qquad b,\ell_v\in F[T],\quad\deg_T\ell_v\le R,
 \qquad \ell_v(0,u)=u_v.
\]
The additive term $b$ has no role in invariance under changes of $x$.
Assume the nodes, group partitions, factor assignments and measurement matrices are over $K$, independent of $u$. Suppose that for every $h\ne g_t$ and every $\xi\in\ker_K C_{t,h}$ one has the formal identity
\[
                  \sum_{v\in S_{t,h}}\ell_v(\alpha_t,u)\xi_v=0
                                                     \quad\text{in }F.
\]
Apply the formal $K$-derivation $\partial/\partial u_j$ to each identity and form the Jacobian matrix
\[
                 J(T)=\left(\frac{\partial\ell_v(T,u)}
                                      {\partial u_j}\right)_{j,v}
                                                   \in F[T]^{n\times n}.
\]
One has $J(0)=I_n$ and $\deg_TJ\le R$. The differentiated identities give precisely the same kernel separation for every row of $J(T)$. A $K$-basis of the measurement kernel remains a basis of its kernel after extension to $F$, so this is also the full kernel condition over $F$. The determinant proof applied over $F$ gives equation~\eqref{eq:C037-dense-bound} again. Formal differentiation of rational functions and $\partial u_v/\partial u_j=\delta_{jv}$ are valid over every characteristic; no evaluation at $u=0$, division by a factorial, or regularity of the rational coefficients at zero is required.

Coefficient-independent nodes and measurement matrices are hypotheses of this sharper kernel-specific derivation. Boolean differences in a fiber need not span the whole linear kernel, so that derivation alone would not cover Boolean-only or coefficient-adaptive summaries. The finite-state argument below does cover those cases, including rationally coefficient-dependent nodes, under the stated generic bound on summary values. The integer measurement fibers of NN-C034 are never being identified with kernels obtained by reducing the measurement matrices modulo a prime. The linear-in-Boolean-coordinates hypothesis of this original argument is removed by NN-C040, which also allows finite separable algebraic coefficient extensions. NN-C040 also excludes rational interpolation of low numerator and denominator degree. Uncontrolled interpolation degree, numerical procedures without the stipulated generic coefficient interpretation, and different summation interfaces remain outside these method-specific obstructions.

\paragraph{Exact consequence for the proposed dense-permanent extension.}
For a symbolic $n\times n$ matrix $A=(a_{ij})$, the classical Ryser identity, in Servedio--Wan Section~2, printed p.~3, is
\[
 \operatorname{per}(A)=(-1)^n
       \sum_{x\in\{0,1\}^n}
       \left(\prod_{j=1}^n(1-2x_j)\right)
       \prod_{i=1}^n\left(\sum_{j=1}^n a_{ij}x_j\right).
\]
There are $2n$ factors, and the independent-incidence total degree over $\mathbb Q[a_{ij}]$ is exactly $d=2n$. Indeed the product of the leading homogeneous parts is a nonzero polynomial in the independent incidence and coefficient variables. Each dense row factor nevertheless has scope all $n$ variables. The rank obstruction applies to any one such row, even if the other rows and unary factors are ignored.

E Lemma~4.8 and the integral extension use $r\to\infty$, $\tau\to0$ and $R/Q=O(\log\log n/\sqrt{\log n})$. The last ratio tends to zero, contrary to equation~\eqref{eq:C037-dense-bound}. Thus choosing a larger auxiliary ring, changing from rational to integral Hasse interpolation, or observing that the row factors have input degree one cannot produce that layout. The failure occurs already at the abstract kernel-separation interface, before the construction's bounds on scope, incidence, coefficient height or auxiliary copies are used.

If one retains small groups and few measurement rows but raises $R$ to satisfy the necessary bound, then $R/Q\ge1/2$ for all sufficiently large instances. With the source's actual Ryser degree budget $d=2n$, its truncation parameter is
\[
                    q=\lfloor dR/Q\rfloor\ge n.
\]
Consequently the literal finite-difference enumeration term in E Proposition~3.4 becomes
\[
              C_q=\sum_{j=0}^{\min(q,n)}2^j\binom nj=3^n.
\]
This states what happens to that proved routine and that degree budget; it is not a lower bound on every way to share computations or to prove a smaller degree for a different interpolation polynomial.

Writing one dense linear factor as a sum of unary terms does not turn the product of the factors into a product of unary factors. Directly distributing all row sums produces $n^n$ terms, indexed by one column choice per row, before the cancellations which enforce bijectivity in Ryser's formula. This observation supplies no lower bound for a more efficient aggregate procedure. It pinpoints the missing step: a new way to sum dense linear aggregates which avoids the one-group/all-vector-kernel condition, or a different parameter-preserving reduction, would be needed to derive a dense-permanent circuit from this route.

The obstruction is architectural rather than a general complexity lower bound. For example, a single dense linear factor is trivial to sum:
\[
        \sum_{x\in\{0,1\}^n}\sum_{j=1}^n u_jx_j
                                  =2^{n-1}\sum_{j=1}^n u_j
                  \qquad(n\ge1).
\]
It admits a polynomial-size integral circuit although the small-$R/Q$ separated layout for its full scope is impossible. The sparse and total-quasilinear-support permanent conclusions of NN-C034 remain intact, and this proof neither improves nor refutes an unrestricted dense-permanent upper bound.\paragraph{Arbitrary Boolean summaries and coefficient-dependent nodes.}
The determinant obstruction is not confined to the all-vector
kernel condition.  The following extension uses only a bound on
the number of Boolean summary states and makes the remaining
linearity hypothesis explicit.

Let $K$ be a field, let $I\subseteq[n]$ have size $w\ge1$, and let
$u_v$ for $v\in I$ be algebraically independent over $K$. Put
$F=K(u_v:v\in I)$. Suppose
\[
 L(T,x;u)=b(T,u)+\sum_{j=1}^n\ell_j(T,u)x_j,\qquad
                  b,\ell_j\in F[T],\quad\deg_T\ell_j\le R,
\]
with
\[
 \ell_j(0,u)=
 \begin{cases}u_j&j\in I,\\0&j\notin I.\end{cases}
\]
At each of $Q$ pairwise distinct nonzero nodes
$\alpha_t\in F$, define the number of distinct formal values
\[
            N_t=\left|\{L(\alpha_t,x;u):x\in\{0,1\}^n\}\right|,
\]
where equality is equality of rational functions in $F$.
No summary partition or its coefficient independence is
assumed. Then
\begin{equation}\label{eq:C037-summary-bound}
 wR\ \ge\ \sum_{t=1}^Q
       \max\!\left(0,w-\lfloor\log_2N_t\rfloor-\delta_t\right),
 \qquad
 \delta_t=
 \begin{cases}
 0,&\alpha_t\in K,\\
 1,&\alpha_t\in F\setminus K.
 \end{cases}
\end{equation}
One can also take $\delta_t=0$ whenever all formal
$u_v$-derivatives of $\alpha_t$ vanish. In particular,
for a full-scope universal linear factor, if every summary
has at most $2^{\varepsilon n}$ states, then
\[
                    R/Q\ge1-\varepsilon-1/n.
\]
The $1/n$ term can be omitted for coefficient-independent nodes.
Thus subexponentially many summary states still force
$R/Q\ge1-o(1)$, even for rationally coefficient-dependent nodes
and rationally coefficient-dependent linear deformations.

We prove the extension in full. First, if a matrix
$B\in F^{w\times n}$ has rank $k$, then
\begin{equation}\label{eq:C037-cube-image}
                    |B(\{0,1\}^n)|\ge2^k.
\end{equation}
Every nonempty affine fiber of $B$ is a translate of its
$(n-k)$-dimensional kernel. Choose $n-k$ coordinate functions
whose restriction to that kernel is an isomorphism onto
$F^{n-k}$; such a choice follows by choosing free coordinates
in row reduction. The same coordinate projection is injective
on any affine fiber. Its Boolean points therefore inject into
$\{0,1\}^{n-k}$ and number at most $2^{n-k}$. The $2^n$
Boolean points are partitioned by the image values, proving
equation~\eqref{eq:C037-cube-image}. This elementary argument
works over every field, including fields of characteristic two.
Consequently a matrix constant on each part of a partition
of the cube into $N_t$ parts has rank at most
$\lfloor\log_2N_t\rfloor$.

Form the $w\times n$ Jacobian, taking derivatives with $T$
held fixed:
\[
       J(T)=\left(\frac{\partial\ell_j(T,u)}{\partial u_v}
                         \right)_{v\in I,\ 1\le j\le n}.
\]
It has $T$-degree at most $R$, and its $I$-column minor
$A(T)=J(T)|_I$ satisfies $A(0)=I_w$. Thus
$D(T)=\det A(T)$ is nonzero and has degree at most $wR$.
Partition the cube by equality of its formal values
$L(\alpha_t,x;u)$. For $x,x'$ in the same part, subtracting
the two values cancels $b$. Differentiate the resulting
identity after substituting $\alpha_t(u)$. The formal chain
rule gives
\[
 \left(J(\alpha_t)+v_tq_t\right)(x-x')=0,\qquad
 v_t=\left(\frac{\partial\alpha_t}{\partial u_v}\right)_{v\in I},
 \qquad
 q_t=\left(\frac{\partial\ell_j}{\partial T}
                                (\alpha_t,u)\right)_{1\le j\le n}.
\]
Here $v_t$ is a column and $q_t$ a row, so $v_tq_t$ has
rank at most one, and it is zero when $\alpha_t\in K$.
The matrix $B_t=J(\alpha_t)+v_tq_t$ is constant on every
part of this formal-value partition. The preceding image bound gives
\[
 \operatorname{rank}J(\alpha_t)
 \le\operatorname{rank}B_t+\operatorname{rank}(v_tq_t)
 \le\lfloor\log_2N_t\rfloor+\delta_t.
\]
The same upper bound holds for its column minor
$A(\alpha_t)$. The determinant multiplicity argument already
proved in NN-C037 now implies that $D(T)$ is divisible by
$(T-\alpha_t)^{\max(0,w-\lfloor\log_2N_t\rfloor-\delta_t)}$
for every $t$. These nodes are distinct as elements of $F$,
so degree comparison in $F[T]$ proves
equation~\eqref{eq:C037-summary-bound}. This proof uses
ordinary derivations of the rational function field, never
evaluation at $u=0$ or division by a characteristic-dependent
factorial.

Here is the exact relation to summaries that may themselves
depend on the coefficients. At a fixed $t$, there are only
finitely many formal rational value functions. Remove the
zero sets of their denominators and, for every pair of
distinct such functions, the zero set of the numerator
of their difference. The resulting open subset $U_t$ of
the affine coefficient space is nonempty: the product of
these finitely many nonzero polynomials is nonzero in the
polynomial domain. On this open set, and after any field
extension admitting a point of it, the value partition
is exactly the formal one and has $N_t$ parts. Intersecting
the finitely many $U_t$, and also excluding node collisions
and undefined nodes, still gives a nonempty open subset.
This is an existence argument, not a proposed efficient
construction of this exponentially described open set.

Suppose a possibly coefficient-adaptive summary has at most
$M_t$ states and determines the deformed factor value
pointwise on a Zariski-dense set of coefficient
specializations where these rational expressions are
defined. That dense set meets $U_t$, so $N_t\le M_t$.
In particular a uniform bound on a nonempty generic open
of geometric coefficient space implies the formal bound.
For an infinite base field, it also suffices to have the
bound on all rational points of a nonempty open: these
points are Zariski dense, as follows by successively
choosing coordinates to avoid the finitely many zeros
of nonzero polynomials. The relevant characteristic-zero
coefficient fields satisfy this condition. A bound confined
to the points of one fixed finite field does not imply a
generic bound; that finite set is not Zariski dense.
No regularity or rational formula for the summary itself
is needed here. The rational deformation of the factor
and the generic bound on the number of its possible
summary states are the actual hypotheses.

For the actual E summation interface, a summary for factor
$f$ at node $t$ consists of its assigned group assignment and
the integer measurement vectors. If each group has at most
$2n/r$ bits, there are at most $2\tau n$ measurement rows,
and each measurement takes a value in $[-D,D]\cap\mathbb Z$,
then the number of possible summaries is bounded by
\[
 M_t\le2^{2n/r}(2D+1)^{2\tau n}.
\]
For coefficient-independent summaries the factorization
is a formal identity, so $N_t\le M_t$ directly. For
coefficient-adaptive summaries, the same implication
holds whenever their state bound and factorization hold
on the generic coefficient set just specified.
This is a bound on finite integer fibers, not on kernels
of matrices reduced in a coefficient field. Applying
equation~\eqref{eq:C037-summary-bound} to a full-scope
factor therefore gives
\begin{equation}\label{eq:C037-integer-fiber-bound}
 \frac RQ\ge
       1-\frac2r-2\tau\log_2(2D+1)-\frac1n.
\end{equation}
Again the last term is absent for fixed nodes. The favorable
cost conditions in E and NN-C034 make
$1/r+\tau\log_2(2D+1)=o(1)$. Thus replacing the all-vector
kernel condition by exact constancy on these Boolean
summary fibers does not restore the useful $R/Q=o(1)$ regime.
The stronger kernel-specific estimate of NN-C037 remains
valid when that additional hypothesis holds.

This extension closes both the Boolean-only-fiber escape
and generic coefficient-adaptivity of summaries for these
rational linear deformations. For example, the one-row
formal measurement $\sum_j u_jx_j$ can encode a dense
aggregate, but it takes $2^n$ distinct values in $F$ on
the Boolean cube; its small matrix rank does not give a
small measurement alphabet. Finite arithmetic branching is
covered by the further dense-leaf corollary; numerical
deformations outside its rational-family hypotheses remain
outside the result. Nor are nonlinear deformations in
the Boolean coordinates, nonpolynomial dependence on $T$,
or a different method of summing multiple group contributions
inside one factor. No lower bound for unrestricted permanent
circuits follows.



\paragraph{Finite arithmetic branching on the coefficients.}
The generic obstruction also applies when a finite arithmetic
decision tree supplies the deformation, rather than one global
rational formula. Fix the Boolean size and scope, and let $K$
be an infinite field. Suppose a finite tree is correct on all
$K$-points of a nonempty open subset of affine coefficient
space. Its arithmetic operations are field operations,
including division on branches where denominators are nonzero;
its finitely many discrete tests may select different
deformations and summaries. Assume each leaf produces the
rational, polynomial-in-$T$, linear-in-$x$ deformation
stipulated above, with a common degree bound $R$, a common
number $Q$ of distinct nonzero rational nodes, and at most
$M_t$ summary states determining the value at node $t$.
All these bounds hold pointwise on the leaf's input set.
Then at least one leaf's rational family satisfies
\[
 wR\ge\sum_{t=1}^Q
          \max(0,w-\lfloor\log_2M_t\rfloor-\delta_t),
 \qquad \delta_t\le1,
\]
with $\delta_t=0$ at coefficient-independent nodes of that
leaf. In particular finite branching cannot realize uniformly
forbidden parameters of NN-C037 on that open coefficient set.
Leaf-dependent bounds also give the corresponding inequality
for at least one dense leaf.

Here is the full argument. Let $D_\ell$ be the input set
reaching leaf $\ell$. At least one of these finitely many
sets is Zariski dense. Otherwise, for each leaf choose a
nonzero polynomial vanishing on its proper Zariski closure.
The product of these polynomials is nonzero, since the
coefficient polynomial ring is a domain, but it vanishes
on all points of the given nonempty open set. This is
impossible over an infinite field: its rational points
are Zariski dense, by the elementary polynomial argument
already used above. No algebraicity or constructibility
of the discrete tests themselves is required.

Fix such a dense leaf. Its rational functions are obtained
by composing its field operations. All initial-value
identities hold as identities in the coefficient function
field, because their numerators vanish on this dense
input set. Likewise any coefficient above the claimed
$T$-degree bound vanishes identically if that bound was
specified only pointwise. Its nodes are nonzero and distinct
as rational functions: a zero or coincident pair of such
functions would give the same forbidden equality wherever
the branch was defined. The pointwise state bound holds
on the dense set $D_\ell$, so the formal-value partition
argument already proved gives $N_t\le M_t$. Apply
equation~\eqref{eq:C037-summary-bound} to that leaf to obtain
the asserted inequality.

The finiteness assumption is uniform over the entire
coefficient set at the fixed Boolean size. This does not
automatically cover a bit algorithm whose number of
branches grows without bound with unrestricted coefficient
bit length, operations outside rational arithmetic,
nonlinear deformations in the Boolean coordinates, or
general arithmetic circuits for the permanent. The proof
does not construct the dense leaf efficiently; it is an
obstruction, not a decision-tree simulation algorithm.


\end{proof}
\paragraph{Status and complexity impact.}
Dense-scope determinant bound and coefficient-dependent aggregate extension proved. Complete deduction-proof review by root: passed; imported input theorems are not claimed independently proved. This closes a precise proposed inference from the sparse integral construction to the dense permanent: small input degree does not substitute for locality in E's separated linear interpolation. The obstruction extends to arbitrary Boolean summaries with a generic subexponential state count, rationally coefficient-dependent linear deformations, and rationally moving interpolation nodes. Coefficient-adaptive summaries are covered when their value bound holds on a Zariski-dense coefficient set. Finite arithmetic branching with uniform fixed-size bounds is also covered by the dense-leaf corollary over infinite fields. The remaining route must change the linear deformation or rational-family/finite-tree hypothesis, change how factors involving several groups are summed, exploit a different interpolation invariant, or supply a genuinely parameter-preserving reduction. Enlarging interpolation fields, removing rational denominators, or weakening a field-kernel condition to Boolean fibers alone does not address this obstruction. It is not a lower bound for general arithmetic circuits, a new permanent hardness theorem, or a P/NP or VP/VNP separation.

\paragraph{Geometric and representation dictionary.}
For any field $K$, let $\varphi:\mathcal O_{\mathbb P^1_K}(-R)^w\to\mathcal O_{\mathbb P^1_K}^{w}$ be the vector-bundle map obtained by homogenizing a $w\times w$ polynomial matrix $A(T)$ of degree at most $R$, with $A(0)=I_w$. Its top exterior power is explicitly the determinant section
\[
                   s=\det(\varphi)\in
                     H^0(\mathbb P^1_K,\mathcal O(wR)).
\]
It is nonzero at $[1:0]$. If $\varphi$ has rank at most $k_t$ at the distinct rational points $[1:\alpha_t]$, its effective Cartier zero divisor contains
\[
                        \sum_t(w-k_t)[1:\alpha_t],
                         \qquad \sum_t(w-k_t)\le wR.
\]
This is a standalone geometric form of the obstruction. On the affine chart the local parameter is $T-\alpha_t$; constant row and column reduction makes $w-k_t$ rows divisible by this parameter, proving the claimed local divisor multiplicity. The homogeneous determinant has total degree $wR$; after base change to an algebraic closure, its affine zero multiplicities plus its possible multiplicity at infinity sum to $wR$. Thus the inequality is proved explicitly, without assuming an unverified determinant-functor formula.

For an E factor of scope $I$, restriction to its $w$ coordinate columns gives this map, with $k_t\le\min(w,a_t+c_t)$ because its image at the node factors through the assigned coordinate group and the other groups' measurement quotients. Equation~\eqref{eq:C037-scope-bound} is exactly the resulting divisor-degree inequality. For a rationally coefficient-dependent universal linear factor, apply the same theorem over the generic coefficient field $K(u)$ to the Jacobian matrix proved above. It expresses a genuine obstruction to these families of low-degree vector-bundle degenerations, not a statement that the permanent itself is a determinant.

The exact Grothendieck source interface is EGA~IV4, 21.1.4, printed p.~256, for the divisor of an invertible-sheaf meromorphic section; 21.1.7 and equation~21.1.7.1, p.~257, for effectivity of a regular section; and Proposition~21.2.11 with its proof, p.~262, for the line-section/divisor correspondence. The present local multiplicity and degree calculation supplies the additional quantitative assertion needed here.

Finally, the dense symbolic permanent is still the finite locally free multiplication trace of the displayed Ryser weight in
$\mathbb Z[a_{ij}][x_1,\ldots,x_n]/(x_j^2-x_j)$, up to the scalar $(-1)^n$. EGA~IV4,18.2.1--18.2.2, p.~111 and equation~18.2.1.2 give its trace and base-change interface. The obstruction says that this compactly presented regular trace function cannot acquire NN-C034's useful separated layout merely from its linear factor degrees. It imposes no obstruction on the existence of that finite etale algebra or its trace, and no lower bound on all algorithms for the regular function.\paragraph{Geometric form of the Boolean-summary obstruction.}
Over the generic coefficient field $F$, homogenize the
$I$-column Jacobian minor to obtain
\[
 \varphi:\mathcal O_{\mathbb P^1_F}(-R)^w
               \longrightarrow\mathcal O_{\mathbb P^1_F}^{w},
 \qquad
       s=\det\varphi\in H^0(\mathbb P^1_F,\mathcal O(wR)).
\]
It is nonzero at zero. The partition of the Boolean cube
by the formal values of each deformed linear form
forces the fiber of this Jacobian map at $\alpha_t$ to
have rank at most $\lfloor\log_2N_t\rfloor+\delta_t$.
For moving rational nodes, the additional rank-one map
$v_tq_t$ is exactly the chain-rule correction. Thus the
effective zero divisor of $s$ contains
\[
 \sum_t
 \max(0,w-\lfloor\log_2N_t\rfloor-\delta_t)[1:\alpha_t],
\]
whose degree is bounded by $wR$. A summary with at most
$M_t$ states on a Zariski-dense set of coefficient
specializations bounds $N_t$ by $M_t$, regardless of
coefficient adaptivity of that summary. This is the precise
geometric restatement of equation~\eqref{eq:C037-summary-bound}.
The line-section/divisor conventions and effectivity are
EGA IV4, 21.1.4, 21.1.7 and Proposition21.2.11, printed
pp.~256,257,262; the quantitative multiplicities and the
Boolean fiber count are proved above. The statement concerns
this family of Jacobian degenerations over coefficient space.
It does not equate a finite summary partition with an
algebraic quotient scheme or assert any general obstruction
to existence of the Ryser trace function.


\paragraph{Geometric meaning of finite branching.}
Affine coefficient space is integral. Finitely many input
sets covering a nonempty open cannot all have proper
closed closures, as the polynomial-product proof shows.
On a dense leaf the numerical identities therefore define
rational identities over the generic coefficient field.
Its scope Jacobian minor yields the same map
$\mathcal O_{\mathbb P^1}(-R)^w\to\mathcal O_{\mathbb P^1}^w$,
and the effective zero divisor of its determinant must
contain the rank-defect divisor supplied by the generic
value count. Its degree is at most $wR$, giving the displayed
inequality. The exact divisor/section interface remains
EGA IV4,21.1.4,21.1.7 and Proposition21.2.11, printed
pp.~256,257,262. No quotient scheme is attached to a
discrete test or branch set, and no algebraic structure
of those sets was assumed.




\section{NN-C038 --- Uniform integral circuits for subpolynomial scopes and sparse permanent supports}
\begin{proposition}[NN-C038]
Fix an integer $C\ge1$. There are $n\ge2$ Boolean variables and $m\le Cn$ explicitly listed factors
\[
 \Phi_f(x)=\sum_{j=1}^{\ell_f}y_{f,j}x^{\alpha(f,j)},\qquad
 d_f=\max_j|\alpha(f,j)|,\qquad \sum_f d_f\le Cn,
\]
with independent formal coefficient blocks. Let $1\le w\le n$ bound the declared factor scopes, and let $L$ be the full bit length of their monomial lists, exponents and scopes. Repeated monomials are allowed. Put
\[
 \begin{aligned}
 \lambda&=\lceil\log_2n\rceil,\\
 B&=\left\lceil\log_2\bigl(8(C+1)(w+1)\lambda\bigr)\right\rceil,\\
 z&=\max(2,\lambda/B),\qquad \rho=\lceil\log_2z\rceil,\\
 \theta(n,w)&=\min\left(1,\frac{\rho^{3/2}}{\sqrt z}\right).
 \end{aligned}
\]
The complete formal polynomial
\[
                  Z_\Phi(y)=\sum_{x\in\{0,1\}^n}\prod_f\Phi_f(x)
                                                      \in\mathbb Z[y]
\]
has a deterministically constructible integral division-free circuit of size and bit construction time
\[
                       \operatorname{poly}_C(n,L)\,
                                    2^{O_C(n\theta(n,w))}.
\]
All integer constants have polynomial bit length. Constant-free compilation preserves this bound and has formal degree at most $\max(m,1)$, counting subtraction with the addition degree rule. The identity holds over $\mathbb Z$ and transfers to every commutative coefficient ring.

In particular, for every family with $w=n^{o(1)}$ the additional exponential factor is $2^{o(n)}$, with the original number of Boolean variables retained. For $w\le(\log n)^a$ and fixed $a$, this recovers the bound $2^{O_{C,a}(n(\log\log n)^2/\sqrt{\log n})}$ of NN-C034. For $w\le2^{(\log_2n)^a}$ with fixed $0<a<1$, it gives
\[
 \operatorname{poly}_C(n,L)\,
 2^{O_{C,a}\left(n(\log\log n)^{3/2}/(\log n)^{(1-a)/2}\right)}.
\]
The factor-list degree budget remains essential; arbitrary truth tables on growing scopes are not covered merely by their arity.

For a symbolic $r\times r$ permanent with $S$ allowed entries, put $\lambda_r=\lceil\log_2r\rceil$, $a_r=\max(1,\lceil S/r\rceil)$ and
\[
 \begin{aligned}
 B_{\rm per}&=\left\lceil\log_2\bigl(64(a_r+1)\lambda_r^3\bigr)\right\rceil,\\
 z_{\rm per}&=\max(2,\lambda_r/B_{\rm per}),\\
 \theta_{\rm per}&=
     \min\left(1,
       \frac{\lceil\log_2z_{\rm per}\rceil^{3/2}}{\sqrt{z_{\rm per}}}
         \right).
 \end{aligned}
\]
For $r\ge2$, the formal permanent over $\mathbb Z$ has an integral constant-free circuit of size and construction time
$\operatorname{poly}(r,S)2^{O(r\theta_{\rm per})}$ and formal degree at most $2r$. Thus every support family with $\max(1,S/r)=r^{o(1)}$, equivalently $S\le r^{1+o(1)}$, has a $2^{o(r)}$ construction in the original matrix-order parameter, with no individual row or column bound. This does not include arbitrary dense support.
\end{proposition}
\paragraph{Exact inputs.}
E: The Exponential Time Hypothesis is False, supplied LaTeX. Lemma4.7 general numerical layout interface, including 2(b+w)D\textasciicircum{}2<=M and the exact binomial-tail requirement; Lemma4.8 proof, explicit centered interval and Hoeffding estimate; Section3.1 independent-incidence degree; Proposition3.4 actual q and C\_q.

NN-C034. Integral swap-matrix identity, cyclotomic unit-separated nodes, Hasse interpolation, polynomial coordinate-height proof, coefficient-independent circuit unrolling and constant-free block-degree compilation.

NN-C034 permanent addenda. Ryser factorization with 2r linear factors; exact heavy-row injective expansion and formal-degree accounting.

Servedio--Wan, Computing sparse permanents faster, author manuscript18 June2005. Section2 printed p3 first display: Rysers inclusion-exclusion identity.

NN-C037. Large-scope determinant and generic Boolean-summary obstruction; not contradicted by the smaller scopes treated here.

EGA IV, fourth part, Publications Mathematiques IHES32 (1967). Printed p.111,18.2.1--18.2.2; base-change equation18.2.1.2; primary PDF page110 (zero-based109).

\begin{proof}
The proof reuses the general numerical interface of E Lemma4.7, not its fixed-polylogarithmic parameter corollary outside that corollary's hypotheses. The new point is a uniform choice of parameters depending on the actual logarithmic scope and incidence budget.

If some factor has an empty monomial list, that factor and the partition polynomial are zero. If there are no factors at all ($m=0$), the empty product is one and the output is instead $2^n$, constructible by binary doubling with $O(n)$ constant-free gates of formal degree one. Thereafter all factors may be assumed nonempty. Constant factors can be kept or multiplied outside the counting routine. A bounded initial range of $n$ is expanded directly. Add $n$ dummy degree-zero factors equal to one, with the singleton declared scopes, so that every variable has an incidence. These factors introduce no formal coefficient variables. With $M$ the resulting total incidence count,
\[
                  n\le M\le n+Cn w.
\]
Let
\[
 b_0=\max(1,\lceil M/n\rceil)\lceil\sqrt\lambda\rceil,\qquad
 H=\{v:\deg_{\rm inc}(v)>b_0\}.
\]
Then $|H|\le M/b_0\le n/\sqrt\lambda$. Sum over their Boolean assignments. Keep the remaining declared scopes, even if some listed monomial vanishes, so the choices are independent of all formal coefficients and use no identity oracle.

Each branch has $n'=n-|H|=\Theta(n)$ variables, maximum scope at most $w$, maximum incidence at most $b_0$, and total incidence between $n'$ and $b_0n'$. Its independent-incidence degree remains at most $Cn\le2Cn'$ for sufficiently large $n$. Since
\[
 b_0\le(Cw+2)(\sqrt\lambda+1)
           \le 8(C+1)(w+1)\lambda,
\]
the integer $B$ in the statement satisfies
\[
              b_0,w\le2^B,\qquad B\ge\log_2\lambda'
                 \quad\text{for }\lambda'=\lceil\log_2 n'\rceil.
\]
The extra factors do not change the input-degree budget or the number of genuine coefficient blocks. The branch factor count is at most $(C+1)n$, and all scope data have polynomial length in $n,L$.

We now suppress the prime and construct a layout on $N=n'$ variables, writing $\Lambda=\lceil\log_2N\rceil$. Retain the already known integer $B$. If the numerical checks below fail, expand that branch directly. Otherwise put
\[
 z_0=\Lambda/B,\qquad \rho_0=\lceil\log_2 z_0\rceil,
\]
first requiring $z_0\ge2$. Define $r$ to be the largest nonnegative integer such that
\[
                       r^2B\rho_0^3\le\Lambda.
\]
Require $r\ge2$, put $\ell=\lceil\log_2 r\rceil$, and choose
\[
 \begin{aligned}
 Q&=\left\lfloor\frac{\Lambda}{128\ell}\right\rfloor,\qquad
 R=\max\left(32,\left\lceil\sqrt{64QB}\right\rceil\right),\\
 \tau&=\frac1{r\Lambda},\qquad D=r^{2Q}.
 \end{aligned}
\]
Require $6R\le Q$. All logarithmic ceilings here are computed by comparisons of integers: for example $\rho_0$ is the least integer with $2^{\rho_0}B\ge\Lambda$, and $r,R$ are obtained by integer square-root searches. Thus no real-number oracle is needed.

The exponential dimensions remain polynomial in $N$:
\[
             \log_2D=2Q\log_2r\le\Lambda/64.
\]
In particular the construction in E Lemma4.7, whose cost is polynomial in $N,M,|\mathcal F|,D,Q,r,b,w$, still has polynomial bit cost. Here $b\le b_0$ is the actual remaining incidence bound and $w\le n$; even before imposing the favorable range these data are polynomially bounded in the original input size.

We check every hypothesis of that lemma. For $z_0$ sufficiently large, $B/\Lambda$ is small, so
\[
 2(b+w)D^2\le4\,2^{B+\Lambda/32}
                         \le2^{\Lambda-1}<N\le M.
\]
For example the middle inequality holds once $B\le\Lambda/64$ and $\Lambda$ is sufficiently large. Also
\[
                 2Qr^2\le\Lambda^2/64<N
\]
for all sufficiently large $\Lambda$, since $r^2\le\Lambda$ and $Q\le\Lambda/128$. The construction directly verifies these inequalities, so no unproved small-instance claim is used.

Use the centered interval from the proof of E Lemma4.8:
\[
 \mu=\frac{Q(r-1)}{2r},\qquad
 A_0=\left\lfloor\mu-\frac{R-1}{2}\right\rfloor,\qquad
 \mathcal I=\{A_0,\ldots,A_0+R-1\}.
\]
The checked conditions $R\ge32$ and $6R\le Q$ ensure that this interval is in $\{0,\ldots,Q\}$ and contains all integers within $R/4$ of $\mu$, exactly as in the source proof. Its stated Hoeffding estimate gives
\[
 \Pr\!\left[\operatorname{Bin}\!\left(Q,\frac{r-1}{2r}\right)
                                  \notin\mathcal I\right]
       \le2e^{-R^2/(8Q)}
       \le2e^{-8B}\le2^{-3B}.
\]
On the other hand, $r^2\le\Lambda$, $b\le2^B$ and $\log_2\Lambda\le B$ imply
\[
                  r^2\Lambda b\le2^{3B},
       \qquad 2^{-3B}\le\frac1{r^2\Lambda b}
                               =\frac{\tau}{rb}.
\]
This is exactly the source's required tail bound, with no fixed-polylogarithmic assumption on $b$ or $w$. The exact binomial-tail condition can also be checked by integer arithmetic; its denominator $(2r)^Q$ has $O(\Lambda)$ bits. We therefore have a valid layout whenever the prescribed checks pass.

For completeness, all these checks pass once $z_0=\Lambda/B$ exceeds an absolute constant. As $z_0\to\infty$ one has
\[
 r=\Theta\!\left(\frac{\sqrt{z_0}}{\rho_0^{3/2}}\right),
 \qquad
 Q=\Theta(\Lambda/\ell),\qquad
 \ell=O(\rho_0),\qquad
 R/Q=O\!\left(\sqrt{\rho_0/z_0}\right)=o(1).
\]
Indeed $\sqrt{z_0}/\rho_0^{3/2}\to\infty$, so taking the integer part does not change its order. Since $B\ge\log_2\Lambda$ and $\Lambda\ge z_0$, large $z_0$ also makes $\Lambda$ large. The displayed estimates ensure $r\ge2$, $Q>0$ and $6R\le Q$; the preceding balance and tail inequalities then hold uniformly. Thus if a check fails, $z_0$ lies in a fixed bounded range, not in an uncontrolled exceptional subsequence with $z_0\to\infty$.

We next bound the actual three counting costs. Put $\eta=R/Q$. In the favorable range,
\[
 \eta=O(\sqrt{\rho_0/z_0}),\qquad
 \log_2(2/\eta)=O(\rho_0),\qquad
 1/r=O(\rho_0^{3/2}/\sqrt{z_0}).
\]
For the middle estimate, $R\ge8\sqrt{QB}$ and $Q\le\Lambda/(128\ell)$ give a lower bound $\eta\ge c/\sqrt{z_0}$ with an absolute $c>0$; this bounds the logarithm from above. Also
\[
             \tau\log_2(2D+1)=O(1/r).
\]
Consequently
\begin{equation}\label{eq:C038-parameter-cost}
       \frac RQ\log_2\frac{2Q}{R}
           +\tau\log_2(2D+1)+\frac1r
             =O\!\left(\frac{\rho_0^{3/2}}{\sqrt{z_0}}\right).
\end{equation}

The actual independent-incidence degree of the branch is $d\le2CN$. Take
$q=\lfloor dR/Q\rfloor$ and $J=q+1$, so $QJ>dR$.
E Proposition3.4 uses
\[
 2Qr\,C_q(2D+1)^{2\tau N}2^{2N/r},
       \qquad C_q=\sum_{j=0}^{\min(q,N)}2^j\binom Nj.
\]
The same finite-difference estimate used in NN-C034 gives
\[
 \eta^qC_q\le(1+2\eta)^N,\qquad
 \log_2C_q
      \le q\log_2(1/\eta)+N\log_2(1+2\eta)
      =O_C\!\left(N\eta\log_2(2/\eta)\right).
\]
Equation~\eqref{eq:C038-parameter-cost} therefore bounds the whole exponential factor by
$2^{O_C(N\rho_0^{3/2}/\sqrt{z_0})}$.

The fallback is also covered by the stated uniform bound. If checks fail, $z_0$ is bounded by an absolute constant, and hence
$\min(1,\lceil\log_2\max(2,z_0)\rceil^{3/2}/\sqrt{\max(2,z_0)})$
is bounded below by a positive absolute constant. Direct expansion costs $\operatorname{poly}(n,L)2^N$ and fits the asserted big-O exponent. If the checks pass outside the asymptotically favorable range, the same conclusion follows by enlarging the absolute constant on that bounded $z_0$ range. Thus there is one algorithm and one uniform bound, rather than an unspecified algorithm depending on a subpolynomial growth sequence.

Return to the original parameter. For large $n$, $\lambda/2\le\Lambda\le\lambda$ and $N=\Theta(n)$. The clamped parameter $\max(2,\Lambda/B)$ is between half of $z$ and $z$, up to the harmless common lower bound two, and its logarithmic ceiling is no larger than $\rho$. Its exponent is therefore at most a constant multiple of $n\theta(n,w)$. The heavy-variable branching contributes at most $n/\sqrt\lambda$ to the base-two exponent. This is absorbed: if $z=\lambda/B\ge2$, then $\rho^{3/2}/\sqrt z\ge1/\sqrt\lambda$, while in the clamped range $\theta$ is a positive constant. Summing branches proves the stated original-$n$ bound.

We verify the formal and bit-cost assertions rather than treating an algebra operation as free. The new parameters satisfy $Q,R=O(\log n)$, $D=n^{O(1)}$, $J=O_C(n)$ and $r=n^{O(1)}$. NN-C034's integral swap matrix and cyclotomic algebra use these numerical parameters, not a fixed scope bound. Their explicit coefficient-norm and Hasse recurrence proofs therefore still give polynomial-bit integer coordinates. The auxiliary algebra has rank $O(Q^2)$, with explicit integer multiplication and unit inverses. Its construction and all scalar preprocessing remain polynomial in $n$ and the list length.

After fixing the Boolean bits required by the finite-difference routine, evaluate each listed monomial directly in the truncated algebra, using at most $Cn$ factors or binary powering. There is no enumeration of all $2^{|I_f|}$ local inputs. Before its formal coefficient $y_{f,j}$ is attached, this is a scalar calculation whose coordinate norm grows by at most polynomial bit length, because its number of multiplicative factors is at most $Cn$. Sums over the explicit lists add at most their logarithmic length to this bound. Thus all local scalar preprocessing has polynomial bit cost in $n,L$.

The same coefficient-independent representatives, exact integer measurement fibers, integral Hasse operations and coordinate descent as in NN-C034 give a formal polynomial identity over $\mathbb Z[y]$. Each true factor stays linear in its independent coefficient block. Group products, finite differences, interpolation and branch summation retain total block degree at most $m$. Scalar integer multiplication can be compiled by binary doubling and addition/subtraction, using only $0,1,-1$ and without raising formal degree. The output therefore has a constant-free circuit of formal degree at most $\max(m,1)$. No rational denominator, symbolic pivot, polynomial identity test, or change of characteristic has entered this construction.

The growth conclusions follow directly from the uniform parameter. If $w=n^{o(1)}$, then
$B=O_C(\log(w+1)+\log\log n)=o(\lambda)$, so $z\to\infty$ and $\rho^{3/2}/\sqrt z\to0$. For $w\le(\log n)^a$, $B=O_{C,a}(\log\log n)$ and $\rho=O(\log\log n)$, giving the former exponent. For $w\le2^{(\log_2n)^a}$ with fixed $0<a<1$, $B=O_{C,a}((\log n)^a)$ and $\rho=O(\log\log n)$, giving the displayed stronger-scope bound. The factor $\operatorname{poly}(n,L)$ remains explicit: if an input list itself has exponential length, the theorem does not claim to read it in subexponential time.

\paragraph{Almost linear total permanent support.}
Use Ryser's exact integral identity in the cited Servedio--Wan source. An empty row gives the zero polynomial immediately. Otherwise a minor of order $r'$ with maximum row support $w'$ is the partition polynomial of $r'$ Boolean variables, $r'$ unary sign factors and $r'$ linear row factors. Thus $m=d=2r'$ and the just-proved theorem applies, with $C=2$, regardless of column incidences. This step uses the formal identity over $\mathbb Z$, not a characteristic-two determinant calculation.

For total support $S$, let $a_r=\max(1,\lceil S/r\rceil)$ and $T=a_r\lambda_r^2$. A row is heavy if it has more than $T$ allowed entries. The number $h$ of such rows satisfies
\[
                      h\le S/T\le r/\lambda_r^2.
\]
For large $r$, $r'=r-h$ is between $r/2$ and $r$. Expand the permanent along these heavy rows, over allowed injections $\pi$ into the column set:
\[
 \operatorname{per}A
 =\sum_{\pi}
       \left(\prod_{i\ {\rm heavy}}a_{i,\pi(i)}\right)
            \operatorname{per}A_{V\setminus H,\,[r]\setminus\pi(H)}.
\]
Every permutation restricts uniquely to one such injection, so there is no multiplicity or sign. There are at most $r^h$ branches, adding at most $O(r/\log r)$ to the base-two exponent. The remaining rows have support at most $T$.

For a minor, use the scope bound $\min(T,r')$. Its logarithmic budget in the factor theorem is bounded by $B_{\rm per}$:
\[
 8(2+1)(T+1)\lceil\log_2r'\rceil
       \le24(a_r+1)\lambda_r^3
       \le64(a_r+1)\lambda_r^3.
\]
One may run the proved algorithm with this larger declared scope/incidence budget; its proof required only upper bounds. Since $\lceil\log_2r'\rceil=\Theta(\lambda_r)$, the same comparison of clamped $z$ parameters gives minor cost
$\operatorname{poly}(r,S)2^{O(r\theta_{\rm per})}$.
The outer exponent $r/\log r$ is absorbed by $r\theta_{\rm per}$, as before.

A minor circuit has formal degree at most $2r'$. Multiplying by its $h$ selected entries gives degree at most
$2r'+h\le2r$. Summing branches preserves that bound. All scalar constants retain polynomial bit length and constant-free compilation; the exceptional fixed small sizes are expanded directly. Finally
$\max(1,S/r)=r^{o(1)}$ implies $B_{\rm per}=o(\lambda_r)$ and hence $\theta_{\rm per}=o(1)$. This proves the claimed $2^{o(r)}$ integral formal construction. It does not assert this range for $S=\Theta(r^2)$: then $B_{\rm per}$ is of order $\lambda_r$ and the uniform estimate becomes merely exponential.
\end{proof}
\paragraph{Status and complexity impact.}
Uniform subpolynomial-scope integral construction and permanent corollary proved. Complete deduction-proof review by root: passed; imported input theorems are not claimed independently proved. This broadens the effective integral symbolic trace and permanent classes from fixed polylogarithmic scopes or total support to all subpolynomial scopes with the same linear degree budget, and to total permanent support $r^{1+o(1)}$. It gives one explicit algorithm and a parameter-dependent exponent, retaining the original Boolean-variable or matrix-order parameter, integral constants, polynomial formal degree and all-ring identities. It does not settle dense permanent, polynomial circuit size, or VP/VNP. It does not remove the need to charge the actual degree budget for general growing-arity truth tables. NN-C037's obstruction remains consistent: the present scopes are far smaller than a positive fraction of the variable set, and the new layouts still satisfy the required rank and state-count restrictions.

\paragraph{Geometric and representation dictionary.}
Fix $C\ge1$, $m\le Cn$ explicit polynomial factors with independent coefficients, scope bound $w$ and degree budget $\sum_f\deg\Phi_f\le Cn$, with full list description length $L$. Put $B_0=\mathbb Z[y]$ and
\[
 \mathscr A=B_0[x_1,\ldots,x_n]/(x_i^2-x_i),\qquad
                         u=\prod_f\Phi_f(x).
\]
This finite locally free algebra has rank $2^n$, and its multiplication trace is the regular function
\[
               \operatorname{Tr}_{\mathscr A/B_0}(u)
                 =\sum_{x\in\{0,1\}^n}\prod_f\Phi_f(x)
                                     \quad\text{on }\operatorname{Spec}B_0.
\]
Define
\[
 \begin{aligned}
 \lambda&=\lceil\log_2n\rceil,\\
 B&=\lceil\log_2(8(C+1)(w+1)\lambda)\rceil,\\
 z&=\max(2,\lambda/B),\\
 \theta&=\min(1,\lceil\log_2z\rceil^{3/2}/\sqrt z).
 \end{aligned}
\]
Then
this regular trace function has an integral constant-free circuit of size and deterministic construction time
$\operatorname{poly}_C(n,L)2^{O_C(n\theta)}$ and formal coefficient degree at most $\max(m,1)$. The identity is established before base change and therefore holds over every commutative ring. It is subexponential in the original number of Boolean coordinates whenever the scopes satisfy $w=n^{o(1)}$ and the explicit input length is subexponential. The scheme still has relative dimension zero; no claim of a small explicit idempotent basis is made.

For a bipartite support graph with $r$ vertices in each part and $S\le r^{1+o(1)}$ edges, let a split torus have one specified character for each allowed edge. The actual free torus module
\[
                   \mathcal V_\Omega
                    =\bigoplus_{\pi\ {\rm allowed\ bijection}}
                            \bigotimes_{i=1}^r L_{i,\pi(i)}
\]
has the symbolic permanent as its character. Its class has an integral expression with shared arithmetic subexpressions, allowing virtual differences, of size $2^{o(r)}$ and formal degree at most $2r$, with the full uniform bound given by $B_{\rm per}$ and $\theta_{\rm per}$ in the statement. The heavy-row expansion is an actual direct sum decomposition into selected-edge characters tensored with the corresponding minor modules. Ryser's formula also identifies its character with the multiplication trace of the stated signed weight on the finite Boolean algebra in the original $r$ coordinates.

The exact classical trace and base-change interface is EGA IV4,18.2.1--18.2.2, printed p111, including equation18.2.1.2. Split-torus gradings are SGA3, ExposeI,4.7.2--4.7.3, and direct sum/tensor product in the representation Grothendieck ring are SGA6, ExposeIV,(1.5.2) and Definition2.2, as verified for NN-S019. These statements concern compactly presented regular trace functions and specified effective torus characters; they do not claim arbitrary geometric pushforwards or general GCT multiplicities are polynomial-time computable.



\section{NN-C039 --- Growing-arity integral counting with the actual Boolean degree}
\begin{proposition}[NN-C039]
Fix $C\ge1$. Let $m\le Cn$ factors on $n\ge2$ Boolean variables be explicitly given by monomial lists
\[
               \Phi_f(x)=\sum_j y_{f,j}x^{\alpha(f,j)}
\]
with independent formal coefficient blocks, declared scopes of size at most $1\le w\le n$, and arbitrary nonnegative binary exponents. Let $L$ be the full input bit length. A factor with no listed monomials gives the zero partition polynomial; if $m=0$, the partition polynomial is instead $2^n$. For the nonempty factor lists put
\[
        d=\sum_f\max_j|\{i:\alpha(f,j)_i>0\}|\le mw.
\]
There is no bound on the unreduced exponent sum. Define the explicit parameters
\[
 \begin{aligned}
 \lambda&=\lceil\log_2n\rceil,\\
 B&=\left\lceil\log_2(8(C+1)(w+1)\lambda)\right\rceil,\\
 z&=\max(2,\lambda/B),\qquad \rho=\lceil\log_2z\rceil,\\
 \theta&=\min(1,\rho^{3/2}/\sqrt z),\\
 \vartheta&=\min(1,(1+d/n)\theta).
 \end{aligned}
\]
The complete integral formal weight polynomial
\[
                    Z_\Phi(y)=\sum_{x\in\{0,1\}^n}
                                            \prod_f\Phi_f(x)
\]
has an integral constant-free circuit of size and deterministic bit construction time
\[
                  \operatorname{poly}_C(n,L)\,
                                      2^{O_C(n\vartheta)}.
\]
Its formal coefficient degree is at most $\max(m,1)$, with subtraction assigned the addition degree rule. The output identity holds over $\mathbb Z$ and survives every commutative-ring base change. Very large binary exponents require only their input reading and a zero-versus-positive test.

For $m\le Cn$ full local weight tables on scopes of size at most $w\le\sqrt{\lceil\log_2n\rceil}$, the complete table partition polynomial has size and bit construction time
\[
 \operatorname{poly}_C(n,L_{\rm tab})\,
 2^{O_C\left(n\min\left(1,
           \frac{w(\log\log n)^2}{\sqrt{\log n}}\right)\right)}
\]
and formal degree at most $\max(m,1)$. In particular
$w=o(\sqrt{\log n}/(\log\log n)^2)$ gives a $2^{o(n)}$
integral symbolic construction in the original Boolean-variable parameter.
The full table input is explicitly charged, including the repeated signed weights in its polynomial expansion.

For width-$w$ CNFs with $O(n)$ clauses in the same growing-width range, this constructs the full integral satisfying-assignment polynomial with unary two-rail weights, without verifier variables. If $m$ is the clause count, the circuit has formal degree at most $m+n$. The clause-density restriction is retained; no growing-width sparsification theorem is assumed.
\end{proposition}
\paragraph{Exact inputs.}
NN-C038. Uniform scope-dependent layout, exact checks and fallback, original-variable branching, theta(n,w), all-ring coordinate descent and formal-degree compilation.

E: The Exponential Time Hypothesis is False, supplied LaTeX. Section3.1 independent-incidence degree; complete Proposition3.4 proof with the actual truncation order from degree and interpolation ratio; truncated binomial pair count, including truncation order at least the variable count; polynomial truncation overhead.

NN-C034. Explicit integral Hasse interpolation and coordinate-norm bounds, polynomial in Q,R,J and input description; constant-free block-degree compilation.

NN-C016. Boolean table interpolation formula and its coefficientwise value identity; here applied with growing arity and explicitly charged table expansion.

SGA3, ExposeI. 4.7.2--4.7.3: split-torus representations and character gradings; exact source interface verified for NN-S019.

SGA6, ExposeIV. (1.5.2) and Definition2.2: direct sum, tensor product and the Grothendieck-ring operations; verified for NN-S019.

EGA IV, fourth part, Publications Mathematiques IHES32 (1967). Printed p.111,18.2.1--18.2.2; base-change equation18.2.1.2; primary PDF page110 (zero-based109).

\begin{proof}
The same uniform construction has a useful degree-sensitive form.
Retain the factor-density bound $m\le Cn$, but permit arbitrary
nonnegative binary exponents in the explicitly listed monomials,
without imposing $\sum_f d_f\le Cn$ on those unreduced exponents.
Replace each listed exponent vector by its support indicator
\[
 \bar\alpha(f,j)_i=
       \begin{cases}1&\alpha(f,j)_i>0,\\0&\alpha(f,j)_i=0,\end{cases}
 \qquad
 d=\sum_f\max_j|\bar\alpha(f,j)|.
\]
If one factor has an empty monomial list, the output is zero.
If there are no factors ($m=0$), the output is instead $2^n$,
from the empty product; binary doubling constructs it with
$O(n)$ constant-free gates of formal degree one. Thereafter
all factor monomial lists may be assumed nonempty.
The factors may retain their declared scopes of size at most $w$.
In particular $d\le mw\le Cn^2$. Use exactly the parameter
$\theta(n,w)$ of NN-C038 and put
\[
                  \vartheta(n,w,d)=
                    \min\bigl(1,(1+d/n)\theta(n,w)\bigr).
\]
Then the original formal partition polynomial over $\mathbb Z[y]$
has an integral constant-free circuit of size and deterministic
bit construction time
\begin{equation}\label{eq:C039-actual-degree}
              \operatorname{poly}_C(n,L)\,
                           2^{O_C(n\vartheta(n,w,d))},
\end{equation}
and formal coefficient degree at most $\max(m,1)$.
The dependence on an enormous exponent in the original list
is only through its input bit length and the test whether
it is zero. No arithmetic computation exponential in that
exponent is required.

Here is the proof. At a Boolean value, every positive power
of a coordinate equals that coordinate. Thus
$\Phi_f(x)=\bar\Phi_f(x)$ for all Boolean $x$, as formal
polynomials in the independent weights. Equivalently these
factors represent the same element of the integral Boolean
algebra $\mathbb Z[y][x]/(x_i^2-x_i)$. Reduce each listed
monomial separately and keep its own coefficient symbol,
even when two resulting monomials coincide. This is a
polynomial-time operation on the binary input and uses no
polynomial identity test. It proves the equality of the
two full partition polynomials over $\mathbb Z$ before
any coefficient specialization.

The scope-only branching and layout construction of NN-C038
are unchanged. Its degree estimate on a branch is now
$d_{\rm branch}\le d$, so take the actual
$q=\lfloor d_{\rm branch}R/Q\rfloor$.
With $N$ remaining variables and $\eta=R/Q$, the source
finite-difference estimate gives
\[
 \log_2 C_q
 \le q\log_2(1/\eta)+N\log_2(1+2\eta)
 =O\bigl((N+d)\eta\log_2(2/\eta)\bigr).
\]
The other two exponential costs remain
$O(N/r+N\tau\log_2(2D+1))$. The same parameter proof therefore
bounds the total exponent, including high-incidence
branches, by $O_C((n+d)\theta(n,w))$.
Direct Boolean expansion provides the alternative bound
$\operatorname{poly}_C(n,L)2^{O(n)}$. Choose the corresponding
branch of the algorithm according to whether
$(1+d/n)\theta(n,w)\ge1$. This comparison is computable
exactly by squaring nonnegative quantities and comparing
integers, since $\theta$ is defined by integer logarithmic
ceilings and a rational square root. The original bounded
parameter fallback remains covered because then
$\theta$ is bounded below by a positive absolute constant.
This proves equation~\eqref{eq:C039-actual-degree} uniformly.

The larger degree does not create a hidden bit-cost gate.
After reduction each individual monomial has degree at
most $n$, and the total degree satisfies $d\le Cn^2$.
Hence the truncation length $J=q+1$ is still polynomial
in $n$, and the explicit integral Hasse and coordinate-norm
estimates from NN-C034 remain polynomial in $n,L$.
Each reduced local monomial is evaluated with at most $n$
scalar factors. The unrolling and integral coordinate
descent continue to be formal identities, and every true
factor stays linear in its coefficient block. The formal
degree bound is consequently still $\max(m,1)$; it is
not multiplied by the input-degree budget $d$.

Now suppose the input gives a full table of independent
weights $y_{f,a}$ on every Boolean local assignment, for
$m\le Cn$ scopes of size at most $w$. The factor interpolation
already used in NN-C016 is
\[
 \Phi_f(x)=
       \sum_{a\in\{0,1\}^{I_f}}y_{f,a}
           \prod_{i:a_i=1}x_i
           \prod_{i:a_i=0}(1-x_i).
\]
Its value on a Boolean input is exactly the corresponding
table weight, and its degree is at most $|I_f|$.
Expansion yields at most $3^{|I_f|}$ monomial occurrences.
For the independent-coefficient interface, initially give
each occurrence a separate formal coefficient and then
specialize it to $y_{f,a}$ or $-y_{f,a}$ as appropriate.
All occurrences from the same original factor remain in
one coefficient block. This specialization uses only
constant-free linear expressions, permits repeated weights,
and preserves the formal degree bound. The expanded list
length is polynomial in the original table length:
$3^s\le(2^s)^2$ for every $s\ge0$, and summing these
inequalities gives the same polynomial bound over all factors.
No truth-table oracle is being treated as a unit-cost input.

For $\lambda=\lceil\log_2n\rceil$ and
$1\le w\le\sqrt\lambda$, the scope budget satisfies
$B=O_C(\log\lambda)$ and the established parameter estimate gives
$\theta(n,w)=O_C((\log\log n)^2/\sqrt{\log n})$.
Since $d\le Cnw$, the complete universal table partition
polynomial therefore has integral constant-free size and
bit construction time
\begin{equation}\label{eq:C039-growing-arity}
 \operatorname{poly}_C(n,L_{\rm tab})\,
 2^{O_C\left(
       n\min\!\left(1,
              \frac{w(\log\log n)^2}{\sqrt{\log n}}\right)
          \right)},
\end{equation}
and formal degree at most $\max(m,1)$.
In particular, when
\[
                      w=o\!\left(
                        \frac{\sqrt{\log n}}{(\log\log n)^2}
                              \right),
\]
the full integral symbolic polynomial has a $2^{o(n)}$
construction at the original number of Boolean variables.
The total full-table input length is itself $2^{o(n)}$
in this range when each table entry is one formal input
and there are $O(n)$ factors.

This includes complete weighted counting for linearly many
constraints of that growing arity. For a width-$w$ CNF with
$O(n)$ clauses, specialize its local table to its zero-one
satisfaction indicator and adjoin the $n$ unary two-rail
weight factors. The factor density changes by a constant,
and $d\le O(nw+n)$, so the same bound constructs its full
integral satisfying-assignment polynomial without verifier
variables, with formal degree at most $m+n$ when $m$ is
the original clause count. This is a sparse-CNF growing-width statement;
it does not invoke a sparsification theorem with a growing
width parameter or claim the same range for arbitrary
clause density.


\end{proof}
\paragraph{Status and complexity impact.}
Degree-sensitive construction, binary-exponent reduction and growing-arity table/CNF consequences proved. Complete deduction-proof review by root: passed; imported input theorems are not claimed independently proved. This removes the linear unreduced-input-degree hypothesis from the explicit monomial-list trace algorithm, replacing it with a proved cost depending on the actual Boolean-reduced degree. It also extends complete integral symbolic weighted counting from fixed-arity tables to the stated growing-arity range, with the original Boolean variables and without extra verifier wires. The result remains subexponential rather than polynomial. NN-C042 subsequently handles arbitrary clause density in an explicit slower growing-width range by uniform sparsification and packing. The present theorem does not cover arbitrary clause density throughout its larger growing-width range, a succinct oracle for an exponentially long table, or dense permanent within a subexponential matrix-order bound. The degree in Boolean inputs controls running time, while the formal degree of the output-weight circuit stays bounded by the number of factors; these two degrees are not being identified.

\paragraph{Geometric and representation dictionary.}
For the compactly presented finite locally free Boolean
algebra
\[
 \mathscr A=\mathbb Z[y][x_1,\ldots,x_n]/(x_i^2-x_i),
                \qquad u=\prod_f\Phi_f(x),
\]
the coefficientwise monomial reduction above leaves $u$
unchanged. Its multiplication trace therefore has the
integral constant-free circuit bound
\[
        \operatorname{poly}_C(n,L)\,
           2^{O_C(n\min(1,(1+d/n)\theta(n,w)))}
\]
and formal coefficient degree at most $\max(m,1)$, where
$d$ is the sum of the maximum support sizes of the reduced
factor monomials and $\theta(n,w)$ is the explicitly
defined uniform parameter of NN-C038. This is an identity
of regular trace functions over $\mathbb Z[y]$ and survives
every commutative-ring base change, even when the original
binary exponents are very large. The trace and base-change
interface is EGA IV4,18.2.1--18.2.2, printed p.~111,
including equation18.2.1.2.

For full local tables, let the split torus have one
specified character $y_{f,a}$ for each local assignment.
The actual free module
\[
             \bigoplus_{x\in\{0,1\}^n}
                          \bigotimes_f L_{f,x|I_f}
\]
has that table partition polynomial as its character,
with multiplicities retained if different assignments
have the same character. For $m\le Cn$ and the growing
arity range above, its effective class has an integral
expression with virtual differences and shared arithmetic
subexpressions of size $2^{o(n)}$ and formal degree at
most $\max(m,1)$. For the CNF satisfaction specialization
with its added $n$ unary weights, the character belongs to
the finite etale satisfying-assignment scheme; its circuit
has the same time bound and formal degree at most $m+n$,
with $m$ the clause count. The exact torus and Grothendieck-ring dictionary
is SGA3, ExposeI,4.7.2--4.7.3 and SGA6, ExposeIV,(1.5.2),
Definition2.2. These specified circuit bounds do not assert
polynomial construction of a basis or of general
representation-theoretic multiplicities.




\section{NN-C040 --- A jet determinant excludes arbitrary Boolean deformations with few generic values}
\begin{proposition}[NN-C040]
Let $K$ be any field, let $u_1,\ldots,u_w$ be algebraically independent over $K$, and put $F=K(u_1,\ldots,u_w)$ and $N=2^w$, where $w\ge1$. For each Boolean vector $x\in\{0,1\}^w$, choose an arbitrary polynomial
\[
 L_x(T)\in F[T],\qquad \deg_TL_x\le R,
 \qquad L_x(0)=\sum_{j=1}^w u_jx_j.
\]
There is no linearity, polynomial degree, or circuit-size hypothesis on the dependence of $L_x$ on $x$. Let $\alpha_1,\ldots,\alpha_Q\in F^\times$ be distinct. At node $t$, partition the Boolean cube into the $M_t$ classes on which the formal value $L_x(\alpha_t)\in F$ is equal. Set
\[
 \begin{aligned}
 H_w(k)&=\sum_{j=0}^{k}\binom wj,\\
 A_w(k)&=\sum_{j=0}^{k}(w-j+1)\binom wj
                      \le(w+1)H_w(k)\qquad(0\le k<w).
 \end{aligned}
\]
Use $D_t(k)=H_w(k)$ if every partial derivative of $\alpha_t$ with respect to the $u_j$ is zero, and use $D_t(k)=A_w(k)$ otherwise. Then
\begin{equation}\label{eq:C040-exact}
 Rw2^{w-1}\ \ge\
  \sum_{t=1}^Q\sum_{k=0}^{w-1}
                        \max(0,2^w-M_tD_t(k)).
\end{equation}
More precisely, if the value classes at node $t$ have sizes $m_{t,g}$, the corresponding summand may be replaced by
$\sum_g\max(0,m_{t,g}-D_t(k))$.

Consequently, in any family with $w\to\infty$ and
\[
                   \max_{1\le t\le Q}\log_2 M_t=o(w),
\]
one necessarily has $R/Q\ge1-o(1)$. This holds for moving rational nodes, for completely arbitrary Boolean-coordinate deformations, and over every characteristic. The same statement applies when an arbitrary coefficient-adaptive summary has at most $M_t$ states determining $L_x(\alpha_t)$ on a Zariski-dense set of coefficient specializations over an infinite field. It also applies to finite rational arithmetic decision trees with uniform node, degree and state bounds on their leaves, by restriction to a dense leaf.

Thus a universal dense linear factor cannot be deformed by low-degree polynomials in the interpolation variable into only $2^{o(w)}$ formal values at every node while retaining $R/Q=o(1)$, even after allowing arbitrary nonlinear dependence on the Boolean coordinates. This is a necessary condition for the specified interpolation method, not a permanent circuit lower bound.

All these inequalities remain valid over any finite separable coefficient extension $F'/F$, permitting coefficients and nodes in $F'$ and using the uniquely extended coefficient derivations. In characteristic zero this includes every finite algebraic coefficient extension. A pointwise summary claim in this version must be generic on an integral coefficient model with function field $F'$.

For rational deformations $L_x(T)=A_x(T)/B_x(T)$ regular at zero and the nodes, with reduced numerator and denominator degrees at most $R$, the exact inequality instead has left side $2w2^wR$. Subexponential state counts therefore force $R/Q\ge1/4-o(1)$. With a common denominator, numerator degree $a$ alone satisfies the original polynomial inequality and $a/Q\ge1-o(1)$, regardless of denominator degree.
\end{proposition}
\paragraph{Exact inputs.}
NN-C037. Formal-value partitions, generic-open specialization and finite rational arithmetic-tree corollary; the present proof removes its linear-in-Boolean-coordinates hypothesis.

E: The Exponential Time Hypothesis is False, supplied LaTeX. Definition3.1 and Proposition3.4: polynomial interpolation degree, distinct interpolation nodes and enumeration of integer measurement states; used only to compare the proved obstruction with that architecture.

EGA IV, fourth part, Publications Mathematiques IHES32 (1967). 21.1.4 printed p256; 21.1.7 and equation21.1.7.1 p257; Proposition21.2.11 with full proof p262: effective divisor of a regular invertible-sheaf section and section/divisor correspondence.

Keith Conrad, Separability II, author notes. Theorem1.5 printed p1 and complete proof in Section4 p6; TheoremB.8 and full proof pp11--12, CorollaryB.10 and ExampleB.9 p12.

\begin{proof}
Differentiate only in the independent coefficient parameters. For each Boolean vector define
\[
 P_x(T)=\left(\frac{\partial L_x(T)}{\partial u_1},\ldots,
                    \frac{\partial L_x(T)}{\partial u_w}\right)
                    \in F[T]^w.
\]
The usual derivations on the rational function field give
$P_x(0)=x$ and $\deg_TP_{x,j}\le R$ in every characteristic. No division by a factorial is involved.

Form the $N\times N$ matrix whose rows are indexed by subsets $S\subseteq[w]$ and whose columns are indexed by Boolean vectors:
\[
         V(T)_{S,x}=\prod_{j\in S}P_{x,j}(T).
\]
At $T=0$, its entry is one exactly when $S\subseteq\operatorname{supp}(x)$. Order rows and columns by the same linear extension of inclusion among subsets. The resulting Boolean inclusion matrix is triangular with diagonal entries one. Hence
\[
       \Delta(T)=\det V(T)\ne0,\qquad \Delta(0)=1,
       \qquad \deg_T\Delta\le
                R\sum_{S\subseteq[w]}|S|=Rw2^{w-1}.
\]
This determinant is a proof device; no algorithm is asked to build its exponentially large matrix.

\paragraph{The coefficient vectors available near one value class.}
Fix a node $\alpha=\alpha_t$, put $z=T-\alpha$, and set
$v=\nabla_u\alpha\in F^w$. If the common value of one class is $\lambda_g\in F$, the chain rule gives, for every $x$ in that class,
\[
 \nabla_u\lambda_g=P_x(\alpha)+v\,
                         \partial_TL_x(\alpha).
\]
Thus there are a vector $c_g\in F^w$ common to the class and scalars $a_x\in F$ such that
\[
 P_x(\alpha)=c_g+v a_x,
 \qquad
 P_x(\alpha+z)=c_g+v a_x+z h_x(z),
 \qquad h_x(z)\in F[z]^w.
\]
If $v=0$, these initial vectors are all the same point. Otherwise they lie on one affine line. The same direction $v$ works for every class, although the centers may differ.

Write $\operatorname{ev}(p)=(\prod_{j\in S}p_j)_{S\subseteq[w]}$ for the squarefree monomial evaluation vector. For every subset $B\subseteq[w]$ define the vector
\[
 [v_B(p)]_S=
 \begin{cases}
       \displaystyle\prod_{j\in S\setminus B}p_j,&B\subseteq S,\\
       0,&B\not\subseteq S.
 \end{cases}
\]
Expanding each product gives the identity
\begin{equation}\label{eq:C040-Hasse}
 \operatorname{ev}(p+z h(z))=
  \sum_{B\subseteq[w]}z^{|B|}
                  \left(\prod_{j\in B}h_j(z)\right)v_B(p).
\end{equation}
It is an identity with integer coefficients and therefore holds in every characteristic. In particular, for $p=c_g$, every coefficient vector of order at most $k$ lies in the span of the $H_w(k)$ vectors $v_B(c_g)$ with $|B|\le k$.

For $p=c_g+v a$, the vector $v_B(c_g+va)$ is polynomial in $a$ of degree at most $w-|B|$. Write its coefficient vectors as $v_{B,\ell}(c_g,v)$, for $0\le\ell\le w-|B|$. Applying equation~\eqref{eq:C040-Hasse} shows that all coefficient vectors of orders at most $k$, from every column in this value class, lie in the common vector space
\[
 W_{g,k}=\operatorname{span}_F
  \{v_{B,\ell}(c_g,v):|B|\le k,\ 0\le\ell\le w-|B|\}.
\]
Its dimension is at most $A_w(k)$. Coefficients of the scalar products of the $h_{x,j}(z)$, and the scalars $a_x$, merely form linear combinations within this same space. When $v=0$, the preceding smaller dimension bound $H_w(k)$ applies. This proves the asserted bound $D_t(k)$ without any restriction on the original dependence on $x$.

\paragraph{Collision multiplicities.}
Expand every column of $V(\alpha+z)$ as a polynomial in $z$ and expand the determinant by multilinearity in its columns. A nonzero determinant term chooses one coefficient vector of some order $d_x\ge0$ from each column. Fix $k$. Within a class of size $m_{t,g}$, at most $D_t(k)$ of the selected orders can be at most $k$: more such vectors would lie in the same vector space of dimension at most $D_t(k)$ and would be linearly dependent, forcing that determinant term to vanish. Therefore every nonzero determinant term satisfies
\[
 \begin{aligned}
 |\{x:d_x>k\}|&\ge
       \sum_g\max(0,m_{t,g}-D_t(k))\\
 &\ge \max(0,N-M_tD_t(k)).
 \end{aligned}
\]
All these inequalities hold simultaneously for the same term. Since
$\sum_x d_x=\sum_{k\ge0}|\{x:d_x>k\}|$, its total exponent of $z$ is at least the sum of either displayed lower bound over $0\le k<w$. Cancellation among determinant terms can only increase the first nonzero exponent. It follows that
\[
 \operatorname{ord}_{T=\alpha_t}\Delta\ge
  \sum_{k=0}^{w-1}\sum_g\max(0,m_{t,g}-D_t(k))
 \ge \sum_{k=0}^{w-1}\max(0,N-M_tD_t(k)).
\]
The factors $T-\alpha_t$ are pairwise coprime in $F[T]$. Adding their multiplicities and comparing with the previously proved degree of the nonzero determinant proves equation~\eqref{eq:C040-exact} and its class-size refinement. In particular, for any integer $0\le k<w$ the weaker but useful bound is
\begin{equation}\label{eq:C040-threshold}
 Rw2^{w-1}\ge
       (k+1)\sum_{t=1}^Q\max(0,N-M_tD_t(k)),
\end{equation}
because $D_t(j)\le D_t(k)$ when $j\le k$.

\paragraph{Uniform subexponential value counts.}
Fix a real number $0<\beta<1/2$ and put $k=\lfloor\beta w\rfloor$. The elementary binomial estimate needed here follows directly from the binomial theorem: with $a=\beta/(1-\beta)<1$,
\[
 H_w(k)\le a^{-k}(1+a)^w
       \le a^{-\beta w}(1+a)^w
       =2^{h_2(\beta)w},
\]
where $h_2(\beta)=-\beta\log_2\beta-(1-\beta)\log_2(1-\beta)<1$. The first inequality uses $a^j\ge a^k$ for $j\le k$. Thus, uniformly in the nodes,
\[
 \frac{M_tD_t(k)}{2^w}
   \le (w+1)\,2^{\max_s\log_2M_s-(1-h_2(\beta))w}
                                  \longrightarrow0.
\]
Dividing equation~\eqref{eq:C040-threshold} by $Qw2^{w-1}$ gives
$R/Q\ge 2(k+1)(1-o(1))/w=2\beta-o(1)$.
For every $\varepsilon>0$, choose a fixed $\beta<1/2$ with $2\beta>1-\varepsilon/2$ and then take $w$ large. This proves $R/Q\ge1-\varepsilon$ eventually, which is exactly the claimed $1-o(1)$ lower bound. No hidden bound on $Q$ is needed because the state-count hypothesis is uniform over its nodes.

For coefficient-independent nodes, the bound is occasionally exact: if $M_t=1$ for every node, then
\[
 \sum_{k=0}^{w-1}(2^w-H_w(k))
     =\sum_{j=0}^w j\binom wj=w2^{w-1},
\]
so $R\ge Q$. The asymptotic statement above also covers moving nodes, for which a polynomially larger jet space has the same exponential scale.

\paragraph{Adaptive summaries and the precise algorithmic scope.}
The formal-value argument of NN-C037 uses only a finite list of rational functions, not linearity in the Boolean variables. In the present setting that list consists of the $L_x(\alpha_t)$. For each node, invert all their denominators and all nonzero pairwise differences. On the resulting nonempty coefficient open set, different formal values remain different after specialization. Any summary with at most $M_t$ states determining the specialized value on a Zariski-dense coefficient set therefore bounds the formal number of values by $M_t$. Applying the proved inequality gives the claimed summary obstruction. The same dense-leaf argument from NN-C037 applies to a finite rational arithmetic decision tree: identities holding on the selected dense leaf hold in $F$, and its generic state bounds transfer to this formal list.

These applications require generic coefficient correctness or the formal rational-function hypothesis. A fixed finite set of numerical coefficient vectors is not a Zariski-dense substitute. If a factor has scope $w$ inside a larger Boolean cube, fix all coordinates outside that scope before applying the theorem; the number of values cannot increase, and the original linear factor on that restricted cube is still the stated universal scope-$w$ factor. The theorem permits arbitrary values of the deformed factor on all Boolean vectors, including values produced by a nonlinear Boolean polynomial or circuit. It retains polynomial dependence on $T$ of degree at most $R$, distinct rational interpolation nodes, and exact dependence on universal coefficients. Methods using nonpolynomial interpolation or a different summation interface remain outside its claim. Nothing in the proof lower-bounds unrestricted arithmetic circuits for the permanent or other counting functions.

\paragraph{Separable algebraic coefficient operations.}
All the conclusions of NN-C040 continue to hold after replacing
$F=K(u_1,\ldots,u_w)$ by an arbitrary finite separable field extension
$F'/F$, permitting both the coefficients of the $L_x(T)$ and the
interpolation nodes to lie in $F'$. In the definition of $D_t(k)$,
use the unique extensions to $F'$ of the coefficient derivations
$\partial/\partial u_j$. No bound on the extension degree, nor an
algorithmic presentation of this field, is required for this necessary
condition. In characteristic zero, this includes all coefficients and
nodes obtained by finitely many algebraic operations from $K(u)$.

Indeed, the classical extension theorem for derivations supplies each
operator $\partial_j:F'\to F'$, retaining
$\partial_j(u_i)=\delta_{ij}$. Its exact source is Keith Conrad,
\emph{Separability II}, Theorem~1.5 on printed p.~1 and its complete
proof in Section~4 on p.~6; the explicit simple-extension formula and
proof are Theorem~B.8, pp.~11--12, with Corollary~B.10 on p.~12.
Apply these extended derivations coefficientwise to $L_x(T)$, treating
$T$ as constant. The resulting vectors $P_x(T)\in F'[T]^w$ still have
degree at most $R$ and value $x$ at zero. The ordinary product rule in
$F'$ gives
\[
 \partial_j\bigl(L_x(\alpha_t)\bigr)
   =\bigl(\partial_jL_x\bigr)(\alpha_t)
        +\partial_TL_x(\alpha_t)\partial_j(\alpha_t).
\]
Therefore the formal-value classes again constrain their initial
vectors to points or affine lines, exactly as in the main proof.
All Hasse expansions, column-dimension estimates, determinant
multiplicities and degree comparisons are field-algebra identities,
so they hold over $F'$ unchanged. This proves both the exact inequality
and $R/Q\ge1-o(1)$ for uniformly subexponential value counts.
Only first derivations are used; no additional integrability or
commutation claim is needed.

A pointwise summary bound in this version must hold on a Zariski-dense
set of an integral coefficient model with function field $F'$, on which
the chosen algebraic coefficients and nodes are actual rational
functions. On such a model, deleting their finitely many poles and
nonzero pairwise value differences gives the same nonempty-open
argument as before. No claim about an unspecified numerical root
selection over the base coefficient space is substituted for this
hypothesis. In positive characteristic the separability condition is
retained: Conrad's Example~B.9 on p.~12 explicitly shows that the
$u$-derivation on $\mathbb F_p(u)$ cannot extend across
$\mathbb F_p(u^{1/p})$.



\paragraph{Rational dependence on the interpolation variable.}
The same small-state strategy cannot be restored merely by replacing
polynomial deformation by rational deformation of low numerator and
denominator degree. Retain the coefficient field $F$, or a finite
separable extension $F'$ as above, and write that field as $E_0$.
Suppose each Boolean vector is assigned
\[
 L_x(T)=\frac{A_x(T)}{B_x(T)}\in E_0(T),
 \qquad \deg A_x,\deg B_x\le R,
\]
where $R\ge0$, the numerator and denominator are relatively prime,
and $L_x$ is regular at zero and at all the distinct interpolation
nodes $\alpha_t$. Assume the same initial values
$L_x(0)=\sum_j u_jx_j$. A zero numerator is allowed, with
$B_x=1$. The reduced presentation and regularity imply
$B_x(0)\ne0$ and $B_x(\alpha_t)\ne0$.
Define the formal-value classes, their counts $M_t$, and the functions
$D_t(k)$ exactly as in NN-C040. Then
\begin{equation}\label{eq:C040-rational}
 2w2^w R\ \ge\
    \sum_{t=1}^Q\sum_{k=0}^{w-1}
                           \max(0,2^w-M_tD_t(k)),
\end{equation}
again with the class-size refinement from the main theorem.
In particular uniformly $M_t=2^{o(w)}$ forces
\[
                       R/Q\ge\tfrac14-o(1).
\]
No sharpness of this constant is claimed. Its positive value is
sufficient to exclude $R/Q=o(1)$ in this rational-degree model.

To prove this, use the extended coefficient derivations when needed
and put
\[
 N_{x,j}=(\partial_jA_x)B_x-A_x(\partial_jB_x),
 \qquad P_{x,j}=\frac{N_{x,j}}{B_x^2}.
\]
The numerator degree is at most $2R$, and $P_x(0)=x$.
The squarefree evaluation matrix $V(T)$ now has rational-function
entries. Clear its $x$-column by the factor $B_x(T)^{2w}$;
the resulting polynomial matrix has entries
\[
 \widetilde V(T)_{S,x}
    =B_x(T)^{2(w-|S|)}\prod_{j\in S}N_{x,j}(T).
\]
Every entry has degree at most $2wR$, so
$\deg\det\widetilde V\le2w2^wR$. Its value at zero is
$\prod_x B_x(0)^{2w}\ne0$, because $V(0)$ is still the Boolean
inclusion matrix. The determinant is therefore nonzero.

At any chosen node, all the $B_x(\alpha_t)$ are nonzero.
Consequently column clearing multiplies the determinant by a unit
in the local power-series ring and does not alter its order of
vanishing there. The chain rule for the rational $L_x$ gives the
same point or affine-line description of $P_x(\alpha_t)$ in each
formal-value class. Write
$P_x(\alpha_t+z)=c_g+v a_x+z h_x(z)$ with
$h_x(z)\in E_0[[z]]^w$. The Hasse product identity and the
low-order vector-space estimates from the main proof hold for
these formal series as well: each coefficient of a fixed order
uses only finitely many coefficients of the $h_x$.
In the determinant expansion, every coefficient likewise involves
only finitely many choices of column orders. Thus the same
simultaneous filtration argument gives all the previously proved
local multiplicity lower bounds. Their sum is at most the degree
of $\det\widetilde V$, proving
\eqref{eq:C040-rational}.

For the asymptotic consequence, choose any fixed $\beta<1/2$ and
$k=\lfloor\beta w\rfloor$. The main theorem's binomial estimate
again makes $M_tD_t(k)/2^w=o(1)$ uniformly. The first $k+1$
terms of the local sum therefore give
$2w2^wR\ge(k+1)Q2^w(1-o(1))$, hence
$R/Q\ge\beta/2-o(1)$. Letting the fixed choice of $\beta$
approach $1/2$ proves the claimed $1/4-o(1)$ bound.

There is a sharper statement when the denominators can be chosen
common to all Boolean vectors. Suppose
$L_x(T)=A_x(T)/B(T)$, with $B$ nonzero at zero and the nodes,
and set $a=\max_x\deg A_x$. Then
\[
                 \widehat L_x(T)=A_x(T)/B(0)
\]
is a polynomial deformation of degree at most $a$, has the same
initial linear factor, and has exactly the same formal-value
partition at each node: its values there equal those of $L_x$
multiplied by the common nonzero scalar $B(\alpha_t)/B(0)$.
NN-C040 therefore applies directly with degree $a$, giving its
stronger exact inequality and $a/Q\ge1-o(1)$ for subexponential
state counts. This common-denominator conclusion requires no
bound at all on $\deg B$.

All generic coefficient and finite arithmetic-tree qualifications
remain those of NN-C040. This addendum concerns the actual degrees
of reduced rational functions in $T$, not merely the size of a
circuit describing them; a small circuit can describe a rational
function of very large degree. Rational deformations of degree
comparable to the number of nodes, nonalgebraic families, and
summation algorithms not using this value-state interface remain
outside the excluded low-degree route. No general arithmetic
circuit lower bound follows.


\end{proof}
\paragraph{Status and complexity impact.}
The complete nonlinear Boolean deformation, jet-divisor separable-extension and rational-interpolation arguments are proved. Complete deduction-proof review by root: passed; imported input theorems are not claimed independently proved. This closes the nonlinear-in-Boolean-coordinate escape from NN-C037 within polynomial interpolation families over the rational coefficient field or any finite separable coefficient extension. Arbitrary adaptive summaries with uniformly subexponential generic state counts still force the interpolation degree to be asymptotically at least the number of nodes. In particular that route cannot extend the E small-state method to a universal factor of linear scope while preserving its small interpolation-degree ratio. Low-degree rational interpolation is also excluded, with a proved positive lower bound on its degree-to-node ratio; a common denominator retains the sharper polynomial bound. This is a method-specific obstruction and leaves the general dense permanent algorithm and polynomial circuit gates open. The proof imposes no computational representation on the deformed Boolean functions and does not claim that its exponentially large determinant is efficiently constructible.

\paragraph{Geometric and representation dictionary.}
Let $F=K(u_1,\ldots,u_w)$ and let
\[
             Y=\operatorname{Spec}
             F[x_1,\ldots,x_w]/(x_j^2-x_j:1\le j\le w)
\]
be the finite etale Boolean scheme of rank $N=2^w$. For each $j$, the coefficient derivative $P_j(T)$ of the assigned family $L(T)\in\Gamma(Y,\mathcal O_Y)[T]$ is an arbitrary section of that algebra with $P_j(0)=x_j$ and $T$-degree at most $R$. There is no claim that the $P_j(T)$ remain idempotent or define an algebra endomorphism of $Y$.

On $\mathbb P^1_F$, squarefree multiplication of these sections and evaluation on the $N$ points of $Y$ define a morphism of vector bundles
\[
 \varphi:\bigoplus_{S\subseteq[w]}
                  \mathcal O_{\mathbb P^1}(-R|S|)
                 \longrightarrow\mathcal O_{\mathbb P^1}^{N}.
\]
Its affine-coordinate matrix is $V(T)^{\mathsf T}$, with column $S$ of degree at most $R|S|$; homogenization defines the indicated global map. At $T=0$ it is invertible. Its determinant is therefore a nonzero global section of
$\mathcal O_{\mathbb P^1}(Rw2^{w-1})$ and defines an effective divisor of that degree. At a rational node $\alpha_t$, coefficient differentiation confines the initial evaluation points of each formal-value class to one point if $\nabla\alpha_t=0$, or to one affine line otherwise. The explicit Hasse expansion in equation~\eqref{eq:C040-Hasse} bounds its order-$k$ coefficient-vector space by $H_w(k)$ or $A_w(k)$, respectively. Consequently this determinant divisor contains
\[
 \sum_t\left(\sum_{k=0}^{w-1}
            \max(0,2^w-M_tD_t(k))\right)[\alpha_t],
\]
with the stronger class-size multiplicities proved above. Comparison of effective-divisor degrees proves equation~\eqref{eq:C040-exact}, and subexponential value multiplicities force $R/Q\ge1-o(1)$.

The section/divisor interface is EGA IV4,21.1.4 p256,21.1.7 and equation21.1.7.1 p257, and Proposition21.2.11 with full proof p262, verified from the original scan. The bundle, its determinant degree and all local multiplicities have been constructed and proved here. The value partition is used only to constrain jets of sections; it is not asserted to define a geometric quotient, an algebra endomorphism, or a new representation multiplicity. This is an explicit geometric obstruction for low-degree families of regular functions on a finite etale scheme, with an honest limitation to that interpolation interface.

After the finite separable change of coefficient field $F\subset F'$,
replace $Y$ and $\mathbb P^1_F$ by $Y_{F'}$ and $\mathbb P^1_{F'}$.
The determinant section again belongs to
$\mathcal O_{\mathbb P^1_{F'}}(Rw2^{w-1})$, is nonzero at zero, and
has exactly the proved lower bounds on its divisor multiplicities at
the $F'$-rational interpolation nodes. Thus finite separable coefficient
extensions preserve this effective-divisor obstruction, including the
$R/Q\ge1-o(1)$ consequence. The extension of coefficient derivations is
the stated classical separability theorem, and the divisor interface
remains EGA IV4,21.1.4,21.1.7 and Proposition21.2.11 at the original
pages cited in NN-C040. This is a base change of the explicitly
constructed bundle calculation; it does not assert an efficient
representation of algebraic coefficients or a general circuit lower
bound.


The rational coefficient-derivative sections on
$Y_{E_0}\times\mathbb P^1_{E_0}$ give a meromorphic squarefree
evaluation matrix. Clearing its Boolean-component column $x$ by
$B_x^{2w}$ constructs a global bundle map
\[
 \mathcal O_{\mathbb P^1}(-2wR)^{2^w}
       \longrightarrow \mathcal O_{\mathbb P^1}^{2^w}.
\]
Its determinant is a nonzero section of
$\mathcal O_{\mathbb P^1}(2w2^wR)$ and has precisely the proved
lower bounds on its orders at the nodes, because every clearing
factor is a unit there. Its effective divisor therefore yields
\eqref{eq:C040-rational}. With a common denominator, multiplication
by $B(T)/B(0)$ reduces the value partitions to the earlier polynomial
bundle statement of degree $a w2^{w-1}$. The same verified EGA IV4
section/divisor locators apply. Thus allowing controlled poles does
not remove this effective-divisor obstruction; neither efficient
construction of the exponentially large matrix nor a geometric
quotient of the summary classes is asserted.




\section{NN-C041 --- Integral permanent circuits for sparse residuals plus a sublinear supplied rank}
\begin{proposition}[NN-C041]
Let $n\ge2$. Suppose an $n\times n$ matrix is supplied in the form
\[
                 A=S+UV,\qquad U\in R^{n\times d},
                 \quad V\in R^{d\times n},
\]
where the allowed support of $S$ has size $s$, the displayed
factorization is part of the input, and $R$ is any commutative
ring. Equivalently take all allowed entries of $S,U,V$ to be
independent variables over $\mathbb Z$. Put
$H=\binom{n+d}{d}$, and let $\theta_{\rm per}(n,s)$ be the
explicit clamped parameter in NN-C038, with its support parameter
equal to $s$. There is a deterministically constructible integral
division-free circuit for $\operatorname{per}(S+UV)$ of size and
bit construction time
\begin{equation}\label{eq:C041-cost}
       \operatorname{poly}(n,d,s)\,H^{O(1)}
                       2^{O(n\theta_{\rm per}(n,s))}.
\end{equation}
The construction can be constant-free, with formal degree at
most $2n^2+2n$, under the convention that constant leaves have
degree one and subtraction has the addition degree rule.
All identities hold over $\mathbb Z$ and specialize to every
commutative ring, including rings with nilpotents.

In particular, if $s\le n^{1+o(1)}$ and $d=o(n)$, the circuit
and its construction have size $2^{o(n)}$, up to the explicit
polynomial input cost. The matrix itself can be dense and of
full rank. No algorithm for finding this decomposition from
an arbitrary matrix is asserted. Numerical evaluation has
the corresponding number of ring operations; a bit-complexity
claim additionally retains the chosen representation and cost
of those ring operations.
\end{proposition}
\paragraph{Exact inputs.}
NN-C038. Uniform integral explicitly listed factor construction, fixed density and Boolean degree; full parameter theta\_per and heavy-row support argument.

E: The Exponential Time Hypothesis is False, supplied LaTeX. Proposition3.4 and Lemma4.7, as fully developed into the all-ring coefficient-independent construction NN-C034/038.

Servedio--Wan, Computing sparse permanents faster, author manuscript18 June2005. Section2 printed p3 first display: Ryser inclusion-exclusion identity; exact formula previously verified in the primary PDF.

EGA IV, fourth part, Publications Mathematiques IHES32 (1967). Printed p.111,18.2.1--18.2.2; base-change equation18.2.1.2; primary PDF page110 (zero-based109).

SGA3, ExposeI. 4.7.2--4.7.3: split-torus representations and character gradings, previously verified for NN-S019.

\begin{proof}
The counting engine is the fully reviewed NN-C038 consequence
of E: a product of $O(n)$ explicitly listed factors with total
Boolean degree $O(n)$ and scopes bounded by $w$ has an integral
formal partition circuit of size
$\operatorname{poly}(n,L)2^{O(n\theta(n,w))}$.
Its circuit identity holds over every coefficient ring, and its
construction depends only on the declared Boolean scopes and
monomial lists. We will use it over an explicitly presented
finite free integer algebra and charge for every coordinate.

\paragraph{An integral moment algebra.}
Let $\mathcal I=\{\alpha\in\mathbb N^d:|\alpha|\le n\}$,
so $|\mathcal I|=H$. The conventions for $d=0$ give one empty
multi-index and $H=1$. Form the free abelian group with basis
\[
             t^\alpha z^{[\beta]},\qquad
                    \alpha,\beta\in\mathcal I.
\]
Give it the multiplication
\begin{equation}\label{eq:C041-algebra}
 (t^\alpha z^{[\beta]})(t^{\alpha'}z^{[\beta']})
  =\binom{\beta+\beta'}{\beta}
           t^{\alpha+\alpha'}z^{[\beta+\beta']}
\end{equation}
when $|\alpha+\alpha'|,|\beta+\beta'|\le n$, and zero otherwise;
the multi-index binomial is the product of the coordinate
binomials. This is a commutative unital algebra $B$ of rank
$H^2$. Associativity follows from
\[
 \binom{\beta+\beta'}{\beta}
 \binom{\beta+\beta'+\beta''}{\beta+\beta'}
 =\binom{\beta'+\beta''}{\beta'}
 \binom{\beta+\beta'+\beta''}{\beta},
\]
an identity of integers, and from the fact that a truncated
total degree can never decrease in a further multiplication.
Thus no division by a factorial is performed. This is the
tensor product of an ordinary truncated polynomial algebra
and a truncated divided-power algebra, but the displayed
rule and its verification suffice as its definition here.

For a vector $v=(v_1,\ldots,v_d)$ over any coefficient ring,
define the element
\[
                  E(v)=\sum_{\beta\in\mathcal I}
                               v^\beta z^{[\beta]}\in B_R.
\]
The ordinary multivariate binomial formula and
\eqref{eq:C041-algebra} give the exact truncated identity
\begin{equation}\label{eq:C041-exponential}
                         E(v)E(w)=E(v+w).
\end{equation}
It is valid in every characteristic. Each coefficient of
degree at most $n$ is the full binomial convolution; all
higher coefficients vanish by the defined truncation.

Define an integer-linear functional, extended $R$-linearly,
\begin{equation}\label{eq:C041-pairing}
 \ell:B_R\longrightarrow R,
 \qquad
 \ell(t^\alpha z^{[\beta]})=
       \begin{cases}1,&\alpha=\beta,\\0,&\alpha\ne\beta.\end{cases}
\end{equation}
This functional is not an algebra homomorphism or an algebra
trace. If $P(t)=\sum_{|\alpha|\le n}c_\alpha t^\alpha$, then
\begin{equation}\label{eq:C041-substitution}
                             \ell(P(t)E(v))=P(v).
\end{equation}
The formula follows immediately from the two basis expansions.

\paragraph{The Boolean partition identity.}
For $x\in\{0,1\}^n$, put
\[
 P_x(t)=\prod_{i=1}^n
       \left(\sum_j S_{ij}x_j+\sum_{a=1}^d U_{ia}t_a\right),
 \qquad v(x)_a=\sum_j V_{aj}x_j.
\]
The polynomial $P_x$ has total $t$-degree at most $n$.
For the $j$th column $v_j$ of $V$, use the unary factor
\[
             G_j(x_j)=1-(1+E(v_j))x_j\in B_R[x_j].
\]
It has value $1$ at zero and $-E(v_j)$ at one. Consequently
\eqref{eq:C041-exponential} gives
\[
       \prod_{j=1}^nG_j(x_j)=(-1)^{|x|}E(v(x)).
\]
Taking the Boolean sum over $B_R$ and then applying $\ell$
therefore yields
\begin{align}\label{eq:C041-ryser}
 (-1)^n\ell\left(\sum_{x\in\{0,1\}^n}
                         P_x(t)\prod_jG_j(x_j)\right)
 &=\sum_x(-1)^{n-|x|}P_x(v(x))\\
 &=\sum_x(-1)^{n-|x|}
               \prod_i\left(\sum_j A_{ij}x_j\right)
   =\operatorname{per}A.\nonumber
\end{align}
The last identity is the classical Ryser formula, in the already
checked Servedio--Wan, \emph{Computing sparse permanents faster},
Section~2, printed p.~3, first display, at
\url{https://www.cs.columbia.edu/~atw12/papers/perm.pdf}.
The following short expansion only checks its specialization to
arbitrary commutative rings, rather than claiming a new identity.
Expand the product as a sum over maps $f:[n]\to[n]$.
The term for $f$ has Boolean sum
$\sum_{X\supseteq\operatorname{im}f}(-1)^{n-|X|}$,
which is zero unless $\operatorname{im}f=[n]$, when it is one.
A surjection between these two $n$-element sets is a permutation.
This proves the identity without division or a characteristic
restriction. It is the same Ryser interface used in NN-C038.

Each of the $n$ row factors in $P_x(t)$ has Boolean degree at
most one and scope equal to the allowed support of its row of
$S$. The $n$ factors $G_j$ are unary and have Boolean degree
one. The dense matrices $U,V$ occur only in coefficients of
$B_R$. Thus the total Boolean degree and factor count are at
most $2n$, regardless of $d$ or the density of $UV$.
Set $w$ to the maximum of one and the residual row support
sizes; constant row factors may be assigned a harmless singleton
declared scope. With this $w$, NN-C038 applies with
its fixed constant $C=2$. Its universal coefficient variables
are then specialized to the stated elements of $B_R$.
Repeated or coincident coefficient values cause no difficulty:
the universal circuit is an integer polynomial identity.

\paragraph{Charging the coefficient algebra.}
Enumerating the basis and its multiplication rule takes
$\operatorname{poly}(n,d,H)$ bit operations. Each nonzero
structure integer in \eqref{eq:C041-algebra} has at most
$O(n)$ bits, since
$\binom{\beta+\beta'}{\beta}\le2^{|\beta+\beta'|}\le2^n$.
An addition in $B_R$ uses $H^2$ coordinate additions, and
a naive multiplication uses $H^{O(1)}$ ring operations and
integer scalar additions. Multiplication by any of these
integers is implemented by doubling and addition, so it
introduces no arbitrary constant leaf or extra formal degree.

The row coefficients have coordinate degree at most one in
the original entries. The unary coefficients use only
monomials $v_j^\beta$ with $|\beta|\le n$; their coordinate
circuits have degree at most $n$ and construction cost
$\operatorname{poly}(n,d,H)$. A coordinate circuit for any
coefficient has constant-free formal degree at most $n$ for
$n\ge2$, including its integer constant coordinates.
Substituting these coordinate circuits into the NN-C038
universal circuit of formal degree at most $2n$, and replacing
algebra operations by the coordinate operations just described,
gives formal degree at most $2n^2$. Applying $\ell$ is only a
sum of selected coordinates. This proves the cost
$\operatorname{poly}(n,d,s)H^{O(1)}2^{O(n\theta(n,w))}$
for bounded residual row scope. The output may have much
smaller ordinary degree; the stated formal bound is sufficient
and does not depend on cancellation to become polynomial.

\paragraph{Removing the maximum-row restriction.}
Let $\lambda=\lceil\log_2 n\rceil$ and
$a=\max(1,\lceil s/n\rceil)$, and declare a residual row heavy
when it contains more than $a\lambda^2$ allowed entries of
$S$. The number $h$ of heavy rows obeys
$h\le s/(a\lambda^2)\le n/\lambda^2$.
Expand the permanent of the actual matrix $A$ along those
rows over all injections into its $n$ columns. This gives
at most $n^h$ summands. The selected entries are actual
$A_{ij}=S_{ij}+\sum_a U_{ia}V_{aj}$; they are not restricted
to the support of $S$.

In each remaining minor of order $k=n-h$, the residual rows
have support at most $a\lambda^2$, and the dense part is still
the displayed product $U_{\rm remaining\ rows}V_{\rm remaining\ columns}$.
Use the preceding construction with truncation degree $k$;
its basis count is $\binom{k+d}{d}\le H$. For large $n$,
$n/2\le k\le n$. Exactly the logarithmic-budget comparison
in the fully proved NN-C038 heavy-row argument bounds every
minor by the parameter $\theta_{\rm per}(n,s)$:
\[
 8(2+1)(a\lambda^2+1)\lceil\log_2 k\rceil
           \le64(a+1)\lambda^3.
\]
The factor $n^h$ contributes $O(n/\log n)$ to the exponent,
which is absorbed by $O(n\theta_{\rm per})$ as in NN-C038.
Multiplying a minor circuit by the $h$ chosen entries adds
at most $2h$ to its formal degree; constructing each entry
uses $O(d+1)$ gates of degree at most two. The total degree
is at most $2k^2+2h\le2n^2+2n$. Fixed small $n$ and empty
minors are evaluated directly, with permanent of order zero
equal to one. This proves \eqref{eq:C041-cost} uniformly.

Finally $s\le n^{1+o(1)}$ implies
$\theta_{\rm per}(n,s)=o(1)$ by NN-C038. For $1\le d=o(n)$,
\[
 \log H\le d\log\left(\frac{e(n+d)}d\right)=o(n),
\]
because $r\log(e(1+r)/r)\to0$ as $r=d/n\to0$.
For $d=0$ the conclusion is immediate. Thus all charged
algebra operations and all heavy-row branches together cost
$2^{o(n)}$ in the asserted mixed regime.
\end{proof}
\paragraph{Status and complexity impact.}
Root-developed complete deduction and complexity-agent whole-proof review passed. Complete deduction-proof review by complexity agent (root-developed proof): passed; imported input theorems are not claimed independently proved. This is a new use of the E-based sparse factor construction
with a fully charged integral moment algebra. It is not a
claim that low rank alone has newly become tractable, and it
does not recognize the displayed decomposition. A diagonal
residual can already have full rank, and adding a rank-one
product can make every matrix entry nonzero. The matrix-order
parameter, supplied decomposition, residual support and rank
cost are all retained. The construction avoids the C040
small-value hypothesis: the global dense sums are stored in
coefficient moments, rather than collapsed to subexponentially
many formal values at interpolation nodes.

\paragraph{Geometric and representation dictionary.}
Over the universal coefficient ring
$R=\mathbb Z[S_{ij},U_{ia},V_{aj}]$, the explicitly constructed
algebra $B_R$ is finite free of rank $H^2$. Its spectrum is a
finite locally free, generally nonreduced thickening of the
coefficient base. Let
\[
 Y=\operatorname{Spec}
 B_R[x_1,\ldots,x_n]/(x_j^2-x_j:1\le j\le n).
\]
This Boolean scheme is finite etale of rank $2^n$ over
$\operatorname{Spec}B_R$. The sum inside
\eqref{eq:C041-ryser} is its multiplication trace of the
specified product section. EGA IV4,18.2.1--18.2.2, printed
p.~111, is the already checked finite-free trace and base-change
interface. The new integer algebra and the section have been
given explicitly, including all of their ranks and costs.

Give $t^\alpha z^{[\beta]}$ the torus character
$\alpha-\beta\in\mathbb Z^d$. The multiplication preserves
this grading. Under SGA3, Expose I,4.7.2--4.7.3, it is the
corresponding diagonal-torus action on $B_R$. The functional
$\ell$ projects to the weight-zero free summand and then
sums its basis coordinates. It is an explicitly specified
linear contraction, not the trace of the nonreduced algebra
$B_R$ and not an asserted ring homomorphism. In particular
nilpotent moment coordinates are retained rather than erased
by an ordinary algebra trace. The final universal polynomial
is the character of the usual free permutation-indexed
module with each entry expanded as $S_{ij}+\sum_aU_{ia}V_{aj}$;
the signed circuit computes that positive final character.
No canonical GCT object or effective extraction of arbitrary
representation multiplicities is being supplied.



\section{NN-C042 --- Growing-locality inversion with arbitrary polynomially many observations}
\begin{proposition}[NN-C042]
Put
\[
 \begin{aligned}
 \lambda&=\lceil\log_2n\rceil,\qquad
 \nu=\max(2,\lceil\log_2\lambda\rceil),\\
 t&=8k+2\lceil\log_2\nu\rceil,\qquad M=2^t,\\
 B&=2\sum_{j=0}^{k-1}M^j
 \end{aligned}
\]
for integers $n\ge2$ and $2\le k\le n$. The following is an explicit sufficient parameter condition:
\begin{equation}\label{eq:C042-gate}
                       ((k+1)B)^4\nu^{12}\le\lambda.
\end{equation}
For every width-at-most-$k$ CNF $F$ on $n$ declared variables with arbitrary explicit clause count and full description length $L$, this condition implies a deterministic integral constant-free construction of its full two-rail satisfying-assignment polynomial
\[
 A_F(z)=\sum_{x\models F}\prod_{i=1}^n z_{i,x_i}
\]
in size and bit construction time
\begin{equation}\label{eq:C042-time}
             \operatorname{poly}(n+L)\,2^{O(n/\nu)},
\end{equation}
with an absolute constant in the exponent and formal degree at most $2n$. The polynomial identity holds over $\mathbb Z$ and every commutative ring. Exact integer counting has the same time bound. No restriction on clause density is imposed.

Condition~\eqref{eq:C042-gate} holds eventually for every family with
$k=o(\sqrt{\log\log n})$. More explicitly it holds when
$\nu\ge2^{16}$ and $2\le k\le\lfloor\sqrt\nu/8\rfloor$.
Thus this is a genuinely growing-width statement, with all width dependence exposed rather than absorbed in a fixed-width constant.

The following stronger threshold choice preserves all the conclusions above. Define
\[
 \begin{aligned}
 c_0&=1,\qquad c_{j+1}=\sum_{i=0}^{j}c_i c_{j-i},\\
 \beta_j&=4c_{j-1}(12a)^{j-1}\quad(1\le j<k),\\
 C_a&=\sum_{j=1}^{k-1}\beta_j,\qquad
 h_a=\lceil\log_2(4aC_a)\rceil.
 \end{aligned}
\]
Choose the first power of two $a\ge k$ satisfying $k\nu h_a\le a$,
and put $B_a=2(1+aC_a)$. This search terminates with
$a<16(k\nu)^4$. The alternative sufficient condition is
\[
                         ((k+1)B_a)^4\nu^{12}\le\lambda.
\]
It holds eventually for every
$k=o(\log\log n/\log\log\log n)$, and explicitly whenever
\[
 \nu\ge2^{16},\qquad
 2\le k\le
   \left\lfloor\frac{\nu}{64\lceil\log_2\nu\rceil}\right\rfloor.
\]
The preceding geometric-threshold statement remains a valid weaker bound.

Let $f:\{0,1\}^n\to\{0,1\}^m$ be given by scopes of size at most $k$ and explicit local truth tables, with polynomial description length and arbitrary polynomially bounded output length $m$. In the same growing-locality range, every target fiber can be counted exactly, image membership decided, and a preimage found when one exists, in deterministic time
\[
                              2^{O(n/\log\log n)}.
\]
The parameter remains the original $n$ seed bits, regardless of $m$. If $m>n$, image membership distinguishes $f(U_n)$ from $U_m$ with advantage at least $1-2^{n-m}$ in this time. These conclusions exclude exponential-in-$n$ security in this specified local model, while leaving polynomial-time one-wayness and pseudorandomness open.
\end{proposition}
\paragraph{Exact inputs.}
E: The Exponential Time Hypothesis is False, supplied LaTeX. AppendixA in full: LemmasA.1--A.3, exact branching priority, insertion bound, Gamma and C formulas, tree-node bound and occurrence bound B; lines3636--3885 of supplied main\_full\_version.tex. Lemma6.3 is used through these explicit formulas, not through hidden fixed-k constants.

NN-C039. Uniform integral actual-degree construction for at most Cn explicit monomial-list factors; here C=2 is fixed while the grouped scope and degree vary.

NN-C022. Forbidden-local-assignment CNF encoding, prefix-count inversion and image-support distinguishing, retaining original seed variables.

Calabro--Impagliazzo--Paturi, A Duality between Clause Width and Clause Density for SAT (CCC 2006). Sections6.1--6.3, sixth page of the primary author PDF: arbitrary threshold priority and the improved convolution analysis. The exact constants, search, disjoint counting interpretation and uniform growing-width estimates are proved here.

EGA IV, fourth part, Publications Mathematiques IHES32 (1967). Printed p.111,18.2.1--18.2.2; base-change equation18.2.1.2; primary PDF page110 (zero-based109).

\begin{proof}
We expose the width dependence in E's AppendixA before applying any counting construction. Empty clauses give the zero polynomial, and a formula with no clauses gives $\prod_i(z_{i,0}+z_{i,1})$. Retain all $n$ declared variables throughout; do not propagate away unit-constrained variables. The case of locality below two is covered by using $k=2$.

\paragraph{Uniform sparsification parameters.}
Use the branching rule and normalization of E LemmasA.1--A.3 with the explicitly chosen integer $M=2^t$. Define
\[
 C=4\sum_{j=0}^{k-2}(3M)^j,
 \qquad \Gamma=2(3^{k-1}-1),
 \qquad h=\lceil\log_2(4CM)\rceil.
\]
Those lemmas and the subsequent node-encoding argument in AppendixA give depth at most $Cn$, at most $\Gamma n/M$ choices of the designated second child on a path, and at most
\[
                 (Cn+1)^2\,2^{\Gamma h n/M}
\]
nodes whenever $M\ge\Gamma$. The leaf satisfying sets are pairwise disjoint and partition the parent's satisfying set by LemmaA.1. Each variable occurs at most $B=2\sum_{j=0}^{k-1}M^j$ times in every leaf. These are the exact displayed estimates in the supplied proof; they hold for the present integer $k$ without an unspecified constant depending on $k$.

Here $\Gamma\le2^{2k}$ and
$C\le8(3M)^{k-2}$ for $k\ge2$ and $M\ge2$. Consequently
\[
 h\le k(t+4),\qquad
 t\le8k+2\log_2\nu+2,
 \qquad M\ge2^{8k}\nu^2.
\]
In particular $M\ge\Gamma$. Since $\log_2\nu\le\nu$ and $\nu\ge2$,
\[
 \frac{\Gamma h\nu}{M}
 \le\frac{k(8k+2\log_2\nu+6)}{2^{6k}\nu}
 \le\frac{k(4k+5)}{2^{6k}}\le1.
\]
The final elementary inequality holds for every integer $k\ge2$, for example from $k\le2^k$ and $4k+5\le9\cdot2^k$.
Thus the node bound is at most $(Cn+1)^2 2^{n/\nu}$, with the exact rational choice $\sigma=1/\nu$. Computing $t,M,B,C,\Gamma,h$ requires polynomially many bit operations in $k+\log\nu$ and their explicit bit lengths; their powers and comparisons are integer computations.

For completeness we also charge the node work uniformly in $k$. There are at most
$U=(2n+1)^{k+1}$ possible clauses of width at most $k$ after removing repeated literals and tautologies. A direct implementation can store the current list, enumerate at most $2^k$ subsets of each clause to form the $(c,H)$ counters, scan for the priority choice, and restore the antichain by pairwise subset tests. The threshold lengths are $O(kt)$ bits, and counters have $O(k\log(n+2))$ bits. These operations use
\[
                  U^{O(1)}\operatorname{poly}(n,k,t)
\]
bit operations per node, with an absolute polynomial exponent; this bound also covers deterministic tie breaking. Initial normalization takes a polynomial in the supplied length $L$. No appeal to an unspecified $\operatorname{poly}_k$ running-time constant is needed.

\paragraph{Packing a leaf into a fixed number of factors per variable.}
A nonempty leaf has at most $Bn$ clauses. Divide its clause list deterministically into at most $n$ groups, each containing at most $B$ clauses. For a clause $C$, let
\[
 v_C(x)=\prod_{\ell\in C}(1-\ell(x)),
 \qquad I_C(x)=1-v_C(x)
                         \quad\hbox{over }\mathbb Z,
\]
where a positive literal has affine value $x_i$ and a negative literal has affine value $1-x_i$. On the Boolean cube, $I_C$ is exactly its zero-one satisfaction indicator. For each group $G$ insert the factor
$\Phi_G(x)=\prod_{C\in G} I_C(x)$.
It has declared scope and ordinary degree at most $kB$, and its product with the other group factors is the exact leaf indicator.

Expand each group explicitly as a monomial list. Each $I_C$ has at most $1+2^k$ signed monomial occurrences, so a group has at most
\[
                    (1+2^k)^B\le2^{(k+1)B}
\]
occurrences. Keep repeated occurrences instead of requiring cancellation or symbolic equality tests; each listed coefficient is $1$ or $-1$. Exponents are at most $B$ and are stored in binary. Introducing an independent coefficient for every occurrence and then specializing it to its known sign fits NN-C039 exactly, including coincident monomials.

Add the $n$ unary factors
\[
                   W_i(x_i)=(1-x_i)z_{i,0}+x_i z_{i,1}.
\]
There are at most $2n$ factors. With
$w=\min(n,kB)$, their scopes are at most $\max(1,w)$ and their actual Boolean-reduced degree sum is at most $n(w+1)$. Their partition polynomial is the leaf's full two-rail satisfying-assignment polynomial over $\mathbb Z$. NN-C039 applies with its density constant $2$, which is fixed. Its universal coefficient inputs become signs or linear two-rail expressions, so its formal degree bound remains at most $2n$. Summing over the sparsification leaves gives exactly $A_F$, because those leaves partition the satisfying assignments. No assignment multiplicities are introduced by the sum.

\paragraph{The uniform cost bound.}
Under condition~\eqref{eq:C042-gate},
\[
             (k+1)B\le\lambda^{1/4}/\nu^3,
             \qquad w\le\lambda^{1/4}.
\]
The expanded factor input length is therefore at most
$\operatorname{poly}(n,k,t,B)2^{(k+1)B}=n^{O(1)}$ for all sufficiently large $n$, with an absolute exponent. The same condition gives $k\le\lambda^{1/4}$ and $\log C=O(k^2+k\log\nu)$, so the logarithm of the node-processing factor, of the depth polynomial $(Cn+1)^2$, and of all parameter work is
\[
                   O(k\lambda+k^2+k\log\nu+\lambda)
                            =O(\lambda^{5/4})
                            =o(n/\nu).
\]
The finitely many smaller admissible sizes are absorbed by an absolute constant in the claimed bound.

For $w\le\sqrt\lambda$, NN-C039's explicit parameter satisfies
$\theta(n,w)=O(\nu^2/\sqrt\lambda)$ with an absolute constant at density two. The actual-degree multiplier is at most $w+2$. Hence the logarithm of its exponential cost per leaf is at most
\[
 O\left(\frac{n(w+2)\nu^2}{\sqrt\lambda}\right)
  =O\left(\frac{n}{\nu\lambda^{1/4}}
                 +\frac{n\nu^2}{\sqrt\lambda}\right)
  =O(n/\nu).
\]
Together with the at most $(Cn+1)^2 2^{n/\nu}$ nodes and the input-normalization cost, this proves~\eqref{eq:C042-time}. All constants here are uniform in the admissible $k$; the density supplied to NN-C039 never grows with $k$.

To extract the exact count, specialize the two-rail weights to one and evaluate the resulting integer circuit modulo $2^{n+1}$. The true count lies between zero and $2^n$, so its residue in this range is unique. Each operation then has polynomial bit cost in $n$. This avoids any unstated bound on the sizes of intermediate integer circuit values. The circuit identity itself remains an exact integral identity, not only a congruence.

\paragraph{An explicit growing-locality range.}
The displayed choice of $B$ satisfies
\[
 \log_2((k+1)B)
 \le\log_2(k+1)+2+8k(k-1)
                          +2(k-1)\lceil\log_2\nu\rceil.
\]
If $k=o(\sqrt\nu)$, the right side is $o(\nu)$, and condition~\eqref{eq:C042-gate} follows eventually because
$\log_2\lambda\ge\nu-1$ once $\nu>2$.
For the stated explicit constant range $k\le\sqrt\nu/8$, a simpler upper bound on the logarithm of the left side of~\eqref{eq:C042-gate} is
\[
       8+\nu/2+\sqrt\nu(\log_2\nu+1)+16\log_2\nu.
\]
For $\nu\ge2^{16}$, the last two terms together with eight are at most $\nu/4$: separately,
$\sqrt\nu(\log_2\nu+1)\le\nu/8$ and
$16\log_2\nu+8\le\nu/8$.
These inequalities hold at $2^{16}$ and remain valid as $\nu$ grows since the relevant logarithm-to-power ratios decrease there. The whole bound is at most $3\nu/4\le\nu-1\le\log_2\lambda$. This proves the explicit sufficient range. An implementation can check~\eqref{eq:C042-gate} exactly, using integer powers, and use exhaustive search outside it; no unproved assertion about a fixed-width constant is involved.

\paragraph{Preimages, search and image support.}
Given a local map and target $y$, exclude every local table assignment whose output differs from $y_j$ by its single forbidden-assignment clause, exactly as in NN-C022. This uses at most $m2^k$ clauses of width at most $k$, and exactly the original $n$ variables. The satisfying set is precisely $f^{-1}(y)$. The expanded description is polynomially bounded whenever the supplied local tables and output length are polynomially bounded; their explicit reading cost is retained in~\eqref{eq:C042-time}. Thus the exact fiber count and image membership have the claimed original-seed bound.

For inversion, add a unit constraint fixing the next seed bit to zero, count that prefix, and retain it if the count is positive; otherwise fix the bit to one. Keep all $n$ variables declared even late in the search, so the locality gate is checked at the original $n$, not at a shrinking residual dimension. There are at most $n+1$ count calls. All added constraints have width one, and no density restriction was imposed. This proves deterministic inversion of every nonempty fiber. Finally, the exact image indicator accepts $f(U_n)$ with probability one and accepts $U_m$ with probability $|\operatorname{im}f|/2^m\le2^{n-m}$. Their probability difference proves the stated distinguishing advantage. These are subexponential-time algorithms and do not imply polynomial-time attacks.

\paragraph{A stronger uniform sparsification threshold.}
The growing-locality range in NN-C042 can be enlarged by changing
its threshold sequence. Retain $\lambda,\nu,n,k$ from that theorem.
Define positive integers $c_j$ by
\[
 c_0=1,\qquad c_{j+1}=\sum_{i=0}^{j}c_i c_{j-i}.
\]
For a power of two $a\ge k$, put
\[
 \begin{aligned}
 \beta_j&=4c_{j-1}(12a)^{j-1}\quad(1\le j<k),\\
 \theta_0&=1,\qquad \theta_j=a\beta_j\quad(1\le j<k),\\
 C_a&=\sum_{j=1}^{k-1}\beta_j,\qquad
 h_a=\lceil\log_2(4aC_a)\rceil.
 \end{aligned}
\]
Take the first power of two $a\ge k$ satisfying the exact integer test
\begin{equation}\label{eq:C042-Catalan-search}
                         k\nu h_a\le a.
\end{equation}
Set
\[
                         B_a=2(1+aC_a).
\]
All conclusions of NN-C042 hold with its parameter condition replaced by
\begin{equation}\label{eq:C042-Catalan-gate}
                       ((k+1)B_a)^4\nu^{12}\le\lambda.
\end{equation}
The construction and integer search are uniform in $k$. They satisfy
\[
 a<16(k\nu)^4,
 \qquad B_a\le(48a)^{k-1},
 \qquad
 \log_2 B_a\le(k-1)(10+4\log_2(k\nu)).
\]
Consequently condition~\eqref{eq:C042-Catalan-gate} holds eventually
for every $k=o(\nu/\log\nu)$. An explicit sufficient range is
\begin{equation}\label{eq:C042-Catalan-range}
 \nu\ge2^{16},\qquad
 2\le k\le
        \left\lfloor\frac{\nu}{64\lceil\log_2\nu\rceil}\right\rfloor.
\end{equation}
Thus the full integral arbitrary-density counting construction and
the original-seed inversion/image test apply in this stronger
growing-locality range, in time $\operatorname{poly}(n+L)2^{O(n/\nu)}$
and with the same formal degree bound $2n$.

\paragraph{Why the marked-degree argument remains valid.}
Use exactly E's normalized clause families and deterministic priority:
choose the smallest width $c$ with a branch, then the largest size
of its shared nonempty proper literal subset $H$. The branching
condition is now
\[
              \deg_c(H)\ge\theta_{c-|H|}.
\]
The two children are still $G\wedge H$ and
$G\wedge\bigwedge P_i\wedge\bigwedge_{\ell\in H}\{\neg\ell\}$.
Their satisfying sets partition the parent's, on the same declared
variables. E LemmaA.1's insertion and survival proof applies here:
it uses only that at least $\theta_j$ distinct width-$j$ petals are
selected and that no old clause is contained in a new strict subset,
not the formula $M^j$ for the threshold. In particular the second
child inserts at least $\theta_j$ new width-$j$ clauses when its
petal width is $j$; the added units do not subsume these petals.

We verify the required varying-threshold marked-degree bound directly.
For $1\le i<j<k$ and a literal set $S$ of size $i$, fewer than
$2\theta_{j-i}$ currently present marked width-$j$ clauses contain
$S$. Initially there are no marked clauses. Before a step which
inserts width-$j$ clauses, the selected width is larger than $j$.
The smallest-width priority implies that fewer than $\theta_{j-i}$
current width-$j$ clauses contain $S$. If the first child is selected,
at most one is added, leaving fewer than $2\theta_{j-i}$.
For the second child, suppose some new width-$j$ petal contains $S$.
Then $S$ is disjoint from $H$, and $H\cup S$ is a larger proper
subset of the selected width-$c$ clause. The largest-shared-subset
priority therefore implies
\[
 \#\{P_i:S\subseteq P_i\}
       =\deg_c(H\cup S)<\theta_{c-|H\cup S|}
                          =\theta_{j-i}.
\]
The other inserted clauses are units. Adding the old and new bounds
proves the asserted marked-degree bound, and deletion can only lower
it. This proof needs positive integer thresholds but no multiplicative
relation among them.

Let $T_j$ count all width-$j$ insertions along a path prefix, including
those later deleted. Units never disappear before a terminal empty
clause, so $T_1\le2n$. For $j\ge2$, summing the marked-degree bound
at singleton literals gives at most $4n\theta_{j-1}$ current marked
width-$j$ clauses. Charge every deletion of a marked width-$j$ clause
to a newly inserted strict subset of width $i<j$. The same bound
shows that one such insertion receives fewer than $2\theta_{j-i}$
charges. Hence the exact recurrence needed is
\begin{equation}\label{eq:C042-Catalan-recurrence}
       T_j\le4n\theta_{j-1}
                       +2\sum_{i=1}^{j-1}\theta_{j-i}T_i.
\end{equation}
No arbitrary-width extension of E's fixed-threshold recurrence has
been assumed: the preceding marking and deletion argument proves it.

\paragraph{Bounding insertions and the number of branches.}
The recursively defined $c_j$ count plane full binary trees with $j$
internal vertices, by splitting a tree at its root. A preorder record
of its internal versus leaf vertices has length $2j+1$ and ends in a
leaf; deleting that final symbol embeds this set into binary words
of length $2j$. Thus $c_j\le4^j$. Also $c_{j-1}\le c_j$ by the
recurrence, for $j\ge1$.

We claim $T_j\le\beta_j n$. For $j=1$ this follows from
$T_1\le2n\le4n$. Assuming it for smaller widths, equation
\eqref{eq:C042-Catalan-recurrence} gives, for $j\ge2$,
\[
 \begin{aligned}
 T_j/n
 &\le4a\beta_{j-1}
                       +2a\sum_{i=1}^{j-1}\beta_{j-i}\beta_i\\
 &=16a(12a)^{j-2}c_{j-2}
       +32a(12a)^{j-2}
                     \sum_{i=1}^{j-1}c_{j-i-1}c_{i-1}\\
 &\le48a(12a)^{j-2}c_{j-1}
   =\beta_j.
 \end{aligned}
\]
The convolution is exactly $c_{j-1}$ by its definition.
Every branch inserts a clause of width below $k$, so depth is at
most $C_a n$. A second-child branch with petal width $j$ spends at
least $\theta_j=a\beta_j$ width-$j$ insertions. There are at most
$n/a$ such choices for each $j$, and at most $kn/a$ in all.
Since $a\ge k$, this is at most $n\le C_a n/4$.

A node is determined by its depth and the positions of the
second-child choices. The same elementary node encoding used in E's
AppendixA therefore bounds all nodes by
\[
             (C_a n+1)^2\,2^{(kn/a)h_a}
                       \le(C_a n+1)^2\,2^{n/\nu},
\]
where the last inequality is the exact test
\eqref{eq:C042-Catalan-search}. If $kn/a<1$, only the all-first-child
path is possible and this bound remains true. At a leaf, one literal
occurs in at most one unit clause and in fewer than $\theta_{c-1}$
width-$c$ clauses for every $c\ge2$. Thus each variable has at most
$2(1+\sum_{j=1}^{k-1}\theta_j)=B_a$ occurrences. This proves the
required disjoint sparsification with its full explicit parameters.

\paragraph{The integer parameter search is bounded.}
Since $c_j\le4^j$,
\[
 C_a\le4\sum_{j=0}^{k-2}(48a)^j
          \le8(48a)^{k-2},\qquad
 h_a\le k(\log_2a+6).
\]
The logarithm of a power of two is an integer. Put $z=k\nu\ge4$ and
consider the candidate
$a_0=2^{4\lceil\log_2z\rceil}$, so
$z^4\le a_0<16z^4$. At this candidate,
\[
 k\nu h_{a_0}
   \le k^2\nu(4\log_2z+10)
   \le \frac{z^2}{2}(4\log_2z+10)
   \le z^4\le a_0.
\]
The penultimate inequality holds for $z\ge4$; it follows by checking
$4\log_2z+10\le2z^2$ at four and noting that the right-minus-left
function increases thereafter. Also $a_0\ge k$.
Thus the exact search terminates by this candidate and its first
successful value satisfies the asserted bound on $a$.
It examines $O(\log(k\nu))$ powers. The integers $c_j,\beta_j,C_a,B_a$
have $O(k\log(k\nu))$ bits, so direct recurrence and repeated
multiplication implement the search in $\operatorname{poly}(k,\log\nu)$
bit operations.

Furthermore
\[
 B_a=2(1+aC_a)\le2+16a(48a)^{k-2}
                          \le(48a)^{k-1}.
\]
Combined with $a<16(k\nu)^4$, this gives the stated logarithmic
bound on $B_a$. If $k=o(\nu/\log\nu)$, then $k<\nu$ eventually
and $k\log(k\nu)=o(\nu)$, so
$4\log_2((k+1)B_a)+12\log_2\nu=o(\nu)$.
Condition~\eqref{eq:C042-Catalan-gate} follows eventually.

For the explicit range~\eqref{eq:C042-Catalan-range}, put
$L_\nu=\lceil\log_2\nu\rceil$. Then $k+1\le\nu$ and
$\log_2(k\nu)\le2\log_2\nu$, whence
\[
 \begin{aligned}
 4\log_2((k+1)B_a)+12\log_2\nu
 &\le16\log_2\nu+40k+32k\log_2\nu\\
 &\le\nu/16+5\nu/128+\nu/2
  =77\nu/128\le\nu-1\le\log_2\lambda.
 \end{aligned}
\]
Here $\nu\ge2^{16}$ gives $L_\nu\ge16$ and
$16\log_2\nu\le\nu/16$. This proves the sufficient range using
only the exact selected parameters.

\paragraph{Packing, counting and original-seed consequences.}
Each resulting leaf has at most $B_a n$ clauses. Pack it into at most
$n$ products of at most $B_a$ exact integer clause indicators, with
scope and degree at most $kB_a$ and at most $2^{(k+1)B_a}$ signed
monomial occurrences per group, and add the $n$ unary weight factors.
This is exactly the construction in NN-C042 with $B_a$ in place of
$B$. Under~\eqref{eq:C042-Catalan-gate}, its lists are polynomial in
$n$, its density is two, and NN-C039 gives the same
$2^{O(n/\nu)}$ integral circuit bound and formal degree $2n$.
The node work is still $(2n+1)^{O(k)}$ times a polynomial in the
explicit parameter lengths; on the integer gate its logarithm is
$o(n/\nu)$ as before. All constants are therefore uniform in the
new admissible widths. Disjointness, exact integer trace recovery,
forbidden-assignment clauses, and prefix constraints retain the
original $n$ seed bits exactly as proved in NN-C042.

The threshold convolution is motivated by Calabro--Impagliazzo--Paturi,
\emph{A Duality between Clause Width and Clause Density for SAT},
Section6.1--6.3, sixth page of the primary author PDF. That source
sketches the improved convolution analysis. The complete
varying-threshold proof, integer constants, parameter search and
counting-compatible interpretation used here have been supplied
explicitly; no unspecified ``sufficiently small'' parameter bound
from that sketch is used.


\end{proof}
\paragraph{Status and complexity impact.}
Uniform growing-width sparsification, integral counting and original-seed inversion are proved. Complete deduction-proof review by root: passed; imported input theorems are not claimed independently proved. This removes the fixed-locality restriction for arbitrary polynomially many local observations. The current proved range is $k=o(\log\log n/\log\log\log n)$, with the exact integer Catalan-threshold gate and explicit constant range above. The original seed length remains the exponential parameter, with no verifier-wire or output-length blowup. The theorem also gives the full integral weighted polynomial, not merely a parity value. It excludes exponential-in-seed security for these growing-locality maps while leaving ordinary polynomial-time security, logarithmic-locality arbitrary predicates, and unrestricted Boolean circuit inversion open. It uses E and its integral symbolic consequences, without A or M.

\paragraph{Geometric and representation dictionary.}
Let
\[
 X_n=\operatorname{Spec}
       \mathbb Z[x_1,\ldots,x_n]/(x_i^2-x_i:1\le i\le n).
\]
A local Boolean map gives a morphism $X_n\to X_m$ by integer multilinear idempotents, and a Boolean target section $y$ has finite split fiber $Y_y$ of rank $|f^{-1}(y)|$. For the universal two-rail section over the coefficient base $\mathbb Z[z]$,
\[
            w=\prod_i((1-x_i)z_{i,0}+x_i z_{i,1}),
\]
the regular function
$\operatorname{Tr}_{(\mathcal O(Y_y)\otimes_{\mathbb Z}\mathbb Z[z])/\mathbb Z[z]}(w)$ is its full preimage polynomial. Under either the original exact parameter condition~\eqref{eq:C042-gate} or the stronger Catalan-threshold condition~\eqref{eq:C042-Catalan-gate}, this function has an integral constant-free circuit of size~\eqref{eq:C042-time} and formal degree at most $2n$, compatible with every commutative-ring base change. In particular the conclusion holds for $k=o(\log\log n/\log\log\log n)$, or in the explicit range~\eqref{eq:C042-Catalan-range}, with polynomially many output coordinates.

E's disjoint sparsification is a disjoint-union decomposition of this Boolean fiber on the original $n$ coordinates. Grouping clauses changes only the presentation of its acceptance idempotent; it does not change the scheme or introduce new witness coordinates. Trace adds over the disjoint components and is then computed by NN-C039's integral construction. The rank can be recovered exactly from the computed integer trace, while reduction modulo two retains only its parity. The source-interface locators are EGA IV4,18.2.1--18.2.2, printed p111, with trace base change in equation18.2.1.2. The image test concerns the support of a morphism between finite split schemes; no purity, zeta, or arbitrary representation-multiplicity extraction claim is made. This is an effective finite-scheme theorem for the specified growing local presentations, and not a polynomial-time consequence of cohomological purity.


The finite split fibers, acceptance idempotents and two-rail trace
functions of NN-C042 are unchanged. The new threshold tree gives a
finer effective presentation by disjoint Boolean subschemes on the
same original coordinates. Under~\eqref{eq:C042-Catalan-gate}, and
in particular in range~\eqref{eq:C042-Catalan-range}, their integral
trace functions have the same uniform constant-free circuit and
formal-degree bounds. Exact rank extraction and image support of
$X_n\to X_m$ therefore have the stronger growing-locality scope.
The exact trace/base-change interface remains EGA IV4,18.2.1--18.2.2,
p111, equation18.2.1.2. This alters the effective presentation bound,
not the weight-zero purity of the finite fibers, and supplies no
polynomial-time or general GCT multiplicity algorithm.




\section{NN-C043 --- Unconditional bounded-output simulation with rational radical constants of fixed prime degree}
\begin{proposition}[NN-C043]
Let $C$ be a division-free arithmetic circuit with Boolean inputs and operations $+,-,\times$. Every constant leaf is an explicitly written rational number or an explicitly signed principal square root of an explicitly written rational number. The principal square root of a negative rational means $i$ times the positive square root of its absolute value. Constant expressions formed from such leaves may themselves be given by arithmetic subcircuits, counted in the full description length $L$.

Suppose that on every Boolean input $x\in\{0,1\}^n$ the output is a rational integer $g(x)$ satisfying $|g(x)|\le2^w$, where $w$ is a supplied nonnegative integer. Given the circuit, its literal constant descriptions, the bound $w$, and $x$, one can compute $g(x)$ deterministically in $\operatorname{poly}(L+n+w)$ bit operations. No formal-degree bound, global number-field degree bound, rational-squareclass independence promise, integer factorization, or Riemann-hypothesis input is required.

Consequently polynomial-size, polynomial-description circuit families of this kind with polynomial output bit bounds give functions in $\mathrm{FP}/\mathrm{poly}$; polynomial-time constructible descriptions give functions in $\mathrm{FP}$. If all $\#\mathrm P$ counting families have such descriptions, then $\#\mathrm P\subseteq\mathrm{FP}/\mathrm{poly}$, with $\#\mathrm P=\mathrm{FP}$ in the uniform case. The statement concerns promised integer outputs; it is not an algorithm for comparing general sums of square roots.

\emph{Fixed odd-prime version.} Fix an odd prime $\ell$. Replace the square-root leaves of NN-C043
by explicitly chosen $\ell$th roots of explicitly written rational
numbers. A root is encoded by the rational radicand and an exponent
$j\in\{0,\ldots,\ell-1\}$: its value is
$\zeta_\ell^j$ times the unique real $\ell$th root, where
$\zeta_\ell=\exp(2\pi i/\ell)$. Rational leaves and a leaf
$\zeta_\ell$ are also allowed. Constant expressions use the same
division-free operations and count toward the total description
length $L$.

If the circuit output on every Boolean input is a rational integer
$g(x)$ with $|g(x)|\le 2^w$, the output on a given input is computable
deterministically in $\operatorname{poly}_\ell(L+n+w)$ bit operations.
The constructed coefficient field has degree $(\ell-1)\ell^r$,
exponential in the number $r$ of independent radical generators.
There is no independence promise, factorization oracle, formal-degree
bound, or Riemann-hypothesis assumption. The uniform and advice-family
consequences of NN-C043 hold with this enlarged, fixed-$\ell$
presentation class.
\end{proposition}
\paragraph{Exact inputs.}
NN-C020. Background comparison: bounded scalar-output simulation from a polynomial-rank explicit coefficient algebra; the new theorem allows exponential global degree.

Rosser--Schoenfeld, Approximate formulas for some functions of prime numbers, Illinois J. Math.6 (1962),64--94. Printed p69, Corollary to Theorem3, equation(3.13): the nth prime is less than n(log n+log log n) for n>=6; original page image directly read.

J.S.Milne, Fields and Galois Theory, version5.10 (September2022). Theorem5.30 with complete proof, printed pp75--76; Example5.31(b),p76: rational squareclasses. Fixed odd-prime extension also uses cyclotomic Lemma5.9,p65 and Theorem5.10,p66, and finite-field Propositions4.19--4.20,p53.

Keith Conrad, Hensel's Lemma. Theorem2.1,pp1--3, including the complete digit-lifting proof and equation(2.4); Remark2.3,p3 identifies the unique finite-precision lift.

D.J.Bernstein, Factoring into coprimes in essentially linear time, Journal of Algorithms54 (2005),1--30. Abstract; Algorithm18.1 and Theorem18.2,p24 of the 2004.04.04 author manuscript, retrieved from indexed primary text. Only polynomial-time coprime refinement is needed; its elementary size invariant is given below. Full PDF retrieval returned HTTP502, so no whole-paper source reading is claimed.

Stacks Project, Local structure of etale ring maps. Definition10.144.1,tag00UB; Lemma10.144.2(1)--(2),tag00UC: monic simple-root presentations are etale and remain so under base change.

EGA IV, fourth part, Publications Mathematiques IHES32 (1967). Printed p.111,18.2.1--18.2.2; base-change equation18.2.1.2; primary PDF page110 (zero-based109).

\begin{proof}
All roots and rational constants in this proof are those explicitly encoded in the input. The algorithm does not recover a presentation from an arbitrary complex constant. Zero radical leaves are replaced by zero immediately.

\paragraph{A consistent radical presentation without factorization.}
Write each nonzero radical radicand as $A_i/B_i$, where $A_i$ is a nonzero integer and $B_i$ is a positive integer. Put $N_i=|A_i|B_i$. Its encoded root is, up to its supplied sign,
\[
 \sqrt{A_i/B_i}=
 \begin{cases}
  \sqrt{N_i}/B_i,&A_i>0,\\
  i\sqrt{N_i}/B_i,&A_i<0.
 \end{cases}
\]
The total bit length of the $N_i$ and denominators is $O(L)$.

Coprime refinement is classical; the cited Bernstein algorithm gives a stronger running time than is needed here. To make the present representation cost explicit, maintain a set of distinct active integers greater than one, initially the distinct $N_i>1$, and exponent vectors expressing every $N_i$ as a product of their powers. Whenever two distinct active values $u,v$ have $d=\gcd(u,v)>1$, replace them by $d,u/d,v/d$, discarding ones and merging equal values. An old factor $u^{e_u}v^{e_v}$ becomes
\[
                 d^{e_u+e_v}(u/d)^{e_u}(v/d)^{e_v}.
\]
Thus the exponent vectors are updated exactly, including any merged columns. The product of the distinct active values decreases by a factor at least two at every such refinement: before merging, the replacing product is $uv/d$, and merging only decreases it further. Its initial logarithm is $O(L)$. Hence there are only $O(L)$ refinements, at most $O(L)$ active values, and all integer values have $O(L)$ bits. Every exponent is at most $\log_2N_i$, by the product identity in its row. Gcd computation, exact division and vector updates therefore take polynomial time. At termination the active values are pairwise coprime.

Test each active value for being a perfect square by exact integer square root. Absorb square atoms into rational integer factors. For each remaining atom $b_j$, split every exponent into its even and odd parts. This constructs, in polynomial time, exact identities
\[
        N_i=c_i^2\prod_{j=1}^r b_j^{\epsilon_{ij}},
        \qquad c_i\in\mathbb Z_{>0},\quad
        \epsilon_{ij}\in\{0,1\},
\]
where the $b_j>1$ are pairwise coprime nonsquares and all displayed integers have polynomial bit length. The equality of positive real roots follows because both sides are positive. Negative radicands are handled by one additional generator with square $-1$.

The classes of the remaining positive $b_j$ are independent in $\mathbb Q^\times/\mathbb Q^{\times2}$: each nonsquare has an odd prime valuation, and pairwise coprimality prevents any other atom from cancelling that valuation. The class of $-1$, when used, is independent of them because a nonzero rational square is positive. This proof of independence is a consequence of the constructed coprime presentation; it does not assume prime factors are computed. Denote all the resulting radicands, including $-1$ if needed, by $b_1,\ldots,b_r$ after renumbering.

Let $D$ be the product of the positive denominators of the rational leaves and of the $B_i$. It has $O(L)$ bits. Every encoded constant is now represented by a polynomial-size expression in rational constants from $\mathbb Z[D^{-1}]$ and formal generators $u_j$ with $u_j^2=b_j$. Define
\[
 \mathcal O=\mathbb Z[u_1,\ldots,u_r]/(u_j^2-b_j:1\le j\le r).
\]
The monomials $\prod_j u_j^{\delta_j}$, with $\delta_j\in\{0,1\}$, form a free integral basis. By Kummer theory, specifically Milne Theorem5.30 and Example5.31(b), the field generated by the corresponding complex square roots has degree $2^r$ over $\mathbb Q$. Thus the map from $\mathbb Q\otimes\mathcal O$ to that field is an isomorphism, and the chosen complex interpretation of $\mathcal O[D^{-1}]$ is injective. This basis and its exponential multiplication table are used only to prove injectivity; the algorithm never constructs them.

For a Boolean input $x$, the circuit evaluates symbolically to some $v_x\in\mathcal O[D^{-1}]$. The promised integer output and injectivity imply the exact algebra identity
\begin{equation}\label{eq:C043-scalar-identity}
                             v_x=g(x)\cdot1
                           \quad\hbox{in }\mathcal O[D^{-1}].
\end{equation}
Therefore every compatible specialization of the generators preserves that integer value, even if the selected residue roots do not preserve the original real signs. This is where a consistent global radical presentation is needed.

\paragraph{A small prime with no splitting condition.}
Form
\[
                    W=2D\prod_i N_i.
\]
It is a positive integer with $O(L)$ bits. Write $H$ for its binary bit length and put $s=\max(6,H+1)$. The integer $W$ has at most $H-1$ distinct prime divisors. The classical bound in Rosser--Schoenfeld, printed p69, equation(3.13), gives
\[
                    p_s<s(\log s+\log\log s)<s^2
                         \qquad(s\ge6).
\]
The last inequality follows from $\log\log s<\log s$ and $2\log s<s$ in this range. Sieve the integers through $s^2$ and select the first prime $p$ not dividing $W$. Such a prime exists among the first $s$ primes. This is deterministic polynomial work in $L$, and $p$ has polynomial numerical value, not just polynomial bit length. Since $2\mid W$, it is odd; since every positive atom divides some $N_i$, it avoids every $b_j$ as well as $D$. No prime splitting theorem or factorization of $W$ is used.

\paragraph{All unit square roots in one rank-two algebra.}
Choose a quadratic nonresidue $\eta\in\{1,\ldots,p-1\}$ by enumerating the nonzero squares modulo $p$. They form a subgroup of index two, so such an $\eta$ exists. Put
\[
             e=w+2,\qquad Q=p^e,
             \qquad A_e=(\mathbb Z/Q\mathbb Z)[T]/(T^2-\eta).
\]
This algebra is free with the two coordinates $1,T$. For every radicand $b_j$, including a possible $-1$, select $c\in\{1,\eta\}$ and a nonzero residue $r_1$ with $cr_1^2\equiv b_j\pmod p$. This can be done by at most $2p$ trials because the square subgroup and its nonsquare coset exhaust the units.

Use the simple-root Hensel lift from Conrad Theorem2.1, with its finite-precision interpretation in Remark2.3. Explicitly, if $cr_h^2\equiv b_j\pmod{p^h}$, choose
\[
 t_h\equiv
 \frac{b_j-cr_h^2}{p^h}(2cr_h)^{-1}\pmod p,
 \qquad r_{h+1}=r_h+p^h t_h.
\]
The displayed division by $p^h$ is exact in the integers, and $2cr_h$ is a unit modulo $p$. This gives $cr_{h+1}^2\equiv b_j\pmod{p^{h+1}}$. After $e-1$ steps set $\rho_j=r_e$ if $c=1$, and $\rho_j=r_eT$ otherwise. Then $\rho_j^2=b_j$ in $A_e$. All these operations involve integers with $O(L+e\log p)$ bits and a polynomial number of trials and lifts. This finite algorithm is the only Hensel computation used.

Because $p\nmid D$, the relations and denominator inverses define a ring homomorphism
\begin{equation}\label{eq:C043-local-map}
             \mathcal O[D^{-1}]\longrightarrow A_e,
                         \qquad u_j\longmapsto\rho_j.
\end{equation}
This is a map from the localized order, not a characteristic-changing map from its number field. In particular no denominator divisible by $p$ is inverted in its source.

\paragraph{Evaluation, recovery and cost.}
Represent every gate value by a pair $(a,b)$ modulo $Q$. Multiplication is
\[
          (a+bT)(c+dT)=(ac+\eta bd)+(ad+bc)T.
\]
Evaluate the preprocessed constant expressions and then every circuit gate using these formulas and the supplied Boolean input. There are polynomially many gates in $L+n$, and all stored integers have $O((w+2)\log p)$ bits. No intermediate value over characteristic zero is computed. Arbitrarily many repeated-squaring gates cause no additional height or degree parameter beyond their encoded gate count.

By~\eqref{eq:C043-scalar-identity} and the homomorphism property, the final pair is exactly $(g(x)\bmod Q,0)$. Since $p\ge3$,
\[
                          Q=p^{w+2}>2^{w+1},
\]
so there is a unique integer in $[-2^w,2^w]$ with that residue. Its balanced representative recovers $g(x)$. Every preprocessing, prime search, root lift, denominator inverse and gate evaluation uses polynomially many bit operations in $L+n+w$. This proves the single-instance assertion and the uniform and advice-family consequences.

Apply this simulation to the bounded counting outputs in NN-C015 to obtain the stated $\#\mathrm P$ consequences under their explicit circuit-presentation hypothesis. This does not assert that arbitrary complex circuits admit these radical descriptions. Internal division by a circuit-computed algebraic value is outside the theorem: such a value need not map to a unit modulo the selected prime, and the proof provides no avoidance certificate for it. The algorithm also does not verify its integer-output promise or compare nonrational square-root sums.

\paragraph{Fixed odd-prime extension.}
We give the extensions of all presentation and local-algebra steps;
in particular, we do not independently select residue roots of
possibly dependent input radicals.

\paragraph{An injective global presentation.}
Delete zero radicals. For a nonzero radicand $A_i/B_i$, with
$B_i>0$, set $N_i=|A_i|B_i^{\ell-1}$. Its chosen root is
\[
       \zeta_\ell^{j_i}\operatorname{sgn}(A_i)
                         N_i^{1/\ell}/B_i,
\]
where the positive real root is used for $N_i$.
Because $\ell$ is fixed, the total binary size of these integers
is $O_\ell(L)$.
Apply the gcd refinement already proved in NN-C043 to the $N_i$.
Its decreasing product of distinct active atoms gives
$O_\ell(L)$ refinements and polynomial-size exponent vectors.
For each resulting pairwise coprime atom, test by integer root
extraction whether it is an $\ell$th power. Absorb those that
are, and reduce all other exponents modulo $\ell$.
This produces exact identities
\[
 N_i=c_i^\ell\prod_{a=1}^r b_a^{e_{ia}},\qquad
 0\le e_{ia}<\ell,
\]
with positive pairwise coprime $b_a$ that are not $\ell$th powers.
All displayed integer sizes and expression-construction costs
are polynomial in $L$. Integer root testing can be done by
binary search, using fixed-degree integer powers.

The classes of the $b_a$ are independent in
$\mathbb Q^\times/\mathbb Q^{\times\ell}$.
Indeed, each $b_a$ has a prime valuation not divisible by $\ell$;
coprimality makes that valuation occur in no other atom.
Multiplication by any nonzero exponent modulo the prime $\ell$
cannot annihilate it. This uses existence of the valuations,
and does not require computing prime factors.

Put $K=\mathbb Q(\zeta_\ell)$.
Cyclotomic irreducibility gives $[K:\mathbb Q]=\ell-1$.
The same classes remain independent over $K$:
if $q=\prod_a b_a^{h_a}=\alpha^\ell$ for $\alpha\in K^\times$,
then taking the field norm gives
\[
             q^{\ell-1}=\operatorname{Nm}_{K/\mathbb Q}
                                      (\alpha)^\ell.
\]
Since $\gcd(\ell,\ell-1)=1$, every rational prime valuation
of $q$ is divisible by $\ell$. Thus $q$ is a rational $\ell$th
power and every $h_a$ is zero modulo $\ell$.
Kummer theory therefore gives
\[
 [K(b_1^{1/\ell},\ldots,b_r^{1/\ell}):K]=\ell^r.
\]
The exact classical locators are Milne, \emph{Fields and Galois
Theory}, version 5.10 (September 2022), Lemma 5.9, printed p.~65,
Theorem 5.10, p.~66, and Theorem 5.30 with its degree-index
formula and complete proof, pp.~75--76.

Let $D$ be the product of the positive denominators of the rational
leaves and the $B_i$. Define the localized monic algebra
\[
 \mathcal O_\ell=
 \mathbb Z[D^{-1}][Z,U_1,\ldots,U_r]
  /(\Phi_\ell(Z),U_1^\ell-b_1,\ldots,U_r^\ell-b_r).
\]
It is free of rank $(\ell-1)\ell^r$ over $\mathbb Z[D^{-1}]$.
Its characteristic-zero interpretation in the chosen complex field
is surjective and has equal source and target dimensions, by the
degree computation just made. It is consequently injective;
freeness also gives injectivity before tensoring with $\mathbb Q$.
No exponential basis is constructed by the algorithm.
For a Boolean input $x$, the bounded integer-output promise implies
the identity
\[
                 C(x)=g(x)\cdot1\quad\hbox{in }\mathcal O_\ell.
                                                               \tag{C043.P1}
\]

\paragraph{One fixed-degree residue field contains all roots.}
Set
\[
        d=\ell(\ell-1),\qquad
        W=2\ell(\ell-1)D\prod_i N_i.
\]
Use the small-prime construction in NN-C043 on this $W$.
It deterministically finds a prime $p\nmid W$ of polynomial
numerical size in $L$. In particular $p\ne\ell$ and every
remaining radicand and denominator is a unit modulo $p$.

The cyclic group $\mathbb F_{p^d}^\times$ contains the subgroup
$\mathbb F_p^\times$ of order $p-1$. Moreover,
\[
                    \ell\mid\frac{p^d-1}{p-1}.                 \tag{C043.P2}
\]
If $p\equiv1\pmod\ell$, the quotient is a sum of $d$ terms,
each congruent to one, and $\ell\mid d$.
Otherwise $p^{\ell-1}\equiv1\pmod\ell$ and $p-1$ is
invertible modulo $\ell$, which proves the same divisibility.
Choose a generator of the cyclic group for this argument only.
Its power $(p^d-1)/(p-1)$ generates $\mathbb F_p^\times$;
(C043.P2) therefore shows that every element of
$\mathbb F_p^\times$ is an $\ell$th power in
$\mathbb F_{p^d}^\times$. This field also contains a primitive
$\ell$th root of unity.

The residue field is constructible without randomness or a
splitting-prime search. Enumerate monic degree-$d$ polynomials
over $\mathbb F_p$, testing each for a monic factor of degree
between one and $\lfloor d/2\rfloor$ by exhaustive trial.
An irreducible polynomial exists, by existence and simplicity
of finite fields. The exact references are Milne, loc.~cit.,
Propositions 4.19 and 4.20 with their proofs, printed p.~53;
the proof of 4.19 also states multiplicative cyclicity.
The enumeration uses $p^{O(d)}$ operations, which is polynomial
in $L$ because $d$ is fixed and $p$ has polynomial numerical size.
If $\bar f$ is the polynomial found, construct
$\mathbb F_p[T]/(\bar f)$ with its $d$ explicit coordinates.
Enumerate its $p^d$ elements to find a root of each
$V^\ell-b_a$ and a nonidentity root of $V^\ell-1$.
The latter has order $\ell$ since $\ell$ is prime.
These searches are deterministic polynomial work for fixed $\ell$.

\paragraph{Explicit lifting in constant rank.}
Lift the coefficients of $\bar f$ to integers in
$\{0,\ldots,p-1\}$, obtaining a monic polynomial $f$.
For $e=w+2$ put
\[
       A_e=(\mathbb Z/p^e\mathbb Z)[T]/(f).
\]
It has the explicit free basis $1,T,\ldots,T^{d-1}$.
An element is a unit if its reduction modulo $p$ is nonzero:
lift an inverse modulo $p$ and invert the remaining factor
$1+pc$ by the finite geometric series of length $e$.
The irreducible polynomial $\bar f$ over the finite field is
separable, so $f'$ is a unit in $A_e$ by this criterion.
Thus $A_e$ is finite etale as well as free of the fixed rank $d$.

Here is the root-lifting algorithm, including its meaning in
this finite algebra. Let $F(V)$ be $V^\ell-b_a$ or
$V^\ell-1$, and start with its selected nonzero residue root.
If $a_h$ is a root modulo $p^h$, represent it by $d$ integer
coordinates. Compute $F(a_h)$ in $\mathbb Z[T]/(f)$.
Every coordinate is divisible by $p^h$. In the residue field set
\[
 t_h=-\bigl(F(a_h)/p^h\bigr)F'(a_h)^{-1},\qquad
                     a_{h+1}=a_h+p^h t_h.
\]
The quotient is coordinatewise exact before reducing modulo $p$;
the derivative is a unit because $p\ne\ell$ and $a_h$ is nonzero
modulo $p$. Expansion of the fixed-degree polynomial shows
$F(a_{h+1})=0\pmod{p^{h+1}}$: all terms beyond the linear term
are divisible by $p^{2h}$, and $2h\ge h+1$ for $h\ge1$.
This proves each lifting step directly. It is the simple-root
digit method of Conrad, \emph{Hensel's Lemma}, Theorem 2.1 and
Remark 2.3, applied in the explicitly presented unramified algebra.

Write $u_a$ for the lifted root of $V^\ell-b_a$, and $z$ for
the lift of the chosen nonidentity root of $V^\ell-1$.
Since $z-1$ is a unit and
$(z-1)\Phi_\ell(z)=z^\ell-1=0$, one has $\Phi_\ell(z)=0$.
As $D$ is invertible modulo $p^e$, these assignments define
an actual ring map
\[
               \mathcal O_\ell\longrightarrow A_e,
                  \qquad Z\longmapsto z,\quad U_a\longmapsto u_a.
                                                               \tag{C043.P3}
\]
Its source is the localized integral order. No map from a
characteristic-zero field to a positive-characteristic ring
is being asserted.

Every root step involves $d$ coordinates, fixed-degree powers,
and integers of $O_\ell(L+e\log p)$ bits. The number of roots
and the number of lifts are polynomial in $L+w$.
Evaluation of each original circuit gate then uses only
addition and multiplication in the fixed-rank algebra $A_e$,
with all coordinates reduced modulo $p^e$ after every gate.
This costs $\operatorname{poly}_\ell(L+n+w)$ bit operations,
regardless of the sizes of characteristic-zero intermediate
values or the circuit's formal degree.

By (C043.P1)--(C043.P3), the final coordinate vector is
$(g(x)\bmod p^e,0,\ldots,0)$. Since $p\ge3$ and $e=w+2$,
$p^e>2^{w+1}$. Its first coordinate has a unique lift to the
interval $[-2^w,2^w]$, which is the required output.
This also proves the asserted uniform and nonuniform family
consequences by the same description-length argument as NN-C043.
\end{proof}
\paragraph{Status and complexity impact.}
Complete fixed-prime radical deduction proved; quadratic self-check and root-authored odd-prime self-check passed. Complete deduction-proof review by root (quadratic base); complexity agent (root-developed fixed odd-prime extension): passed; imported input theorems are not claimed independently proved. This removes the polynomial global-rank restriction of NN-C020 for explicitly presented rational radicals of any one fixed prime degree. Quadratic generators allow global degree $2^r$ and local rank two; generators of a fixed odd prime degree $\ell$ allow global degree $(\ell-1)\ell^r$ and local rank $\ell(\ell-1)$. Consistent coprime preprocessing is deterministic and does not require integer factorization. The bounded integer-output promise, explicit binary radicands and division-free arithmetic remain essential hypotheses. General complex constants, growing or mixed unrelated root degrees, unrestricted nested radicals, and circuit-computed divisions remain outside this result. The theorem is unconditional and uses neither E nor A nor M; the new point is the bounded-output simulation and its exact effective presentation scope, not the classical coprime-refinement or Kummer theorems.

\paragraph{Geometric and representation dictionary.}
Let $b_1,\ldots,b_r$ be the independent radicands constructed in the proof and let
\[
 R=\mathbb Z[(2D\prod_j b_j)^{-1}],\qquad
 X=\operatorname{Spec}
     R[u_1,\ldots,u_r]/(u_j^2-b_j:1\le j\le r).
\]
Then $X\to\operatorname{Spec}R$ is a finite etale Kummer cover of rank $2^r$, with the independent sign action of $(\mathbb Z/2\mathbb Z)^r$. Each quadratic relation has invertible derivative $2u_j$ because $2$ and $b_j$ are units. The exact standard-etale criterion is Stacks Definition10.144.1 and Lemma10.144.2(1)--(2), tags00UB and00UC; the generic field and character group are given by Milne Theorem5.30 and Example5.31(b). The monomial basis is the corresponding sign-character decomposition over the generic fiber.

There is a deterministically constructible prime $p=\operatorname{poly}(L)$, and for every requested precision $e$ a finite etale rank-two algebra
\[
 A_e=(\mathbb Z/p^e\mathbb Z)[T]/(T^2-\eta)
\]
with an explicit $\mathbb Z/p^e\mathbb Z$-morphism
\[
       \operatorname{Spec}A_e
              \longrightarrow X\times_R\operatorname{Spec}(\mathbb Z/p^e\mathbb Z).
\]
Thus a local point of residue degree at most two suffices even though the global cover has exponential degree. No degree-one point or completely split prime is required. The Hensel construction explicitly lifts that point through every indicated nilpotent thickening.

For any encoded arithmetic section whose values on the Boolean input components descend to integers $g(x)$ with $|g(x)|\le2^w$, pulling back along this point gives $g(x)\cdot1$ in $A_{w+2}$. Its normalized finite-algebra trace is
\[
              \frac12\operatorname{Tr}_{A_{w+2}/(\mathbb Z/p^{w+2}\mathbb Z)}
                           (g(x)\cdot1)
                    =g(x)\pmod{p^{w+2}}.
\]
The inverse of two exists, and the trace is literally the trace of the two-by-two multiplication matrix. EGA IV4,18.2.1--18.2.2,p111, with base change in equation18.2.1.2, supplies the exact classical trace interface. The magnitude bound recovers the unique integer lift in polynomial time. This is an effective local-specialization and descent theorem for the given Kummer presentations, not a purity-to-multiplicity assertion or a new GCT obstruction.

Invert $2\ell(\ell-1)D\prod_a b_a$ in $\mathbb Z$, giving $R$.
The corresponding scheme $X=\operatorname{Spec}\mathcal O_\ell
\times_{\operatorname{Spec}\mathbb Z[D^{-1}]}\operatorname{Spec}R$
is finite etale of rank $(\ell-1)\ell^r$ over $R$.
Indeed the derivatives of the radical relations are units.
For the cyclotomic relation, $Z-1$ is a unit because
$\Phi_\ell(1)=\ell$ is a unit, and differentiation of
$(Z-1)\Phi_\ell(Z)=Z^\ell-1$ gives the same derivative criterion.
After the cyclotomic base change it is a $\mu_\ell^r$-torsor.
This is the standard-etale criterion in Stacks
Definition 10.144.1 and Lemma 10.144.2, tags 00UB and 00UC.

The construction gives a point over the finite etale rank-$d$
algebra $A_e$ on every indicated reduction of $X$ modulo $p^e$.
Its lifted scalar section has normalized trace
\[
       d^{-1}\operatorname{Tr}_{A_e/(\mathbb Z/p^e\mathbb Z)}
                          (g(x)\cdot1)=g(x)\pmod{p^e};
\]
$d$ is invertible by the chosen avoidance factor in $W$.
The finite-algebra trace and base-change dictionary is
EGA IV$_4$, 18.2.1--18.2.2, printed p.~111, especially
equation 18.2.1.2, as already checked in NN-C020 and NN-C043.
The global cover has exponential rank, while this explicit
local coefficient algebra has the fixed rank $\ell(\ell-1)$.

This extends a proved coefficient-presentation class relevant
to the bounded counting-output transfer. It does not construct
such presentations for arbitrary complex arithmetic circuits.
The prime $\ell$ is fixed throughout; neither a growing prime
nor mixed unrelated radical degrees nor unrestricted nested
radicals is covered by the cost and injectivity proof above.
The integer-output promise and the exclusion of computed
division remain exactly as in NN-C043. No new positivity or
occurrence obstruction follows from this local trace identity.



\section{NN-C044 --- Polynomial advice from logarithmic individual algebraic degrees and bounded integer outputs}
\begin{proposition}[NN-C044]
Let $C_n$ be a family of division-free arithmetic circuits over $\mathbb C$, with Boolean inputs, operations $+,-,\times$, and polynomially many gates. Every constant leaf is algebraic over $\mathbb Q$. Suppose that its primitive irreducible polynomial over $\mathbb Z$ has degree at most $d(n)$, and that the sum $L(n)$ of the dense binary lengths of these polynomials is polynomially bounded. Put
\[
                    D(n)=\operatorname{lcm}(1,2,\ldots,d(n)).
\]
Assume that $D(n)$ is polynomially bounded. If every Boolean output is a rational integer $g_n(x)$ satisfying $|g_n(x)|\le2^{w(n)}$ for a polynomial $w$, then $g\in\mathrm{FP}/\mathrm{poly}$. In particular, the degree condition holds whenever $d(n)=O(\log(n+2))$.

More quantitatively, a circuit with $s$ gates, $n$ input bits, constant-polynomial description length $L$, output bound $2^w$, and the above integer $D$ has advice of length and evaluation time $\operatorname{poly}(s+n+L+w+D)$. Neither its formal degree nor the degree of the field generated by all constants occurs in this bound. The advice contains a finite residue-field presentation and a coherent tuple of local constant values; they are not asserted uniformly constructible from the separate minimal polynomials.

For Boolean outputs this coefficient model has exactly the Boolean power of $\mathrm P/\mathrm{poly}$. If every $\#\mathrm P$ counting family has circuits with these bounds, then $\#\mathrm P\subseteq\mathrm{FP}/\mathrm{poly}$. These implications are unconditional and require no Riemann-hypothesis input.
\end{proposition}
\paragraph{Exact inputs.}
NN-C020; NN-C043. Background comparison only: explicit polynomial-rank bounded-output simulation and uniform fixed-prime radical simulation; the present theorem replaces global presentation by polynomial advice.

J.S. Milne, Fields and Galois Theory, version 5.10 (September 2022). Proposition 2.12 and full proof, printed pp29--30; Corollary 2.13,p30; finite-field Propositions 4.19--4.21,p53, Corollary 4.22 and Proposition 4.23,p54; resultant definition and Proposition 4.35,p58.

Rosser--Schoenfeld, Approximate formulas for some functions of prime numbers, Illinois J. Math. 6 (1962),64--94. Printed p69, Corollary to Theorem 3, equation (3.13), nth-prime bound; printed p71, Theorem 12, equation (3.35), psi(x)<1.03883x for x>0; both original page images read.

Keith Conrad, Hensel's Lemma. Theorem 2.1 and complete digit-lifting proof, pp1--3, with finite-precision Remark 2.3,p3; the coordinate version used here is proved explicitly.

Stacks Project, Local structure of etale ring maps. Definition 10.144.1,tag00UB; Lemma 10.144.2(1)--(2),tag00UC, standard etale presentation and base change.

D. Ierardi and D. Kozen, Parallel Resultant Computation, author chapter manuscript dated June 10,1994. Section 20.2.1, printed pp11--12: equation(20.1), Sylvester determinant definition, Theorem20.7 and complete proof, followed by resultant-polynomial specialization; original page images read.

EGA IV, fourth part, Publications Mathematiques IHES32 (1967). Printed p.111,18.2.1--18.2.2; base-change equation18.2.1.2; primary PDF page110 (zero-based109).

\begin{proof}
Fix one input length and suppress the subscript $n$. The circuit structure is ordinary nonuniform advice. Let its algebraic constants be $\alpha_1,\ldots,\alpha_r$. Constants may repeat; including repeated polynomial descriptions only increases $L$. If there are no constant leaves, use $d=D=1$ and omit the corresponding products. The degree and height assumptions concern each literal algebraic leaf; arithmetic subcircuits producing new constants are counted as gates and need no separate height bound.

\paragraph{Polynomially many bits of bad-prime data.}
Write the primitive irreducible polynomial of $\alpha_i$ as
\[
 f_i(T)=a_iT^{d_i}+\cdots\in\mathbb Z[T],
 \qquad a_i>0,\quad 1\le d_i\le d,
\]
and let $H_i\ge1$ bound the binary lengths of its coefficients. It is separable in characteristic zero. Therefore
\[
              R_i=\operatorname{Res}(f_i,f_i')\ne0.
\]
Resultants are integer polynomials in the coefficients and commute with reduction. For an explicit height bound, their Sylvester determinant has order $2d_i-1$ and entries bounded in absolute value by $d_i2^{H_i}$. Thus
\[
 |R_i|\le(2d_i-1)!\,(d_i2^{H_i})^{2d_i-1},
 \qquad
 \log_2|R_i|=O\bigl(d_i(H_i+\log(d_i+1))\bigr).
\]
This determinant and its specialization property are the classical interface in Ierardi--Kozen, Section20.2.1, equation(20.1), Theorem20.7 and the following resultant-polynomial paragraph, printed pp11--12; the root formula is also Milne Proposition4.35. The determinant bound is a height estimate, not an algorithm that enumerates its permutations. Exact integer determinant computation, or the usual resultant algorithm, uses polynomial time in these explicit polynomial descriptions; uniform construction of this data is not needed for the advice conclusion.

Set
\[
                   W=2D\prod_{i=1}^r a_i|R_i|.
\]
Its bit length $H$ is polynomial in $L+\log D$. Put $t=\max(6,H+1)$. The integer $W$ has at most $H-1$ distinct prime divisors. Rosser--Schoenfeld, p69, equation (3.13), gives
\[
                 p_t<t(\log t+\log\log t)<t^2.
\]
Consequently there is a prime $p<t^2$ not dividing $W$. It can even be selected by sieving to $t^2$ and testing divisibility; no splitting-prime theorem is used. In particular $p$ has polynomial numerical value, is odd, and divides none of $D,a_i,R_i$.

Normalize $P_i=f_i/a_i\in\mathbb Z_p[T]$. Its reduction is monic of degree $d_i$ and has no multiple root, since the leading coefficient and resultant are nonzero modulo $p$. Each irreducible factor of this reduction has degree at most $d_i\le d$, hence degree dividing $D$. The finite-field results in Milne, Corollary 4.22 and Proposition 4.23, imply that every $\overline P_i$ splits into distinct roots in $\mathbb F_{p^D}$.

\paragraph{A small local field containing all conjugates.}
There exists a monic irreducible $\overline G\in\mathbb F_p[T]$ of degree $D$, by the finite-field existence and simple-extension theorems. Choose one such polynomial as advice and lift its coefficients to $\{0,\ldots,p-1\}$, obtaining $G\in\mathbb Z[T]$. Its presentation uses $O(D\log p)$ bits. This is an existence and advice step; enumerating $p^D$ candidate elements or polynomials is not a polynomial-time construction when $D$ grows, and no such enumeration is used here.

Define
\[
 A_\infty=\mathbb Z_p[T]/(G),\qquad
 K_p=A_\infty[1/p],\qquad
 A_e=(\mathbb Z/p^e\mathbb Z)[T]/(G).
\]
The monic relation gives the free basis $1,T,\ldots,T^{D-1}$, so $A_\infty$ is complete and separated for its $p$-adic topology and has residue field $\mathbb F_{p^D}$. An element with nonzero residue is a unit: lift its residue inverse, then invert the resulting $1+pz$ by the convergent geometric series. Every nonzero element of $A_\infty$ is $p^k$ times a unit, by taking the smallest $p$-adic valuation of its finitely many coordinates. Therefore $A_\infty$ is a domain and $K_p$ is a field of degree $D$ over $\mathbb Q_p$. This establishes the needed field property without presuming that a chosen integral lift is irreducible over $\mathbb Q$. Moreover, $\overline G$ is separable because finite fields are perfect, so $G'$ is a unit in $A_e$; the standard-etale criterion makes every $A_e$ finite etale over $\mathbb Z/p^e\mathbb Z$.

Every simple root of $\overline P_i$ lifts uniquely to a root of $P_i$ in $A_\infty$. Here is the relevant coordinate form of the classical Hensel argument. Given $u_h$ with $P_i(u_h)\equiv0\pmod{p^h}$, choose the unique residue $v_h\in\mathbb F_{p^D}$ satisfying
\[
 v_hP_i'(u_h)\equiv-P_i(u_h)/p^h\pmod p,
 \qquad u_{h+1}=u_h+p^hv_h.
\]
The derivative is a unit, and the polynomial Taylor identity shows $P_i(u_{h+1})\equiv0\pmod{p^{h+1}}$. All coefficients of $P_i$ are $p$-integral because $p\nmid a_i$. Passing to the coordinatewise limit gives a root. The same correction formula gives uniqueness among lifts of the chosen residue. The distinct residue roots produce $d_i$ distinct roots, so $f_i$ splits completely in $K_p$. This is the digit-lifting interface of Conrad Theorem 2.1, with the displayed argument supplying its finite free algebra version.

Let $K=\mathbb Q(\alpha_1,\ldots,\alpha_r)$, interpreted using the circuit's actual complex constants. It may have exponential degree. It is generated by roots of the product $\prod_i f_i$, which splits in $K_p$. Milne Proposition 2.12 therefore supplies a $\mathbb Q$-embedding
\[
                              \sigma:K\hookrightarrow K_p.
\]
Every $\sigma(\alpha_i)$ lies in $A_\infty$, because all roots of $f_i$ there are the integral Hensel lifts just constructed. The tuple is coherent: all algebraic relations among the actual constants are preserved simultaneously. Independently choosing one modular root of each $f_i$ need not have this property. The embedding is used to prove existence of correct advice; the proof does not compute it or recognize its image among all root tuples.

\paragraph{Polynomial advice, gate simulation, and exact recovery.}
Set $e=w+2$. In addition to the circuit structure, use as advice $p,G$ and the $D$ coordinates of
\[
                   \beta_i=\sigma(\alpha_i)\bmod p^e\in A_e
                        \quad(1\le i\le r).
\]
The total additional advice length is
\[
                       O\bigl((r+1)D(w+2)\log p\bigr).
\]
This bound is polynomial in the theorem's parameters. Neither the minimal polynomials nor a basis of $K$ need be evaluated or expanded by the Boolean simulator.

Evaluate every arithmetic gate in the coordinate algebra $A_e$. Addition and subtraction are coordinatewise. Multiplication is polynomial multiplication followed by reduction by the monic degree-$D$ polynomial $G$, so a direct implementation uses polynomially many operations on integers of $O(e\log p+\log D)$ bits per gate. There is no division by a circuit-computed value. This computes the reduction of $\sigma(C(x))$ because the constants are a coherent tuple and every gate operation commutes with the homomorphism on the integral subring generated by them. The map to $A_e$ is from that $p$-integral subring, not from the characteristic-zero field as a whole.

By the promised scalar output, $C(x)=g(x)\in\mathbb Z$, so the output coordinates are exactly
\[
                         (g(x)\bmod p^e,0,\ldots,0).
\]
Since $p\ge3$ and $e=w+2$, one has $p^e>2^{w+1}$. The balanced representative recovers the unique integer in $[-2^w,2^w]$ with this residue. All stored residues and coordinate operations have the stated polynomial cost, regardless of the size or degree of the characteristic-zero intermediate values. This proves the quantitative advice statement and its family consequence.

\paragraph{Logarithmic individual degrees and the remaining uniform gate.}
The identity
\[
 \log\operatorname{lcm}(1,\ldots,d)
   =\sum_{p^a\le d}\log p=\psi(d)
\]
follows by taking the largest prime power dividing each integer in the least common multiple. Rosser--Schoenfeld Theorem 12, p71, equation (3.35), gives $\psi(d)<1.03883d$. Therefore $d=O(\log(n+2))$ implies $D=(n+2)^{O(1)}$. The hypothesis allows a polynomial number of unrelated algebraic constants and imposes no bound on their joint field degree.

Conversely, Boolean circuits are simulated arithmetically by $\neg a=1-a$, $a\wedge b=ab$, and $a\vee b=a+b-ab$. Their constants have degree one and bounded height. This proves the asserted equality with $\mathrm P/\mathrm{poly}$ for Boolean outputs. Applying the integer-output simulation separately to each promised counting family proves the $\#\mathrm P$ conditional consequence.

No polynomial-time algorithm for selecting a coherent tuple follows from the existence argument. Even a polynomial-rank residue-field presentation was supplied as advice. NN-C043 proves uniform special cases by explicit radical preprocessing; that extra structure is absent here. Arbitrary transcendental constants, algebraic leaves of unrestricted individual degree or height, and internal divisions by circuit-computed values remain outside this theorem. In particular this is not an unconditional conversion of arbitrary complex Valiant circuits into Boolean circuits.
\end{proof}
\paragraph{Status and complexity impact.}
Complete deduction written and self-reviewed; root and analytic-agent complete proof readings passed. Complete deduction-proof review by root: passed; imported input theorems are not claimed independently proved. The global polynomial-rank and explicit compositum-presentation conditions in NN-C020 can be removed for nonuniform bounded-integer-output simulation when every algebraic constant has polynomial height and individual degree $O(\log n)$. The small-prime argument is unconditional; the coherence problem moves into polynomial advice instead of being silently treated as a uniform root-selection algorithm. This is a substantive coefficient-model strengthening of the Cook--Valiant transfer interface, while the unrestricted complex-constant theorem and its hypotheses remain separate. The result uses neither E nor A nor M.

\paragraph{Geometric and representation dictionary.}
Let $f_i,a_i,R_i,D,W$ be as in the proof, put $R=\mathbb Z[1/W]$, and define the finite cover
\[
 Y=\operatorname{Spec}
       R[T_1,\ldots,T_r]/\bigl(f_i(T_i)/a_i:1\le i\le r\bigr).
\]
It is finite etale of rank $\prod_i d_i$: the equations are monic, and inversion of the resultants makes each derivative a unit in its corresponding factor algebra. This is precisely the standard-etale criterion of Stacks Definition 10.144.1 and Lemma 10.144.2, tags00UB and00UC. The chosen actual tuple $(\alpha_i)$ selects the generic field component $K$ of $Y_{\mathbb Q}$; no algorithm for computing this component or its degree is claimed.

There is a prime $p$ of polynomial numerical size, with $p\nmid W$, and a degree-$D$ finite etale algebra
\[
       A_e=(\mathbb Z/p^e\mathbb Z)[T]/(G)
\]
for every $e\ge1$, together with compatible morphisms
\[
       \operatorname{Spec}A_e\longrightarrow
           Y\times_R\operatorname{Spec}(\mathbb Z/p^e\mathbb Z).
\]
These points preserve every relation over $R$ among the coordinates of the selected generic component: they arise from a single embedding $K\hookrightarrow A_\infty[1/p]$ whose integral constant values lie in $A_\infty$. The local rank is $D$, even when the generic component has much larger degree. The existence uses complete splitting of the residue polynomials in $\mathbb F_{p^D}$ and Hensel lifting; it does not assert a degree-one point, a completely split prime over $\mathbb Q$, or a canonical selection of these local points.

For an arithmetic section whose restriction to every Boolean input component in the selected generic field is an integer $g(x)$ with $|g(x)|\le2^w$, pullback to $A_{w+2}$ is the scalar $g(x)1$. Because $p\nmid D$, normalized finite-algebra trace gives
\[
 D^{-1}\operatorname{Tr}_{A_{w+2}/(\mathbb Z/p^{w+2}\mathbb Z)}(g(x)1)
                  =g(x)\pmod{p^{w+2}}.
\]
The trace is the trace of the multiplication matrix in the displayed free basis. Its exact classical construction and base change are EGA IV4,18.2.1--18.2.2,p111, equation18.2.1.2. The output bound recovers $g(x)$ uniquely. If $D$ and the remaining descriptions are polynomially bounded, the field presentation and compatible local point are polynomial advice and each pullback is polynomial-time computable. This is a finite-etale specialization theorem with an explicit nonuniform computational conclusion; it supplies no purity-to-multiplicity implication and no new representation-positivity criterion.



\section{NN-G001: the arithmetic zeta of the actual $D_5$ plethysm algebra}
\paragraph{Statement.}
Let
\[
 X=\operatorname{Spec}\mathbb Z[U,V,Z]/(Z^2-UV^2-4U^4),
 \qquad \deg(U,V,Z)=(2,3,4).
\]
For every finite field $\mathbb F_q$, including characteristic two,
\[
 |X(\mathbb F_q)|=q^2,\qquad
 Z(X_{\mathbb F_q},T)=(1-q^2T)^{-1},\qquad
 \zeta_X(s)=\zeta(s-2).
 \tag{G001.1}
\]
The grading nevertheless has Hilbert series
\[
 H_X(t)=\frac{1+t^4}{(1-t^2)(1-t^3)},\qquad
 \dim X_s=\lfloor s/3\rfloor+\mathbf1_{2\mid s},
 \tag{G001.2}
\]
where $X_s$ in this formula means the degree-$s$ coordinate space. Its
weight support is $S=\{0,2,3,4,\ldots\}$, its saturation in its group
$\mathbb Z$ is $\mathbb N$, and its conductor is $2$. Thus the ungraded
arithmetic zeta agrees with that of $\mathbb A^2_{\mathbb Z}$ while the
actual plethysm grading has a hole and has no homogeneous rational
Ehrhart model. Under A the nontrivial zeros in (G001.1) lie on
$\operatorname{Re}s=5/2$; this supplies no information recovering that
hole or an Ehrhart model.

\paragraph{Inputs and locators.}
ALPOGE\_GCT, PLE-F5 and PLE-F8, respectively
\path{work/plethysm.tex}, the subsection beginning ``The actual
plethysm ray'', and \path{work/positivity_actual_plethysm.tex}:
the complex graded algebra is the highest-weight multiplicity algebra
for $S^{(7s,5s)}\mathbb C^2$ in
$\operatorname{Sym}^{4s}(\operatorname{Sym}^3\mathbb C^2)$.
Kahle--Micha\l{}ek, \emph{Obstructions to combinatorial formulas for
plethysm}, Theorem 5, supplies the non-Ehrhart assertion. A is used
only for the final zero-location assertion.

\paragraph{Proof.}
Suppose first that $q$ is odd and let $\chi$ be its quadratic character,
extended by $\chi(0)=0$. For $U=0$ there are $q$ points: $V$ is arbitrary
and $Z=0$. For $U=u\ne0$ the number is
\[
 q+\sum_{v\in\mathbb F_q}\chi(uv^2+4u^4)
 =q+\chi(u)\sum_v\chi(v^2+4u^3)=q-\chi(u).
\]
Here, for every $a\ne0$, the equation $y^2-x^2=a$ has $q-1$
solutions: its two factors $y-x,y+x$ are a freely chosen nonzero
element and its prescribed reciprocal multiple. It also has
$q+\sum_x\chi(x^2+a)$ solutions. Therefore that character sum is
$-1$. Summing $q-\chi(u)$ over $u\ne0$ and using
$\sum_{u\ne0}\chi(u)=0$ proves $q^2$.
If $q$ is even, squaring is a bijection of $\mathbb F_q$, so each
pair $(U,V)$ admits exactly one $Z$. This proves the count for all
finite fields, hence for every extension of each residue field.
The defining exponential formula for local zeta functions gives
$\exp(\sum_{r\ge1}q^{2r}T^r/r)=(1-q^2T)^{-1}$, and the product
over rational primes gives (G001.1).
The relation is monic in $Z$, so $1,Z$ is a free basis over
$\mathbb Z[U,V]$. This proves flatness over $\mathbb Z$ and the
displayed Hilbert series. Its coefficient formula is the supplied
PLE-F5 calculation; directly it counts $U^iV^jZ^\epsilon$ with
$\epsilon=0,1$. Degrees $2,3$ generate precisely $S$; $1$ is absent,
$2\cdot1\in S$, and every integer at least two is present. The
remaining conclusions now follow from the stated inputs.\hfill$\square$

\paragraph{Algebraic geometry and Grothendieck dictionary.}
This is an actual flat affine scheme over $\mathbb Z$, together with a
$\mathbb G_m$ action. The action, equivalently its grading, is the
extra datum in SGA 3, I, Proposition 4.7.3 and Corollary 4.7.3.1.
Forgetting it loses the Hilbert function. SGA 5, III, 4.7 is the
Lefschetz--Verdier trace framework for the local zeta; equality of
zeta functions alone is not an identification of cohomological
degrees, motives, or canonical bases. EGA II, Proposition 2.4.7(i),
additionally shows that taking a Veronese subring preserves
$\operatorname{Proj}$ as a scheme while changing the recorded grading.
There is no implication from the normality of this complex $D_5$
surface to saturation of the support of an arbitrary torus grading:
the example itself refutes that inference.

\paragraph{Complexity impact.}
An arithmetic realization with the same zeta as affine space cannot
by itself establish PH1, PH2, PH3, or a positivity algorithm. This is
an explicit obstruction to the proposed recovery from the classical
zeta sheet, rather than a complexity-class separation.



\section{NN-G002: the exact output obstruction from the exponential poles}
\paragraph{Statement and proof.}
On the P-F31 cubic family, put $M=2^{n-1}$ and let $f_n$ be its
stretching quasi-polynomial. A full constituent table for $f_n$ has
at least $M$ entries. Every denominator presentation
\[
 H_n(t)=Q(t)\prod_j(1-t^{a_j})^{-b_j},\qquad a_j,b_j>0,
 \tag{G002.1}
\]
has some $a_j$ divisible by $M$; its expanded denominator has degree
at least $M$. Indeed P-F31 gives a genuine pole at every primitive
$M$th root. At one such root a denominator factor in (G002.1) must
vanish, which is exactly $M\mid a_j$. A quasi-polynomial of period
$L$ has a generating function whose poles are $L$th roots of unity:
sum its polynomial constituents over arithmetic progressions and
differentiate geometric series. Thus $M\mid L$. This proves both
output assertions.\hfill$\square$

\paragraph{Scope and dictionary.}
The input is P-F31, \path{work/positivity_actual_cubic_pole.tex},
the theorem ``actual cubic pole'' and its denominator consequence;
P-F56 makes the two-adic valuation of a quasi-period exact. This is
an obstruction to an \emph{expanded} Hilbert-series representation.
The integer $M$ itself has only $n$ binary digits, so this proves no
exponential lower bound for arithmetic circuits, succinct rational
expressions, evaluation, or positivity. In geometric language these
are Hilbert characters of the graded invariant coordinate rings;
the grading/action dictionary is SGA 3, I, 4.7.3.1 and the projective
construction is EGA II, 2.4.1--2.4.7. No statement about the size of
a defining system or a Quot/Hilbert parameter space follows from
the length of this particular output format.



\section{NN-G003: arbitrarily large strict indices with no support holes}
\paragraph{Statement.}
For every $L\ge2$ define
\[
 f_L(n)=L\lfloor n/L\rfloor+1\quad(n\ge0).
\]
Its support is all of $\mathbb N$, while the positivity and strict
saturation indices in GCT VI, Definitions 1.2.2 and 1.2.4, are both
$L-2$. It has a positive witness description of size polynomial in
$\log L+\log(n+1)$, and
\[
 \sum_{n\ge0}f_L(n)t^n=
 \frac{1+(L-1)t^L}{(1-t)(1-t^L)}.
 \tag{G003.1}
\]

\paragraph{Proof.}
The constituent on $n\equiv r\pmod L$ is $X-r+1$.
For $f_L(n+p)$ the constituent on $n\equiv r$ is
$X+p+1-((r+p)\bmod L)$. Its least constant coefficient is
$p-L+2$. All coefficients are nonnegative exactly when $p\ge L-2$.
All these constituents have slope one; each is positive at every
integer $X\ge1$ exactly when its value at one is positive. The least
such value is $p-L+3$, and strict positivity again means
$p\ge L-2$. No constituent is identically zero. The count is always
at least one. It counts a distinguished object and the pairs
$(i,j)$ with $1\le i\le\lfloor n/L\rfloor$ and $1\le j\le L$.
Binary comparisons and division verify this set with polynomial
resources. Finally
$f_L(n)=1+L\sum_{k\ge1}\mathbf1_{n\ge kL}$ gives (G003.1).
\hfill$\square$

\paragraph{What this resolves.}
The inference ``exponential strict saturation index implies an
exponentially distant realized hole'' is false even for compactly
described positive counts. This is a logical countermodel, not a
new plethysm family. P-F57 supplies actual plethysm examples with
large strict indices; its off-residue evaluation must not silently
be replaced by an actual vanishing coefficient. The supplied
positive-start assertion is retained at statement level only.

\paragraph{Geometric translation.}
The series is the Hilbert series of the graded
$k[x,y]$-module
$k[x,y]\oplus k[x,y](-L)^{\oplus(L-1)}$, with
$\deg x=1$, $\deg y=L$. Thus the same failure of inference occurs
for effective coherent graded modules, equivalently equivariant
modules under SGA 3, I, 4.7.3. The associated sheaf construction is
EGA II, Proposition 2.5.2. This uses a module, not a normal domain
or an actual GCT multiplicity algebra; no stronger geometric
realization is asserted.



\section{NN-G004: cancellation and precision on the cubic pole formula}
\paragraph{Statement.}
Fix $0<c<1/4$. At the witness $x_n$ in P-F57, write the specific
pole expansion as $f_n(x_n)=\sum_\rho a_{\rho,n}$, with the dyadic
wave $W_n(x_n)$ among its grouped summands. If $f_n(x_n)>0$, its
relative summation condition number satisfies
\[
 \kappa_n:=\frac{\sum_\rho|a_{\rho,n}|}{f_n(x_n)}
 \ \ge\ 2^{(1/24-c/6)n^4-O(n^3\log n)}.
 \tag{G004.1}
\]
If $f_n(x_n)=0$ and the nonzero wave is retained, the relative
condition number is infinite. The lower bound is for this formula;
it gives a loss of $\Omega(n^4)$ bits of relative accuracy, not an
exponential number of required bits.

\paragraph{Proof.}
P-F57 gives
$2W_n(x_n)\ge2^{n^4/24-2n^3-1}$ and
$f_n(x_n)\le2^{(c/6)n^4+O(n^3\log n)}$.
The sum of the absolute values of the pole summands is at least
the absolute value of their dyadic subsum, hence at least
$|W_n(x_n)|$. Division gives (G004.1).
For a finite sum $F=\sum a_j\ne0$, relative perturbations
$|\Delta a_j|\le\eta|a_j|$ change it by at most
$\eta\sum|a_j|$; equality in the bound is attained by aligning
the perturbations with one complex direction. This proves the
stated relative condition number. A relative error of size
$2^{-b}$ in the terms can therefore cause relative error as large
as $2^{-b}\kappa_n$ in the sum. Taking logarithms proves the bit
statement.\hfill$\square$

\paragraph{Inputs, geometry and complexity.}
The inputs are ALPOGE\_GCT P-F56d and P-F57, with the exact wave
and dimension estimates in the supplied intake part 4 and the
paired-wave/strict-index proof files. A cyclotomic partial-fraction
expansion is a character computation in a localized representation
ring; the tensor and restriction operations themselves appear in
cote 2, scan 17, and symmetric-group operations in cote 133,
scan 28. A virtual-character cancellation is not a proof that
no different effective realization exists. The inequality gives
no Boolean or algebraic circuit lower bound and no obstruction
to exact integer arithmetic or symbolic cancellation. It makes
the ``near-supported zero'' precise only relative to a specified
expansion and norm.



\section{NN-G005: the universal support valuation and its exact linearization}
\paragraph{Statement.}
Let $L$ be a finite distributive lattice with least element $0$,
and let $P$ be its poset of join-irreducibles. The universal abelian
group for normalized inclusion--exclusion valuations is
\[
 V(L):=\mathbb Z[L]/\langle[0],\ [x]+[y]-[x\vee y]-[x\wedge y]\rangle
 \simeq\mathbb Z^P.                                      \tag{G005.1}
\]
Under the identification $L\simeq\mathcal I(P)$ with order ideals,
the universal valuation is $I\mapsto\sum_{p\in I}e_p$.
In contrast, the abelian group completion of the join monoid
$(L,\vee,0)$ is zero. For a Boolean lattice on $r$ atoms,
the kernel of the augmentation $V(L)\to\mathbb Z$ is $A_{r-1}$
with its usual integral coordinates; this is an extra valuation
construction, not the group completion of idempotent addition.

\paragraph{Proof.}
For completeness, every element of a finite lattice is the join
of the join-irreducibles below it, by induction on the interval
below it. In a distributive lattice a join-irreducible $p$ is
join-prime: if $p\le x\vee y$, then
$p=(p\wedge x)\vee(p\wedge y)$, so $p\le x$ or $p\le y$.
Consequently $x\mapsto\{p:p\le x\}$ identifies $L$ with the
order ideals in $P$; the inverse takes joins. Characteristic
functions of ideals satisfy inclusion--exclusion, giving a map
$V(L)\to\mathbb Z^P$. Define the inverse on a basis vector by
\[
 e_p\longmapsto[\mathord\downarrow p]
                    -[\mathord\downarrow p\setminus\{p\}].
\]
If $p$ is maximal in an ideal $I$, use inclusion--exclusion for
$I\setminus\{p\}$ and $\mathord\downarrow p$ to obtain
\[
 [I]=[I\setminus\{p\}]+[\mathord\downarrow p]
              -[\mathord\downarrow p\setminus\{p\}].
\]
Induction on $|I|$ proves the inverse assertion and the universal
property. A homomorphism from the join monoid to an abelian group
sends $x\vee x=x$ to $2a=a$, hence every $x$ to zero.
For an antichain $P$ with $r$ elements, augmentation is
$(n_1,\ldots,n_r)\mapsto\sum_i n_i$; its kernel is precisely the
integral root lattice with basis $e_i-e_r$, $i<r$.
\hfill$\square$

\paragraph{Scope and geometric translation.}
This corrects the meaning of ``support linearization'' in
\path{03_split_zero_defect_preprint.tex}, its support section,
and \path{04_primitive_support_defects_preprint.tex}, the
primitive-support construction. A finite set of locally closed
strata gives an inclusion--exclusion valuation to its constructible
functions; assigning an Euler characteristic, a compactly supported
cohomology class, or a rank to each stratum is a further choice of
the weights $e_p\mapsto a_p$. For finite disjoint rational points
this is literal $K_0$ of their vector spaces. The external/direct
sum and tensor operations are the framework of cote 2, scan 17;
SGA 5, III, 4.7 concerns the subsequent trace realization, not a
canonical trace on the bare lattice. This theorem creates neither
witness multiplicities nor an effective counting algorithm from
Boolean support.



\section{NN-G006: faithful Artin--Tate zeta encoding without RH}
\paragraph{Statement.}
Let
$\mathcal A=\mathbb Z[q^{\pm1},\varepsilon]/(\varepsilon^2-1)$ and
$\chi=\chi_{-4}$. The assignment
\[
 a=\sum_j(a_j+b_j\varepsilon)q^j\longmapsto
 Z_a^{(2)}(s)=\prod_j\zeta^{(2)}(s-j)^{a_j}
                         L(s-j,\chi)^{b_j}                 \tag{G006.1}
\]
is injective, even as an identity of meromorphic functions. Here
$\zeta^{(2)}=(1-2^{-s})\zeta$ and all products are finite. It
identifies multiplication of coefficients with the Witt tensor
product of local Euler factors. No RH is needed for injectivity.
For the full completed products the correct factors are
\[
 \Lambda_\zeta(s)=\pi^{-s/2}\Gamma(s/2)\zeta(s),\qquad
 \Lambda_\chi(s)=(4/\pi)^{(s+1)/2}\Gamma((s+1)/2)L(s,\chi).
 \tag{G006.2}
\]
The signed pole order of
$\prod_j\Lambda_\zeta(s-j)^{a_j}\Lambda_\chi(s-j)^{b_j}$ at an
integer $m$ is $a_m+a_{m-1}$. Consequently these signed orders
recover the finitely supported sequence $(a_j)$. The conventional
entire $\xi(s)=\tfrac12s(s-1)\Lambda_\zeta(s)$ cannot be used
in this pole-recovery formula.

\paragraph{Proof.}
All factors in (G006.1) have absolutely convergent Dirichlet
series in a common sufficiently far right half-plane, even with
integer negative exponents. If their product is one, its logarithm
there has zero coefficients by uniqueness of absolutely convergent
Dirichlet series. A quick proof of that uniqueness takes the least
index with nonzero coefficient, multiplies by its $s$th power,
and lets real $s\to+\infty$; absolute convergence at a fixed
abscissa dominates the tail. At the coefficient indexed by $p^n$
this gives
\[
 \sum_j\{a_j+\chi(p)^n b_j\}p^{nj}=0\quad(n\ge1).
\]
At $p=5$, the distinct nonzero numbers $5^j$ and the Vandermonde
matrix for their first finitely many powers give $a_j+b_j=0$
for every $j$. At $p=3$, use odd $n=2r+1$ and the distinct bases
$3^{2j}$ to obtain $a_j-b_j=0$. Thus $a_j=b_j=0$. The same proof
works formally from the Euler factors.
Locally the Witt product sends
$(1-\alpha T)^{-1}$ and $(1-\beta T)^{-1}$ to
$(1-\alpha\beta T)^{-1}$. Because $\chi(p)^2=1$ at odd primes,
this realizes exactly multiplication in $\mathcal A$.
The classical functional equations (DLMF 25.4.1--25.4.4 and
25.15.5, with the elementary Gauss sum for $\chi_{-4}$) give
(G006.2) simple poles for
$\Lambda_\zeta$ at $0,1$, and no others, whereas $\Lambda_\chi$
is entire and nonzero on the real axis. The latter nonvanishing
follows for positive real $s$ by pairing successive terms of the
alternating beta series, and for $s\le0$ by its functional equation;
the value at zero follows from that at one. Similarly the completed
zeta is finite and nonzero at real points other than $0,1$.
Only shifts $j=m,m-1$ therefore contribute to the signed pole order
at $m$. Starting beyond the support and proceeding downwards
recovers $a_j$. Multiplying by $s(s-1)/2$ removes the poles of
$\Lambda_\zeta$, proving the warning about $\xi$.
\hfill$\square$

\paragraph{Inputs and exact dictionary.}
This repairs and strengthens the injectivity/pole statements in
\path{OURS_CONSULT_a157422a_COUNTING_CONSTRUCTION.md},
Theorem 4(a),(d),(e). The zeta and primitive odd Dirichlet functional
equations have the exact classical locators DLMF 25.4.1--25.4.4
and 25.15.5 (Apostol, \emph{Introduction to Analytic Number
Theory}, equations (12)--(16), pp.\ 259--263); A is required only when
one additionally locates their nontrivial zeros on the lines
$\operatorname{Re}s=j+1/2$. The geometric objects are Tate twists
of the trivial and Gaussian quadratic compatible rank-one systems.
For each $\ell$ their local systems are lisse on
$\operatorname{Spec}\mathbb Z[1/2\ell]$; the factor at $\ell$
is read using a different auxiliary prime. Their representation
ring is $\mathcal A$; SGA 3, I, 4.7.2--4.7.3 describes the
comodule/graded representation formalism. The Witt operations
are explicitly present in cote 43, scan 37. The local factor at
$2$ is omitted for tensor multiplicativity; restoring it as an
ordinary scalar correction does not make an unramified sign
character exist at the ramified prime. Faithful encoding is not
efficient decoding and does not prove PH1 or a counting bound.



\section{NN-G007: all effective lifts, and the two different effective cones}
\paragraph{Statement.}
For a Laurent polynomial $c=\sum_j c_jq^j\in\mathbb Z[q^{\pm1}]$,
write $c=c_+-c_-$ coefficientwise with disjoint nonnegative parts.
Every effective class $\widehat c$ in $\mathcal A$ with
$\widehat c|_{\varepsilon=-1}=c$ is uniquely of the form
\[
 \widehat c=c_++h+\varepsilon(c_-+h),\quad
 h\in\mathbb N[q^{\pm1}],\qquad
 d:=\widehat c|_{\varepsilon=1}=|c|+2h.             \tag{G007.1}
\]
The Burnside linearization for $C_2$ is an integral ring
isomorphism
\[
 B(C_2)=\mathbb Z[x]/(x^2-2x)\ \simeq\
 R(C_2)=\mathbb Z[\varepsilon]/(\varepsilon^2-1),
 \qquad x\mapsto1+\varepsilon.                    \tag{G007.2}
\]
Its effective-set cone maps to $a+b\varepsilon$ with
$a\ge b\ge0$, strictly smaller than the effective-representation
cone $a,b\ge0$. The nonzero regular class $1+\varepsilon$ is not
an idempotent: its square is $2(1+\varepsilon)$. It maps to the
supported zero under the quotient of the \emph{effective semiring}
\[
 \mathbb N[\varepsilon]/(\varepsilon^2-1)\longrightarrow G(\mathbb Z),
 \quad a+b\varepsilon\longmapsto
 \begin{cases}\tau,&a=b=0,\\a-b\text{ supported},&a+b>0.\end{cases}
\]
This is not a quotient homomorphism from the whole representation ring.

\paragraph{Proof.}
For each coefficient solve $a-b=c_j$ with $a,b\ge0$. If $c_j\ge0$
the solution is $(c_j+h_j,h_j)$, and if $c_j<0$ it is
$(h_j,-c_j+h_j)$. This proves (G007.1), including uniqueness.
A finite $C_2$-set is a disjoint union of fixed points and free
two-point orbits. The latter square is two free orbits, proving
the Burnside presentation. Substitution $\varepsilon=x-1$
is inverse to (G007.2). A set with $r$ fixed points and $s$ free
orbits has permutation representation $(r+s)+s\varepsilon$,
which gives exactly the asserted cone. The sign representation
$\varepsilon$ is effective but lies outside that cone. The
algebraic calculation of $(1+\varepsilon)^2$ proves nonidempotence.
The displayed map preserves the operations because evaluation at
$\varepsilon=-1$ does so, while a sum or product of two effective
classes is zero only in the usual empty cases. It is surjective:
use $a=n,b=0$ for positive $n$, $a=0,b=-n$ for negative $n$,
and $a=b=1$ for the supported zero.\hfill$\square$

\paragraph{Geometric consequences.}
At an inert odd prime the free orbit is
$\operatorname{Spec}\mathbb F_{p^2}$: it has no $\mathbb F_p$
points and two $\mathbb F_{p^2}$ points. Its geometric support,
its dimension, and its rational-point count are therefore distinct
invariants. Globally its Artin representation comes from
$\operatorname{Spec}\mathbb Z[i,1/2]$ and has zeta
$\zeta^{(2)}(s)L(s,\chi)$.
Changing $h$ in (G007.1) changes the Euler product by
\[
 \prod_j\zeta_{\mathbb Q(i)}^{(2)}(s-j)^{h_j}.
\]
A puts the nontrivial zeros of all these choices on their
appropriate lines, but selects none of the choices. Negative
Tate powers are legitimate twists of an $\ell$-adic object;
$q^{-1}$ itself is not the point count of a scheme over
$\mathbb F_q$. Thus an effective Artin--Tate representation is
not automatically an effective scheme count.
The exact archive comparison is cote 133, scan 19, where passing
from equivariant sets to representations must respect this
distinction; its subset/exterior-power formula already needs the
correction recorded in ALPOGE\_GCT PL-01. The group-scheme
representation formalism is SGA 3, I, 4.7.2. This gives no way to
recover satisfiability from a virtual cancellation or an
arbitrarily chosen dimension lift.



\section{NN-G008: rational descent does not imply a finite Tate count}
\paragraph{Statement and proof.}
The all-rank $\mathbb Q(D)$ descent in the counting consult does
not imply Laurent-polynomial structure constants in every
$\mathbb Q(D)$-rational basis. Already take the two-dimensional
algebra with basis $1,P$ and relation $P^2=DP$, and replace
$P$ by $b=P/(D^2+1)$. Then
\[
 b^2=\frac{D}{D^2+1}\,b.                            \tag{G008.1}
\]
This rational function is not a Laurent polynomial in $q$ for
$D=q+2+q^{-1}$. Indeed a Laurent polynomial invariant under
$q\leftrightarrow q^{-1}$ belongs to $\mathbb Q[q+q^{-1}]
=\mathbb Q[D]$: induction removes the highest pair
$q^m+q^{-m}$, itself a polynomial in $q+q^{-1}$. But
$D/(D^2+1)$ is not a polynomial in $D$, by division in
$\mathbb Q[D]$.\hfill$\square$

\paragraph{Correction and dictionary.}
Theorem 1(c) of the consult correctly distinguishes rational
descent from integrality; its Theorem 4(c) and final rank-four
summary do not preserve that distinction. A rational function
can also have an even-order pole at $v=1$, so ``even vanishing
order'' requires regularity there. For example rescaling by
$D-4$ gives poles. In the geometric dictionary a generic fibre
and a rational change of basis do not choose an integral lattice
or an effective class in $K_0$; SGA 3, I, 4.7.3 is a grading
equivalence, not an integrality theorem. Exterior powers and
effective direct summands, discussed in cote 133, scan 28 and
SGA 6, VI, Theorem 3.2, require actual vector-bundle classes.
The remaining rank-four gate includes a defined canonical basis,
integral Laurent constants, and an effective compatible lift.
The formal substitution $v=\sqrt q$ constructs neither a CM
Hecke character nor a compatible arithmetic family.



\section{NN-G009: a quantitative size obstruction for all-word Weil realizations}
\paragraph{Statement.}
Fix a prime power $q>1$ and $k\ge1$. If a finite-dimensional
Weil object $V_k$ has traces
\[
 \operatorname{Tr}(F^n\mid V_k)=
 40(q^{n/2}+q^{-n/2})^{2k}\quad(n\ge1),
 \tag{G009.1}
\]
then its semisimplification has eigenvalue $q^a$ with multiplicity
$40\binom{2k}{k+a}$ for every $-k\le a\le k$, and
$\dim V_k=40\cdot4^k$. If instead (G009.1) is an inserted trace
$\operatorname{Tr}(A_kF^n\mid V)$ on a fixed $V$, at least
$2k+1$ distinct nonzero eigenvalues of $F$ are required. In
particular no fixed finite $V$ carries the traces for all $k$.

\paragraph{Proof.}
The binomial expansion gives
\[
 40(q^{n/2}+q^{-n/2})^{2k}
 =\sum_{a=-k}^k40\binom{2k}{k+a}(q^a)^n.
\]
Distinct nonzero exponential sequences are linearly independent:
their values at the first $r$ positive integers form an invertible
Vandermonde matrix after scaling columns. Ordinary traces are
power sums of eigenvalues with their algebraic multiplicities,
even with Jordan blocks, so uniqueness proves the first assertion.
Their sum is $40\cdot2^{2k}$. For an inserted trace, a Jordan
decomposition expresses it as a sum of polynomials in $n$ times
$\lambda^n$ over eigenvalues $\lambda$ of $F$. These
exponential-polynomial sequences are independent: applying
successive difference operators $E-\lambda$ isolates each
base, or equivalently their generating functions have poles at
distinct reciprocals. Hence every base $q^a$ must occur among
the eigenvalues of $F$, proving the dimension lower bound.
\hfill$\square$

\paragraph{Inputs, geometry and complexity.}
The word-trace identity is GCT\_NSRH F42 and counting consult
Theorem 4(c). The new assertion quantifies the necessary size;
it does not rebrand F42 as new. A realization is $40$ copies of
$H^*((\mathbb P^1)^{2k})(k)$, with all cohomological degrees even.
It has a short tensor description even though its total rank is
exponential. Tensor products and trace are the operations of
cote 2, scan 17 and SGA 5, III, 4.7. This explains why a finite
list of multiplication objects can coexist with the absence of
a single bounded-rank Weil object. Neither a short tensor
description nor A gives a polynomial-size positive polytope
formula or rank-four canonical basis.



\section{NN-G010: the precise category supplied by a dimension ring}
\paragraph{Statement.}
Let $B$ be a finite based unital algebra over
$\mathbb Z[q^{\pm1}]$, with its unit among the basis elements.
Suppose its constants $c_{ab}^z$ and an associative unital
nonnegative table $d_{ab}^z$ satisfy
$d\ge|c|$ coefficientwise and $d\equiv c\pmod2$.
Then the effective table
\[
 \widehat c_{ab}^z=(d_{ab}^z+c_{ab}^z)/2
                +\varepsilon(d_{ab}^z-c_{ab}^z)/2
 \tag{G010.1}
\]
is realized by a monoidal additive category whose \emph{split}
Grothendieck ring is free over $\mathcal A$ on this basis. The
construction need not be idempotent complete, abelian, rigid,
semisimple as a category, or have a simple tensor unit. It
supplies no canonical basis or uniform complexity bound.

\paragraph{Proof.}
The map $\mathcal A\hookrightarrow
\mathbb Z[q^{\pm1}]\times\mathbb Z[q^{\pm1}]$ given by
evaluation at $\varepsilon=1,-1$ is injective, with image the
pairs congruent modulo two. Thus (G010.1) is associative and
unital. Let $\mathcal T$ be finite-dimensional
$\mathbb Z\times C_2$-graded vector spaces over a characteristic
zero field, with simple objects $q^j\varepsilon^a$ and their
ordinary tensor product. Its isomorphism classes are exactly
the effective elements of $\mathcal A$.
In $\operatorname{Mat}_{|B|}(\mathcal T)$ use matrix convolution
as tensor product. For each $a\in B$ choose the matrix object
$L_a$ with $(z,b)$ entry of class $\widehat c_{ab}^z$.
Associativity of the table gives an isomorphism
\[
 L_a\otimes L_b\simeq\bigoplus_z
        \widehat C_{ab}^z\otimes L_z,
\]
because objects of $\mathcal T$ are determined by their classes.
Take the replete full additive subcategory consisting of finite
sums of shifts $q^j\varepsilon^aL_b$. It is tensor closed and
inherits the coherent associativity constraints of matrices.
The unit matrix is $L_1$.
To compute its split $K_0$, inspect the column indexed by the
unit basis element: the unit law makes that column of $L_b$
the basis vector indexed by $b$. Therefore two sums of shifted
$L_b$ are isomorphic exactly when their effective coefficients
agree. Their additive monoid is free over the effective
coefficient monoid, and group completion gives the stated
free split $K_0$.
In the ambient full matrix category the endomorphisms of the
unit are $|B|$ independent scalar diagonal entries. Their
idempotents need not split within the chosen subcategory.
Thus neither idempotent completeness nor a simple unit follows;
taking a Karoubi or abelian completion changes the question.
\hfill$\square$

\paragraph{Geometric translation.}
For $X$ a disjoint union of $|B|$ rational points, matrices of
Tate--Artin objects are kernels on $X\times X$. Convolution is
$p_{13!}(p_{12}^*(-)\otimes p_{23}^*(-))$. This is the finite
zero-dimensional instance of the external product and projection
formula on cote 2, scan 17; Artin gradings use SGA 3, I, 4.7.2.
The construction refines counting consult Theorem 5 and
Corollary 8 by stating precisely which Grothendieck group is
proved. It uses the already known structure constants to build
the objects. It cannot establish their positivity without those
constants, nor establish PH0 or PH1 for a canonical geometric
construction.



\section{NN-G011: a smaller composing rank-three dimension ring}
\paragraph{Statement.}
The supplied rank-three algebra admits a dimension ring with
the same 132 nonzero positions and a total of $7,276$
coefficientwise cancelling pairs, of which $6,612$ occur in
the bottom-cell coefficients. This improves the supplied
search-window value $7,952$. No claim of global minimality is
made. The parameters are
\[
 \begin{gathered}
 D=q+2+q^{-1},\quad\delta=(D-1)(D-2),\quad W=7D-8,\\
 \Lambda_1=D,\quad\Lambda_2=3D^2-10D+4,\quad
 G_1=2W\delta,\quad G_2=W\delta,\quad S=2W\delta^2,\\
 U_1=13D-16,\quad U_2=3D-4,\quad
 V_1=14D^2-17D,\quad V_2=2D^2+D-2.
 \end{gathered}                                             \tag{G011.1}
\]
For $x,y\in\{1,12,2,21\}$ write the primitive word reduction
$P_xP_y=D^\alpha\Lambda_i^\beta P_z$, where $\beta=0,1$.
Keep $\mu$ as unit and define
\[
 \begin{aligned}
 \Sigma^2&=S\Sigma,&
 \Sigma\gamma_x^i=\gamma_x^i\Sigma&=D^{|x|-1}G_i\Sigma,\\
 \gamma_x^i\gamma_y^j&=D^{|x|+|y|-2}W\Sigma&&(i\ne j),\\
 \gamma_x^i\gamma_y^i&=\delta D^\alpha\Lambda_i^\beta\gamma_z^i
       +\begin{cases}
          D^{|z|+\alpha-2}U_i\Sigma,&\beta=0,\\
          D^{|z|+\alpha-1}V_i\Sigma,&\beta=1.
        \end{cases}
 \end{aligned}                                             \tag{G011.2}
\]
The exponent in the $\beta=0$ case is nonnegative for the
eight word pairs where it occurs.

\paragraph{Proof of associativity.}
Over $\mathbb Q(D)$ take two $2\times2$ matrix blocks with
\[
 P_1=\begin{pmatrix}D&\Lambda_i\\0&0\end{pmatrix},\qquad
 P_2=\begin{pmatrix}0&0\\1&D\end{pmatrix}.
\]
They obey $P_j^2=DP_j$ and $P_1P_2P_1=\Lambda_iP_1$,
$P_2P_1P_2=\Lambda_iP_2$. The four $P_x$ are linearly
independent generically because
$\Lambda_i(D^2-\Lambda_i)\ne0$. In the direct product of
these blocks and two scalar blocks send $\gamma_x^i$ to
$\delta P_x$ in its own matrix block, to zero in the other,
to $D^{|x|-1}G_i$ in the first scalar block and to zero in
the last. Send $\Sigma$ to $S$ in the first scalar block
and zero elsewhere, and $\mu$ to the identity in every
block. These ten images are independent, successively by
the last scalar, the two matrix blocks, and the first scalar.
They are closed under products because
\[
 \frac{G_1G_2}{S}=W,\quad
 \frac{G_i(G_i-\delta D)}S=U_i,\quad
 \frac{G_i(DG_i-\delta\Lambda_i)}S=V_i.
\]
The first ratio follows immediately from (G011.1); the next
two equal $2W-D,(W-D)/2$ and
$2DW-\Lambda_1,(DW-\Lambda_2)/2$, respectively.
They are exactly the polynomials displayed there. Evaluating
the scalar block gives (G011.2). The embedded algebra is
associative over $\mathbb Q(D)$, hence all its polynomial
structure identities hold over $\mathbb Z[D]$.

\paragraph{Proof of effectivity, domination and parity.}
Write $u=q+q^{-1}=D-2$. The true cell parameters are
$\lambda_1=D$, $\lambda_2=u^2$ and
$\delta_2=-\delta_1=(D-1)(D-4)$; the latter has Laurent
coefficients $q^2-q-q^{-1}+q^{-2}$.
Our $\delta$ is its coefficientwise absolute value plus $2$.
Moreover
\[
 \Lambda_2-\lambda_2=2(q^2+q^{-2})+2(q+q^{-1})\ge0.
\]
The true bottom parameters are
$W_0=U_{1,0}=U_{2,0}=D$,
$V_{1,0}=5D-4$, $V_{2,0}=2D^2-5D+4$,
$G_{1,0}=4D(D-1)$, $G_{2,0}=D^2(D-1)$, and
$S_0=4D^2(D-1)^2$. Their differences from (G011.1) are
\[
 \begin{array}{c|l}
 W-W_0&6u+4\\
 U_1-U_{1,0}&12u+8\\
 U_2-U_{2,0}&2u\\
 V_1-V_{1,0}&14u^2+34u+16\\
 V_2-V_{2,0}&6u+6\\
 G_1-G_{1,0}&14(q^3+q^{-3})+22(q^2+q^{-2})+42u+36\\
 G_2-G_{2,0}&6(q^3+q^{-3})+8(q^2+q^{-2})+16u+12\\
 S-S_0&(u+1)^2\{14(q^3+q^{-3})+8(q^2+q^{-2})+26u\}.
 \end{array}                                               \tag{G011.3}
\]
Every difference is coefficientwise nonnegative and even.
All true bottom parameters are nonnegative Laurent
polynomials, so the new ones are also. Multiplication by
nonnegative powers of $D$ preserves these properties.
For within-cell coefficients use
$|fg|\le|f||g|$ coefficientwise, together with the inequalities
for $\delta$ and $\Lambda_i$. Congruence modulo two is
preserved by sums and products and is unaffected by the
opposite signs of the two true $\delta_i$. This proves every
dimension-ring condition for all 132 positions.

\paragraph{Exact score.}
The nineteen unit positions contribute no excess. In the
true table the sum of absolute Laurent coefficients of the
bottom constants is $4096$. This is obtained at $D=4$ from
\[
 S_0+20(G_{1,0}+G_{2,0})+200W_0
           +20(U_{1,0}+U_{2,0}+V_{1,0}+V_{2,0})=4096.
\]
For within-cell constants the types
$(\alpha,\beta)=(0,0),(1,0),(0,1),(1,1)$ have multiplicities
$2,6,6,2$. The absolute coefficient norms in cell one are
$4,8,8,24$, and in cell two $4,8,16,12$, by direct expansion
of $\delta_iD^\alpha\lambda_i^\beta$. Their weighted sum is
$328$. Thus the total true absolute norm is $4424$.
For the new table, at $D=4$ the values are
\[
 (\delta,W,G_1,G_2,S,U_1,U_2,V_1,V_2,\Lambda_1,\Lambda_2)
 =(6,20,240,120,1440,36,8,156,34,4,12).
\]
The same bottom formula gives $17320$, and the within-cell
sum is $6(26+14\cdot4)+6(26+14\cdot12)=1656$.
Hence the excess divided by two is
$(18976-4424)/2=7276$, with bottom part
$(17320-4096)/2=6612$.\hfill$\square$

\paragraph{Inputs, certificates and meaning.}
The exact input is counting consult Theorem 2 and
\path{OURS_E4_scripts/e3_constants.json}.
\path{verify_dimension_7276.py} independently builds the
table with integer Laurent dictionaries, checks all 1000
associativity triples, and verifies all coefficient inequalities
against the supplied 132 constants. It does not import the
supplied dimension-ring builder. The full certificate is
\path{dimension_certificate_7276.json}.
NN-G010 converts the result into a finite-kernel monoidal
category, with the same SGA 3/cote 2 dictionary. This lowers
an explicit upper bound for the noncanonical cancellation
cost. It neither selects a canonical lift, proves global
optimality, nor decides PH0/PH1 at rank four. No unimodality
claim is made for this particular smaller lift.



\section{NN-G012: why the sign-point trace symmetry reappears}
\paragraph{Statement.}
In the specialization $q=-1$, $D=0$ of the supplied based
rank-three algebra, let $\gamma_x=\gamma_x^2$ and set
\[
 \beta_1=\gamma_1+\Sigma/4,\quad
 \beta_2=\gamma_2+\Sigma/4,\quad
 \beta_{12}=\gamma_{12},\quad\beta_{21}=\gamma_{21}.
\]
Their span is a matrix subalgebra over $\mathbb Q$, identified
with $M_2$ by
\[
 \beta_1\mapsto8E_{12},\quad\beta_2\mapsto8E_{21},\quad
 \beta_{12}\mapsto16E_{11},\quad\beta_{21}\mapsto16E_{22}.
 \tag{G012.1}
\]
Here $\Sigma^2=0$ and $\Sigma$ annihilates this subalgebra.
The matrix-trace Gram matrix, in the order $1,12,21,2$, is
\[
 \begin{pmatrix}0&0&0&64\\0&256&0&0\\0&0&256&0\\64&0&0&0\end{pmatrix}.
 \tag{G012.2}
\]
Thus the cyclic slot relation is exactly the ordinary cyclic
trace identity in a rescaled matrix-unit basis. This explains
the monomial pairing at this specialization without asserting
a symmetry of the generic table or of its six unimodality
failures.

\paragraph{Proof.}
At $D=0$ the true parameters are $\delta_2=4$,
$\lambda_2=4$; every $\Sigma\gamma_x$ and $\Sigma^2$ is zero.
Use $P_1=2E_{12}$, $P_2=2E_{21}$. The within-cell
coefficients agree with the products of $4P_x$, and the
only nonzero bottom coefficients in its sixteen products
are $4\Sigma$ in the four products whose cell output is
$16\gamma_1$ or $16\gamma_2$. Replacing those outputs by
$16\beta_1$ or $16\beta_2$ absorbs exactly $4\Sigma$;
all other products require no correction. This proves closure
and (G012.1). The four images are independent, so the map is
an isomorphism. Direct matrix multiplication gives (G012.2).
For the supplied tensor-Hecke trace, direct evaluation
with its exact arithmetic (\path{check_signpoint_Hecke_trace.py},
reusing the supplied implementation) shows that it vanishes on $\Sigma,\gamma_1,
\gamma_2$ and takes value $16$ on $\gamma_{12},\gamma_{21}$;
it therefore restricts to this matrix trace. Finally
$\operatorname{tr}(ABC)=\operatorname{tr}(BCA)$ is a direct
index calculation.\hfill$\square$

\paragraph{Geometry and scope.}
The full sixteen-product certificate is
\path{verify_signpoint_matrix_block.py} and
\path{signpoint_matrix_certificate.json}. The four-dimensional
kernel algebra on a two-point scheme is again the convolution
algebra of cote 2, scan 17; (G012.1) is its explicit basis change.
The trace cyclicity is the finite-dimensional case of SGA 6,
I, section 8 and SGA 5, III. At the sign point, the cyclic
permutation mixes four failed generic coefficients with two
shortest coefficients, as E4 section 6b records. Accordingly
this explains the restored trace symmetry, not generic
unimodality or the existence of a rank-four canonical basis.



\section{NN-G013: the missing operation in positive-formula isolation}
\paragraph{Statement.}
Suppose a nonnegative multiplicity $g(x)$ has a uniform witness
formula $g(x)=|W_x|$ with $W_x\subseteq\{0,1\}^{m(x)}$,
$m$ polynomially bounded, and polynomial-time verification.
If the parity of $|W_x\cap h^{-1}(0)|$ can be decided in
deterministic polynomial time for every affine Boolean hash
$h$, then positivity of $g$ is in RP. If this applies to one
of the unary NP-hard plethysm or Kronecker positivity families,
then $\mathrm{NP}=\mathrm{RP}$. The positive witness formula
alone only places positivity in NP.

\paragraph{Proof.}
Choose $k$ uniformly from $1,\ldots,m+2$ and an affine map
$h:\mathbb F_2^m\to\mathbb F_2^k$ uniformly. For distinct
inputs its outputs are independent and uniform. If
$N=|W_x|>0$, some such $k$ makes
$\mu=N2^{-k}\in[1/4,1/2]$. Put
$X=|W_x\cap h^{-1}(0)|$. Then
$\mathbb E X=\mu$ and
$\mathbb E X(X-1)=N(N-1)2^{-2k}\le\mu^2$.
Since $\mathbf1_{X=1}\ge X-X(X-1)$ for every nonnegative
integer $X$, the probability of a singleton is at least
$\mu-\mu^2\ge3/16$. A random $k$ therefore succeeds with
probability at least $3/(16(m+2))$. On that event parity is
one. On an empty witness set parity is always zero, so
polynomial repetition gives an RP algorithm. NP-hardness
gives $\mathrm{NP}\subseteq\mathrm{RP}$, and the usual
polynomial random string is an NP certificate for an RP
acceptance, proving equality.\hfill$\square$

\paragraph{Exact inputs, signed failure and complexity scope.}
The unary hardness locators are Fischer--Ikenmeyer,
Theorem 3.1, label \path{thm:plethysm-positivity-np-hard},
and Ikenmeyer--Mulmuley--Walter, \texttt{tkronnp} in the
supplied main source. The latter's positive formula for
an NP-hard subclass is \texttt{tkronpos}, so positivity of
a formula is not itself a tractability theorem.
For a signed sum isolation tests the presence of summands:
on $S_2$, $g((1,1),(1,1),(1,1))=(1-1)/2=0$, although
both character summands exist. A and M do not supply the
hashed parity oracle. E can evaluate an explicitly encoded
witness count after a parsimonious bounded-width CNF
encoding, but its exponent is measured in that encoding's
variables, including verifier and hash auxiliary variables;
the word ``polynomial'' does not preserve a linear exponent.
There is no derivation of NP=RP from the stated inputs.

\paragraph{Geometric translation.}
The witness relation and affine hashes are equations over
$\mathbb F_2$ on Boolean coordinates; verification constraints
give a finite affine scheme whose rational points are the
witnesses after parsimonious encoding. Its point-count
parity is additional information beyond support of its
terms. SGA 5, III, 4.7 interprets counts as alternating
Frobenius traces but does not evaluate them in polynomial
time. The finite-set/representation distinction is
NN-G007 and cote 133, scan 19. An arbitrary split-zero
coefficient formula is not such a witness scheme until
this identification is constructed.



\section{NN-G014: what a unique extremal localization actually proves}
\paragraph{Statement and proof.}
Let $K$ be a field with a valuation $\nu$ and let
$F=\sum_{i=1}^rF_i$. If exactly one summand $F_j$ has
least valuation, then
$\nu(F)=\nu(F_j)$ and its leading residue equals that of
$F_j$ in the associated graded ring. To see this divide
by $F_j$: every other ratio has positive valuation, so
$F/F_j$ has residue one. This proves $F\ne0$ as a
function, not $F(1)\ne0$ after a specialization: $1-t$
is the elementary counterexample.
For $K=\mathbb Q((t))$ and an actual coefficient formula
$g=\operatorname{CT}_tF$,
the stronger hypotheses $\nu(F_i)\ge0$ for all $i$, a unique
term of valuation zero, and positive residue for that term
prove $g>0$. More generally a unique exposed weight in a
multigraded expansion isolates its coefficient, provided
the target coefficient is exactly that exposed weight.
\hfill$\square$

\paragraph{Reduced gate and dictionary.}
For a GCT positivity algorithm this leaves three precise
tasks: a coefficient-preserving localization, polynomially
bounded certificates for the target exposure and residue,
and a procedure finding such a certificate for every yes
instance. A deterministic polynomial-time complete procedure
would yield P=NP on the known hard families; randomized
one-sided completeness yields NP=RP. A bare unique smallest
signed term does neither. The source is the question left
by the GCT support consult, section on P-F30/P-F23 and
reverse isolation. Geometrically a one-parameter subgroup
of a torus filters an equivariant module by weights; the
associated graded residue retains the exposed piece. The
exact action/graded-module dictionary is SGA 3, I, 4.7.3;
the projective/graded localization construction is EGA II,
2.2--2.5. No degeneration has been asserted to preserve
the desired multiplicity without the coefficient identity.



\section{NN-G015: support inclusion cannot decide multiplicity dominance}
\paragraph{Statement.}
Comparing two strictly positive succinct witness counts
is already PP-hard, even when their Boolean supports are
identically nonempty. Consequently the support inclusion
supplied by BIP cannot, solely as support information,
decide the multiplicity inequalities that remain possible
obstructions. This is a limitation of the information,
not a PP-hardness theorem for the actual GCT multiplicities.

\paragraph{Proof.}
Given a Boolean circuit on $n\ge1$ input bits take
$a(x)=\#\operatorname{CircuitSAT}(x)+1$ and
$b(x)=2^{n-1}+1$. Both have positive polynomially
verifiable witness counts: adjoin one distinguished
witness to the satisfying assignments, respectively to
all strings of length $n-1$. Then
$a(x)>b(x)$ exactly when more than half the assignments
satisfy the circuit. Majority circuit acceptance is the
defining PP computation with its polynomial-time verifier
unrolled into a polynomial-size Boolean circuit. The support bits of both counts are
always one and contain none of that comparison data.
For orbit closures $Z\subseteq\Omega$, restriction gives
the graded equivariant surjection
$\mathbb C[\Omega]\twoheadrightarrow\mathbb C[Z]$.
Complete reducibility degree by degree implies
$m_Z(\lambda)\le m_\Omega(\lambda)$. Hence a strict
reverse inequality obstructs containment, whereas
support inclusion alone asserts only that a positive
left multiplicity has a positive right multiplicity.
\hfill$\square$

\paragraph{Inputs, border scope and geometry.}
BIP, \emph{No occurrence obstructions in geometric
complexity theory}, Theorem 1.1, gives its stated
$n\ge m^{25}$ support inclusion for padded permanent
and determinant orbit closure. It does not prove
containment and does not eliminate multiplicity
obstructions or other equations. Orbit closure is the
geometric object for border complexity; confusing it
with the orbit or the support set would discard the
defining ideal. In Grothendieck language the needed
object is the equivariant closed subscheme, its
quotient coordinate algebra, and its graded locally
free representation pieces. Cote 133, scan 28 supplies
the Schur/tensor operations on their $K_0$ classes;
cote 2, scan 17 supplies restriction and tensor products.
No FGA representability theorem by itself computes the
ideal or the multiplicity comparison. The example in
this proof can be realized by finite Boolean witness
schemes with trivial group action, using succinct
equations; its explicit vector space may have
exponentially many basis vectors. Thus it makes no
unproved reduction into the BIP orbit-closure family.



\section{NN-G016: a cutoff projection repairs the finite-prime boundary calculation}
\paragraph{Statement and exact model.}
Fix $\ell>0$. Give $\mathbb Z\cup\{\infty\}$ the topology
in which the integers are isolated and $m\to\infty$ means
$m\to+\infty$. Put
\[
Y=(\mathbb Z\cup\{\infty\})\times\mathbb R,
\qquad T(m,y)=(m+1,y+\ell),\quad
T(\infty,y)=(\infty,y+\ell).
\]
For the crossed-product extension
\[
0\longrightarrow I=C_0(\mathbb Z\times\mathbb R)\rtimes_T\mathbb Z
\longrightarrow E=C_0(Y)\rtimes_T\mathbb Z
\longrightarrow Q=C_0(\mathbb R)\rtimes_{+\ell}\mathbb Z
\longrightarrow0,
\]
the connecting map $K_0(Q)\simeq\mathbb Z\to K_1(I)\simeq\mathbb Z$
is multiplication by $-1$ when the boundary rank-one class
is positive and increasing argument is the positive winding convention.
In particular it is an isomorphism. The proof supplies the
missing rank-one projection and a self-adjoint lift in $E$.

\paragraph{Proof: an explicit rank-one corner.}
Write $UfU^{-1}=\alpha_1(f)$, where
$\alpha_n(f)(m,y)=f(m-n,y-n\ell)$. Define the real cutoff
\[
c(y)=\begin{cases}\cos(\pi y/(2\ell)),& |y|\le\ell,\\0,&|y|>\ell.\end{cases}
\]
On every interval of length $\ell$ the two possible nonzero
terms are a sine and a cosine, so
$\sum_{n\in\mathbb Z}c(y-n\ell)^2=1$. Regarding $c$ as a
bounded continuous function on $Y$, independent of $m$, put
\[
P=\sum_{n\in\mathbb Z}c(y)c(y-n\ell)U^n\ \in M(E).
\]
This sum has only the three terms $n=-1,0,1$. It is
self-adjoint by the crossed-product involution. The coefficient
of $U^k$ in $P^2$ is
\[
\sum_n c(y)c(y-n\ell)^2c(y-k\ell)
=c(y)c(y-k\ell),
\]
so $P$ is a projection. Its boundary restriction $p$ belongs
to $Q$, because each coefficient is compactly supported in $y$.
In the boundary orbit representation indexed by
$\theta\in\mathbb R/\ell\mathbb Z$, its matrix is
\[
p_\theta(i,j)=c(\theta+i\ell)c(\theta+j\ell).
\]
It is the rank-one projection onto the unit vector with
entries $c(\theta+i\ell)$. The action on $\mathbb R$ is free
and proper. Its crossed product is locally the compact operators
over the circle; a rank-one projection nonzero on every fiber is
full. Thus $pQp\simeq C(\mathbb R/\ell\mathbb Z)$, and $[p]$
is the positive generator of $K_0(Q)$. More directly, compression
is scalar on every rank-one fiber, and multiplying $p$ by an
arbitrary continuous $\ell$-periodic function realizes every
scalar continuous section. This also exhibits the unit of the corner.

\paragraph{Proof: the lift and its winding.}
Set
\[
\varphi(t)=\tfrac12-\pi^{-1}\arctan t,
\qquad h(m,y)=\varphi(y-m\ell),\qquad h(\infty,y)=1.
\]
This is bounded, continuous and $T$-invariant. Hence $h$ is a
central multiplier of $E$. It is not asserted to belong to $E$.
Nevertheless $H=hP$ does belong to $E$: each of its three
coefficients has bounded $y$ support; as $m\to-\infty$ it tends
uniformly to zero there; as $m\to+\infty$ it tends uniformly to
its specified boundary coefficient. These are precisely the
vanishing-at-infinity and continuity conditions on this support.
Thus $H=H^*$ is a self-adjoint lift of $p$.

The exponential formula for the connecting map now gives
\[
\partial[p]=[u],\qquad
u=\exp(2\pi iH)=1+(\exp(2\pi ih)-1)P.
\]
On the open part use $t=y-m\ell$. This identifies its transformation
group with the product of $\mathbb R_t$ and the translation action
of $\mathbb Z$ on itself. Accordingly
$I\simeq C_0(\mathbb R_t,\mathcal K(\ell^2\mathbb Z))$.
In this description
\[
P_t=|v_t\rangle\langle v_t|,\quad
v_t(m)=c(t+m\ell),\quad \|v_t\|=1,
\qquad
u(t)=1+(w(t)-1)P_t,
\]
where
\[
w(t)=\exp(2\pi i\varphi(t))=\frac{t+i}{t-i}.
\]
The vectors $v_t$ are norm-continuous: locally only finitely many
coordinates can be nonzero. They trivialize a full rank-one
corner of $I$. The Morita map therefore sends $[u]$ to the scalar
class $[w]\in K_1(C_0(\mathbb R))$. In particular
$u-1\in I$, also seen directly from $|w(t)-1|\to0$ at both ends.
The continuous argument $2\pi\varphi(t)$ decreases from $2\pi$
to $0$ as $t$ increases from $-\infty$ to $+\infty$.
The resulting loop has winding number $-1$, which proves the claim.
\hfill$\square$

\paragraph{Source correction, scope and geometry.}
The explicit model is the one in corpus paper 03, lines 850--998,
and paper 04, lines 358--504, with $\ell=\log p$. Their scalar
$h$ cannot by itself lift a rank-one projection: its boundary
value is a multiplier unit in a nonunital algebra. The replacement
$hP$ removes precisely that gap. The classical analytic inputs
are Blackadar, \emph{K-Theory for Operator Algebras}, second edition,
\S9.3.2 (printed page 68), for the exponential connecting map,
and Williams, \emph{Crossed Products of C*-Algebras}, author draft
version 3.1 (6 September 2006), Corollary 4.11 and Remark 4.12
(printed page 126), for the free proper action Morita equivalence.
Both locators were checked in the primary author PDFs.

The group action on all of $Y$ is also free and proper: compact
sets have bounded $y$ coordinates, so only finitely many translates
can meet them. Thus the same free proper action theorem applies
to $E$, compatibly with its open and closed pieces. Transferring
the calculation to a larger arithmetic quotient still requires the
claimed identification of its edge extension with this exact model.
No verification of that global identification is smuggled into
the present model theorem, and no purity or complexity bound follows.

An honest algebraic-geometric comparison is available for the
scalar loop: $(t+i)/(t-i)$ is a rational function on
$\mathbb P^1_{\mathbb C}$ with divisor $[-i]-[i]$ and unit norm
on the real locus. Equivalently its real coordinates are
$((t^2-1)/(t^2+1),2t/(t^2+1))$ on the norm-one torus
$x^2+y^2=1$; the removed degree-two divisor $t^2+1=0$ has no
real points. The zero principal divisor class in the Picard
group and the nonzero winding on the real locus are different
invariants. In Grothendieck language a rank-one corner is an
actual rank-one object before taking $K_0$, not the literal unit
of a multiplier algebra. The supplied archive cote 2, scan 17
provides the analogous pullback/tensor formalism for geometric
objects; it does not identify this analytic $K_1$ class with an
\'{e}tale trace or Weil weight. There is no claimed EGA/SGA theorem
turning the crossed-product model into an algebraic scheme.



\section{NN-G017: fixed smooth arithmetic geometry permits exponentially distant actual grading holes}
\paragraph{Statement.}
For every integer $L\ge2$ equip the same smooth scheme
$X=\mathbb A^2_{\mathbb Z}$ with the $\mathbb G_m$ action
\[
t\cdot(x,y)=(t^Lx,t^{L+1}y).
\]
Its graded coordinate algebra has support
$S_L=\langle L,L+1\rangle$, largest missing nonnegative degree
$L^2-L-1$, conductor $L(L-1)$, and exactly $L(L-1)/2$ holes.
These are actual missing graded pieces, rather than negative
off-residue evaluations of quasipolynomial constituents.
Nevertheless for every $L$ the underlying scheme has point counts
$q^{2r}$ over $\mathbb F_{q^r}$, global zeta $\zeta(s-2)$,
and compactly supported Weil cohomology consisting of
$\mathbb Q_\ell(-2)$ in degree four. Even the induced
$\mathbb G_m$ action on that cohomology is trivial for every $L$.
Thus these arithmetic and cohomological data cannot bound the
grading conductor polynomially in the binary size of the weights.

\paragraph{Proof of the exact holes and the easy counting rule.}
The monomials $x^ay^b$ are a homogeneous basis, so
\[
H_L(z)=\frac1{(1-z^L)(1-z^{L+1})},\qquad
h_L(n)=\#\{(a,b)\in\mathbb N^2:La+(L+1)b=n\}.
\]
Write $n=QL+r$ with $0\le r<L$. Every possible $b$ has the
form $r+jL$ with $j\ge0$, and its associated $a$ is
$Q-r-j(L+1)$. Consequently
\[
h_L(n)=
\begin{cases}
0,&Q<r,\\
1+\left\lfloor\dfrac{Q-r}{L+1}\right\rfloor,&Q\ge r.
\end{cases}                                                \tag{G017.1}
\]
In residue class $r$ the missing degrees are exactly
$r,\ L+r,\ldots,(r-1)L+r$ if $r>0$, and there are none
if $r=0$. There are therefore $\sum_{r=1}^{L-1}r=L(L-1)/2$
holes. Their largest element is attained for $r=L-1$ and
is $(L-2)L+(L-1)=L^2-L-1$. Every larger degree is present,
which proves the conductor assertion. Formula (G017.1)
decides membership and computes the dimension using a constant
number of integer divisions on $O(\log(n+1)+\log L)$ bits.
It also gives a positive two-variable polytope formula with
coefficients of polynomial bitlength. A large hole depth here
does not make this positivity problem hard.

\paragraph{Proof of arithmetic and cohomological indistinguishability.}
The underlying scheme is independent of the action; its point
count is $q^{2r}$. Taking the exponential generating series
gives $(1-q^2T)^{-1}$, and multiplying over primes gives
$\zeta(s-2)$. The standard compact-support calculation for
affine space gives the displayed Tate object. To make the
classical inputs explicit, compute $H^*(\mathbb P^1)$ over
$\mathbb C$ by SGA 4, XVI, Theorem 4.1, specialize the constant
prime-to-characteristic coefficients by XVI, Corollary 2.2,
and apply the open--closed exact sequence for
$\mathbb A^1\subset\mathbb P^1$ with complement one point.
The result is $H_c^2(\mathbb A^1)=\mathbb Q_\ell(-1)$ and
zero in other degrees. SGA 4, XVII, Theorem 5.4.3 then gives
the affine-plane calculation by K\"unneth. The Tate
normalization is exactly SGA 4, XVIII, (2.8.1). The paragraph
following that formula states that connected algebraic groups
acting on affine space act trivially on this top trace sheaf;
in particular every one of the above $\mathbb G_m$ actions
induces the same trivial action on $H_c^4$. All these primary
locators were opened in this run. Since specifying $(L,L+1)$ takes $\Theta(\log L)$ bits, its quadratic conductor is
exponential in the grading description length.\hfill$\square$

\paragraph{Geometric dictionary and consequences.}
This is already a theorem about a fixed smooth normal scheme
over $\mathbb Z$ with varying diagonal group-scheme actions.
SGA 3, I, Proposition 4.7.3 and Corollary 4.7.3.1 identify
the actions with the coordinate-ring gradings. The saturated
support is $\mathbb N$: the group generated by $L,L+1$ is
$\mathbb Z$ and their rational cone is $\mathbb R_{\ge0}$.
For the associated semigroup curve
$\operatorname{Spec}\mathbb Z[t^L,t^{L+1}]$, the finite
normalization is $\operatorname{Spec}\mathbb Z[t]$ and the
conductor ideal in $\mathbb Z[t]$ is
$t^{L(L-1)}\mathbb Z[t]$. Indeed $t$ is integral because
$t^L$ belongs to the subring, the two fraction fields agree
because $t=t^{L+1}/t^L$, and $\mathbb Z[t]$ is integrally
closed. The conductor calculation follows termwise from the
already proved description of the holes. This semigroup
curve is an associated encoding of the support; it is not
the smooth scheme $X$ itself.

A makes the nontrivial zeros of the common arithmetic zeta
lie on $\operatorname{Re}s=5/2$, but it supplies no missing
weight spaces and changes none of the holes. A fixed finite
pure Weil realization of an \emph{ungraded} count, even with
its induced connected-group action, therefore cannot stand
in for the grade-indexed trace or witness objects needed by
NSRH and PH1. The two-variable formula here already has
the desired positive counting form; this example refutes a
proposed inference about conductor bounds, not PH1 itself
and not a complexity conjecture about actual plethysm.



\section{NN-G018: the cubic pole forces exponential rank in any stretching-as-iteration realization}
\paragraph{Statement.}
Let $f_n(r)$ be the ordinary cubic plethysm stretching sequence
in P-F31, with $n\ge4$ and $M=2^{n-1}$. Suppose that over a
characteristic-zero field there are a finite-dimensional vector
space $V_n$ and endomorphisms $A_n,F_n$ such that
\[
f_n(r)=\operatorname{Tr}(A_nF_n^r\mid V_n)
\]
for all sufficiently large integers $r$. Then
\[
\dim V_n\ \ge\ \varphi(M)=2^{n-2}.
\tag{G018.1}
\]
No commutation between $A_n$ and $F_n$ is assumed, and Jordan
blocks are allowed. The same bound applies to the sum of the
dimensions of a finite complex when alternating traces with
insertions are used. This is a rank bound for interpreting the
stretching index as iteration of one operator; it is not a
bound on descriptions of a family of separate graded objects.

\paragraph{Proof and a sharp general linear-algebra formulation.}
For any $d$-dimensional $V$ one has the formal identity
\[
\sum_{r\ge0}\operatorname{Tr}(AF^r)t^r
=\operatorname{Tr}\bigl(A(I-tF)^{-1}\bigr).
\]
The adjugate formula shows that its denominator divides
$\det(I-tF)$, which has degree at most $d$. Altering finitely
many initial coefficients adds a polynomial and changes no
finite poles or their orders. Therefore the total degree of
the reduced denominator of any rational sequence realized
eventually in this manner is at most $d$.

For completeness, this is the exact minimum dimension for
eventual realizations over an algebraically closed
characteristic-zero field. Write the partial fraction expansion
of the rational generating function after discarding its
polynomial part. A pole of order $e$ at $t=\lambda^{-1}$
contributes a sequence $\lambda^rP(r)$ with $\deg P=e-1$.
Use one Jordan block $F=\lambda I+N$ of size $e$. The identity
\[
F^r=\sum_{j=0}^{e-1}\binom rj\lambda^{r-j}N^j
\]
shows that the linear functionals $\operatorname{Tr}(AF^r)$
span all such sequences: choose entries of $A$ on the $j$th
lower diagonal to prescribe $\operatorname{Tr}(AN^j)$, and
use that the polynomials $\binom rj$ form a basis. These
diagonals are disjoint, so the choices are independent.
Taking the direct sum over poles gives dimension equal to
the sum of their orders. Together with the lower bound this
proves exactness for eventual, inserted-trace realizations.

P-F31 asserts a genuine simple pole at every primitive
$M$th root. There are $\varphi(M)=M/2$ such roots, so the
reduced denominator has at least that degree and (G018.1)
follows. One could equivalently start from a single primitive
pole, since the Hilbert series has rational coefficients and
its Galois conjugates are also poles. For an alternating
inserted trace replace the collection of spaces by their
direct sum, $F$ by its block diagonal operator and $A$ by
the block diagonal insertions with their signs. The same
denominator argument applies.\hfill$\square$

\paragraph{Ordinary traces are more restrictive.}
For $A=I$, Jordan blocks contribute only their sizes times
$\lambda^r$ to the trace. Thus a finite sum, difference or
cohomological alternating sum of ordinary traces has only
simple poles in its generating series. Any quasipolynomial
having a nonconstant polynomial constituent has a pole of
order greater than one at some root of unity: decompose its
periodic coefficient functions into their finite Fourier
series, obtaining $\sum_\rho P_\rho(r)\rho^r$; a nonconstant
constituent forces some $P_\rho$ to have positive degree.
Its pole of order $1+\deg P_\rho$ cannot cancel against
terms based at other roots. Such a sequence cannot be an
ordinary trace sequence on any fixed finite complex.
This last statement is a general criterion; the rank bound
above needs only the specifically cited P-F31 simple poles.

\paragraph{Geometric dictionary and exact limit of the conclusion.}
The input is P-F31, theorem \texttt{actualcubicpole}, part two,
in \path{positivity_actual_cubic_pole.tex}; its family is
the Cartan multiplicity algebra of cubic $SL_n$ invariants.
Its variable $t$ records the grading in SGA 3, I, 4.7.3.1.
To reinterpret $f_n(r)$ as an extension-field trace in the
SGA 5, III, 4.7 framework requires a new identification of
that grading index with the power of Frobenius. If such an
identification is proposed using one finite cohomological
carrier, (G018.1) bounds its total rank. Neither a Tate twist,
an inserted correspondence nor virtual cancellation removes
the pole count; insertions and signs were already allowed.
It is essential to distinguish this from the Tate--Artin
realizations of individual Laurent-polynomial structure
constants in NN-G006--NN-G011, where the field size and
the basis-word or stretching parameter are separate variables.

The rank obstruction is unconditional relative to the cited
refutation and is unaffected by A or conditional M. It
rules out a polynomial-rank iteration carrier on this actual
plethysm family, but leaves short tensor descriptions,
varying graded carriers, succinct Hilbert rational functions,
and PH1 positive witness formulas open. No arithmetic-circuit
or Boolean lower bound follows from it alone.



\section{NN-G019: a finite optimum theorem for the integral-ratio dimension-ring class}
\paragraph{Precisely delimited class and statement.}
Consider dimension rings with the table shape (G011.2), but
allow the parameters to vary subject to the following conditions:
\[
\begin{gathered}
\delta_1=\delta_2=\delta\in\mathbb Z[D]
   \text{ is monic of degree two},\qquad \deg W\le2,\\
G_i=k_iW\delta,\qquad S=k_1k_2W\delta^2,
\qquad k_1,k_2\in\mathbb Z_{>0},\\
\Lambda_i\in\mathbb Z[D],\qquad
\Lambda_i\ge\lambda_i,\quad\Lambda_i\equiv\lambda_i\pmod2
\quad\text{coefficientwise in }q,
\end{gathered}                                              \tag{G019.1}
\]
where $D=q+2+q^{-1}$, $\lambda_1=D$ and
$\lambda_2=(D-2)^2$. As always, every table coefficient
must dominate the absolute value of the original coefficient
and have the same parity. There is no imposed degree or
coefficient cutoff on $\Lambda_i$ or on any other parameter.
Within this class the minimum number of cancelling pairs is
exactly $7276$. There are exactly two minimizing parameter
sets. Both have
\[
k_1=2,\quad k_2=1,\quad
\delta=D^2-3D+2,\quad
\Lambda_1=D,\quad\Lambda_2=3D^2-10D+4,
\]
and their two choices of $W$ are
\[
W=7D-8\qquad\text{or}\qquad W=4D^2-13D+8.
\tag{G019.2}
\]
The lower bound is a finite exact-arithmetic certificate
after the range reduction below. It is not a global optimum
over all dimension rings or over all polynomial block parameters.

\paragraph{Proof: algebraic constraints and finite bounds.}
The scalar-block identities force, with $\{i,j\}=\{1,2\}$,
\[
U_i=\frac{k_iW-D}{k_j},\qquad
V_i=\frac{k_iDW-\Lambda_i}{k_j}.                 \tag{G019.3}
\]
The coefficient constraints include
$W\ge D$, $\delta\ge |(D-1)(D-4)|$, $U_i\ge D$,
$V_i\ge K_i$, where $K_1=5D-4$ and
$K_2=2D^2-5D+4$, always with parity congruence.
Modulo two, the equation $k_2V_1=k_1DW-\Lambda_1$
therefore becomes
$k_2D=k_1D^2+D$. The Laurent polynomials $D$ and $D^2$
are linearly independent over $\mathbb F_2$, so
$k_1$ is even and $k_2$ is odd.

Put $w=W(4)$ and $d=\delta(4)$, which are the sums of
the Laurent coefficients. Let $T$ be the sum of all
nonunit dimension coefficients, evaluated at $q=1$.
The word-pair multiplicities computed in NN-G011 give
\[
\begin{split}
T={}&k_1k_2wd^2+20wd(k_1+k_2)+200w\\
 &+20\bigl(U_1(4)+U_2(4)+V_1(4)+V_2(4)\bigr)\\
 &+d\bigl(52+14\Lambda_1(4)+14\Lambda_2(4)\bigr).
\end{split}                                                \tag{G019.4}
\]
The cancelling-pair score is $(T-4424)/2$. The already
proved incumbent NN-G011 means that, to prove the optimum,
it suffices to examine $T\le18976$. Since $w,d\ge4$,
$U_i(4)\ge4$, $V_i(4)\ge16$, and $\Lambda_i(4)\ge4$,
\[
T\ge k_1k_2wd^2+20wd(k_1+k_2)+200w+800+164d.
\tag{G019.5}
\]
In particular $T\ge472w+1456$ and
$T\ge8d^2+404d+1600$. Both $w,d$ are even, so
$w\le36$ and $d\le26$. Also
$T\ge320(k_1+k_2)+2256$ gives $k_1+k_2\le52$.
These bounds are consequences of the desired score;
they are not experimental search cutoffs.

All possible $\delta$ and $W$ are now given by the finite
nonnegative integer parameters
\[
\begin{aligned}
\delta&=(q^2+q^{-2})+(1+2a)(q+q^{-1})+2b,
       &4+4a+2b&\le26,\\
W&=c_2(q^2+q^{-2})+c_1(q+q^{-1})+c_0,
       &c_0+2c_1+2c_2&\le36,
\end{aligned}
\]
where $a,b\ge0$, $c_0\ge2$ is even, $c_1\ge1$ is
odd, and $c_2\ge0$ is even. These formulas cover every
monic quadratic $\delta$ and every degree-at-most-two $W$
with their required Laurent coefficient inequalities and parity.

\paragraph{Proof: eliminate the unbounded $\Lambda_i$.}
Fix the bounded parameters and consider cell $i$.
Let $C_i=k_iDW-k_jK_i$. The requirement
$V_i\ge K_i$ with $V_i\equiv K_i\pmod2$ is equivalent to
\[
\Lambda_i\le C_i,\qquad
\Lambda_i\equiv C_i\pmod{2k_j}
\quad\text{coefficientwise}.
\]
For each exponent $e$, put $l_e=[q^e]\lambda_i$ and
$c_e=[q^e]C_i$. A possible coefficient must be at least
$l_e$, so its unique smallest value is
\[
L_e=l_e+((c_e-l_e)\bmod 2k_j),
\tag{G019.6}
\]
where the residue is in $\{0,\ldots,2k_j-1\}$.
There is no feasible value if the parity differs from
$l_e$ or if $L_e>c_e$. Otherwise the coefficients $L_e$
give a symmetric integral Laurent polynomial and hence an
element of $\mathbb Z[D]$. Coefficients outside the support
of $C_i$ and $\lambda_i$ must be zero; thus this step also
proves the required degree bound, instead of assuming one.

Increasing $\Lambda_i(4)$ while fixing the other bounded
parameters changes $T$ with strictly positive slope
$14d-20/k_j\ge36$. Consequently only the componentwise
smallest choice (G019.6) can minimize the score. This
eliminates all remaining unbounded choices.

\paragraph{Proof: the exact finite certificate.}
Enumerate the nonnegative integer tuples just displayed,
the positive even $k_1$ and positive odd $k_2$ with sum
at most $52$, and retain (G019.5)$\le18976$. There are
$1744$ such tuples. For each, divide (G019.3) coefficientwise
and test integrality, domination and parity for $U_i$;
test $G_i$ and $S$ against the true parameters; then apply
(G019.6). Exactly the following $21$ tuples survive those
tests, with $\Lambda_i$ chosen minimally:
\[
\begin{array}{rr|rr|rrr|r}
k_1&k_2&a&b&c_0&c_1&c_2&\text{pairs}\\\hline
2&1&0&1&6&3&4&7276\\
2&1&0&1&6&7&0&7276\\
2&3&0&0&10&5&0&7656\\
2&1&0&0&14&7&0&8288\\
2&1&0&1&10&3&4&9040\\
2&1&0&1&10&7&0&9040\\
2&1&0&2&6&3&4&9312\\
2&1&0&2&6&7&0&9312\\
2&1&1&0&6&3&4&9312\\
2&1&1&0&6&7&0&9312\\
2&1&0&0&10&7&4&9732\\
2&1&0&0&10&11&0&9732\\
2&1&0&0&18&3&4&9732\\
2&1&0&0&18&7&0&9732\\
2&1&0&0&6&11&4&11176\\
2&1&0&0&6&15&0&11176\\
2&1&0&0&14&3&8&11176\\
2&1&0&0&14&7&4&11176\\
2&1&0&0&14&11&0&11176\\
2&1&0&0&22&3&4&11176\\
2&1&0&0&22&7&0&11176
\end{array}
\]
This finite calculation is reproduced by
\path{verify_dimension_ratio_optimum.py}, whose complete
source is included. It uses only integer Laurent dictionaries
and imports no supplied search or table implementation. Its
JSON certificate lists every surviving tuple, both minimizing
tables, their actual $\Lambda_i$, and the bottom and cell
masses. A separate search using the supplied polynomial
arithmetic gives the same counts and minimum.

The first minimizing table is the new $W$ in (G019.2),
and the second is NN-G011. For both, (G019.6) gives the
displayed $\Lambda_i$. The table conditions were checked
coefficientwise, and associativity follows from the same
block embedding as NN-G011: the four images in each matrix
block are independent since $\Lambda_i(D^2-\Lambda_i)\ne0$,
and (G019.3) gives exactly the required scalar-block products.
Thus both minima are attained by actual dimension rings.
Strict positivity of the slopes above proves that no
nonminimal $\Lambda_i$ gives another minimizer. This proves
the stated optimum and the two-parameter-set classification.
\hfill$\square$

\paragraph{Geometric translation and smaller remaining gate.}
By NN-G010 both minimizing rings give additive monoidal
categories of kernels on a fixed ten-point scheme with
Tate--Artin coefficient objects. The score counts the added
regular $C_2$ pairs in the structure-constant objects,
coefficient by coefficient. The relevant actual-object and
split-$K_0$ distinction is SGA 3, I, 4.7.2 and archive cote 2,
scan 17; effectivity is not equality in a virtual representation
ring. This optimization selects no canonical geometric functor.

Any improvement below $7276$ must now leave at least one
specific part of (G019.1): use unequal cell $\delta_i$,
a nonmonic or higher-degree $\delta$, $\deg W>2$,
nonintegral or nonconstant ratios $G_i/(W\delta)$,
allow $\Lambda_i$ below the true $\lambda_i$ in some
coefficients while compensating elsewhere in the table,
or change the table shape. These are genuine remaining
possibilities. The result narrows the search gate but
does not prove global minimality, canonical positivity,
PH1, or anything about rank four by itself.



\section{NN-G020: normal toric rays with conductor at most two and hard first coefficient}
\paragraph{Statement and input representation.}
There is a polynomial-time construction from a $3$-CNF formula $F$
on $n\geq1$ declared variables and $m$ clauses to an integral weight
matrix $A_F$, a primitive character $b$, and the diagonal action of
the split torus $T=\mathbb G_m^{n+m}$ on
\[
 S=\mathbb Z[X_1,Y_1,\ldots,X_n,Y_n,Z_1,\ldots,Z_m].
\]
The ray algebra
\[
 R_F=\bigoplus_{k\geq0}S_{kb}
 \tag{G020.1}
\]
is a normal, finitely generated, flat $\mathbb Z$-algebra and is
exactly the invariant ring $S^{\ker\chi_b}$. Its graded pieces have
canonical monomial bases and satisfy
\[
 \operatorname{rank}_{\mathbb Z}(R_F)_1=\#F,
 \qquad
 \{k:(R_F)_k\ne0\}=
 \begin{cases}
  \mathbb N,&F\text{ is satisfiable},\\
  \{0,2,3,4,\ldots\},&F\text{ is unsatisfiable}.
 \end{cases}
 \tag{G020.2}
\]
Here $\mathbb N$ includes zero. In particular the \emph{actual ray
conductor}, defined as the least $c\geq0$ with every integer $k\geq c$
in the support, is either zero or two. A rational point of the
degree-one polytope and an integral point in every degree $k\geq2$
are explicitly known independently of satisfiability.

The input for the complexity statement is this polynomial-size
\emph{diagonal-torus presentation} $(A_F,b)$, or equivalently the
inequalities below. No polynomial bound on a displayed generating
set or relations for the invariant algebra is asserted. In this
presentation, deciding whether $(R_F)_1\ne0$ is NP-complete, and
computing its rank is \#P-complete. E gives the exact subexponential
counting bounds stated below, with the original number $n$ of
Boolean variables. A places the nontrivial zeros of the common
ambient arithmetic zeta on one shifted line, while that zeta
contains no information about (G020.2).

\paragraph{Construction and homogeneous lattice-point model.}
Write every clause as three literal positions; repeating a literal
pads a shorter clause without changing any assignment. For each
clause $C$, let
\[
 L_C(x,y)=\sum_{\text{positions in }C}
 \begin{cases}x_i,&\text{the position is the literal }x_i,\\
 y_i,&\text{the position is the literal }\neg x_i.
 \end{cases}
\]
Repeated occurrences are counted with multiplicity. Define $A_F$
by the equations
\[
 x_i+y_i=k\quad(1\leq i\leq n),\qquad
 L_C(x,y)-s_C=k\quad(1\leq C\leq m),
 \tag{G020.3}
\]
and put $b=(1,\ldots,1)\in\mathbb Z^{n+m}$. The columns of $A_F$
are the weights of $X_i,Y_i,Z_C$, in that order. Its entries have
absolute value at most three, and its construction is polynomial
in the formula encoding. In the group-scheme action, a monomial
with exponent vector $z=(x,y,s)$ has character $A_Fz$; this is a
statement about comodules, not merely about rational points of
the torus over a small finite field.

At $k=1$, nonnegative integral solutions of the first equations
have $x_i\in\{0,1\}$ and $y_i=1-x_i$. The remaining equations
uniquely force $s_C=L_C(x,y)-1$. This is nonnegative exactly
when clause $C$ is true. Hence the map from satisfying assignments
to exponent vectors is a bijection. It proves the first equality
in (G020.2), including multiplicities and declared variables that
do not occur in a clause.

Set
\[
 P_F=\{z\in\mathbb R_{\geq0}^{2n+m}:A_Fz=b\}.
\]
This is a bounded rational polytope: $0\leq x_i,y_i\leq1$ and
$0\leq s_C\leq2$. It always contains the explicitly specified
point $x_i=y_i=s_C=1/2$. Its dilations give the exact homogeneous
formula
\[
 \operatorname{rank}_{\mathbb Z}(R_F)_k
   =|kP_F\cap\mathbb Z^{2n+m}|\quad(k\geq0).
 \tag{G020.4}
\]
At $k=0$, the first equations force $x=y=0$ and then $s=0$,
so the formula holds there as well. In particular $b$ is in the
weight cone for every formula; cone feasibility is already settled.

For every integer $k\geq2$, choose
\[
 x_i=\lfloor k/2\rfloor,\qquad
 y_i=\lceil k/2\rceil,\qquad
 s_C=L_C(x,y)-k.
 \tag{G020.5}
\]
Every clause has three positions, so
$L_C(x,y)\geq3\lfloor k/2\rfloor\geq k$ for $k\geq2$.
Thus these are integral nonnegative solutions for every such
degree. Combining this with the degree-zero and degree-one
calculations proves the full support formula and its conductor.

\paragraph{Exact invariant-ring identification and normality.}
The linear functional $\beta$ on the character lattice takes the
value $3m+1$ on every variable-equation row and $-1$ on every
clause-equation row. Each column for $X_i$ or $Y_i$ therefore has
$\beta$-value $3m+1$ minus its total number of literal occurrences,
which is at least one. Each slack column has $\beta$-value one.
Moreover
\[
 \beta(b)=n(3m+1)-m>0.
 \tag{G020.6}
\]
Consequently no nonzero monomial has weight zero, and no monomial
has weight $kb$ for $k<0$. This also proves finite rank of every
graded piece directly from a positive bound on its exponents.

Since $b$ is primitive, it extends to an integral basis of the
character lattice $M=\mathbb Z^{n+m}$. Thus
$K=\ker(\chi_b:T\to\mathbb G_m)$ is a split connected torus and
its character lattice is $M/\mathbb Zb$. Restricting the monomial
grading to $K$ gives, over $\mathbb Z$ itself,
\[
 S^K=\bigoplus_{k\in\mathbb Z}S_{kb}
     =\bigoplus_{k\geq0}S_{kb}=R_F.
 \tag{G020.7}
\]
The second equality uses (G020.6); omitting that step would leave
the proposed ray/invariant-ring identification unjustified.

For completeness, finite generation is the elementary rational-cone
argument. Let $H$ be the set of nonnegative integral exponent vectors
whose weights lie in $\mathbb Zb$. It is the intersection of a
rational polyhedral cone with a lattice. Choose finitely many
integral vectors generating that cone over $\mathbb R_{\geq0}$.
For a lattice vector in the cone, subtract the integer parts of
its coefficients in a nonnegative linear combination of those
vectors. The remainder is a lattice vector in the bounded region
where all those coefficients lie in $[0,1)$. There are finitely
many such remainders. They, together with the chosen vectors,
generate $H$. Thus $R_F=\mathbb Z[H]$ is finitely generated.
This is Gordan's lemma in its lattice form; see Cox--Little--Schenck,
\emph{Toric Varieties}, Proposition 1.2.17, printed p.~30.

To prove normality without an unmentioned quotient hypothesis,
let $f\in\operatorname{Frac}(S^K)$ be integral over $S^K$.
Its monic integral equation also has coefficients in $S$, so it
is integral over $S$. The polynomial ring $S$ over $\mathbb Z$
is integrally closed, hence $f\in S$. As a ratio of two invariant
elements, $f$ is invariant as a rational function, and therefore
$f\in S^K$. This proves that $R_F$ is integrally closed. Finally,
its monomial basis makes it a free $\mathbb Z$-module, so it is
flat. Normality of this algebra is entirely consistent with the
single missing degree one: saturation of its exponent semigroup
does not mean saturation of its image under the degree map.

\paragraph{Complexity consequences, with E's original parameter.}
An exponent vector of degree one has a polynomial-size encoding,
and (G020.3) verifies it in polynomial time. The established
bijection is a parsimonious reduction from \#3SAT and its
positivity version. Thus the stated NP- and \#P-completeness follow
from the classical verifier encoding, which can also be made
explicit here. Compile a polynomial-time Boolean verifier into a
polynomial-size circuit with AND and NOT gates. For each AND gate
$z=x\wedge y$ impose
$(\neg z\vee x)$, $(\neg z\vee y)$, and
$(z\vee\neg x\vee\neg y)$; for each NOT gate $z=\neg x$ impose
$(z\vee x)$ and $(\neg z\vee\neg x)$. Impose the output bit by
a unit clause, and repeat literals to pad each clause to three
positions. For every input witness, all gate bits have exactly
one consistent extension, accepted precisely when the verifier
accepts. Thus counting these assignments is parsimoniously as
hard as counting the witnesses of any \#P verifier, and the same
construction proves NP-hardness of positivity. Membership in
\#P and NP uses the degree-one witness already given. This proof
does not bound the new formula's variable count by the number
of input witness bits alone; its gate bits are counted as well.

More concretely, a deterministic polynomial-time positivity algorithm
on these presented normal toric rays would imply $P=NP$.
Polynomial-size certificates for their degree-one vanishing,
checkable in deterministic polynomial time and complete on all
these inputs, would imply $NP=coNP$: the same certificates would
put $3$-UNSAT in NP. These are implications, not assertions of
the separations $P\ne NP$ or $NP\ne coNP$.

Applying the designated E Theorem 1.1 directly to $F$, not to an
inflated auxiliary-variable formula, computes the degree-one
rank in
\[
 (n+m)^{O(1)}2^{O(n\log\log n/\sqrt{\log n})}
 \quad\text{bit operations and polynomial space}.
 \tag{G020.8}
\]
For fixed $D$ and $m\leq Dn$ the exponent improves to
$O_D(n/\sqrt{\log n})$, with polynomial prefactor and space.
The theorem is stated for $n\geq20$; the finitely many smaller
values are handled by direct enumeration. There is no hidden
replacement of $n$ by $2n+m$ in this use of E.

\paragraph{Geometric dictionary and the exact failure of zeta recovery.}
The torus action and its multigraded comodule are the objects of
SGA 3, I, 4.7.2--4.7.3 and Corollary 4.7.3.1. Equation (G020.7)
is an actual affine quotient by the indicated split subtorus;
the finite generation and normality proofs above make its scheme
meaning explicit. The ray grading also gives $\operatorname{Proj}R_F$
and its associated graded-module sheaves in the sense of EGA II,
2.4--2.5, especially 2.4.7 and 2.5.2. No identification with a
plethysm or Kronecker multiplicity algebra is asserted.

Put $N=2n+m$. The underlying \emph{ambient} scheme
$\operatorname{Spec}S=\mathbb A^N_{\mathbb Z}$ is independent of
the formula once $n,m$ are fixed. For every prime power $q$ and
integer $r\geq1$ it has $q^{Nr}$ points over $\mathbb F_{q^r}$.
Its only nonzero compact-support $\ell$-adic cohomology group
is $H_c^{2N}=\mathbb Q_\ell(-N)$, and the connected torus action
induces the trivial action on that group. The precise trace and
connected-action locator is SGA 4, XVIII, 2.8.1 and the immediately
following paragraph; the affine-space calculation also follows
from XVII, 5.4.3 and the curve trace calculation XVIII, 1.1.6--1.1.9.
Consequently
\[
 Z(\mathbb A^N_{\mathbb F_q},T)=(1-q^NT)^{-1},\qquad
 \zeta_{\mathbb A^N_{\mathbb Z}}(s)=\zeta(s-N).
 \tag{G020.9}
\]
Under A its nontrivial zeros lie on $\Re s=N+1/2$. These common
point counts and this rank-one pure Weil object do not recover
the degree-one rank, despite the normal ray algebra, known rational
half-point, positive monomial rule and conductor bound two.
For an explicit smallest comparison, the formulas
$(x\vee x\vee x)\wedge(x\vee x\vee x)$ and
$(x\vee x\vee x)\wedge(\neg x\vee\neg x\vee\neg x)$
have degree-one ranks one and zero and the same ambient
$\mathbb A^4_{\mathbb Z}$.

This last assertion concerns the ungraded ambient realization;
no formula-independent zeta or cohomology is claimed for
$\operatorname{Spec}R_F$ itself, or for grade-indexed trace objects.
Likewise the conductor here is an actual support conductor,
not GCT VI's strict off-residue saturation index. The construction
does not contradict a conjecture confined to actual plethysm or
Kronecker families. It proves that normality, a homogeneous
lattice-point model, known moment-cone membership, a uniformly
tiny actual conductor, and a finite pure ambient realization
still leave an NP-hard first-coefficient gate in this explicit
geometric family. Neither A nor the unrelated conditional M
removes that gate; E supplies precisely (G020.8).
\hfill$\square$



\section{NN-G021: signed auxiliary parameters lower the exact rank-three optimum}
\paragraph{Statement and the larger class.}
The auxiliary cell parameter $\Lambda_i$ is not itself a structure
constant of the rank-three table. Requiring it to dominate the
corresponding trace parameter is therefore an extra restriction.
Removing that restriction gives an associative effective dimension
ring with
\[
 7064\text{ cancelling pairs in total},\qquad
 6652\text{ in the bottom cell}.                         \tag{G021.1}
\]
An explicit choice keeps $\delta,W,G_1,G_2,S,U_1,U_2$ from NN-G011,
and replaces the four remaining parameters by
\[
 \Lambda_1=D-2,\quad \Lambda_2=3D^2-10D,\qquad
 V_1=14D^2-17D+2,\quad V_2=2D^2+D.                      \tag{G021.2}
\]
Although $\Lambda_2$ has a negative constant Laurent coefficient,
all actual multiplication coefficients are effective.

Moreover $7064$ is the exact minimum in the following class:
the same closed table shape, a common monic quadratic
$\delta\in\mathbb Z[D]$, a polynomial $W\in\mathbb Z[D]$ of
degree at most two, and arbitrary \emph{positive rational constant}
ratios $G_i/(W\delta)$; the $\Lambda_i$ may be arbitrary signed
polynomials in $\mathbb Z[D]$. Only the actual table coefficients
must be integral, dominate the absolute trace coefficients, and
have the required parity. There are eight minimizing parameter
sets in this class, described below. No optimum over unrestricted
dimension rings is claimed.

\paragraph{The explicit improvement.}
Use Laurent coefficient lists indexed by nonnegative exponents,
with negative exponents supplied by symmetry. Thus
\[
 \delta=(2,1,1),\quad D-2=(0,1),\quad
 3D^2-10D=(-2,2,3).
\]
The actual cell products are
\[
 \delta\Lambda_1=(2,3,1,1),\qquad
 \delta\Lambda_2=(6,7,6,5,3).                           \tag{G021.3}
\]
The absolute trace targets for these products are, respectively,
\[
 X_1=(2,1,1,1),\qquad X_2=(2,3,2,1,1).                 \tag{G021.4}
\]
They are dominated with the correct parity. The unchanged $G_i,S,W,U_i$
already pass their coefficient conditions by NN-G011. The new
$V_1,V_2$ equal the previous ones plus two and still dominate their
trace targets with unchanged parity. The closed table identities
remain valid because $V_i=(k_iDW-\Lambda_i)/k_j$ with
$(k_1,k_2)=(2,1)$, so decreasing $\Lambda_1$ by two and
$\Lambda_2$ by four increases both $V_i$ by two.

All the table's within-cell coefficients are $\delta D^\alpha$
or $\delta D^\alpha\Lambda_i$; its bottom coefficients are
nonnegative powers of $D$ times the parameters just checked.
Thus these base checks prove coefficientwise domination, parity
and effectivity everywhere. Associativity follows from the same
scalar-block and two-matrix-block embedding used in NN-G011;
the conditions $\Lambda_i(D^2-\Lambda_i)\ne0$ hold for (G021.2).
Alternatively the included verifier checks all $1000$ triples
directly with integer Laurent dictionaries.

The total nonunit coefficient mass is now $18552$, divided into
bottom mass $17400$ and cell mass $1152$. The true absolute masses
are $4096$ and $328$. Consequently the cancelling-pair counts are
$(18552-4424)/2=7064$ and $(17400-4096)/2=6652$.
The bottom count rises by forty while the total count falls by
$212$: forcing positivity on an auxiliary parameter had concealed
a less expensive effective multiplication table.

\paragraph{Why the larger optimization class has a finite bound.}
Write $W=gW_0$, where $g>0$ is its coefficient content and $W_0$
is primitive. Positivity and $W\equiv D\pmod2$ imply that $g$ is
odd and that
\[
 W_0=c_0+c_1(q+q^{-1})+c_2(q^2+q^{-2}),
 \quad c_0\ge2\text{ even},\ c_1\ge1\text{ odd},\
 c_2\ge0\text{ even},\ \gcd(c_0,c_1,c_2)=1.
 \tag{G021.5}
\]
The common quadratic has the exact form
\[
 \delta=2b+(1+2a)(q+q^{-1})+(q^2+q^{-2}),\qquad a,b\ge0.
 \tag{G021.6}
\]
Set $d=\delta(4)=4+4a+2b$ and $w_0=W_0(4)=c_0+2c_1+2c_2$;
evaluation at $D=4$ is evaluation at $q=1$.

The polynomial $q^2\delta$ is monic and primitive. Gauss's content
lemma shows that $W_0\delta$ and $W_0\delta^2$, after harmless
Laurent shifts, are primitive. Integrality of $G_i=k_iW\delta$
therefore forces $a_i:=gk_i$ to be a positive integer. Integrality
of $S=G_1G_2/W$ forces $g\mid a_1a_2$. The trace congruences
for $G_1,G_2$ then force $a_1$ even and $a_2$ odd: modulo two,
$W_0\delta=D^2(D-1)$ is nonzero, the first trace target is zero,
and the second equals that product. Thus
\[
 \begin{aligned}
 G_i&=a_iW_0\delta,\\
 S&=(a_1a_2/g)W_0\delta^2,\\
 U_i&=g(a_iW_0-D)/a_j,\\
 V_i&=g(a_iDW_0-\Lambda_i)/a_j\quad(j\ne i),
 \end{aligned}                                         \tag{G021.7}
\]
where $a_1a_2/g$ is a positive even integer. This accounts for
all rational constant ratios in the stated class, including
nonintegral ones.

Let $T$ be the total nonunit dimension mass. The exact table
count used in NN-G019 remains valid even for signed $\Lambda_i$:
\[
 T=S(4)+20(G_1(4)+G_2(4))+200gw_0
   +20\sum_i(U_i(4)+V_i(4))
   +52d+14d\sum_i\Lambda_i(4).
 \tag{G021.8}
\]
Since $\delta\Lambda_i\ge X_i$ and the masses of (G021.4)
are eight and sixteen, the last term is at least $336$.
Also $\sum_i(U_i(4)+V_i(4))\ge40$. Any score at most the
already proved $7276$ therefore satisfies $T\le18976$ and
\[
 (a_1a_2/g)w_0d^2+20w_0d(a_1+a_2)
       +200gw_0+1136+52d\le18976.                     \tag{G021.9}
\]
This implies the complete finite ranges
\[
 w_0\le36,\quad d\le30,\quad a_1+a_2\le52,\quad g\le23.
 \tag{G021.10}
\]
For the first bound the left side is at least $472w_0+1344$,
and $w_0$ is even. For the second it is at least
$8d^2+292d+1936$, which exceeds $18976$ at $d=32$.
For the third it is at least $320(a_1+a_2)+2272$.
The last follows already from $200gw_0\ge800g$.
All four lower bounds use only positive terms and the minimum
values $d,w_0\ge4$, $a_1+a_2\ge3$ and $a_1a_2/g\ge2$.

\paragraph{Exact minimization over every signed $\Lambda_i$.}
No degree restriction on $\Lambda_i$ has been hidden. If its
degree exceeded three, its highest Laurent coefficient would
occur with the same sign at the top of $\delta\Lambda_i$ and
with the opposite sign at the top of $V_i$, since
$\deg(DW_0)\le3$. Both cannot be effective. Hence
$\deg\Lambda_i\le3$.

For each $i$, put $a_j'=a_j/\gcd(g,a_j)$. Integrality and parity
of $V_i$ are exactly the coefficientwise congruences
\[
 \Lambda_i\equiv a_iDW_0-a_j'K_i\pmod{2a_j'},\qquad
 K_1=5D-4,\quad K_2=2D^2-5D+4,                      \tag{G021.11}
\]
and domination of $K_i$ gives the upper bounds
\[
 (\Lambda_i)_e\le
 \left\lfloor\frac{g a_i(DW_0)_e-a_j(K_i)_e}{g}\right\rfloor.
 \tag{G021.12}
\]
To check (G021.11), divide the original congruence
$g(a_iDW_0-\Lambda_i)\equiv a_jK_i\pmod{2a_j}$ by
$\gcd(g,a_j)$. The remaining coefficient of $\Lambda_i$
is odd and coprime to the modulus; its inverse is odd, so its
product with $a_j'$ equals $a_j'$ modulo $2a_j'$.

Write $h=2b$, $u=1+2a$ and
$\Lambda_i=L_0+\sum_{e=1}^3L_e(q^e+q^{-e})$.
The requirement $\delta\Lambda_i\ge X_i$, coefficientwise,
is the following complete list, with absent $X_i$ coefficients
interpreted as zero:
\[
 \begin{array}{rcl}
 L_3&\ge&0,\\
 L_2+uL_3&\ge&(X_i)_4,\\
 L_1+uL_2+hL_3&\ge&(X_i)_3,\\
 L_0+uL_1+hL_2+uL_3&\ge&(X_i)_2,\\
 uL_0+(h+1)L_1+uL_2+L_3&\ge&(X_i)_1,\\
 hL_0+2uL_1+2L_2&\ge&(X_i)_0.
 \end{array}                                           \tag{G021.13}
\]
The congruences give $\Lambda_1\equiv D$ and
$\Lambda_2\equiv D^2\pmod2$, so they enforce the required
product parity as well. Conversely that parity forces these
congruences modulo two since $\delta\not\equiv0\pmod2$.
These inequalities bound $L_3$,
then $L_2$, then $L_1$, and finally $L_0$ from below;
(G021.12) bounds each from above. When $h=0$, the last
inequality is simply tested without dividing by $h$.

For fixed $L_3,L_2,L_1$, the coefficient of $L_0$ in the score
mass is $14d-20g/a_j$. Therefore only the least allowed $L_0$
is needed when that coefficient is positive, and only the
greatest when it is negative. If it is zero every allowed
$L_0$ has the same score. This endpoint rule is used only
after all six inequalities and the congruences are imposed.
It is not an assumption that the auxiliary polynomial is
coefficientwise nonnegative.

The complete finite enumeration of (G021.5)--(G021.13),
implemented in \path{search_dimension_signed_lambda.py},
has $2572$ tuples after (G021.9), $99$ after the structural
$U,G,S$ tests, and $33636$ checked $\Lambda_i$ endpoints.
Twenty-eight structural tuples admit feasible cell parameters.
Their minimum scores are displayed in the included JSON
certificate; the only ones at or below $7276$ are
\[
 \begin{array}{c|c|c|c|c|c}
 (a_1,a_2,g)&(a,b)&(c_0,c_1,c_2)&\text{score}
       &\#\Lambda_1\text{ choices}&\#\Lambda_2\text{ choices}\\\hline
 (2,1,1)&(0,1)&(6,3,4)&7064&2&3\\
 (2,1,1)&(0,1)&(6,7,0)&7064&1&2
 \end{array}                                           \tag{G021.14}
\]
For either row the first available $\Lambda_1$ is $D-2$.
For the quadratic-$W$ row only, there is also the symmetric
Laurent list $(4,-1,-2,2)$. Both rows allow
$\Lambda_2=(6,-2,3,0)$ or $(-2,2,3,0)$; the quadratic row
also allows $(6,-2,-1,4)$. These yield $2\cdot3+1\cdot2=8$
parameter sets. At every minimum the two slopes are strictly
positive (their numerators in (G021.7) are $64$ and $148$),
so the endpoint rule loses no additional minimizing choices.
The arithmetic source and complete survivor data are included;
the calculation imports no supplied search implementation.

All eight displayed parameter sets have nonzero
$\Lambda_i(D^2-\Lambda_i)$, so the block embedding proves
associativity for their tables. Each dominates the trace
coefficients by the exact tests above. Thus the computed
minimum is attained by actual associative dimension rings.
The explicit table (G021.2) is independently checked by
\path{verify_dimension_7064.py}, which compares all $132$
positions with the supplied trace constants and checks every
one of the $1000$ associativity triples. This proves both
the improvement and the exact optimum in the stated class.
\hfill$\square$

\paragraph{Geometric meaning and the remaining gate.}
NN-G010 turns each of these effective multiplication tables
into the same type of additive monoidal kernel category on
a ten-point scheme, with effective Artin--Tate multiplicity
objects. Only the actual structure constants specify those
objects. An auxiliary factor such as $\Lambda_2$ need not
be the class of an effective object even when its product
with $\delta$ is. This is the precise distinction between
an effective tensor structure and a chosen polynomial
factorization of its coefficients; no cancellation history
can be read from a supported-zero tag alone.
The comodule dictionary is SGA 3, I, 4.7.2; kernel products
are the construction recorded in archive cote 2, scan 17.
Neither source canonically selects these tables.

The class bound now allows signed $\Lambda_i$ and rational
constant ratios. A score below $7064$ must leave another
specified restriction: a common monic quadratic $\delta$,
degree at most two for $W$, constant ratios $G_i/(W\delta)$,
or the closed table shape. This is a smaller global
optimization gate than NN-G019. It remains a question
about noncanonical realizations, with no consequence for
PH0, PH1 or a rank-four canonical basis.



\section{NN-G022: unequal monic quadratic cells force at least 7958 pairs}
\paragraph{Statement.}
Keep the closed rank-three table shape, but allow two distinct monic
quadratic cell parameters $\delta_1,\delta_2\in\mathbb Z[D]$.
Let $W\in\mathbb Z[D]$ have \emph{arbitrary degree}, and suppose
\[
 G_i=k_iW\delta_i\quad(k_i\in\mathbb Q_{>0}),\qquad
 S=G_1G_2/W.
 \tag{G022.1}
\]
Allow arbitrary signed auxiliary $\Lambda_i$. If all actual
structure constants are integral Laurent polynomials, dominate
their absolute trace targets, and have the trace parity, then
the table has at least $7958$ cancelling pairs.

In particular no table of this type with unequal monic quadratic
cells can improve NN-G021's $7064$ construction, even by raising
the degree of $W$. Combined with NN-G021, the exact minimum is
$7064$ when both cells are allowed to be arbitrary monic quadratics,
$\deg W\le2$, and the ratios in (G022.1) are positive rational
constants. This is an enlargement of that result, not an assertion
about unrestricted dimension rings.

\paragraph{Integral normalization and the parity constraints.}
Every bar-invariant integer Laurent polynomial is in $\mathbb Z[D]$:
write it in the basis $1,q^j+q^{-j}$ and use the recurrence
$C_0=2$, $C_1=D-2$, $C_{j+1}=(D-2)C_j-C_{j-1}$.
Thus an actual table coefficient that is a rational function of
$D$ is an integer polynomial in $D$ as soon as it is an integral
Laurent polynomial. In particular the two $U_i$ below belong
to $\mathbb Z[D]$.

Write $W=gW_0$ with positive integer coefficient content $g$
and primitive $W_0$. The direct mixed-cell coefficient gives
$W\ge D$ and $W\equiv D\pmod2$, so $g$ is odd,
$W_0\in\mathbb Z[D]$, and $W_0\equiv D\pmod2$.
Each $q^2\delta_i$ is monic and primitive. Gauss's content lemma
in $\mathbb Z[q^{\pm1}]$ therefore gives integers
\[
 a_i=gk_i>0,\qquad g\mid a_1a_2.
\]
The actual coefficient $\delta_i$ has the same parity as
$\delta_{\rm tr}=(D-1)(D-4)$, namely $D(D-1)$ modulo two.
Consequently $W_0\delta_i\equiv D^2(D-1)\ne0$ in
$\mathbb F_2[D]$. The first $G$ trace is $4D(D-1)$ and the
second is $D^2(D-1)$. Their parities force
\[
 a_1\text{ even},\qquad a_2\text{ odd},\qquad
 a_1a_2/g\text{ a positive even integer}.             \tag{G022.2}
\]
The closed table identities now read
\[
 \begin{aligned}
 G_i&=a_iW_0\delta_i,&
 S&=(a_1a_2/g)W_0\delta_1\delta_2,\\
 U_i&=\frac{g\delta_i(a_iW_0-D)}{a_j\delta_j},&
 V_i&=\frac{g\delta_i(a_iDW_0-\Lambda_i)}{a_j\delta_j}
       \quad(j\ne i).
 \end{aligned}                                        \tag{G022.3}
\]
In particular $U_i\ge D$ and $U_i\equiv D\pmod2$.

\paragraph{The two quadratics must share an odd integer root.}
If $\delta_1$ and $\delta_2$ were coprime in $\mathbb Q[D]$,
polynomiality of $U_1$ in (G022.3) would imply
$\delta_2\mid a_1W_0-D$. Division by a monic integer polynomial
has integer quotient and remainder, so this divisibility holds
in $\mathbb Z[D]$. Reducing the quotient identity modulo two
would say that the degree-two monic polynomial $D(D-1)$ divides
$D$ in $\mathbb F_2[D]$. This is impossible. Notice that this
argument places no bound on $\deg W$.

Distinct monic quadratics sharing a root over $\mathbb Q$ have
a common linear factor; its root is rational and hence an integer.
We may therefore write
\[
 \delta_i=(D-r)(D-s_i),\qquad r,s_1,s_2\in\mathbb Z,
 \qquad s_1\ne s_2.
 \tag{G022.4}
\]
After cancelling $D-r$ in (G022.3), the two remaining linear
factors are coprime. Polynomiality of $U_1,U_2$ gives
\[
 a_1W_0(s_2)=s_2,\qquad a_2W_0(s_1)=s_1.             \tag{G022.5}
\]
The two roots of each $\delta_i$ have opposite parity, since
its reduction is $D(D-1)$. If $r$ were even, $s_2$ would be odd,
contradicting the first equality of (G022.5), whose left side
is even. Thus $r$ is odd and both $s_i$ are even.

The coefficient at $q$ in (G022.4) is $4-r-s_i$ and the
constant Laurent coefficient is $6-2(r+s_i)+rs_i$.
Their domination conditions give
\[
 r+s_i\le3,\qquad 6-2(r+s_i)+rs_i\ge0.              \tag{G022.6}
\]
Also $\delta_i(4)=(4-r)(4-s_i)>0$. Neither root equals four;
they must both be below four, because being both above four
would contradict their sum bound. Hence $r\le3$ and $s_i\le2$.
If $r=3$, the constant coefficient in (G022.6) equals $s_i$,
while the sum bound gives $s_i\le0$. Thus both $s_i$ would
be zero, contrary to their distinctness. We have proved
\[
 r\le1,\qquad r\text{ odd},\qquad s_i\le2\text{ even}.
 \tag{G022.7}
\]

\paragraph{All exceptional small roots are excluded.}
First suppose $s_2=0$. Then the first equation of (G022.5)
gives $W_0(0)=0$, so write $W_0=DH$ with $H\in\mathbb Z[D]$.
Since $s_1\ne0$, the second equation gives $a_2H(s_1)=1$.
Therefore $a_2=1$, $H(s_1)=1$, and $H-1=(D-s_1)Q$ for an
integer polynomial $Q$. The second $U$ coefficient is then
\[
 U_2=(g/a_1)D^2Q(D).
\]
As $U_2\in\mathbb Z[D]$, it is divisible by $D^2$ in that
ring, regardless of the rational scalar displayed in this
expression. Its parity therefore cannot equal $D$. This
contradicts the required parity of $U_2$, and excludes $s_2=0$.

For any even integer $s$, the congruence $W_0\equiv D\pmod2$
implies that $W_0(s)$ is even. Thus (G022.5) says that the
nonzero even integer $s_2$ is divisible by $2a_1\ge4$.
Together with $s_2\le2$, this proves $s_2\le-4$.

If $s_1=0$, then $W_0=DH$ as before. Since $s_2\ne0$,
the first equation of (G022.5) would give $a_1H(s_2)=1$,
impossible because $a_1$ is a positive even integer.
If instead $s_1=2$, the second equation and oddness of $a_2$
give $a_2=1$ and $W_0(2)=2$. Put $s_2=-a_1t$; here
$t=-W_0(s_2)>0$ is an even integer. An integer polynomial's
value difference $W_0(s_2)-W_0(2)=-t-2$ must be divisible by
$s_2-2=-a_1t-2$. But
\[
 0<t+2<a_1t+2,
\]
which makes that divisibility impossible. These cases exhaust
the even possibilities above $-2$. We have proved
\[
 s_1\le-2,\qquad s_2\le-4,\qquad
 d_1:=\delta_1(4)\ge18,\quad d_2:=\delta_2(4)\ge24.
 \tag{G022.8}
\]

\paragraph{The primitive gluing polynomial has mass at least eight.}
Let $w_0=W_0(4)$, the sum of its Laurent coefficients. The
conditions $W_0\ge0$, $gW_0\ge D$ and
$W_0\equiv D\pmod2$ show that $w_0$ is even and at least four.
If it is four, $W_0=D$; if it is six, $W_0=D+2$.
Indeed every additional noncentral symmetric pair has even
coefficients and contributes at least four to the mass.

The first possibility contradicts
$a_1W_0(s_2)=s_2$ because $a_1\ge2$ and $s_2\ne0$.
In the second, $a_2(s_1+2)=s_1$. The case $a_2=1$ is
impossible. For odd $a_2\ge3$ this requires
$s_1=-2a_2/(a_2-1)\in[-3,-2)$, which contains no even integer.
Thus
\[
 w_0\ge8.                                             \tag{G022.9}
\]
This argument again allows $W$ of arbitrary degree.

\paragraph{The score lower bound.}
The exact nonunit coefficient mass for unequal cells is
\[
 T=S(4)+20(G_1(4)+G_2(4))+200W(4)
 +20\sum_i(U_i(4)+V_i(4))
 +26(d_1+d_2)+14\sum_i(\delta_i\Lambda_i)(4).
\]
The base table targets give $U_i(4)\ge4$, $V_i(4)\ge16$
and $(\delta_1\Lambda_1)(4)\ge8$,
$(\delta_2\Lambda_2)(4)\ge16$, as in (G021.4).
Using (G022.2), (G022.8) and (G022.9), all terms are bounded
below by
\[
 \begin{aligned}
 T&\ge2w_0d_1d_2+20w_0(2d_1+d_2)+200w_0
             +26(d_1+d_2)+1136\\
  &\ge2\cdot8\cdot18\cdot24
       +20\cdot8(2\cdot18+24)+200\cdot8
       +26(18+24)+1136\\
  &=20340.
 \end{aligned}                                        \tag{G022.10}
\]
Subtracting the true absolute nonunit mass $4424$ and dividing
by two gives at least $7958$ cancelling pairs. This lower bound
does not assume $\Lambda_i$ effective, make an optimization-window
assumption, or use a numerical search. It proves the statement.
\hfill$\square$

\paragraph{Geometric interpretation and the reduced remaining gate.}
This is a rigidity consequence of the parity of actual effective
Artin--Tate multiplicity objects and integrality of the closed
kernel-composition table. The equality modulo two is the condition
that dimension and signed trace arise from one representation,
as in NN-G007 and SGA 3, I, 4.7.2. The integer roots here are
roots of auxiliary parameter polynomials; they are not asserted
to be Frobenius eigenvalues or points of an arithmetic scheme.
Archive cote 2, scan 17 supplies the honest kernel-product
dictionary through NN-G010, not an arithmetic interpretation
of those roots.

Within monic quadratic cells and positive rational constant
ratios, any realization competitive with $7064$ must have
$\delta_1=\delta_2$, even for higher-degree $W$. The remaining
optimization possibilities are higher-degree $W$ in that common
quadratic class, nonmonic or nonquadratic cells, nonconstant
ratios, or a different table shape. This improves the exact
gate without introducing a canonical construction or a
nonstandard purity theorem.



\section{NN-G023: the exact quadratic-cell optimum with arbitrary gluing degree}
\paragraph{Statement.}
The minimum cancelling-pair score is $7064$ in the following
entire class: the closed rank-three table shape; monic quadratic
$\delta_1,\delta_2\in\mathbb Z[D]$; arbitrary-degree
$W\in\mathbb Z[D]$; positive rational constant ratios
$G_i/(W\delta_i)$; and arbitrary rational functions
$\Lambda_i\in\mathbb Q(D)$. All actual multiplication
coefficients must be effective, dominate their absolute trace
targets, and have the correct parity.

Thus the degree restriction on $W$ in NN-G021 is removed.
The proof does not infer any bound on Laurent support locations
from a bound on coefficient mass. It leaves those locations
arbitrary and eliminates the only three potentially better
families uniformly in the location. Classification of all
minimizers in this enlarged class is not asserted.

\paragraph{Reduction to a finite mass-and-congruence relaxation.}
NN-G022 excludes unequal quadratic cells from any score below
$7064$, so put $\delta_1=\delta_2=\delta$. Normalize
$W=gW_0$ and $G_i=a_iW_0\delta$ exactly as in NN-G021.
Integrality of the auxiliary parameters is now automatic.
Each $V_i$ is an actual integral Laurent polynomial and
is a rational function of $D$, so it is in $\mathbb Z[D]$
by the bar-invariance argument in NN-G022. The identity
$\Lambda_i=k_iDW-k_jV_i$ first puts $\Lambda_i$ in
$\mathbb Q[D]$. Its actual cell coefficient
$\delta\Lambda_i$ is in $\mathbb Z[D]$; division by the
monic polynomial $\delta$ then puts $\Lambda_i$ in
$\mathbb Z[D]$. Thus rational auxiliary factors create
no additional case. NN-G022's unequal-cell bound never
used integrality or polynomiality of $\Lambda_i$.

The content and parity arguments there do not use $\deg W\le2$:
$g$ is odd, $a_1$ is positive even, $a_2$ positive odd,
and $g\mid a_1a_2$. Similarly its mass bound (G021.9)
does not use that degree hypothesis. Thus every proposed
improvement still satisfies
\[
 d:=\delta(4)\le30,\quad w:=W_0(4)\le36,
 \quad a_1+a_2\le52,\quad g\le23.
 \tag{G023.1}
\]
The quadratic is (G021.6), with $d=4+4a+2b$.

Write the primitive gluing polynomial as
\[
 W_0=c_0+c_1(q+q^{-1})+
          \sum_{j\ge2}c_j(q^j+q^{-j}),
\]
where $c_0\ge2$ is even, $c_1\ge1$ is odd and the other
$c_j$ are nonnegative even integers of finite support.
The two conditions $U_i\in D+2\mathbb N[q^{\pm1}]$ imply
\[
 M\mid c_j\quad(j\ge2),\qquad
 M=\operatorname{lcm}\left(
 \frac{2a_2}{\gcd(ga_1,2a_2)},
 \frac{2a_1}{\gcd(ga_2,2a_1)}\right).
 \tag{G023.2}
\]
Indeed the coefficient at such an exponent in $U_i$ is
$ga_ic_j/a_j$, which must be an even nonnegative integer.
The second factor in this least common multiple retains
the full two-primary part of $2a_1$, so $4\mid M$.
Put $h=\sum_{j\ge2}c_j/M$. Then
\[
 w=c_0+2c_1+2Mh.                                     \tag{G023.3}
\]
This bounds the \emph{total coefficient mass} in the outer
positions; it does not bound the indices $j$.

The coefficient at zero and the coefficient at one of each
$U_i$ must also have their exact parity and domination.
Explicitly, for $(c,t)=(c_0,2)$ and $(c_1,1)$ the two numbers
\[
 g(a_ic-t)/a_j\quad(j\ne i)
 \tag{G023.4}
\]
must be integers at least $t$ and congruent to $t$ modulo two.
Primitivity requires $\gcd(c_0,c_1)=1$ when $h=0$, and
requires $\gcd(c_0,c_1,M)=1$ when $h>0$. The latter is
used only as a necessary condition.

For an actual table set $\ell_i=\Lambda_i(4)$. Its two
effective cell products give $d\ell_1\ge8$ and
$d\ell_2\ge16$. Evaluating the $V_i$ congruences and
domination gives, with $a_j'=a_j/\gcd(g,a_j)$,
\[
 \begin{gathered}
 \ell_i\equiv4a_iw\pmod{2a_j'},\\
 \left\lceil X_i(4)/d\right\rceil\le\ell_i\le
 \left\lfloor\frac{4ga_iw-16a_j}{g}\right\rfloor,
 \qquad (X_1(4),X_2(4))=(8,16).
 \end{gathered}                                       \tag{G023.5}
\]
No degree bound on $\Lambda_i$ is used here.
The exact mass expression is (G021.8). For fixed other
parameters it is affine in $\ell_i$, with coefficient
$14d-20g/a_j$. Minimize it over the one-dimensional
congruence interval (G023.5) by taking the first allowed
integer if that coefficient is nonnegative, and the last
if it is negative. This gives a necessary lower bound on
the score, not a claim that the chosen masses arise from
actual Laurent polynomials.

All the remaining ranges are finite: positive even
$a_1\le50$, positive odd $a_2\le51$ with sum at most $52$;
odd $g\le23$ dividing $a_1a_2$; $a,b\ge0$ with
$4+4a+2b\le30$; even $2\le c_0\le36$; odd
$1\le c_1\le17$; and $0\le h\le(36-c_0-2c_1)/(2M)$.
In addition impose $G_i(4)\ge48$ and $S(4)\ge576$,
which are necessary actual coefficient-mass inequalities.
These are precisely the loops and tests in
\path{scalar_dimension_bound.py}; it uses only integer
arithmetic and records the whole relaxation.

The complete calculation gives $42$ rows whose relaxed
score is strictly below $7064$. Thirty-nine have $h=0$,
so $\deg W\le1$ and NN-G021 already excludes an actual
score below $7064$. Exactly three rows have $h>0$:
\[
 \begin{array}{c|c|c|c|c|c}
 (a_1,a_2,g)&(a,b)&(c_0,c_1,M,h)&w&d&
       \text{relaxed score}\\\hline
 (2,1,1)&(0,0)&(6,3,4,1)&20&4&5180\\
 (2,1,1)&(0,0)&(10,3,4,1)&24&4&6624\\
 (2,1,1)&(0,1)&(6,3,4,1)&20&6&6916
 \end{array}                                         \tag{G023.6}
\]
Because $h=1$, every such row has
\[
 W=c_0+3(q+q^{-1})+4(q^j+q^{-j})
 \quad\text{for one arbitrary integer }j\ge2.
 \tag{G023.7}
\]
The complete $42$-row JSON certificate is included. The
three remaining families, rather than a bounded list of
support indices, are now handled by exact algebra.

\paragraph{Excluding the two mass-four cell families.}
For $d=4$ the quadratic is
$\delta=q^2+q+q^{-1}+q^{-2}$. The nonnegative-index
Laurent coefficient list of $\delta^2$ is
$(4,2,1,2,1)$, and its coefficients at indices $j\ge2$
are at most two. Since $S=2W\delta^2$, (G023.7) gives
\[
 [q^0]S=8c_0+24+16[q^j]\delta^2\le8c_0+56.
\]
For $c_0=6$ or $10$ this is at most $104$ or $136$.
But the true bottom coefficient
$S_{\rm tr}=4D^2(D-1)^2$ has constant Laurent coefficient
$144$. These two families violate coefficientwise
domination for every integer $j\ge2$.

\paragraph{The remaining mass-six family and the fourth-root test.}
Here $\delta=(D-1)(D-2)$ and
$W=3D+4(q^j+q^{-j})$. The parity condition on
$V_2=(DW-\Lambda_2)/2$ gives
\[
 \Lambda_2\equiv DW-2K_2
  \equiv -D^2+10D-8=:C\pmod4.
 \tag{G023.8}
\]
Thus $\ell_2\equiv C(4)=16\equiv0\pmod4$.
Its effective product $A_2=\delta\Lambda_2$ has mass
$6\ell_2\ge16$, so the only possibility below eight
would be $\ell_2=4$.

We exclude that value without referring to $j$.
The coefficient list of $\delta C$ is
$(22,23,10,5,-1)$. Combining this congruence modulo four
with $A_2\ge X_2=(2,3,2,1,1)$ shows that
\[
 A_2=B+4H,\qquad
 B=(2,3,2,1,3),\qquad
 H\in\mathbb N[q^{\pm1}]\text{ symmetric}.
 \tag{G023.9}
\]
The mass of $B$ is $20$. If $\ell_2=4$, the mass of
$A_2$ is $24$, so $H$ has mass one. A symmetric effective
integer Laurent polynomial of mass one is the constant
one. Hence $A_2$ would have coefficient list
$(6,3,2,1,3)$.

At $q=i$ this last polynomial has value $6-4+6=8$.
On the other hand $\delta(i)=0$, so the identity
$A_2=\delta\Lambda_2$ requires $A_2(i)=0$.
This contradiction proves $\ell_2\ge8$.
Also $\ell_1$ is even and $6\ell_1\ge8$, so
$\ell_1\ge2$. Substituting the fixed masses of this
family into the exact score expression gives
\[
 T=17832+64\ell_1+74\ell_2
       \ge17832+64\cdot2+74\cdot8=18552.
\]
Its score is therefore at least $7064$, rather than
the unrealizable scalar-relaxation value $6916$.

Every potential improvement in (G023.6) has now been
excluded for arbitrary outer location. NN-G021 supplies
an actual table of score $7064$, and NN-G022 excludes
unequal quadratics. This proves the exact optimum in
the full class stated above.\hfill$\square$

\paragraph{Geometric dictionary and remaining global scope.}
The fourth-root evaluation is the trace of a finite
character of the grading: restrict a $\mathbb G_m$
representation to $\mu_4$ and evaluate its character at
$i$. This is the diagonalizable-group/comodule formalism
of SGA 3, I, 4.7.2--4.7.3. It is not a finite-field
cardinality, a Frobenius eigenvalue, or a zero of an
arithmetic zeta function. Effective graded objects may
have zero trace at this character. What fails in the
hypothetical small-score table is the required divisibility
of one effective character by the specified cell character.
The actual effective multiplicity objects and their
kernel products remain those of NN-G010 and archive
cote 2, scan 17.

The arbitrary-degree gluing gate is closed for monic
quadratic cells and rational constant ratios. A smaller
score must use nonmonic or nonquadratic cells, nonconstant
ratios, or another
table shape. None of these possibilities is dismissed.
The optimum remains noncanonical and has no implication
for a rank-four canonical basis, PH0, PH1 or the missing
nonstandard-purity bounds.



\section{NN-G024: nonprimitive common cells give the exact 5932-pair optimum}
\paragraph{Statement.}
There is an effective associative rank-three dimension ring with
$5932$ cancelling pairs, of which $5208$ are in the bottom cell.
Moreover $5932$ is the exact minimum in the entire following
class: the closed rank-three table shape, a common cell polynomial
$\delta\in\mathbb Z[D]$ of arbitrary degree and coefficient content,
arbitrary-degree $W\in\mathbb Z[D]$, positive rational constant
ratios $G_i/(W\delta)$, and arbitrary rational auxiliary functions
$\Lambda_i\in\mathbb Q(D)$. Only the actual multiplication
coefficients are required to be integral, effective, coefficientwise
dominant over the absolute trace targets, and congruent to those
targets modulo two.

This class includes all common-cell constructions in NN-G019,
NN-G021 and NN-G023 and removes their monicity and cell-degree
restrictions. The new optimum is attained by a cell polynomial
of content three and a nonintegral auxiliary $\Lambda_2$.
The theorem does not cover unequal general cells, nonconstant
ratios, or other table shapes.

\paragraph{An explicit integral table with a rational auxiliary factor.}
Put $\delta_0=(D-1)(D-2)$ and choose
\[
 \begin{aligned}
 \delta&=3\delta_0,& W&=3D,\\
 G_1&=6D\delta_0,&G_2&=9D\delta_0,&S&=18D\delta_0^2,\\
 U_1&=D,&U_2&=3D,\\
 \Lambda_1&=D-2,&\Lambda_2&=(D^2-2D-4)/3,\\
 V_1&=2D^2-D+2,&V_2&=4D^2+D+2.
 \end{aligned}                                                   \tag{G024.1}
\]
The ratios are $(k_1,k_2)=(2/3,1)$. The identities
\[
 S=G_1G_2/W,\qquad
 U_i=(k_iW-D)/k_j,\qquad
 V_i=(k_iDW-\Lambda_i)/k_j\quad(j\ne i)
\]
hold exactly. The two matrix blocks and two scalar blocks in
NN-G011 therefore give an injective associative realization
of the closed table over $\mathbb Q(D)$: neither
$\Lambda_i(D^2-\Lambda_i)$ vanishes identically.

For effectivity and integrality it is the actual product
$\delta\Lambda_i$ that matters. In symmetric Laurent
coefficient lists indexed by nonnegative exponents, the base
coefficients of this table and their absolute trace targets are
\[
\begin{array}{c|l|l}
 &\text{dimension coefficient}&\text{absolute trace target}\\\hline
 \delta &(6,3,3)&(0,1,1)\\
 \delta\Lambda_1&(6,9,3,3)&(2,1,1,1)\\
 \delta\Lambda_2&(2,5,2,3,1)&(2,3,2,1,1)\\
 W&(6,3)&(2,1)\\
 U_1&(2,1)&(2,1)\\
 U_2&(6,3)&(2,1)\\
 V_1&(12,7,2)&(6,5)\\
 V_2&(28,17,4)&(6,3,2)\\
 G_1&(36,30,18,6)&(16,12,4)\\
 G_2&(54,45,27,9)&(14,11,5,1)\\
 S&(504,450,324,180,72,18)&(144,120,68,24,4)
\end{array}                                                       \tag{G024.2}
\]
Every left list is nonnegative, coefficientwise at least the
right list, and congruent to it modulo two. The closed table
uses these coefficients multiplied by nonnegative powers of
$D$. For signed trace cell coefficients use
$|D^\alpha f|\le D^\alpha|f|$ coefficientwise. These facts
prove every dimension-ring condition. In particular the
nonintegral $\Lambda_2$ is not itself a structure constant;
its product with the nonprimitive cell is integral.

At $D=4$ the bottom mass is
\[
 2592+20(144+216)+200\cdot12+20(4+12+30+70)=14512.
\]
The two cell masses sum to
$52\cdot18+14\cdot6(6+4)=1776$, so $T=16288$.
Subtracting the true absolute masses gives
\[
 (16288-4424)/2=5932,\qquad (14512-4096)/2=5208.
 \tag{G024.3}
\]
The independent integer-Laurent verifier
\path{verify_dimension_5932.py} reconstructs all $132$
positions and checks all $1000$ associativity triples. It
handles the actual integer products directly and does not
round or formally clear an unverified denominator.

\paragraph{Content normalization for the entire common-cell class.}
Write
\[
 \delta=c\delta_0,\quad W=gW_0,
\]
where $c,g>0$ are coefficient contents and $\delta_0,W_0$
are primitive symmetric integral Laurent polynomials.
They are also in $\mathbb Z[D]$, because symmetric integral
Laurent polynomials have an integral triangular expansion
in $q^j+q^{-j}$ and hence in $D=q+2+q^{-1}$.
The actual parity conditions imply that $c,g$ are odd,
\[
 \delta_0\equiv D(D-1),\qquad W_0\equiv D\pmod2.
\]
Both polynomials are effective. In particular
$d:=\delta_0(4)\ge4$ and $w:=W_0(4)\ge4$ are even.

Gauss's content lemma applied after Laurent shifts gives
positive integers $a_i$ such that
\[
 \begin{aligned}
 G_i&=a_iW_0\delta_0,&
 S&=(a_1a_2/g)W_0\delta_0^2,\\
 U_i&=g(a_iW_0-cD)/a_j,&
 V_i&=g(a_iDW_0-\Psi_i)/a_j,
 \end{aligned}                                                     \tag{G024.4}
\]
where $\Psi_i=c\Lambda_i$ and $j\ne i$.
Here $a_i=k_i g c$, $a_1$ is even, $a_2$ is odd, and
$g\mid a_1a_2$. Indeed the primitive products force the
integrality statements, and reduction modulo two of the
$G_i$ gives their different scalar parities.

Even though $\Lambda_i$ may be rational, $\Psi_i$ is in
$\mathbb Z[D]$. First (G024.4), with integral actual $V_i$,
puts $\Psi_i$ in $\mathbb Q[D]$. The actual within-cell
coefficient is $\delta\Lambda_i=\delta_0\Psi_i\in\mathbb Z[D]$.
Primitivity of $\delta_0$ and Gauss's lemma force
$\Psi_i\in\mathbb Z[D]$. This is the exact integrality
statement; it does not force division by the content $c$.
Modulo two these actual cell products also give
$\Psi_1\equiv D$ and $\Psi_2\equiv D^2$.

Set $p_i=\Psi_i(4)$. The exact mass expression is
\[
 \begin{split}
 T={}&(a_1a_2/g)wd^2+20wd(a_1+a_2)+200gw\\
 &+20\sum_i\{U_i(4)+V_i(4)\}+52cd+14d(p_1+p_2).
 \end{split}                                                       \tag{G024.5}
\]
The two effective cell products imply $dp_1\ge8$ and
$dp_2\ge16$. Also $U_i(4)\ge4$ and $V_i(4)\ge16$.
Thus the score bound of NN-G021 applies with an extra
nonnegative $52(c-1)d$ term. For any proposed improvement
on (G024.3) it gives the following generous complete bounds:
\[
 d\le30,\quad w\le36,\quad a_1+a_2\le52,\quad g\le23.
 \tag{G024.6}
\]
Furthermore the right side of (G024.5) is at least
$3024+208c$: use $a_1a_2/g\ge2$, $a_1+a_2\ge3$,
$d,w\ge4$, and the $1136$ lower bound for the last
cell-product term and the four $U,V$ terms. Since a
proposed improvement has $T<16288<18552$, the weaker
bound $c\le73$, with $c$ odd, is valid. None of these
bounds asserts that a Laurent support index is bounded.

\paragraph{A complete finite relaxation, retaining arbitrary support locations.}
Write
\[
 W_0=c_0+c_1(q+q^{-1})+\sum_{j\ge2}c_j(q^j+q^{-j}).
\]
Here $c_0\ge2$ is even, $c_1\ge1$ is odd, and $c_j$
is nonnegative even for $j\ge2$. The same outer-coefficient
argument as in (G023.2) gives $M\mid c_j$, where
\[
 M=\operatorname{lcm}\left(
 \frac{2a_2}{\gcd(ga_1,2a_2)},
 \frac{2a_1}{\gcd(ga_2,2a_1)}\right),\quad 4\mid M.
\]
Put $h=\sum_{j\ge2}c_j/M$, so $w=c_0+2c_1+2Mh$.
For $(b,t)=(c_0,2),(c_1,1)$ the necessary zero and
first coefficient conditions on $U_i$ are
\[
 g(a_i b-ct)/a_j\in t+2\mathbb N.                     \tag{G024.7}
\]
Primitivity requires $\gcd(c_0,c_1)=1$ for $h=0$ and
$\gcd(c_0,c_1,M)=1$ for $h>0$.

Let $a_j'=a_j/\gcd(g,a_j)$. Evaluating the $V_i$
congruences and domination gives the integer intervals
\[
 \begin{gathered}
 p_i\equiv4a_iw\pmod{2a_j'},\\
 \left\lceil X_i(4)/d\right\rceil\le p_i\le
 \left\lfloor(4ga_iw-16a_j)/g\right\rfloor,
 \quad(X_1(4),X_2(4))=(8,16).
 \end{gathered}                                                       \tag{G024.8}
\]
For fixed other data the coefficient of $p_i$ in (G024.5)
is $14d-20g/a_j$. Its minimum over (G024.8) is therefore
at the first allowed integer when this coefficient is
nonnegative and at the last otherwise. This is a lower
bound using necessary conditions, not an assertion that
such a polynomial $\Psi_i$ exists.

There is one additional necessary leading-coefficient
condition when $h=0$. Then $\deg W_0=1$. If $\deg\Psi_2>2$,
its highest Laurent coefficient would have opposite signs
in the effective polynomials $\delta_0\Psi_2$ and $V_2$.
The leading coefficient of $\delta_0$ is positive, so this
is impossible. Since $\Psi_2\equiv D^2\pmod2$, its
coefficient $L$ at $q^2$ must therefore be positive odd.
The coefficient of $V_2$ at $q^2$ is required to be an
even integer at least two. Consequently
\[
 L\ge1,\quad L\equiv a_2c_1-2a_1'\pmod{2a_1'},\quad
 L\le\left\lfloor(g a_2c_1-2a_1)/g\right\rfloor.
 \tag{G024.9}
\]
This test uses no cell-degree restriction.

For reproducibility the whole finite range is as follows:
even $2\le a_1\le50$, odd $1\le a_2\le51$ with sum
at most $52$; odd $1\le g\le23$ dividing $a_1a_2$;
odd $1\le c\le73$; even $4\le d\le30$;
even $2\le c_0\le36$; odd $1\le c_1\le17$;
and $0\le h\le(36-c_0-2c_1)/(2M)$.
Impose (G024.7), the primitive-content necessary tests,
$G_i(4)\ge48$, $S(4)\ge576$, and the endpoint lower
bound from (G024.8). The included integer-only script
\path{verify_common_cell_content_optimum.py}
records every survivor below $T=16288$.
There are seventeen before (G024.9), and exactly five
after it:
\[
\begin{array}{c|c|c|c|c|r}
 (a_1,a_2,g)&c&d&(c_0,c_1,M,h)&w&\text{relaxed score}\\\hline
 (6,9,3)&3&4&(2,1,4,0)&4&3716\\
 (2,1,1)&1&4&(6,3,4,1)&20&5180\\
 (2,1,1)&1&4&(6,7,4,0)&20&5180\\
 (10,5,5)&3&4&(2,1,4,0)&4&5268\\
 (6,15,3)&1&4&(2,1,20,0)&4&5372
\end{array}                                                         \tag{G024.10}
\]
No outer index has been enumerated or bounded.

\paragraph{Uniform exclusion of every survivor.}
All five rows have $d=4$. The effective parity-correct
primitive polynomial of that mass is unique:
\[
 \delta_0=q^2+q+q^{-1}+q^{-2}=D(D-3).
\]
Indeed these four odd coefficients already use the entire
mass. In the first, fourth and fifth rows $W_0=D$ and
$[q^0](D\delta_0)=2$. The first and fifth rows therefore
have $[q^0]G_1=12<16$, while the fourth has
$[q^0]G_2=10<14$. Each violates an actual trace target.

For the third row the coefficient list of $\delta_0^2$
is $(4,2,1,2,1)$ and $S=2W_0\delta_0^2$, so
$[q^0]S=2(6\cdot4+2\cdot7\cdot2)=104<144$.
For the second row $h=1$ means
\[
 W_0=6+3(q+q^{-1})+4(q^j+q^{-j})\quad(j\ge2)
\]
for an arbitrary integer $j$. The same list yields
\[
 [q^0]S=48+24+16[q^j]\delta_0^2\le104<144.
\]
Thus this entire unbounded-location family also fails.
All necessary-relaxation survivors below $5932$ have
been excluded. The actual table (G024.1) attains $5932$,
proving the asserted optimum.\hfill$\square$

\paragraph{Geometric translation and the remaining gate.}
NN-G010 realizes the integral effective multiplicities of
(G024.2) by actual graded sign representations in an
additive monoidal kernel category on a ten-point scheme.
The exact dictionaries remain SGA 3, I, 4.7.2--4.7.3
for graded representations and archive cote 2, scan 17
for finite-kernel composition. There is no object of
fractional dimension attached to $\Lambda_2$; the rational
factor is merely a coordinate factorization in $\mathbb Q(D)$.
The structure constant $\delta\Lambda_2$ is the actual
integral object. The content-three factor is three copies
of its effective cell object, not division by three in
an effective Grothendieck semiring.

For the equal-cell, constant-ratio tables realized by these
chosen kernel generators, the minimum sum over the $132$
coefficient positions and grading pieces of half the excess of dimension over
absolute signed trace is exactly $5932$. This is a minimum
for the specified character-valued multiplication data and
its NN-G010 realizations, not for all monoidal categories
or all generating sets with the same Grothendieck ring.

The same supplied trace algebra therefore admits smaller
composing effective lifts after a purely auxiliary
integrality restriction is removed. This demonstrates
why formal factor parameters cannot be identified with
canonical count objects. It establishes neither a
canonical family nor a rank-four basis, PH0, PH1 or
nonstandard purity. The precise remaining lower-score
gates are unequal general cells, nonconstant ratios,
or a different multiplication-table shape; arbitrary
common-cell degree and content are now covered.



\section{NN-G025: the exact 4520-pair optimum for proportional cells}
\paragraph{Statement.}
The exact minimum is $4520$ cancelling pairs in the closed
rank-three table class with arbitrary proportional effective
cell polynomials $\delta_1,\delta_2\in\mathbb Z[D]$,
arbitrary-degree $W\in\mathbb Z[D]$, positive rational
constant ratios $G_i/(W\delta_i)$, and arbitrary rational
auxiliary functions $\Lambda_i$. All actual coefficients
must obey the usual integrality, domination, effectivity
and parity conditions. Proportional means
$\delta_1/\delta_2\in\mathbb Q_{>0}$; the contents need
not agree and the common primitive polynomial has no
degree restriction.

An attainer has $3532$ bottom-cell cancelling pairs.
This enlarges NN-G024's equal-cell class, whose exact
minimum $5932$ remains valid. The new global upper
bound is $4520$; global minimality without the
proportional-cell and constant-ratio restrictions is
not asserted.

\paragraph{The attainer and its exact score.}
Let $f=(D-1)(D-2)$. Set
\[
 \begin{aligned}
 \delta_1&=f,&\delta_2&=3f,&W&=3D,\\
 G_1&=6Df,&G_2&=5Df,&S&=10Df^2,\\
 U_1&=3D,&U_2&=D,\\
 \Lambda_1&=6D^2-25D+20,&
 \Lambda_2&=(D^2-2D-4)/3,\\
 V_1&=15D-12,&V_2&=2D^2+D+2.
 \end{aligned}                                                     \tag{G025.1}
\]
Here $G_i=k_iW\delta_i$ with $(k_1,k_2)=(2,5/9)$.
Direct substitution gives
\[
 \frac{G_1G_2}{S}=W,\qquad
 \frac{G_i(G_i-D\delta_i)}{S}=U_i,\qquad
 \frac{G_i(DG_i-\delta_i\Lambda_i)}{S}=V_i.
\]
The injective two-matrix-block/two-scalar-block construction
of NN-G011 works with the two different $\delta_i$ and
these nonzero $\Lambda_i(D^2-\Lambda_i)$, proving
associativity.

The changed effective base coefficients, in symmetric
nonnegative-index Laurent lists, are
\[
\begin{array}{c|l|l}
 &\text{dimension coefficient}&\text{absolute trace target}\\\hline
 \delta_1&(2,1,1)&(0,1,1)\\
 \delta_2&(6,3,3)&(0,1,1)\\
 \delta_1\Lambda_1&(22,9,17,5,6)&(2,1,1,1)\\
 \delta_2\Lambda_2&(2,5,2,3,1)&(2,3,2,1,1)\\
 G_1&(36,30,18,6)&(16,12,4)\\
 G_2&(30,25,15,5)&(14,11,5,1)\\
 S&(280,250,180,100,40,10)&(144,120,68,24,4)\\
 V_1&(18,15)&(6,5)\\
 V_2&(16,9,2)&(6,3,2)
\end{array}                                                        \tag{G025.2}
\]
Also $W=U_1=3D$ and $U_2=D$. These lists prove effectivity,
domination and parity for all base coefficients; powers
of $D$ propagate them to the full table exactly as in
NN-G024. No actual coefficient has a denominator.
The independent integer-Laurent checker verifies all
$132$ positions and all $1000$ associativity triples.

At $D=4$ the bottom mass is
\[
 1440+20(144+120)+200\cdot12+20(12+4+48+38)=11160.
\]
The cell mass is
$26(6+18)+14(6\cdot16+6\cdot4)=2304$.
Thus $T=13464$ and
\[
 (13464-4424)/2=4520,\qquad(11160-4096)/2=3532.
 \tag{G025.3}
\]

\paragraph{Normalization and complete necessary bounds.}
Write $\delta_i=c_i f$ with $f$ primitive effective and
$c_i$ positive odd integers, and $W=gW_0$ with $g$
positive odd and $W_0$ primitive effective. Put
$d=f(4)$ and $w=W_0(4)$. The same Gauss and parity
arguments as NN-G024 give positive integers $a_1$ even
and $a_2$ odd, with $g\mid a_1a_2$, and
\[
 \begin{aligned}
 G_i&=a_iW_0f,&S&=(a_1a_2/g)W_0f^2,\\
 U_i&=g(a_iW_0-c_iD)/a_j,&
 V_i&=g(a_iDW_0-\Psi_i)/a_j,
 \end{aligned}                                                      \tag{G025.4}
\]
where $\Psi_i=c_i\Lambda_i\in\mathbb Z[D]$.
Integrality of $\Psi_i$ follows from polynomiality of
$V_i$ and the primitive product $f\Psi_i$, exactly
as before. It does not require $c_1=c_2$.
The actual cell-product masses are $dp_i$ with
$p_i=\Psi_i(4)$, and $dp_1\ge8$, $dp_2\ge16$.
The exact total mass is
\[
 \begin{split}
 T={}&(a_1a_2/g)wd^2+20wd(a_1+a_2)+200gw\\
 &+20\sum_i\{U_i(4)+V_i(4)\}
       +26(c_1+c_2)d+14d(p_1+p_2).
 \end{split}                                                        \tag{G025.5}
\]
For a proposed improvement on $4520$, the complete
bounds (G024.6) still hold, since $c_1+c_2\ge2$.
Moreover
\[
 T\ge3024+104(c_1+c_2).
\]
Using even the weaker cap $T<16288$ gives
$c_1+c_2\le126$, so each positive odd $c_i\le125$.
This generous bound suffices; no degree bound on $f$
or on outer Laurent support locations is being used.

Use $b_0,b_1$ for the zero and first coefficients of
$W_0$, and $M,h$ as in (G024.7), so
$w=b_0+2b_1+2Mh$.
The complete scalar ranges are those of NN-G024,
with the one content variable replaced by
$1\le c_i\le125$, both odd. For each $i$ impose
\[
 g(a_i b-c_i t)/a_j\in t+2\mathbb N
 \quad\text{for }(b,t)=(b_0,2),(b_1,1).               \tag{G025.6}
\]
The bound from the first coefficient also permits
$c_i\le a_i b_1$ without removing any actual case.
The outer-coefficient divisibility, primitive-content
necessary tests, $S(4)\ge576$, and $G_i(4)\ge48$ are
unchanged. The intervals and endpoint rule (G024.8)
for $p_i$ are unchanged as well. When $h=0$, the
leading-coefficient test (G024.9) is still necessary.

For fixed other parameters, the coefficient of $c_i$
in (G025.5) is $26d-80g/a_j$. Thus minimize the
content contribution at the least admissible $c_i$
when that coefficient is nonnegative, and at the
greatest otherwise. The two contents can be minimized
independently for a necessary relaxation: ignoring
their joint upper bound only enlarges the possible
set. Their individual bounds are imposed. An endpoint
row is a lower bound for its whole class of content
choices, not a claim that all those choices coincide.

The exact integer-only script
\path{verify_proportional_cells_optimum.py}
uses precisely these ranges, congruences and endpoint
rules. It leaves just four rows with relaxed score
below $4520$:
\[
\begin{array}{c|c|c|c|c|r}
(a_1,a_2,g)&(c_1,c_2)_{\rm endpoint}&d&(b_0,b_1,h)&
 (p_1,p_2)_{\rm endpoint}&\text{relaxed score}\\\hline
(6,5,3)&(1,3)&4&(2,1,0)&(6,4)&2636\\
(4,9,3)&(1,5)&4&(2,1,0)&(4,8)&3480\\
(6,9,3)&(3,3)&4&(2,1,0)&(6,4)&3716\\
(6,5,3)&(1,3)&6&(2,1,0)&(6,4)&4160
\end{array}                                                        \tag{G025.7}
\]
All have $W_0=D$. In each of the first three,
$f$ has mass four and is consequently $D(D-3)$.
Its product with $D$ has central coefficient two.
Thus their $G_1$ central coefficients are respectively
$12,8,12$, all below the trace target $16$.
These exclusions are independent of the choice of
contents within the row.

\paragraph{The last row and an exact residue obstruction.}
In the remaining row the effective parity-correct
primitive polynomial of mass six is uniquely
$f=(D-1)(D-2)$: relative to the mass-four polynomial,
the only possible addition of mass two is a central
coefficient two. Condition (G025.6) at the first
coefficient forces $c_1=1$ and $c_2=3$. Indeed
$3(6-c_1)/5$ must be a positive odd integer, giving
$c_1=1$, and $(5-c_2)/2$ must be a positive odd
integer, giving $c_2=3$.

Since $W_0=D$, effectivity of $f\Psi_i$ and $V_i$
forces $\deg\Psi_i\le2$ by their opposite leading
signs above degree two. The congruence for $V_1$ is
\[
 \Psi_1\equiv6D^2-5(5D-4)\pmod{10}.
\]
Writing $\Psi_1=L_0+L_1(q+q^{-1})+L_2(q^2+q^{-2})$
gives
\[
 L_0\equiv6,\quad L_1\equiv9,\quad L_2\equiv6
       \pmod{10}.                                      \tag{G025.8}
\]
The highest product coefficient forces $L_2\ge0$,
and $V_1=(3/5)(6D^2-\Psi_1)$ forces $L_2\le6$.
Hence $L_2=6$. The mass $p_1=L_0+2L_1+12$ is
congruent to six modulo ten; $6p_1\ge8$ makes
six its smallest possible positive value.

Suppose $p_1=6$. Then $L_0=-6-2L_1$.
The central coefficient of $f\Psi_1$ is
$2L_0+2L_1+12=-2L_1$ and must be at least two,
so $L_1\le-1$. Its coefficient at $q^3$ is
$L_1+6$ and must be at least one, so $L_1\ge-5$.
Together with (G025.8) this forces $L_1=-1$ and
$L_0=-4$. But its coefficient at $q$ would then be
$L_0+3L_1+6=-1$, below the target one.
Thus $p_1\ne6$, and consequently $p_1\ge16$.
The second mass is a multiple of four and satisfies
$6p_2\ge16$, so $p_2\ge4$.

Substitution in the exact mass identity for this row gives
\[
 T=12016+72p_1+74p_2
   \ge12016+72\cdot16+74\cdot4=13464.
\]
It cannot improve on $4520$. All four rows of the
complete necessary relaxation have now been excluded
below that score, and (G025.1) attains it. This proves
the exact minimum in the stated class.\hfill$\square$

\paragraph{Geometric meaning and residual scope.}
The two cells have the same primitive grading character
and different positive multiplicities. Their actual
coefficient objects still lie in the effective graded
sign-representation semiring, and NN-G010 supplies the
same finite-kernel monoidal realization. The precise
classical dictionaries are SGA 3, I, 4.7.2--4.7.3 and
archive cote 2, scan 17. Rational $\Lambda_2$ remains
an auxiliary factor, not a fractional effective object.
Different contents of the two cell characters do not
alter the supplied trace algebra whose effective lift
is being sought; they change the chosen dimension data.

Among the NN-G010 kernel realizations whose two cell characters
are positive integral multiples of the same primitive character
and whose specified ratios are constant, the least total
half-excess of dimension over absolute signed trace is $4520$,
with $3532$ in the bottom-cell coefficients for (G025.1).
These numerical invariants refer to the chosen table and
generators; they are not invariants of every categorification
of the abstract trace algebra.

The exact common-cell minimum and the exact proportional-cell
minimum differ, so equality of cell characters is a
substantive optimization restriction, not a consequence
of the trace constraints. A smaller score must now use
nonproportional cell polynomials, nonconstant ratios,
or another table shape. No canonicality, generic
rank-four positivity, PH0/PH1 or nonstandard purity
conclusion is inferred from this noncanonical optimum.



\section{NN-G026: the exact optimum for all constant-ratio closed tables}
\paragraph{Statement.}
The exact cancelling-pair minimum is $4520$ for the closed
rank-three table shape with arbitrary effective cell polynomials
$\delta_1,\delta_2\in\mathbb Z[D]$, arbitrary-degree
$W\in\mathbb Z[D]$, positive rational constant ratios
$G_i/(W\delta_i)$, and arbitrary rational auxiliary functions
$\Lambda_i\in\mathbb Q(D)$. The cells need not be proportional;
there are no degree, monicity, or content restrictions.
Every actual structure constant must still be integral,
effective, coefficientwise dominant over its absolute trace
target, and parity-correct.

The attainer is NN-G025. The new assertion is the exclusion
of every smaller table in this enlarged class. Only the
constant-ratio and closed-shape assumptions remain as
optimization restrictions; an unrestricted dimension-ring
optimum is not asserted.

\paragraph{Normalization without proportionality.}
Write $\delta_i=c_i f_i$, where $f_i\in\mathbb Z[D]$
is primitive effective, and $W=gW_0$ with $W_0$ primitive
effective. All contents $c_i,g$ are positive odd integers.
The reduction of each $f_i$ modulo two is
$\sigma=D(D-1)$, and $W_0\equiv D$.
Gauss's content lemma gives positive integers $a_1$ even,
$a_2$ odd, with $g\mid a_1a_2$, such that
\[
 \begin{aligned}
 G_i&=a_iW_0f_i,&S&=(a_1a_2/g)W_0f_1f_2,\\
 U_i&=\frac{g f_i(a_iW_0-c_iD)}{a_j f_j},&
 V_i&=\frac{g(a_iDW_0f_i-B_i)}{a_j f_j},
 \qquad B_i=\delta_i\Lambda_i.
 \end{aligned}                                                       \tag{G026.1}
\]
Here $B_i$ is the actual integral effective cell coefficient.
Unlike in the proportional case, no integrality of
$c_i\Lambda_i$ is asserted or needed.
Put $d_i=f_i(4)$, $w=W_0(4)$ and $b_i=B_i(4)$.
The masses $d_i,w$ are even and at least four, while
$b_1\ge8$ and $b_2\ge16$.

The exact nonunit mass is
\[
 \begin{split}
 T={}&(a_1a_2/g)wd_1d_2+20w(a_1d_1+a_2d_2)+200gw\\
 &+20\sum_i\{U_i(4)+V_i(4)\}
       +26(c_1d_1+c_2d_2)+14(b_1+b_2).
 \end{split}                                                        \tag{G026.2}
\]
Any improvement would have $T<13464$. The base bounds
$U_i(4)\ge4$, $V_i(4)\ge16$, and $b_1+b_2\ge24$
give the following complete finite ranges:
\[
 \begin{gathered}
 4\le w\le24,\quad4\le d_1\le50,\quad4\le d_2\le78
       \quad\text{all even},\\
 a_1\ge2\text{ even},\quad a_2\ge1\text{ odd},\quad
 a_1+a_2\le33,\quad1\le g\le15\text{ odd},\quad g\mid a_1a_2,\\
 1\le c_i\le125\quad\text{odd}.
 \end{gathered}                                                       \tag{G026.3}
\]
For clarity, these are score bounds, not degree bounds.
Indeed $T\ge472w+1344$ and
$T\ge320(a_1+a_2)+2272$ give the first and scalar-sum
bounds, while $T\ge800g$ gives the content bound for $W$.
Also
\[
 T\ge8d_1d_2+186d_1+106d_2+1936.
\]
Setting $d_2\ge4$ gives $T\ge218d_1+2360$;
setting $d_1\ge4$ gives $T\ge138d_2+2680$.
Finally $T\ge3024+104(c_1+c_2)$ gives an even stronger
bound than the generous individual cap $125$ used here.

\paragraph{An exact scalar relaxation for arbitrary cells and support indices.}
For fixed $a_i,g,w,d_i$, enumerate positive odd $c_i$
within (G026.3) and the necessary bound $c_i\le a_iw/4$.
Impose that
\[
 u_i:=\frac{g d_i(a_iw-4c_i)}{a_jd_j}
              \text{ is an even integer at least four}.             \tag{G026.4}
\]
The coefficient of $c_i$ in (G026.2) is
$d_i(26-80g/(a_jd_j))$. Choose the least admissible content
if that coefficient is nonnegative, and the greatest
otherwise. This gives a lower bound for all contents in
that row. Ignoring their joint bound only weakens the
relaxation, and does not exclude any actual table.

For $b_i$ the necessary condition that $V_i(4)$ be an
even integer at least sixteen is exactly
\[
 \begin{gathered}
 b_i\equiv4a_iwd_i\pmod{M_i},\qquad
 M_i=\frac{2a_jd_j}{\gcd(g,2a_jd_j)},\\
 X_i\le b_i\le
 \left\lfloor\frac{4g a_iwd_i-16a_jd_j}{g}\right\rfloor,
 \qquad (X_1,X_2)=(8,16).
 \end{gathered}                                                       \tag{G026.5}
\]
Its coefficient in (G026.2) is $14-20g/(a_jd_j)$,
so the same endpoint rule minimizes it on this finite
congruence interval. Also impose $G_i(4)\ge48$ and
$S(4)\ge576$. These are all necessary conditions;
no divisibility of $B_i$ by $f_i$ is assumed.

The complete integer-only enumeration
\path{verify_constant_ratio_global_optimum.py}
has $1776$ scalar tuples after the coarse mass bounds,
and exactly $55$ endpoint rows below $T=13464$.
The full JSON certificate lists every row. Thirty-one
have $w=4$. In all remaining twenty-four rows
\[
 (a_1,g)=(2,1),\qquad a_2\in\{1,3\}.
\]
Their mass pairs are exactly
\[
\begin{array}{c|c|l}
 a_2&w&(d_1,d_2)\\\hline
 1&8&(4,12),(6,12),(6,18),(8,16)\\
 1&10&(4,8),(6,8),(6,12),(4,16),(6,16)\\
 1&12&(4,6),(4,8),(4,10),(6,6),(4,12),(6,12)\\
 1&14&(4,8)\\
 1&16&(4,8)\\
 3&8&(4,4),(6,4),(6,6),(10,4)\\
 3&10&(6,4)\\
 3&12&(4,4),(6,4)
\end{array}                                                         \tag{G026.6}
\]
The endpoint contents in the certificate need not equal
those of an actual hypothetical table. The exclusions
below either apply to every content in the row or derive
the actual contents anew from (G026.4).

\paragraph{Three elementary polynomial observations.}
First, coprime $f_1,f_2$ are impossible in (G026.1).
If they were coprime over $\mathbb Q$, polynomiality of
$U_1$ would force $f_2\mid a_1W_0-c_1D$ over $\mathbb Q[D]$.
Primitivity and Gauss's lemma make the quotient integral.
Reducing modulo two would say that $D(D-1)$ divides $D$,
which is false.

Second, when $w=4$, effectivity and parity force $W_0=D$,
and the two cells must be proportional. To see this,
write $f_1=Hu_1$, $f_2=Hu_2$ with primitive integer
polynomials and coprime $u_1,u_2$ over $\mathbb Q$.
The nonzero factors $a_iW_0-c_iD=(a_i-c_i)D$ in
(G026.1) force both $u_1,u_2$ to divide $D$.
A zero factor would make $U_i=0$, which is forbidden.
Since $f_1\equiv f_2\equiv\sigma\ne0$ modulo two,
the nonzero reduction of $H$ gives $u_1\equiv u_2$.
The choices $1$ and $D$ cannot be different, and they
cannot both be $D$ because $u_1,u_2$ are coprime.
Thus both are constant and $f_1=f_2$. NN-G025 excludes
all thirty-one $w=4$ rows below $4520$.

Third, the only effective parity-correct primitive
polynomials of masses four and six are respectively
\[
 f^{(4)}=D(D-3),\qquad f^{(6)}=(D-1)(D-2).
\]
They are coprime. Every such polynomial of mass eight
has the form
\[
 f^{(4)}+4\quad\text{or}\quad
 f^{(4)}+2(q^j+q^{-j})\quad(j\ge1).                   \tag{G026.7}
\]
No polynomial in (G026.7) vanishes at $D=0$ or $D=3$.
For $D=0$, choose $q=-1$, so the added pair has value
$4(-1)^j\ne0$. For $D=3$, choose $q=e^{i\pi/3}$;
the pair contributes $4\cos(j\pi/3)$, which is one
of $4,2,-2,-4$. The constant addition four also
never vanishes. Hence every mass-eight polynomial
is coprime to $f^{(4)}$, uniformly in its unbounded
support location $j$.

\paragraph{Parity modulo four eliminates the nonmultiples of four.}
For the remaining rows $a_1=2$ and $g=1$. Define
\[
 Q=(a_2W_0-c_2D)/2\in\mathbb Z[D],
\]
which is integral because $a_2,c_2$ are odd and
$W_0\equiv D$. The second $U$ identity is
$f_1U_2=f_2Q$. Reducing it modulo two and cancelling
$\sigma\ne0$ gives $Q\equiv D$. Therefore
\[
 a_2W_0\equiv(c_2+2)D\pmod4.                        \tag{G026.8}
\]
At $D=4$ this implies $w\equiv0\pmod4$.
Thus all $w=10$ and $w=14$ rows in (G026.6) are
excluded, regardless of their contents.

\paragraph{The rows with $a_2=1$.}
For $w=8$, positivity of $U_2(4)$ forces $c_2<2$,
so $c_2=1$. Every mass pair listed in this group has
$d_2\ge2d_1$. Then $U_1(4)\ge4$ in (G026.4) forces
$c_1\le2$, hence $c_1=1$ as well.
Equation (G026.8) gives $W_0\equiv3D\pmod4$.
At total mass eight there is exactly one effective
such polynomial:
\[
 W_0=3D-4,
\]
with coefficient list $(2,3)$; any additional outer
pair or increase of a noncentral coefficient modulo
four would already use at least eight further units
of mass.

Write $f_i=Hu_i$ as above. Polynomiality of the two
$U_i$ now forces
\[
 u_1\mid D-2,\qquad u_2\mid5D-8.
\]
Their reductions modulo two agree. Thus either both
are constant, which is the proportional case, or
$u_1=D-2$ and $u_2=5D-8$, up to the already normalized
positive primitive units. In the latter case the
first $U$ identity gives $U_1=D-2$. Its central
Laurent coefficient is zero, below the required
coefficient two of $D$. This excludes all four
$w=8$ rows.

For $w=12$, positivity of $U_2(4)$ gives $c_2<3$,
so again $c_2=1$. Equation (G026.8) requires
$W_0\equiv3D\pmod4$. An effective polynomial of
mass twelve with this residue is uniquely $3D$:
the minimum coefficients are $(2,3)$ of mass eight,
and the remaining four units can only increase the
central coefficient by four. But $3D$ has content
three, contradicting primitivity of $W_0$. All six
rows in this group are excluded.

The only $w=16$ row has mass pair $(4,8)$, whose
primitive cells are coprime by (G026.7). It is
excluded by the first polynomial observation.

\paragraph{The rows with $a_2=3$.}
After the modulo-four test, these have $w=8$ or
$w=12$. Mass pairs $(4,4)$ and $(6,6)$ force identical
primitive cells and are covered by NN-G025; $(6,4)$
is coprime and impossible. The only remaining pair
is $(d_1,d_2)=(10,4)$ at $w=8$.
In this row
\[
 U_1(4)=(10/12)(16-4c_1),\qquad
 U_2(4)=(24-4c_2)/5.
\]
Requiring both to be at least four and the contents
to be positive odd gives $c_1=c_2=1$.
Equation (G026.8) gives $W_0\equiv D\pmod4$.
At mass eight this uniquely gives $W_0=D+4$.
Now $f_2=D(D-3)$ and polynomiality of
\[
 U_1=\frac{f_1(D+8)}{3f_2}
\]
forces $f_2\mid f_1$, because $f_2$ is coprime
to $D+8$. The quotient is integral by Gauss's lemma.
Its value at $D=4$ would be $d_1/d_2=10/4$, which
is not an integer. This last row is impossible.

All fifty-five necessary-relaxation rows have been
excluded below $4520$. The NN-G025 attainer belongs
to the stated class, so the minimum is exactly
$4520$.\hfill$\square$

\paragraph{Geometric dictionary and the smaller remaining gate.}
The parity and modulo-four arguments are constraints
on actual effective graded sign representations:
dimension and signed trace agree modulo two, while
the composing table imposes further integral divisions.
SGA 3, I, 4.7.2--4.7.3 provides the diagonalizable-group
character dictionary, and the NN-G010 realization is
the finite-kernel composition recorded in archive
cote 2, scan 17. The evaluations in (G026.7) restrict
the grading character to finite roots of unity;
they are not Frobenius counts or arithmetic zeta zeros.
No rational auxiliary factor is promoted to an
effective object unless its actual coefficient has
been verified integral and nonnegative.

In geometric terms, all effective graded multiplicity objects
for the two cells may now vary independently, subject to this
closed kernel-table presentation and its constant rational
ratio identities. Among the resulting NN-G010 realizations,
the least total half-excess of dimension over the absolute
signed trace coefficients is exactly $4520$. The trace
specialization is fixed, but no canonical choice of objects,
coherence data, or generators is claimed. This optimum is
presentation-dependent, not a universal categorical invariant.

Arbitrary nonproportional cell polynomials and
unbounded support locations have now been treated
within the constant-ratio closed-table class. A
smaller score must use nonconstant ratios
$G_i/(W\delta_i)$ or a different table shape.
This is an exact realization theorem with explicit
scope, not a canonical-basis theorem or a consequence
for PH0, PH1, rank four, or nonstandard purity.



\section{NN-G027: nonconstant ratios produce smaller explicit realizations}
\paragraph{Statement.}
There is a closed-table effective associative dimension ring
with $3452$ cancelling pairs, of which $2852$ are in the
bottom cell. The addendum below improves these numbers to
$3352$ and $2752$, respectively. Both constructions have
nonconstant rational functions
$G_i/(W\delta_i)$, so both lie outside the exact $4520$
class of NN-G026. Thus the constant-ratio restriction is
substantive. No exact optimum over the enlarged rational-function
class is asserted.

\paragraph{Construction and associativity.}
Put
\[
 \sigma=D(D-3),\qquad h=\sigma+1,\qquad B=5\sigma+4.
\]
Choose
\[
 \begin{aligned}
 \delta_1&=\sigma+2,&\delta_2&=3\sigma+2,&W&=3D,\\
 G_1&=6Dh,&G_2&=DB,&S&=2DhB,\\
 U_1&=3D,&U_2&=D,&V_1&=15D-6,&V_2&=2D^2+D,\\
 A_1&=6D^2h-B(5D-2),&
 A_2&=D^2B-2h(2D^2+D),&
 \Lambda_i&=A_i/\delta_i.
 \end{aligned}                                                       \tag{G027.1}
\]
The actual within-cell coefficients involving $\Lambda_i$
are $A_i$ multiplied by nonnegative powers of $D$.
They are integral polynomials even though the auxiliary
$\Lambda_i$ are rational functions.

All three scalar-block identities follow directly:
\[
 \frac{G_1G_2}{S}=3D=W,\qquad
 G_1-D\delta_1=G_2,\qquad
 G_2-D\delta_2=G_1/3.
\]
Consequently $G_i(G_i-D\delta_i)/S=U_i$.
The two definitions of $A_i$ give
\[
 \frac{G_1(DG_1-A_1)}S=3(5D-2)=V_1,\qquad
 \frac{G_2(DG_2-A_2)}S=V_2.
\]
Also $\Lambda_i(D^2-\Lambda_i)\ne0$: the leading
terms of $\Lambda_1,\Lambda_2$ are respectively
$6D^2$ and $D^2/3$. The injective generic block
construction of NN-G011 therefore proves associativity
over $\mathbb Q(D)$. As all actual structure constants
are integral, its identities hold over $\mathbb Z[D]$.
This argument makes no semisimplicity claim at the
special values where an auxiliary factor has a pole.

\paragraph{Coefficient verification and exact score.}
The following symmetric Laurent lists give every base
coefficient needed for the effectivity, domination and
parity checks:
\[
\begin{array}{c|l|l}
 &\text{dimension coefficient}&\text{absolute trace target}\\\hline
 \delta_1&(2,1,1)&(0,1,1)\\
 \delta_2&(2,3,3)&(0,1,1)\\
 A_1&(14,5,1,5,6)&(2,1,1,1)\\
 A_2&(2,3,4,3,1)&(2,3,2,1,1)\\
 G_1&(24,24,18,6)&(16,12,4)\\
 G_2&(18,19,15,5)&(14,11,5,1)\\
 S&(172,152,114,78,40,10)&(144,120,68,24,4)\\
 V_1&(24,15)&(6,5)\\
 V_2&(14,9,2)&(6,3,2)
\end{array}                                                         \tag{G027.2}
\]
The remaining base coefficients are $W=U_1=3D$ and
$U_2=D$. Every entry in the middle column dominates
the right column coefficientwise with the same parity.
The closed multiplication table only multiplies these
base coefficients by nonnegative powers of $D$;
for signed traces use the coefficientwise triangle
inequality. Hence all $132$ positions satisfy the
required conditions. The included independent
integer-Laurent verifier also checks all $1000$
associativity triples and reports zero failures.

At $D=4$ the bottom mass is
\[
 960+20(120+96)+200\cdot12+20(12+4+54+36)=9800.
\]
The two cell characters have masses $6,14$, and the
actual $A_1,A_2$ have masses $48,24$. Thus their
contribution is $26(6+14)+14(48+24)=1528$.
The total is $T=11328$, yielding
\[
 (11328-4424)/2=3452,\qquad (9800-4096)/2=2852.
 \tag{G027.3}
\]
Finally
\[
 \frac{G_1}{W\delta_1}=\frac{2(\sigma+1)}{\sigma+2},
 \qquad
 \frac{G_2}{W\delta_2}=\frac{5\sigma+4}{3(3\sigma+2)},
\]
which are both nonconstant. The lower score is therefore
consistent with NN-G026, and proves that extending
its optimum beyond constant ratios would be false.
\hfill$\square$

\paragraph{A verified refinement to 3352 pairs.}
The same identities permit a smaller realization. Keep
$H=D^2-3D+1$, and now put
\[
 J=5D^2-17D+10.
\]
Define
\[
 \begin{aligned}
 \delta_1&=6H-J=D^2-D-4,&
 \delta_2&=J-2H=3D^2-11D+8,\\
 G_1&=6DH,&G_2&=DJ,&S&=2DHJ,\\
 W&=3D,&U_1&=3D,&U_2&=D,\\
 V_1&=15D,&V_2&=2D^2+D-4,\\
 A_1&=6D^2H-5DJ,&
 A_2&=D^2J-2H(2D^2+D-4),&\Lambda_i&=A_i/\delta_i.
 \end{aligned}                                                       \tag{G027.4}
\]
Again $G_1-D\delta_1=G_2$ and
$G_2-D\delta_2=G_1/3$, while the definitions of $A_i$
give the two $V_i$ identities. Neither
$\Lambda_i(D^2-\Lambda_i)$ is identically zero, so the
same injective block construction proves associativity.
The required base coefficients are
\[
\begin{array}{c|l|l}
 &\text{dimension coefficient}&\text{absolute trace target}\\\hline
 \delta_1&(0,3,1)&(0,1,1)\\
 \delta_2&(4,1,3)&(0,1,1)\\
 A_1&(6,5,1,5,6)&(2,1,1,1)\\
 A_2&(6,5,6,1,1)&(2,3,2,1,1)\\
 G_1&(24,24,18,6)&(16,12,4)\\
 G_2&(18,17,13,5)&(14,11,5,1)\\
 S&(156,140,106,70,36,10)&(144,120,68,24,4)\\
 V_1&(30,15)&(6,5)\\
 V_2&(10,9,2)&(6,3,2)
\end{array}                                                         \tag{G027.5}
\]
They prove every base condition; powers of $D$ propagate
them to the full table. The independent verifier
\path{verify_dimension_3352.py} checks all $132$
positions and all $1000$ triples with zero failures.

The bottom mass is
\[
 880+20(120+88)+200\cdot12+20(12+4+60+32)=9600.
\]
The cell mass is $26(8+12)+14(40+32)=1528$.
Consequently
\[
 T=11128,\qquad
 (11128-4424)/2=3352,\qquad(9600-4096)/2=2752.
 \tag{G027.6}
\]
The two ratios are $2H/(6H-J)$ and
$J/(3(J-2H))$, both nonconstant. This strengthens the
explicit upper bound to $3352$, retaining the original
$3452$ example as the first checked construction.
NN-G028 proves an exact optimum for the indicated
fixed-bottom subclass; no unrestricted optimum is
claimed here.\hfill$\square$

\paragraph{Geometric meaning and residual scope.}
The actual polynomials in (G027.2) are effective
Artin--Tate multiplicity objects and define another
NN-G010 finite-kernel monoidal realization of the
same trace table. With its chosen generators, the
sum over multiplication positions and grading pieces
of half the excess of dimension over absolute signed
trace is $3452$ for the first construction, and $3352$
for (G027.4), with bottom-cell contributions $2852$
and $2752$, respectively.
This is an explicit upper bound for a presentation-dependent
quantity, not a universal invariant of all categories
realizing the abstract algebra.

The rational $\Lambda_i$ and the two nonconstant ratios
are generic coordinate factorizations; no object of
fractional dimension or meromorphic sheaf is inferred.
Only the verified actual coefficients define objects.
The precise dictionary is SGA 3, I, 4.7.2--4.7.3 and
archive cote 2, scan 17 through NN-G010. The example
leaves the exact optimum for arbitrary rational ratios
and for other table shapes open. It does not give a
canonical family, a rank-four basis, PH0/PH1 or
nonstandard purity.



\section{NN-G028: fixed bottom coefficients force an exact 3352-pair optimum}
\paragraph{Statement.}
Among all effective associative closed rank-three tables with
\[
 W=3D,\qquad U_1=3D,\qquad U_2=D,                    \tag{G028.1}
\]
the exact cancelling-pair minimum is $3352$.
Cell polynomials, contents and degrees are arbitrary;
$G_i/(W\delta_i)$ and $\Lambda_i$ may be arbitrary
rational functions of $D$. Only the actual table
coefficients are required to be integral, effective,
dominant and parity-correct. There are exactly two
minimizing parameter sets within this closed presentation.
One is (G027.4); the other replaces its $V_2$ by
$2D^2-D+4$ and adjusts $A_2$ accordingly.

The fixed coefficients (G028.1) are a substantive
restriction. This theorem closes a rational-function
optimization problem inside that restriction, and
makes no lower-bound assertion when $W$ or the two
$U_i$ vary.

\paragraph{A forced two-polynomial parameterization.}
The scalar identities used in NN-G011 are necessary
for associativity in this closed table: associativity
with $\Sigma$ gives
\[
 WS=G_1G_2,\qquad U_iS=G_i(G_i-D\delta_i),\qquad
 V_iS=G_i(DG_i-A_i),\quad A_i=\delta_i\Lambda_i.
\]
For example, $\Sigma(\gamma_1^1\gamma_1^2)$ gives
the first identity; multiplication by $\Sigma$ of
a same-cell two-letter product gives the second,
and of a three-letter reduction gives the third.
Thus no generic-embedding assumption is added in
the lower-bound argument.

Substituting (G028.1) yields
\[
 G_1-D\delta_1=G_2,\qquad
 3(G_2-D\delta_2)=G_1.
\]
Consequently
\[
 \begin{gathered}
 H=(\delta_1+\delta_2)/4,\qquad
 J=(\delta_1+3\delta_2)/2,\\
 G_1=6DH,\quad G_2=DJ,\quad S=2DHJ,\qquad
 \delta_1=6H-J,\quad\delta_2=J-2H.                 \end{gathered}\tag{G028.2}
\]
Both $H,J$ are effective integral symmetric Laurent
polynomials, hence belong to $\mathbb Z[D]$.
For $J$, integrality follows directly from
$\delta_1\equiv\delta_2$ modulo two. To prove it for
$H$, first put $K=(\delta_1+\delta_2)/2\in\mathbb Z[D]$.
The actual coefficient $G_1=3DK$ is zero modulo two,
so cancellation in $\mathbb F_2[D]$ gives $K\equiv0$;
therefore $H=K/2$ is integral. Effectivity follows
from their displayed positive combinations of the cells.
The $G_2$ parity gives $J\equiv D(D-1)$ modulo two.

Let $\sigma=q^2+q+q^{-1}+q^{-2}=D(D-3)$ denote
the absolute cell trace coefficient. The actual cell
bounds in (G028.2) are exactly
\[
 2H+\sigma\le J\le6H-\sigma                       \tag{G028.3}
\]
coefficientwise. In particular $H$ has coefficients
at least one at each of $q^{\pm1},q^{\pm2}$.
Write
\[
 m=H(4)\ge4,\qquad n=J(4)\ge2m+4.
\]
Here $m$ is an integer and $n$ is a positive even integer.
No degree or Laurent-support bound has been imposed.

\paragraph{A mass inequality removes every $m\ge6$.}
Let $v_i=V_i(4)$ and $a_i=A_i(4)$ in this paragraph;
these symbols are masses, not the ratio numerators of
NN-G026. The actual coefficients obey
\[
 a_1=96m-(n/3)v_1\ge8,\qquad
 a_2=16n-2mv_2\ge16.
\]
The exact total mass becomes
\[
 \begin{split}
 T={}&8mn+584m+80n+2720+
       \frac{5760m}{n}+\frac{160n}{m}\\
    &+\left(14-\frac{60}{n}\right)a_1
       +\left(14-\frac{10}{m}\right)a_2.
 \end{split}                                                        \tag{G028.4}
\]
Both coefficients on the last line are positive,
since $m\ge4$ and $n\ge12$. Therefore
\[
 T\ge L(m,n):=8mn+584m+80n+3056
       +\frac{5760m-480}{n}+\frac{160n-160}{m}.       \tag{G028.5}
\]
For $m\ge6$ and $n\ge2m+4$, this is minimized at
$(m,n)=(6,16)$. Here is a direct verification of that
global assertion, rather than a bounded numerical test.
The derivative with respect to $n$ is
\[
 8m+80+160/m-(5760m-480)/n^2.
\]
Its value at $n=2m+4$, multiplied by $m(m+2)^2$,
is
\[
 8\{m^2(m^2+14m-116)+135m+80\}>0
 \quad(m\ge6).
\]
It is increasing in $n$, so $L(m,n)\ge L(m,2m+4)$.
The latter expression is
\[
 16m^2+776m+6576-\frac{6000}{m+2}+\frac{480}{m},
\]
whose derivative
$32m+776+6000/(m+2)^2-480/m^2$ is positive for
$m\ge6$. Thus
\[
 T\ge L(6,16)=11138>11128.                          \tag{G028.6}
\]
The new attainer has total mass $11128$, so every
competing minimizer has $m\le5$.

If $m=4$, (G028.3) forces $H=\sigma$.
Then $[q^0]G_1=[q^0](6D\sigma)=12<16$, violating
its trace target. If $m=5$, the only possible addition
to those four mandatory coefficients is one central
unit, so
\[
 H=(1,1,1)=D^2-3D+1.                                \tag{G028.7}
\]
This proves a degree reduction from the score and
actual coefficient inequalities; it does not infer
a degree bound from mass alone.

\paragraph{Only four $J$ polynomials remain.}
By (G028.3) the support of $J$ is contained in that
of $H$, so write its symmetric coefficient list as
$(j_0,j_1,j_2)$. Its bounds and parity give
\[
 j_0\in\{2,4,6\},\qquad j_1,j_2\in\{3,5\}.
\]
Effectivity of $A_i=DG_i-(G_j/W)V_i$ forces
$\deg V_i\le2$: above that degree its highest term
would have negative sign, since $H,J,V_i$ have
positive leading coefficients. The coefficient of
$V_2$ at $q^2$ is even and at least two. The top
coefficient of $A_2=D^2J-2HV_2$ is therefore at most
$j_2-4$, forcing $j_2=5$; its $V_2$ coefficient is
then exactly two.

The central coefficient of $S=2DHJ$ equals
$8j_0+16j_1+60$ and must be at least $144$.
This leaves precisely
\[
 (j_0,j_1,j_2)=(2,5,5),(4,5,5),(6,3,5),(6,5,5).
 \tag{G028.8}
\]
Each $J$ in this list is primitive. Since
$A_1=6D^2H-(J/3)V_1$ is integral, Gauss's lemma
forces $V_1=3C$ with $C$ an integral effective
polynomial. The parity and target of $V_1$ give
$C_0\ge2$ even, $C_1\ge3$ odd, and $C_2$ even.
The leading inequality $6-5C_2\ge0$ forces
$C_2=0$. Thus $C=(x,y)$ with $x\ge2$ even and
$y\ge3$ odd.

\paragraph{Complete coefficient optimization of the two remaining factors.}
The symmetric list of $6D^2H$ is
$(96,90,66,30,6)$. Subtracting $JC$ and comparing
with $X_1=(2,1,1,1)$ gives exactly
\[
 \begin{gathered}
 y\in\{3,5\},\qquad
 5x+j_1y\le65,\\
 j_1x+(j_0+5)y\le89,\qquad
 j_0x+2j_1y\le94.                                  \end{gathered}\tag{G028.9}
\]
For fixed $y$, the largest even $x$ satisfying these
three bounds maximizes $C(4)=x+2y$.

Write $V_2=(z,t,2)$, where $z\ge6$ is even and
$t\ge3$ is odd. Expanding $D^2J-2HV_2$ and comparing
with $X_2=(2,3,2,1,1)$ gives exactly
\[
 \begin{gathered}
 t\le(j_1+15)/2,\qquad
 z+2t\le3j_0+4j_1,\\
 z+2t\le2j_0+(7j_1+13)/2,\qquad
 z+t\le j_0/2+2j_1+12.                             \end{gathered}\tag{G028.10}
\]
Here $t$ ranges only over $3,5,7,9$. For fixed $t$
take the largest allowed even $z$. These lists
include every coefficient condition: the remaining
top coefficients are $6$ and $1$, and parity is
already enforced. The total mass is strictly decreasing
in $v_1$ and $v_2$, with coefficients
$20-14n/3$ and $20-28m$, so maximizing these masses
independently is the exact minimizing rule.

The finite arithmetic is displayed in full by
\path{verify_fixed_bottom_optimum.py}; its four
rows are
\[
\begin{array}{c|c|c|c|c|r}
 J&\max C(4)&\max V_2(4)&a_1&a_2&\text{score}\\\hline
 (2,5,5)&18&30&84&52&3720\\
 (4,5,5)&18&36&48&24&3452\\
 (6,3,5)&20&32&40&32&3352\\
 (6,5,5)&16&38&64&36&3728
\end{array}                                                         \tag{G028.11}
\]
There is only one minimizing $J$. In that row (G028.9)
forces the maximizing $C=(10,5)$, so $V_1=15D$.
Equation (G028.10) has exactly two maximizing choices,
\[
 V_2=(10,9,2)=2D^2+D-4,
 \quad\text{or}\quad
 V_2=(14,7,2)=2D^2-D+4.
\]
The first is the verified construction (G027.4).
For the second the actual $A_2$ list is
$(6,5,2,5,1)$; it has the same mass and satisfies
all required conditions. The other data are unchanged,
and the generic block construction again proves
associativity. This proves both attainment and the
complete two-parameter-set minimum in the class
(G028.1).\hfill$\square$

\paragraph{Geometric translation and the next exact gate.}
Fixing (G028.1) fixes three actual multiplicity objects
in the NN-G010 finite-kernel presentation. Among all
its effective graded sign-representation tables with
those fixed objects and the specified trace specialization,
the sum over multiplication positions and grading pieces
of half the dimension-minus-absolute-trace excess is
exactly $3352$. There are two minimizing character-valued
parameter sets in this presentation; no uniqueness of
categories or coherence data is asserted. This is an
optimization of chosen multiplication data, not an
intrinsic minimum for every categorification of the
abstract trace algebra.

The exact source dictionary remains SGA 3, I,
4.7.2--4.7.3 for effective grading characters and
archive cote 2, scan 17 for kernel composition.
The mass reduction respects actual coefficient
objects and never treats a rational auxiliary factor
as an effective object. A smaller closed-table score
must change at least one of $W,U_1,U_2$ in (G028.1);
other table shapes are a further separate possibility.
No canonicality, rank-four, PH0/PH1 or purity theorem
is inferred from this exact restricted optimum.



\section{NN-G029: an existence-preserving degree cutoff makes the unrestricted optimum decidable}
\paragraph{Statement.}
Fix a finite based trace algebra with bar-invariant Laurent
structure constants $c_{ij}^k(q)\in\mathbb Z[q^{\pm1}]$,
including a specified unit. An effective symmetric dimension
lift consists of bar-invariant $d_{ij}^k(q)\in\mathbb N[q^{\pm1}]$
that are associative and unital, dominate $|c_{ij}^k|$
coefficientwise, and have the same coefficients modulo two.
Extra positions with zero trace are permitted unless explicitly
forbidden. Its cancelling-pair cost is
\[
 C(d)=\frac12\sum_{i,j,k,e}
             \bigl(d_{ij,e}^k-|c_{ij,e}^k|\bigr).
\]
Suppose one such lift of cost at most $C_0$ is known.
Put
\[
 N_0=\sum_{i,j,k,e}|c_{ij,e}^k|,\qquad N=N_0+2C_0,
 \qquad K=\max\bigl(1,\max\{|e|:c_{ij,e}^k\ne0\}\bigr),
\]
Let $E$ consist of the positive trace exponents together with one,
put $h=|E|$, and set $n_* =\lfloor N/2\rfloor+h$.
Every lift of cost at most $C_0$ has another lift of exactly
its cost, with the same coefficient lists and the same nonzero
structure-constant positions, whose Laurent degrees are at most
\[
 B=n_*(K+1)4^{\,n_*-1}.                              \tag{G029.1}
\]
Consequently a finite terminating search computes the true
minimum over all symmetric effective associative lifts in
this presentation, without a closed-table or constant-ratio
restriction.

This is an existence-preserving compression statement.
It does \emph{not} bound the degrees of every existing lift,
assert a practical search algorithm, or identify the numerical
unrestricted optimum.

For the supplied rank-three trace table, $N_0=4443$ and $K=5$.
The NN-G027 incumbent has $C_0=3352$, so $N=11147$ and
$n_*=5578$. Every competing lift therefore has an equal-score
representative of Laurent degree
\[
 \le33468\,4^{5577}<2^{11170}.                       \tag{G029.2}
\]
One may search all $10^3$ structure-constant positions, including
new even coefficients at positions of zero trace, or retain any
specified forbidden positions. Both variants are finite and
complete for their stated presentation.

\paragraph{Elementary small-integer-point lemma.}
Let $P\subset\mathbb R^n$ be a nonempty polyhedron defined by
finitely many integer inequalities $a_r\cdot x\ge b_r$.
Equalities may be represented by their two inequalities.
Assume $P\subseteq\{x_i\ge1\}$, every row has Euclidean
norm at most four, and $|b_r|\le K$ with $K\ge1$.
If $P$ contains an integer point, it contains one satisfying
\[
 1\le x_i\le n(K+1)4^{n-1}\quad\text{for every }i.
 \tag{G029.3}
\]

Here is a complete polyhedral proof. Include the inequalities
$x_i\ge1$ in the presentation and homogenize:
\[
 \mathcal C=\{(x,t):t\ge0,\ a_r\cdot x\ge b_rt\}.
\]
Because $x_i\ge t\ge0$, this cone lies in the nonnegative
orthant and has no nonzero line. Its section
$\sum_i x_i+t=1$ is a compact polytope. Every point of
this section is a convex combination of its vertices:
if its smallest face has positive dimension, choose a
nonzero direction in that face's affine span, follow it
in both directions to the two boundary endpoints of
the compact face, and write the point as their convex
combination. The endpoints lie in faces of smaller
dimension, so induction terminates. There are only
finitely many vertices, since each is determined by
a full-rank subset of the finite active constraints.
Thus $\mathcal C$ is generated by finitely many extreme
rays.

An extreme ray with $t>0$, rescaled to $t=1$, has the
form $(v,1)$ with $v$ a vertex of $P$: a nontrivial
segment through $v$ in $P$ would contradict extremality
of that ray. An extreme ray with $t=0$ has the form
$(r,0)$ with $r$ an extreme ray of the recession cone
\[
 \operatorname{rec}(P)=\{r:a_r\cdot r\ge0\}.
\]
In particular $r_i\ge0$. Applying the cone decomposition
to an existing integer point $(x,1)$ gives
\[
 x=\sum_j\lambda_jv_j+\sum_s\mu_s r_s,
 \qquad \lambda_j,\mu_s\ge0,\quad\sum_j\lambda_j=1.
 \tag{G029.4}
\]
This proves the needed vertex-plus-recession decomposition
also when $P$ has lower affine dimension. It explains
explicitly the rays on the $t=0$ face; none is lost.

A vertex $v$ has $n$ linearly independent active integer
rows. Otherwise a small positive and negative displacement
in their common kernel would remain feasible, contradicting
that it is a vertex. Orthogonalizing the rows by Gram--Schmidt shows that the
absolute determinant is the product of the orthogonalized
row lengths, each at most its original length; this proves
the determinant bound used below. Cramer's rule then gives
\[
 |v_i|\le nK4^{n-1}:                                  \tag{G029.5}
\]
the denominator is a nonzero integer of absolute value at
least one, and each numerator is a sum of $n$ numbers
$b_r$ times cofactors bounded by $4^{n-1}$.

Likewise an extreme recession ray has $n-1$ independent
active homogeneous rows. Otherwise its active common
kernel would have a direction transverse to the ray,
and small displacements would split it inside the cone.
Its integer cofactor vector spans the ray. Choose the
nonnegative orientation and divide by the greatest common
divisor of its entries to obtain a primitive integer
ray $r$, with
\[
 0\le r_i\le4^{n-1}.                                 \tag{G029.6}
\]
The assertion also covers $n=1$, using the empty determinant
one. If the recession cone is zero there are simply no
such rays.

The ray part of (G029.4) needs at most $n$ rays. If more
are used, they are linearly dependent. A nonzero linear
relation among nonnegative nonzero rays has coefficients
of both signs. Subtract a suitable multiple of that
relation from the nonnegative coefficients until one
coefficient becomes zero, preserving all the others'
nonnegativity, and repeat. The vertex convex combination
is left unchanged.

Now replace each of these at most $n$ coefficients $\mu_s$
by its fractional part. Equivalently set
\[
 x'=x-\sum_s\lfloor\mu_s\rfloor r_s.
\]
This is an integer vector, because $x$ and the chosen
primitive rays are integral. It remains in $P$, because
it is the same convex combination of vertices plus a
nonnegative recession combination. The coefficients in
that combination are less than one, so (G029.5)--(G029.6)
give
\[
 x'_i\le nK4^{n-1}+n4^{n-1}=n(K+1)4^{n-1}.
\]
Since $P\subseteq\{x_i\ge1\}$, this proves (G029.3).
Notice that no rounding of the vertex convex combination
was performed. Integer flooring was applied only to
integer recession directions.\hfill$\square$

\paragraph{Encoding every relevant exponent collision and fixed anchor.}
Take any lift $d$ of cost at most $C_0$. Its total
coefficient mass is $N_0+2C(d)\le N$. Keep every
coefficient and every structure-constant position fixed.
For each occurrence of a positive exponent in a symmetric
pair $a(q^e+q^{-e})$, introduce one positive integer
variable $x_\nu$, and replace the pair's exponents by
$x_\nu,-x_\nu$. There are at most $\lfloor N/2\rfloor$
such variables: each pair uses coefficient mass at least
two. Central coefficients need no variable. Distinct
occurrences may use separate variables even if their
original numerical exponents agree.

Adjoin one anchor variable $y_r$ for each $r\in E$ and require
\[
 y_r=r\quad(r\in E).
\]
Thus the total number of variables is at most $n_*$.
There are $h$ such variables. These are affine integer equalities with constants
bounded by $K$; they will not be rescaled. Every recession
direction has zero anchor coordinates, so the flooring
step preserves them literally.

Form the finite collection $\mathcal L$ of all linear
forms $0$, all positive and negative individual variables,
and all sums of two such signed variables. Repetitions
such as $x_\nu+x_\nu$ are allowed. Each form has coefficient
$\ell^1$ norm at most two. It includes every exponent
appearing in a product of two structure constants,
as well as every single exponent and every fixed trace
weight $\pm r$ through its anchor.

Use the original integer exponent assignment to record
the complete order pattern of these forms. If two forms
have the same value, impose their equality. If one is
larger, impose that their difference be at least one.
Also impose positivity of every variable and the fixed
anchor equations. These are finitely many integer linear
constraints, and the original assignment is an integer
solution. Each non-anchor row has $\ell^1$ norm at most
four and right side zero or one; an anchor row has norm
one and right side at most $K$ in absolute value.
Consequently the small-integer-point lemma supplies
another solution with all coordinates at most (G029.1).

This construction preserves the exact order pattern,
not only a set of selected equalities. In particular it
prevents previously different sums from colliding and
prevents a non-anchor monomial from newly occupying a
fixed trace weight. The constraints also retain equalities
between distinct occurrences of the same original weight.

\paragraph{Why the new exponent assignment is another lift of the same cost.}
Use the small solution to redefine the monomial exponents,
leaving their coefficients unchanged. Positivity of the
variables keeps each symmetric pair away from exponent
zero; their negatives still give bar invariance. The
single-exponent order constraints keep distinct monomials
within each coefficient distinct. Thus each coefficient
list, each nonzero table position, and the total mass are
unchanged.

Domination and parity at the prescribed trace weights
are unchanged, because equality or inequality with every
anchor $\pm y_r$ and with zero was preserved and the
anchors still equal the original integers $r$.
Every monomial outside the fixed trace support stays
outside it and retains its even coefficient. The specified
unit has only its unchanged constant coefficient, so
unitality is preserved.

For associativity, fix three basis elements and an output
basis element. Each side of
\[
 \sum_h d_{ij}^h d_{hk}^{\ell}
       =\sum_h d_{jk}^h d_{ih}^{\ell}
\]
is a sum of monomial products. Their numerical coefficients
are products of the original unchanged coefficients.
Their exponents are forms in $\mathcal L$. The original
identity says that the coefficient totals agree in every
exponent-collision class. Exactly the same classes occur
after compression, so every coefficient total still
agrees. This verifies all associativity identities.
The trace coefficient norm is fixed and the dimension
mass is unchanged, hence the cancelling-pair cost is
exactly the same. This proves (G029.1).

\paragraph{The finite algorithm and its exact scope.}
There are only finitely many arrays of nonnegative integer
Laurent coefficients with total mass at most $N$ and
exponents in $[-B,B]$. Enumerate them, impose bar symmetry,
the specified unit, coefficient domination and parity,
and test associativity by exact finite polynomial
arithmetic. Optional prescribed zero positions are
retained throughout. The known incumbent ensures the
accepted set is nonempty. Minimize the cost over it.

Any better unbounded-degree lift would compress to an
accepted lift of the same cost in this finite set.
Therefore the computed minimum is the unrestricted
symmetric-lift minimum for the chosen basis and unit,
not merely a degree-window minimum. The search bound is
immense: neither polynomial time nor any useful running
time estimate follows from this argument.

For the numerical specialization, NN-G011 proves the
absolute nonunit trace mass $4424$, and the nineteen
unit positions add $19$, giving $N_0=4443$.
The supplied trace coefficient table has maximum
Laurent exponent five. The checked NN-G027 lift has
mass $11128+19=11147$. Hence $n_*\le5578$ and
$n_*(K+1)4^{n_*-1}\le33468\,4^{5577}<2^{11170}$.
The small certificate
\path{verify_degree_cutoff_constants.py} checks
these input norms and constants; it does not perform
the enormous finite search or substitute for the proof.

\paragraph{Sparse anchors give polynomial-size certificates without an incumbent.}
The anchor set need only contain the positive exponents
actually present in the trace, together with one; all
other trace coefficients are zero. With
\[
 E=\{1\}\cup\{e>0:c_{ij,e}^k\ne0\text{ for some }i,j,k\},
 \qquad h=|E|,
\]
the proof uses $h$ anchor variables, not $K$ variables.
The same small-point lemma applies with right-hand-side
bound $K=\max E$, and gives
\[
 n\le\lfloor N/2\rfloor+h,\qquad
 \log_2 B=O(N+\log K).                              \tag{G029.7}
\]
Indeed $h\le\lfloor N_0/2\rfloor+1\le N+1$ whenever
the trace norm is at most the allowed dimension mass.
Preserving comparison with every exponent in $E$ is
sufficient: a monomial at a nonzero trace weight stays
there, and a monomial outside all those weights remains
outside the nonzero trace support. Adding the weight-one
anchor fixes the marked character used in the geometric
interpretation. The rank-three numerical bound is unchanged,
since its positive trace exponents are exactly $1,\ldots,5$.

This proves a uniform NP-membership statement. The input
is a finite based trace table encoded sparsely: rank and
basis indices, integer coefficients and signed exponents
are written in binary, and a \emph{total coefficient-mass
budget $N$} is written in unary. Equivalently one may
consider an input regime in which $N$ is polynomially
bounded by the input length. The decision problem asks
whether a symmetric effective associative unital lift
of total mass at most $N$ exists. Trace-algebra identities
and normalization of its sparse coefficient lists can
be checked as part of the input verification.

No known incumbent is needed for this decision result.
If the answer is yes, apply the compression proof to
that yes-witness with the prescribed budget $N$.
A unit on rank $r$ already uses at least $2r-1$ coefficient
units, so a yes-instance has $r\le(N+1)/2$. The trace
norm must also be at most $N$; otherwise the answer is
immediately no. These observations make both the rank
and the sparse anchor count polynomially bounded.

The compressed witness has at most $N$ nonzero monomials.
Their coefficients have $O(\log N)$ bits, their table
positions have $O(\log r)$ bits, and (G029.7) gives
$O(N+\log K)$ bits per exponent. Sorting the monomials
inside each table coefficient gives a canonical sparse
encoding of polynomial length. A verifier checks total
mass, bar symmetry, the unit, domination and parity by
sparse comparisons with the input table. It checks all
associativity identities by multiplying sparse polynomials,
adding their binary exponents, and collecting equal sums.
There are polynomially many such products and terms in
$r,N$; all coefficients and exponent sums have polynomial
bitlength. Thus the full verification is polynomial in
the input and witness lengths, proving membership in NP.

This is a certificate theorem for the existence of an
effective associative lift under a mass budget. It is
not NP-hardness, a polynomial-time algorithm, or PH1 for
actual plethysm or Kronecker multiplicities. With a binary
budget that may be exponentially large, the stated
polynomial certificate bound does not follow.

\paragraph{Geometric translation and the reduced optimization gate.}
The exponent variables are characters in a common finite
presentation of a diagonalizable group. More explicitly,
start with a free character group on these symbols and
one marked symbol $e_0$. Impose the homogeneous character
relations corresponding to the exponent collisions and
the anchor relations $y_r-r e_0=0$. An integral cocharacter
to this diagonalizable group, with value one on $e_0$,
assigns the integer grading weights. The strict order
conditions select an open rational polyhedral region
in that affine cocharacter slice. The original and the
compressed solutions are two integral cocharacters in
the same finite collision region, both fixing the marked
trace weights.

SGA 3, I, 4.7.2--4.7.3 identifies these character
presentations with diagonalizable-group comodules.
Pulling back their finite multiplicity objects along
the two cocharacters preserves the displayed finite
multiplication relations; NN-G010 and archive cote 2,
scan 17 then give the corresponding finite-kernel
realizations. Thus the actual conclusion is existence
of an equal-cost realization with bounded grading
weights and the same specified trace table. No equivalence
of the two categories or canonical choice of grading
is asserted.

The reassignment is not a homomorphism of the original
whole weight lattice $\mathbb Z$ to itself, and is not
a single Veronese scaling: the weight-one anchor remains
one. It preserves the finite character relations needed
for this multiplication presentation. The unrestricted
symmetric-lift optimum is now a finite effective search
problem, and unary-mass-budget existence has polynomial-size
verifiable certificates. Its value and any practical
optimization method remain open. Nothing here computes actual GCT multiplicities,
proves PH0/PH1, or transfers a nonstandard purity statement.



\section{NN-G030: associativity forces modulo-four constraints beyond trace parity}
\paragraph{Statement.}
For every integral associative closed rank-three table with
the specified trace parities, put $\sigma=D(D-1)$ and
$A_2=\delta_2\Lambda_2$, the actual second cell-product
coefficient. No constant-ratio, degree, or auxiliary-integrality
hypothesis is needed. Then the following are polynomial
congruences in $\mathbb Z[D]$:
\[
 S\equiv\sigma G_1,\qquad
 G_2-D\delta_2\equiv G_1,\qquad
 DG_2-A_2\equiv G_1\pmod4.                           \tag{G030.1}
\]
In fact
\[
 \begin{aligned}
 G_1&\in(4,2D),\\
 G_2,A_2&\in(4,2D,D^2),\\
 S&\in(4,2D^2(D-1)).
 \end{aligned}                                                       \tag{G030.2}
\]
Thus $G_1,G_2,S,A_2$ take values divisible by four
at every even integer $D$, and $S$ does so at every
integer $D$. In particular their ordinary dimension
masses at $D=4$ are multiples of four.
These necessary conditions use associativity and
integral trace parity alone, before effectivity or
domination is tested.

\paragraph{Proof.}
The supplied trace table gives
\[
 \begin{gathered}
 W\equiv D,\quad\delta_2\equiv\sigma,\quad
 G_1\equiv S\equiv0,\quad G_2\equiv D\sigma,\\
 U_2\equiv V_2\equiv D,\qquad A_2\equiv D^2\sigma
       \pmod2.
 \end{gathered}
\]
Write $G_1=2H$ and $S=2K$, and define the integral
polynomials
\[
 P=(G_2-D\delta_2)/2,\qquad
 Q=(DG_2-A_2)/2.
\]
The three necessary associativity identities reduce to
\[
 WK=HG_2,\qquad U_2K=G_2P,\qquad V_2K=G_2Q.
\]
Reducing the first modulo two and cancelling the
nonzero polynomial $D$ in $\mathbb F_2[D]$ gives
$\overline K=\overline H\sigma$.
The second and third then give
\[
 D\overline H\sigma=D\sigma\overline P
                       =D\sigma\overline Q.
\]
Cancelling $D\sigma\ne0$ gives
$\overline P=\overline Q=\overline H$.
Multiplying these statements by two proves (G030.1).
This does not require $\overline H$ itself to be nonzero.

At $D=0$, (G030.1) says
\[
 S(0)\equiv0,\quad G_2(0)\equiv G_1(0),\quad
 A_2(0)\equiv-G_1(0)\pmod4.
\]
The integers $G_1(0),G_2(0)$ are even, and $W(0)$
is even. Since $WS=G_1G_2$, their product at zero
is divisible by eight. If their common residue
modulo four were two, their product would instead
be four modulo eight. Hence $G_1(0),G_2(0),A_2(0)$
are all divisible by four.

Every coefficient of $G_1$ is even, so its divisible-by-four
constant term gives $G_1\in(4,2D)$. The parities
$G_2\equiv D^2(D-1)$ and
$A_2\equiv D^3(D-1)$ show that their linear
coefficients are even; their constant terms are
multiples of four. This gives their ideal membership
in (G030.2). Finally substitute $G_1=4a+2DB$
in the first congruence of (G030.1):
\[
 S\equiv2D^2(D-1)B\pmod4,
\]
which proves the last ideal membership. Evaluation
at even integers gives the claimed divisibilities.
For arbitrary integer $D$, the product $D^2(D-1)$
is even, proving the additional assertion for $S$.
\hfill$\square$

\paragraph{A low-mass scalar candidate that these constraints exclude.}
Ordinary mass parity and the scalar associativity
identities alone admit the following hypothetical data:
\[
 \begin{array}{c|rrrrrrrrrrrr}
 &w&g_1&g_2&s&d_1&d_2&u_1&u_2&v_1&v_2&a_1&a_2\\\hline
 \text{values}&10&70&90&630&4&12&6&6&30&48&10&24
 \end{array}
\]
Here lower-case symbols denote evaluations at $D=4$,
with $a_i=A_i(4)$. They obey all target mass lower
bounds and are even. Directly,
\[
 ws=g_1g_2,
 \qquad u_i s=g_i(g_i-4d_i),\qquad
 v_i s=g_i(4g_i-a_i),
\]
For example both $u_i s$ equal $3780$;
$v_1s=18900$ and $v_2s=30240$.
Their relaxed total mass would be $8522$, or
$(8522-4424)/2=2049$ pairs. But $g_1=70$, $g_2=90$
and $s=630$ are not multiples of four. Thus no integral
polynomial table with the required trace parity can
realize these scalar data. This is a strict exclusion
from the scalar relaxation, not an assertion that the
hypothetical data were an actual lower-cost lift.

The constraints do not say that every coefficient
of $G_1$ or $S$ is divisible by four. In the actual
$3352$ attainer the coefficient of $D$ in $G_1$
is six and the coefficient of $D^2$ in $S$ is
$-94$. Both are two modulo four. The polynomial
ideals in (G030.2), rather than coefficientwise
divisibility by four, are the proved statements.
The small exact checker verifies the displayed
congruences on that attainer and verifies all scalar
identities in the excluded mass example.

\paragraph{Geometric translation and algorithmic scope.}
At the level of integer congruences, individual effective
graded sign representations require dimension and signed
trace to agree modulo two.
Associativity of their finite-kernel composition imposes
the additional congruences (G030.1). In the symmetric
character ring $\mathbb Z[D]$, they say that the listed
coefficient characters vanish in the explicit quotient
rings $\mathbb Z[D]/(4,2D)$,
$\mathbb Z[D]/(4,2D,D^2)$, and
$\mathbb Z[D]/(4,2D^2(D-1))$, as appropriate.
These are exact scheme-theoretic congruence statements
over the nonreduced characteristic-four base, not
purity or point-counting assertions.

The character interpretation is SGA 3, I, 4.7.2--4.7.3;
$D=4$ is evaluation at the identity grading character,
while $D=0$ is evaluation at $q=-1$.
The composing objects remain the NN-G010 kernel
realization of archive cote 2, scan 17. These conditions
provide inexpensive necessary tests before a bounded
search or NP verification, but are not sufficient for
an effective lift and do not determine its optimum,
canonicality, rank-four positivity, or PH0/PH1.



\section{NN-G031: a competitive closed table must have gluing mass at least eight}
\paragraph{Statement.}
Every effective associative symmetric closed table of the specified trace
with at most $3352$ cancelling pairs has $W(4)\ge8$.
This permits arbitrary cell degrees and contents and arbitrary
rational ratios and auxiliary parameters, subject only to
integrality and effectivity of the actual coefficients.
In particular the entire $W=D+2$ branch in this presentation is excluded, without
enumerating support locations or assuming a degree bound.

\paragraph{Mass notation and immediate bounds.}
Put $\psi=D(D-3)=q^2+q+q^{-1}+q^{-2}$, the absolute
cell trace, and use lower-case letters for evaluation at $D=4$.
For $A_i=\delta_i\Lambda_i$, write $a_i=A_i(4)$.
All these masses are positive even integers. The trace imposes
\[
 d_i\ge4,\quad u_i\ge4,\quad v_i\ge16,\qquad
 a_1\ge8,\quad a_2\ge16.
\]
The nonunit mass, with the unit contributions cancelled from
the pair count, is
\[
 T=s+20(g_1+g_2)+200w+20(u_1+u_2+v_1+v_2)
       +26(d_1+d_2)+14(a_1+a_2).                    \tag{G031.1}
\]
The proposed score cap is exactly $T\le11128$.
Associativity gives
\[
 ws=g_1g_2,\qquad u_i g_j=w(g_i-4d_i)\quad(i\ne j).
                                                               \tag{G031.2}
\]
Consequently $u_1u_2<w^2$. Since $u_i\ge4$, $w=4$
is impossible. Parity and domination give $w\ge4$ even,
so it remains to exclude $w=6$.

An effective symmetric polynomial of mass six dominating
$D=q+2+q^{-1}$ and congruent to it modulo two is uniquely
$W=D+2$: its surplus mass two must be central.
Similarly $u_i=4$ forces $U_i=D$, and $u_i=6$ forces
$U_i=D+2=W$.

\paragraph{A useful divisibility consequence of $U_i=D$.}
For this $W$, the identities are
\[
 U_1G_2=W(G_1-D\delta_1),\qquad
 U_2G_1=W(G_2-D\delta_2).                           \tag{G031.3}
\]
If $U_1=D$, reduction of the first at $D=0$ implies
$D\mid G_1$, and coprimality of the monic polynomials
$D,W$ implies $W\mid G_2$. Write $G_1=DR$, $G_2=WT$;
all quotients are integral. The second identity at $D=0$
then implies $T(0)=0$, because $W(0)^2=4$.
Thus
\[
 G_2=DW C_2,\qquad G_1=D\delta_1+D^2C_2.
\]
The parity $G_1\equiv0$ and
$\delta_1\equiv D(D-1)$ in $\mathbb F_2[D]$ gives
$C_2\equiv D-1$ after cancellation of $D^2$.
Therefore $g_2/24$ is an odd integer.
The same argument with indices interchanged shows that
$U_2=D$ implies
\[
 G_1=DW C_1,\qquad G_2=D\delta_2+D^2C_1,
 \qquad C_1\equiv0\pmod2,                          \tag{G031.4}
\]
where the last parity follows from $G_1\equiv0$.
Hence $48\mid g_1$. None of these quotient polynomials
is asserted to be effective.

\paragraph{The five possible $u$ pairs.}
Since $u_1,u_2$ are even and at least four, their product
less than $36$ leaves only
\[
 (4,4),\ (4,6),\ (6,4),\ (4,8),\ (8,4).
\]
Solving (G031.2) at $w=6$ gives
\[
 g_1=\frac{24(6d_1+u_1d_2)}{36-u_1u_2},\qquad
 g_2=\frac{24(u_2d_1+6d_2)}{36-u_1u_2}.             \tag{G031.5}
\]
In either pair containing eight, $g_1+g_2\ge576$,
so $20(g_1+g_2)>11128$ already.

For $(u_1,u_2)=(6,4)$, the actual identities with
$U_1=W,U_2=D$ yield
\[
 G_1=DW K,\qquad K=(\delta_1+\delta_2)/2.
\]
The polynomial $K$ is integral and effective, with
$K\ge\psi$ coefficientwise. Since $G_1\equiv0$ and
$DW\equiv D^2$ modulo two, $K$ has even coefficients.
Its coefficients at $q^{\pm1},q^{\pm2}$ are therefore
at least two. Hence $t=d_1+d_2=2K(4)\ge16$.
Formula (G031.5) gives $g_1=12t$ and
$g_2=8t+4d_2\ge8t+16$. Consequently
\[
 T\ge s+20(g_1+g_2)+1200
    \ge16t^2+432t+1520
    \ge12528>11128.
\]

For $(u_1,u_2)=(4,6)$, the identities instead give
$G_2=DW K$ with the same effective
$K=(\delta_1+\delta_2)/2\ge\psi$.
Now $G_2\equiv D^2(D-1)$ implies $K\equiv D-1$.
In Laurent coefficients its central coefficient is
odd and at least one, its coefficients at $q^{\pm1}$
are odd and at least one, and those at $q^{\pm2}$
are even and at least two. Therefore $K(4)\ge7$,
so $t\ge14$. The preceding symmetric scalar estimate,
including the remaining lower bounds in (G031.1), gives
\[
 T\ge16t^2+458t+2696
    \ge12244>11128.
\]

\paragraph{The remaining pair $(4,4)$.}
Here (G031.5) becomes
\[
 g_1=\frac{12}{5}(3d_1+2d_2),\qquad
 g_2=\frac{12}{5}(2d_1+3d_2),\qquad
 g_1+g_2=12t.
\]
Also
\[
 s=\frac{24}{25}(6t^2+d_1d_2)\ge\frac{144}{25}t^2.
\]
Using the remaining mass bounds in (G031.1),
\[
 T\ge\frac{144}{25}t^2+266t+2336.                  \tag{G031.6}
\]
At $t=24$ this exceeds $11128$, and it is increasing
for $t\ge0$. Thus $t<24$.

Write $d_1=2a,d_2=2b$. The integrality of $g_1,g_2$
gives $b=a+5k$ for some integer $k$. The stronger
conditions $48\mid g_1$ and $g_2/24$ odd show
that $a$ is even and $k$ is odd:
\[
 g_1=24a+48k,\qquad g_2=24a+72k.
\]
With $a,b\ge2$ and $d_1+d_2<24$, the only possibilities
are $(d_1,d_2)=(4,14)$ and $(16,6)$.
For completeness, $k\ge1$ forces $t=4a+10k\ge18$;
if $k\ge3$ it is at least $38$, while $k=1$ forces
$a=2$. If $k\le-1$, the inequality $b\ge2$ gives
$a\ge2-5k$; evenness improves this to $a\ge8$ when
$k=-1$, leaving $a=8$, and $k\le-3$ forces $t\ge42$.

For $(4,14)$ one has $(g_1,g_2)=(96,120)$ and
$\delta_1=\psi$. In (G031.4), write $C_1=2H$;
then $H$ is integral and $H(4)=g_1/48=2$.
Since both $U_i=D$, the first identity gives
\[
 \delta_1=W C_1-D C_2=2W H-D C_2.
\]
As $D\mid\delta_1$ and $W(0)=2$, this forces
$H(0)=0$. Thus $D\mid H$ over $\mathbb Z[D]$, and
$H(4)$ is a multiple of four, contradicting $H(4)=2$.

For $(16,6)$ one has $(g_1,g_2,s)=(144,120,2880)$.
The remaining identity gives
\[
 v_1=\frac{576-a_1}{20}.
\]
Since $v_1$ is even, $a_1\equiv16\pmod{40}$;
with $a_1\ge8$ it follows that $a_1\ge16$.
Hence $14a_1+20v_1=576+13a_1\ge784$.
The other terms of (G031.1) then imply
\[
 T\ge2880+5280+1200+160+572+784+544
       =11420>11128.
\]
This excludes the last branch and proves the statement.
\hfill$\square$

\paragraph{Geometric translation and scope.}
In the NN-G010 realization by graded sign-representation
kernels, $w$ is the ordinary dimension of the gluing
coefficient object. The theorem says that any such
closed-table realization costing at most $3352$
regular-sign pairs must use a gluing object of dimension
at least eight. A dimension-six gluing object has the
unique character $D+2$; its associative kernel products
would force the incompatible character congruences and
mass bounds proved above. This is an actual dimension
obstruction for the effective realization, not a statement
about Frobenius weights or point counts. The character
dictionary is SGA 3, I, 4.7.2--4.7.3; the composition
dictionary is archive cote 2, scan 17, as realized in
NN-G010. The result shrinks the remaining optimization
gate but gives no global optimum or canonical GCT basis.




\section{NN-G032: bounded-mass lift existence is fixed-parameter tractable}
\paragraph{Statement.}
For the symmetric effective associative lift problem of NN-G029,
let $N$ be a bound on the total coefficient mass, including
the unit, and let $L$ be the bit length of the sparse input
trace table, its rank and its specified unit and forbidden
positions, and the budget. There are a computable function
$F$ and an absolute constant $c$ such that lift existence
with mass at most $N$ is decidable in time
\[
                    F(N)(L+1)^c.                  \tag{G032.1}
\]
Thus the problem is fixed-parameter tractable in the numerical
mass budget. This is valid even when the budget is encoded
in binary; it is not a polynomial-time claim for unbounded
$N$. A lift can also be produced in the same type of bound.
The claim concerns the unrestricted symmetric table problem;
it does not identify that optimization with a canonical GCT
basis or imply any positive multiplicity formula.

The proof eliminates integer exponent variables symbolically.
The new consequence is the parameterized algorithm for this
lift problem; the integer-linear quantifier elimination used
below is a standard method, proved here self-containedly.
It never searches through all integers up to the largest input
trace exponent. Its constants can be extremely large, so it
does not make the $N=11147$ example a practical computation.

\paragraph{Finite discrete patterns and integer formulas.}
First canonically combine repeated sparse monomials, and
check the prescribed symmetry, unit and input format. These
operations take polynomial time in $L$. If the absolute trace
mass exceeds $N$, reject. A unital rank-$r$ lift contains at
least $2r-1$ unit coefficients, so reject also if $2r-1>N$.
Consequently every nonrejected input has $r\le(N+1)/2$,
at most $N$ nonzero trace monomials, and trace coefficient
absolute values at most $N$.

The following discrete information consequently has only
a number of possibilities bounded by a computable function
of $N$: the rank; the occupied trace positions and their
integer coefficients; the forbidden positions; and the
coefficient-and-position patterns of all possible effective
lift tables of total mass at most $N$.
For a lift pattern, every noncentral symmetric monomial pair
costs at least two units of mass. Assign a distinct formal
positive integer variable to each occurrence of such a pair,
and retain the central coefficients as fixed integers.
There are at most $\lfloor N/2\rfloor$ exponent variables.
Within each position, require that separate pairs have
distinct exponents. Coefficients are selected as part of
the finite pattern, so there are no unknown coefficient
products in what follows.

Let $k_1,\ldots,k_h$ be the distinct positive exponents of
the input trace. They are free integer parameters; their
binary magnitudes are not included in the discrete pattern.
One has $h\le\lfloor N/2\rfloor$.
As in NN-G029, form zero, the signed individual exponent
variables and trace parameters, and all sums of two signed
individual forms. Enumerate their possible total order
and equality patterns. Each chosen pattern is expressed
by a conjunction of linear equalities and strict
inequalities with integer coefficients bounded in terms
of $N$. Since the variables are integers, a strict
inequality is an ordinary weak inequality with an added
constant one.

For a fixed collision pattern, domination and parity are
comparisons of fixed coefficients, and every associativity
coefficient is a fixed integer sum of fixed integer products.
Discard a pattern if these tests fail, or if it violates a
unit or forbidden-position condition. Every remaining
pattern yields a formula
\[
 \exists x_1,\ldots,x_m\in\mathbb Z\quad
       \Phi(x_1,\ldots,x_m;k_1,\ldots,k_h),          \tag{G032.2}
\]
where $m\le N/2$ and $\Phi$ is a finite conjunction of
integer linear inequalities, including $x_i\ge1$.
The numbers of patterns, atoms and variables, and all
coefficients and constants in these formulas, are bounded
by functions of $N$ alone. The input trace exponents
appear only as free parameters $k_j$.

Conversely, a satisfying integer assignment to any retained
formula gives an actual lift: its monomials have the selected
positive exponents; its constants dominate the trace with
the correct parity; and all coefficient identities for
associativity hold by the selected collision partition.
Thus the finite disjunction of (G032.2) is equivalent to
the desired bounded-mass existence question. This step
does not require an incumbent or the numerical cutoff
from NN-G029.

\paragraph{Elementary elimination of one integer variable.}
We prove the needed elimination rule completely. Consider
any finite Boolean combination of integer linear inequalities
and congruences in one quantified integer $y$, other integer
variables, and free parameters. All coefficients and moduli
are fixed integers. Equalities are pairs of inequalities.
Negating an inequality changes it to an inequality with
constant shift one. A negated congruence modulo a fixed
positive $m$ is a disjunction over its other $m-1$ residues.
After distributing conjunctions and disjunctions, it is
enough to eliminate $y$ from one conjunction of inequalities
and congruences. This may greatly increase size, but the
increase depends only on the formula, not on parameter values.

Let $B\ge1$ be the least common multiple of the absolute
values of all nonzero coefficients of $y$ in that conjunction.
If none occurs, set $B=1$. Substitute $z=By$ and add
$z\equiv0\pmod B$.
An atom $ay+t\ge0$, where $t$ does not involve $y$,
becomes, for $a\ne0$,
\[
 \operatorname{sgn}(a)z+(B/|a|)t\ge0.
\]
A congruence $ay+t\equiv0\pmod m$ becomes
\[
 \operatorname{sgn}(a)z+(B/|a|)t
                   \equiv0\pmod{mB/|a|}.
\]
These transformations are equivalent under $z=By$:
the inequality is multiplied by a positive integer, and
$m\mid u$ is equivalent to $cm\mid cu$ for every
positive integer $c$. Atoms with $a=0$ are unchanged.
Conversely $z\equiv0\pmod B$ recovers an integer $y$.

All remaining atoms involving $z$ are therefore lower bounds
$z\ge\ell_i$, upper bounds $z\le u_j$, or congruences
$z\equiv r_a\pmod{m_a}$, where the bounds and residues
are integer linear forms in the other variables and free
parameters. Let $M$ be the least common multiple of the
fixed positive moduli. Enumerate $a=0,\ldots,M-1$ and
impose the compatibility atoms
\[
                      a\equiv r_b\pmod{m_b}
                 \quad\hbox{for every }b.
\]
These atoms exactly say that the chosen residue class
$z\equiv a\pmod M$ satisfies every congruence.

If either the lower-bound list or the upper-bound list is
empty, that residue class has an integer satisfying every
remaining one-sided bound. Thus no further condition is
needed in that case. If both lists are nonempty, enumerate
an index $i$ for a largest lower bound and a number
$b\in\{0,\ldots,M-1\}$, and impose
\[
 \ell_i\ge\ell_{i'}\ (\hbox{all }i'),\qquad
 \ell_i+b\equiv a\pmod M,\qquad
 \ell_i+b\le u_j\ (\hbox{all }j).
\]
These conditions are necessary and sufficient: the smallest
integer in the chosen residue class at least $\ell_i$ is
exactly $\ell_i+b$ for its unique such $b$.
The resulting finite disjunction contains no $z$ or $y$
and consists only of integer linear inequalities and
congruences. This proves elimination of one integer
existential quantifier, with an explicit finite procedure.

\paragraph{Runtime and witness recovery.}
Apply this elimination successively to all the exponent
variables in every retained formula (G032.2). Its initial
size, coefficients, moduli and number of variables depend
only on $N$. At every step the resulting size, coefficients
and moduli still depend only on $N$, even though they may
grow rapidly. The result is a finite Boolean combination
of integer linear inequalities and fixed-modulus
congruences in the input exponents $k_j$ alone.
Formula construction and elimination therefore cost some
computable function of $N$, independently of the binary
magnitudes of the $k_j$.

Substitute the given binary exponents into this formula.
There are only $F_0(N)$ atoms and coefficients, each of
bit length at most $F_0(N)$ for a computable $F_0$.
Integer addition, multiplication and division by these
fixed coefficients and moduli can all be carried out by
elementary bit algorithms in time bounded by
$F_1(N)(L+1)^2$. The finite initial parsing and sorting
contribute another fixed polynomial in $L$.
This proves (G032.1), with an absolute polynomial exponent.

One may also retain the selected disjunct at every elimination
step. Recover variables in reverse order, selecting a true
disjunct after substituting the free parameters and the variables
already recovered at that stage. With two-sided
bounds use the recorded integer $z=\ell_i+b$; with only
lower bounds choose the least compatible residue representative
above their maximum; with only upper bounds choose the
largest representative below their minimum; and with no
bounds choose the selected residue itself. Then set $y=z/B$.
The compatibility atoms ensure divisibility by $B$.
Repeating this construction yields the integer exponents
of a lift. The number of steps and every arithmetic
coefficient are bounded by functions of $N$, so recovery
obeys the same type of parameterized time and bit-length
bound. In particular every recovered bit length is
$O_N(L+1)$, since each recovery step uses a bounded number
of linear operations and fixed-modulus divisions with
coefficients depending only on $N$.\hfill$\square$

\paragraph{Geometric translation and limitations.}
For a fixed coefficient pattern, the grading weights lie
in the cocharacter lattice of a finitely presented
diagonalizable group, with the associativity collisions
imposing integral linear relations. Domination and the
chosen order of weights specify finitely many lattice
chambers. Eliminating exponent variables projects these
integer lattice conditions; the resulting congruences
record finite-index lattice obstructions. Their character
duals are finite diagonalizable groups such as $\mu_m$.
The actual conclusion is that, for bounded total dimension
mass, admissibility of the prescribed trace weights is
decidable by finitely many such linear inequalities and
congruences, computed uniformly from the mass budget.

This is an arithmetic statement about cocharacter lattices,
not a claim that inequalities define Zariski-open subsets
or that an effective lift is a canonical geometric object.
The grading dictionary is SGA 3, I, 4.7.2--4.7.3, and
the actual composing kernel realization is NN-G010,
using archive cote 2, scan 17. The algorithm gives no
polynomial bound on the mass budget needed by varying
GCT inputs, no practical small value of $F$, and no
rank-four positivity, purity, PH0 or PH1 theorem.




\section{NN-G033: an effective realization with 2744 cancelling pairs}
\paragraph{Statement.}
The same rank-three trace table has an effective associative
symmetric closed-table realization with exactly $2744$
cancelling pairs, including $1296$ in the bottom cell.
This improves the unrestricted symmetric realization upper
bound from $3352$. All actual structure coefficients are
integral; the second auxiliary $\Lambda$ remains rational.
No global optimum is asserted in this construction theorem.

\paragraph{Construction and proof of associativity.}
As throughout, $D=q+2+q^{-1}$. Put
\[
 C=D-1,\qquad H=2D(D-2),\qquad
 W=2D^2-7D+4=D+2(q^2+q^{-2}).
\]
Define the actual base coefficients by
\[
 \begin{aligned}
 \delta_1&=D(D-3),&\delta_2&=WC-H,\\
 G_1&=DH,&G_2&=DWC,&S&=D^2CH,\\
 U_1&=D,&U_2&=W,&V_1&=2D^2-3D,&V_2&=3W,\\
 A_1&=D^2(D-3),&A_2&=D^2WC-3DH,&
 \Lambda_i&=A_i/\delta_i.
 \end{aligned}                                                       \tag{G033.1}
\]
The symbol $\Lambda_i$ is only a rational coordinate
factor; the actual cell products use $\delta_i$ or
$A_i$, multiplied by nonnegative powers of $D$.
In particular $\Lambda_1=D$.

Here $H-\delta_1=DC$. Directly,
\[
 \frac{G_1G_2}{S}=W,\qquad
 G_1-D\delta_1=D^2C,\qquad
 G_2-D\delta_2=DH=G_1.
\]
It follows that
\[
 \frac{G_1(G_1-D\delta_1)}S=D=U_1,
 \qquad
 \frac{G_2(G_2-D\delta_2)}S=W=U_2.
\]
The remaining identities are
\[
 DG_1-A_1=DCV_1,\qquad DG_2-A_2=3DH,
\]
and therefore $G_i(DG_i-A_i)/S=V_i$ for both $i$.
Neither $\Lambda_i$ nor $D^2-\Lambda_i$ is the zero
rational function. For the second index this follows
explicitly from
\[
 D^2\delta_2-A_2=DH(3-D)\ne0.
\]
The injective generic block realization proved in
NN-G011 thus proves associativity over $\mathbb Q(D)$.
All actual coefficients are integral, so these identities
hold in $\mathbb Z[D]$. No regularity of the auxiliary
rational coordinate at every specialization is used.

\paragraph{Every coefficient condition and the exact score.}
A list $(a_0,\ldots,a_m)$ below denotes
$a_0+\sum_{j=1}^m a_j(q^j+q^{-j})$.
The base coefficient checks are completely recorded by
\[
\begin{array}{c|l|l}
 &\text{dimension coefficient}&\text{absolute trace target}\\\hline
 W&(2,1,2)&(2,1)\\
 U_1&(2,1)&(2,1)\\
 U_2&(2,1,2)&(2,1)\\
 \delta_1&(0,1,1)&(0,1,1)\\
 \delta_2&(0,1,1,2)&(0,1,1)\\
 G_1&(16,14,8,2)&(16,12,4)\\
 G_2&(18,17,13,7,2)&(14,11,5,1)\\
 S&(164,144,96,46,14,2)&(144,120,68,24,4)\\
 V_1&(6,5,2)&(6,5)\\
 V_2&(6,3,6)&(6,3,2)\\
 A_1&(2,3,3,1)&(2,1,1,1)\\
 A_2&(22,23,26,23,11,2)&(2,3,2,1,1)
\end{array}                                                         \tag{G033.2}
\]
Each left entry dominates the right entry coefficientwise
and has the same parity. Multiplication by the required
nonnegative powers of $D$ preserves parity and domination
of the absolute trace, by the coefficientwise triangle
inequality. This proves the conditions for the full
closed multiplication table. The included independent
integer-Laurent verifier checks all $132$ positions and
all $1000$ associativity triples with zero failures.

At $D=4$, the masses are
\[
 (w,u_1,u_2,v_1,v_2)=(8,4,8,20,24),\qquad
 (g_1,g_2,s)=(64,96,768),
\]
and $(d_1,d_2,a_1,a_2)=(4,8,16,192)$.
Hence the bottom mass is
\[
 768+20(64+96)+200\cdot8+20(4+8+20+24)=6688.
\]
The cell mass is $26(4+8)+14(16+192)=3224$.
The nonunit total is $9912$, giving
\[
 \frac{9912-4424}{2}=2744,\qquad
 \frac{6688-4096}{2}=1296.                         \tag{G033.3}
\]
Thus there are $1448$ pairs outside the bottom cell.
For comparison, the $3352$ realization had $2752$
bottom pairs and $600$ elsewhere. The new construction
reduces the total while moving more of the cancellation
into the two nonbottom cells; the objective is the total
over the specified presentation.

The ratios
\[
 \frac{G_1}{W\delta_1}
   =\frac{2D(D-2)}{W(D-3)},\qquad
 \frac{G_2}{W\delta_2}=\frac{D(D-1)}{WC-H}
\]
are nonconstant rational functions. Thus this realization
continues to lie outside the constant-ratio class of
NN-G026 and does not contradict its exact $4520$ optimum.
\hfill$\square$

\paragraph{A smaller complete search cutoff.}
Including the $19$ unit coefficients, the new total mass
is $N=9931=4443+2\cdot2744$. There are still five positive
trace anchors $1,\ldots,5$, so NN-G029 gives
\[
 n\le\lfloor9931/2\rfloor+5=4970,\qquad
 \deg d\le29820\cdot4^{4969}<2^{9953}.              \tag{G033.4}
\]
Accordingly, any equal-or-better unrestricted symmetric
realization has an equal-score representative below
this new cutoff, allowing extra zero-trace positions.
This improves the earlier $2^{11170}$ existence cutoff;
it does not bound the degree of every existing realization
or make the finite search practical.

\paragraph{Geometric translation and limitations.}
The effective Artin--Tate coefficient objects in
(G033.2), with the NN-G010 kernel-composition realization,
give an actual additive monoidal model of the same
graded sign trace. The excess of ordinary dimension
over absolute signed trace is precisely twice $2744$
regular-sign pairs in its chosen multiplication
presentation, with $1296$ pairs in the bottom cell.
The gluing character has dimension eight and includes
weight-two characters absent from $D$.

This construction concerns actual coefficient objects,
not formal $K_0$ cancellation or a count obtained by
evaluating a virtual class. Its class is governed by
SGA 3, I, 4.7.2--4.7.3 and the archive cote 2, scan 17
composition dictionary implemented in NN-G010.
The rational auxiliary $\Lambda_2$ does not represent
a fractional-dimensional object. The global minimum,
canonicality, rank-four extensions and PH0/PH1 remain
separate questions. The older exact subclass bounds
and the NN-G031 implication for all costs at most
$3352$ remain valid with their stated hypotheses.




\section{NN-G034: the exact optimum with one fixed first-cell realization is 2744}
\paragraph{Statement and class.}
In the symmetric closed-table presentation, impose only
\[
 \delta_1=D(D-3),\qquad G_1=2D^2(D-2),\qquad U_1=D.   \tag{G034.1}
\]
Allow all remaining actual coefficient polynomials to vary,
with arbitrary degrees and contents, and allow rational
auxiliary $\Lambda_i$ and nonconstant ratios.
Among effective associative lifts in this class, the exact
minimum is $2744$ cancelling pairs. There is exactly one
minimizing set of actual base polynomials, namely NN-G033.
This is an exact class theorem, not the unrestricted
$1000$-position optimum and not the optimum over all closed
tables with their first-cell data free.

\paragraph{Forced parameters and the first-cell optimum.}
Put $C=D-1$ and $H=2D(D-2)$, so $G_1=DH$.
The identities $U_1G_2=W(G_1-D\delta_1)$ and
$WS=G_1G_2$ force
\[
                  G_2=DWC,\qquad S=D^2CH.            \tag{G034.2}
\]
Thus $(g_1,g_2,s)=(64,12w,768)$ at $D=4$.
The first $V$ identity forces
\[
 A_1=D^2H-DCV_1,\qquad a_1=256-12v_1.
\]
Since $v_1$ is even and $a_1\ge8$, one has $v_1\le20$.
The first-cell contribution satisfies
\[
 20v_1+14a_1=3584-148v_1\ge624,
\]
with equality exactly when $v_1=20,a_1=16$.

There is exactly one allowable $V_1$ at this equality.
It dominates the trace character $5D-4=(6,5)$,
has the same parity, and has surplus mass four.
Since $A_1$ is effective and its first term $D^2H$
has degree four, $V_1$ has degree at most two.
The only possibilities are
\[
 5D,\qquad7D-8,\qquad5D-4+2(q^2+q^{-2})=2D^2-3D.
\]
The respective central coefficients of $A_1$ are
$-10,-6,2$. The last choice gives
$A_1=D^2(D-3)=(2,3,3,1)$ and satisfies all conditions;
the first two fail effectivity.

\paragraph{A complete mass reduction.}
For a competitor with score at most $2744$, the nonunit
mass is $T\le9912$. Formula (G031.1) and the preceding
first-cell bound give
\[
 \begin{aligned}
 T&\ge2856+440w+20u_2+20v_2+26d_2+14a_2,\\
 u_2&=\frac{w(3w-d_2)}{16},\qquad
 v_2=\frac{w(48w-a_2)}{64}.                         \end{aligned} \tag{G034.3}
\]
The omitted amount is exactly $148(20-v_1)\ge0$.
Using $u_2\ge4,v_2\ge16,d_2\ge4,a_2\ge16$ gives
$T\ge3584+440w$, hence $w\le14$.
NN-G031 rules out $w=4,6$ at this score.
Eliminating $u_2,v_2$ in (G034.3) gives
\[
 T\ge2856+440w+\frac{75}{4}w^2
       +\left(26-\frac{5w}{4}\right)d_2
       +\left(14-\frac{5w}{16}\right)a_2.           \tag{G034.4}
\]
For $w=12,14$ both final coefficients are positive,
and even the first three terms are at least
$10836>9912$. Thus only $w=8$ or $w=10$ remain.

Write $T_j=q^j+q^{-j}$, an integral polynomial in $D$.
Every effective symmetric $W\equiv D\pmod2$
dominating $D$ has the following exhaustive small-mass forms:
\[
 \begin{array}{c|l}
 w=8&D+4\quad\hbox{or}\quad D+2T_j\quad(j\ge1),\\
 w=10&D+6\quad\hbox{or}\quad D+2+2T_j\quad(j\ge1).
 \end{array}                                                       \tag{G034.5}
\]
Indeed the even surplus has mass four or six; a noncentral
surplus pair costs four, leaving only the displayed central
remainder. Each polynomial in (G034.5) is primitive in
$\mathbb Z[D]$, since its Laurent coefficients have gcd one.

\paragraph{An equal-mass divisibility lemma.}
Suppose $W,U$ both belong to the same row of (G034.5),
and $U=WR$ for $R\in\mathbb Z[D]$ with $R\equiv1\pmod2$.
Then $U=W$.
If $\deg W\ge2$, its leading coefficient is two.
A nonconstant quotient would require $\deg U\ge2$,
also with leading coefficient two; thus $R$ would have
leading coefficient one, contradicting $R\equiv1$
in positive degree. A constant quotient has $R(4)=1$
and is one. For the degree-one choice with leading
coefficient three, a nonconstant integral quotient would
have to have leading coefficient $2/3$, impossible.

It remains to consider $W=D+\ell$ with $\ell=4$ or six.
A nonconstant quotient requires
$U=D+\ell-4+2T_j$ with $j\ge2$.
At $D=-\ell$, divisibility would give $T_j(-\ell)=2$.
This is impossible. To see it directly, set $x=\ell+2\ge6$;
the positive sequence
$a_j=(-1)^jT_j(-\ell)$ has $a_0=2,a_1=x$ and
$a_{j+1}=xa_j-a_{j-1}$. Induction gives
$a_{j+1}>a_j>2$ for $j\ge1$.
This completes the lemma without bounding $j$.

\paragraph{The complete $w=8$ branch.}
Each $W$ in its row of (G034.5) is coprime over
$\mathbb Q[D]$ to $D(D-2)$: at $D=0$ its value is
$4$ or $\pm4$, and at $D=2$ its value is among
$6,2,-2$. The two second-cell identities therefore
force
\[
 U_2=WR,\qquad V_2=WZ\qquad(R,Z\in\mathbb Z[D]).
\]
Here polynomial divisibility follows first over
$\mathbb Q[D]$ and then integrally by primitivity.
Cancelling $D$ modulo two gives $R\equiv Z\equiv1$.
Formula (G034.3) shows $4\le u_2\le10$, so the
positive odd integer $R(4)=u_2/8$ must be one.
The equal-mass lemma gives $U_2=W$, and then
\[
                 \delta_2=WC-H,\qquad d_2=8.         \tag{G034.6}
\]
Likewise $16\le v_2\le46$, so $Z(4)$ is three or five.
The corresponding $a_2$ are $192$ and $64$.

For $W=D+4$, (G034.6) has negative leading coefficient,
so it is not effective. For $W=D+2T_j$, one has
\[
             \delta_2=-D(D-3)+2(D-1)T_j.
\]
Domination of $D(D-3)$ forces $j\le2$: when $j\ge3$
the product $(D-1)T_j$ has no coefficient at $q^1$,
whereas $D(D-3)$ does. Thus only $W=3D-4$ and
$W=D+2T_2$ remain.

The second $V$ identity is now
\[
              A_2=D^2\bigl(WC-2(D-2)Z\bigr).         \tag{G034.7}
\]
If $W=3D-4$, effectivity of $V_2=WZ$ makes the leading
coefficient of a nonconstant $Z$ positive. Its parity
makes that coefficient even. For $\deg Z=1$,
the leading coefficient inside (G034.7) is at most
$3-4<0$, and higher degrees also give a negative
leading coefficient. Thus $Z$ must be constant.
But then $V_2$ has degree one and cannot dominate
the trace coefficient two at $q^2$. This excludes
the degree-one choice.

For $W=D+2T_2=(2,1,2)$, the same leading-coefficient
argument gives $\deg Z\le1$: degree two would have
leading coefficient at least two, against the leading
coefficient two of $WC$.
Write $Z=(z-2a)+a(q+q^{-1})$, where $z=Z(4)\in\{3,5\}$.
Here $a\ge0$ is even, while the central coefficient
is allowed to be negative. The central coefficient
of $V_2=WZ$ is $2z-2a$ and must be at least six.
If $z=3$, this forces $a=0$ and $Z=3$.
If $z=5$, it permits only $a=0,2$.
The central coefficient of $A_2$ is
$70-16z+4a$, which is negative in both $z=5$ cases.
Consequently only $Z=3$ survives. Its $a_2=192$
gives, from (G034.4),
\[
                  T\ge7576+16\cdot8+\frac{23}{2}\cdot192
                    =9912.                         \tag{G034.8}
\]
Equality forces the unique first-cell optimum already
proved, and gives exactly the parameters of NN-G033.

\paragraph{The complete $w=10$ branch: scalar constraints.}
Here (G034.3) reads
\[
 u_2=\frac58(30-d_2),\qquad
 v_2=\frac5{32}(480-a_2).
\]
Evenness and positivity force $d_2\equiv14\pmod{16}$
and $d_2<30$, hence $d_2=14,u_2=10$.
They also force $a_2\equiv32\pmod{64}$.
The choice $a_2=32$ gives the lower bound $T\ge9668$,
whereas the next choice $a_2=96$ already gives
$T\ge10364>9912$. Thus a competitor must have
\[
                  d_2=14,\quad u_2=10,\quad
                  a_2=32,\quad v_2=70.               \tag{G034.9}
\]

\paragraph{Coprime polynomials in the $w=10$ branch.}
No $W$ in this row is divisible by $D$.
It is divisible by $D-2$ exactly when
$W=D+2+2T_j$ with $j\equiv2\pmod4$, since
$T_j(2)$ is zero for odd $j$ and
$2(-1)^{j/2}$ for even $j$.
First suppose $W$ is coprime to $D(D-2)$.
Then $U_2=WR,V_2=WZ$ integrally, with
$R\equiv Z\equiv1$, and $R(4)=1$.
The equal-mass lemma gives $U_2=W$ and
$\delta_2=WC-H$.

The choice $W=D+6$ again gives a negative leading
coefficient in $\delta_2$. For
$W=D+2+2T_j$, domination of $\delta_2$ requires
\[
                       (D-1)(1+T_j)\ge D(D-3).
\]
The coefficient at $q^2$ forces $j\le3$.
The coprime possibilities are $j=1$ and $j=3$.
For $j=1$, $W=3D-2$ and the same leading-coefficient
argument in (G034.7) forces $Z$ constant, so $V_2$
cannot dominate its required $q^2$ coefficient.

For $j=3$, effectivity in (G034.7) forces
$\deg Z\le2$: a degree-three leading coefficient
would be positive even and produce a negative leading
term, and larger degrees are also impossible.
Write $Z=(z_0,z_1,z_2)$ in symmetric Laurent notation.
By (G034.9), $Z(4)=7$.
The coefficient of $q^2$ in $D^2WC$ is $54$.
Since $G_1=(16,14,8,2)$, the coefficient of $q^2$
in $G_1Z$ is
\[
        8z_0+16z_1+16z_2=8Z(4)=56.
\]
Thus $A_2=D^2WC-G_1Z$ has coefficient $-2$ at $q^2$,
contradicting effectivity. This excludes every coprime
$w=10$ choice, including all unbounded locations $j$.

\paragraph{The $D-2$ factor in the $w=10$ branch.}
It remains to consider $W=D+2+2T_j$ with
$j\equiv2\pmod4$. Write $W=(D-2)B$.
The polynomial $B$ is primitive, $B(0)=-3$, and
$B(2)=1$: for even $j$, $T_j$ is an even polynomial
in $D-2$, so $W'(2)=1$.
Thus $B$ is coprime to $D(D-2)$; modulo two,
$B\equiv1$. The identities now force
\[
 U_2=BR,\quad V_2=BZ,\qquad
 R\equiv Z\equiv D\pmod2,
 \qquad R(4)=2,\quad Z(4)=14.                       \tag{G034.10}
\]
The leading coefficient of $B$ is two and
$\deg B=j-1$. If $R$ were constant, its parity
would contradict $R\equiv D$.
Otherwise $\deg U_2\ge j\ge2$, so the mass-ten
classification gives leading coefficient two for $U_2$.
Hence $R$ has leading coefficient one, and its parity
forces $\deg R=1$. The only mass-ten polynomial of
degree $j$ in (G034.5) is $W$ itself. Therefore
$U_2=W$, $R=D-2$ and $\delta_2=WC-H$.
The preceding $q^2$ domination argument gives $j\le3$;
thus necessarily $j=2$.

We have reached
\[
 W=2D^2-7D+6=(D-2)(2D-3),\qquad B=2D-3.
\]
The second $V$ identity reduces to
\[
                    A_2=D^2(WC-2Z).                 \tag{G034.11}
\]
Effectivity forces $\deg Z\le2$. Indeed its leading
coefficient is positive when nonconstant; degree three
would make it positive even by (G034.10), exceeding
the leading coefficient two of $WC$ after the factor
two in (G034.11), and higher degrees also fail.
Write $Z=(z_0,z_1,z_2)$. Its parity says that $z_2$
is even, $z_1$ is odd, and $z_0$ is even; its mass says
$z_0=14-2z_1-2z_2$.
The relevant coefficients of (G034.11) are
\[
 \begin{array}{c|l}
 q^1&-25+2z_1+8z_2,\\
 q^2&32-4z_1-8z_2,\\
 q^4&11-2z_2.
 \end{array}
\]
Domination of their trace targets $3,2,1$ implies
\[
 z_1+4z_2\ge14,\qquad z_1+2z_2\le7,\qquad z_2\le5.
\]
The first two inequalities force $z_2\ge7/2$.
Evenness then gives $z_2=4$, and oddness gives
$z_1=-1$, hence $z_0=8$.
But the central coefficient of $V_2=BZ$ is then
$z_0+4z_1=4$, below its required value six.
This excludes the final $w=10$ branch.

The $w=8$ equality realization of NN-G033 is therefore
the unique minimizer in (G034.1), proving the theorem.
\hfill$\square$

\paragraph{Geometric translation and remaining gate.}
Fixing (G034.1) fixes the first-cell coefficient object,
its bottom action character, and one gluing product in
the NN-G010 category of graded sign-representation kernels.
Among all effective associative completions in this closed
presentation, the least total excess of dimension over
absolute trace is exactly $2\cdot2744$, achieved by the
single set of coefficient characters of NN-G033.
The proof allows arbitrary grading locations in the other
objects; the mass-ten $D-2$ factor is treated by an actual
integral divisibility argument, not discarded as generic.

The finite-character evaluations $D=0$ and $D=2$
are restrictions at $q=-1$ and $q=i$, under SGA 3,
I, 4.7.2--4.7.3. They are not assertions about Frobenius
eigenvalues. The monoidal composition dictionary is
archive cote 2, scan 17, via NN-G010.
Any smaller closed-table score must change at least one
of the three first-cell data in (G034.1); other table
shapes remain outside this class. This is a strict
remaining gate, with no canonical-basis, rank-four,
PH0/PH1 or purity conclusion.




\section{NN-G035: a complete first-cell mass gate for competitive closed tables}
\paragraph{Statement.}
Every symmetric closed-table effective associative lift with
at most $2744$ cancelling pairs lies in the complete necessary
mass relaxation below. It has $w\in\{8,10,12\}$ and $s\ge588$.
Its first-cell mass triple $(g_1,d_1,u_1)$ belongs to the
following set of $21$ triples, grouped by $g_1$:
\[
\begin{array}{c|l}
 g_1&\text{possible }(d_1,u_1)\\\hline
 64&(4,4)\\
 72&(4,6),(6,4)\\
 76&(4,6)\\
 80&(4,8),(8,4),(8,6),(12,4)\\
 84&(6,6),(10,6),(14,4)\\
 88&(8,8)\\
 96&(4,10),(4,12),(6,6),(12,4),(14,6)\\
 100&(4,10),(6,10),(10,6)\\
 128&(12,8)
\end{array}                                                         \tag{G035.1}
\]
The full relaxation has $43$ base rows and $467$ possible
auxiliary increment pairs, all explicitly listed in the
certificate. Its minimum relaxed score is $1700$.
Consequently the minimum over all symmetric closed tables
is between $1700$ and $2744$.
Surviving mass tuples are not asserted realizable; no
Laurent degree or support-location bound is inferred.
This theorem does not apply to the broader unrestricted
$1000$-position problem of NN-G029.

\paragraph{Complete finite bounds.}
Use the actual-coefficient mass notation of (G031.1).
The known lower bounds are
\[
 s\ge576,\quad g_i\ge48,\quad d_i\ge4,\quad
 u_i\ge4,\quad v_i\ge16,\quad a_1\ge8,\quad a_2\ge16.
\]
All are even; NN-G030 also makes $g_1,g_2,s,a_2$
multiples of four. The score cap is $T\le9912$.
NN-G031 gives $w\ge8$, and the mass formula implies
\[
              20(g_1+g_2)\le9912-3520,
              \qquad g_1+g_2\le316.                 \tag{G035.2}
\]
Thus each $g_i\le268$.
If $w\ge20$, the relation $g_1g_2=ws\ge576w$
and the arithmetic-geometric mean inequality give
$g_1+g_2\ge48\sqrt{20}>214$.
Since the sum is divisible by four, it is at least
$216$. The mass formula would then imply
\[
 T\ge576+20\cdot216+200\cdot20+800+208+336
       =10240>9912.
\]
Hence $w$ is one of $8,10,12,14,16,18$.

Fix one such $w$ and multiples of four $g_i\ge48$
with sum at most $316$. Then $s=g_1g_2/w$ must be
an integer multiple of four at least $576$.
For each index $i$, with $j\ne i$, enumerate the
even integers $4\le d_i\le\lfloor g_i/4\rfloor$,
retaining exactly those for which
\[
             u_i=\frac{w(g_i-4d_i)}{g_j}
\]
is an even integer at least four. These finite bounds
contain every possible actual cell mass.

The auxiliary masses satisfy
$v_i=w(4g_i-a_i)/g_j$. Set $t_1=2,t_2=4$, and
\[
 M_i=\operatorname{lcm}\left(t_i,
                   \frac{2g_j}{\gcd(w,2g_j)}\right).
\]
Their necessary congruence is exactly
$a_i\equiv4g_i\pmod{M_i}$, together with
$a_1\ge8,a_2\ge16$.
Let $a_i^0$ be the least integer satisfying this lower
bound and congruence, and put
$v_i^0=w(4g_i-a_i^0)/g_j$. If $v_i^0<16$, discard
the tuple, since increasing $a_i$ only decreases $v_i$.
Otherwise every actual auxiliary mass is
\[
 a_i=a_i^0+M_i k_i,\qquad
 v_i=v_i^0-\frac{wM_i}{g_j}k_i,\qquad k_i\ge0.       \tag{G035.3}
\]
The mass increment per unit $k_i$ is
\[
 \Delta_i=14M_i-20\frac{wM_i}{g_j}>0,
\]
because $w\le18$ and $g_j\ge48$ make
$14-20w/g_j\ge13/2>0$.
Thus the least auxiliary masses give a rigorous lower
bound $T_0$ for the tuple, and all remaining choices
satisfy the explicit finite inequalities
\[
 \Delta_1k_1+\Delta_2k_2\le9912-T_0,\qquad
 \frac{wM_i}{g_j}k_i\le v_i^0-16.                   \tag{G035.4}
\]

\paragraph{The extra $U_i=D$ tests for masses eight and ten.}
For $w=8,10$, (G034.5) shows that $W$ is primitive
and coprime to $D$.
If $u_1=4$, then $U_1=D$. The proof of (G031.4)
gives $G_2=DWC$ for an integral $C\equiv D-1\pmod2$;
therefore $g_2/(4w)$ is an odd integer.
If $u_2=4$, the analogous quotient is even, so
$8w\mid g_1$. These tests use only the stated mass
of $W$ and integral associativity, not a constant-ratio
or a degree assumption.

\paragraph{Why the minimal bottom polynomial is impossible.}
Suppose $s=576$. Coefficientwise domination and equality
of masses force
\[
                  S=4D^2(D-1)^2.                    \tag{G035.5}
\]
Write $C=D-1$ temporarily. From $WS=G_1G_2$ and
$G_1\equiv0,G_2\equiv D^2C\pmod2$, reduction after
dividing by two shows that $G_1$ is divisible by four.
Writing $G_1=4J$ gives
\[
             JG_2=WD^2C^2,\qquad J\equiv DC\pmod2.
\]
Since $W\equiv D\pmod2$, $W(1)$ is odd, so $C$
does not divide $W$. Neither $J$ nor $G_2$ can be
divisible by $C^2$, by their displayed reductions
modulo two. Their product contains exactly $C^2$;
therefore each contains exactly one factor $C$.
Write
\[
          J=CA,\qquad G_2=CB,\qquad
          AB=WD^2,\qquad A\equiv D, B\equiv D^2\pmod2.
\]
The actual $D$-adic orders of $A,B,W$ are at most
$1,2,1$, respectively, again by these nonzero reductions.
If $D\nmid W$, their first two orders sum to two;
if $D\mid W$, they sum to three. In either case
at least one is maximal: either $A=DA_1$ with
$A_1\equiv1\pmod2$, or $B=D^2B_1$ with
$B_1\equiv1\pmod2$.
Thus at least one of $g_1,g_2$ is $48$ times a
positive odd integer.

It cannot equal $48$: using $g_1g_2=576w$ in
the corresponding $u_i$ identity gives
$u_i=4-d_i/3<4$.
Hence some $g_i\ge144$.
For that index, eliminate $v_i$ using $s=576$:
\[
 20v_i+14a_i=\frac{5g_i^2}{36}
                   +\left(14-\frac{5g_i}{144}\right)a_i.
\]
The second coefficient is positive because $g_i\le268$.
Since $a_i\ge8$, the right side is at least
$5g_i^2/36+112-5g_i/18$, an increasing function for
$g_i\ge144$, and hence at least $2952$.
For the other index $20v_j+14a_j\ge432$.
Also $g_1+g_2\ge192$ and $w=g_1g_2/576\ge12$.
Together these estimates imply
\[
 T\ge576+3840+2400+160+208+2952+432
       =10568>9912.
\]
This contradiction excludes every minimal-bottom case;
no finite computation was used for this polynomial step.

\paragraph{Exact enumeration and completeness of the certificate.}
The verifier enumerates precisely the finite ranges in
(G035.2)--(G035.4), using integer gcd and lcm arithmetic.
It first tests $s$, the lower mass bound, and the two
$U_i=D$ divisibility tests, then enumerates every $d_i$
and uses the least admissible $a_i^0$.
There are $1445$ coarse $(w,g_1,g_2)$ tuples before
the cell tests and $55$ surviving base rows with
$T_0\le9912$. Exactly $12$ have $s=576$ and are
excluded by the preceding proof. The remaining $43$
rows have exactly $467$ increment pairs (G035.4).
The certificate lists every row and every pair.
It explicitly includes the actual NN-G033 mass tuple.

Projecting these rows gives (G035.1), $w\in\{8,10,12\}$,
and least $s=588$. Their least $T_0$ is $7824$, giving
the relaxed lower bound $(7824-4424)/2=1700$.
Every actual realization at or below the incumbent
maps to one of these rows and increments by the
proved equations and bounds. Larger scores are already
above $1700$. This proves the stated closed-table
lower bound and all finite necessary restrictions.
\hfill$\square$

\paragraph{Geometric translation and scope.}
For the actual graded sign-representation kernels of
NN-G010, this theorem bounds the possible ordinary
dimensions of the first-cell coefficient object, its
bottom action object, and its correction product by
the $21$ listed triples. Their grading locations are
still unrestricted. The bottom coefficient object
cannot have exactly the absolute trace character
$4D^2(D-1)^2$ in any competitive realization; integral
character factorization and parity force a dimension
contradiction. The finite list records necessary
dimension data, not actual objects for every row.

The dictionary is SGA 3, I, 4.7.2--4.7.3 and archive
cote 2, scan 17 via NN-G010. No equality of formal
$K_0$ classes is treated as an effective realization.
The result shrinks the closed-table search to explicit
mass data, while NN-G029--NN-G032 separately address
the wider unrestricted presentation. Neither statement
supplies a canonical GCT basis, rank-four positivity,
or PH0/PH1.




\section{NN-G036: the exact unique 2744-pair optimum for the symmetric closed table}
\paragraph{Statement.}
The minimum cancelling-pair score in the specified symmetric
closed-table presentation is exactly
\[
                                2744.                \tag{G036.1}
\]
The actual polynomial parameter realization attaining it is
unique within that presentation and is NN-G033.
All coefficient degrees and contents, all grading support
locations, and all rational ratios and auxiliary parameters
are allowed, subject to integrality and effectivity of the
actual structure constants. This is not a numerical lower
bound for unrestricted extra-position multiplication tables.

The proof retains its checked stages in their original order.
The first stage forces every competitive $g_1=64$ table into
NN-G034, which has the unique $2744$ attainer. The remaining
ordinary-dimension rows are then eliminated by uniform
integral polynomial arguments. Intermediate intervals and
row counts below describe successive stages, not the final
gate. The last proof and complete projection leave no other
competitive row.

\paragraph{Uniform irreducibility of the small gluing characters.}
Retain $T_j=q^j+q^{-j}$ as an integral polynomial in $D$.
Every mass-eight character in (G034.5) is irreducible
over $\mathbb Q[D]$.
Every mass-ten character there is irreducible except
\[
 W=D+2+2T_j\quad(j\equiv2\pmod4),
\]
in which case
\[
               W=(D-2)B,\qquad B\text{ primitive and irreducible},
               \qquad B\equiv1\pmod2,\quad B(4)=5.   \tag{G036.2}
\]
These statements include all support indices $j$.

Here is an elementary proof. The degree-one choices
are immediate. Every remaining $W$ is primitive,
has leading coefficient two, and reduces to $D$
modulo two. In any nontrivial primitive integral
factorization, one factor has leading coefficient
$\pm1$ and the other $\pm2$. The first factor's
degree is unchanged upon reduction modulo two;
since the product of the reductions is $D$, it
must have degree one. Up to a sign it is $D-r$,
with $r$ even. The other factor reduces to one,
so its constant term is odd.

For mass eight, $W(0)=\pm4$, hence $r=\pm4$.
But $W(4)=8$, and $W(-4)=-4+2T_j(-4)\ne0$.
The last assertion follows from the recurrence
used in NN-G034: at $D=-4$, the absolute values
of $T_j$ are at least six and strictly increase.
There is therefore no such factor.

For mass ten, $W(0)$ is $6$ or $-2$, hence
$r\in\{2,-2,6,-6\}$.
At $D=-2$ the value of $W$ is $2T_j(-2)\ne0$;
at $D=6$ it is positive; and at $D=-6$ the
same recurrence, with initial absolute value eight,
makes $-4+2T_j(-6)$ nonzero.
At $D=2$ the value is $4+2T_j(2)$, which is
zero exactly for $j\equiv2\pmod4$.
This proves that $D-2$ is the only possible factor.
After division by it, $B$ is primitive with leading
coefficient two and reduction one modulo two.
If $B$ factored nontrivially, a factor with odd
leading coefficient would retain positive degree
modulo two, impossible when the product reduces
to one. Thus $B$ is irreducible.
Finally $B(4)=W(4)/2=5$.

\paragraph{The consequence of $g_1=64$.}
The complete mass list NN-G035 gives
$d_1=4,u_1=4$ and $w=8$ or ten in this branch.
Therefore
\[
                 \delta_1=\psi:=D(D-3),\qquad U_1=D.
\]
The primitive $W$ is coprime to $D$ in both small-mass
families. The integral divisibility argument of
NN-G031 gives polynomials $C,Q\in\mathbb Z[D]$ with
\[
 \begin{aligned}
 G_2&=DWC,&G_1&=D\delta_1+D^2C=D^2Q,\\
 C&\equiv D-1\pmod2,&Q&=C+D-3\equiv0\pmod2,\\
 C(4)&=3,&Q(4)&=4.
 \end{aligned}                                                       \tag{G036.3}
\]
No effectivity of $C$ or $Q$ is asserted here.

If $W$ is irreducible, it cannot divide $Q$:
primitivity would make the quotient integral,
so $W(4)=8$ or ten would divide $Q(4)=4$.
Thus $W$ is coprime to $G_1$ and the second-cell
identity implies $U_2=WR$ integrally, with
$R\equiv1\pmod2$.
For $w=8$, the scalar relation
$u_2=12-d_2/2$ gives $4\le u_2\le10$.
Its positive odd ratio $u_2/8$ is therefore one.
The equal-mass lemma of NN-G034 implies $U_2=W$
and then $d_2=8$.
For $w=10$, evenness and positivity of
$u_2=5(30-d_2)/8$ give $d_2=14,u_2=10$.
The same equal-mass lemma gives $U_2=W$.

In the exceptional mass-ten case (G036.2), the
irreducible $B$ cannot divide $Q$, since $B(4)=5$
does not divide four. It is also coprime to $D$.
Thus $U_2=BR$ integrally; modulo two, $R\equiv D$.
The scalar relation still gives $u_2=10$.
Since $\deg B=j-1$ and its leading coefficient is
two, the same leading-coefficient argument as in
NN-G034 forces $\deg R=1$: $R$ cannot be constant
by its parity, and the mass-ten $U_2$ then has
leading coefficient two, so $R$ has leading
coefficient one, which its parity allows only
in degree one.
Therefore $U_2$ has degree $j$ and is the unique
mass-ten character of that degree, namely $W$.
We have proved $U_2=W$ in every case, with no
generic coprimality assumption left unexamined.

\paragraph{Small mass forces the polynomial $C$ itself.}
The two first identities now imply
\[
 \delta_1+\delta_2=(W-D)C.
\]
Put $F=(\delta_1+\delta_2)/2$.
It is an effective symmetric integral polynomial,
dominates $\psi$ coefficientwise, and has mass
six for $w=8$ or nine for $w=10$.
The character $C$ has mass three and satisfies
$C\equiv D-1\pmod2$.

If $W=D+4$ or $D+6$, then $F=2C$ or $3C$.
In particular $C$ is effective. Since $F\ge\psi$,
its coefficients at $q^{\pm1},q^{\pm2}$ force
those of $C$ to be at least one. This would give
$C(4)\ge4$, a contradiction.

For the remaining mass-eight characters,
$F=T_jC$. Its leading coefficient is positive.
If $\deg C>1$, that coefficient in $C$ is even
and at least two by parity. The degree of $F$
then exceeds two, so its outer pair contributes
at least four mass units disjointly from the four
units already required by $\psi$. This contradicts
$F(4)=6$. Thus $\deg C=1$.

For mass ten, $F=(1+T_j)C$. If $\deg C>1$,
the same argument produces an outer pair of
coefficient at least two outside $\psi$.
Since $F(4)=9$, the only possibility is exactly
\[
                         F=\psi+1+2T_k
\]
for some $k>2$: four mass units come from $\psi$,
four from that outer pair, and the remaining
unit must be central. This is impossible modulo
two at $D=1$. Its displayed value is odd there,
whereas $F=(1+T_j)C$ is even because $C(1)$ is
even. Therefore $\deg C=1$ also in this case.

Write $C=aD+b$, with $a$ positive odd and
$4a+b=3$. Its symmetric central coefficient is
$3-2a$ and its $q^1$ coefficient is $a$.
If $w=8,j=1$, the $q^1$ coefficient of $T_1C$
is $3-2a$ and must be at least one, forcing $a=1$.
If $w=10,j=1$, the corresponding coefficient of
$(1+T_1)C$ is $3-a$, again forcing the positive
odd $a$ to be one.
For $j\ge2$, the leading pair is outside $\psi$;
an $a\ge3$ would give total mass at least
$4+2a\ge10$, exceeding six or nine.
Thus in every case
\[
                  C=D-1,\qquad Q=2D-4,\qquad
                  G_1=2D^2(D-2).                    \tag{G036.4}
\]
The three fixed first-cell data of NN-G034 are now
forced, rather than assumed. That theorem proves
that the sole competitive realization in this
branch is NN-G033, of cost exactly $2744$.

\paragraph{The improved finite gate and lower bound.}
Remove from the NN-G035 certificate the five base
rows with $g_1=64$. Every strict improvement must
lie in the remaining $38$ rows, with $348$
auxiliary increment pairs and $20$ first-cell
mass triples. The least remaining relaxed score
is $1932$; the exact projection checker verifies
these four numbers. The removed branch has the
unique actual score $2744$ just proved. This establishes
the intermediate lower bound $1932$; the combined
refinement below strengthens it to $2572$.
\hfill$\square$

\paragraph{Geometric translation and scope.}
Within the NN-G010 realization by graded sign-representation
kernels, a competitive first-cell bottom-action object
of ordinary dimension $64$ determines its entire character
and, in turn, the unique $2744$ completion in this closed
presentation. An improvement requires dimension at least
$72$ for that object and initially obeys $20$
dimension triples, reduced to $15$ by the addendum below. This statement has no bound on the
grading support of the other surviving objects.
The irreducibility argument is a statement in the
integral grading-character ring, not a claim of
irreducibility of the representations themselves.

The dictionary is SGA 3, I, 4.7.2--4.7.3; the character
evaluation at $q=-1$ and the parity reductions are
integral restrictions of this ring. Actual composition
uses archive cote 2, scan 17 through NN-G010.
The proof neither identifies a canonical GCT basis nor
applies this lower bound to the general extra-position
problem. At this intermediate stage only a lower bound was established.
The later exclusions below close the entire closed-table gate;
the extra-position optimization remains separate.


\paragraph{Combined refinement: the 2572 lower bound for closed tables.}
This independently reviewed addendum strengthens the
preceding intermediate bound. The full argument below
excludes twelve more necessary base rows and proves the
closed-table interval $[2572,2744]$.
Its first step excludes the two
NN-G035 base rows
\[
 (w,g_1,g_2,d_1,d_2,u_1,u_2)
   =(8,72,96,6,6,4,8),\quad(8,80,96,8,4,4,8),
\]
with every allowed auxiliary increment, not only their
least auxiliary masses. The mass-ten rows with first-cell
triple $(80,8,4)$ are not excluded by this argument.

In either row, $U_1=D$ and the earlier integral quotient
argument gives $G_1=DH$, $G_2=DWC$,
$H=\delta_1+DC$, $C\equiv D-1\pmod2$, and $C(4)=3$.
The mass-eight $W$ is irreducible and coprime to $D$.
It cannot divide $H$, whose mass is respectively $18$
or $20$, because eight divides neither number.
Thus $W$ is coprime to $G_1$ and both second-cell
identities give $U_2=WR,V_2=WZ$ integrally.
Their parities give $R\equiv Z\equiv1$.
The equal-mass lemma forces $U_2=W$, and hence
\[
             \delta_1+\delta_2=(W-D)C.               \tag{G036.R1}
\]

For the first row, both cells have their unique
mass-six character $\delta_i=\psi+2$, where
$\psi=D(D-3)$. If $W=D+4$, (G036.R1) would give
$C=(\psi+2)/2$, which is not an integral polynomial.
Otherwise $W=D+2T_j$ and
$T_jC=\psi+2=(D-1)(D-2)$.
Since $C\equiv D-1$ has positive degree, this forces
$j=1$, $C=D-1$, $W=3D-4$ and $H=2(D-1)^2$.
Now
\[
                 A_2=D^2WC-DHZ.
\]
If $Z$ has degree one, its leading coefficient is
positive even, so the leading coefficient of $A_2$
is at most $3-4<0$. Higher degrees also fail.
Thus $Z$ is constant, but then $V_2=WZ$ has degree
one and cannot dominate its required $q^2$ coefficient
two. This excludes the first row.

For the second row, $\delta_2=\psi$ and
$F=(\delta_1+\delta_2)/2$ is effective, has mass six,
and dominates $\psi$. The mass-eight part of the
small-mass argument in NN-G036 forces $C=D-1$;
the only possible $W$ are $3D-4$ and $D+2T_2$.
For $W=3D-4$ one has $\delta_1=\psi+4$ and
$H=2D(D-2)+4$. Its degree and leading coefficient
give the same contradiction for $Z$ as in the first row.

For $W=D+2T_2$, the forced polynomials are
\[
 \delta_1=\psi+2T_3,\qquad
 H=2D(D-2)+2T_3,\qquad C=D-1.
\]
Here $D^2WC$ has degree five and leading coefficient
two, while $DH$ has degree four and leading coefficient
two. A nonconstant $Z\equiv1\pmod2$ would again give
a negative leading coefficient in $A_2$.
Consequently $Z$ is a positive odd integer.
The scalar relations give $v_2=8Z$ and
$a_2=384-10v_2$. Since $v_2\ge16$ and $a_2\ge16$,
the only possible odd value is $Z=3$, so
$v_2=24,a_2=144$.

The first auxiliary relation is $a_1=320-12v_1$.
It gives $v_1\le26$; and substitution of every other
now-fixed mass into the total gives
\[
                    T=9760+148(26-v_1).
\]
Since $v_1$ is even, the cap $T\le9912$ forces
$v_1=26$ and $a_1=8$.
We show that these masses cannot satisfy coefficientwise
domination.

Effectivity of $A_1=D^2H-D(D-1)V_1$ gives
$\deg V_1\le3$. Write $V_1=(d,a,b,c)$ in symmetric
Laurent notation. Domination and parity require
$d\ge6$ even, $a\ge5$ odd, $b\ge0$ even, and
$c\ge0$ even. The leading coefficient of $A_1$
also gives $c\le2$.
The mass condition is $d+2a+2b+2c=26$.
The relevant coefficients of $A_1$, after eliminating
$d$, are
\[
 \begin{array}{c|l}
 q^0&-44+2a+6b+8c,\\
 q^1&-24+a+3b+5c,\\
 q^2&14-a-2b-c,\\
 q^3&24-a-3b-4c.
 \end{array}
\]
The target lower bounds at degrees zero and three,
respectively two and one, force
$a+3b+4c=23$. The degree-one target then gives
$c-1\ge1$, so $c=2$ and $a+3b=15$.
The degree-two target gives $b-3\ge1$, hence $b\ge4$.
This implies $a\le3$, contradicting $a\ge5$.
The second row is therefore impossible as well.
\hfill$\square$

Deleting these two rows alone leaves $36$ base rows,
$311$ auxiliary increment pairs, and $19$ first-cell
mass triples, with least relaxed score $1956$.
The following uniform arguments exclude three more
mass blocks.

\paragraph{A common-factor observation.}
The identities used throughout are
\[
 WS=G_1G_2,\qquad U_1G_2=W(G_1-D\delta_1),\qquad
 U_2G_1=W(G_2-D\delta_2).                            \tag{G036.R2}
\]
If a primitive irreducible factor $P$ of $W$ divides
neither $U_1$ nor $U_2$, then it divides both $G_1$
and $G_2$. Indeed it divides at least one from the
first identity, and the appropriate one of the next
two forces it into the other. All divisibilities
first obtained over $\mathbb Q[D]$ are integral by
primitivity. In particular divisibility by $P$ forces
$P(4)$ to divide the corresponding ordinary mass.

\paragraph{The entire mass block $(w,g_1,g_2)=(10,80,80)$.}
Here $u_i=10-d_i/2$, so $u_i\in\{4,6,8\}$.
If $W$ is irreducible, its mass ten divides neither
$u_i$, and (G036.R2) gives $W\mid G_1,G_2$.
If $W=(D-2)B$ is the exceptional case of (G036.2),
the primitive irreducible $B$ has mass five, again
dividing neither $u_i$. Thus write $G_i=BP_i$.
The second and third identities give
\[
 U_1P_2=(D-2)(BP_1-D\delta_1),\qquad
 U_2P_1=(D-2)(BP_2-D\delta_2).
\]
An effective character of mass four, six or eight,
dominating and congruent to $D$, has nonzero value
at $D=2$: these values are $2$, $4$, or one of
$6,2,-2$, respectively. Therefore $D-2$ divides
both $P_i$. We have proved $W\mid G_1,G_2$ also
in this exceptional case.

Now write $G_i=WP_i$ in every case. Then
\[
 P_1(4)=P_2(4)=8,\qquad P_1\equiv0,\quad
 P_2\equiv D(D-1)\pmod2.
\]
The polynomial $P_2$ cannot be divisible by $D$:
its integral quotient would have value two at $D=4$
but would reduce to $D-1$ modulo two, whose value
there is odd. Thus $P_2(0)\ne0$.
The divided identities at $D=0$ show also
$P_1(0)\ne0$, because $W(0)\ne0$.
Multiplying them gives
\[
                        W(0)^2=U_1(0)U_2(0).        \tag{G036.R3}
\]
For the mass-ten characters, $W(0)$ is six or
minus two. For the possible $U_i$, the values are
zero, two, four or minus four. Equality (G036.R3)
is possible only when $W(0)=-2$ and both $U_i(0)=2$.
Consequently
\[
 U_1=U_2=D+2,\qquad d_1=d_2=8,\qquad
 W=D+2+2T_j\quad\hbox{with }j\hbox{ odd}.             \tag{G036.R4}
\]
Each cell character is then $\delta_i=\psi+2R_i$,
where $R_i$ has mass two and is either the constant
two or some $T_k$.

Subtracting the two divided identities in (G036.R2),
and putting $E=R_1-R_2$, gives
\[
 A_0(P_1-P_2)=E,\qquad
 A_0=\frac{D+2+T_j}{D}.                             \tag{G036.R5}
\]
The quotient is integral since $j$ is odd and
$T_j(0)=-2$. The difference $P_1-P_2$ is nonzero,
as its reduction is $D(D-1)$, so $E\ne0$.
Every nonzero difference between two members of
$\{2,T_1,T_2,\ldots\}$ has leading coefficient
$1$ or $-1$, and its value at $D=0$ has absolute
value at most four.

If $j=1$, $A_0=2$, contradicting this odd leading
coefficient of $E$.
For odd $j\ge3$, $A_0$ is monic of degree $j-1$.
Thus $P_1-P_2$ has leading coefficient $\pm1$;
its reduction $D(D-1)$ forces its actual degree
to be exactly two. Write it $aD^2+bD+c$, with
$a=\pm1$, $b$ odd and $c$ even.
The recurrence $T_{j+1}=(D-2)T_j-T_{j-1}$ gives
\[
 T_j'(0)=(-1)^{j-1}j^2,\qquad A_0(0)=1+j^2.
\]
For completeness, the derivative recurrence is
$r_{j+1}=2(-1)^j-2r_j-r_{j-1}$ with
$r_0=0,r_1=1$; the displayed formula satisfies it
by direct substitution.
Evaluating (G036.R5) at zero now gives
$|(1+j^2)c|\le4$, so $c=0$.
But $P_1(4)=P_2(4)$ then implies
$16a+4b=0$, making $b=-4a$ even, a contradiction.
This excludes the entire mass block, including
all six base rows and all their auxiliary choices.

\paragraph{The entire mass block $(w,g_1,g_2)=(8,96,96)$.}
Here $u_i=8-d_i/3$ is four or six.
The mass-eight $W$ is irreducible, and its mass
divides neither $u_i$, so (G036.R2) again gives
$G_i=WP_i$. Now $P_i(4)=12$, with
$P_1\equiv0$ and $P_2\equiv D(D-1)$ modulo two.
The polynomial $P_1$ cannot be divisible by $D$:
the quotient would have even coefficients but
ordinary mass $12/4=3$.
The divided identities at zero consequently
force both $P_i(0)\ne0$ and give (G036.R3).
But $W(0)^2=16$, whereas $U_i(0)$ are zero or two,
whose products are zero or four. This contradiction
excludes both base rows in this block.

\paragraph{The two asymmetric mass-ten blocks.}
The finite certificate has one base row for
$(w,g_1,g_2)=(10,76,100)$, with
$(d_1,d_2,u_1,u_2)=(4,6,6,10)$, and its index-swapped
row with $(g_1,g_2)=(100,76)$.
Let $G_-$ denote the smaller action character and
$G_+$ the larger. Their cell characters are
respectively $\psi$ and $\psi+2$, and the smaller
index has $U_-=D+2$.

If $W$ is irreducible, it cannot divide $G_-$,
since ten does not divide $76$.
It therefore divides $G_+$ and then $U_+$ by
(G036.R2). Since $u_+=10$, the equal-mass lemma
forces $U_+=W$.
If $W=(D-2)B$, the irreducible mass-five $B$ cannot
divide $G_-$, so the same argument gives $B\mid U_+$.
Writing $U_+=BR$, its reduction is $R\equiv D$.
The leading-coefficient argument from NN-G036 forces
$\deg R=1$ and then the mass-ten classification gives
$U_+=W$ in this case as well.

The larger-index identity now gives
$G_+=G_-+D(\psi+2)$.
Substitution into the smaller-index identity yields
\[
 (W-D-2)G_-=D\bigl(W\psi+(D+2)(\psi+2)\bigr).       \tag{G036.R6}
\]
If $W=D+6$, its denominator is four and the resulting
$G_-$ has leading coefficient $1/2$, contradicting
integrality.
For $W=D+2+2T_j$, equation (G036.R6) gives
\[
 G_-=D\psi+\frac{N(D)}{T_j(D)},\qquad
 N(D)=D(D+2)(D^2-3D+1).
\]
Integrality forces the monic $T_j$ to divide $N$.
Thus $j\le4$. None does:
$N(2)=-8$ excludes $T_1=D-2$ and also $T_3$,
which contains $T_1$;
the remainder modulo $T_2=D^2-4D+2$ is $16D-10$;
and $T_4=D^4-8D^3+20D^2-16D+2$ is not the
monic degree-four polynomial $N$.
This excludes both asymmetric blocks, uniformly
over all original support indices.
\hfill$\square$

\paragraph{The resulting finite dimension gate.}
In addition to the five rows already removed in
NN-G036, this addendum removes two isolated mass-eight
rows, six rows in the $(10,80,80)$ block, two in the
$(8,96,96)$ block, and the two asymmetric rows.
The remaining complete necessary projection has
$26$ base rows, $42$ auxiliary increment pairs,
and $15$ first-cell mass triples. Its least relaxed
score is $2572$. The proved exact attainer still
has score $2744$, so the closed-table minimum lies
in $[2572,2744]$.

In the NN-G010 geometry, these are exclusions of
ordinary-dimension patterns for actual composing
graded sign-representation objects. The tests at
$D=0$ are integral character evaluations at $q=-1$;
the derivative there is a character-polynomial
calculation, not a Weil-weight assertion.
They preserve the SGA 3, I, 4.7.2--4.7.3 and
archive cote 2, scan 17 dictionary of NN-G036.
The surviving rows are not declared realizable,
and this lower bound is not transferred to arbitrary
extra-position lifts or canonical GCT multiplicities.


% Further NN-G036 refinement: root whole-proof PASS 2026-11-22T00:03:47Z.
\paragraph{The last mass-eight competitor.}
The remaining mass-eight base row has
\[
 (g_1,g_2,d_1,d_2,u_1,u_2)=(128,80,12,4,8,4).
\]
Thus $\delta_2=\psi$ and $U_2=D$. The integral divisibility
argument already used in NN-G035 gives
\[
 G_1=DWC,\qquad G_2=D(\psi+DC),\qquad
 C\equiv0\pmod2,\quad C(4)=4.
\]
For clarity, $DG_1=W(G_2-D\psi)$ and $\gcd(W,D)=1$
first give $G_1=WQ$ and $G_2=D(Q+\psi)$ integrally.
The other cell identity at zero forces $G_1(0)=0$,
so $Q(0)=0$ and $Q=DC$. Reducing $G_1=DWC$ modulo
two gives the displayed parity of $C$.
The irreducible mass-eight $W$ cannot divide $G_2$:
it is coprime to $D$, whereas the integral polynomial
$\psi+DC$ has value twenty at four, not divisible by eight.
Consequently $U_1G_2=W(G_1-D\delta_1)$ forces $W\mid U_1$.
The equal-mass lemma gives $U_1=W$. Cancelling now yields
\[
 \delta_1+\delta_2=(W-D)C.
\]
Write $C=2H$ with $H\in\mathbb Z[D]$ and $H(4)=2$.
The effective character
$F=(\delta_1+\delta_2)/2=(W-D)H$ has mass eight and
dominates $\psi$.
If $W=D+4$, then $F=4H$, so $H$ is effective;
domination forces its coefficients at both positive
exponents one and two to be at least one. Its mass
would therefore be at least four, contradicting two.
Otherwise $W=D+2T_j$ and $F=2T_jH$ has all Laurent
coefficients even. Its coefficients at exponents
$1,-1,2,-2$ must all be at least two. These exhaust
the mass eight, hence $F=2\psi$ and $T_jH=\psi$.
Since $T_j$ is monic of degree $j$, this requires
$j\le2$. But $T_1=D-2$ does not divide $\psi$, and
$T_2=D^2-4D+2$ is not the monic quadratic $\psi$.
This excludes the last mass-eight competitor.
Together with the preceding exclusions and NN-G034,
every competitive mass-eight closed table is the unique
$2744$ attainer. No bound on its initial support indices
was imposed.

\paragraph{The mass-ten block with equal action masses one hundred.}
The row is
\[
 (g_1,g_2,d_1,d_2,u_1,u_2)=(100,100,10,10,6,6).
\]
Thus both $U_i=D+2$. If $W$ is irreducible, its mass
ten divides neither $u_i$, so (G036.R2) forces
$W\mid G_1,G_2$. In the exceptional factorization
$W=(D-2)B$, the irreducible $B$ has mass five and
again divides neither $U_i$; it divides both $G_i$.
After cancelling $B$ in the cell identities, evaluation
at $D=2$, where $U_i(2)=4\ne0$, forces the remaining
factor $D-2$ into both $G_i$. Thus in every case
\[
 G_i=WP_i,\qquad P_i(4)=10,\qquad
 P_1\equiv0,\quad P_2\equiv D(D-1)\pmod2.
\]
The polynomial $P_2$ cannot be divisible by $D$,
since its quotient would have nonintegral value $10/4$
at four. The divided cell identities at zero then
force both $P_i(0)\ne0$ and give $W(0)^2=4$.
Consequently $W=D+2+2T_j$ with $j$ odd.
Write $\delta_i=\psi+2R_i$. Now $R_i$ is an effective
character of mass three, hence is either three or
$1+T_k$. A nonzero difference $E=R_1-R_2$ still
has leading coefficient $\pm1$ and $|E(0)|\le4$.
It is nonzero because the parity difference of
$P_1,P_2$ is nonzero.
Subtracting the divided identities gives exactly
\[
 \frac{D+2+T_j}{D}(P_1-P_2)=E.
\]
For $j=1$, the factor is two, contradicting the odd
leading coefficient of $E$. For odd $j\ge3$, it is
monic, with value $1+j^2$ at zero. Hence
$P_1-P_2=aD^2+bD+c$, where $a=\pm1$, $b$ is odd,
and $c$ is even: the degree follows from its leading
coefficient and its reduction $D(D-1)$.
The bound on $E(0)$ forces $c=0$.
Equal values of the two $P_i$ at four then give
$16a+4b=0$, contradicting that $b$ is odd.
The row is impossible.

\paragraph{The two mass-ten blocks with action masses eighty-four and one hundred.}
In either orientation the smaller action has
$\delta_-=\psi+2$, $U_-=D+2$, while the larger
has $\delta_+=\psi$ and $u_+=10$.
An irreducible $W$ of mass ten cannot divide $G_-$
of mass eighty-four; (G036.R2) consequently forces
$W\mid G_+$ and $W\mid U_+$, whence $U_+=W$
by the equal-mass lemma.
If $W=(D-2)B$, the mass-five irreducible factor
$B$ likewise cannot divide $G_-$. Thus $B\mid U_+$.
Writing $U_+=BR$, reduction gives $R\equiv D\pmod2$.
The leading coefficient argument in NN-G036 forces
$\deg R=1$; the mass-ten classification then gives
$U_+=W$ also in this case.
The larger-index identity now says $G_+=G_-+D\psi$.
Substituting in the smaller-index identity gives
\[
 (W-D-2)G_-=D\bigl(W(\psi+2)+(D+2)\psi\bigr).
\]
For $W=D+6$ this would make the leading coefficient
of $G_-$ equal to one half. Otherwise
$W=D+2+2T_j$ and
\[
 G_-=D(\psi+2)+\frac{N(D)}{T_j(D)},\qquad
 N=D(D+2)(D^2-3D+1).
\]
The same integral divisibility obstruction as in
(G036.R6) applies: $T_j\mid N$ requires $j\le4$,
but $N(2)=-8$, $T_1\mid T_3$,
$N\bmod T_2=16D-10$, and $N\ne T_4$.
Both orientations are excluded, with every auxiliary
increment allowed by their scalar relaxation.
\hfill$\square$

\paragraph{Scope and geometric content.}
These four base-row exclusions concern the ordinary
dimensions of actual graded $C_2$-representation
structure objects in the matrix-kernel category of
NN-G010, with the same symmetric closed multiplication
table and character constraints. In particular, within
the mass-eight gluing stratum the actual smallest
effective lift is the $2744$-pair object of NN-G033,
uniquely within that presentation. The integral torus
character interpretation is SGA~3, I, 4.7.2--4.7.3;
kernel composition is the supplied archive transcription,
cote~2, scan~17. No new claim about a canonical GCT
representation family or a different multiplication
table follows.

\paragraph{A mass-twelve gluing lemma.}
Throughout, $D=q+2+q^{-1}$ and $T_j=q^j+q^{-j}$.
Every effective $W\ge D$ with $W\equiv D\pmod2$ and $W(4)=12$
is one of
\[
 D+8,\quad D+4+2T_j,\quad D+4T_j,\quad
 D+2T_i+2T_j\ (i<j).
\]
Indeed $(W-D)/2$ is an effective symmetric character of mass four,
and these are its complete partitions into a central coefficient
and paired nonzero exponents. Hence $W(0)\in\{-8,0,8\}$.
If $W(0)=0$, write $W=DH$. Then
\[
 H\in\mathbb Z[D],\quad H(4)=3,\quad H\equiv1\pmod2,
 \quad H(0)\ne1,-1.
\]
Either $H=3$ (the case $W=3D$), or $H$ is primitive and
irreducible over $\mathbb Q[D]$, with leading coefficient two.
For the last assertion, every nonconstant case has leading
coefficient two and odd constant coefficient. In a nontrivial
primitive integral factorization, one factor has odd leading
coefficient and would keep its positive degree modulo two,
contradicting that the product reduces to one.
For the assertion at zero, the zero-valued cases in the list are
$D+4+2T_o$ with $o$ odd and $D+2T_o+2T_e$ with $o$ odd and
$e$ even. The derivative identity from (G036.R5) gives
\[
 H(0)=1+2o^2\quad\hbox{or}\quad 1+2(o^2-e^2).
\]
Neither is one or minus one: equality would require equal
squares of opposite parity or $e^2-o^2=1$, impossible for
positive distinct integers $e,o$. This also includes $o=1$
and the constant $H=3$.

\paragraph{The nonzero-value branch.}
For any effective $U\ge D$ with $U\equiv D\pmod2$ and mass $u$,
\[
 |U(0)|\le u-4,
\]
because $(U-D)/2$ is effective and its absolute value at $q=-1$
is at most its mass. In every remaining mass-twelve row of
NN-G036, $(u_1-4)(u_2-4)<64$, as follows directly from the
complete listed mass tuples. If $W(0)=8$ or $-8$, the identities
\[
 WG_1-U_1G_2=DW\delta_1,\qquad
 WG_2-U_2G_1=DW\delta_2
\]
have an invertible coefficient matrix at $D=0$ over $\mathbb Q$.
Thus $D$ divides both $G_i$ in $\mathbb Z[D]$.
Using $G_1\equiv0$ and $G_2\equiv D^2(D-1)\pmod2$, their
quotients by $D$ have even evaluation at $D=4$. Consequently
both $g_i$ are divisible by eight. This excludes any mass row
containing $84$ in this branch.
If also $\delta_1=\delta_2=\psi=D(D-3)$, divide the same
identities by $D$ and apply the invertible matrix at zero once
again. Both $G_i$ are divisible by $D^2$. Since $WS=G_1G_2$
and $W(0)\ne0$, the polynomial $S$ is divisible by $D^4$.
In particular, the mass block $(g_1,g_2)=(80,96)$ or $(96,80)$,
where $s=640$, is impossible: $4^4$ does not divide $640$.

\paragraph{The two blocks with action masses $84,88$.}
Their action masses are $84,88$ in either order, with associated
$u$-masses six and eight. The preceding paragraph leaves only
$W=DH$. If $H$ is nonconstant, it cannot divide the mass-six
character $D+2$ (the integral leading coefficient would have
to divide one), nor a mass-eight character (its value three
would have to divide eight). The common-factor argument of
(G036.R2) forces $H$ to divide both $G_i$, contradicting that
three does not divide $88$.
If $H=3$, reduction modulo three supplies the same argument
in the integral domain $\mathbb F_3[D]$. Both $U_i$ have
nonzero reduction: one is $D+2$, and the other has mass eight.
The three identities (G036.R2) then force both $G_i$ to have
all coefficients divisible by three, again contradicting
mass $88$. Thus both base rows are impossible.

\paragraph{The block with both action masses $84$.}
Here $U_1=U_2=D$ and $d_1=d_2=14$. Again only $W=DH$
remains. The common-factor argument (or reduction modulo
three when $H=3$) gives $G_i=HP_i$ integrally. Their masses
are $P_i(4)=28$, and their reductions are
$P_1\equiv0$, $P_2\equiv D^2(D-1)\pmod2$.
Cancelling $DH$ in the two cell identities gives
\[
 P_2=HP_1-D\delta_1,\qquad
 P_1=HP_2-D\delta_2.
\]
Since $H(0)^2\ne1$, evaluating at zero forces
$P_1(0)=P_2(0)=0$. In particular $P_1/D$ is an integral
polynomial with every coefficient even, but its evaluation
at four is $28/4=7$. This is a contradiction.

\paragraph{The two smallest mass-twelve rows.}
Suppose $(g_1,g_2)$ is $(80,96)$ or $(96,80)$ and both
cells are $\psi$. Let $G_-$ have mass eighty and $G_+$
have mass ninety-six; their $U$-masses are eight and twelve.
The nonzero-value branch was already excluded, so $W=DH$.
For nonconstant $H$, its value three cannot divide eighty,
so $H\nmid G_-$. The identities (G036.R2) force
$H\mid G_+$ and $H\mid U_+$. Every mass-twelve $U_+$
has value at zero equal to zero or $\pm8$. Since $H(0)$ is
odd and is neither one nor minus one, its divisibility
forces $U_+(0)=0$. Write $U_+=DJ$. The classification above
applies to $J$: it is either three or primitive irreducible
of leading coefficient two. Since $H\mid J$, the latter
case is forced and $J=H$. Thus $U_+=W$.
For $H=3$, reduce $U_+G_-=3D(G_+-D\psi)$ modulo three.
Since $G_-$ has mass eighty, its reduction is nonzero, so
all coefficients of $U_+$ are divisible by three. Its odd
coefficients at $q,q^{-1}$ are at least three each, and its
even central coefficient is at least six. These exhaust its
mass twelve, so $U_+=3D=W$ also in this case.

The larger-cell identity now gives $G_+=G_-+D\psi$.
The other identity becomes
\[
             (W-U_-)G_-=D\psi(W+U_-).
\]
The mass-eight classification gives $U_-(0)=4$ or $-4$,
while $W(0)=0$. Hence $D^2\mid G_-$.
If $G_-=G_1$, its quotient by $D^2$ has even coefficients
but mass $80/16=5$, a contradiction.
If $G_-=G_2$, use the necessary mod-four identity
$G_2-D\delta_2\equiv G_1\pmod4$ from NN-G030.
Here $\delta_2=\psi$ and $G_1=G_2+D\psi$, so it would
require $2D\psi\equiv0\pmod4$, which is false.
This excludes both rows, including every auxiliary increment.

These five exclusions use neither a degree bound nor
positivity of any quotient. They retain precisely the
closed-table character dictionary and scope of NN-G036.

\paragraph{Intermediate combined gate at 00:03:47 UTC.}
The additional four small-gluing and five mass-twelve exclusions
leave exactly $17$ necessary base rows, $26$ auxiliary increment
pairs, and $8$ first-cell mass triples. The least remaining
relaxed score is $2604$, so the intermediate closed-table
interval at this stage is $[2604,2744]$. Every strict improvement has $w=10$
or $12$; in the mass-ten case its first data must be
$(g_1,g_2,d_1,u_1)=(80,120,8,4)$, with
$d_2\in\{6,10,14,18\}$. The earlier numerical intervals in
this subsection are successive intermediate bounds.
The exact finite projection is
\path{verify_second_competitive_refinement.py};
the preceding polynomial arguments justify every deletion
without a support-degree bound.

\paragraph{Geometric conclusion at the intermediate stage.}
Within the specified closed presentation of composing graded
$C_2$-representation objects from NN-G010, every effective
lift has at least $2604$ cancelling trivial/sign pairs and
one with $2744$ pairs exists. In the gluing-dimension-eight
stratum, that $2744$ object is the unique minimum in this
presentation. A strict improvement can occur only in the
seventeen displayed ordinary-dimension rows, whose actual
character realization remains a further condition.
The bounds concern genuine effective representation objects,
using SGA~3, I, 4.7.2--4.7.3 and the archive kernel-product
transcription at cote~2, scan~17. They are not statements
about other multiplication-table shapes or canonical GCT
multiplicities.


% NN-G036 further uniform refinement; root whole-proof PASS 2026-11-22T00:06:53Z.
\paragraph{The entire remaining mass-ten gluing branch.}
The reviewed finite projection leaves four mass-ten base rows:
\[
 \begin{aligned}
 &g_1=80,\quad g_2=120,\quad d_1=8,\quad U_1=D,\\
 &(d_2,u_2)\in\{(6,12),(10,10),(14,8),(18,6)\}.
 \end{aligned}
\]
Every allowed auxiliary increment in these rows has
\[
 a_1=8,\qquad v_1=26,\qquad v_2\in\{56,58\}.
\]
Equality with the first auxiliary's absolute trace mass
therefore fixes the entire polynomial, not only its mass:
\[
 A_1=D(D-2)(D-3)
   =2+q+q^{-1}+q^2+q^{-2}+q^3+q^{-3}.                 \tag{G036.T1}
\]
We exclude all four rows, with arbitrary support locations.

The integral $U_1=D$ argument gives
\[
 \begin{aligned}
 G_2&=DWC,\qquad G_1=DH,\qquad H=\delta_1+DC,\\
 C&\equiv D-1\pmod2,\qquad H\equiv0\pmod2,\\
 C(4)&=3,\qquad H(4)=20.
 \end{aligned}                                                   \tag{G036.T2}
\]
The mass-eight cell is necessarily
$\delta_1=\psi+2R$, where $R$ is either two or $T_k$.
In particular $\delta_1(0)$ is four or minus four.

First suppose $W$ is irreducible. If it did not divide
$H$, it would be coprime to $G_1=DH$, because $W(0)\ne0$.
The second auxiliary identity
$V_2G_1=W(DG_2-A_2)$ would force $W\mid V_2$ integrally.
Evaluation at four would then give $10\mid v_2$, contrary
to $v_2=56$ or $58$. Thus $W\mid H$.

In the exceptional case $W=(D-2)B$ of (G036.2),
$B$ is primitive irreducible, $B(4)=5$,
$B\equiv1\pmod2$, and $B(0)=-3$.
If $B\nmid H$, the same auxiliary identity gives
$B\mid V_2$, contradicting $5\nmid v_2$.
Therefore $H=2BK$ for an integral polynomial $K$;
the factor two follows by reducing $H$ modulo two.
But then $H(0)$ is divisible by six, whereas
$H(0)=\delta_1(0)=\pm4$. The exceptional case is impossible.

We are left with the irreducible case, where
\[
 H=2WK,\qquad K\in\mathbb Z[D],\qquad K(4)=1.
\]
Here the integral quotient by $W$ is even because
$\overline W=D$ and $\overline H=0$ in the domain
$\mathbb F_2[D]$. If $W(0)=6$, the identity
$H(0)=2W(0)K(0)=\pm4$ is impossible. Hence
\[
 W=D+2+2T_j,\quad j\ {\rm odd},\qquad W(0)=-2,
 \qquad K(0)\in\{1,-1\}.
\]
Every integer polynomial satisfies $K(4)\equiv K(0)\pmod4$.
Since $K(4)=1$, this forces $K(0)=1$.
Consequently $\delta_1(0)=H(0)=-4$, so
$R=T_k$ with $k$ odd. At $D=2$ the Chebyshev
evaluations are therefore
\[
 W(2)=4,\qquad \delta_1(2)=\psi(2)+2T_k(2)=-2.
                                                                  \tag{G036.T3}
\]

The first auxiliary identity, after cancelling $W$
and then $D$, is
\[
 C V_1=2DWK-A_1/D.
\]
By (G036.T1), evaluation at two gives
\[
                  C(2)V_1(2)=16K(2).               \tag{G036.T4}
\]
The integer $C(2)$ is odd by (G036.T2).
Therefore $16$ divides $V_1(2)$.
The coefficientwise lower bound for $V_1$ is the
effective character $V_{\min}=6+5(q+q^{-1})$,
of mass sixteen. Its parity and mass twenty-six imply
$V_1=V_{\min}+2P$ for an effective symmetric integral
character $P$ of mass five. At $D=2$, equivalently
$q=i$, the triangle inequality gives
$|P(i)|\le5$ and hence
\[
                         -4\le V_1(2)\le16.
\]
Its only multiples of sixteen in this interval are
zero and sixteen.
If $V_1(2)=0$, (G036.T4) gives $K(2)=0$.
The monic polynomial $D-2$ would divide $K$ integrally,
so its value two at four would divide $K(4)=1$,
which is impossible.
Thus $V_1(2)=16$, and (G036.T4) gives $K(2)=C(2)$.
But (G036.T2)--(G036.T3) also give
\[
 C(2)=\frac{2W(2)K(2)-\delta_1(2)}2
      =4K(2)+1.
\]
This would require $-3K(2)=1$, contradicting integrality.
Every mass-ten competitor is excluded.
\hfill$\square$

\paragraph{Meaning of this refinement.}
The closed-table optimum is consequently attained at
gluing dimension eight unless a strictly smaller object
occurs in the remaining dimension-twelve stratum.
The proof uses restriction of integral torus characters
at $q=-1$ and $q=i$, together with ordinary dimensions
at $q=1$; it assumes no Weil-weight estimate and no
effectivity of $C,H/W$, or $K$.
Its geometric objects and source dictionary remain the
actual composing graded sign representations of NN-G010,
SGA~3, I, 4.7.2--4.7.3, and the supplied kernel-product
transcription at archive cote~2, scan~17.

\paragraph{Projection after the mass-ten exclusion.}
Deleting its four base rows and six auxiliary increments leaves
$13$ necessary base rows, $20$ auxiliary increments and
$7$ first-cell mass triples. Every strict improvement on
$2744$ now has gluing dimension twelve. The least relaxed
score remains $2604$. The exact row projection is certified
separately by \path{verify_mass10_final_exclusion.py}.


% Further NN-G036 refinement; root whole-proof PASS 2026-11-22T00:12:35Z.
\paragraph{A mass-six action gives integral divisibility.}
Suppose $W\in\mathbb Z[D]$ satisfies $W\equiv D\pmod2$
and $w=W(4)$ is divisible by four.
Write $W=D+2R$. Then $R(4)=(w-4)/2$ is even,
and $R(-2)\equiv R(4)\pmod2$. Therefore
\[
                         W(-2)\equiv2\pmod4.        \tag{G036.U1}
\]
In particular $W$ and $D+2$ are coprime over
$\mathbb Q[D]$.
If $(D+2)G=W F$ for integral polynomials, then
$W\mid G$ in $\mathbb Z[D]$, even when $W$ is
not primitive. To check the content explicitly,
write $W=cW_0$ with $W_0$ primitive and $c>0$.
Gauss's content identity and primitivity of $D+2$
first give $c\mid G$. Write $G=cG_0$.
Coprimality gives $W_0\mid G_0$ over $\mathbb Q[D]$,
and primitivity makes this quotient integral.

Consequently, in any of the closed tables under
consideration, $u_i=6$ forces $U_i=D+2$, and
\[
                 4\mid w,\ u_i=6
                 \quad\Longrightarrow\quad w\mid g_j
                 \qquad(i\ne j).                  \tag{G036.U2}
\]
At $w=12$ this excludes both $(g_1,g_2)=(72,112)$
rows, since $u_1=6$ but $12\nmid112$.
It also excludes the three $(80,96)$ rows with
$d_2=14,u_2=6$, and their three index-reversed rows:
in each case the opposite action has mass eighty.
Thus eight entire base rows, with all their auxiliary
increments, are impossible.

\paragraph{The forward row with action masses eighty and ninety-six.}
Consider
\[
 (w,g_1,g_2,d_1,d_2,u_1,u_2)
                    =(12,80,96,12,4,4,12).         \tag{G036.U3}
\]
Here $U_1=D$ and $\delta_2=\psi$.
If $W(0)\ne0$, the integral quotient argument for
$U_1=D$ gives
$G_2=DWC$ with $C\equiv D-1\pmod2$.
It applies to nonprimitive $W$ as well by the same
content argument just proved. Its value would be
$C(4)=96/(4\cdot12)=2$, whereas its displayed
parity requires an odd value. This is impossible.

It remains that $W=DH$, with the mass-twelve
classification proved above: $H(4)=3$,
$H\equiv1\pmod2$, and $H(0)\ne\pm1$.
Either $H=3$, or $H$ is primitive irreducible
of leading coefficient two.
Since three does not divide $g_1=80$, in the
nonconstant case $H\nmid G_1$.
The second cell identity then forces $H\mid U_2$.
As in the earlier mass-twelve proof, $H(0)$ is
odd and not $\pm1$, whereas $U_2(0)$ is zero
or $\pm8$; therefore $U_2(0)=0$.
Writing $U_2=DJ$, both $H,J$ are primitive
irreducible of leading coefficient two and mass
three, so $J=H$ and $U_2=W$.
For $H=3$, reducing the same identity modulo
three instead gives $3\mid U_2$ coefficientwise;
its effective mass twelve and parity force
$U_2=3D=W$. Thus $U_2=W$ in every case.

The two cell identities now imply
\[
 G_2=G_1+D\psi,\qquad
                    (H-1)G_1=D(H\delta_1+\psi).     \tag{G036.U4}
\]
Write $\delta_1=\psi+2R$, where $R$ is effective
of mass four. Then $\delta_1(0)\in\{-8,0,8\}$.
If $\delta_1(0)=0$, both $\delta_1$ and $\psi$
are divisible by $D$. Since $H(0)\ne1$,
(G036.U4) forces $D^2\mid G_1$.
The quotient has even coefficients but value
$80/16=5$ at four, a contradiction.

If $\delta_1(0)=\pm8$, first $D\mid G_1$.
Divide (G036.U4) by $D$ and evaluate at zero.
Coprimality of the consecutive integers
$H(0),H(0)-1$ shows that $H(0)-1$ divides eight.
The explicit root classification gives
\[
 H(0)=1+2o^2\quad\hbox{or}\quad
 H(0)=1+2(o^2-e^2),
\]
with $o$ positive odd and $e$ positive even.
In particular $H(0)\equiv3\pmod8$.
The only such odd integer with $H(0)-1\mid8$
is three. The first formula then forces $o=1$;
the second would require $o^2-e^2=1$,
impossible for positive $e$.
Consequently $H=3$ and $W=3D$.

Equation (G036.U4) becomes
$G_1=D(2\psi+3R)$.
Its even parity forces $R=2P$ integrally, where
$P$ is effective of mass two, hence either two
or some $T_k$. Thus
\[
 \delta_1=\psi+4P,\qquad
 G_1=D(2\psi+6P),\qquad
 G_2=D(3\psi+6P).
\]
The central Laurent coefficient of $2D\psi$
is four. For $P=T_k$ with $k\ge2$, that of
$6DP$ is zero, whereas $G_1$ must have central
coefficient at least sixteen.
Therefore $P=2$ or $T_1$.

Put $L=2\psi+6P$, an integral quadratic with
leading coefficient two and $L(4)=20$.
The second auxiliary identity gives
$V_2L=3(DG_2-A_2)$.
Modulo three, $L$ reduces to $2\psi\ne0$,
so $V_2=3Z$ integrally.
Since $V_2$ is effective and $V_2\equiv D\pmod2$,
$Z$ is effective and $Z\equiv D\pmod2$.
The actual second auxiliary coefficient is
\[
              A_2=D^2(3\psi+6P)-Z(2\psi+6P).       \tag{G036.U5}
\]
The first term has degree four and leading
coefficient three. A $Z$ of degree greater
than two makes the leading coefficient of
$A_2$ negative. If $\deg Z=2$, its leading
coefficient is positive even and at least two,
so the degree-four coefficient is at most
$3-2\cdot2<0$.
Thus effectivity forces $\deg Z\le1$.
But then $V_2=3Z$ has no $q^2$ coefficient,
contradicting its required coefficient at least two.
The forward row (G036.U3) is excluded.
\hfill$\square$

\paragraph{Geometric scope.}
The divisibility lemma removes ordinary-dimension
patterns for effective graded sign-representation
objects in the specified closed kernel presentation.
It does not assume positive polynomial quotients:
content and integrality alone supply the character
divisibilities. The forward-row exclusion further
uses actual coefficientwise effectivity of the
structure objects. These are the same mathematical
objects of NN-G010 and SGA~3, I, 4.7.2--4.7.3,
with composition in the supplied archive
cote~2, scan~17 transcription. No change of
multiplication-table presentation is permitted
in this conclusion.

\paragraph{The two mass-twelve rows of relaxed cost $2604$.}
Let $G_-$ have mass eighty and $G_+$ have mass ninety-six,
in either orientation of the closed table. The two rows now considered have
\[
 W(4)=12,\qquad \delta_- =\psi+2R,\quad U_-=D+2,
 \qquad \delta_+=\psi,\quad U_+(4)=12,
\]
where $\psi=D(D-3)$ and $R$ is either the constant two or $T_k$,
with $k\ge1$. Every allowed auxiliary increment is included.
In particular $\delta_-(0)=4$ or $-4$.

We first record the divisibility interface used in both orientations.
Write $W=D+2B$ with $B\in\mathbb Z[D]$. Since $W(4)=12$,
$B(4)=4$; therefore $B(-2)$ is even, and
$W(-2)\equiv2\pmod4$. Thus $W$ is coprime to $D+2$
over $\mathbb Q[D]$. From the minus-cell identity
\[
                 (D+2)G_+=W(G_--D\delta_-)
\]
it follows that $W$ divides $G_+$ in $\mathbb Z[D]$.
Here integrality follows either by Gauss's lemma and content,
or directly as follows. Write $W=cW_0$, with $W_0$ primitive.
The rational coprimality gives $G_+=W_0Q$ with $Q$ integral.
After cancellation, $c$ divides every coefficient of $(D+2)Q$;
the product-content identity and primitivity of $D+2$ imply
that $c$ divides every coefficient of $Q$.

If $W(0)\ne0$, the preceding mass-twelve classification gives
$W(0)=\pm8$ and $U_+(0)\in\{-8,0,8\}$.
The two cell identities at zero force both $G_i(0)=0$:
their coefficient determinant is
$W(0)^2-U_-(0)U_+(0)=64-2U_+(0)\ne0$.
Write $G_\pm=DP_\pm$ integrally. Dividing the identities by
$D$ and evaluating at zero yields
\[
 wP_-(0)-2P_+(0)=w\delta_-(0),\qquad
 wP_+(0)-uP_-(0)=0,
 \quad w=W(0),\ u=U_+(0).
\]
Consequently
\[
 P_-(0)=\frac{64\delta_-(0)}{64-2u}.
\]
For $u=8$ or $-8$ this is respectively
$4\delta_-(0)/3$ or $4\delta_-(0)/5$, neither an integer.
Hence $u=0$ and $P_+(0)=0$, so $D^2$ divides $G_+$.
On the other hand $G_+=WQ$ with $Q(4)=96/12=8$.
Since $W(0)\ne0$, $D^2$ divides $Q$ integrally.
Evaluation at four would make sixteen divide eight, a contradiction.

It remains to treat $W(0)=0$, so $W=DH$ with $H$ as in
the preceding mass-twelve gluing lemma. If $H$ is nonconstant,
it is primitive irreducible with leading coefficient two,
$H(4)=3$, and odd $H(0)\ne1,-1$.
It cannot divide $G_-$, since three does not divide eighty.
The plus-cell identity therefore forces $H\mid U_+$.
Divisibility at zero, together with $U_+(0)\in\{-8,0,8\}$,
forces $U_+(0)=0$. Write $U_+=DJ$.
The same classification makes $J$ either three or a primitive
irreducible polynomial of leading coefficient two and value three
at four. Since $H\mid J$, the constant case is impossible and
$J=H$. Thus $U_+=W$.
If $H=3$, reduce the plus-cell identity modulo three.
The reduction of $G_-$ is nonzero because its mass is eighty,
so all coefficients of $U_+$ are divisible by three.
Its odd coefficients at $q,q^{-1}$ are at least three each,
and its even central coefficient at least six. These exhaust
mass twelve, whence again $U_+=3D=W$.

Cancelling $W$ in the plus-cell identity gives
$G_+=G_-+D\psi$. The other cell identity becomes
\[
 (W-D-2)G_-=D\bigl(W\delta_-+(D+2)\psi\bigr).
\]
Its right side is divisible by $D^2$, whereas the first factor
on the left has value $-2$ at zero. Therefore $D^2\mid G_-$.
If $G_-=G_1$, its quotient by $D^2$ has every coefficient even
but evaluation $80/16=5$ at four, a contradiction.
If $G_-=G_2$, the necessary congruence of NN-G030 is
\[
 G_--D\delta_-\equiv G_+\pmod4.
\]
Substituting $G_+=G_-+D\psi$ gives
$2D(\psi+R)\equiv0\pmod4$, so $R\equiv\psi\pmod2$.
This is impossible. The constant two reduces to zero; if
$R=T_k$, its monic degree is $k$, so equality with the degree-two
polynomial $\psi$ would require $k=2$. But
$T_2\equiv D^2$ and $\psi\equiv D^2+D$ modulo two.
Both orientations are excluded for arbitrary grading support.
\hfill$\square$

The objects and scope are those of NN-G036: actual effective graded
sign representations composing in the symmetric closed table.
Restriction at $q=-1$, ordinary dimension at $q=1$, and integral
character divisibility are operations in the torus representation
ring of SGA~3, I, 4.7.2--4.7.3. No quotient is assumed effective,
and no statement about extra-position tables is inferred.

\paragraph{The reverse mass-twelve row with a mass-four action.}
We now exclude the remaining orientation
\[
 (w,g_1,g_2,d_1,d_2,u_1,u_2)
                  =(12,96,80,4,12,12,4).
\]
Write $G_+=G_1$, $G_-=G_2$, $U_+=U_1$, $U_-=U_2=D$,
$\delta_+=\psi=D(D-3)$, and $\delta_-=\psi+2R$.
Here $R$ is an effective symmetric character of mass four.
In particular $|\delta_-(0)|\le8$ and $|U_+(0)|\le8$.

Suppose first $W(0)\ne0$, so $W(0)=\pm8$.
The minus-cell identity
$DG_+=W(G_--D\delta_-)$ and coprimality of $W,D$
give $G_+=WQ$ with $Q$ integral, including by the content
argument if $W$ is nonprimitive.
Then $G_-=D(Q+\delta_-)$.
The plus-cell identity at zero implies $Q(0)=0$, because
its left side is $U_+(0)G_-(0)=0$ and its right side is
$W(0)^2Q(0)$.
Thus $Q=DC$ and
\[
 G_+=DWC,\qquad G_-=D\delta_-+D^2C,\qquad C(4)=2.
\]
The trace parity $G_1\equiv0$, together with $W\equiv D$
modulo two, forces $C\equiv0\pmod2$.
Dividing the plus-cell identity by $D$ now gives
\[
 U_+(\delta_-+DC)=W(WC-\psi).
\]
At zero this says
$64C(0)=U_+(0)\delta_-(0)$.
The right side has absolute value at most sixty-four.
Therefore $|C(0)|\le1$; its even parity forces $C(0)=0$.
The integral polynomial $C$ would be divisible by $D$,
contradicting $C(4)=2$.

Suppose next $W(0)=0$, so $W=DH$ with the previously
proved mass-twelve classification. Since $H(4)=3$ and
$G_-(4)=80$, the same irreducibility/content argument
as in the preceding mass-twelve rows forces $U_+=W$.
More explicitly, a nonconstant $H$ is primitive irreducible
and cannot divide $G_-$; hence it divides $U_+$.
Its odd value $H(0)\ne\pm1$ forces $U_+(0)=0$,
and the classification of $U_+/D$ gives $U_+=W$.
When $H=3$, reduction modulo three and the mass-twelve
effectivity/parity bound give the same conclusion.
Consequently
\[
 G_+=G_-+D\psi,\qquad
                  (H-1)G_-=D(H\delta_-+\psi).
\]
The necessary NN-G030 congruence is
$G_--D\delta_-\equiv G_+\pmod4$.
It forces $R\equiv\psi\pmod2$.
The effective character $R$ has mass four, while
$\psi=T_1+T_2$ already has four odd Laurent coefficients
and mass four. Hence coefficientwise effectivity and
the mass bound force $R=\psi$ exactly. In particular
$\delta_-=3\psi$ and
\[
                    (H-1)G_-=D\psi(3H+1).
\]
Put $J=(H-1)/2\in\mathbb Z[D]$.
Then $J(4)=1$. It is either the constant one (when $H=3$)
or a monic nonconstant polynomial, since the leading
coefficient of a nonconstant $H$ is two.
The preceding identity becomes
\[
                    J(G_--3D\psi)=2D\psi.
\]
Thus the monic $J$ divides $2D^2(D-3)$ in $\mathbb Z[D]$.
Its only possibilities over $\mathbb Q[D]$ are
$D^a(D-3)^b$ with $0\le a\le2$, $0\le b\le1$;
monicity fixes the scalar. Its value one at four
excludes $a>0$. Therefore $J=1$ or $J=D-3$.

If $J=1$, the identities force
\[
 W=3D,\quad G_-=5D\psi,\quad G_+=6D\psi,\quad
 S=10D^3(D-3)^2.
\]
If $J=D-3$, they force
\[
 W=D(2D-5),\quad G_-=D^2(3D-7),\quad
 G_+=2D^2(2D-5),\quad S=2D^3(3D-7).
\]
In each case $S=G_+G_-/W$ is the actual bottom
structure coefficient. Its prescribed coefficientwise
lower bound is $4D^2(D-1)^2$, whose central Laurent
coefficient is $144$.
Using $[q^0]D^h=\binom{2h}{h}$, the two forced
central coefficients of $S$ are respectively
\[
 10(252-6\cdot70+9\cdot20)=120,\qquad
                    6\cdot70-14\cdot20=140.
\]
Both are smaller than $144$, contradicting effectivity
and the given trace domination. This excludes the reverse
row and every permitted auxiliary increment, without a
support-degree restriction.
\hfill$\square$

This is an integral-character obstruction in the same
closed kernel presentation as NN-G036. The final inequality
compares actual multiplicities of the trivial torus character;
the dictionary is SGA~3, I, 4.7.2--4.7.3 and NN-G010.
No conclusion about a different multiplication table follows.

% Final NN-G036 row exclusion and exact conclusion; root whole-proof PASS 2026-11-22T00:21:30Z.
\paragraph{The last mass-twelve row.}
The only row not covered by the preceding arguments is
\[
 w=12,\quad g_1=g_2=96,\quad s=768,\quad
 \delta_1=\delta_2=\psi,\quad u_1=u_2=10.
\]
Every effective $U_i\ge D$ of mass ten with the required
parity is $D+6$ or $D+2+2T_j$. In particular
$U_i(0)\in\{6,-2\}$, so both values are nonzero.

If $W(0)\ne0$, it is $8$ or $-8$ by the mass-twelve
classification. In the pair of cell identities
\[
 WG_1-U_1G_2=DW\psi,\qquad
 WG_2-U_2G_1=DW\psi,
\]
the coefficient determinant at zero is
$64-U_1(0)U_2(0)\ne0$, because the possible products
are $36,-12,4$. Hence both $G_i$ are divisible by $D$.
After cancelling this factor, the right sides are
$W\psi$, still divisible by $D$.
Applying the same invertible matrix at zero again
shows that both $G_i$ are divisible by $D^2$.
Since $WS=G_1G_2$ and $W(0)\ne0$, it follows that
$D^4\mid S$ integrally.
The quotient has even coefficients because $S\equiv0$
modulo two, but its value at four is
$768/4^4=3$. This is impossible.

It remains that $W=DH$, where $H(4)=3$,
$H\equiv1\pmod2$, and $H$ is either three
or primitive irreducible with leading coefficient two.
If $H$ is nonconstant, it cannot divide either
$U_i$, since three does not divide ten.
The common-factor argument of (G036.R2) gives
$H\mid G_1,G_2$ integrally.
If $H=3$, reduce both cell identities modulo three.
The reductions of the $U_i$ are nonzero, since their
masses are ten; the integral domain $\mathbb F_3[D]$
therefore forces the reductions of both $G_i$ to
vanish. Thus $H\mid G_1,G_2$ also in this case.
Write $G_i=HP_i$. Then
\[
 P_i(4)=32,\qquad
 P_1\equiv0,\quad P_2\equiv D^2(D-1)\pmod2,
\]
and the divided cell identities are
\[
 DH P_1-U_1P_2=D^2\psi,\qquad
 DH P_2-U_2P_1=D^2\psi.
\]
Because each $U_i(0)\ne0$, evaluation at zero
forces both $P_i(0)=0$.
Writing $P_i=DQ_i$ and cancelling one $D$ gives
\[
 DH Q_1-U_1Q_2=D\psi,\qquad
 DH Q_2-U_2Q_1=D\psi.
\]
Evaluation at zero again forces both $Q_i(0)=0$.
Consequently $P_i=D^2R_i$ with integral $R_i$.
Their masses are $R_i(4)=32/16=2$.
But $P_2\equiv D^2(D-1)$ forces
$R_2\equiv D-1\pmod2$, whose evaluation at four
is odd. This contradicts $R_2(4)=2$.
The last row is excluded for every grading support.
\hfill$\square$

\paragraph{Exact closed-table optimum and geometric conclusion.}
The complete necessary mass gate of NN-G035,
followed by all the uniform exclusions of NN-G036,
leaves exactly the unique $2744$-pair realization
of NN-G033--NN-G034. Therefore the minimum number
of cancelling trivial/sign pairs in the specified
symmetric closed-table presentation is exactly
$2744$, and its polynomial parameter realization
is unique within that presentation.
All cell degrees and contents, all grading support
locations, and all rational auxiliary parameters
were allowed throughout, subject only to integrality
and effectivity of the actual structure constants.
The final certificate projects every surviving
ordinary-dimension row; it does not enumerate
bounded grading degrees to justify the exclusions.

In the NN-G010 category this is an exact minimum
for actual effective graded sign-representation
objects composing by the fixed kernel table.
The grading is the diagonalizable-group dictionary
of SGA~3, I, 4.7.2--4.7.3, and composition is
the supplied archive cote~2, scan~17 transcription.
The theorem concerns this presentation, not all
possible effective multiplication tables.
The unrestricted extra-position optimization still
has the actual upper bound $2744$, the finite
equal-score cutoff of NN-G029/NN-G033, and the
total-mass FPT algorithm of NN-G032; its numerical
minimum remains open.




\section{NN-G037: an unrestricted 104-pair lower bound with shared quadratic products}
\paragraph{Statement.}
Every effective associative dimension lift of the supplied
rank-ten trace table, with the prescribed basis element
$\mu$ as normalized unit, has at least $104$ cancelling pairs.
All $810$ nonunit-input structure positions are allowed,
including positions with identically zero trace.
No symmetry of the grading and no closed-table shape
is assumed. Together with NN-G033 this gives the genuine
unrestricted interval
\[
                              104\le m_{\rm all}\le2744.             \tag{G037.1}
\]
The lower bound is modest and does not assert attainment.
The initial certificate below proves $61$ pairs. The shared-quadratic
addendum then strengthens it to $104$ with an exact integer certificate.

\paragraph{Ordinary masses and the proposed small budget.}
Index the ten basis elements in the order
\[
 \mu,\Sigma,\gamma_{1,1},\gamma_{1,12},\gamma_{1,2},
 \gamma_{1,21},\gamma_{2,1},\gamma_{2,12},
 \gamma_{2,2},\gamma_{2,21}.
\]
Let $c_{ab}^z(q)$ be the given trace structure coefficients
and put $A_{ab}^z=\sum_e|[q^e]c_{ab}^z|$.
The supplied table has
$\sum_{a,b,z}A_{ab}^z=4443$, including the nineteen
unit positions.
For an effective lift, let $d_{ab}^z$ be its ordinary
dimension, obtained by evaluating its dimension character
at $q=1$. Write
\[
 E_{ab}^z=d_{ab}^z-A_{ab}^z.
\]
Each $E_{ab}^z$ is a nonnegative even integer;
the total number of cancelling pairs is
$\frac12\sum E_{ab}^z$.
The normalized unit fixes $E_{0b}^z=E_{a0}^z=0$.
Thus there are $9\cdot9\cdot10=810$ free positions.
Suppose for contradiction that the pair count is at most
sixty, so their total added mass satisfies
\[
                         \sum E_{ab}^z\le M:=120.                  \tag{G037.2}
\]

\paragraph{Necessary inequalities with integer coefficients.}
For each triple of nonunit inputs and each output, evaluate
associativity at $q=1$:
\[
 \sum_y(A_{ab}^y+E_{ab}^y)(A_{yc}^z+E_{yc}^z)
 =
 \sum_y(A_{bc}^y+E_{bc}^y)(A_{ay}^z+E_{ay}^z).
\]
Choose either orientation and call its sides $L,R$.
Expand them as
$L=L_0+\ell_L(E)+Q_L(E)$ and
$R=R_0+\ell_R(E)+Q_R(E)$,
where the three parts are constant, linear, and quadratic.
All quadratic coefficients are nonnegative.
For any two distinct free variables $x,y$, (G037.2) gives
\[
 xy\le\frac{(x+y)^2}{4}\le\frac M4(x+y);
 \qquad x^2\le Mx.
\]
Apply these bounds separately to the monomials of $Q_R$,
and write the resulting upper bound as
$Q_R(E)\le(M/4)\sum_e r_eE_e$.
Here each occurrence of a product of distinct variables
adds one to each of their $r_e$, while a square adds four.
Terms containing a fixed unit-position variable vanish.
Since $Q_L(E)\ge0$, associativity implies the necessary
integer linear inequality
\[
 4(L_0-R_0)
 \le \sum_e
 \bigl(4(\ell_{R,e}-\ell_{L,e})+M r_e\bigr)E_e.       \tag{G037.3}
\]
This remains valid when its left side is negative.

\paragraph{The complete finite certificate.}
The file
\path{unrestricted_mass_lower_bound_certificate.json}
lists seventy-four nonnegative integer multipliers,
each at most $565$, together with the five indices
$(a,b,c,z,\text{orientation})$ selecting (G037.3).
The independent exact verifier reconstructs $A$ directly
from the original supplied
\path{OURS_E4_scripts/e3_constants.json}, expands
those seventy-four inequalities, and checks that their
weighted sum has left side
\[
 B=12{,}209{,}088
\]
and that every one of the $810$ weighted coefficients
on its right side is at most
\[
 C=101{,}136.
\]
All variables are nonnegative, so (G037.2)--(G037.3)
would force
\[
 12{,}209{,}088=B
 \le C\sum_eE_e
 \le120C=12{,}136{,}320.
\]
The strict gap is $72{,}768$, a contradiction.
Hence $\sum E>120$. Since the sum is even, it is at
least $122$, proving at least $61$ cancelling pairs.
The corresponding total dimension mass is at least
$4443+122=4565$.
\hfill$\square$

The exploratory linear-programming solver only found the
integer multipliers; no numerical optimization assertion
is an input to the proof.
The verifier uses exact integers and the original trace
table, imports no optimization library, and does not
reuse the exploratory inequality builder.

\paragraph{Geometric translation and limits.}
For the graded sign-representation structure objects of
NN-G010, forgetting both grading and the sign action
gives their ordinary nonnegative integer dimensions.
Tensor composition and direct sum make these dimensions
an associative multiplication table.
The initial certificate forces at least $61$ trivial/sign cancellation pairs
among the structure objects of every such effective lift,
even when zero-trace structure positions are populated.
This is a lower bound on actual effective multiplicities,
not on virtual Grothendieck-group representatives.
The grading dictionary is SGA~3, I, 4.7.2--4.7.3;
the kernel-composition interpretation is the supplied
archive cote~2, scan~17 transcription.
No categorical rigidity, canonical-basis property, rank-four
result, or complexity-class separation is inferred.
The shared-quadratic refinement below improves this lower bound;
the unrestricted optimum is not asserted.

\paragraph{Shared quadratic products strengthen the unrestricted bound to $104$.}
Every normalized-unit effective lift in the full $810$-position problem
has at least $104$ cancelling pairs.  No original-support, grading
symmetry, or closed-table hypothesis is imposed.  The proof retains
shared quadratic products across different associativity equations,
instead of bounding the positive part of each equation separately.
This is a refinement of NN-G037, not a new support-class theorem.

Use the ordinary masses $A$ and the $810$ free nonnegative even
integers $E_i$ already defined in NN-G037.  Suppose that at most $103$
pairs suffice, and set
\[
                         \sum_{i=1}^{810}E_i\le M=206. \tag{G037.4}
\]
For each of the $7290$ nonunit-input associators, expand its exact
ordinary-dimension equality in the form
\[
 L_k(E,Q)=b_k,\qquad
 b_k=R_k(A)-L_k(A),\qquad Q_{ij}=E_iE_j=Q_{ji}. \tag{G037.5}
\]
Here $L_k(E,Q)$ is the difference of the linear and quadratic parts
of the left and right products; it has integral coefficients.
The symbol $L_k$ in this displayed linearized expression does not
include the constant term.  Positions with a unit input remain
fixed, including their possible contribution to its linear part.

Let $\mathcal Q$ consist of every unordered variable pair with a
nonzero quadratic coefficient in at least one of these associators,
together with all $810$ diagonal pairs.  Its size for the supplied
table is $112302$.  For each variable $i$, let
$\mathcal Q_i=\{\{i,j\}\in\mathcal Q\}$, so a diagonal pair occurs
only once.  Under (G037.4), the following are nonnegative:
\[
\begin{aligned}
 g_i^{\rm diag}&=Q_{ii}-2E_i,\\
 g_{ij}^{\rm pair}&=M(E_i+E_j)-4Q_{ij}\qquad(i\ne j),\\
 g_i^{\rm budget}&=ME_i-\sum_{\{i,j\}\in\mathcal Q_i}Q_{ij}.
\end{aligned}
\tag{G037.6}
\]
The first assertion uses that a nonzero even integer is at least two.
For the second, $4E_iE_j\le(E_i+E_j)^2\le M(E_i+E_j)$.
For the third, the sum contains only distinct terms from
$E_i\sum_j E_j$, and all omitted terms are nonnegative.
No independence of the products and no probabilistic moment
interpretation is used.

\paragraph{An exact weighted contradiction.}
The certificate
\path{unrestricted_shared_quadratic_104_certificate.json}
lists $334$ nonzero signed integer weights $\lambda_k$ for (G037.5)
and $1324$ positive integer weights $\mu_r$ for selected bounds
in (G037.6).  Put
\[
 B=\sum_k\lambda_k b_k,\qquad
 \sum_k\lambda_k L_k(E,Q)+\sum_r\mu_r g_r(E,Q)
       =\sum_i a_iE_i+\sum_{\{i,j\}\in\mathcal Q}c_{ij}Q_{ij}.
\tag{G037.7}
\]
The exact equalities and nonnegative bounds imply that the last
expression is at least $B$.
Discarding each term with $c_{ij}<0$ increases it.  For the remaining
positive coefficients use
\[
 Q_{ij}\le\frac M4(E_i+E_j)\quad(i\ne j),\qquad
 Q_{ii}\le ME_i.
\]
Consequently
\[
 4B\le\sum_i C_iE_i,\qquad
 C_i=4a_i+
       M\!\sum_{j\ne i}\max(c_{ij},0)
       +4M\max(c_{ii},0). \tag{G037.8}
\]
These are exact integer formulas; negative values of $a_i$ cause
no problem because the final comparison uses $E_i\ge0$.

The independent verifier
\path{verify_unrestricted_shared_quadratic.py}
reads the original supplied trace table and reconstructs all
source masses, each selected associator, the full graph $\mathcal Q$,
and every weighted bound.  It does not import an optimizer or the
exploratory matrix builder.  Its exact result is
\[
\begin{aligned}
 4B&=8{,}241{,}282{,}112,\\
 \max_i C_i&=40{,}001{,}488,\\
 M\max_iC_i&=8{,}240{,}306{,}528,\\
 4B-M\max_iC_i&=975{,}584>0.
\end{aligned}
\tag{G037.9}
\]
There are $40$ positive quadratic remainders in (G037.8), none
diagonal; all are explicitly accounted for in the corrected
coefficients.  Equations (G037.4) and (G037.8) would force
$4B\le M\max_iC_i$, contradicting (G037.9).
The added mass is therefore greater than $206$ and, being even,
at least $208$.  This proves the asserted $104$-pair lower bound.

The floating solver was used only to locate candidate weights.
The recovery script rounds those candidates to integers; the
certificate above contains only integers, and its strict contradiction
is checked independently.  The finder used a smaller product graph.
The verifier deliberately uses the entire graph $\mathcal Q$:
the extra terms in any budget bound enter with negative coefficients,
so they remain valid and cannot weaken the final upper bound.
No inference from a numerical infeasibility report is needed.

\paragraph{Geometric conclusion and remaining unrestricted gate.}
For every effective graded sign-representation lift of the prescribed
trace, ordinary rank gives an associative nonnegative multiplication
table.  The polynomial map $E\mapsto(E_i,E_iE_j)$ places this table
in the corresponding quadratic graph in affine space; its associator
equations are linear after adjoining these product coordinates.
Nonnegative even ranks and the proposed total-rank cap imply
(G037.6), whose exact weighted consequence excludes that entire
quadratic-graph slice.  Thus actual structure objects across all
$810$ allowed positions require at least $104$ cancelling pairs,
even if zero-trace positions are populated.  The SGA~3 I4.7.2--4.7.3
grading dictionary and NN-G010 actual composing-kernel interpretation
(supplied archive cote~2, scan~17 transcription) are unchanged.
The new unrestricted interval is $[104,2744]$.  The stronger
original-$132$-support interval $[900,2744]$ and the exact symmetric
closed-table minimum $2744$ concern different loci and are not
used to prove this unrestricted bound.




\section{NN-G038: a complete 19-parameter classification for the original 132-position support}
\paragraph{Statement and coefficient fields.}
Let $R=\mathbb Z[q,q^{-1}]$ and $K=\mathbb Q(q)$.
For symmetric Laurent coefficients one may instead use
$R^{\rm inv}=\mathbb Z[D]$ and $K^{\rm inv}=\mathbb Q(D)$,
where $D=q+2+q^{-1}$.
Every associative normalized-unit table having exactly the
original $132$ nonzero positions is uniquely determined over
$K$ by nineteen nonzero scalar parameters:
five for each four-element cell, eight bottom-action values,
and the bottom square value.
Conversely the formulas below give an associative table
for every such parameter tuple; the required nonzero
bottom coefficients select an explicit open subset.

Consequently the entire effective dimension-lift problem
with this support is equivalent to finitely many explicit
rational-function conditions on those nineteen parameters:
every resulting coefficient must lie in $R$, dominate the
absolute trace coefficientwise, and agree with the trace
modulo two. For symmetric lifts the parameters and all
coefficients are required to be invariant under $q\mapsto q^{-1}$.
This removes associativity as a further equation-solving
gate in this whole support class. It does not solve the
remaining integral positivity optimization.

\paragraph{The precise support pattern.}
Write the basis as
$\mu,\Sigma,e_{i,rs}$, with $i\in\{1,2\}$ and
$r,s\in\{1,2\}$, where
\[
 e_{i,11}=\gamma^i_1,\quad e_{i,12}=\gamma^i_{12},
 \quad e_{i,21}=\gamma^i_{21},\quad e_{i,22}=\gamma^i_2.
\]
The unit products are fixed. A product involving $\Sigma$
and another nonunit has only a $\Sigma$ output.
Products from distinct cells have only a $\Sigma$ output.
Within cell $i$ the product of $e_{i,rs}$ and $e_{i,tv}$
has exactly two possible outputs: $e_{i,rv}$ and $\Sigma$.
All the specified $132$ coefficients are nonzero.
The original trace table has this exact pattern; the
independent interface checker verifies all its positions.

\paragraph{The bottom summand is forced to split over the field.}
Write $\Sigma^2=S\Sigma$, with $S\ne0$, and initially allow
$\Sigma b=\ell_b\Sigma$, $b\Sigma=r_b\Sigma$.
Associativity of $(\Sigma b)\Sigma=\Sigma(b\Sigma)$ gives
$\ell_bS=r_bS$. Hence $\ell_b=r_b=:\theta_b$.
For the eight cell basis elements these are the eight
bottom-action parameters; $\theta_\mu=1$ and
$\theta_\Sigma=S$ are fixed by normalization.
The identity $\Sigma(ab)=(\Sigma a)b$ now gives
\[
 \theta_a\theta_b
   =c_{ab}\theta_{\operatorname{out}(a,b)}+S B_{ab}
       \quad\hbox{within a cell},                  \tag{G038.1}
\]
where $c_{ab}$ is its cell coefficient and $B_{ab}$
its bottom coefficient. In distinct cells this reads
$\theta_a\theta_b=SB_{ab}$.
Thus
\[
 B_{ab}=
 \begin{cases}
 (\theta_a\theta_b-c_{ab}\theta_{\operatorname{out}(a,b)})/S,
       &a,b\text{ in one cell},\\
 \theta_a\theta_b/S,&a,b\text{ in distinct cells}.
 \end{cases}                                                     \tag{G038.2}
\]

For the actual splitting, put $e=\Sigma/S$.
It is a central idempotent, and
$\widetilde b=b-\theta_b e$ satisfies
$e\widetilde b=\widetilde b e=0$.
Equation (G038.1) gives
$\widetilde a\,\widetilde b
 =c_{ab}\widetilde{\operatorname{out}(a,b)}$
inside one cell, and zero between cells.
The complementary normalized unit is $\mu-e$.
No division by $S$ is claimed to preserve the original
integral or effective lattice: this splitting takes place
only over the stated fraction field.

\paragraph{Five parameters classify a four-element cell.}
Suppress the cell index and write
$a=e_{11},b=e_{12},c=e_{21},d=e_{22}$.
Its multiplication has the form
$e_{rs}e_{tv}=c_{rs,tv}e_{rv}$, with every coefficient nonzero.
Define the five parameters by actual table entries:
\[
 a^2=u a,\quad ac=x a,\quad ba=y a,\quad
 bc=z a,\quad cb=p d.                              \tag{G038.3}
\]
They are nonzero and therefore invertible in $K$.
Associativity of $a^2b=a(ab)$ forces $ab=u b$:
if $ab=\lambda b$, then $u\lambda=\lambda^2$
and $\lambda\ne0$. Likewise $ca=u c$.
More generally these two identities are exactly the
instances needed below of
$e_{r1}a=u e_{r1}$ and $ae_{1s}=u e_{1s}$.

Define
\[
 E_{rs}=e_{r1}e_{1s}=s_{rs}e_{rs},\qquad
 s_{11}=s_{12}=s_{21}=u,\quad s_{22}=p,
 \qquad M=\begin{pmatrix}u&x\\y&z\end{pmatrix}.
\]
Associativity gives the complete product directly:
\[
 \begin{aligned}
 E_{rs}E_{tv}
 &=e_{r1}(e_{1s}e_{t1})e_{1v}\\
 &=u M_{st}\,e_{r1}e_{1v}
  =u M_{st}E_{rv}.
 \end{aligned}
\]
Here the four coefficients of $e_{1s}e_{t1}$
are precisely the four entries of $M$ in (G038.3).
Returning to the original basis yields
\[
                   c_{rs,tv}
       =\frac{u M_{st}s_{rv}}{s_{rs}s_{tv}}.          \tag{G038.4}
\]
This proves necessity and uniqueness of all sixteen cell
coefficients from the five selected entries.

Conversely, for any five nonzero parameters, the product
$E_{rs}E_{tv}=uM_{st}E_{rv}$ is associative:
both associations of three basis elements have coefficient
$u^2M_{st}M_{vw}$ and the same output.
Equivalently it is the matrix product $X\circ Y=X(uM)Y$.
Rescaling by the nonzero $s_{rs}$ gives (G038.4).
No invertibility of $\det M$ is needed; no semisimplicity
or matrix-algebra isomorphism is being silently assumed.

\paragraph{Sufficiency for the full table and restoration of the unit.}
Take the two cell algebras from (G038.4).
Form their direct sum with zero cross-products and then
its external unitalization $C$, of dimension nine.
In the associative algebra $C\times K$, set
\[
 \mu=(1_C,1),\qquad \Sigma=(0,S),\qquad
 e_{i,rs}=(\overline e_{i,rs},\theta_{i,rs}).
\]
These ten vectors form a basis: their projections to $C$
are its unit and eight independent cell vectors, while
$S\ne0$ supplies the remaining scalar direction.
The vector $\mu$ is exactly the normalized unit.
Multiplication gives $\Sigma^2=S\Sigma$,
$\Sigma b=b\Sigma=\theta_b\Sigma$, the cell coefficients
(G038.4), and exactly the bottom coefficients (G038.2).
This proves all associativity identities at once and
establishes sufficiency.

\paragraph{Exact integral and effective conditions.}
The nineteen parameters are recovered from actual
coefficients by (G038.3), $\theta_b$, and $S$.
Thus a required dimension lift is obtained if and only if
the following finite list holds for their formulas.

\begin{enumerate}
\item The unit entries are exactly those of $\mu$;
all unspecified positions are zero.
\item The $113$ specified nonunit-input coefficients
given by (G038.2)--(G038.4), $S$, and the two copies of
each $\theta_b$ belong to $\mathbb Z[q,q^{-1}]$.
\item For each specified position and every Laurent exponent,
the resulting coefficient is at least the absolute value
of the prescribed trace coefficient and is congruent
to that trace coefficient modulo two.
\item When the symmetric problem is intended, the parameters,
equivalently the recovered table, are fixed by
$q\mapsto q^{-1}$.
\end{enumerate}

Condition three ensures effectivity and the exact nonzero
support, because every specified trace coefficient is
a nonzero Laurent polynomial.
Admissibility also forces each of the nineteen recovered
parameters to be an integral effective Laurent polynomial:
each parameter is itself an actual structure coefficient.
It does not require positivity of any quotient appearing
in the fraction-field parametrization.
All rational expressions are explicit; testing them for a
given finite tuple is different from finding an optimizing
tuple of unbounded degree.

For completeness, keep the prescribed associative trace
table $t$ fixed. For each admissible dimension coefficient $d$,
the lift
\[
                  \widehat d=(d+t)/2+\epsilon(d-t)/2
\]
has nonnegative integral Laurent coefficients by domination
and parity. Its associativity follows by evaluating at
$\epsilon=1,-1$: these recover the associative tables $d,t$,
and the two evaluations together are injective on
$R[\epsilon]/(\epsilon^2-1)$ because $R$ has no two-torsion.
The normalized unit also lifts to the normalized unit.

\paragraph{This class strictly contains the fixed-$D$ closed presentation.}
Start with the actual NN-G033 lift and take any odd integer
$k\ge3$. Give $\gamma^1_1,\gamma^1_2$ weight two,
all other nonunit basis elements weight one, and $\mu$
weight zero. In the dimension table alone, replace each
basis vector $b$ over $K$ by $k^{w_b}b$, while keeping
the prescribed trace table fixed. Every nonunit dimension coefficient is
multiplied by
$k^{w_a+w_b-w_z}$, whose exponent is nonnegative
because input weights are at least one and output weights
are at most two.
The factor is an odd positive integer, so all coefficient
domination and parity conditions survive, and the support
and normalized unit remain unchanged.
Associativity survives the basis change.
The reconstruction just proved then supplies its effective
$C_2$-lift with the original trace table.
But the first-cell square coefficient becomes
$k^2D\delta_1$, whereas its base cross coefficient
becomes $k^3\delta_1$.
Their ratio is $D/k$, rather than $D$.
Hence this effective lift lies outside the presentation
optimized in NN-G036. No smaller score is claimed.

\paragraph{Algebraic-geometric translation and the remaining gate.}
On the locus where all $113$ specified nonunit coefficients
are invertible, the associativity scheme is explicitly an
open subset of the nineteen-dimensional split torus:
the nineteen recovered coordinates and formulas
(G038.2)--(G038.4) are mutually inverse.
The additional open conditions invert the thirty-two
within-cell bottom numerators; all other required
coefficients are products or ratios of already invertible
coordinates.
This assertion is valid over arbitrary commutative base
rings on that invertible locus, since every cancellation
used in the proof is by an explicitly inverted element.
In affine terms its coordinate ring is a localization of
$\mathbb Z[t_1^{\pm1},\ldots,t_{19}^{\pm1}]$.
An integral lift gives a $K$-point whose reconstructed
coefficients lie in $R$; it need not give an $R$-point
of this localized scheme, since the nineteen coefficients
need not be units in $R$.

The integral grading-character interpretation is SGA~3,
I, 4.7.2--4.7.3; actual effective coefficients give the
matrix-kernel objects of NN-G010, with composition in
the supplied archive cote~2, scan~17 transcription.
The open torus describes the algebraic associativity
parameters, not the coefficientwise positive integer
lattice. The exact $2744$ optimum from NN-G036 applies
to its smaller fixed-$D$ subfamily only.
The remaining support-preserving optimization is now
precisely the nineteen-parameter integral-effectivity
problem above. Tables populating any additional position
remain a separate part of the full $810$-position problem.

\paragraph{A fixed-anchor obstruction inside the full support class.}
This refinement uses the nineteen-parameter classification, not the
fixed-$D$ closed presentation. Set
\[
 C=D-1,\qquad K=D-2,\qquad W=2D^2-7D+4,\qquad H=2DK.
\]
Consider a symmetric effective lift with the original support, and
fix only the following three recovered parameters to their values
in NN-G033:
\[
 S=D^2CH,\qquad \theta_{2,11}=\theta_{2,22}=G=DCW.
                                                               \tag{G038.5}
\]
There is no assumption on the first cell, either word action, or
any of the five second-cell multiplication parameters.
Then both word actions have ordinary dimension at least $384$.
Moreover the second-cell square coefficient $u=c_{11,11}$ is
already forced by these three anchors:
\[
             u=D(CW-H).                                        \tag{G038.6}
\]
Thus one cannot lower either large word-action dimension of the
incumbent by varying all the remaining parameters while retaining
just (G038.5). This is a local obstruction, not a lower bound for
the entire support-preserving optimization.

\paragraph{Forcing the square coefficient.}
The characters $W$ and $DH=2D^2K$ are relatively prime over
$\mathbb Q[D]$, and $W$ is primitive. Hence
$Wf/(DH)\in\mathbb Z[D]$ implies $f/(DH)\in\mathbb Z[D]$:
Gauss's lemma cancels the primitive coprime polynomial factors,
and multiplication by the primitive $W$ cannot remove a remaining
rational denominator from an integral polynomial.
The bottom coefficient of $a^2$, with $a=e_{2,11}$, is
\[
 B_{aa}=\frac{G(G-u)}S=\frac{W(G-u)}{DH}=WJ,
 \qquad G-u=DHJ,\quad J\in\mathbb Z[D].             \tag{G038.7}
\]
Its target is $D$, of mass $4$, whereas the absolute target of
$u$ is $DK(D-3)$, of mass $8$.
At $D=4$ the anchors give $G(4)=96$, $(DH)(4)=64$,
and $W(4)=8$. Therefore
\[
       8J(4)\ge4,\qquad 96-64J(4)\ge8,
\]
so $J(4)=1$. Reduction of $B_{aa}=WJ$ modulo two gives
$DJ\equiv D$, hence $J\equiv1\pmod2$ in $\mathbb F_2[D]$.
Both $WJ$ and $u=G-DHJ$ are effective Laurent polynomials.
If $\deg J>1$, their leading coefficients have opposite signs,
since $\deg G=4$ and $\deg(DH)=3$.
If $\deg J=1$, its leading coefficient is positive and even,
whereas the leading coefficient of $u$ is
$2-2\operatorname{lc}(J)<0$.
Thus $J$ is constant, and $J(4)=1$ gives $J=1$.
This proves (G038.6).

\paragraph{The word action reduces to a positive character remainder.}
Write $b=e_{2,12}$, $c=e_{2,21}$, and $d=e_{2,22}$.
The classification forces $ab=ub$ and $ca=uc$ at the cell level.
For $F=\theta_b$, its $ab$ bottom coefficient is therefore
\[
       P=\frac{(G-u)F}{S}=\frac F{DC}.                         \tag{G038.8}
\]
The identical argument uses $ca$ for $F=\theta_c$.
In both cases the prescribed bottom coefficient is $D^2$.
Consequently $P\in\mathbb Z[D]$ and, as Laurent characters,
\[
          P=D^2+2R,\qquad R\in\mathbb N[q,q^{-1}],
          \quad R(q)=R(q^{-1}).                                \tag{G038.9}
\]
Suppose $F(4)<384$. Since $(DC)(4)=12$, we have
$P(4)<32$ and $r:=R(4)<8$.

Let $\delta$ be the cell coefficient of $ad$ with output $b$,
or of $da$ with output $c$, as appropriate to the chosen $F$.
It is an actual integral coefficient, and
$\delta\equiv D(D-3)\pmod2$.
The corresponding bottom coefficient is
\[
       \frac{G^2-\delta F}{S}
          =\frac{DCW^2-\delta P}{DH}\in\mathbb Z[D].             \tag{G038.10}
\]
If $P(0)\ne0$, divisibility by $D$ first forces $\delta(0)=0$.
Write $\delta=DQ$; its parity gives $Q\equiv D-3\pmod2$.
Divisibility by the second factor of $D$ then gives
\[
                   Q(0)P(0)=-16.
\]
Since $Q(0)$ is odd, it equals $1$ or $-1$ and
$P(0)=\mp16$. But $D=0$ corresponds to $q=-1$, so
$|P(0)|=2|R(-1)|\le2R(1)=2r<16$, a contradiction.
Hence $P(0)=R(-1)=0$.

Expand the symmetric effective character as
$R=c_0+\sum_{j\ge1}c_j(q^j+q^{-j})$ with $c_j\ge0$.
The equation $R(-1)=0$ gives
\[
              r=4\sum_{j\ {
m odd}}c_j.
\]
Thus $r=0$ or $4$. If $r=4$, the expansion is either
$R=2+q^o+q^{-o}$, or
$R=q^o+q^{-o}+q^e+q^{-e}$, where $o$ is positive odd
and $e$ is positive even. At $D=2$, equivalently $q=i$,
these possibilities give $R(2)=2$ or $-2$, hence
$P(2)=8$ or $0$. But (G038.10) at $D=2$ forces
\[
                       \delta(2)P(2)=8,
\]
which is impossible: $\delta(2)$ is even by its parity.
This excludes every support location $o,e$, not merely a
finite degree window.

It remains to exclude $r=0$, namely $P=D^2$ and $F(4)=192$.
For the selected word $b$, use its $ba$ bottom coefficient;
for $c$, use its $ac$ bottom coefficient. If $x$ is the
corresponding actual cell coefficient, each of these is
\[
       B=\frac{G(F-x)}S=WZ,\qquad
              F-x=DHZ,\quad Z\in\mathbb Z[D],                  \tag{G038.11}
\]
by the same primitive divisibility argument.
Its target is $V=2D^2-5D+4$, an effective character of mass $16$.
The absolute target of $x$ has mass $16$.
Hence, at $D=4$,
\[
                  16\le B(4)=\frac{192-x(4)}8\le22.
\]
But $B(4)=8Z(4)$ is an integral multiple of $8$, so
$B(4)=16$. Coefficientwise domination of the effective
character $V$ and equality of their masses force $B=V$.
This contradicts $W\mid B$, because $W$ and $V$ are distinct
quadratics with the same leading coefficient; equality would
be the only possible integral polynomial multiple.
All cases with $P(4)<32$ are excluded, proving $F(4)\ge384$
for each word separately.

\paragraph{Geometric meaning and limitation of the anchor result.}
Within the nineteen-dimensional parameter torus of NN-G038,
(G038.5) specifies a three-coordinate fiber over $\mathbb Q(D)$.
On the integral effective-character points of this fiber,
the square coordinate is forced to (G038.6), and both word-action
coefficient objects have rank at least $384$ after forgetting
the grading. Equality of the two word actions was never assumed.
The restrictions $q=-1,i$ are restrictions of grading characters
to $\mu_2$ and $\mu_4$, in the SGA~3 I4.7.2--4.7.3 dictionary;
actual coefficient objects compose as in NN-G010 and the supplied
archive cote~2, scan~17 transcription. This fiber result does not
restrict other choices of $S$ or the two generator actions, and
it asserts neither a global optimum nor uniqueness of the word
characters of mass $384$.

\paragraph{A global lower bound for the entire original-support class.}
Every normalized-unit effective lift with exactly the original
$132$ nonzero positions has at least $344$ cancelling pairs.
This includes arbitrary finite Laurent gradings: neither
$q\leftrightarrow q^{-1}$ symmetry nor any of the three anchors
in (G038.5) is assumed. Together with NN-G033, this first global step gives the interval
$[344,2744]$, strengthened to $[560,2744]$ in the following
full-cell addendum.
The separate all-$810$-position problem still has only the
proved lower bound $61$ from NN-G037.

Here is a finite exact certificate and its complete reduction.
All masses in this argument are evaluations at $q=1$ of actual
effective dimension characters. The prescribed trace has total
absolute mass $4443$, of which $19$ comes from the normalized
unit. If $T$ denotes the mass in the $113$ specified nonunit-input
positions, its cancelling-pair score is $(T-4424)/2$.
Write $s=S(1)$ and write the four action masses in the two
cells as $x_1,\ldots,x_4$ and $y_1,\ldots,y_4$.
The two generator actions in each cell have target mass $48$;
the two word actions have target mass $192$.
Every mass is an even integer, by the source trace parities.
The source constants used below are independently reconstructed
from the original Laurent table by the verifier.

\paragraph{The bottom square has mass divisible by four.}
This restriction uses the full characters, not just their masses.
Use a generator from each cell and put
$G_1=\theta_{1,11}$, $G_2=\theta_{2,11}$, and let $W_0$
be their cross-cell bottom coefficient. The classification gives
$SW_0=G_1G_2$. In $\mathbb F_2[q,q^{-1}]$ the prescribed
parities are
\[
 S\equiv G_1\equiv0,\qquad
 G_2\equiv D^2(D-1),\qquad W_0\equiv D.
                                                               \tag{G038.12}
\]
Let $k$ be the minimum two-adic valuation of the coefficients
of $G_1$. It is finite and at least one. Since the Laurent
ring modulo two is a domain and neither $G_2$ nor $W_0$
vanishes in it, the coefficient valuation of $S$ is also $k$.
Divide the product identity by $2^k$ and reduce modulo two.
Cancellation of the nonzero $D$ gives
\[
          S/2^k\equiv D(D-1)(G_1/2^k)\pmod2.
\]
Evaluation at $q=1$ makes its right side zero modulo two.
Hence $2^{k+1}\mid s$, and in particular $4\mid s$.
No factorization or symmetry hypothesis on $S$ was used.

\paragraph{Cross-cell integrality reduces the action lattice to divisors.}
Put $A=\sum_i x_i$ and $B=\sum_j y_j$.
Every cross-cell bottom mass is $x_i y_j/s$, and it is an
even integer. Thus $2s\mid x_i y_j$ for every pair.
Let $g=\gcd(x_1,\ldots,x_4)$ and $a=\gcd(g,2s)$.
Then $a\mid2s$, all $x_i$ are multiples of $a$, and all
$y_j$ are multiples of $b=2s/a$.
Indeed $2s\mid gy_j$ by Bezout, and
$\gcd(g/a,b)=1$. Consequently it suffices to consider all
positive divisors $a$ of $2s$, assigning the action steps
\[
          \operatorname{lcm}(a,2),\qquad
          \operatorname{lcm}(2s/a,2)                              \tag{G038.13}
\]
to the two cells. This is a necessary covering of all action
tuples, not a claim that every divisor choice is realizable.
The total mass of the $32$ cross-cell coefficients is exactly
$2AB/s$.

\paragraph{Self-square costs and exhaustive budget bounds.}
For any cell basis vector with action mass $x$, write $c$ for
its cell output mass in its square and $b_0'$ for its bottom
output mass. Associativity with $\Sigma$ gives
\[
                x^2=cx+s b_0'.                                 \tag{G038.14}
\]
The two output masses are even. Their lower bounds $(c_0,b_0)$
are $(8,4)$ for a generator in either cell, $(8,64)$ for a word
in the first cell, and $(16,64)$ for a word in the second cell.
For fixed $s,x$, define the necessary minimum self-square cost
\[
 E_{s,c_0,b_0}(x)=
 \min\{c+b_0':c\ge c_0,\ b_0'\ge b_0,
             \ c,b_0'\in2\mathbb Z,\ x^2=cx+s b_0'\}.             \tag{G038.15}
\]
An empty set excludes that action value. The minimum is
computed by the direct finite enumeration of even
$c_0\le c<x$, retaining precisely those $c$ for which
$(x^2-cx)/s$ is an even integer at least $b_0$.
There is no nonlinear numerical optimization in this calculation.

Suppose there are at most $343$ pairs. Then $T\le5110$,
and the total excess over the absolute trace masses is at most
$686$. Thus
\[
 s\in\{576,580,\ldots,1260\},\qquad
 R_s=686-(s-576).
\]
Each action value with target $m\in\{48,192\}$ lies between
$m$ and $m+\lfloor R_s/2\rfloor$, because it occurs twice
among the structure coefficients. The sum of the four actions
in either cell is at most $480+\lfloor R_s/2\rfloor$.
These bounds deliberately ignore the excess already spent
on other positions, so no feasible tuple is lost.

For each $s$ and divisor $a\mid2s$, form a finite dynamic
program for each cell, with its action step from (G038.13).
Start with the entry $0\mapsto0$. Process two generator roles
and then two word roles. For each allowed action $x$ with
finite value (G038.15), add $x$ to the sum and its self-square
cost to the accumulated cost, retaining the smallest cost
at every sum. Discard sums above
$480+\lfloor R_s/2\rfloor$.
Call the resulting arrays $E_1(A)$ and $E_2(B)$.
This considers the four roles independently; no equality
between generator actions or between word actions is imposed.
Induction on the processed roles proves that each retained
entry is exactly the least sum of their self-square costs
at that total action mass.

The absolute trace masses of all the remaining, non-square
within-cell coefficients sum to $776$. The original eight
self-square costs sum to $352$; the cross-cell trace masses
sum to $800$. These agree with the full accounting
$576+1920+352+800+776=4424$.
It follows for every feasible divisor row and retained pair
of action sums that
\[
       T\ge s+2(A+B)+\frac{2AB}{s}+776+E_1(A)+E_2(B).
                                                               \tag{G038.16}
\]
The right side is an integer: (G038.13) already ensures the
cross-cell divisibilities. Extra within-cell associativity
constraints are not imposed, so this remains a necessary
relaxation for the full original-support problem.

\paragraph{The exact finite certificate.}
The independent verifier
\path{verify_support132_lower_bound.py} reads the supplied
\path{e3_constants.json}, checks every mass and parity
interface above, and executes exactly these integer loops.
There are $3020$ divisor rows, of which $272$ have nonempty
arrays on both sides, and $5059$ retained pairs of cell sums.
The minimum of (G038.16) is $5112$. Its unique minimizing
row of aggregate data is
\[
 (s,a,b,A,B,E_1,E_2)=(600,20,60,520,600,192,264).
\]
For an additional audit of that row, its minimizing action
values are $60,60,200,200$ in the first cell and
$60,60,240,240$ in the second. The generator self-square
cost is $20+4=24$; the first-cell word cost is $8+64=72$;
the second-cell word cost is $20+88=108$.
Thus its total is
\[
          600+2240+1040+776+192+264=5112.
\]
The verifier checks all rows, not just this minimizer, using
exact integers and direct enumeration of $c$ in (G038.15).
The separate exploratory finder used a closed congruence
formula; none of that finder code or data is an input to this
certificate, and no floating computation is used. The full certificate is saved as
\path{support132_lower_bound_certificate.json}.
Since $5112>5110$, every assumed lift with at most $343$
pairs is excluded. The proof establishes the lower bound
$344$; the aggregate minimizing row is not asserted to lift
to an associative Laurent table.

\paragraph{Geometric translation and exact scope of the global bound.}
The coordinate open torus from NN-G038 parametrizes all
associative tables with this support. On all of its integral
effective-character points with the prescribed trace, the
ordinary-rank functional now has excess at least $688$
over the absolute trace objects, hence at least $344$
copies of the cancelling pair $1+\epsilon$ in total.
The proof first uses the parity characters over
$\mathbb F_2[q,q^{-1}]$ and then the ordinary-dimension
specialization of the actual coefficient objects. SGA~3
I4.7.2--4.7.3 and NN-G010, with the supplied archive cote~2,
scan~17 transcription, give respectively the grading and
composing-kernel dictionary. The scalar dynamic program
is a necessary numerical image of this effective locus;
it neither identifies that image exactly nor replaces the
objects by virtual $K_0$ cancellation. The $344$ bound applies
to every point with the original support, including nonsymmetric
gradings and arbitrary parameter values. It has no asserted
extension to tables using additional support positions.

\paragraph{Full within-cell associativity strengthens the support bound to $560$.}
The preceding $344$ bound used only self-square identities within
cells. Imposing all their ordinary-dimension associativity raises
the bound to $560$ for the same entire original-support class,
still without symmetry or fixed anchors. More precisely, the
integer relaxation consisting of the original support, normalized
unit, coefficientwise ordinary-mass lower bounds, even nonunit
masses, associativity, and the necessary restriction $4\mid s$
has exact optimum $560$ pairs. Its scalar attainer below is not
asserted to admit a Laurent grading with the prescribed trace.
At this stage the original-support interval is $[560,2744]$,
strengthened below to $[836,2744]$;
the separate all-$810$-position interval remains $[104,2744]$.

\paragraph{Only one divisor remains after the four self-square coefficients.}
In one cell use the order $(a,b,c,d)=(e_{11},e_{12},e_{21},e_{22})$.
Write its self-square coefficients as $u,y,x,h$ in this order.
They are positive even integers in the scalar relaxation.
The five-parameter classification of NN-G038 forces its complete
matrix of cell coefficients to be
\[
 \begin{pmatrix}
 u&u&x&ux/p\\
 y&y&hp/u&h\\
 u&p&x&x\\
 uy/p&y&h&h
 \end{pmatrix},                                                  \tag{G038.17}
\]
where the output of the entry at $(e_{rs},e_{tv})$ is $e_{rv}$.
Here $p$ is the coefficient of $cb$ with output $d$.
Indeed (G038.4) gives $h=uz/p$, hence $z=hp/u$;
substitution gives every displayed entry.
Every admissible choice is covered by choosing the even
positive divisor $\delta=ux/p$ of $ux$, setting
$p=ux/\delta$, and checking that $p$, $hp/u$, and $uy/p$
are positive even integers. The absolute trace lower bounds
are checked for all sixteen entries. This is a finite
exhaustive parametrization, with no search cutoff on $p$
being guessed.

For a chosen action vector
$\boldsymbol\theta=(\theta_a,\theta_b,\theta_c,\theta_d)$,
the bottom coefficient corresponding to each entry $c_{ij}$
is necessarily
\[
       \frac{\theta_i\theta_j-c_{ij}
                    \theta_{\operatorname{out}(i,j)}}s.          \tag{G038.18}
\]
It too must be an even integer at least its original
absolute trace mass. Conversely the cell table (G038.17),
with (G038.18), obeys all cell and bottom associativity by
the classification. The checker imposes these tests directly,
not through a numerical equation solver.

\paragraph{Complete ranges and safe pruning under the $559$-pair budget.}
Suppose a lift has at most $559$ pairs, so $T\le5542$ and
its added mass is at most $1118$. Repeat the proved divisor
and self-square reduction with
\[
 s\in\{576,580,\ldots,1692\},\qquad
 R_s=1118-(s-576).
\]
The individual action ranges and cell-sum bounds are exactly
as before, replacing $686$ by $1118$.
For each divisor pair $a b=2s$, first compute the weak
self-square arrays $E_1(A),E_2(B)$ by direct enumeration
in (G038.15). Retain only pairs satisfying
\[
 \mathcal B(A,B)+E_1(A)+E_2(B)\le5542,
 \qquad
 \mathcal B(A,B)=s+2(A+B)+2AB/s+776.                \tag{G038.19}
\]
For each retained first-cell sum $A$, put
\[
 L_1(A)=\max_{B\ {
m retained\ with}\ A}
                \{5542-\mathcal B(A,B)-E_2(B)\},
\]
and define $L_2(B)$ with the roles exchanged.
These are upper bounds for the respective actual sums
of self-square costs in every competitive table, since
the other cell costs at least its weak minimum and the
remaining within-cell positions cost at least $776$.
Taking the maximum over partners keeps every possibility.

Now enumerate the four action values independently, with
the appropriate divisor step, target lower bounds, and fixed
sum $A$ or $B$. For each action tuple enumerate every four
self-square coefficient choices from (G038.14) whose summed
cost is at most $L_i$. Finally enumerate the divisor
$\delta$ in (G038.17) and apply all the cell and bottom
integer, parity, and target-mass tests just described.
Let $C_i(A)$ be the least full mass of the $32$ coefficients
in that cell at action sum $A$, or infinity if none passes.
Every actual table in the budget survives these stages,
so its total is bounded below by
\[
                s+2(A+B)+2AB/s+C_1(A)+C_2(B).                   \tag{G038.20}
\]
For fixed $s,A,B$ the cross-cell total depends only on the
action sums. Minimizing the two cell costs separately thus
provides a valid necessary lower bound even though the
minimizers need not arise from graded coefficient objects.

\paragraph{Independent exact certificate and a scalar attainer.}
The self-contained verifier
\path{verify_support132_full_scalar.py} again reads the
original trace table and reconstructs all its mass and parity
interfaces. It uses direct enumeration of the even self-square
coefficient $c$, then chooses $\delta$ in (G038.17).
The separate exploratory finder instead chose $p$ as a divisor
of $\gcd(ux,uy)$; the verifier reads none of its survivor files.
It regenerates all rows from the $559$-pair budget.
The exact audit counts are:
\[
 \begin{aligned}
 &5026\ \text{divisor rows},\qquad
 679\ \text{rows with both weak arrays nonempty},\\
 &34817\ \text{weak action-sum pairs},\qquad
 65\ \text{competitive divisor rows},\\
 &174\ \text{competitive action-sum pairs},\qquad
 12473\ \text{action tuples},\\
 &953\ \text{self-admissible action tuples},\qquad
 17542\ \text{self-coefficient tuples},\\
 &257562\ \text{tested divisor choices},\qquad
 727\ \text{passing full cells}.
 \end{aligned}
\]
For every retained global combination, (G038.20) is at least
$5544$, so none satisfies the assumed $5542$ budget.
The certificate is
\path{support132_full_scalar_certificate.json}; it includes
all competitive weak rows and the minimizing aggregate data.

The matching scalar table has
\[
 \begin{aligned}
 s&=648,\qquad
 \boldsymbol\theta=(72,216,216,72)\quad\text{in each cell},\\
 (u,x,y,z,p)&=(18,18,18,54,54)\quad\text{in each cell}.
 \end{aligned}                                                 \tag{G038.21}
\]
Its cell matrix is
\[
 \begin{pmatrix}
 18&18&18&6\\18&18&54&18\\18&54&18&18\\6&18&18&18
 \end{pmatrix}.
\]
Both cells have total mass $784$, both action sums equal
$576$, and the total nonunit mass is
$648+2304+1024+1568=5544$.
The verifier independently reconstructs the whole $132$-position
scalar table, checks all absolute-mass and parity bounds,
and checks every one of its $1000$ associativity triples by
direct integer multiplication. This supplies an actual attainer
for the stated ordinary-dimension relaxation, with
$(5544-4424)/2=560$ pairs. It does not supply grading
characters or assert a trace-preserving Laurent lift at that cost.

\paragraph{Geometric meaning and the smaller remaining gate.}
On every integral effective-character point of the full
nineteen-parameter original-support locus, the ordinary-rank
excess is at least $1120$, or $560$ cancelling pairs after
forgetting the grading. The proof uses the generic associative
cell multiplication together with the exact ordinary-rank
images of its coefficients; it imposes no anchor fiber and
no Laurent involution. The scalar table (G038.21) is a point
of the ordinary-dimension associativity locus satisfying these
numerical constraints. Whether it lies in the image of the
integral graded effective locus is a further lifting question.
The grading and actual composing-kernel interpretation remains
SGA~3 I4.7.2--4.7.3, NN-G010, and the supplied archive cote~2,
scan~17 transcription. Formal $K_0$ subtraction is not being
used to realize that missing graded point. This moves the
original-support lower bound from $344$ to $560$ while leaving
the all-$810$-position bound and the $2744$ graded attainer unchanged.

\paragraph{A character-factor obstruction invisible to the scalar attainer.}
The $560$-pair scalar table cannot be a graded lift of the
prescribed trace. More generally, in either cell of every
original-support Laurent lift, a generator self-square mass
and a word self-square mass cannot both be congruent to two
modulo four. Consequently, writing the four self-square masses
in the order $(u,y,x,h)$ of (G038.17), one necessarily has
\[
       (4\mid u\ \hbox{and}\ 4\mid h)
        \quad\hbox{or}\quad
       (4\mid x\ \hbox{and}\ 4\mid y).                           \tag{G038.22}
\]
Here the symbols in this divisibility statement denote ordinary
masses. The proof uses actual grading characters before passing
to their masses; it is not a property of arbitrary associative
integer scalar tables.

\paragraph{Factor-order lemma over the integral Laurent ring.}
Let $R=\mathbb Z[q,q^{-1}]$. For a Laurent polynomial $f$ with
nonzero reduction $\bar f$ modulo two, write
\[
       r(f)=\operatorname{ord}_{q=1}\bar f
       \quad\text{in }\mathbb F_2[q,q^{-1}].
\]
Suppose $f_1f_2=g_1g_2$ in $R$, each of the four reductions
is nonzero, and all four integer evaluations at $q=1$ have
two-adic valuation exactly one. Then
\[
                       \{r(f_1),r(f_2)\}
                         =\{r(g_1),r(g_2)\}                    \tag{G038.23}
\]
as multisets.

To prove this, factor each Laurent polynomial in the UFD $R$,
obtained by localizing $\mathbb Z[q]$ at the powers of $q$.
No polynomial has a factor two, since its reduction is nonzero.
The valuation of its evaluation is one, so precisely one
irreducible factor, occurring once, has even value at $q=1$;
that factor's value also has valuation one. All its other
irreducible factors have odd value there. Their reductions
therefore have order zero at $q=1$, and the unique even-valued
factor has reduction order exactly $r(f)$.
The polynomial equality makes the multisets of those
irreducible even-valued factors coincide on its two sides.
Taking their reduction orders proves (G038.23).
The units $\pm q^j$ have odd evaluation and reduction order
zero, so arbitrary Laurent shifts and unit signs cause no
exception. No positivity of these auxiliary irreducible
factors is claimed or needed.
The precise UFD inputs are the Stacks Project, Algebra,
Lemmas 10.120.7 and 10.120.10 (localization and polynomial
unique factorization), section tag 034O and polynomial-lemma
tag 0BC1: \url{https://stacks.math.columbia.edu/tag/034O}.
Their complete statements and proofs were independently read
by the root reviewer; this is a source check distinct from the
finite arithmetic verification.

\paragraph{Application to the actual two cells.}
Use uppercase letters for the actual coefficient characters
whose masses are $u,y,x,h$, and put
$\Delta_1=UX/P$, $\Delta_2=UY/P$, $Z=HP/U$ in the notation
of (G038.17). All these quotients are actual coefficient
characters of an admissible lift, hence lie in $R$.
The four product identities are
\[
 UX=\Delta_1P,\quad UY=\Delta_2P,\quad
 HX=\Delta_1Z,\quad HY=\Delta_2Z.                               \tag{G038.24}
\]
Their source parities give the following exact orders:
\[
 \begin{array}{c|ccc}
 &U,H&X,Y&\Delta_1,\Delta_2\quad;\quad P,Z\\\hline
 \text{first cell}&4&4&2\quad;\quad6\\
 \text{second cell}&4&6&2\quad;\quad8
 \end{array}
\]
Indeed $D=(q+1)^2/q$ modulo two and $D-1$ has order zero
at $q=1$. The generator squares reduce to $D^2(D-1)$;
the word squares reduce to $D^2(D-1)$ in the first cell
and $D^3(D-1)$ in the second. The cross coefficients
$\Delta_i$ reduce to $D(D-1)$, while $P,Z$ reduce to
$D^3(D-1)$ or $D^4(D-1)$ respectively.
All these reductions are nonzero, independently of the
chosen lift.

Suppose, for example, that $u$ and $x$ both equal two modulo
four. Then $U(1)X(1)$ has valuation two.
The other two actual coefficients in $UX=\Delta_1P$
have positive even masses, so each has valuation exactly one.
The factor-order lemma would give
$\{4,4\}=\{2,6\}$ in the first cell, or
$\{4,6\}=\{2,8\}$ in the second, both impossible.
The other three identities in (G038.24) give the same
contradiction for every generator/word pair. If either
generator square fails divisibility by four, both word
squares must have it; this is exactly (G038.22).
In the scalar attainer (G038.21), all four masses equal
$18$, so the obstruction excludes its graded realization.
It does not by itself classify all scalar minimizers.

\paragraph{Geometric interpretation.}
The order $r(f)$ is the multiplicity at the characteristic-two
point $q=1$ of the reduction of a nonzero character polynomial
on the grading torus. The two-adic valuation of $f(1)$ records
its specialization to the identity section. Unique factorization
shows that when this valuation is one, a unique integral
irreducible divisor carries that specialization, so product
identities preserve its reduction-order multiset. Applied to
the actual cell coefficient identities, this proves the ordinary
rank alternatives (G038.22) for every point of the integral
effective-character locus. The SGA~3 I4.7.2--4.7.3 grading
interpretation and NN-G010 composing-kernel dictionary remain
as before. The auxiliary factors are polynomial divisors and
need not be effective representation objects; no tensor-factor
or motivic decomposition is inferred from them.

\paragraph{The complete-support interval improves to $[836,2744]$.}
Insert the necessary alternative (G038.22) into the full scalar
cell enumeration just proved. Under an assumed budget of $835$
pairs, its mass cap is $6094$, its added-mass cap is $1670$,
and the bottom-square range is
\[
                         s=576,580,\ldots,2244.
\]
All weak arrays, action bounds, and maximum-partner self-cost
caps are defined exactly as in (G038.19), with $6094$ in place
of $5542$. When a four-tuple of self-square coefficients has
been selected, discard it if it fails (G038.22), before
examining its divisor choices. The uniform factor lemma proves
that this cannot discard a Laurent lift.

The independent exact verifier
\path{verify_support132_character_filtered.py} regenerates
these ranges from the budget and checks all eight source
reduction orders in each cell directly. It retains direct
self-square enumeration and the $\delta$-divisor parametrization
from the previous verifier. Its complete audit gives
\[
 \begin{aligned}
 &7724\ \text{divisor rows},\quad
 1359\ \text{rows with both weak arrays nonempty},\\
 &151120\ \text{weak action-sum pairs},\quad
 301\ \text{competitive divisor rows},\\
 &2061\ \text{competitive action-sum pairs},\quad
 339645\ \text{action tuples},\\
 &33565\ \text{self-admissible action tuples},\quad
 3617401\ \text{self-coefficient tuples},\\
 &1193063\ \text{tuples rejected by (G038.22)},\quad
 45488060\ \text{divisor choices},\\
 &55898\ \text{passing full cells}.
 \end{aligned}
\]
Every retained global combination has mass at least $6096$.
The certificate is
\path{support132_character_filtered_certificate.json}.
Since $6096>6094$, every original-support lift has at least
$836$ pairs. No Laurent degree bound is involved.

For precision about the remaining numerical loss, this
character-filtered scalar relaxation has a matching scalar
attainer. Its bottom square is $624$, its two action vectors
are $(96,216,216,96)$ and $(104,208,208,104)$, and its cell
parameters $(u,x,y,z,p)$ are respectively
\[
                  (44,8,8,44,44),\qquad(8,16,16,32,32).
                                                               \tag{G038.25}
\]
Both action sums are $624$, the cell masses are $936$ and $792$,
and the nonunit total is
\[
                624+2496+1248+936+792=6096.
\]
The verifier independently reconstructs this $132$-position
integer table, checks every target-mass and parity bound,
checks (G038.22), and directly verifies all $1000$ associativity
triples. Thus $836$ is the exact optimum of this stated
filtered scalar relaxation. A trace-preserving graded lift at
that value remains unconstructed and is not implied by the
scalar attainer.

\paragraph{Geometric conclusion and remaining gate.}
The geometric conclusion is now that every integral
effective-character point in the full nineteen-parameter
original-support locus has ordinary-rank excess at least
$1672$, or $836$ cancelling pairs. The improved obstruction
comes from characteristic-two divisor multiplicities before
ordinary dimension is taken. The preceding $344$ and $560$
bounds are historical stages. At this stage the original-support interval is $[836,2744]$,
strengthened below to $[900,2744]$, while the all-$810$-position
minimum remains in $[104,2744]$ and the smaller symmetric
closed presentation retains its exact unique $2744$ optimum.

\paragraph{Saturated coefficients give finite factor-allocation tests.}

This refinement concerns the complete original 132-position support class
of NN-G038, with arbitrary Laurent grading and no inversion symmetry.
It does not concern the 810-position problem.  The numerical conclusion below has an independently checked exact certificate.

Put $R=\mathbb Z[q,q^{-1}]$.  If $f\in R$ has nonzero reduction
$\bar f\in\mathbb F_2[q,q^{-1}]$, let
$r(f)=\operatorname{ord}_{q=1}\bar f$.  Laurent monomials are units
and have order zero.  Write $m(f)=|f(1)|$ whenever $f(1)\ne0$.
For such an $f$,
\[
                 m(f)\text{ is odd}\quad\Longleftrightarrow\quad r(f)=0.
\tag{G038.SF1}
\]
Indeed reduction followed by evaluation at $q=1$ is reduction modulo two
of the integer $f(1)$.  This elementary equivalence does not require
coefficientwise positivity of $f$.

\paragraph{Finite divisor-allocation lemma.}
Suppose $F,U,\Delta,P\in R$ have nonzero reductions, nonzero evaluations
at one, and $FU=\Delta P$.  Fix the entire polynomial $F$, and let
\[
 \mathcal D(F)=
 \{(m(G),r(G)):G\text{ divides }F\text{ in }R\}.
\]
This is a finite set because the units $\pm q^j$ do not change either
statistic.  There must exist $(g,r_G)\in\mathcal D(F)$ for which the
positive numbers
\[
 h=\frac{m(\Delta)}g,\qquad j=\frac{m(U)}h
\]
are integers and, on putting
\[
 r_H=r(\Delta)-r_G,\qquad r_J=r(U)-r_H,
\]
all the following necessary conditions hold:
\[
 \begin{aligned}
 &r_H,r_J\ge0,\qquad \frac{m(F)}g j=m(P),\\
 &r(F)-r_G+r_J=r(P),\\
 &(h\text{ odd})\Longleftrightarrow(r_H=0),\qquad
   (j\text{ odd})\Longleftrightarrow(r_J=0).
 \end{aligned}
\tag{G038.SF2}
\]
The same test can be applied with any of the four factors held fixed,
by interchanging the two sides or the two factors on a side.

\begin{proof}
The ring $R$ is a UFD: Gauss's lemma gives this for $\mathbb Z[q]$,
and localization at powers of $q$ preserves unique factorization.
Choose a gcd $G$ of $F$ and $\Delta$, and write
$F=GF_0$, $\Delta=GH$ with $\gcd(F_0,H)=1$.  The identity
$F_0U=HP$ implies $H\mid U$, so $U=HJ$ and $P=F_0J$.
All factors have nonzero reductions and evaluations, since their
products do.  Absolute integer evaluation gives the displayed
formulas for $g,h,j$; orders add in the domain
$\mathbb F_2[q,q^{-1}]$.  Applying (G038.SF1) to $H$ and $J$
proves every condition in (G038.SF2).  Signs and arbitrary Laurent
unit factors do not alter this argument.
\end{proof}

\paragraph{Why a coefficient at its minimum is fixed.}
For any coefficient $d$ of an effective $C_2$ lift of prescribed
trace $t$, coefficientwise $d\ge |t|$.  If $d(1)=|t|(1)$, the
nonnegative Laurent polynomial $d-|t|$ has coefficient sum zero,
so $d=|t|$.  This turns an ordinary-dimension equality into an
exact fixed polynomial, without bounding the support of any
unsaturated coefficient.

Let $D=q+2+q^{-1}$, $C=D-1$, $K=D-2$, and
$\psi=D(D-3)$.  The possible fixed coefficients in the four
identities
\[
 UX=\Delta_1P,\quad UY=\Delta_2P,\quad
 HX=\Delta_1Z,\quad HY=\Delta_2Z
\tag{G038.SF3}
\]
are exactly the following source absolute-trace polynomials at their
respective minimum masses:
\[
\begin{array}{c|c|c|c}
\text{coefficient role}&F&m(F)&r(F)\\ \hline
\text{generator square, or cell-one word square}&\psi K&8&4\\
\text{cross coefficient }\Delta&\psi&4&2\\
\text{cell-two word square}&\psi K^2&16&6\\
\text{cell-one }P,Z&\psi(D^2-2D-2)&24&6\\
\text{cell-two }P,Z&C(q^2+q^{-2})^2&12&8
\end{array}
\tag{G038.SF4}
\]
Their complete factorizations, ignoring Laurent units, are
\[
\begin{aligned}
 \psi K&\sim\Phi_2^2\Phi_6\Phi_4,&
 \psi&\sim\Phi_2^2\Phi_6,\\
 \psi K^2&\sim\Phi_2^2\Phi_6\Phi_4^2,&
 \psi(D^2-2D-2)&\sim\Phi_2^2\Phi_6 B,\\
 C(q^2+q^{-2})^2&\sim\Phi_3\Phi_8^2,&
 B&=q^4+2q^3+2q+1.
\end{aligned}
\tag{G038.SF5}
\]
Here the notation is explicit:
$\Phi_2=q+1$, $\Phi_3=q^2+q+1$,
$\Phi_4=q^2+1$, $\Phi_6=q^2-q+1$, and $\Phi_8=q^4+1$.
The three quadratics are irreducible by their negative discriminants;
$\Phi_8(q+1)$ is Eisenstein at two.
The monic quartic $B$ has no rational linear factor, since neither
$1$ nor $-1$ is a root.  In a monic integral quadratic factorization,
the constant terms are both one or both minus one.
In the first case the factors are $q^2+a q+1$ and $q^2+b q+1$,
requiring $a+b=2$, $ab=-2$, which has no integral solution.
In the second case the coefficients of $q^3$ and $q$ are opposites,
contrary to $B$.  This proves irreducibility of $B$.
The relevant factor statistics are
\[
\begin{array}{c|rrrrrr}
 F&\Phi_2&\Phi_3&\Phi_4&\Phi_6&\Phi_8&B\\ \hline
 m(F)&2&3&2&1&2&6\\
 r(F)&1&0&2&0&4&4
\end{array}
\tag{G038.SF6}
\]
Therefore every set $\mathcal D(F)$ in (G038.SF4) is obtained by
choosing factor exponents between zero and the multiplicities in
(G038.SF5), multiplying the first statistics and adding the second.
The value-one, order-zero factor $\Phi_6$ may be omitted from the
statistics enumeration, even though it remains a genuine factor.

\paragraph{The displayed 836-pair scalar attainer fails this test.}
In its first cell, $m(U)=44$, $m(X)=8$, $m(\Delta_1)=8$,
and $m(P)=44$.  The factor $X$ is saturated and hence equals
$\psi K$.  Apply the lemma to $XU=\Delta_1P$.
The root orders are respectively $4,4,2,6$.
Write $a=v_2(g)$ for a divisor statistic of $X$.
Then $v_2(h)=3-a$ and $v_2(j)=a-1$.
If $a=1$, then $j$ is odd, so $r_J=0$, whereas
$r_J=4-(2-r_G)=2+r_G>0$.
If $a\ge2$, the restriction $r_G\le2$ forces
$G$ to have the same statistics as $\Phi_2^2$, namely
$(g,r_G)=(4,2)$: using $\Phi_4$ with another even-valued
factor gives order at least three.  This gives $r_H=0$ but
$h=2$ even, again impossible.  The case $a=0$ would make
$j$ nonintegral.  Thus this particular scalar attainer has no
Laurent lift.  This assertion does not exclude all scalar
attainers at that mass, and by itself gives no improved global
numerical bound.

\paragraph{Geometric translation and limitation.}
The saturated coefficient is a fixed effective character of the
grading torus.  Its integral character polynomial has a finite
prime-divisor decomposition; allocation of those divisors across
the multiplication identity must respect both specialization at
the identity and the order of the characteristic-two divisor there.
These are necessary conditions on all integral character points
of the NN-G038 associativity scheme, with no degree cutoff.
The irreducible polynomial factors used in the proof are not
asserted to be effective characters, subrepresentations, tensor
factors, or geometric summands.  Their use is solely in the UFD
of integral Laurent polynomials.  The geometric conclusion is an
actual exclusion of the displayed ordinary-dimension point from
the integral character-point image, not an exclusion of a virtual
$K_0$ class or a statement about the full 810-position scheme.

\paragraph{The complete original-support lower bound rises to $900$.}
Add the necessary tests (G038.SF2) to the previous full scalar
enumeration.  For each cell, test each of the four product identities
(G038.SF3), and within each identity test every factor whose mass equals
its minimum in (G038.SF4).  A failed test excludes the scalar cell.
The finite divisor-allocation lemma proves that no Laurent lift can
be removed.  The tests impose no support bound, inversion symmetry,
or equality between the four action values.

To exclude $899$ pairs, the nonunit mass cap is $6222$ and the
added-mass cap is $1798$.  The bottom square therefore ranges over
\[
                         s=576,580,\ldots,2372. \tag{G038.SF7}
\]
For each $s$ and divisor split $ab=2s$, retain the weak action arrays
and safe maximum-partner self-cost caps of (G038.19), with this new
mass cap.  Enumerate the four independent action values, all four
self-square coefficients, and the positive even divisors
$\Delta_1\mid ux$ exactly as in the independently checked full-cell
verifier.  Keep the earlier necessary alternative (G038.22).
The new tests are applied to the reconstructed coefficients before
computing the full cell cost.  Every final full cell is also checked
at each of its sixteen within-cell coefficients and sixteen bottom
coefficients.  The proof of completeness of the ranges, partner
pruning, and divisor parametrization is unchanged: the new test only
removes cells by a proved necessary condition.

The self-contained verifier
\path{verify_support132_saturated_filtered.py} reads the original
trace table directly, regenerates all weak ranges, and never reads
the exploratory search output.  It also independently verifies all
sixteen source absolute-polynomial identities, all five displayed
Laurent factorizations, and all source reduction orders.  Divisor
statistics are computed from the displayed polynomial factors,
rather than imported from a precomputed survivor list.
Its complete exact audit is
\[
\begin{aligned}
 &8323\ \text{divisor rows},\quad
 1527\ \text{rows with both weak arrays nonempty},\\
 &194563\ \text{weak action-sum pairs},\quad
 357\ \text{competitive divisor rows},\\
 &2966\ \text{competitive action-sum pairs},\quad
 577588\ \text{action tuples},\\
 &60988\ \text{self-admissible action tuples},\quad
 8974103\ \text{self-coefficient tuples},\\
 &3077158\ \text{tuples rejected by (G038.22)},\\
 &113396168\ \text{divisor choices},\quad
 2653320\ \text{choices rejected by (G038.SF2)},\\
 &108653\ \text{passing full cells}.
\end{aligned}
\tag{G038.SF8}
\]
Every retained global combination has nonunit mass at least $6224$,
strictly above the $6222$ assumed cap.  Hence every original-support
integral Laurent lift has at least $900$ cancelling pairs.
The full certificate is
\path{support132_saturated_filtered_certificate.json}.

This strengthened scalar relaxation again has a matching scalar
attainer.  It has
\[
\begin{aligned}
 s&=616,\\
 \boldsymbol\theta_1&=(88,220,220,88),&
 (u,x,y,z,p)_1&=(32,24,24,32,32),\\
 \boldsymbol\theta_2&=(56,224,224,112),&
 (u,x,y,z,p)_2&=(12,26,26,104,52).
\end{aligned}
\tag{G038.SF9}
\]
The two cell costs are $952$ and $960$; both action sums are $616$.
Thus the total nonunit mass is
\[
                      616+2464+1232+952+960=6224.
\]
The verifier reconstructs all $132$ scalar coefficients, checks
every ordinary target-mass and parity condition, both kinds of
character-derived scalar filter, and all $1000$ associativity
triples by direct multiplication.  This proves that $900$ is the
exact optimum of the stated filtered ordinary-dimension relaxation.
It does not provide a graded lift at $900$ or a new upper bound.

\paragraph{Geometric conclusion and the remaining image problem.}
The ordinary-rank excess on every integral effective-character point
of the nineteen-parameter original-support locus is now at least
$1800$, or $900$ cancelling pairs.  This conclusion holds for arbitrary
Laurent gradings.  The saturated strata have fixed integral divisors
on the grading torus, and their finite factor allocations must be
compatible with specialization at the identity and the multiplicity
of the characteristic-two divisor there.  Applying this constraint
before taking dimensions strengthens the numerical image bound.
The scalar table (G038.SF9) belongs to the ordinary associative locus
and passes the stated necessary filters, but whether it lies in the
image of the effective integral character locus remains open.
Consequently the original-support interval becomes $[900,2744]$;
the all-$810$-position interval remains $[104,2744]$ and the symmetric
closed presentation retains its exact unique $2744$ optimum.
The grading dictionary is SGA~3 I4.7.2--4.7.3; the actual
composing-kernel interpretation is NN-G010 and the supplied archive
cote~2, scan~17 transcription.  The UFD input is precisely the Stacks
Project, Algebra, Lemmas 10.120.7 and 10.120.10,
\url{https://stacks.math.columbia.edu/tag/034O}; it supplies no
effectivity assertion about the auxiliary polynomial factors.



\section{NN-S001: fixed fibres, rather than an automatic Boolean fixed part}
\label{sec:NN-S001}

\paragraph{Inputs.} This is a general form of corpus 01, Section 3
(semimodules as diagrams of modules), corpus 02's corresponding repair,
and corpus 16, Theorem \texttt{thm:fixed-locus-unit}. No E, A, or M is used.

\begin{theorem}[NN-S001]
Let $R$ be a nonzero commutative unital ring, $S=G(R)$, and $M$ an
$S$-semimodule. Write $e=0_R$, $B=eM$, and $M_b=\{m:em=b\}$.
Then $B$ is a join semilattice, $M$ is the disjoint union of the
$R$-modules $M_b$, whose respective zero elements are $b$, and
\[
 \operatorname{Fix}_M(u)=\coprod_{b\in B}\ker(u-1:M_b\longrightarrow M_b)
 \qquad(u\in R^\times).
\]
In particular $\operatorname{Fix}_M(-1)=\coprod_b M_b[2]$.
This is exactly the skeleton $B$ if and only if every fibre has no
nonzero $2$-torsion. The assertion need not hold for an arbitrary
$G(R)$-semimodule even when $R=\mathbb Z$.
\end{theorem}

\begin{proof}
The scalar identities $e+e=e$, $1+e=1$, and $1+(-1)=e$ imply
$b+b=b$, $m+em=m$, and $m+(-1)m=em$. Thus $B$ is an idempotent
commutative monoid, with join its addition. In $M_b$, addition is closed,
$b$ is an identity, and $(-1)m$ is the inverse of $m$. Multiplication
by a supported scalar $r\in R$ preserves $M_b$ because $er=e$.
The semimodule axioms now give precisely the $R$-module axioms in
each fibre, including the action of $0_R$ as its zero element $b$.
The displayed formula is the equation $um=m$ inside this module.
For a counterexample let $M=\mathbb Z/2\mathbb Z$ acted on through
$G(\mathbb Z)\to\mathbb Z$. It has the single skeleton element $0$,
but both module elements are sign-fixed. In contrast every fibre of
$G(R)^d$ is a coordinate module $R^I$; its sign-fixed part is its
Boolean skeleton when $2$ is a non-zero-divisor in $R$.
\end{proof}

\paragraph{Consequence and remaining gate.}
Replacing $\mathbb Q_\ell$ by arbitrary $G(R)$-semimodules in the
sign-only argument does not automatically make the entire fixed part
Boolean. One must specify torsion-free linear fibres and then verify
the trace and effective-spectrum hypotheses on each chosen realization.
This corrects the unrestricted coefficient assertion in prompt 10(a);
it yields no parity or satisfiability algorithm.

\paragraph{Geometric dictionary for NN-S001.}
The skeleton indexes a diagram of ordinary modules, hence of
quasi-coherent sheaves when the modules are localized over $\operatorname{Spec}R$.
The fixed fibre is the kernel of the endomorphism $u-1$ of that sheaf.
This is an exact module calculation, not a claim that the skeleton is
an additional vector-space summand. The ordinary affine-module
dictionary is EGA I, Section 1.3; exact edition locators are recorded
in the source register. The nonadditive skeleton lies outside that
ordinary quasi-coherent category.



\section{NN-S002: the sign-fixed scheme has a characteristic-two fibre}
\label{sec:NN-S002}

\begin{theorem}[NN-S002]
The fixed-point scheme of sign on $\mathbb A^d_{\mathbb Z}$ is
\[
 F_d=\operatorname{Spec}\mathbb Z[x_1,\ldots,x_d]/(2x_1,\ldots,2x_d).
\]
For $d>0$ it is not flat over $\mathbb Z$. Its characteristic-zero
fibre is the origin, whereas its characteristic-two fibre is the
whole affine space. Consequently the reduction of
$\operatorname{Fix}_{G(\mathbb Z)}(-1)=\{\tau,e\}$ is a proper subset
of $\operatorname{Fix}_{G(\mathbb F_2)}(-1)=G(\mathbb F_2)$.
\end{theorem}

\begin{proof}
For every commutative ring $T$, sign fixes a vector $(a_i)$ exactly
when $2a_i=0$ for all $i$. The displayed coordinate algebra represents
this functor. Tensoring it with $\mathbb Q$ or $\mathbb F_2$ gives
the two claimed fibres. The nonzero element $x_1$ in the integral
coordinate algebra is killed by $2$; a flat $\mathbb Z$-module has
no such torsion, since tensoring $0\to\mathbb Z\xrightarrow{2}\mathbb Z$
must preserve injectivity. Over $\mathbb Z$ the scalar fixed equation
has only amplitude zero, while over $\mathbb F_2$ sign is identity.
The external element $\tau$ is fixed in both cases, proving the
last assertion.
\end{proof}

\paragraph{Geometric dictionary and scope.}
This theorem is already a scheme-theoretic statement: it is the
equalizer of identity and the order-two automorphism of the vector
group. The equalizer itself \emph{does} commute with base change,
as its displayed equations show. What fails is replacing it by the
zero section using its characteristic-zero points and then reducing.
The coordinate calculation provides the full proof. The rank-one
example is already recorded in SGA 3, Expos\'e I, Remark 4.7.1.2,
printed page 29 of the corrected edition dated 13 October 2024
(the editorial addition is identified by its note 39). The new
application here is to the split-semimodule claim; the classical
fixed-scheme phenomenon is not claimed as new. Equalizers and base
change use SGA 3, I, Sections 1.2 and 1.6.
For an $R$-module the analogous issue is the kernel of multiplication
by $2$; its change after tensoring is measured by the associated
$\operatorname{Tor}_1$ sequence. No assertion about a finite-field
Frobenius trace follows merely from the sign-fixed Boolean pair.



\section{NN-S003: additive linearization erases the support skeleton}
\label{sec:NN-S003}

\begin{theorem}[NN-S003]
Every additive map from an idempotent commutative monoid to an
abelian group is zero. Thus the universal ring reflection of $G(R)$
is $R$, and extension of scalars along $G(R)\to R$ kills every
semimodule on which $e$ acts as identity. In particular
\[
 R\otimes_{G(R)}\mathbb B=0.
\]
Even on free rank-one split modules, Boolean Frobenius data do not
determine a Weil weight: for each integer $j$, multiplication by
$q^j$ on $G(\mathbb Q_\ell)$ has the same Boolean action, but its
amplitude realization has eigenvalue $q^j$ and weight $2j$.
\end{theorem}

\begin{proof}
If $x+x=x$, then $f(x)+f(x)=f(x)$ in the group, so $f(x)=0$.
Any map from $G(R)$ to a ring therefore identifies $e$ with the
additive identity, which also receives $\tau$. The remaining
supported scalars have exactly the ring operations of $R$, proving
the universal property. In the tensor product, for $em=m$ one has
$r\otimes m=r\otimes em=re\otimes m=0$, since $e$ maps to zero
in $R$. Finally $q^j$ is a nonzero scalar, so its Boolean support
action is identity for every $j$. On the ordinary one-dimensional
realization every complex embedding of its eigenvalue has absolute
value $q^j$. This gives distinct weights with identical skeletons.
\end{proof}

\paragraph{Purity consequence.}
If a finite Boolean set has a Frobenius permutation, its separate
free linearization on the atoms is a permutation representation;
its eigenvalues are roots of unity and it has weight zero. This
linearization is not additive for union: the vector attached to
$A\cup A$ equals that attached to $A$, not twice that vector.
It therefore does not contradict the theorem. Nor is scalar
group completion the same construction as the split $K_0$ of a
category of finite Boolean objects, whose direct-sum ranks can survive.

The correct conditional transfer is fibrewise: supply a finite
effective Frobenius spectrum, the path-count identities and bounds
of GCT\_NSRH F03 (or F04/F17), at every required embedding, and the
dual spectrum. That cited criterion then gives its stated weight
bound and, with duality, purity. The existence of $G(R)$, its fixed
pair, and A provide none of those missing realizations. The same
support is compatible with every Tate weight, so weight zero on
the skeleton cannot force PH1 or positivity of a canonical formula.

\paragraph{Geometric dictionary for NN-S003.}
The amplitude rank-one examples are Tate Weil objects; their
Frobenius eigenvalues give the standard weights of Weil II,
Definition 1.2.2. The distinction between additive group completion
and the Grothendieck group of an exact category is essential in
SGA 6, Expos\'e IV (the definition of $K_0$ and its functorial
operations). No realization functor to that exact category is
constructed by a semiring support map.



\section{NN-S004: arbitrarily large support geometry with the same zeta}
\label{sec:NN-S004}

\paragraph{Background.} The construction and ideal classification
of $G_L(R)$ are corpus 16, Theorems \texttt{thm:lattice-normal-form},
\path{thm:lattice-ideal-classification}, and
\path{thm:lattice-prime-classification}. They are inputs to the
following zeta consequence, not new constructions in this run.

\begin{theorem}[NN-S004]
Let $L$ be a nontrivial bounded distributive lattice and $R$ an
infinite commutative unital ring. In $G_L(R)$ give an ideal its Bourne quotient norm.
Every finite-norm ideal is an amplitude ideal
$I_J=Z_L\cup\{(r,1_L):r\in J\}$, and its quotient is $R/J$.
Conversely each finite quotient $R/J$ gives this ideal and norm.
Moreover $I_JI_K=I_{JK}$. In particular for every number field $K$
and every such $L$,
\[
 \zeta^{\rm cl}_{G_L(\mathcal O_K)}(s)=\zeta_K(s).
\]
For finite $L$ of cardinality $\ell$, the separated quotient has
cardinality $N(J)+\ell-1$ and the separated sheet is
\[
 \sum_{0\ne J\subseteq\mathcal O_K}(N(J)+\ell-1)^{-s}.
\]
Thus the classical sheet cannot reconstruct the cardinality, height,
or branching of the support lattice, even when its RH is supplied by A.
\end{theorem}

\begin{proof}
A pure-support ideal lies in the kernel of the amplitude map to
$R$. That map therefore descends to its Bourne quotient and is
still surjective. Its quotient norm is infinite. For $I_J$, the
amplitude map modulo $J$ gives a surjection onto $R/J$.
All zero-fibre elements become equivalent to the external zero:
they lie in $I_J$, hence are Bourne-equivalent to zero by adding
the element itself on the other side. Supported scalars differing
by $J$ are equivalent by the ordinary additive equations, and the
map to $R/J$ shows that no other pairs are equivalent. This proves
the quotient claim and exhausts finite-norm ideals by the cited
classification. Products of amplitude ideals have amplitude $JK$;
they contain every $z_\lambda$ because $e_Lz_\lambda=z_\lambda$
with both factors in the two ideals. Conversely every product
amplitude is in $JK$, proving equality. For $R=\mathcal O_K$ the
finite-index ring ideals are exactly the nonzero ideals, with
the usual multiplicative norms. The classical Dirichlet series
is therefore identical term by term to $\zeta_K$.

In a separated quotient $G_L(R/J)$ the zero-amplitude fibre has
$\ell$ elements and each of the other $|R/J|-1$ amplitudes has
one, proving its cardinality. Finally choose chains of arbitrarily
large finite length, or Boolean lattices of arbitrarily large rank.
These have different heights or branching but give the identical
classical Dirichlet series. For the chain with $n+1$ elements,
corpus 16, Corollary \path{cor:lattice-split-dimension}, gives
$\dim G_L(\mathbb Z)=n+1$, which is unbounded.
\end{proof}

\paragraph{Geometric dictionary and complexity impact.}
The arithmetic scheme seen by these finite Bourne quotients is
exactly $\operatorname{Spec}\mathcal O_K$; ring reflection kills
the full idempotent lattice. The extra lattice-prime stratum in
corpus 16's spectral ordinal sum is a semiring spectrum, not an
extra stratum of the ordinary EGA scheme. For abelian $K$, A and
abelian Artin factorization give RH for the common classical
sheet. This equality of analytic invariants imposes no reconstruction
of the arbitrarily chosen support lattice. This is an explicit
failure of a proposed inference from the invariant, not a logical
independence proof about P versus NP. It supplies no efficient
algorithm for an implicitly presented support lattice.


\section{NN-S005: special values of every shifted Dedekind sheet}
\label{sec:NN-S005}

\paragraph{Inputs.} Corpus 01, Sections 5--8, and corpus 02,
Section 14, give the shift mechanism. The residue calculation below
extends the displayed Hurwitz special-value formula to every number
field and every real shift $t\geq0$. It uses the meromorphic
continuation of the Dedekind zeta function, with its sole simple
pole at one. It uses neither A nor M.

\begin{theorem}[NN-S005]
For a number field $K$, set $\kappa_K=\operatorname{Res}_{s=1}\zeta_K(s)$
and
\[
 F_{K,t}(s)=\sum_{0\ne J\subseteq\mathcal O_K}(N(J)+t)^{-s},
 \qquad t\geq0,\quad \Re s>1.
\]
This function continues meromorphically to $\mathbb C$, with only a
simple pole at one and residue $\kappa_K$. For every integer $m\geq0$,
\[
 \boxed{F_{K,t}(-m)=
 \sum_{k=0}^m\binom mk t^k\zeta_K(-m+k)
       -\frac{\kappa_K}{m+1}t^{m+1}.}                 \tag{S005.1}
\]
In particular $F_{K,t}(0)=\zeta_K(0)-\kappa_Kt$.
The last term is essential: substituting $s=-m$ before resolving
the zero times pole in the shift expansion would omit it.
\end{theorem}

\begin{proof}
Write $\zeta_K(s)=\sum_{n\geq1}a_n n^{-s}$ and choose an integer
$M>t$. Separate the finitely many terms with $n\leq M$. The binomial
series, absolutely convergent for $t/n<1$, gives
\[
 F_{K,t}(s)=\sum_{n\leq M}a_n(n+t)^{-s}
  +\sum_{k\geq0}\frac{(-t)^k(s)_k}{k!}
       \left(\zeta_K(s+k)-\sum_{n\leq M}a_n n^{-s-k}\right).       \tag{S005.2}
\]
First this equality holds in $\Re s>1$ by absolute convergence.
For a compact set of $s$-values, $(s)_k/k!$ is bounded by a fixed
power of $k$ times a constant. Choose a fixed $c$ large enough that
$\sum a_n n^{-\Re s-c}$ converges uniformly on that compact set.
For $k\geq c$ the parenthesis in (S005.2), expressed by its
convergent tail series, is bounded by a constant times
$(M+1)^{-k+c}$. Thus the $k$-series converges normally away from
its initially displayed poles. The term indexed by $k\geq1$ has
an apparent pole at $s=1-k$, which is cancelled by its simple
factor $s+k-1$. Only the $k=0$ pole at one remains, with the stated
residue. This also justifies evaluation of the continued sum.

At $s=-m$, terms with $k\leq m$ have coefficient
$(-t)^k(-m)_k/k!=\binom mk t^k$. The term $k=m+1$ has limit
\[
 \frac{(-t)^{m+1}}{(m+1)!}
 \left(\prod_{j=0}^{m-1}(-m+j)\right)\kappa_K
       =-\frac{\kappa_K t^{m+1}}{m+1}.
\]
Every term $k>m+1$ vanishes because the rising factorial is zero
and the zeta factor is regular. The finite subtractions in
(S005.2) cancel its initial finite sum by
$(n+t)^m=\sum_{k=0}^m\binom mk t^k n^{m-k}$.
This proves (S005.1), including $m=0$ with the empty product equal
to one.
\end{proof}

\paragraph{Applications and precise limits.}
NN-S004 identifies the separated sheet for a support lattice of
size $\ell$ with $t=\ell-1$, so (S005.1) applies for arbitrary
finite support size. For $K=\mathbb Q$ and $t=1$ it agrees with
$\zeta(s)-1$; this particular identity does not extend to a
general $K$. For example $\zeta_{\mathbb Q(i)}(s)=\zeta(s)\beta(s)$,
$\beta(0)=1/2$, and $\kappa_{\mathbb Q(i)}=\pi/4$ give
$F_{\mathbb Q(i),1}(0)=-(1+\pi)/4$. These values do not assert
the location of the zeros of the shifted sheet. Neither A nor M
provides its missing functional equation.

\paragraph{Geometric dictionary for NN-S005.}
Each finite amplitude quotient is the arithmetic closed
subscheme $\operatorname{Spec}(\mathcal O_K/J)$; the lattice
changes its chosen quotient-cardinality statistic from $N(J)$ to
$N(J)+\ell-1$. The theorem is an analytic statement about that
statistic. It does not identify it with an Euler product or with
a trace of a finite Weil object. In the ordinary scheme dictionary
of EGA I, Section 4 (closed subschemes), $J$ remains the same
ideal sheaf; adjoining support classes is extra semiring data.
The residue term is analytic boundary data, not an additional
effective summand in $K_0$.



\section{NN-S006: all local phase atoms forced by reflection}
\label{sec:NN-S006}

\paragraph{Inputs and attribution.}
We use A to put the nontrivial zeros on the critical line and
corpus 10, Theorem \texttt{thm:local-phase-packet} and Proposition
\path{prop:packet-symmetry-local}, for the reduced logarithmic
germ. The phase formalism retains the source's attribution:
J.~Tobley, \emph{Unpublished correspondence on orthogonal-walk
phase decompositions, Taylor-series atomization, and affine
cyclotomic charts}, ``project source transcripts communicated with
the manuscript materials, May 2026.'' The following
all-orders restriction is the consequence of the supplied theorem
A; it assumes no simplicity.

\begin{theorem}[NN-S006]
Let $H$ be holomorphic near a point $\rho$ on $\Re s=1/2$, and
suppose $H(s)=\overline{H(1-\bar s)}$ there. If its zero at $\rho$
has order $m$, write
\[
 H(\rho+z)=\lambda z^m\exp h(z),\qquad
 h(0)=0,\quad h(z)=\sum_{n\geq1}c_nz^n.
\]
Then
\[
 \lambda=(-1)^m\bar\lambda,\qquad
 c_n=(-1)^n\bar c_n.                                  \tag{S006.1}
\]
In particular all tangent phase atoms
$\operatorname{Im}(c_ni^n)$ vanish. The normal logarithmic modulus
of the reduced germ is even in the real displacement, and its
normal phase, using the logarithm with value zero at the origin,
is odd.
This applies to every zero of $\xi$ under A. For a primitive
Dirichlet character with root number $\varepsilon$, it applies
to $H_\chi=\varepsilon^{-1/2}\Lambda(s,\chi)$ at every nontrivial
zero under A. For a Dedekind zeta completion it applies under A
when $K/\mathbb Q$ is abelian, and under M, conditionally, for
general $K$.
\end{theorem}

\begin{proof}
Since $1-\bar\rho=\rho$, reflection says
$H(\rho+z)=\overline{H(\rho-\bar z)}$. Comparing the leading
terms gives the first identity. Dividing by the leading terms
gives $\psi(z)=\overline{\psi(-\bar z)}$ for the reduced germ
$\psi=H(\rho+z)/(\lambda z^m)$ with $\psi(0)=1$.
The two holomorphic logarithms with value zero at the origin
are therefore equal: $h(z)=\overline{h(-\bar z)}$.
Coefficient comparison proves the second identity.
Thus even coefficients are real and odd coefficients purely
imaginary. Multiplication by $i^n$ makes either type real, giving
the tangent assertion; evaluating at real $z$ gives the other two.

For $\xi$ the reflection follows from its functional equation
and real coefficients. For primitive $\chi$ the two standard
identities are
$\Lambda(s,\chi)=\varepsilon\Lambda(1-s,\bar\chi)$ and
$\overline{\Lambda(\bar s,\chi)}=\Lambda(s,\bar\chi)$.
Choose a unit square root $\alpha=\varepsilon^{1/2}$.
Since $\bar\alpha=\alpha^{-1}$, the normalized function
$\alpha^{-1}\Lambda(s,\chi)$ has the stated reflection.
The Dedekind completion is self-dual with real coefficients.
The appropriate RH input is used only to put every zero in the
fixed line; the local argument itself applies to any zero already
known to lie there. It does not assert that a zero is simple or
that an individual deformation velocity is nonzero.
\end{proof}

\paragraph{Geometric and GCT dictionary for NN-S006.}
At each finite jet order $r$, the completed local coordinate
algebra is $\mathbb C[z]/(z^{r+1})$ with the semilinear involution
$z\mapsto-z$, conjugating scalars. With $w=-iz$ its invariant
real form is $\mathbb R[w]/(w^{r+1})$, and scalar extension to
$\mathbb C$ recovers the original algebra. Identity (S006.1)
is exactly descent of the reduced logarithmic section to this
real Artin scheme. The proof is coefficientwise descent, so no
unproved effectivity assertion is needed. The infinitesimal
neighbourhood convention is EGA I, Section 4, and finite locally
free descent is the usual affine faithfully flat descent
dictionary. This is a real form and a reflection symmetry;
it supplies no finite-field Frobenius operator and no positive
canonical basis for a nonstandard quantum group.



\section{NN-S007: the literal zero divisor and tail are not finite Weil objects}
\label{sec:NN-S007}

\begin{theorem}[NN-S007]
The literal infinite zero-support and tail constructions in
corpus 05--10 cannot themselves be finite constructible
algebraic Weil objects by the following two direct routes.
First, the set of nontrivial complex zeta zeros is not a
Zariski-constructible subset of $\mathbb A^1_{\mathbb C}$.
Second, if $V_n$ are nonzero vector spaces over a characteristic-zero
field for infinitely many $n$, then
\[
 \prod_nV_n\big/\bigoplus_nV_n
\]
has infinite dimension. A puts the zero support on a line but
removes neither obstruction.
\end{theorem}

\begin{proof}
A closed proper algebraic subset of the affine line over
$\mathbb C$ is finite, since a nonzero one-variable polynomial
has finitely many roots. A constructible subset is a finite union
of locally closed subsets, hence is finite or cofinite. The
nontrivial zero set is infinite, while its complement contains
every real $s>1$ and is infinite; it is therefore not constructible.

Choose indices $i_1<i_2<\cdots$ with $V_{i_n}\ne0$ and
choose $v_{i_n}\ne0$. For distinct positive integers
$a_1,a_2,\ldots$, put $a_j^nv_{i_n}$ in factor $i_n$ and zero
in every other factor, and consider the resulting classes. If a finite linear
combination is zero in the quotient, then
$\sum_{j=1}^rc_ja_j^n=0$ for all sufficiently large $n$.
Taking $r$ consecutive such $n$ gives an invertible Vandermonde
matrix, multiplied by a nonsingular diagonal matrix. Hence all
$c_j=0$. Infinitely many independent classes result. The same
argument applies if the original index set has further factors.
\end{proof}

\paragraph{Exact scope of the obstruction.}
Corpus 06, Theorem 6.8, and corpus 07's vector-tail proposition
identify the relevant tail as this product modulo finite support.
There is no objection to the analytic coherent sheaf
$\mathcal O_{\mathbb C}/(\xi)$ on the analytic affine line.
The assertion is about its proposed direct identification with a
finite constructible algebraic object, or with one finite-dimensional
Weil representation. A different finite realization of selected
data is not ruled out; it must be constructed and shown to retain
the desired structure constants. That is the precise surviving
gate in prompt 10(b).

\paragraph{Geometric dictionary for NN-S007.}
Constructibility requires a finite algebraic stratification and
finite-rank stalks in the SGA 4 constructible category. The first
proof excludes the proposed literal support; the second excludes
the proposed literal stalk at infinity. Ordinary finite-dimensional
Weil representations and the weight definition of Weil II 1.2.2
cannot be assigned to that infinite tail without an additional
finite quotient or realization. The Boolean skeleton's scalar
linearization instead vanishes by NN-S003. These are distinct
operations, neither of which creates the missing finite object.


\section{NN-S008: a corrected spectral bridge and the scope of RH symmetry}
\label{sec:NN-S008}

\begin{proposition}[NN-S008]
The equality in corpus 09 and corpus 10, Theorem
\texttt{thm:bridge}, must be restricted to
\[
 \left\{z:\ |\operatorname{Im}z|<\tfrac12,
               \ \zeta(\tfrac12-iz)=0\right\},
\]
or written using $\xi$ in place of $\zeta$. Without the restriction
it includes trivial zeros and is false for the intended nontrivial
spectral support. Under A, the involutions $c(s)=\bar s$ and
$w(s)=1-s$ have identical actions on the reduced nontrivial zero
support. Their Klein-four action has image of order two and kernel
$\{1,cw\}$. This statement about support does not identify their
actions on the local jet algebras.
\end{proposition}

\begin{proof}
The trivial zero $\rho=-2$ maps under $z=i(\rho-1/2)$ to
$-5i/2$. It satisfies the unrestricted displayed zeta equation
and is outside the intended strip. Conversely the affine inverse
$s=1/2-iz$ sends the displayed strip exactly to $0<\Re s<1$,
so the restricted equality follows from the stated Connes--Consani
bridge. The completed entire $\xi$ already removes the trivial
zeros and poles. On the critical line $\bar\rho=1-\rho$, giving
$c=w$ on the support under A. Complex conjugation is nontrivial
there because there are nonreal zeros, so the image has order two.
At a fixed zero the composite still acts semilinearly on the
local coordinate by $z\mapsto-\bar z$; NN-S006 describes that
action on jets. The reduced-point calculation cannot erase it.
\end{proof}

\paragraph{Spectral and geometric dictionary.}
The Connes--Consani input is \emph{Hochschild homology, trace map
and $\zeta$-cycles}, equation (2.2), the common-zero assertion in
Section 2, and Proposition 4.2. Its spectral values become
$-\gamma^2-1/4$ under A. Real negative eigenvalues alone do not
prove self-adjointness or semisimplicity: a two-by-two Jordan
block with a real negative eigenvalue is a counterexample.
In geometric language the reduced support loses infinitesimal
and equivariant data; the real Artin schemes of NN-S006 retain
the latter. The passage from a scheme to its reduction is EGA I,
Section 5.1; no finite-field realization or canonical-basis
positivity follows from this reduction.



\section{NN-S009: a supported zero does not retain cancellation history}
\label{sec:NN-S009}

\begin{theorem}[NN-S009]
For any nonzero commutative ring $R$, the value in $G(R)$ of a
formal supported expression cannot distinguish a zero supplied
as a constant from a zero obtained by cancellation. The same
failure holds for a representation ring $R(\Gamma)$, even with
the grade or other ambient coordinate specified. Thus
``no supported zero'' is not an intrinsic equivalent of a
positive formula unless the encoding of absence and the allowed
formulas are specified.
\end{theorem}

\begin{proof}
The supported constant $e=0_R$ and the expression
$1\mathbin\oplus(-1)$ have the same value $e$ and the same
support bit $1$. In a representation ring the zero object and
the virtual expression $[V]-[V]$ likewise have the identical
supported value. More generally a class $r$ equals
$(r+a)-a$, so value and support cannot reveal whether an
expression used subtraction. Prescribing a grade does not change
any of these equalities; it only records that the value is to be
placed at that grade. Therefore any predicate of the resulting
$G(R)$-value takes the same answer on these pairs.

For comparison, with a \emph{specified witness encoding} the sum
$\bigoplus_{y\in W}1$ equals $\tau$ exactly when $W$ is empty.
Indeed a nonempty finite sum of supported terms remains supported,
even if its amplitude vanishes. This is a statement about that
encoding, and does not give a method of computing a sum over an
implicitly defined witness set.
\end{proof}

\paragraph{Application to GCT and moonshine.}
Corpus 11, its introductory four-way distinction, the two-zero-fibre
lemma, and the support-complete character inversion theorem correctly
separate a missing coordinate from a chosen zero coordinate.
They do not separate the last two constructions in the proof.
In particular a PH0 claim concerns effectivity of structure
constants in a chosen canonical basis or positive formula.
It is not a consequence of calling a coefficient supported.
PH1 requires further controlled counting geometry. The theorem
does not refute either correctly specified assertion.

\paragraph{Geometric dictionary for NN-S009.}
In the Grothendieck group of representations or vector bundles,
$[V]-[V]=0=[0]$ by its defining relations. The effective cone is
the image of actual objects and does not record a chosen
presentation of a virtual class. SGA 6, Expos\'e IV, Definition
equation (1.5.2) and Definition 2.2, provide this exact group-completion
dictionary. The support annotation is extra grading or
stratification data; it does not restore a lost exact-sequence
or expression history. The conclusion concerns both orbit-ring
multiplicities and representation-valued moonshine coefficients.



\section{NN-S010: the universal joint quotient for support and residue}
\label{sec:NN-S010}

\begin{theorem}[NN-S010]
For every integer $m\geq2$, the map
\[
 n\longmapsto(n\bmod m,\mathbf1_{n>0})
 \quad\text{from }\mathbb N\text{ to }(\mathbb Z/m\mathbb Z)\times\mathbb B
\]
has image isomorphic to $G(\mathbb Z/m\mathbb Z)$, with $m+1$
elements. It is the coarsest semiring quotient retaining both
residue and support: its congruence identifies precisely the
positive integers congruent modulo $m$, and leaves zero alone.
Neither residue alone nor support alone determines the other.
For $m=2$, addition is exactly the pair (XOR, OR).
\end{theorem}

\begin{proof}
Both coordinates are semiring homomorphisms: for nonnegative
integers positivity of a sum is disjunction and positivity of a
product is conjunction. Their joint image consists of $(0,0)$
and all $(a,1)$, since every residue has a positive representative.
The induced addition and multiplication are the pair normal form
of $G(\mathbb Z/m\mathbb Z)$. The kernel equivalence has precisely
the claimed classes. If another quotient retains both maps, two
integers identified in it must have identical joint images, so
it factors onto this image. This proves the universal claim.
Zero and $m$ have equal residue and different support; $1$ and
$m$ have equal support and different residues. In characteristic
two the residue sum is XOR and the second coordinate is OR.
\end{proof}

\paragraph{Trace and complexity consequence.}
A count divisible by $p$ is either zero or a positive multiple
of $p$. A mod-$p$ trace distinguishes neither case until support
is independently available. Conversely Boolean existence does
not compute a residue. For the natural one element there is
no unital semiring map $\mathbb B\to\mathbb F_p$, since
$1+1=1$ would force $1=0$; there is no unital map in the other
direction since summing $p$ copies of $1$ would send $0$ to $1$.
E supplies exact counts for its stated instance classes and
therefore both coordinates within its stated running time.
That time is subexponential, not polynomial; A supplies no
further general count-recovery procedure from these two maps.

\paragraph{Geometric dictionary for NN-S010.}
For a finite reduced scheme of $n$ rational points, the two
statistics are nonemptiness and the trace $n\bmod p$ on its
degree-zero coefficient space. The theorem identifies the
smallest semiring of joint values. The finite-fibre trace is
the sum over fibres in SGA 4, Expos\'e XVII, Theorem 6.2.3;
the coherent mod-$p$ point-count relation used elsewhere in
this catalogue is SGA 7 II, Expos\'e XXII, Theorem 3.1 and
Corollary 3.2. The abstract quotient itself neither gives a
polynomial-size cohomology presentation nor computes its trace.



\section{NN-S011: split determinant support is efficiently evaluable but not multiplicative}
\label{sec:NN-S011}

\paragraph{Inputs.}
Corpus 16, Theorem \texttt{thm:split-determinant}, identifies the
amplitude with the determinant and the support with the Boolean
permanent. The following computation and counterexample isolate
what that pair does and does not imply for prompt 12.

\begin{theorem}[NN-S011]
For a matrix with entries $\tau$ or $1$ in $G(\mathbb F_2)$,
its split determinant can be computed in polynomial time: its
amplitude is the parity of perfect matchings and its support is
their existence. The split determinant on all matrices is not
multiplicative, even in dimension two.
\end{theorem}

\begin{proof}
Compute the determinant of the amplitude matrix over $\mathbb F_2$
by Gaussian elimination. In this characteristic the determinant
is the permanent modulo two. Compute whether the bipartite graph
of supported entries has a perfect matching by the augmenting-path
algorithm. This algorithm starts at the empty matching and finds
an augmenting path by graph search, increasing cardinality each
time. If none exists, the symmetric difference with any putative
larger matching would contain such a path, so the current matching
is maximum. At most the matrix dimension many searches suffice.
The pair of answers is one of $(0,0),(0,1),(1,1)$ and determines
the split determinant. These are polynomially many field and
graph operations.

For nonmultiplicativity take
\[
 A=\begin{pmatrix}1&\tau\\1&\tau\end{pmatrix},\qquad
 B=\begin{pmatrix}1&1\\\tau&\tau\end{pmatrix}.
\]
Both split determinants are $\tau$, because both graphs lack a
perfect matching. Their matrix product has all four entries $1$.
Its two supported determinant terms sum to $e$, the supported
zero in characteristic two. Hence
$\det_{\rm sp}(AB)=e\ne\tau=
\det_{\rm sp}(A)\det_{\rm sp}(B)$.
\end{proof}

\paragraph{Geometric dictionary and circuit consequence.}
The ordinary determinant is the action on the top exterior power
of a free module; its multiplicativity is an identity in that
additive category. The Boolean permanent records a different
incidence relation, and the counterexample proves that this
support refinement is not a determinant functor with the same
composition axiom. The exterior-power operations are SGA 6,
Expos\'e V, Section 2. The specialized evaluation algorithm above
does not give a small monotone arithmetic circuit for the
universal support polynomial; it uses a graph algorithm. Thus
monotone circuit lower bounds for matching are compatible with
efficient evaluation of this supported-zero pair. No analogous
algorithm for arbitrary sparse quadratic satisfiability follows.


\section{NN-S012: three different supported-zero circuit restrictions}
\label{sec:root-circuit-models}

\paragraph{The models and exact statement.}
Let $R$ be a nonzero commutative unital ring. Circuits have fan-in two
primitive gates $\oplus,\odot$ in $G(R)$, constants in $G(R)$, and
Boolean inputs encoded by $0\mapsto\tau$, $1\mapsto1_R$.
Write $e=0_R$ for the supported zero, $\pi:G(R)\to R$ for amplitude,
and $\sigma:G(R)\to\mathbb B$ for support. A circuit computes a
Boolean function by amplitude if its final amplitude is always the
element $0_R$ or $1_R$ representing that function. Distinguish the
following three restrictions:
\begin{enumerate}
\item Reading only final support $\sigma$ imposes no restriction on
      intermediate amplitudes.
\item Requiring the final value to be exactly $\tau$ or $1_R$ forbids
      a supported zero at the output, but permits it at internal gates.
\item Requiring that no gate, including a constant gate, takes the
      value $e$ on any Boolean input forbids supported zeros at every
      primitive gate.
\end{enumerate}
In the third model the final amplitude is again required to be Boolean.
Circuit sizes count primitive gates, with a fixed constant cost for each
ring constant.

For every $R$, models 1 and 2 compute precisely the monotone Boolean
functions, and their minimum circuit sizes are within constant factors
of monotone Boolean circuit size. The same holds for model 3 whenever
$R\not\cong\mathbb F_2$. Over $\mathbb F_2$, however, model 3 computes
only constants and conjunctions of subsets of the input variables.
In particular, two-variable OR has no such circuit over $G(\mathbb F_2)$,
although it has a constant-size such circuit over $G(\mathbb F_4)$.

If $R$ is fixed and finite, unrestricted amplitude computation over
$G(R)$ is equivalent, up to constant factors in circuit size, to ordinary
Boolean computation. The finite-ring hypothesis is needed for this
automatic simulation: no claim about unit-cost arbitrary integer or
complex constants is implicit.

\paragraph{Proof of the support and output statements.}
At every gate the pair $(\pi,\sigma)$ evolves by ring addition and
multiplication in the first coordinate and OR and AND in the second.
The support of each constant is fixed. Induction on the circuit therefore
gives a monotone circuit of the same size for every gate's support.
This proves the lower bounds and monotonicity in model 1. In model 2,
the support of the output equals its Boolean amplitude, proving the same
lower bound and monotonicity. In model 3 this equality also holds at the
Boolean output.

For the upper bound in model 1, replace Boolean OR by $\oplus$ and AND
by $\odot$; only the support is read. For model 2, use multiplication
for AND and the following constant-size circuit for OR on inputs
$a,b\in\{\tau,1_R\}$:
\[
 (a\oplus b)\oplus((-1_R)\odot(a\odot b)).             \tag{S012.1}
\]
Its amplitude is $a+b-ab$, with $\tau$ sent to zero. Its support is
the Boolean OR of the two inputs. Checking their four possible values
shows that its final value is $\tau$ for $(\tau,\tau)$ and $1_R$
otherwise. Thus each Boolean gate is replaced by a fixed number of
primitive gates and its output has the required encoding. Internal
supported zeros have not been excluded in this construction.

\paragraph{Proof for every-gate absence.}
First let $R=\mathbb F_2$ and suppose a circuit satisfies model 3.
Every gate takes values in $\{\tau,1\}$. If both inputs of an addition
gate have nonempty support as functions of the input bits, their
monotone support functions both take value one at the all-ones input.
The two input values there are both $1$, so the addition gate takes
the forbidden value $e=1\oplus1$. Hence at least one input of each
addition gate is identically $\tau$. Delete that gate and retain its
other input. This can also be done effectively by evaluating support
at the all-ones input: a monotone support is identically zero exactly
when it vanishes there. What remains has only multiplication gates,
the constants $\tau,1$, and input variables. It computes a constant
or a conjunction. Conversely all those functions have such circuits.

If $2\ne0$ in $R$, circuit (S012.1) already satisfies model 3.
The first sum has nonzero amplitude $1$ or $2$ whenever supported;
the product and its negative have amplitude $1$ and $-1$ whenever
supported, and the output has just been checked. All constants used
are nonzero. This proves the monotone upper bound in that case.

Suppose instead that $2=0$ and $R\not\cong\mathbb F_2$.
Choose any $\lambda\in R\setminus\{0,1\}$; no unit hypothesis is
needed. The following circuit computes OR without a supported zero
at any primitive gate:
\[
 s=a\oplus(\lambda\odot b),\qquad
 h=(1+\lambda)\oplus a,\qquad
 t=b\odot h,\qquad s\oplus t.                         \tag{S012.2}
\]
The nonzero supported values of $s$ belong to
$\{1,\lambda,1+\lambda\}$. The gate $h$ is always supported and has
value $1+\lambda$ or $\lambda$, both nonzero. The gate $t$ is absent
when $b=\tau$, and otherwise has one of these same two nonzero values.
The final value is $\tau,1,1,1$ on the ordered inputs
$(\tau,\tau),(1,\tau),(\tau,1),(1,1)$, respectively. All constants
and gates therefore satisfy the restriction. AND is multiplication
on $\{\tau,1\}$, so substituting these circuits at monotone gates
proves the upper bound for every remaining nonzero ring. The lower
bound was the support shadow. This completes the classification.

\paragraph{Proof of the finite-ring amplitude simulation.}
Ordinary AND and NOT are represented in amplitude by $ab$ and $1-a$.
Lifting them to $G(R)$ preserves these amplitude identities by the
homomorphism $\pi$, even when a false intermediate value is $e$ rather
than $\tau$. Thus every Boolean circuit has a constant-factor
arithmetic lift. Conversely, when $R$ is fixed and finite, every
element of $G(R)$ has a fixed-length binary encoding. Each addition
or multiplication is a fixed finite lookup table and has a Boolean
circuit of constant size. Encode all gates and extract the final
amplitude bit. This gives the reverse constant-factor simulation.
The elementary proof uses no E, A, M, or unproved complexity separation.

\paragraph{What is corrected in the supplied circuit dictionary.}
The source is \path{OURS_TRANSLATION_SPLIT_ZERO_20261121.md},
Section 9, items 1--6. Its support-shadow assertion is correct.
The phrase ``no supported zeros'' must specify model 2 or model 3.
In characteristic two the displayed OR formula
$x\oplus y\oplus xy$ has an internal supported zero at $x=y=1$.
Thus its proposed monotone equivalence is correct for output-only
absence, but false for absence at every primitive gate over
$\mathbb F_2$. Passing to $\mathbb F_4$ repairs that precise gate
restriction with (S012.2). It does not prove a Boolean circuit lower
bound: unrestricted computation over either fixed finite field still
has the same Boolean circuit complexity.

For a monotone family $f_n$, let $M_n$ be its monotone circuit size
and $C_n$ its general Boolean circuit size. The exact corrected
comparison is between $M_n$ and $C_n$, or the equivalent split models
specified above. In particular, $\mathrm{CLIQUE}\in\mathrm P/\mathrm{poly}$
is equivalent to a polynomial upper bound in the unrestricted amplitude
model. A statement merely that the ratio $M_n/C_n$ is exponential
does not provide that upper bound: as numerical growth conditions,
$M_n=2^{n^a}$ and $C_n=2^{n^a/2}$ have an exponential ratio and neither
is polynomial. This example is an explanation of the missing logical
premise, not an assertion of these actual CLIQUE complexities.
Any claimed equivalence using ``exponential savings'' must include
the polynomial endpoint, or supply an additional theorem proving it.
The same caution applies to E: its proved subexponential running times
give subexponential circuits for the encoded computations, not a
polynomial endpoint. NN-C001 and NN-C019 give the exact supported-zero
removal and instance-dependent isolation consequences available from E.

\paragraph{Geometric formulation and proof dictionary.}
Over $\mathbb F_2$, put
\[
 X=\operatorname{Spec}
 \mathbb F_2[x_1,\ldots,x_n]/(x_1^2-x_1,\ldots,x_n^2-x_n).
\]
This is the finite disjoint union of the $2^n$ rational points indexed
by Boolean assignments. The amplitude at a gate is a section of its
finite free function algebra; its support is an additional upward-closed
subset of that indexing Boolean poset. In model 3 the support is exactly
the nonvanishing locus of the section. Two nonempty upward-closed
supports contain the all-ones point, and their sections both equal one
there. Their sum vanishes there in characteristic two while the support
tag remains present. This is precisely the addition obstruction proved
above. It is a statement about the actual presentation by functions and
support tags, not about a virtual Grothendieck class or a Frobenius weight.

The finite-free algebra and its base-change trace are the setting of
EGA IV$_4$, 18.2.1, especially equation (18.2.1.2), and 18.2.2,
original printed page 111. The original page image was read in this
run. Base change to $\mathbb F_4$ keeps the same split finite scheme
and supplies the additional constant $\lambda$ used in (S012.2).
Thus every-gate zero avoidance can change under a finite scalar
extension even though the Boolean function and its support poset do
not change. The support tags are extra data beyond this algebra;
NN-S003 explains why additive ringization forgets them. For a general
ring, absence of the zero element at an $R$-valued input is not
silently identified with invertibility or with a base-change-stable
open condition. The general-ring theorem was proved directly in
$G(R)$, and the displayed scheme dictionary is specifically the
finite-field case. Neither purity nor A or M imposes the missing
polynomial upper or lower circuit bound.
\hfill$\square$


\section{NN-S013: zero avoidance, inversion, and necessary cancellation in generators}

\paragraph{Statement.}
Use the circuit models of NN-S012, and give multi-output circuits a
polynomial number of output bits. The following refinements hold.
\begin{enumerate}
\item Every map given by a circuit over $G(\mathbb F_2)$ that avoids
      supported zeros at every primitive gate can be inverted in
      deterministic polynomial time on its image, given the circuit.
      The algorithm also decides whether a proposed output has a preimage.
\item Ordinary one-way functions exist if and only if there are one-way
      functions with uniform polynomial-size circuits over
      $G(\mathbb F_4)$ avoiding supported zeros at every primitive gate.
      They also exist if and only if there are such functions over
      $G(\mathbb F_2)$ avoiding supported zeros only at the outputs.
      These equivalences retain the usual uniform or nonuniform
      adversary convention on both sides.
\item Let a polynomial-size circuit over $G(R)$, for a fixed finite
      nonzero ring $R$, have $n$ Boolean inputs and at least $n+1$ Boolean amplitude
      outputs. If its amplitude map is a pseudorandom generator secure
      against nonuniform polynomial-size distinguishers, then among
      its first $n+1$ outputs some output takes the supported zero on
      at least $\Omega(2^n/n^2)$ Boolean inputs, for all sufficiently
      large $n$. The selected output may depend on $n$.
\end{enumerate}
The third conclusion concerns the final gates themselves, not merely
some unspecified internal cancellation. There is also an unconditional
quantitative inequality, proved below, for every such map whether or
not it is pseudorandom. None of the assertions uses A or M.

The imported classical results are Goldreich--Izsak,
\emph{Monotone Circuits: One-Way Functions versus Pseudorandom
Generators}, Theory of Computing 8 (2012), 231--238:
Theorem 1.1, printed p.~232, and its proof in Section 3;
Theorem 1.2, printed p.~232, with the one-coordinate/two-coordinate
distinguishers constructed in Section 4.2, pp.~236--237.
Their Theorem 3.1 and proof, p.~234, provide the precise published
slice-circuit result attributed there to Berkowitz's 1982 unpublished
manuscript. No unlocated theorem number in that manuscript is assumed.

\paragraph{Proof of efficient inversion.}
The effective deletion in NN-S012 turns the promised circuit into
multiplication gates, constants and input variables. Associate to each
nonzero gate the subset of variables in its conjunction; associate a
separate zero marker to a gate that is identically $\tau$.
Multiplication takes unions unless one factor has the zero marker.
The constant one corresponds to the empty subset. These sets are
computed with polynomially many bit-set operations.

For a proposed output $y$, first reject if it assigns one to a constant
zero output or zero to a constant one output. Let $A$ be the union of
all conjunction subsets belonging to outputs where $y_i=1$. Set
$x_j=1$ for $j\in A$ and $x_j=0$ otherwise. Every proposed preimage
must have all bits in $A$ set to one, so this is the smallest possible
preimage in the coordinatewise order. It satisfies every output that
is required to be one. If a remaining nonconstant output is required
to be zero, with conjunction subset $S_i$, then it vanishes at this
$x$ exactly when $S_i\not\subseteq A$. If instead $S_i\subseteq A$,
every possible preimage would make that output one, so no preimage
exists. Testing these inclusions therefore either returns a correct
preimage or proves emptiness. The entire computation is polynomial
in the circuit and input/output sizes. For a uniform family the
circuit is itself polynomial-time constructible; for nonuniform
families it can be supplied as polynomial advice. Thus this family
cannot furnish one-wayness in the corresponding convention.

\paragraph{Proof of the one-way-function equivalences.}
By Goldreich--Izsak Theorem 1.1, existence of one-way functions implies
existence of such functions computed by uniform polynomial-size
monotone Boolean circuits. Replace their AND and OR gates by the
explicit every-gate zero-free $G(\mathbb F_4)$ gadgets of NN-S012.
This changes neither their Boolean input/output map nor their
inversion problem and only multiplies circuit size by a constant.
Alternatively use the output-only $G(\mathbb F_2)$ gadgets of that
finding. Their internal supported zeros are allowed in this second
model. Both conversions are uniform. Conversely, any alleged
one-way map in either finite-ring model is an ordinary efficiently
computable Boolean map by NN-S012's finite lookup-table simulation.
The same adversary and input distribution give ordinary one-wayness.
This proves both equivalences. The conclusion is conditional only
on existence of ordinary one-way functions; it does not assert that
they exist or that E preserves any stronger exponential security.

\paragraph{Proof of the necessary output-cancellation bound.}
Retain the first $n+1$ outputs. Let $F(x)\in\{0,1\}^{n+1}$ be their
amplitude vector and $S(x)\in\{0,1\}^{n+1}$ their support vector.
The latter is a monotone map, regardless of any internal cancellations.
For each output $i$ put
\[
 B_i=\{x:C_i(x)=e\}.
\]
Since the amplitude outputs are Boolean, the only possible split
values are $\tau,e,1$. Hence $F_i(x)\ne S_i(x)$ if and only if
$x\in B_i$.

Goldreich--Izsak's quantitative Theorem 1.2 and its explicit proof
give an absolute $c>0$, for every sufficiently large $n$, a set
$J\subseteq\{1,\ldots,n+1\}$ of size one or two, and a monotone
Boolean test $D$ depending only on those coordinates, such that
\[
 |\mathbb E D(S(U_n))-\mathbb E D(U_{n+1})|\ge c/n^2.
\]
The test is either a coordinate or the AND of two coordinates.
Couple $F(U_n)$ and $S(U_n)$ using the same input. Their $D$-values
agree outside $\bigcup_{i\in J}B_i$. The triangle inequality gives
the unconditional estimate
\[
 |\mathbb E D(F(U_n))-\mathbb E D(U_{n+1})|
 \ge \frac{c}{n^2}
       -\Pr\left[U_n\in\bigcup_{i\in J}B_i\right].       \tag{S013.1}
\]
For a pseudorandom generator the left side is negligible, since the
chosen tests form a family of constant-size nonuniform circuits.
Consequently the probability on the right is at least
$c/n^2-\operatorname{negl}(n)$. The union bound and $|J|\le2$
give, for at least one selected output,
\[
 \Pr[C_i(U_n)=e]\ge c/(2n^2)-\operatorname{negl}(n).
\]
Multiplication by $2^n$ proves the stated cardinality bound.
Truncating a longer pseudorandom output to its first $n+1$ bits
preserves security, so the argument applies to every positive stretch.
The quantitative statement is deliberately made with nonuniform
security; (S013.1) itself has no security or uniformity hypothesis.

\paragraph{Slice elimination and the complexity gate.}
Goldreich--Izsak Theorem 3.1 gives a polynomial-time construction of
a monotone circuit for a slice function from any Boolean circuit for
it. Combining that theorem with NN-S012 yields polynomial-cost
elimination of supported zeros at the output over $G(\mathbb F_2)$,
and at every primitive gate over $G(\mathbb F_4)$. It does not give
every-gate elimination over $G(\mathbb F_2)$: even the one-slice
extension of two-variable OR is OR itself and violates the conjunction
classification. Thus the slice theorem, the one-way-function theorem,
and the pseudorandom-generator obstruction have a consistent precise
translation, but use different notions of zero avoidance.

E supplies the sparse preimage-counting algorithms of NN-C002 and
NN-C011, and the instance-dependent isolating construction NN-C019.
Those theorems retain their sparsity and witness-size parameters.
The polynomial monotone circuits furnished by the classical one-way
theorem need not have linearly many bounded-arity factors in their
original input dimension, and a subexponential bound would in any
case not be polynomial inversion. No conclusion that E destroys
one-way functions is drawn. The cancellation lower bound (S013.1)
is a distributional necessary condition for pseudorandomness, not a
general Boolean or algebraic circuit-size lower bound.

\paragraph{Geometric formulation and exact dictionary.}
On the split Boolean finite scheme $X$ of NN-S012, an output fibre
is a clopen union of Boolean points. The inversion proof computes
emptiness and a smallest point of fibres of maps whose coordinate
pullbacks are conjunction monomials. A Boolean amplitude map has
a corresponding morphism between the split Boolean schemes, while
the support map retains an additional clopen tag at each coordinate.
The locus $B_i$ is the clopen intersection where that tag is present
and the amplitude section is zero. Uniform input measure is the
normalized counting measure on the $2^n$ rational points. Thus
(S013.1) is equivalently a lower bound on the rank of the finite
coordinate algebra of one such $B_i$, divided by $2^n$.

Finite-free trace and its base change are exactly EGA IV$_4$,
18.2.1--18.2.2, original printed p.~111; the finite-fibre sheaf trace
is SGA 4, XVII, 6.2.3. Every such finite rational-point scheme has
Frobenius acting as the identity on its degree-zero cohomology.
That purity does not construct its potentially exponentially large
point-indexed basis from the circuit. The missing effectiveness is
the same distinction as NN-C006. The finite extension from
$\mathbb F_2$ to $\mathbb F_4$ changes the available zero-free circuit
presentations; it leaves the underlying Boolean finite set and its
counting distribution unchanged. Neither the trace identity nor
the sign-fixed skeleton supplies a polynomial inverter for the
monotone maps allowed over $\mathbb F_4$.
\hfill$\square$


\section{NN-S014: a logarithmic cancellation bound and a uniform distinguisher}

\paragraph{Statement and exact inputs.}
There is an absolute constant $c>0$ such that every monotone map
$S:\{0,1\}^n\to\{0,1\}^{n+1}$, for all sufficiently large $n$,
has a test depending on at most two output coordinates with distinguishing
advantage at least $c\log n/n^2$ against the uniform distribution.
The test is a coordinate or the AND of two coordinates, with its
acceptance convention possibly complemented. For a uniformly
polynomial-time computable family $S$, a uniform randomized polynomial-time
algorithm selects such a test to within an additive $O(n^{-3})$ loss,
with failure probability at most $2(n+2)^2e^{-2n^2}$.

Consequently, for a uniformly constructible polynomial-size $G(R)$ circuit
family over any fixed finite nonzero commutative ring, with Boolean
amplitude output and positive stretch, ordinary security against
\emph{uniform randomized polynomial-time} distinguishers implies that
some final output equals the supported zero on
\[
                 \Omega(2^n\log n/n^2)
\]
seed inputs. The same bound holds for nonuniform-secure generators without
the uniform circuit-construction hypothesis. This sharpens the quantitative
bound and removes the nonuniform-adversary restriction of NN-S013 for
uniform circuit families. It is a necessary condition for pseudorandomness,
not an existence assertion or a circuit-size lower bound.

The classical inputs are Goldreich--Izsak, \emph{Monotone Circuits:
One-Way Functions versus Pseudorandom Generators}, Theory of Computing
8 (2012), 231--238: Corollary~4.4, pp.~235--236, and Theorem~4.5,
p.~236, followed by the proof of Theorem~1.2 in Section~4.2,
pp.~236--237. The logarithm is already present in their covariance
estimate; the following proof changes the marginal-bias threshold.
Thus no priority claim is made for a new classical correlation theorem.
The new split-circuit and uniform-selection consequences are proved below.
The primary source is
\url{https://theoryofcomputing.org/articles/v008a010/v008a010.pdf}.

\paragraph{Proof of the sharper two-coordinate estimate.}
Write $f_i$ for the $i$th coordinate of $S$ and $I_i(j)$ for its
influence at input coordinate $j$. The KKL estimate in the cited
Corollary~4.4 gives an absolute $a>0$ such that every function with bias
$o(\log n/n)$ has some influence at least $a\log n/n$, for all
sufficiently large $n$. The constant can be chosen uniformly here:
changing a function on a fraction $\delta$ of inputs makes it exactly
balanced and changes each influence by at most $2\delta$, as proved
immediately before the cited corollary.

Suppose that every marginal bias is at most
$\delta=c_0\log n/n^2$, where the absolute $c_0>0$ will be fixed below.
Choose one such influential input for each of the $n+1$ outputs.
Two distinct outputs, say $f_i,f_k$, share a chosen input. The
Talagrand inequality in the cited Theorem~4.5 gives
\[
 \operatorname{Cov}(f_i,f_k)
 \ge c_T\frac{T}{\log(e/T)},\qquad
 T=\sum_j I_i(j)I_k(j)
 \ge a^2\frac{(\log n)^2}{n^2}.
\]
The expression on the right is increasing for $0<T\le1$ and is
at least $c_1\log n/n^2$ for an absolute $c_1>0$ and all large $n$.
For completeness, $T\le1$: for a monotone Boolean function in the
uniform cube, the absolute value of its first Fourier coefficient is one half of its
influence, so Parseval gives $\sum_j I_i(j)^2\le4\operatorname{Var}(f_i)
\le1$; Cauchy--Schwarz proves the bound for $T$.
The displayed lower estimate also follows directly by substituting
$a^2(\log n)^2/n^2$ and bounding its logarithmic denominator by
$3\log n$ for sufficiently large $n$.

Since both marginal means are at least $1/2-\delta$,
\[
 \mathbb E(f_if_k)-\frac14
 \ge c_1\frac{\log n}{n^2}-\delta+\delta^2.
\]
Fix $c_0=c_1/4$. If a marginal has bias greater than $\delta$,
that coordinate gives advantage greater than $\delta$. Otherwise
the AND of the two coordinates above has advantage at least
$3c_1\log n/(4n^2)$. Decreasing the universal constant if needed
proves the assertion for every sufficiently large $n$. The proof
uses no computation of the influential coordinates.

\paragraph{Uniform selection of the test.}
Let $\mathcal T$ consist of the $n+1$ coordinate tests and all
two-coordinate AND tests. Its size is at most $(n+2)^2$.
For each $D\in\mathcal T$, its uniform mean $u_D$ is known exactly:
it is $1/2$ or $1/4$. Generate $L=n^8$ independent samples
$S(U_n)$ and compute every empirical mean $\widehat\mu_D$.
Put $\varepsilon=n^{-3}$. The elementary Bernoulli exponential
bound and a union bound give
\[
 \Pr\left[\max_{D\in\mathcal T}
       |\widehat\mu_D-\mathbb E D(S(U_n))|>\varepsilon\right]
 \le 2(n+2)^2e^{-2n^2}=:\eta_n.                 \tag{S014.1}
\]
Here the Bernoulli bound follows by exponential Markov applied to
independent centred samples and the inequality
$\mathbb E e^{s(X-\mathbb EX)}\le e^{s^2/8}$ for $0\le X\le1$.
One proof of this latter inequality differentiates the logarithmic
moment-generating function twice: its second derivative is the variance
under an exponential tilt and is at most $1/4$; its value and first
derivative vanish at zero. Optimizing $s$ gives each tail
$e^{-2L\varepsilon^2}$, so no additional probabilistic premise is hidden.

Choose a test maximizing $|\widehat\mu_D-u_D|$. Use $D$ itself
when the empirical bias is nonnegative and $1-D$ otherwise.
On the event complementary to (S014.1), the chosen, oriented test
$D_*$ satisfies
\[
 \mathbb E D_*(S(U_n))-\mathbb E D_*(U_{n+1})
 \ge c\frac{\log n}{n^2}-2\varepsilon.           \tag{S014.2}
\]
Indeed a genuinely best test loses at most $\varepsilon$ on
estimation, empirical maximization cannot decrease that value,
and the chosen test loses at most a further $\varepsilon$.
The training samples are independent of the challenge on which
the test is evaluated. The entire selection and evaluation algorithm
is uniform and polynomial time when $S$ is uniform and polynomial time.
Neither a hard-coded influential coordinate nor nonuniform advice is
required. The exponent eight is only a convenient fixed polynomial;
no optimal sample complexity is asserted.

\paragraph{Transfer to supported zeros.}
Let $F$ be the first $n+1$ Boolean amplitude outputs and $S$ their
support outputs. As in NN-S013, $S$ is monotone and
$F_i(x)\ne S_i(x)$ holds exactly on
$B_i=\{x:C_i(x)=e\}$. Put
$b_n=\max_{i\le n+1}\Pr[U_n\in B_i]$.
Every selected test depends on at most two coordinates. For each fixed
training outcome, coupling $F$ and $S$ on the same seed changes its
expectation by at most $2b_n$; complementing the test does not change
this bound. Average (S014.2) over the independent training samples.
On the exceptional event any bias is at least $-1$. Therefore the
uniform randomized distinguisher satisfies the unconditional inequality
\[
 \mathbb E D_*(F(U_n))-\mathbb E D_*(U_{n+1})
 \ge c\frac{\log n}{n^2}-2b_n-2n^{-3}-2\eta_n.  \tag{S014.3}
\]
If $F$ is secure against uniform randomized polynomial time, the
absolute value of the left side is negligible. Thus
\[
 b_n\ge\frac{c\log n}{2n^2}-n^{-3}-\eta_n
                   -\operatorname{negl}(n)
      =\Omega(\log n/n^2).
\]
Multiplying by $2^n$ gives the cardinality claim. For nonuniform security,
one may simply select the test from the first part as advice, obtaining
the same conclusion without the training step. Truncation preserves
pseudorandomness, so every polynomial output length greater than $n$
is covered. No use of E, A, or M occurs in this distributional theorem.
\hfill$\square$

\paragraph{Geometric and complexity dictionary.}
For each input length, the Boolean input algebra is the coordinate
algebra of the finite split scheme with $2^n$ points. Each $B_i$ is a
clopen subscheme, and its finite free rank is $\#B_i$.
The conclusion is therefore a lower bound on the rank of at least one
supported-zero output fibre, with probability measured by rank divided
by $2^n$. The finite-algebra trace convention and its base-change
identity are EGA~IV$_4$, Section~18.2.1, especially equation~(18.2.1.2),
and Section~18.2.2, printed p.~111. This is an explicit finite-scheme
dictionary; no l-adic weight or purity theorem is used to obtain the
probability bound.

The construction of the clopen fibre is implicit in the circuit;
its exponentially many points need not be enumerated. The theorem
improves a necessary cancellation condition and supplies a uniform
attack when that condition fails. It gives neither a lower bound on
the total number of circuit gates nor a proof of P versus NP or
VP versus VNP. E's separate subexponential inversion consequences
retain their original sparsity and locality hypotheses.


\section{NN-S015: parity in an infinitesimal Frobenius extension invisible to zeta}

\paragraph{Statement and inputs.}
For every 3-CNF $\Phi$, let $F_\Phi$ be the degree-$r$ form in
$r$ variables of NN-C023, with odd $r\ge3$, and put
\[
 c_\Phi=[z_1\cdots z_r]F_\Phi
       =\#\operatorname{SAT}(\Phi)\pmod2,\qquad
 H=\sum_{i=1}^r z_i^r.
\]
Over the fixed dual-number base
$B=\mathbb F_2[\epsilon]/(\epsilon^2)$, define
\[
 \pi_\Phi:\mathcal X_\Phi
   =V(H+\epsilon F_\Phi)\subset\mathbb P^{r-1}_B
                  \longrightarrow\operatorname{Spec}B.
\]
This is a smooth projective relative hypersurface, constructed in
polynomial time by a compact equation. Its reduction is the fixed
Fermat hypersurface $X_0=V(H)$ for every input with the same $r$.
The based coherent line
\[
 M_\Phi=H^{r-2}(\mathcal X_\Phi,\mathcal O)=B e
\]
has absolute Frobenius given on its underlying two-dimensional
$\mathbb F_2$-space, with ordered basis $(e,\epsilon e)$, by
\[
       N_\Phi=
       \begin{pmatrix}0&0\\c_\Phi&0\end{pmatrix}.       \tag{S015.1}
\]
Thus computing its rank, testing whether it is zero, or deciding whether
the resulting extension of two zero-Frobenius lines splits is
parity-counting complete on this class.

All positive-power traces of $N_\Phi$, its characteristic polynomial,
and its class in the Grothendieck group of finite-dimensional
$\mathbb F_2$-spaces with an endomorphism are independent of $c_\Phi$.
Moreover every finite-field point count and the entire arithmetic zeta
function of $\mathcal X_\Phi$ are identical to those of $X_0$.
For each odd prime $\ell$, their geometric $\ell$-adic cohomology groups
are canonically isomorphic with Frobenius. These are pure because $X_0$
is smooth and projective. Here relative smoothness is over the
nonreduced base $B$; the total $\mathbb F_2$-scheme is not claimed smooth.

The inputs are the coefficient encoding and complete coherent calculation
of NN-C023, its smooth-pencil proof NN-C024, SGA~1, Expos\'e~II,
Theorem~4.10(iii) and Corollary~4.12, and the topological invariance
theorem SGA~4, Expos\'e~VIII, Theorem~1.1, Corollary~1.2 and
Example~1.3(a). The latter have exact locators in the primary re-edition
\url{https://www.normalesup.org/~forgogozo/SGA4/08/08.pdf},
pp.~1--2, retaining original page markers 367--369.
Purity is Deligne, \emph{La conjecture de Weil I}, Theorem~1.6,
printed p.~276. The exact-category Grothendieck-group convention is
SGA~6, Expos\'e~IV, equation~(1.5.2) and Definition~2.2.
No use of A or M is made.

\paragraph{Proof of smoothness and the based coherent line.}
The matrix in NN-C024 that expresses selected multiples of the partial
derivatives in a common monomial degree is $I+uA$. At $u=\epsilon$,
it has the explicit inverse $I-\epsilon A$. Thus all monomials of
that degree lie in the partial-derivative ideal over $B$ itself.
At every projective point some homogeneous coordinate is invertible,
so these partial derivatives generate the unit ideal locally.
Euler's identity, with odd degree $r$, converts this to the
nonvanishing of a relative chart derivative on the hypersurface.
The Jacobian criterion in SGA~1 II~4.10(iii)/4.12 gives smoothness
over $B$. Projectivity follows from the displayed closed immersion.

Multiplication by $H+\epsilon F_\Phi$ is injective on the polynomial
ring over $B$. Indeed writing a putative annihilator as
$A_0+\epsilon A_1$ gives $HA_0=0$ and then $HA_1=0$ in the
polynomial ring over $\mathbb F_2$, so both vanish. This gives
the same exact hypersurface sequence as NN-C023, now over $B$:
\[
 0\longrightarrow\mathcal O_{\mathbb P^{r-1}_B}(-r)
   \xrightarrow{\ H+\epsilon F_\Phi\ }
     \mathcal O_{\mathbb P^{r-1}_B}
   \longrightarrow\mathcal O_{\mathcal X_\Phi}\longrightarrow0.
\]
The coordinate-cover monomial computation works over any coefficient
ring: the sole top Laurent monomial of degree $-r$ is
$1/(z_1\cdots z_r)$, and all intermediate cohomology vanishes.
Hence the boundary map identifies $M_\Phi$ with the free line
on this monomial. It also gives $H^0(\mathcal O_{\mathcal X_\Phi})=B$
and no other nonzero coherent groups. This is an actual based module,
not merely a rank assertion inferred from a special fibre.

The natural Frobenius diagram from NN-C023 applies over $B$.
Absolute Frobenius is semilinear with respect to the square map on $B$;
in the displayed basis it satisfies
\[
       \varphi(b e)=b^2\epsilon c_\Phi e.
\]
The coefficient of $z_1\cdots z_r$ in $H$ is zero, which gives
this formula. For $b=b_0+\epsilon b_1$ one has $b^2=b_0$, with
$b_i\in\mathbb F_2$. Thus $\varphi(e)=\epsilon c_\Phi e$ and
$\varphi(\epsilon e)=0$, proving (S015.1). The linearized relative
Frobenius between the base-pulled coherent line and $M_\Phi$ is
multiplication by $\epsilon c_\Phi$. In this particular example the
underlying absolute map happens also to be $B$-linear, because
$(b^2-b)\epsilon=0$; no such assertion is made for general bases.

\paragraph{Proof of the extension and information-loss statements.}
Let $T$ act as $N_\Phi$, and set $k=\mathbb F_2$.
The invariant subspace $k\epsilon e$ gives an exact sequence of
$k[T]$-modules
\[
 0\longrightarrow(k,T=0)\longrightarrow(M_\Phi,N_\Phi)
     \longrightarrow(k,T=0)\longrightarrow0.          \tag{S015.2}
\]
When $c_\Phi=0$, it splits. When $c_\Phi=1$, any lift of the
quotient basis is $e+a\epsilon e$ and is sent by $T$ to
$\epsilon e\ne0$, so no invariant complementary zero-action
line exists. Thus it does not split.
More explicitly, applying $\operatorname{Hom}_{k[T]}(-,k)$ to
the free resolution
\[
 0\longrightarrow k[T]\xrightarrow{T}k[T]
       \longrightarrow k\longrightarrow0
\]
gives $\operatorname{Ext}^1_{k[T]}(k,k)=k$, since the induced
differential is zero. In these fixed bases, the class of
(S015.2) is exactly $c_\Phi$.

Both matrices in (S015.1) have characteristic polynomial $u^2$,
all positive-power traces zero, and the same associated graded
zero-action object. Relation (S015.2) gives the same class
$2[(k,0)]$ in the exact-category Grothendieck group. Their ranks,
however, are zero and one respectively. Hence trace data and the
additive Grothendieck class cannot reconstruct this extension class.
This is an explicit failure of determination, independent of any
unproved computational lower bound.

The two alternatives already occur with the same $r=3$ in the
encoding: one unconstrained Boolean variable has two assignments
and gives $c_\Phi=0$, while its positive unit clause has one assignment
and gives $c_\Phi=1$. Both have one factor--variable incidence after
the tautological-factor convention of NN-C023. Thus the common
special fibre can be the very same Fermat cubic; changing the
input size is not responsible for the lost information.

There is also an exact orbit-closure interpretation. Over any extension
of $k$, write $N=N_1$. For $t\ne0$,
\[
 \operatorname{diag}(1,t)N\operatorname{diag}(1,t)^{-1}=tN.
\]
The polynomial family $tN$ specializes to zero at $t=0$. Thus zero
lies in the closure of the algebraic $\operatorname{GL}_2$-orbit of
$N$, and every regular polynomial function invariant under this
algebraic group has the same value on the two matrices. Indeed its
value on $tN$ is constant for invertible $t$ and hence, as a polynomial
identity, is the same at zero. This concerns algebraic-group invariants,
not just invariants under the finite group $\operatorname{GL}_2(k)$.
The rank distinction is lost by this invariant quotient as well as by
traces and $K_0$. Computing the rank of an explicitly supplied
two-by-two matrix remains trivial; the parity-complete operation is
recovering its Frobenius action from the compact hypersurface equation.

For a finite extension $K/\mathbb F_2$, every map
$\operatorname{Spec}K\to\mathcal X_\Phi$ sends $\epsilon$ to zero.
It therefore factors uniquely through $X_0$. This proves the
equality of point counts over every finite extension, and hence
of the zeta functions. Equivalently the reduction immersion has
the same closed points and residue fields. This equality is for
the full arithmetic zeta function, not just its reduction modulo two.

The closed immersion $X_0\hookrightarrow\mathcal X_\Phi$ is defined
by a nilpotent ideal. SGA~4 VIII~1.1,1.2,1.3(a) gives the
equivalence of their small \'etale topoi and the corresponding
cohomology isomorphisms. Apply this to each finite
$\mathbb Z/\ell^n$ coefficient sheaf and then pass to the inverse
limit and tensor with $\mathbb Q_\ell$. Naturality makes the
isomorphism compatible with arithmetic Frobenius. The special
fibre $X_0$ is smooth projective because its homogeneous derivatives
are $z_i^{r-1}$; Deligne's theorem gives the asserted purity.
Thus even the entire finite Weil-cohomology system does not
determine the coherent extension in (S015.2).

The reduction $\Phi\mapsto\mathcal X_\Phi$ has polynomial size,
and rank one of $N_\Phi$ is equivalent to odd parity-SAT. Membership
in parity-P on this encoded class is the monomial-choice computation
in NN-C023; complementing the parity convention covers a zero or
splitting test. This proves the completeness assertion.
A deterministic polynomial algorithm for these tasks would give
$\mathrm P=\oplus\mathrm P$ and then $\mathrm{NP}=\mathrm{RP}$
by the same isolation argument as NN-C006. E instead supplies the
subexponential evaluation through the original formula, with the
same input-parameter bounds as NN-C023. No unconditional circuit
lower bound follows from completeness alone.
\hfill$\square$

\paragraph{Grothendieck formulation and the remaining gate.}
The finding is a theorem about a smooth projective first-order
deformation over a length-two local base, its relatively free
coherent cohomology line, and the extension class of its absolute
Frobenius after forgetting to the residue field. The exact
dictionary is the based hypersurface sequence and (S015.2).
The \'etale invariance theorem explains precisely which realization
forgets this infinitesimal information. The Grothendieck-group
relation forgets the nonsplit extension for a different reason:
it imposes additivity along exact sequences.

All fibres in a fixed dimension have the same pure special-fibre
Weil data, yet computing this coherent extension class is
parity-counting complete from the succinct equation. This refutes
the unrestricted inference that a known pure finite Weil object,
its zeta function, or its $K_0$ class supplies an efficient
computation of all related deformation data. It does not identify
the extension with a GCT canonical-basis multiplicity, nor decide
P versus NP or VP versus VNP. Such an identification and an efficient
coefficient-evaluation construction would be additional hypotheses.


\section{NN-S016: an explicit Chow basis and sphere homotopy type with a counting-hard rank}

\paragraph{Statement and exact inputs.}
Let $X_\Phi/\mathbb Z$ be the smooth sparse cubic of NN-C027, and write
\begin{gather*}
 X_\Phi:\quad t^2-t=\sum_{i=1}^h y_i f_i(x_1,\ldots,x_v),\\
 Y_\Phi=V(f_1,\ldots,f_h)=\coprod_{a\in\mathcal S}\operatorname{Spec}\mathbb Z,
 \qquad |\mathcal S|=N=\#\operatorname{SAT}(\Phi).
\end{gather*}
Here $v,h\ge1$, the retained equation is $f_1=x_1^2-x_1$, and
$\mathcal S$ consists of the uniquely extended satisfying Boolean assignments.
For every field $k$, the variety $X=X_\Phi\times k$ is $k$-rational,
smooth and geometrically integral. For $a\in\mathcal S$ and $e\in\{0,1\}$,
put
\[
 D_{a,e}=V(x-a,t-e)\subset X,\qquad D_{a,e}\simeq\mathbb A^h_k.
\]
These are closed smooth subvarieties of codimension $v$. The full integral
Chow ring is
\begin{equation}\label{eq:S016-chow}
 \operatorname{CH}^*(X)\simeq
 \mathbb Z[u_a:a\in\mathcal S]/(u_a u_b:a,b\in\mathcal S),
 \qquad \deg u_a=v,
\end{equation}
where $u_a=[D_{a,0}]$ and $[D_{a,1}]=-u_a$.
In particular there is no integral torsion, and the only nonzero positive
codimension Chow group is $\operatorname{CH}^v(X)=\mathbb Z^N$.
Every element of this group has an effective cycle representative.
For every $\ell\ne\operatorname{char}k$, the cycle-class map is an isomorphism
\begin{equation}\label{eq:S016-cycle}
 \operatorname{CH}^j(X)\otimes\mathbb Q_\ell
       \xrightarrow{\ \sim\ }
 H^{2j}(X_{\bar k},\mathbb Q_\ell(j))
 \quad\hbox{for all }j\ge0.
\end{equation}
All odd geometric cohomology vanishes. The positive codimension generators
are already defined over the ground field.

Over $\mathbb C$ there is also a homotopy equivalence
\begin{equation}\label{eq:S016-wedge}
                  X_\Phi(\mathbb C)\simeq\bigvee_{a\in\mathcal S}S^{2v},
\end{equation}
where the empty wedge is a point. Thus exact computation of the positive
Chow rank, the nonzero reduced Betti number, or the number of spheres is
\#P-complete on this encoded family. Noncontractibility is NP-complete and
contractibility is coNP-complete on the same family. In addition, $N=0$ if and only if $X_\Phi$ is isomorphic to
$\mathbb A^{v+h}_{\mathbb Z}$, and if and only if its fiber over any
specified field is isomorphic to affine space. Affine-space recognition
is therefore coNP-complete on this family. These assertions concern
the supplied circuit-and-equation encoding; no such upper complexity bound
for arbitrary affine cubics is claimed.

Besides NN-C027, the algebraic inputs are localization for Chow groups
(Stacks Project, Lemma~42.19.3, Tag~02RX), affine-bundle surjectivity
(Lemma~42.32.1, Section Tag~02TS), and homotopy invariance for vector bundles
(Lemma~42.36.3, Tag~02TY). The cohomological inputs are SGA~4, Expos\'e~XVI,
Theorem~3.7 and Remark~3.10(a), and Expos\'e~XVIII, Theorem~3.2.5 and
equation~(3.2.6.2); SGA~5, Expos\'e~VII, Proposition~1.1 and Corollary~1.2
give affine-space invariance. The cycle normalization is Grothendieck,
written by Deligne, SGA~4\,$\frac12$, \emph{Cycle}, Sections~2.2.2 and~2.3.2.
The cycle-map compatibility with rational equivalence is also recorded
precisely in Milne, \emph{Lectures on \`Etale Cohomology}, version~2.21,
Section~23, Theorem~23.4 and the ensuing rational $\ell$-adic construction,
printed pp.~141--142.
The needed topological inputs are located in the proof below.
Neither A nor M is needed for these geometric and complexity statements.

\paragraph{Proof of rationality and generation of the Chow ring.}
On the nonempty principal open $f_1\ne0$ the defining equation solves for
\[
 y_1=\frac{t^2-t-\sum_{i=2}^h y_i f_i(x)}{f_1(x)}.
\]
Consequently this open is $D(f_1)\times\mathbb A^h_k$, with coordinates
$(x,t,y_2,\ldots,y_h)$. This proves rationality over every field, including
fields over which this particular open has no rational points. Smoothness,
geometric integrality and the smooth projection
$\pi:X\to\mathbb A^v_k$ of relative dimension $h$ are NN-C027.

Set $B=\mathbb A^v_k\setminus Y_\Phi$, $D=\coprod_{a,e}D_{a,e}$ and
$U=X\setminus D=\pi^{-1}(B)$. The open sets $D(f_i)$ cover $B$.
On each, solve for $y_i$ as above: the resulting $U\to B$ is Zariski
locally a product with $\mathbb A^h$. No globally chosen right inverse
to the row $(f_i)$ is asserted or needed.

Homotopy invariance gives $\operatorname{CH}^j(\mathbb A^v_k)=0$ for
$j>0$. Localization then gives $\operatorname{CH}^j(B)=0$ for $j>0$.
Affine-bundle surjectivity, which suffices here without any injectivity
claim, gives $\operatorname{CH}^j(U)=0$ for $j>0$. Since each component
of $D$ is an affine $h$-space, its only nonzero Chow group by dimension
is $\operatorname{CH}_h=\mathbb Z$. Localization for $D\subset X$ now
implies
\[
 \operatorname{CH}^j(X)=0\quad(j>0,\ j\ne v),\qquad
 \bigoplus_{a,e}\mathbb Z[D_{a,e}]\twoheadrightarrow
                            \operatorname{CH}^v(X).
\]
This calculation includes $j>v$: the possible source is then a lower
dimensional Chow group of $D$, which is zero, rather than an omitted
localization term.

Flat pullback by $\pi$ carries $[a]\in\operatorname{CH}^v(\mathbb A^v_k)$
to $[D_{a,0}]+[D_{a,1}]$. Both components have multiplicity one since
$t(t-1)$ has distinct simple roots over every field. As $[a]=0$ on
affine $v$-space, we obtain
\begin{equation}\label{eq:S016-opposite}
                     [D_{a,0}]+[D_{a,1}]=0.
\end{equation}
It follows that the $N$ classes $u_a$ generate $\operatorname{CH}^v(X)$.
Geometric integrality gives $\operatorname{CH}^0(X)=\mathbb Z[X]$.
The remaining issue is that the displayed generators have no further
relations, including torsion relations.

\paragraph{Proof of independence and of the cycle-class assertion.}
First pass to an algebraic closure and take $\mathbb Q_\ell$ coefficients.
The ordinary geometric cohomology of $X$, by the compact-support
calculation in NN-C027 and Poincar\'e duality in dimension $v+h$, is
\[
 H^0(X)=\mathbb Q_\ell,\qquad
 H^{2v}(X)=\mathbb Q_\ell(-v)^N,
 \qquad H^i(X)=0\ (i\ne0,2v).
\]
The relative compact-support splitting in NN-C027 is a sheaf identity
obtained from trace and involution projectors. Its geometric base-change
calculation works over every algebraically closed field, including
characteristic two; it does not infer these groups from Frobenius counts.
For $N=0$ the indicated second group is zero. Purity for
$Y_\Phi\subset\mathbb A^v$ and localization give
$H^{2v}(B)=0$: the relevant term to the right is
$H^{2v+1}_{Y_\Phi}(\mathbb A^v)=H^1(Y_\Phi,\mathbb Q_\ell(-v))=0$,
and affine space has no positive ordinary cohomology. On every
trivialization $D(f_i)$, affine-space invariance identifies
$\mathbb Q_\ell$ with $R(\pi|_U)_*\mathbb Q_\ell$.
The identification is the adjunction unit and hence glues; therefore
$H^*(U)=H^*(B)$ and in particular $H^{2v}(U)=0$.

Purity for the smooth closed pair $D\subset X$ and the cohomology
sequence with supports yield a surjection
\[
 \mathbb Q_\ell(-v)^{2N}=H^0(D,\mathbb Q_\ell(-v))
    =H_D^{2v}(X,\mathbb Q_\ell)
       \longrightarrow H^{2v}(X,\mathbb Q_\ell).
\]
The standard basis maps to the cycle classes of $D_{a,e}$, by the
fundamental-class normalization. Equation~\eqref{eq:S016-opposite}
also holds after applying the cycle map. Thus the $N$ images of
$u_a$ alone span the $N$-dimensional space $H^{2v}(X)$ and form a basis.
If an integral relation between the $u_a$ held over the original field
$k$, it would hold after base change and after applying this cycle map.
Its integer coefficients must therefore all be zero. This proves
integral independence, not only independence modulo torsion.

The presentation~\eqref{eq:S016-chow} now follows: a product of two
positive codimension classes lies in $\operatorname{CH}^{2v}(X)=0$
because $v\ge1$. The same bases prove~\eqref{eq:S016-cycle}, and all
other degrees are zero or the unit degree. Finally, an arbitrary
class $\sum_a n_a u_a$ is represented by the effective cycle
\[
 \sum_{n_a\ge0}n_a D_{a,0}
          +\sum_{n_a<0}(-n_a)D_{a,1}.
\]
Thus the monoid of effective codimension-$v$ cycle classes is the
entire group. This last fact uses affineness and rational equivalence;
it must not be treated as a pointed effective cone of a projective variety.

% Addition to root-owned NN-S016; not a separate registered finding.
% Standard amsmath/amssymb/amsthm environments only.
\begin{proposition}[NN-S016: complex homotopy type and the affine-space case]
Retain the sparse quadratic equations $f_1,\ldots,f_h$ of NN-C027,
with $v\geq1$ and
\[
 Y=V(f_1,\ldots,f_h)\simeq\coprod_N\operatorname{Spec}\mathbb Z,
 \qquad
 X=V\left(t^2-t-\sum_{i=1}^h y_i f_i(x)\right).
\]
The complex realization has homotopy type
\[
                      X(\mathbb C)\simeq\bigvee_{a=1}^{N}S^{2v},
\]
where the empty wedge means a point. In particular its integral
cohomology is $\mathbb Z$ in degree zero, $\mathbb Z^N$ in degree $2v$,
and zero otherwise, with every product of positive-degree classes zero.
Moreover,
\[
 N=0\quad\Longleftrightarrow\quad
 X\simeq\mathbb A^{v+h}_{\mathbb Z}
 \quad\Longleftrightarrow\quad
 X_{\mathbb C}\simeq\mathbb A^{v+h}_{\mathbb C}
 \quad\Longleftrightarrow\quad X(\mathbb C)\text{ is contractible}.
\]
The first isomorphism, when it exists, is proved by an explicit
polynomial formula from any identity $\sum_i a_i f_i=1$.
No Quillen--Suslin theorem or effective bound for such an identity is
required.
\end{proposition}

\begin{proof}
Write $\pi:X(\mathbb C)\to\mathbb C^v$ for the projection, and identify
$Y(\mathbb C)$ with its $N$ distinct points. First suppose $N>0$.
We prove the local models and the attaching maps before determining
the homotopy type.

Put $U=\mathbb C^v\setminus Y(\mathbb C)$ and
$A=\pi^{-1}(U)$. For $x\in U$ define the continuous column vector
\[
 a(x)=\frac{(\overline{f_1(x)},\ldots,\overline{f_h(x)})}
                  {\sum_i|f_i(x)|^2},\qquad f(x)\cdot a(x)=1.
\]
Every point of $A$ can uniquely be written
\[
       (x,t,y),\qquad y=a(x)t(t-1)+w,\quad f(x)\cdot w=0.
\]
Thus $A$ is the total space of the kernel vector bundle over
$U\times\mathbb C_t$. More explicitly, with $0\leq s\leq1$, set
\[
 t_s=(1-s)t,\qquad w_s=(1-s)w,\qquad
 y_s=a(x)t_s(t_s-1)+w_s.
\]
This is a strong deformation retraction onto the section $t=0,y=0$,
which is $U$. It is compatible with restriction to every open subset
of $U$. In particular projection identifies all intersection maps
below with the corresponding base-space inclusions. This is a
continuous construction; conjugation is not asserted to be algebraic.

For each $a\in Y(\mathbb C)$, reducedness says that the germs of the
$f_i$ generate the maximal ideal in the regular analytic local ring
of $\mathbb C^v$ at $a$. Their differentials therefore span the
cotangent space. Choose $v$ of them with independent differentials.
After reordering, the holomorphic inverse function theorem makes
\[
                    u_i=f_i(x),\qquad 1\leq i\leq v,
\]
coordinates on a neighborhood of $a$. In particular $h\geq v$.
Choose that neighborhood to map biholomorphically onto a small
Euclidean ball $B_a=\{\|u\|<\delta_a\}$ centered at zero, with these
neighborhoods pairwise disjoint. Each remaining holomorphic function
vanishes at zero and has an expression
\[
 f_j(x(u))=\sum_{i=1}^v b_{ji}(u)u_i\quad(j>v).
\]
For example the coefficients are obtained by integrating the partial
derivatives of $f_j(x(su))$ from $s=0$ to $s=1$, which is holomorphic
on the ball. The invertible triangular change of fiber coordinates
\[
 \eta_i=y_i+\sum_{j>v}b_{ji}(u)y_j\quad(1\leq i\leq v)
\]
turns the equation into
\[
                 \sum_{i=1}^v u_i\eta_i=t(t-1),
\]
with $h-v$ unused affine coordinates. Consequently
$V_a=\pi^{-1}(B_a)$ is analytically isomorphic to
\[
 \left\{(u,\eta,t):\sum_i u_i\eta_i=t(t-1),\ \|u\|<\delta_a\right\}
                   \times\mathbb C^{h-v}.
\]
Here and below $B_a$ also denotes its inverse image in the original
base coordinates.

The restriction $\|u\|<\delta_a$ does not change the homotopy type of
the full affine quadric $Q_{2v}:\sum u_i\eta_i=t(t-1)$.
Indeed take
\[
 c(u)=\min\left(1,\frac{\delta_a}{2\|u\|}\right),\quad c(0)=1,
 \qquad c_s(u)=1-s+s c(u).
\]
The maps $(u,\eta,t)\mapsto(c_s(u)u,c_s(u)^{-1}\eta,t)$ stay on the
quadric. At $s=1$ their image has $\|u\|<\delta_a$, and, if it starts
in that open subset, it remains there throughout the homotopy.
They supply a homotopy inverse to the inclusion of the restricted
quadric. This is a homotopy equivalence, not a claim that the whole
restricted subset is pointwise fixed during the homotopy.

The full quadric is linearly equivalent over $\mathbb C$ to
$\sum_{j=0}^{2v}z_j^2=1$: use $z_0=2t-1$ and, for each $i$,
$z_{2i-1}=\mathrm i(u_i+\eta_i)$, $z_{2i}=u_i-\eta_i$.
Writing $z=p+\mathrm i q$ with real vectors gives
\[
 |p|^2-|q|^2=1,\qquad p\cdot q=0.
\]
It is therefore diffeomorphic to the tangent bundle of $S^{2v}$,
by $(e,q)\mapsto\sqrt{1+|q|^2}\,e+\mathrm i q$ with
$|e|=1$ and $e\cdot q=0$. Scaling $q$ to zero retracts this bundle
onto its zero section. We have proved $V_a\simeq S^{2v}$, including
simple connectivity. On the intersection
\[
             W_a=A\cap V_a=\pi^{-1}(B_a\setminus\{a\}),
\]
the earlier projection retraction gives $W_a\simeq S^{2v-1}$.

Put $d=2v\geq2$. The small linking spheres give a basis of
$H_{d-1}(U;\mathbb Z)$, with no other reduced homology in $U$.
For precision, excise the complement of the balls from the pair
$(\mathbb C^v,U)$. Relative homology becomes the direct sum of the
pairs $(B_a,B_a\setminus\{a\})$. Each summand is $\mathbb Z$ in
degree $d$ and zero otherwise. As $\mathbb C^v$ is contractible,
the boundary map in the relative long exact sequence identifies
these $N$ generators with the linking-sphere basis in $H_{d-1}(U)$.
Naturality and the projection retractions prove that the inclusion
map
\[
     \bigoplus_a H_{d-1}(W_a;\mathbb Z)
                 \longrightarrow H_{d-1}(A;\mathbb Z)
\]
is an isomorphism. This verifies the attaching map on the needed
homology group; a dimension count alone would not suffice.

Apply Mayer--Vietoris to the open cover $X(\mathbb C)=A\cup B$,
where $B=\coprod_a V_a$ and $A\cap B=\coprod_a W_a$.
Its degree-$d$ portion is
\[
 0\longrightarrow\bigoplus_a H_d(V_a;\mathbb Z)
    \longrightarrow H_d(X(\mathbb C);\mathbb Z)
    \longrightarrow\bigoplus_a H_{d-1}(W_a;\mathbb Z)
    \xrightarrow{\ \sim\ }H_{d-1}(A;\mathbb Z).
\]
Thus the local quadric sphere classes give a basis of
$H_d(X(\mathbb C);\mathbb Z)=\mathbb Z^N$.
All other reduced homology groups vanish. At the bottom of the
sequence, including when $d=2$, the map
$H_0(A\cap B)=\mathbb Z^N\to H_0(A)\oplus H_0(B)=\mathbb Z\oplus\mathbb Z^N$
is $e_a\mapsto(1,-e_a)$ and is injective. This accounts for the
potential degree-one boundary term and proves connectivity.

Simple connectivity also requires a check. Add the $V_a$ one at a
time to $A$. They are simply connected and intersect the preceding
union in the path-connected $W_a$, since the $B_a$ are disjoint.
Van Kampen consequently expresses the fundamental group as
$\pi_1(A)$ modulo the normal subgroup generated by the images of
all $\pi_1(W_a)$, with paths to a common basepoint understood.
Projection identifies this quotient with the quotient obtained by
adding the contractible $B_a$ to $U$ in the base. The latter union
is $\mathbb C^v$, whose fundamental group is trivial. Thus
$X(\mathbb C)$ is simply connected. This argument treats $v=1$
without replacing its nontrivial puncture loops by a homology-only
argument.

The spaces involved have CW homotopy type. One may see this directly
from this open cover: $U$ deformation retracts onto a large
rectangular box with the interiors of $N$ small rectangular boxes
removed, by pushing radially away from each puncture and inward
from infinity. Subdividing a finite rectangular grid triangulates
that region. The $V_a$ and $W_a$ have the sphere homotopy types just
proved. The double mapping cylinder for the two-set open cover is
homotopy equivalent to $X(\mathbb C)$; replacing its vertex spaces
by these CW models and its attaching maps by cellular
approximations yields a CW model. Homotopy invariance of this gluing
and the open-cover equivalence are Hatcher, Propositions4G.1--4G.2,
pp.458--459; cellular approximation is Theorem4.8, p.349.

Since the space is simply connected and its homology below degree
$d$ vanishes, successive applications of Hurewicz give
$\pi_i(X(\mathbb C))=0$ for $i<d$ and
$\pi_d(X(\mathbb C))\simeq H_d(X(\mathbb C);\mathbb Z)$.
Represent the displayed $N$ homology basis elements by based sphere
maps and form their wedge. This map from $\bigvee_N S^d$ induces an
isomorphism on all integral homology groups. Both spaces are simply
connected and have CW homotopy type, so the homology form of
Whitehead's theorem makes the map a homotopy equivalence.
The cohomology-ring assertion follows: a product of positive-degree
classes would lie in degree $4v$, whose group vanishes.

If $N=0$, then $U=\mathbb C^v$ and the first explicit deformation
retraction already makes $X(\mathbb C)$ contractible. There is in
fact the stronger arithmetic isomorphism claimed in the statement.
The integral Boolean quotient in NN-C027 is the product of $N$
copies of $\mathbb Z$. If $N=0$ it is the zero ring, so the ideal
$(f_1,\ldots,f_h)$ equals the whole of $R=\mathbb Z[x]$.
Choose a polynomial column $a=(a_1,\ldots,a_h)$ with $f\cdot a=1$.
For a point $(x,t,y)$ of $X$, put
\[
                         z=y+a(2t-t^2).
\]
Then $f\cdot z=t^2-t+2t-t^2=t$. Conversely, from arbitrary
$(x,z)\in\mathbb A^{v+h}_{\mathbb Z}$ set
\[
             t=f(x)\cdot z,\qquad y=z+a(x)(t^2-2t).
\]
These polynomial formulas are inverse over $\mathbb Z$, and
$f\cdot y=t^2-t$ proves that the second lands in $X$.
Thus $X\simeq\mathbb A^{v+h}_{\mathbb Z}$.
If $N>0$, its nonzero degree-$2v$ integral cohomology forbids
contractibility and forbids an isomorphism with complex affine
space. Base change then also forbids the integral isomorphism.
This proves all the equivalences.
\end{proof}

\paragraph{Exact topological source interfaces.}
Allen Hatcher, \emph{Algebraic Topology}, author text
\url{https://pi.math.cornell.edu/~hatcher/AT/AT.pdf}:
Theorem1.20, p.43 (van Kampen);
Theorem2.20, p.119 (excision);
Section2.2, ``Mayer--Vietoris Sequences'', pp.149--150 (unnumbered
exact sequence and its inclusion maps);
Theorem4.32, pp.366--367 (Hurewicz);
Corollary4.33, pp.367--368 (homology Whitehead);
Theorem4.8, p.349 and Propositions4G.1--4G.2, pp.458--459
(CW replacement and open-cover gluing).
The cited statements and the local proof passages used here were
read in the author PDF. The explicit affine-quadric retraction,
local analytic reduction, linking basis and arithmetic polynomial
isomorphism are proved above.

\begin{proposition}[NN-S016: linear-degree integral coordinates in the empty case]
For the Boolean and verifier equations of NN-C027, suppose $N=0$.
There are integral polynomials $a_i(x)$ satisfying
\[
                      \sum_{i=1}^h a_i(x)f_i(x)=1,
 \qquad \deg a_i\leq v.
\]
Their coefficients have $O(v)$ bits. In the polynomial isomorphism
of the preceding proposition, the forward coordinate polynomials have
degree at most $v+2$ and the inverse coordinate polynomials have
degree at most $v+6$, with coefficient bit lengths $O(v+\log(h+1))$.
These are degree and height bounds; a fully expanded coordinate
list can still have exponentially many monomials.
\end{proposition}

\begin{proof}
Every equation has degree at most two and coefficient $\ell^1$-norm
at most three. All $v$ Boolean equations $B_j=x_j^2-x_j$ are retained.
Every non-Boolean verifier equation takes a value in $\{0,1,-1\}$
on a bit assignment. Since there are no satisfying assignments,
for each $b\in\{0,1\}^v$ choose an index $i(b)$ with
$f_{i(b)}(b)\in\{1,-1\}$. Put
\[
 e_b(x)=\prod_{j:b_j=1}x_j\prod_{j:b_j=0}(1-x_j),
 \qquad
 a_i^{(0)}(x)=\sum_{b:i(b)=i}f_i(b)e_b(x).
\]
Then $\deg a_i^{(0)}\leq v$, and evaluation on the Boolean cube gives
$\sum_i a_i^{(0)}(b)f_i(b)=1$. In particular
\[
 P(x)=\sum_i a_i^{(0)}(x)f_i(x)-1
\]
has degree at most $v+2$ and has zero multilinear remainder modulo
the $B_j$. This last assertion holds over $\mathbb Z$: the
multilinear polynomial with prescribed values on the Boolean cube
is uniquely $\sum_b P(b)e_b$, so all its values being zero implies
that it is the zero polynomial.

Here is explicit integral division with controlled degree and
height. For an exponent vector $\alpha$, write
$\epsilon_j=\min(\alpha_j,1)$. Telescope
\[
 x^\alpha-x^\epsilon=
 \sum_j\left(\prod_{l<j}x_l^{\epsilon_l}\right)
       (x_j^{\alpha_j}-x_j^{\epsilon_j})
       \left(\prod_{l>j}x_l^{\alpha_l}\right).
\]
The middle factor is zero for $\alpha_j\leq1$, and otherwise it is
\[
        (x_j^2-x_j)\sum_{r=0}^{\alpha_j-2}x_j^r.
\]
Apply this formula to every monomial of $P$. Since its multilinear
remainder is zero, the result is
\[
                         P=\sum_j q_j B_j,
 \qquad \deg q_j\leq v.
\]
Subtract $q_j$ from the multiplier of the corresponding retained
Boolean equation. The resulting $a_i$ have the asserted degree
and satisfy the ordinary polynomial identity $\sum_i a_i f_i=1$.
This is an identity in the unrestricted integral polynomial ring,
not just in its Boolean quotient.

For a polynomial $Q$, write $\|Q\|_1$ for the sum of the absolute
values of its coefficients. Since $\|e_b\|_1=2^{|\{j:b_j=0\}|}$,
\[
 \sum_i\|a_i^{(0)}\|_1\leq\sum_b\|e_b\|_1=3^v,
 \qquad \|P\|_1\leq3^{v+1}+1.
\]
The division formula introduces at most $|\alpha|\leq v+2$
quotient monomials per input monomial, counted with multiplicity.
Hence
\[
 \sum_j\|q_j\|_1\leq(v+2)(3^{v+1}+1),\qquad
 \|a_i\|_1\leq B_v:=3^v+(v+2)(3^{v+1}+1).
\]
Thus every integer coefficient has $O(v)$ bits.

For the forward isomorphism $z_i=y_i+a_i(2t-t^2)$, the degree is
at most $v+2$ and its coefficient norm is at most $1+3B_v$.
For the inverse, $t=\sum_i f_i(x)z_i$ has degree at most three and
coefficient norm at most $3h$. Therefore
\[
 y_i=z_i+a_i\left(\left(\sum_j f_jz_j\right)^2
                           -2\sum_j f_jz_j\right)
\]
has degree at most $v+6$ and coefficient norm at most
$1+B_v(9h^2+6h)$. These give the claimed bit bounds.
A degree-$v$ polynomial in $v$ variables can have
$\binom{2v}{v}\leq4^v$ monomials; neither the degree nor the
coefficient-height estimate bounds their count polynomially.
\end{proof}

\paragraph{Proof-system and geometric interpretation.}
The multipliers form an integral Nullstellensatz refutation of the
quadratic verifier equations with their Boolean axioms. They also
supply low-degree polynomial coordinates for the same smooth cubic
model's arithmetic affine-space isomorphism. This refines the
Boolean-quotient identity of NN-C009 by explicitly constructing the
ordinary ideal multipliers in this special empty Boolean-scheme
case. No polynomial-size expanded certificate, short arithmetic
circuit, or efficient isomorphism-finding algorithm follows from
this degree calculation alone. NN-C029 separately analyzes what E's
symbolic partition routine supplies for circuit descriptions.

\paragraph{Effectivity and complexity scope.}
On the encoded family of NN-C027, $N$ is the satisfying-assignment
count. The rank of the single positive-degree integral cohomology
group, and the number of spheres in this homotopy type, are therefore
counting-complete quantities; their parity is parity-counting-complete.
Complex contractibility, isomorphism to complex affine space, and
isomorphism to integral affine space are each coNP-complete on this
encoded family, because each is equivalent to unsatisfiability.
Their complementary questions are NP-complete. Every scheme under
discussion is affine; the distinction is whether it is isomorphic
to affine space. The polynomial identity $f\cdot a=1$ and the
resulting isomorphism need not have a polynomial-size description,
so this construction supplies no short coNP certificate. E evaluates
these decisions through the original Boolean formula, while A and M
are not needed for the homotopy or isomorphism statements.
No motivic or $\mathbb A^1$-homotopy equivalence is asserted by this
complex-topological proof.


\paragraph{Every field fiber, and stabilization.}
The integral isomorphism when $N=0$ gives an affine-space isomorphism
after every base change. If $N>0$, equation~\eqref{eq:S016-cycle} gives
nonzero $H^{2v}(X_{\bar k},\mathbb Q_\ell)$ over every field. This
forbids an affine-space isomorphism over that field or over its algebraic
closure. Multiplication by any affine space preserves this cohomology,
so it also forbids stable isomorphism to affine space. Thus the same
unsatisfiability equivalence holds for affine-space recognition over any
fixed field, and for stable affine-space recognition on the encoded family.

\paragraph{Complexity and GCT consequences.}
The algebraic construction has polynomial description length and retains
the original satisfying-assignment count $N$. The explicitly identified
Chow rank, reduced Betti number and sphere count are exactly $N$.
Counting accepting witnesses proves membership in \#P on the specified
encoding, and the construction is parsimonious, proving \#P-hardness.
Equation~\eqref{eq:S016-wedge} makes noncontractibility equivalent to
$N>0$ and contractibility equivalent to $N=0$. These give the NP and
coNP completeness assertions directly, without a general algorithm
for recognizing contractible varieties.

E computes $N$, and hence the compact homotopy description $(2v,N)$,
within its stated original-variable counting bound for the source CNF.
This does not output $N$ individual spheres or cycle generators in that
time when $N$ is exponentially large. For arbitrary verification circuits,
the original circuit-and-witness size replaces the CNF variable parameter
exactly as in NN-C028.

The example strengthens the geometric limit on what purity alone can
compute: the varieties are explicitly rational and simply connected,
all their cohomology is generated by displayed algebraic cycle types,
and the cycle-class map is an isomorphism. The rank still recovers \#SAT.
No unsolved assertion that Tate classes are algebraic is needed in this
family. A polynomial procedure extracting that rank would give
$\mathrm{FP}=\#\mathrm P$ and the familiar Toda collapse; an efficient
noncontractibility test would give $\mathrm P=\mathrm{NP}$. These remain
conditional algorithmic implications, with no unconditional lower bound.
Under A, the zeta factors from NN-C027 have their prescribed shifted
critical lines, while $N$ remains their lower multiplicity. M is unused.
The family supplies an actual cycle realization but does not identify
these cycles with a canonical GCT basis or the multiplicities of the
permanent and determinant orbit closures.

\paragraph{Geometric theorem and Grothendieck dictionary.}
The morphism $X_\Phi\to\operatorname{Spec}\mathbb Z$ is a smooth
geometrically integral rational sparse cubic model. For each field fiber,
the cycle classes of its two affine components above each Boolean section
satisfy exactly one relation, $[D_{a,0}]+[D_{a,1}]=0$, and no others.
Its integral Chow ring is the square-zero extension of $\mathbb Z$ by
the free rank-$N$ group in codimension $v$, and its rational $\ell$-adic
cycle map is an isomorphism. The positive-codimension effective class
monoid equals the whole Chow group; it is not an ordered representation
semiring with a pointed cone. On the complex fiber the corresponding
integral homotopy type is the wedge of $N$ even spheres. The proof is
the localization-and-purity calculation and the explicit local
quadric argument above, rather than an inference of topology from
point counts alone.

This places the statement in Grothendieck's cohomology with supports and
Poincar\'e duality, SGA~4 XVI~3.7/3.10(a), XVIII~3.2.5/3.2.6.2, and
the cycle class of SGA~4\,$\frac12$, \emph{Cycle}~2.2.2/2.3.2.
The primary SGA~4 re-edition is
\url{https://www.normalesup.org/~forgogozo/SGA4/16/16.pdf}, pp.~11--13,
and \url{https://www.normalesup.org/~forgogozo/SGA4/18/18.pdf},
pp.~77--78 (original pp.~585--586 for duality).
In particular the purity twist is $(-v)$; the minus sign is explicit
in the primary re-edition's Remark~3.10(a), p.~13.
The localization dictionary for cycles is Stacks Tags~02RX, 02TS and~02TY.
The completeness assertions refer to succinct input equations and basis
rank extraction, while the individual realized cycle classes have the
explicit geometric descriptions just given.


\section{NN-S017: the counting cubic is an integral wedge of motivic spheres}

\paragraph{Statement and precise scope.}
Retain the actual integral verifier family of NN-C027 and NN-S016,
with $v,h\ge1$:
\[
 X_\Phi=\{t(t-1)=\textstyle\sum_{i=1}^h y_i f_i(x)\},\qquad
 Y_\Phi=V(f_1,\ldots,f_h)=\coprod_{a\in\mathcal S}\operatorname{Spec}\mathbb Z,
 \qquad N=|\mathcal S|.
\]
All $v$ Boolean equations $x_j^2-x_j$ belong to the list. Point
$X_\Phi$ by $x=0,t=0,y=0$; this point exists even when the zero
assignment is not a satisfying assignment. In the pointed unstable
Morel--Voevodsky homotopy category over $\mathbb Z$ there is an isomorphism
\begin{equation}\label{eq:S017-motivic-wedge}
 (X_\Phi,0)\simeq_{\mathbb A^1}
       \bigvee_{a\in\mathcal S}(\mathbb P^1,\infty)^{\wedge v}.
\end{equation}
The empty wedge means the pointed base scheme. The same statement holds
after every base change, with the same assignment index set and normal
framing. In particular this identifies the unstable motivic homotopy
type, not only its realizations, its Grothendieck class, or its ordinary
cohomology.

For every perfect field $k$, its integral geometric motive consequently is
\begin{equation}\label{eq:S017-motive}
 \mathrm M(X_{\Phi,k})\simeq
       \mathbb Z\oplus\mathbb Z(v)[2v]^{\oplus N}.
\end{equation}
For $q,n\ge0$, the entire higher Chow groups are
\begin{equation}\label{eq:S017-higher-chow}
 \operatorname{CH}^{q}(X_{\Phi,k},n)\simeq
 \operatorname{CH}^{q}(k,n)\oplus
       \operatorname{CH}^{q-v}(k,n)^{\oplus N},
\end{equation}
where negative codimension means zero. More generally, for every represented
bigraded motivic cohomology theory $E$ over the chosen base,
\[
 E^{p,q}(X_\Phi)\simeq E^{p,q}(\operatorname{Spec}\mathbb Z)
       \oplus E^{p-2v,q-v}(\operatorname{Spec}\mathbb Z)^{\oplus N},
\]
with the corresponding base changed formula. These are isomorphisms of
groups or modules; no multiplicative assertion is hidden in this formula.

On the specified verifier encoding, the number of nontrivial motivic
sphere summands, equivalently the displayed Tate multiplicity, is
\#P-complete. Recognizing $\mathbb A^1$-contractibility, triviality of the
reduced motive, or the absence of a nontrivial sphere is coNP-complete.
The motive statement is asserted over a fixed perfect field; the
contractibility statement is valid over $\mathbb Z$ and over any fixed
field. These are complexity claims for this effectively constructed
family, not decision bounds for arbitrary varieties or motives.

\paragraph{Imported theorem and attribution.}
The geometric mechanism is classical. We use Asok--Dubouloz--\O stv\ae r,
\emph{Geometric models for algebraic suspensions}, arXiv:2112.08241v2,
Theorem~3.1.1, printed p.~8, and Remark~3.1.3, p.~9
(published in IMRN 2023, pp.~17788--17821).
For a smooth map to a contractible smooth base with a section, an
$\mathbb A^1$-equivalence off a smooth closed subscheme identifies the
total space with the cofiber of the induced map between the two normal
Thom spaces. We verify every hypothesis and this map below. Their
Construction~4.2.1 and Remark~4.3.2 already cover the general geometric
form; the new conclusion here is its precise application to the
parsimonious counting family and its complexity interpretation.
The proof of Theorem~3.1.1 was read in full. Its homotopy-purity input
is Theorem~2.3.1 there, extending Morel--Voevodsky,
\emph{$\mathbb A^1$-homotopy theory of schemes}, Section~3,
Theorem~3.23. This is motivic homotopy purity, distinct from an
inference from finite-field point counts.

\paragraph{Verification of the smooth map and normal framing.}
Put $B=\operatorname{Spec}\mathbb Z$, $W=\mathbb A^v_B$, and let
$\pi:X_\Phi\to W$ be projection to $x$. It admits the section
$s(x)=(x,0,0)$. For the relative equation
$F=t^2-t-\sum y_i f_i$, the partial derivatives in the fiber
coordinates are $2t-1,-f_1,\ldots,-f_h$. In the coordinate ring of $X$,
\[
                   (2t-1)^2-4\sum_i y_i f_i=1.
\]
Thus these derivatives generate the unit ideal. The usual smooth
hypersurface Jacobian criterion proves that $\pi$ is smooth of relative
dimension $h$ over the entire integral base, including characteristic
two. The same criterion is SGA~1, Expos\'e~II, Theorem~4.10(iii) and
Corollary~4.12, already checked for NN-C027.

Let $U=W\setminus Y_\Phi$. On $D(f_i)\subset U$ the equation solves
for $y_i$, leaving the independent fiber coordinates $t$ and the other
$h-1$ variables. Hence $\pi^{-1}(U)\to U$ is a Zariski locally trivial
$\mathbb A^h$-bundle and therefore an $\mathbb A^1$-weak equivalence.
The closed scheme $Y_\Phi$ is smooth over $B$, being a disjoint union of
integral sections. Its embedding in $W$ is a regular immersion of
codimension $v$.

The classes of $b_j=x_j^2-x_j$ frame its conormal bundle. Indeed, on
the component indexed by $a\in\{0,1\}^v$, their differentials are
\[
                  db_j=(2a_j-1)\,dx_j.
\]
Each coefficient is an integral unit. The conormal sequence for the
smooth closed pair therefore identifies these $v$ classes with a basis
of $\Omega^1_{W/B}|_Y$. Equivalently, on a neighborhood of each component
the ideal is generated by $x_j-a_j$ and $b_j=(x_j-a_j)(x_j+a_j-1)$,
whose second factors restrict to units. This also directly verifies
regularity and shows that the framing commutes with every base change.
No assertion that the whole ideal of an arbitrary finite union is
globally generated by these $v$ elements is used.

The fiber over $Y$ is exactly
\[
 X_Y=D_0\amalg D_1,\qquad D_e=\{x\in Y,t=e\}\simeq Y\times\mathbb A^h,
\]
because $t$ and $t-1$ are comaximal over $\mathbb Z$. Flat base change
along $\pi$ identifies its conormal bundle with the pullback of the
framed conormal bundle of $Y\subset W$. Thus both normal bundles have
compatible trivializations of rank $v$. The section $s|_Y$ lands in
$D_0$ as its zero section in the $y$ coordinates.

\paragraph{Computing the cofiber.}
The cited theorem gives the pointed cofiber sequence
\[
 \operatorname{Th}(N_{Y/W})\longrightarrow
 \operatorname{Th}(N_{X_Y/X})\longrightarrow (X,s(0)).
\]
Using the compatible frames, the first term is
$(\mathbb P^1)^{\wedge v}\wedge Y_+$ and the second is
$(\mathbb P^1)^{\wedge v}\wedge(D_0\amalg D_1)_+$.
Contraction of each $\mathbb A^h$ factor identifies this map with the
inclusion
\[
 \bigvee_{a\in\mathcal S}(\mathbb P^1)^{\wedge v}
 \longrightarrow
 \left(\bigvee_{a\in\mathcal S}(\mathbb P^1)^{\wedge v}\right)
 \vee\left(\bigvee_{a\in\mathcal S}(\mathbb P^1)^{\wedge v}\right)
\]
into its first summand. Its cofiber is the second summand. This proves
\eqref{eq:S017-motivic-wedge}, including $N=0$. The computation works
unchanged over every base ring, so no preservation of an unspecified
isomorphism under arbitrary derived base change is required.

\paragraph{Motives, higher cycles, and complexity.}
Applying the reduced motive functor sends a pointed wedge to a direct
sum and sends $(\mathbb P^1,\infty)^{\wedge v}$ to
$\mathbb Z(v)[2v]$. The chosen rational point splits off the unreduced
unit motive. This proves \eqref{eq:S017-motive}. Exact classical
interfaces are Mazza--Voevodsky--Weibel, \emph{Lecture notes on motivic
cohomology}, Example~13.17, Theorem~15.2 and Corollary~15.5 for the
Tate sphere and projective line, and Theorem~19.1 (printed p.~159) for
the higher-Chow comparison. Tate cancellation over perfect fields is
Remark~16.26, with the theorem credited there to Voevodsky's 2002
\emph{Cancellation theorem}. The representable cohomology formula is
simply the image of the finite wedge under the represented functor;
applying motivic cohomology and that comparison gives
\eqref{eq:S017-higher-chow}.

For $N>0$, the already established group
$\operatorname{CH}^v(X_{\Phi,k})=\mathbb Z^N$ proves that the reduced
motive is nonzero over every perfect field. Thus it also rules out
$\mathbb A^1$-contractibility. For an arbitrary field, base change to
its algebraic closure gives the same obstruction; over $\mathbb Z$,
base change to $\mathbb C$ suffices. For $N=0$ the wedge is empty, and
NN-S016 supplies in addition the actual integral affine-space
isomorphism. The source verifier has exactly $N=\#\operatorname{SAT}(\Phi)$
solutions with unique wire extension, so exact extraction of the
displayed multiplicity is parsimoniously \#P-complete, and its
vanishing is exactly UNSAT. These equivalences establish the asserted
coNP membership and hardness on this family. They make no claim that
general motivic isomorphisms have polynomial certificates.

\paragraph{Grothendieck dictionary and implications.}
The ingredients are smooth integral schemes, a framed regular closed
immersion, its normal cone, and the deformation-to-the-normal-cone
purity construction. The normal-cone/intersection-theory interface is
the same geometric framework as SGA~6, Expos\'e~VII, regular immersions;
the exact homotopy theorem used is the later theorem cited above, not
a theorem attributed to SGA~6. The resulting Tate decomposition is
stronger than the SGA~4/SGA~5 realization calculation in NN-C027.
It supplies actual integral motives after passage to a perfect field,
without a Tate conjecture or a deduction of motives from zeta functions.

Existence of a pure Tate realization, rationality, algebraicity of all
classes, and now an explicit unstable motivic sphere type do not supply
an efficient multiplicity extractor. E evaluates this multiplicity at
its source counting bound; it does not make it polynomial-time.
Neither A nor M is used. No canonical GCT multiplicity identification,
polynomial algorithm for SAT, or separation of VP and VNP follows.



\section{NN-S018: a counting-hard integral gamma filtration on a torsion-free Grothendieck ring}

\paragraph{Statement.}
For the same integral cubic $X_\Phi$ and $v\ge1$ of NN-S017, write
$N=\#\operatorname{SAT}(\Phi)$ and $I=\ker(K_0(X_\Phi)\to K_0(\mathbb Z))$,
where the augmentation is pullback along the specified point. Then
\begin{equation}\label{eq:S018-ring}
 K_0(X_\Phi)=\mathbb Z\oplus I,\qquad I\simeq\mathbb Z^N,\qquad I^2=0.
\end{equation}
The same description holds after base change to any field. In compatible
integral sphere bases, every Adams operation satisfies
\[
                  \psi^m(z)=m^v z\qquad (z\in I, m\ge1).
\]
The complete exterior-power and gamma operations on the augmentation
ideal are
\begin{align}
 \lambda_t(z)&=1+z\sum_{m\ge1}(-1)^{m-1}m^{v-1}t^m,
                                                    \label{eq:S018-lambda}\\
 \gamma_t(z)&=1+z\sum_{j=1}^{v}(-1)^{j-1}(j-1)!
                       \left\{\begin{matrix}v\\j\end{matrix}\right\}t^j.
                                                    \label{eq:S018-gamma}
\end{align}
The braces denote Stirling numbers of the second kind. On an arbitrary
element $a+z$, $a\in\mathbb Z$, the exterior-power series is
$(1+t)^a\lambda_t(z)$, so these formulas specify the entire augmented
$\lambda$-ring.

Define
\[
 d_i(v)=\gcd_{i\le j\le v}\left((j-1)!
                    \left\{\begin{matrix}v\\j\end{matrix}\right\}\right)
           \quad(1\le i\le v),\qquad d_{v+1}(v)=0.
\]
For the integral gamma filtration of SGA~6, Expos\'e~V, Section~3.10,
\begin{equation}\label{eq:S018-filtration}
 F_\gamma^0K_0=K_0,\qquad F_\gamma^iK_0=d_i(v)I\ (1\le i\le v),\qquad
 F_\gamma^{v+1}K_0=0.
\end{equation}
Thus for $1\le i<v$ its associated graded groups are
\[
 \operatorname{gr}_\gamma^iK_0(X_\Phi)
       \simeq\left(\mathbb Z/(d_{i+1}(v)/d_i(v))\right)^N,
 \qquad
 \operatorname{gr}_\gamma^vK_0(X_\Phi)=(v-1)!I\simeq\mathbb Z^N.
\]
In particular, if $v\ge3$ is odd,
\begin{equation}\label{eq:S018-torsion}
             \operatorname{gr}_\gamma^{v-1}K_0(X_\Phi)
                        \simeq (\mathbb Z/2)^N.
\end{equation}
The original verifier can be padded parsimoniously so that $v$ is odd
and at least three. On that family the dimension over $\mathbb F_2$ of
this finite graded group is \#P-complete, its nonvanishing is
NP-complete, and its vanishing is coNP-complete. The ring $K_0$ is nevertheless torsion-free; over the perfect
fields covered by NN-S017, every ordinary integral Chow group
is torsion-free as well.
The torsion is in the associated graded of this filtration.

For completeness the higher groups have the additive description
\[
                         K_n(X_\Phi)\simeq K_n(\mathbb Z)^{\oplus(N+1)}
                         \quad(n\ge0),
\]
and likewise with a field in place of $\mathbb Z$. This last formula
does not assert that the higher groups of the base are torsion-free.

\paragraph{Classical interfaces.}
We use the unstable motivic equivalence NN-S017. Representability and
multiplicative Bott periodicity for algebraic $K$-theory on a regular
Noetherian base are imported in their precise standard forms:
Panin--Pimenov--R\"ondigs, \emph{On Voevodsky's algebraic $K$-theory
spectrum BGL}, April~17, 2007 version, Lemma~1.2.6,
Corollaries~1.2.7--1.2.8, and Lemmas~1.3.2--1.3.3,
Proposition~1.3.5, printed pp.~6--8. The source copy is
\url{https://www.math.uni-bielefeld.de/lag/man/251.pdf}.
The regular-base representability statement is also recorded in
Schlichting--Tripathi, \emph{Geometric models for higher
Grothendieck--Witt groups in $\mathbb A^1$-homotopy theory},
introduction, equation~(1.1). We use ordinary algebraic $K$-theory,
whose statement there has no hypothesis that two be invertible;
the later hermitian theorem has a different scope.

The projective-bundle computation is SGA~6, Expos\'e~VI,
Theorem~1.1. Exterior powers on $K_0$ are SGA~6, Expos\'e~V,
Example~2.2(b) and Lemma~2.2.1; the change of variables defining
$\gamma_t$ is Proposition~3.2(ii), and the filtration is (3.10.1).
We apply these classical operations to the new counting family;
no new general theorem about $\lambda$-rings is asserted.

\paragraph{The ring, Bott generator, and Adams action.}
Every motivic sphere $(\mathbb P^1)^{\wedge v}$ is a simplicial
suspension, since $(\mathbb P^1,\infty)\simeq S_s^1\wedge\mathbb G_m$.
A pointed suspension has null reduced diagonal: its pinch map to the
wedge of two copies has both projections homotopic to the identity,
and the composite of that wedge with the smash product is zero.
The same holds for a wedge of suspensions. Products of two reduced
classes in any represented multiplicative cohomology theory thus
vanish. Applying $K$-theory to NN-S017 proves $I^2=0$.
Bott periodicity identifies the reduced $K_n$ of each sphere with
$K_n$ of the base. The rational point splits off the base summand.
This proves \eqref{eq:S018-ring} and the higher-group formula,
since $K_0(\mathbb Z)=K_0(k)=\mathbb Z$ for a field $k$.

Here is also an integral way to compute the operation on a sphere.
On $(\mathbb P^1)^v$ let $u_j=[\mathcal O(1)_j]-1$.
The projective-bundle theorem gives
\[
 K_0((\mathbb P^1)^v)=
       \mathbb Z[u_1,\ldots,u_v]/(u_1^2,\ldots,u_v^2).
\]
The quotient map from the product to the pointed smash product
pulls the reduced Bott generator back to $u_1\cdots u_v$, up to
a choice of sign. Its pullback on the reduced rank-one group is
injective. Adams operations on a line bundle have
$\psi^m([L])=[L]^m$, and are ring endomorphisms; these are the
Newton identities for the exterior-power operations. Therefore
\[
 \psi^m(u_1\cdots u_v)=
       \prod_{j=1}^v((1+u_j)^m-1)=m^v u_1\cdots u_v.
\]
This proves the displayed Adams formula on each sphere and hence on $I$.

The naturality used here is precisely the unstable operation theorem
of Riou, \emph{Algebraic $K$-theory, $\mathbb A^1$-homotopy and
Riemann--Roch theorems}, arXiv:0907.2710v2, September~9, 2009,
Theorems~1.1.2 and 2.3.1, Proposition~2.3.2, and the opening of
Section~2.4, printed pp.~4 and 10--11. Theorem~2.3.1 gives the
unique special $\lambda$-ring structure on $\mathbb Z\times\mathrm{Gr}$
in the unstable motivic category inducing the usual exterior
powers on every smooth scheme over a regular base. Section~2.4
extends it by representability to every motivic space; the
proposition makes it compatible with regular base change.
Consequently the sphere equivalence and the quotient map above
commute with the usual exterior-power and Adams operations.
This is an unstable integral statement; no periodic stable Adams
map requiring inversion of $m$ is used. The primary copy is
\url{https://www.imo.universite-paris-saclay.fr/~joel.riou/doc/operations.pdf}.

\paragraph{Exact integral formulas and the filtration.}
In the torsion-free ring $K_0\otimes\mathbb Q$, the Newton identity is
\[
 \log\lambda_t(z)=\sum_{m\ge1}\frac{(-1)^{m-1}}m\psi^m(z)t^m.
\]
Since every positive coefficient is in the square-zero ideal,
the logarithm is just $\lambda_t(z)-1$. Substituting the Adams
formula gives \eqref{eq:S018-lambda}. All coefficients obtained
are integral, so torsion-freeness proves the identity already
over $\mathbb Z$.

SGA~6, V, Proposition~3.2(ii), gives
$\gamma_t(z)=\lambda_{t/(1-t)}(z)$. Write its scalar factor as
$\gamma_t(z)=1+P_v(t)z$. The series formula shows
\[
 P_1(t)=t,\qquad P_{v+1}(t)=t(1-t)P_v'(t).
\]
The Stirling recurrence
$\left\{\begin{smallmatrix}v+1\\j\end{smallmatrix}\right\}
=j\left\{\begin{smallmatrix}v\\j\end{smallmatrix}\right\}
+\left\{\begin{smallmatrix}v\\j-1\end{smallmatrix}\right\}$
then proves \eqref{eq:S018-gamma} by induction. This proves the
finite polynomial formula, not merely a formal asymptotic expansion.

By definition $F_\gamma^i$ is generated by products
$\gamma^{j_1}(z_1)\cdots\gamma^{j_s}(z_s)$, where every $z_b\in I$
and $\sum j_b\ge i$. Any product with two positive-degree factors
vanishes because $I^2=0$. The remaining generators are exactly
the scalar multiples
\[
 (j-1)!\left\{\begin{matrix}v\\j\end{matrix}\right\}I,
                         \qquad i\le j\le v.
\]
Their sum is the gcd multiple $d_i(v)I$. Multiplication by an
arbitrary element of $K_0=\mathbb Z\oplus I$ adds no further
generators, again because $I^2=0$. This proves
\eqref{eq:S018-filtration}. Taking the indicated quotients gives
all associated graded groups; $d_i(v)$ divides $d_{i+1}(v)$.

When $v$ is odd and at least three,
$\left\{\begin{smallmatrix}v\\v-1\end{smallmatrix}\right\}
=\binom v2$ and
$\left\{\begin{smallmatrix}v\\v\end{smallmatrix}\right\}=1$,
so
\[
 d_{v-1}(v)=\gcd\left(\frac{v(v-1)!}{2},(v-1)!\right)
                   =\frac{(v-1)!}{2},\qquad d_v(v)=(v-1)!.
\]
This proves \eqref{eq:S018-torsion}. For example at $v=3$,
$\gamma_t(z)=1+(t-3t^2+2t^3)z$,
$F_\gamma^1=F_\gamma^2=I$, $F_\gamma^3=2I$, and
$\operatorname{gr}_\gamma^2=I/2I$.

\paragraph{Counting consequence and geometric meaning.}
Append constant-zero verifier variables, with their Boolean equation
and the equation setting them to zero, until $v\ge3$ is odd.
At most three variables are needed; all old assignments extend
uniquely, all equations retain bounded degree and locality, and
the integral verifier scheme remains the same disjoint union of
$N$ sections in its enlarged affine space. Thus the new cubic
still satisfies every hypothesis of NN-C027 and NN-S017.
Formula~\eqref{eq:S018-torsion} now identifies the indicated
$\mathbb F_2$-dimension with the original exact satisfying count.
This proves all complexity claims, including membership on the
specified syntactic family. E supplies its source counting time,
with the padding contributing only polynomial overhead.

In Grothendieck's language this is an explicit augmented
representation-style $\lambda$-ring of vector bundles, with every
Schur operation determined by its exterior powers and with its
integral gamma filtration computed exactly. Its rational graded
pieces have only gamma degree $v$, corresponding to Tate twist
$v$, while lower
integral graded pieces can have finite torsion of counting-hard
rank. The underlying $K_0$ stays free, as do the ordinary Chow groups
over the perfect fields just specified.
The calculation cautions against replacing the integral gamma
filtration by its rational realization when extracting parity
or multiplicity information. It does not identify these groups
with GCT's canonical multiplicities, prove PH1, or settle P versus
NP or VP versus VNP. Neither A nor M is used.



\section{NN-S019: exponential positive arithmetic circuit cost for a local counting character}

\paragraph{Statement.}
Let $G=(V,E_G)$ be a simple graph on $n\ge2$ vertices of
maximum degree at most a fixed $d\ge1$. Suppose that, for every
$S\subseteq V$ with $|S|\le n/2$, its edge boundary satisfies
$|\partial S|\ge h|S|$, where $h>0$ is fixed.
Give each vertex two independent weights $y_{v,0},y_{v,1}$,
and each edge $e=uv$ four independent weights $z_{e,a,b}$.
Choose an orientation of each edge only to order these two bits.
Define the integral polynomial
\[
 P_G(y,z)=\sum_{x\in\{0,1\}^V}
      \prod_{v\in V}y_{v,x_v}
      \prod_{e=uv\in E_G}z_{e,x_u,x_v}.
\]
It has $2^n$ distinct monomials, all with coefficient one,
$2n+4|E_G|=O_d(n)$ variables and degree
$M=n+|E_G|=O_d(n)$. Every binary monotone arithmetic circuit
over $\mathbb R_{\ge0}$ computing it as a formal polynomial
has at least the following number of internal addition or multiplication gates:
\begin{equation}\label{eq:S019-lower}
 2^{\kappa n}\quad\hbox{gates},\qquad
            \kappa=\frac{h}{3(d+1)(d+h)}.
\end{equation}
The circuit model permits arbitrary nonnegative real constants,
addition and multiplication, and sharing of intermediate results.

In contrast, NN-C034 constructs a constant-free integral arithmetic
circuit for the very same polynomial of size
\[
 \operatorname{poly}_d(n)
 2^{O_d(n(\log\log n)^2/\sqrt{\log n})}=2^{o(n)},
\]
and polynomial formal degree. Thus, for these explicitly specified
local counting polynomials on any constant-expansion family,
the positive versus integral arithmetic size ratio is at least
$2^{\kappa n/2}$ for all sufficiently large $n$.
This is a consequence of E combined with the elementary
positive-rectangle argument proved below. It is not a new
general lower-bound technique or a historical-priority claim
for monotone graph polynomials.

\paragraph{The positive circuit gives a small balanced cover.}
There are $M$ weight blocks: one block for each vertex and one
for each edge. Every output monomial contains exactly one
variable from every block. Remove zero gates and gates which
do not contribute to the output. At each remaining gate $g$,
write $P_g$ for its nonzero polynomial. Fix a path from $g$
to the output and a nonzero monomial from every sibling at a
multiplication on this path. Their product is a monomial $c$.
Positivity implies that $pc$ is an output monomial for every
monomial $p$ of $P_g$. Consequently all monomials at $g$ have
the same block multidegree, with each entry zero or one.
Write $A_g$ for its set of occupied blocks and $r_g=|A_g|$.
This reasoning does not assume that the circuit is a formula.
In particular, a positive-degree gate cannot occur twice in
a contributing parse tree, since its block degrees would
then exceed the output degrees.

At an addition gate its two nonzero inputs have the same
block set. At a multiplication gate their block sets are
disjoint and their union is the output block set of that gate.
For a fixed output monomial, follow a contributing parse tree
from degree $M$ downwards. At sums choose a child contributing
the monomial. At products, while a child has degree greater
than $2M/3$, follow such a child. At the first product where
neither child has that degree, the parent still has degree
greater than $2M/3$. Choose the larger child: its degree is
strictly greater than $M/3$ and at most $2M/3$. Since
$h>0$ and $n\ge2$ imply $M\ge3$, this child has degree at
least two and is an internal gate. Thus every output monomial
passes through an internal gate
in the set
\[
             \mathcal B=\{g\text{ internal}:M/3\le r_g\le2M/3\}.
\]

For a fixed $g\in\mathcal B$, sum the products of sibling
polynomials over all paths from $g$ to the output, obtaining
a nonnegative context polynomial $C_g$. This polynomial is
used only for the proof; no efficient construction is needed.
Every term $P_gC_g$ has support contained in the output
support, and $C_g$ occupies exactly the complementary blocks.
The path reasoning just given proves the support cover
\begin{equation}\label{eq:S019-cover}
 \operatorname{supp}(P_G)\subseteq
       \bigcup_{g\in\mathcal B}\operatorname{supp}(P_gC_g).
\end{equation}
There are at most as many products as circuit gates. We need
neither coefficient equality in this cover nor disjointness
between its products.

\paragraph{A balanced product fixes many original bits.}
Consider any nonzero product $AB$ with nonnegative coefficients,
occupying complementary block sets $\mathcal A,\mathcal D$ of
sizes between $M/3$ and $2M/3$, whose support is contained in
$\operatorname{supp}(P_G)$. Let $B_0$ be the vertices whose
incident factor blocks meet both sides of this partition.
Here the unary vertex block is an incident factor, as is every
edge block containing the vertex. Put $b=|B_0|$.

For every $v\in B_0$, all monomials of $A$ and $B$ must
encode the same fixed bit at $v$. Indeed fix a monomial in
$B$. Each monomial in $A$ must be compatible with it, because
their product has positive coefficient and is a valid output
monomial. Since both sides contain a block involving $v$,
this fixes its bit in every monomial of $A$. Fixing one such
monomial of $A$ then fixes it in every monomial of $B$.
Internal consistency of each side also follows from extension
by any nonzero monomial on the other side. Thus the entire
product contains at most $2^{n-b}$ output monomials.

Let $V_A$ and $V_B$ be the vertices all of whose incident
factors belong respectively to $\mathcal A$ and $\mathcal D$.
They partition $V\setminus B_0$. Every factor in
$\mathcal A$ touches $V_A\cup B_0$. Since a vertex touches
at most $d+1$ factors,
\[
 |V_A|\ge\frac{M}{3(d+1)}-b,
 \qquad |V_B|\ge\frac{M}{3(d+1)}-b.
\]
There is no graph edge from $V_A$ to $V_B$: its factor
would have to belong to both sides. Take the smaller of
$V_A,V_B$, of size $a\le n/2$. All edges in its boundary
end in $B_0$, so expansion and the degree bound give
$ha\le db$. Combining with
$a\ge M/(3(d+1))-b$ yields
\[
 b\ge\frac{hM}{3(d+1)(d+h)}\ge\kappa n.
\]
This remains valid if the lower bound for $a$ before
rearrangement is negative. Therefore each product in
\eqref{eq:S019-cover} contains at most $2^{n-\kappa n}$
of the $2^n$ output monomials. If the circuit has $s$ gates,
then $2^n\le s2^{n-\kappa n}$, proving
\eqref{eq:S019-lower}.

\paragraph{The integral upper bound and nonempty family.}
The displayed definition of $P_G$ is exactly a universal
partition polynomial with $n$ Boolean variables, $n$ unary
factors and $|E_G|$ binary factors. Its density is at most
$1+d/2$. Thus every hypothesis of NN-C034 holds without
additional verifier wires, variable copying or a change
of the parameter. Its circuit output is the integral
polynomial above, so the comparison is in characteristic
zero before any reduction. Since the upper exponent is
$o(n)$, the claimed exponential ratio follows directly.

The expander hypothesis supplies an infinite explicit
family. One precise classical source is Reingold--Vadhan--
Wigderson, \emph{Entropy waves, the zig-zag graph product,
and new constant-degree expanders}, August~1, 2001 version,
Theorem~3.3, printed p.~11, and Section~5, pp.~16--18. In the published \emph{Annals of Mathematics}~155 (2002), 157--187 version, these are Theorem~3.3 and its proof on pp.~170--171, with the explicit starting graph in Lemma~5.1 and Proposition~5.2 on pp.~178--179. The readable published version is also arXiv:math/0406038v1. Its constant-degree
graphs have a normalized spectral bound $\lambda<1$
independent of their size. To pass from its allowed
multigraphs to simple graphs, remove loops and merge
parallel edges. For a set of at most half the vertices,
the spectral inequality applied to its centered indicator
gives at least $d(1-\lambda)|S|/2$ crossing edges with
multiplicity. Each unordered edge has multiplicity at
most $d$, so the resulting simple graph has at least
$(1-\lambda)|S|/2$ crossing edges and maximum degree
at most $d$. This gives the required fixed positive $h$.
The original primary version is
\url{https://www.math.ias.edu/~avi/PUBLICATIONS/MYPAPERS/SALIL/EXPANDERS/ANNALS/RVW01.pdf}. The published version is
\url{https://arxiv.org/pdf/math/0406038}.

\paragraph{Representation and Grothendieck translation.}
Let $T$ be the split torus over $\mathbb Z$ with independent
characters indexed by all the weights $y,z$. On the finite
etale Boolean scheme
$Y=(\operatorname{Spec}\mathbb Z\sqcup\operatorname{Spec}\mathbb Z)^V$,
give the free rank-one summand at section $x$ the character
equal to its displayed monomial. Its pushforward to
$\operatorname{Spec}\mathbb Z$ is an actual free $T$-module
$\mathcal V_G$ of rank $2^n$, with character $P_G$.
The dictionary between split-torus representations and
character gradings is SGA~3, Expos\'e~I,
4.7.2--4.7.3; direct sum and tensor product induce the
addition and multiplication of the representation
Grothendieck ring, as in SGA~6, Expos\'e~IV,
(1.5.2) and Definition~2.2.

A circuit building this class solely by direct sums,
tensor products, and nonnegative integer multiplicities
from the specified one-dimensional characters is a
monotone arithmetic circuit. It therefore needs
$2^{\kappa n}$ gates. E and NN-C034 instead give a
$2^{o(n)}$ expression in the representation ring using
virtual differences and shared subexpressions. The
resulting class is still the actual effective class
$[\mathcal V_G]$. This proves an exponential cost for
removing subtraction from these precise representation
expressions, despite their bounded-locality counting origin.

The lower bound concerns this arithmetic model. It does
not cover formulas with arbitrary Schur-functor primitives,
lattice-point-counting oracles, coefficient extraction,
hashing operations, arbitrary dual-character generators, or
the entire finite pushforward as a unit-cost primitive. It does not count intermediate
split-semiring values equal to the supported zero: by
NN-S009, a value alone does not retain cancellation
provenance. It gives no lower bound for unrestricted
Boolean circuits and no P versus NP or VP versus VNP
resolution. In particular, the upper bound remains
subexponential, not polynomial. Neither A nor M is used.


\paragraph{A polynomial monotone Boolean circuit for the same support.}
The formal-polynomial lower bound above persists alongside a much smaller
Boolean support circuit.  For each unary or binary factor tuple, let its
weight be an input bit recording whether that tuple is allowed.  Write
$a=(a_{v,b},a_{e,b,c})$ for all these bits.  Then
\[
 \operatorname{Supp}_G(a)=
 \bigvee_{x\in\{0,1\}^V}
       \left(\bigwedge_v a_{v,x_v}\ \wedge\!
                         \bigwedge_{e=uv}a_{e,x_u,x_v}\right)
\]
has a binary monotone Boolean circuit with $O(n^3+|E_G|)$ gates,
constructible in that time from $G$.  In particular, on the bounded-degree
expander family it has size $O(n^3)$, although every monotone arithmetic
circuit for the full formal counting polynomial needs $2^{\kappa n}$
gates.  No use of E is needed for this Boolean upper bound.

Here is the full construction and proof.  First use input bits
$f_{v,b},f_{e,b,c}$ recording that a tuple is forbidden.  A forbidden
unary tuple $(v,b)$ supplies the clause $x_v\ne b$.  A forbidden edge
tuple $(e=uv,b,c)$ supplies the clause
\[
                       (x_u\ne b)\vee(x_v\ne c).
\]
Thus consistency of the allowed tuples is exactly satisfiability of
this two-CNF.  It is useful that the set of possible clauses is fixed
by $G$: the input bits indicate which clauses are present.

Make an implication graph whose $2n$ vertices are literals
$L_{v,b}$ meaning $x_v=b$, with complement
$\overline{L_{v,b}}=L_{v,1-b}$.  A unary forbidden bit supplies
the edge $L_{v,b}\to L_{v,1-b}$.  A binary forbidden bit supplies
the two edges
\[
 L_{u,b}\longrightarrow L_{v,1-c},\qquad
 L_{v,c}\longrightarrow L_{u,1-b}.
\]
Each adjacency bit is the disjunction of the forbidden-tuple bits
that supply that edge.  The total number of such occurrences is
$2n+8|E_G|$, so all adjacency bits require $O(n+|E_G|)$ binary
OR gates, allowing constants for absent edges.  Every edge
$L\to M$ has its contrapositive edge
$\overline M\to\overline L$.

Enumerate the $2n$ literals as $1,\ldots,2n$.  Compute reachability
using the following Boolean circuits:
\[
 R^{(0)}_{ij}=A_{ij}\vee[i=j],\qquad
 R^{(k)}_{ij}=R^{(k-1)}_{ij}\vee
             \bigl(R^{(k-1)}_{ik}\wedge R^{(k-1)}_{kj}\bigr)
                     \quad(1\le k\le2n).
\]
Induction on $k$ shows that $R^{(k)}_{ij}$ records whether there is
a directed path from $i$ to $j$ whose internal vertices are among
$1,\ldots,k$.  The induction uses a simple path when there is a
path, so that a visit to $k$ splits it into the two required parts.
Consequently $R^{(2n)}$ is the reachability relation.  This circuit
uses $O(n^3)$ binary AND and OR gates.  Its inputs and operations
are all monotone in the forbidden-tuple bits.

The two-CNF is unsatisfiable precisely when
\[
 U_G(f)=\bigvee_{v\in V}
 \left(R^{(2n)}_{L_{v,0},L_{v,1}}\wedge
       R^{(2n)}_{L_{v,1},L_{v,0}}\right)=1.
\]
This is exactly the ordinary two-CNF strong-component criterion of
Aspvall--Plass--Tarjan, Theorem~1, printed p.~122, cited precisely
below.  We import that classical theorem rather than reprove it.

The entire circuit for $U_G$ uses only AND, OR, constants and the
forbidden inputs.  Form its Boolean dual by interchanging AND with
OR, interchanging constants zero with one, and replacing each
forbidden input label by the corresponding allowed input label.
Induction gate by gate, using De Morgan's laws, shows that the
dual circuit computes
\[
                            1-U_G(1-a)=\operatorname{Supp}_G(a).
\]
This is an actual monotone circuit on the allowed bits; no negation
gates are introduced at its inputs or internally.  Dualization
preserves the number of gates.  This proves the claimed bound.

The classical source for the implication graph and its exact
strong-component criterion is Aspvall--Plass--Tarjan,
\emph{A linear-time algorithm for testing the truth of certain
quantified Boolean formulas}, \emph{Information Processing Letters}
8 (1979), 121--123: the graph is defined on printed p.~121, and
Theorem~1 with its complete proof is on printed p.~122.
The primary article is available at
\url{https://mathweb.ucsd.edu/~sbuss/CourseWeb/Math268_2007WS/2SAT.pdf}.
Only the ordinary existential two-CNF theorem is used here; the
quantified extension in its Theorem~2 is not used.  The cubic circuit
construction and monotone dualization are explicit above, so no
circuit bound is being imported from the article's
linear-time random-access-machine algorithm.

\paragraph{Geometric form of the support contrast.}
For each fixed allowed-tuple bit string $a$, let
\[
 B_n=\mathbb Z[x_v:v\in V]/(x_v^2-x_v:v\in V)
                   \simeq\prod_{x\in\{0,1\}^V}\mathbb Z.
\]
For a vertex or edge factor $f$, let $p_{f,a}\in B_n$ be its
Boolean interpolation polynomial, taking value $a_{f,x|I_f}$
on section $x$.  Explicitly the unary interpolation is
$a_{v,0}(1-x_v)+a_{v,1}x_v$, and the binary interpolation is
the corresponding sum of four degree-two Boolean indicator
monomials multiplied by their allowed bits.  Each $p_{f,a}$
is an idempotent.  Put $e_a=\prod_f p_{f,a}$ and
\[
                 Y_a=\operatorname{Spec}(B_n/(1-e_a)).
\]
This is the finite etale scheme over $\mathbb Z$ consisting of
exactly the allowed assignment sections.  Equivalently one may
use all equations $p_{f,a}=1$ together with the Boolean equations:
for idempotents the ideals $(1-e_a)$ and $(1-p_{f,a}:f)$ agree.
One inclusion follows by expanding $1-\prod_f p_{f,a}$ as a
telescoping sum; the other follows from
$(1-p_{f,a})e_a=0$.

For every field $K$, the scheme $Y_a\times\operatorname{Spec}K$
is nonempty if and only if $\operatorname{Supp}_G(a)=1$.
Thus the nonemptiness of this explicitly given finite etale
subscheme is computable by the displayed $O(n^3)$ monotone
Boolean circuit in its allowed-tuple parameters.  Its rank is
$P_G(a)$ as an ordinary nonnegative integer, while the full
torus character before specialization is the exponentially
costly positive arithmetic object described above.  The map
from positive formal polynomials to Boolean support is a
semiring map with idempotent addition.  It is not an evaluation
of invertible torus characters at zero in the representation
ring, and it does not retain multiplicities.  This distinction
is exactly why the arithmetic lower bound and the Boolean
upper bound coexist.


\paragraph{A strengthening with fixed edge constraints and linear Boolean support.}
Let $G=(V,E_G)$ belong to the bounded-degree expander family above,
and put $n=|V|$.  Increase the fixed degree bound to $d\ge3$
if necessary.  Retain only the two unary weights at each vertex:
\[
 J_G(z)=\sum_{I\subseteq V\,\text{independent}}
            \prod_{v\in I}z_{v,1}\prod_{v\notin I}z_{v,0}
                  \in\mathbb Z[z_{v,b}:v\in V, b\in\{0,1\}].
\]
This polynomial is homogeneous of degree $n$, and every monomial
contains exactly one variable from each vertex block
$\{z_{v,0},z_{v,1}\}$.  There is a constant $\gamma>0$, depending
only on the fixed expansion and degree bounds, such that every
binary positive arithmetic circuit over $\mathbb R_{\ge0}$ for
$J_G$ has at least $2^{\gamma n-2}$ internal gates for all
sufficiently large $n$.  Nevertheless NN-C034 supplies an integral
constant-free circuit of size and deterministic bit construction time
\[
 \operatorname{poly}_d(n)
 2^{O_d(n(\log\log n)^2/\sqrt{\log n})}=2^{o(n)},
\]
with polynomial formal degree.  The parameter is the original
number of vertices.  The exponential positive lower bound therefore
persists after every binary edge factor of the preceding universal
polynomial has been fixed to the independent-set constraint.
This assertion requires its own lower-bound argument; it does not
follow merely by specializing a polynomial with a lower bound.

For the same polynomial, Boolean evaluation of the coefficients and
monomials gives the allowed-weight support function
\[
 \operatorname{Supp}^{\rm ind}_G(a)
   =\bigvee_{I\,\text{independent}}
           \left(\bigwedge_{v\in I}a_{v,1}\wedge
                 \bigwedge_{v\notin I}a_{v,0}\right).
\]
It has the explicit monotone formula
\begin{equation}\label{eq:S019-independent-support}
 \operatorname{Supp}^{\rm ind}_G(a)=
       \bigwedge_{v\in V}(a_{v,0}\vee a_{v,1})
       \ \wedge\!
       \bigwedge_{uv\in E_G}(a_{u,0}\vee a_{v,0}).
\end{equation}
In particular it uses at most $2(n+|E_G|)-1$ binary AND and OR
gates.  This upper bound is unconditional and is linear on the
bounded-degree family.  It concerns availability of unary choices,
not the full formal polynomial or its integer evaluation.

\paragraph{The classical lower bound being imported.}
Bova--Capelli--Mengel--Slivovsky, \emph{A Strongly Exponential
Separation of DNNFs from CNF Formulas}, arXiv:1411.1995v3
(19 February 2015), define DNNFs and wire size in Section~2,
printed pp.~2--3, and define the graph CNF
\[
                  F_G(x)=\bigwedge_{uv\in E_G}(x_u\vee x_v)
\]
in equation~(2), printed p.~5.  Their Theorem~4, printed p.~10,
states that, for a $(c,d)$ vertex expander,
\[
 \operatorname{DNNF}(F_G)\ge
      2^{g(c,d)\operatorname{size}(F_G)-1},
 \qquad
 g(c,d)=\frac{c\log_2(1+2^{-d})}{6d^3}>0.
\]
Here size counts wires, and the graph CNF has $3|E_G|$ wires.
We import this theorem; its expander lower-bound proof is not
being presented as a new deduction.  The primary manuscript is
\url{https://arxiv.org/pdf/1411.1995v3}.

To relate the hypotheses, the preceding edge-boundary inequality
and maximum degree $d$ give
$|N(S)\setminus S|\ge(h/d)|S|$ for $|S|\le n/2$.
The graph is connected.  A graph of maximum degree at most two
cannot satisfy this fixed positive edge expansion for unbounded
$n$: a connected such graph is a path or cycle, and a consecutive
set of $\lfloor n/2\rfloor$ vertices has at most two boundary edges.
Thus its actual maximum degree $d'$ lies between $3$ and $d$ for
all sufficiently large instances.  Apply the imported theorem
with $c=h/d$ and $d'$.  Since
$g(h/d,d')\ge g(h/d,d)$ and
$3|E_G|\ge3(n-1)\ge n$ for $n\ge2$, we may take
\[
                 \gamma=\frac{h\log_2(1+2^{-d})}{6d^4}.
\]
The finite excluded range has no effect on the asymptotic statement.

\paragraph{The positive-arithmetic-to-DNNF transfer.}
Let a positive arithmetic circuit with $s$ internal binary gates
compute $J_G$ as a formal polynomial.  Delete zero polynomials
and gates not contributing to the output; this is a semantic
simplification used for a lower-bound proof, and no algorithmic
identity test is required.  For each remaining gate $g$, choose
a path to the output and one nonzero monomial from every sibling
at a multiplication on that path.  Their product is a fixed
context monomial $c$.  Positivity ensures that $pc$ is an output
monomial for every monomial $p$ at $g$.  This statement is valid
for a DAG: if another occurrence of a shared gate is in a sibling,
the ordinary distributive expansion permits its monomial to be
chosen independently in that occurrence.

Every output monomial has block multidegree $(1,\ldots,1)$.
Consequently every monomial at $g$ has the same block multidegree,
with entries zero or one.  Denote its block set by $A_g$.
At a productive multiplication gate the two input block sets
are disjoint, and their union is $A_g$.  At an addition gate
the two nonzero input block sets equal $A_g$.
Since coefficients are nonnegative and zero gates have been
removed, the vertex variables appearing anywhere below a gate
are exactly the union of the blocks in its monomials; no
syntactic input can disappear by cancellation.

Replace addition and multiplication by OR and AND, respectively,
each positive constant by true, and the two input rails by
\[
                         z_{v,1}\mapsto x_v,
                 \qquad z_{v,0}\mapsto\neg x_v.
\]
The disjoint block sets just proved show that every AND gate
has variable-disjoint input subcircuits.  The resulting circuit
is therefore a DNNF, with at most $2s$ wires.  On a Boolean
assignment $x$, a monomial of $J_G$ survives this substitution
exactly when its selected vertex set is $\{v:x_v=1\}$.
Thus the DNNF computes the independent-set predicate
$\bigwedge_{uv}(\neg x_u\vee\neg x_v)$.
Interchanging $x_v$ and $\neg x_v$ at all leaves produces a
DNNF for $F_G$ with the same number of wires; decomposability
is unchanged.  The imported theorem now gives
\[
                       2s\ge2^{\gamma n-1},
             \qquad s\ge2^{\gamma n-2},
\]
as asserted.  Arbitrary nonnegative real constants do not
weaken this argument: only zero versus positivity and formal
monomial support were used.

\paragraph{The integral and Boolean upper bounds.}
There are $n$ Boolean variables, one unary factor
\[
                (1-x_v)z_{v,0}+x_vz_{v,1}
\]
for each vertex, and one fixed binary factor $1-x_ux_v$ for
each edge.  Their partition polynomial is exactly $J_G$.
There are at most $(1+d/2)n$ factors of arity at most two,
so NN-C034 applies at the original parameter $n$.  Its formal
integral identity remains valid after this fixed specialization;
no verification variables or growing-scope factors are introduced.
The integral exponent is $o(n)$, whereas the positive lower
exponent is linear, giving an exponential ratio for large $n$.

For equation~\eqref{eq:S019-independent-support}, necessity of
each unary clause follows because every selected assignment
must use an allowed value at that vertex.  If both endpoints
of an edge have their zero values forbidden, both are forced
to one, so its edge clause is necessary too.  Conversely,
suppose all displayed clauses hold.  Choose
$I=\{v:a_{v,0}=0\}$.  The unary clauses ensure that one is
allowed on $I$; zero is allowed everywhere outside $I$ by
definition.  Every edge clause forbids both endpoints from
lying in $I$, so $I$ is independent.  This gives a permissible
monomial and proves the formula.  Its $n+|E_G|$ binary OR
gates and $n+|E_G|-1$ conjunction gates establish the size bound.
This short formula proof concerns the new allowed-weight
interface, not the imported graph-CNF DNNF lower bound.

\paragraph{Geometric and representation statement.}
The integral scheme
\[
 Y_G=\operatorname{Spec}
 \mathbb Z[x_v:v\in V]/
           (x_v^2-x_v, x_ux_v:uv\in E_G)
\]
is the disjoint union of one copy of $\operatorname{Spec}\mathbb Z$
for every independent set.  Put $B=\mathbb Z[z_{v,b}]$.
On the finite locally free $B$-algebra
$B\otimes_{\mathbb Z}\Gamma(Y_G,\mathcal O)$, multiplication
by $\prod_v((1-x_v)z_{v,0}+x_vz_{v,1})$ has trace $J_G(z)$.
This follows directly from its values in the orthogonal
Boolean-idempotent basis.  The trace and arbitrary base-change
interface is EGA~IV4, 18.2.1--18.2.2, printed p.~111, including
equation~18.2.1.2.  The foregoing circuit bound is a uniform
bound for this regular trace function in the compact graph
presentation; it does not assume that the full idempotent
basis is part of the input.

Equivalently, let $T=(\mathbb G_m)^{2n}$ over $\mathbb Z$,
with the specified rail characters.  The actual free
$T$-module
\[
 \mathcal W_G=
 \bigoplus_{I\,\text{independent}}
       \left(\bigotimes_{v\in I}L_{v,1}\right)
          \otimes
       \left(\bigotimes_{v\notin I}L_{v,0}\right)
\]
has character $J_G$.  In the model with these one-dimensional
generators, direct sums, tensor products, nonnegative integer
multiplicities and shared subexpressions, constructing its
class requires $2^{\Omega(n)}$ positive arithmetic gates.
The integral representation Grothendieck ring admits a
$2^{o(n)}$ expression for that same effective class using
virtual differences, by E and NN-C034.  The classical
dictionary is SGA~3, Expos\'e~I, 4.7.2--4.7.3, for torus
gradings, and SGA~6, Expos\'e~IV, (1.5.2) and Definition~2.2,
for the sum, tensor product and Grothendieck-ring operations.
No arbitrary representation functor, entire finite pushforward,
or coefficient-extraction oracle is counted as a unit-cost gate.

For Boolean unary availability data $a$, discard from $Y_G$
the independent-set sections using a forbidden rail.
This is an open-and-closed integral subscheme, whose nonemptiness
is exactly equation~\eqref{eq:S019-independent-support}.
The linear monotone formula thus decides whether an effective
summand is available, while the full character has the proved
exponential positive arithmetic complexity.  Here Boolean
evaluation is the support map on the positive polynomial
semiring; setting a character weight to zero is not asserted
to define a homomorphism of the torus representation ring.
The result distinguishes these explicit arithmetic, Boolean
and representation models.  It gives neither a polynomial
integral circuit nor a lower bound for unrestricted Boolean
circuits, and uses neither A nor M.
